% 本文件由 LaTeX 排版 Agent 根据有效 translation.json 自动生成。
% author=John Chrysostom; work_id=2101; title=宗徒大事录讲道集

\chapter{宗徒大事录讲道集}

% structure: p0001 source=h1 target=omitted reason=page-title-already-rendered-by-parent

讲道一\newline{}讲道二\newline{}讲道三\newline{}讲道四\newline{}讲道五\newline{}讲道六\newline{}讲道七\newline{}讲道八\newline{}讲道九\newline{}讲道十\newline{}讲道十一\newline{}讲道十二\newline{}讲道十三\newline{}讲道十四\newline{}讲道十五\newline{}讲道十六\newline{}讲道十七\newline{}讲道十八\newline{}讲道十九\newline{}讲道二十\newline{}讲道二十一\newline{}讲道二十二\newline{}讲道二十三\newline{}讲道二十四\newline{}讲道二十五\newline{}讲道二十六\newline{}讲道二十七\newline{}讲道二十八\newline{}讲道二十九\newline{}讲道三十\newline{}讲道三十一\newline{}讲道三十二\newline{}讲道三十三\newline{}讲道三十四\newline{}讲道三十五\newline{}讲道三十六\newline{}讲道三十七\newline{}讲道三十八\newline{}讲道三十九\newline{}讲道四十\newline{}讲道四十一\newline{}讲道四十二\newline{}讲道四十三\newline{}讲道四十四\newline{}讲道四十五\newline{}讲道四十六\newline{}讲道四十七\newline{}讲道四十八\newline{}讲道四十九\newline{}讲道五十\newline{}讲道五十一\newline{}讲道五十二\newline{}讲道五十三\newline{}讲道五十四\newline{}讲道五十五\par

\SourceRecord{Translated by J. Walker, J. Sheppard and H. Browne, and revised by George B. Stevens.；Nicene and Post-Nicene Fathers, First Series；Vol. 11.；Edited by Philip Schaff.；Buffalo, NY: Christian Literature Publishing Co.,；1889.；<http://www.newadvent.org/fathers/2101.htm>.}

% structure: p0001 source=h1 target=section reason=page-title

\section{宗徒大事录讲道词 1}

\BibleRef{宗 1:1-2}\par

\Quote{\Person{德敖斐罗}啊，我在前部书中，论及\Person{耶稣}所行所教的一切，从开始直到祂藉圣神向所选宗徒颁下命令后，被接升天的那一天为止。}\par

对许多人来说，本书及其作者都如此鲜为人知，甚至不知道有这样一本书存在。为此我特地选取这段叙述作为我的讲题，为将此书介绍给不认识它的人，并让如此珍宝不再隐藏不见。因为本书确实对我们有益，不亚于福音书；它充满基督信仰智慧和纯正教义，特别是在有关圣神的论述上。因此，让我们不要轻率地略过它，而要仔细研读。如此，基督在福音中所说的预言，我们在此可见其实际应验；并在事实中看到照耀其中的真理的明证，以及圣神降临后门徒身上发生的巨大变化。例如，他们听到基督说：\BibleQuote{凡信我的，我所做的工程，他也要做，并且还要做比这些更大的工程}\BibleRef{若 14:12}；又当祂预言门徒将要在官长和君王面前受审，在会堂中被鞭打，遭受苦害，却要战胜一切\BibleRef{玛 10:18}，以及福音要传遍天下\BibleRef{玛 24:14}——所有这些，按所说的话如何精确应验，均可在此书中看见，还有更多祂与他们同在时所告诉他们的。在此你还要看到宗徒们自己，仿佛插上翅膀飞越陆地海洋；那些曾经怯懦、愚钝的人，突然之间变得完全不同；他们蔑视财富，超越荣誉、情欲和贪念，总之超越一切此类情感；此外，他们之间现今何等的同心合意；再也没有从前的嫉妒，再也不贪图首位，而是所有德行在他们身上臻于圆满，并由那极显著的——主曾多次嘱咐的——爱德闪耀出来，祂说：\BibleQuote{如果你们之间彼此相亲相爱，世人因此就可认出你们是我的门徒}\BibleRef{若 13:35}。再者，在此还可找到一些教义，若非本书存在，我们无法如此确定地知晓，甚至我们得救的顶点也会被隐藏，无论对于生活实践还是教义而言。\par

然而，本书大部分内容记述\Person{保禄}的事迹，他\BibleQuote{比众人更劳苦}\BibleRef{格前 15:10}。原因是，本书的作者，即蒙福的\Person{路加}，曾是他的同伴；此人的高尚品质，已在许多其他事例中充分显明，尤其在他对导师的坚定依从上，他始终追随。当众人都离弃了\Person{保禄}，有的去了\Place{迦拉达}，有的去了\Place{达尔玛提雅}时，听\Person{保禄}论及这位门徒所说的话：\BibleQuote{只有\Person{路加}与我同在}\BibleRef{弟后 4:10}。他又在致格林多人书函中关于\Person{路加}写道：\BibleQuote{他的颂扬因福音传遍众教会}\BibleRef{格后 8:18}。同样，当他说：\BibleQuote{祂显现给\Person{刻法}，然后显现给那十二位}\BibleRef{格前 15:5}，以及\BibleQuote{按照你们所接受的福音}\BibleRef{格前 15:1}，他指的是\BibleBook{路加}{福音}。那么，为什么他既然与\Person{保禄}在一起直到末了，却没有叙述一切呢？我们可以回答：这里所记的，对那些愿意聆听的人已经足够；圣史的写作总是针对当时最重要的内容，并非以著作书籍为目的；事实上，有许多事情他们是通过口传传授的。虽然此书所载的一切都值得赞叹，但尤其令人敬佩的是宗徒们俯就听众需要的方式；这是圣神所促成的屈尊，祂如此安排，使他们主要论述的主题是关于基督的人性。因为虽然他们谈论基督很多，但对祂的天主性却谈论甚少；他们主要讲论的是人性、受难、复活和升天。因为首先需要的是让人相信祂已经复活并升天。正如基督本人最坚持的是让人知道祂来自圣父，同样，这位作者的主要目的也是宣告基督已从死者中复活、被接升天，并往天主那里去，且来自天主。因为如果祂来自天主这一事实没有首先被相信，那么再加上复活和升天，犹太人会觉得整个教义难以置信。因此，他温和而逐步地将他们引向更高的真理。不仅如此，在\Place{雅典}，\Person{保禄}甚至只是简单地称呼祂为人，没有多说\BibleRef{宗 17:31}。\par

然而，何必提那些犹太人，就连门徒们，每当听到更崇高的道理时，也常常感到困扰和反感。因此，祂也曾对他们说：\Quote{我还有许多事要告诉你们，但你们现在不能承受。}\BibleRef{若 16:12}如果那些与祂同处已久、蒙受许多奥秘、目睹许多奇事的人尚且不能承受，那么，那些刚刚从祭台、偶像、祭献、猫和鳄鱼（因为外邦人崇拜这些东西）以及他们其余邪恶的道路中被拉出来的人，怎能立刻接受教义中更崇高的事理呢？尤其那些犹太人，他们每一天都听着、耳中常常响着：\Quote{上主你的天主是唯一的上主，除祂以外没有别的神}\BibleRef{申 6:4}；他们也亲眼看见祂被钉在十字架上，甚至亲手将祂钉死并埋葬，也没有看见祂复活；当有人告诉他们这同一位就是天主，与圣父平等时，他们怎能不感到震惊和反感呢？因此，他们温和地、逐步地带领他们，以极大的体恤和宽容，俯就他们低下的程度，而他们自己却更丰盛地享受圣神的恩宠，并奉基督的名行了比基督本人更大的事，以便能立刻将他们从卑下的理解中提升起来，并证实基督从死者中复活了。事实上，这正是本书的内容：对复活的证明；一旦这被相信，其余的自会相继而来。那么，本书的主题和整个范围，主要就是我所讲的。现在让我们来听序言本身。\par

\Quote{德敖斐罗！我在第一部书中，已论及耶稣所行所教的一切。}\BibleRef{宗 1:1}他为什么提醒他想起福音呢？为要表明他是如何严格地可信赖的。因为在第一部作品的开头，他说：\Quote{我也认为，从起头详确地考察了一切，按次第给你写下来是很好的。}\BibleRef{路 1:3}他并不满足于自己的见证，而将整个事情引证给宗徒们，说：\Quote{正如那些从开始就亲眼见过，并为圣言服役的人所传授给我们的。}\BibleRef{路 1:2}既然他在前一部作品中已使他的记述可信，那么在这部著作中他就无需再提出他的凭证，因为他的弟子已经一次完全满意了，而通过提及那部前作，他提醒了他要严格信赖他的真实性。因为如果一个人已经表明自己有能力并值得信赖地记述他所听来的事，并且获得了我们的信任，那么当他编写一部不是从别人那里接受来的，而是他亲眼所见、亲耳所闻的记载时，他就更应当得到我们的信任。因为你接受了关于基督的事，那么你更应当接受关于宗徒的事。\par

那么，有人可能会问：这仅仅是历史问题吗？圣神与此无关吗？并非如此。因为如果\Quote{那些从开始就亲眼见过，并为圣言服役的人所传授给我们的}，那么他所讲的，就是\Quote{他们的}话。他为什么不说：\Quote{正如那些被堪当领受圣神的人所传授给我们的}，而说\Quote{那些亲眼见过的人}呢？因为在信仰之事上，恰恰是使人有权被相信的，就是由亲眼所见的人那里学来的；而另一种说法，在愚昧人看来不过是卖弄和虚饰而已。因此，若望也这样说：\Quote{我看见了，并且作证：这就是天主子。}\BibleRef{若 1:34}基督也向迟钝的尼苛德摩同样表达：\Quote{我们讲论我们所知道的，我们见证我们所见过的，但你们却不接受我们的见证。}\BibleRef{若 3:11}因此，祂准许他们在许多细节上凭亲眼所见来作证，当祂说：\Quote{你们也要为我作证，因为你们从起初就与我在一起。}\BibleRef{若 15:27}宗徒们自己后来也常常这样说：\Quote{我们就是这些事的证人，并且圣神，就是天主赐给那些服从祂的人，也是证人。}\BibleRef{宗 2:32}后来，伯多禄为坚定复活的事，又说：\Quote{我们曾与祂一同吃喝。}\BibleRef{宗 10:41}因为人们更容易接受与祂同行之人的见证，因为当时关于圣神的概念对他们来说还很遥远。因此，若望那时在他的福音中论及血和水时，说他\Emph{亲眼看见}了，对他们而言，他将亲眼所见的事实等同于最高的见证，虽然圣神的见证比亲眼所见更为确实，但对于不信者并非如此。路加是圣神的参与者，这是十分清楚的，既从现今仍发生的奇迹，也从当时连普通人也被赐予圣神的事实，又来自保禄的见证，他说：\Quote{他的赞美在福音中}\BibleRef{格后 8:18}，以及他所蒙受的任命：因为说了这话后，宗徒又加上：\Quote{不但如此，他也被众教会选派，与我们同行，为这项由我们管理的恩宠服务。}\par

现在请注意他是何等谦逊。他没有说\InnerQuote{我以前所传的福音}，却说\Quote{我作了前部书}；认为福音的称号对他过于崇高；尽管正是为此，宗徒才赞扬他，说道：\Quote{他的称赞是在福音里。}但他自己谦逊地说：\Quote{我作了前部书——\Person{德敖斐罗}啊，论及耶稣所开始行和教的一切}——不是简单地说\Quote{一切}，而是从开始到结束；他说：\Quote{直到祂被接升天的那一天。}然而若望说，不可能把所有的事都写下：因为他说：\Quote{我想，所写的书，就是世界也容不下了。}\BibleRef{若 21:25} 那么，福音作者在这里怎么说\Quote{一切}呢？他不是说\Quote{一切}，而是说\Quote{论及一切}，意思是\Quote{以概括的方式，笼统地}，以及\Quote{论及主要而紧急的一切}。然后他告诉我们他在什么意义上说\Emph{一切}，当他加上：\Quote{耶稣所开始行和教的一切}；意指祂的奇迹和教训；不仅如此，还暗示祂的行为也是一种教训。\par

但现在请看作者那充满慈爱和宗徒精神的情感：他为了一位个人的缘故，不辞辛劳地为他写了整部福音。他说：\Quote{为使你认清你所受的教导是确实的。}\BibleRef{路 1:4} 事实上，他曾听见基督说：\Quote{我父的旨意是要这些小子中一个也不丧失。}\BibleRef{玛 18:14} 为什么他不把这些写成一部书，寄给一个人\Person{德敖斐罗}，却分为两部分呢？为了清楚，并让兄弟有休息的间歇。此外，这两部书在主题上是不同的。\par

但请思考基督如何借着行动来证实祂的话。于是祂说：\Quote{你们跟我学吧，因为我是良善心谦的。}\BibleRef{玛 11:29} 祂教导人要贫穷，并用行动表明：\Quote{人子没有枕头的地方。}\BibleRef{玛 8:20} 同样，祂命令人爱仇敌，并在十字架上教导同样的功课，祂为那些钉死祂的人祈祷。祂说：\Quote{如果有人想控告你，拿你的内衣，你就连外衣也给他。}\BibleRef{玛 5:40} 而\Emph{祂}不仅给了自己的衣服，还流了自己的血。祂就这样吩咐别人教导。因此保禄也说：\Quote{你们仿效我们吧。}\BibleRef{斐 3:17} 没有什么比一位教师只嘴上显哲学更冷漠了：这不是教师的角色，而是伪君子的角色。因此，宗徒们先以行为教导，然后以言语；更确切地说，当他们的行为如此响亮时，他们不需要言语。说基督的受难是行动也不算错，因为在受苦中，祂完成了那伟大而奇妙的行动，借着它摧毁了死亡，并为我们成就了其他一切。\par

\Quote{直到祂藉圣神嘱咐了所选的宗徒之后，被接上升的那一天。祂藉圣神赐下诫命之后}\BibleRef{宗 1:2}；即祂对他们所说的是属神的话，没有任何人性的；要么是这个意思，要么是祂藉圣神赐给他们诫命。你看到他以何等卑微的言辞谈论基督，正如基督自己谈论自己一样？\Quote{我若靠天主的圣神驱魔}\BibleRef{玛 12:28}；因为圣神确实在那个圣殿中运行。那么祂命令了什么？祂说：\Quote{所以你们要去，使万民成为门徒，因父及子及圣神之名给他们授洗，教训他们遵守我所吩咐你们的一切。}\BibleRef{玛 28:19-20}这是对宗徒们的高度赞扬：将这样的重任托付给他们，即世界的救恩！充满圣神的话语！作者在\Quote{藉圣神}这一表达中暗示了这一点（主说：\Quote{我对你们所说的话，就是圣神。}\BibleRef{若 6:63}）；从而引导听者渴望知道这些命令是什么，并建立宗徒的权威，因为他们即将说出圣神的话和基督的诫命。他说：\Quote{祂赐下诫命之后，被接升天。}他没有说\Quote{升天}，而是仍像谈论一个人那样说。看来祂在复活后也教导了门徒，但关于这段时间，没有人向我们详细叙述全部。圣若望确实和本书作者一样，比其他人都更详细地谈论这个话题，但没有人清楚地叙述一切（因为他们急于转向其他事）；然而，我们通过这些事从宗徒那里学到了，因为他们听到了什么，就告诉了。祂也向他们显示自己活着。他先讲了升天，然后提到复活；因为你已被告知\Quote{祂被接升天}，所以怕你猜测祂是被他人接上去的，他就加上：\Quote{祂向他们显示自己活着。}因为如果祂在更大的事上显示了自己，当然也在较小的事上显示。你看到他是如何随意而不觉察地沿途播下这些伟大教义的种子吗？\par

\BibleQuote{四十天之久向他们显现。}祂如今不再像复活以前那样常与他们在一起。因为作者不是简单地说\Quote{四十天}，而是说\Quote{四十天之久}。祂往来显现，又隐没；借此引导他们达到更高的认识，不再让他们以从前的方式对待祂，而是采取有效措施，既使祂复活的事实得以相信，又使祂自己被看为超越凡人的。同时，这两件事是相反的；为了相信祂的复活，需要许多显于人性的作为，而为另一目的，则恰恰相反。然而，两者都达到了效果，各在其适宜的时候。\par

那么，为什么祂不向众人显现，只向宗徒们显现呢？因为对大多数人来说，这恐怕只像是幻影，因为他们不明白这奥迹的奥秘。因为如果门徒们自己起初尚且怀疑不安，需要亲手触摸和祂同食的证据，那么大众多半会如何呢？因此，祂藉着（宗徒所行的）奇迹，使祂复活的证据确凿无疑，以致不仅当时的人——这是亲眼见证的结果——而且以后的所有人都能确信祂复活的事实。基于此，我们也与不信者辩论。如果祂没有复活，而是仍旧死亡，那么宗徒们如何因祂的名行奇迹呢？可是你说，他们没有行奇迹？那么我们的宗教(\OriginalTerm{民族}{\GreekText{ἔθνος}})是如何建立的呢？对于这一点，他们肯定无法辩驳或否认我们亲眼所见的事实。因此，当他们说没有发生奇迹时，反而给自己造成了更重的伤害。因为最大的奇迹莫过于：没有任何奇迹，全世界却热切地投入十二个贫穷文盲的网中。因为这些渔夫取胜，并非靠金钱财富，也非靠言辞智慧，更非任何其他此类事物；因此反对者即使不情愿，也必须承认这些人身上有神圣的力量，因为人的力量绝不可能成就如此伟大的功绩。为此，祂当时在地上停留了四十天，用这段时间提供他们亲眼见到祂本人的可靠证据，免得他们以为所见的是幻影。不仅如此，祂还增添了与他们在桌前同食的证据：为了表明这点，作者补充说：\BibleQuote{祂与他们一同吃饭时，祂命令他们。}\BibleRef{宗 1:4} 宗徒们自己也常提出这一情境作为复活的确证；比如他们所说：\BibleQuote{凡与祂同饮共食的人。}\BibleRef{宗 10:41}\par

而当祂在那四十天里向他们显现时，祂做了什么呢？是的，作者说，祂同他们交谈，\BibleQuote{讲论天主国的事。}\BibleRef{宗 1:3} 因为门徒们对已发生的事既悲伤又困惑，又即将出去面对巨大的困难，祂用关于未来的讲论使他们恢复勇气。\BibleQuote{祂吩咐他们不要离开耶路撒冷，但要等候父的恩许。}\BibleRef{宗 1:4} 首先，祂带他们到加里肋亚，他们又恐惧又颤抖，以便他们能安全地听祂的话。后来，当他们听了，并与祂一起过了四十天，\BibleQuote{祂命令他们不要离开耶路撒冷。}为什么呢？正如当士兵要冲入人群时，没有人会让他们在没有武装之前就出击，又如马匹在没有御者之前不会被允许从栏中出发；同样，基督不让这些人在圣神降临之前出现在战场上，以免他们容易被众人击败和俘虏。这还不是唯一的原因，还有在耶路撒冷有许多人将要相信。再者，以免有人说，他们离开自己的熟人，去陌生人面前炫耀；因此，在那些把基督置于死地的人中间，在那些钉死祂和埋葬祂的人中间，在犯下这邪恶行为的同一城中，他们展示祂复活的证据；从而堵住所有外来反对者的口。因为当那些钉死祂的人甚至成为信徒时，这显然既证明了被钉的事实和行为的邪恶，又提供了复活的有力证据。此外，免得宗徒们说，我们怎么可能在邪恶血腥的人中间生存，他们人数众多，我们却如此稀少卑微，请看祂如何消除他们的恐惧和忧虑，用这些话：\BibleQuote{但要等候父的恩许，就是你们听我说过的。}\BibleRef{宗 1:4} 你要说，他们什么时候听过这话？就是当祂说：\BibleQuote{我去为你们有益，因为我若不去，护慰者便不会到你们这里来。}\BibleRef{若 16:7} 又说：\BibleQuote{我要请求父，祂必会赐给你们另一位护慰者，使祂永远与你们同在。}\BibleRef{若 14:16}\par

但为什么圣神降临在他们身上，不是在基督还在的时候，也不是在他离开之后立即，而是基督在第四十日升天，而圣神在\BibleQuote{五旬节日，即第五十日，到了}的时候降临？\BibleRef{宗 2:1} 而且，如果圣神还没有来，祂怎么说：\BibleQuote{你们领受圣神吧？}\BibleRef{若 20:22} 为了使他们配得并适于接受祂。因为如果达尼尔在看见天使时就昏倒了\BibleRef{达 8:17}，那么这些将要领受如此浩大恩宠的人，岂不更甚？或者可以说，基督是将来之事当作已成之事来说的，正如祂说：\BibleQuote{你们要践踏蛇和蝎子，并胜过魔鬼的一切能力。}\BibleRef{路 10:19} 但为什么圣神还没有来呢？应当先让他们渴望这件事，然后才领受恩宠。为此，基督亲自离开，然后圣神降临。因为如果祂还在那里，他们就不会像后来那样热切地期待圣神。因此，圣神也不是在基督升天之后立即降临，而是过了八或九天。我们也一样；当我们感到需要时，我们对天主的渴望最为强烈。因此，若翰选择在他自己被囚禁的时候，派遣他的门徒到基督那里去，那时他们更感到需要耶稣。此外，应当先让我们的本性在天上被看见，和好得以完成，然后圣神降临，喜乐才会纯全无杂。因为，如果圣神已经降临，基督然后离开，圣神留下，那么安慰就不会像现在这样大。事实上，他们紧紧抓住祂，不忍与祂分离；所以为了安慰他们，祂说：\BibleQuote{我去为你们有益。}\BibleRef{若 16:7} 为此，祂也等候了中间的那些日子，好让他们先沮丧一会儿，如我所说，感到需要祂，然后获得完满而纯全的喜乐。但如果圣神次于圣子，那么安慰就不够充分；祂又怎能说：\BibleQuote{我去为你们有益？} 因此，教义中更重大的事保留给了圣神，免得门徒以为祂是次等的。\par

也要想想祂如何使他们必须留在耶路撒冷，应许赐给他们圣神。以免他们在祂升天后再次逃散，祂用这个期待像绳索一样将他们拴在那个地方。但祂说了：\BibleQuote{等候圣父的应许，就是你们从我听过的}，然后祂又加说：\BibleQuote{若翰确实以水施洗，但你们要在不多几日以后，以圣神受洗。}（宗 1:4-5）因为现在祂确实让他们清楚地看到祂与若翰之间的区别，而不是像先前那样隐晦地暗示；实际上，祂以前说得非常隐晦，当祂说：\BibleQuote{然而，在天国里最小的比他还大}，但现在祂清楚地说：\BibleQuote{若翰以水施洗，但你们要以圣神受洗。}\BibleRef{玛 11:11} 而且祂不再使用见证，只是提及若翰这个人，提醒门徒祂所说的话，并向他们表明他们现在已比若翰更大，因为他们也要以圣神施洗。再者，祂没有说：我以圣神给你们施洗，而是说：\Quote{你们要以圣神受洗}，以此教训我们谦卑。因为若翰的见证已经足够清楚，是基督自己要施洗：\BibleQuote{祂要以圣神和火给你们施洗}\BibleRef{路 3:16}；因此祂也提及了若翰。\par

因此，\WorkTitle{福音书}是基督所言所行的历史；而\WorkTitle{宗徒大事录}则是那位\Quote{另一位护慰者}所言所行的历史。并非圣神在\WorkTitle{福音书}中没有行许多事；正如基督在此\WorkTitle{宗徒大事录}中仍然在人内工作，如同祂在\WorkTitle{福音书}中一样：只是那时圣神通过圣殿工作，现在通过宗徒；那时，祂进入童贞玛利亚的子宫，塑造了圣殿；现在，进入宗徒的灵魂；那时，以鸽子的形象；现在，以火的样子。为什么？显示主在那里的温和，但在这里也显示祂的惩罚，现在祂也提醒他们审判。因为当需要宽恕时，需要许多温和；但现在我们获得了恩赐，从此是审判和考量的时期。\par

但爲何基督說：\Quote{你們將要受洗}，而事實上樓上並沒有水？因爲聖洗聖事中更本質的部分是聖神，水正是藉着祂而發揮作用；同樣，我們的主也被稱爲受傅者，並非因爲祂曾被油傅抹，而是因爲祂領受了聖神。此外，我們確實發現他們先後領受了水的洗禮和聖神的洗禮，且在不同時刻。在我們身上，兩者發生在同一行動中，但那時它們是分開的。起初他們由若翰施洗；因爲如果娼妓和稅吏都去受那洗禮，那麼那些日後要受聖神洗禮的人更應如此。爲使宗徒們不致說，他們總是在應許中領受它\BibleRef{若 14:15}（因爲基督已多次論及聖神，免得他們以爲那是\OriginalTerm{非位格的能力或運作}{\GreekText{ἐνέργειαν} \GreekText{ἀνυπόστατον}}；爲使他們不致如此說，祂便加上：\Quote{不多幾日以後}。祂沒有說明具體時間，好使他們時常警醒；但祂告訴他們這事將很快發生，免得他們喪氣；然而祂沒有加上確切時間，好使他們常保持警醒。祂不僅藉此安慰他們，即時間的短暫，更藉着說：\Quote{你們從我所聽見的應許}。因爲祂說：這不是我唯一一次告訴你們，我早已應許了，我必將成就。那麼，當這個如此臨近的日子祂都不願揭示時，祂不指明末日的時間又有何奇怪呢？這是有充分理由的：爲使他們常常警醒，處於期待和專注的狀態。\par

因爲人若不警醒，決不能，決不能享受恩寵的益處。你豈不見厄里亞對其門徒所說的話？\Quote{你若見我被接去}\BibleRef{列下 2:10}，你所求的必爲你成就。基督也常對那些來到祂面前的人說：\Quote{你信嗎？}因爲如果我們不將自己歸屬於所賜之恩，我們也不會深切感受其益處。保祿的情形也是如此：恩寵並非立刻臨到他，而是有三天的間隔，期間他眼瞎，同時被淨化，並因恐懼而預備。正如那些染紫色的人在布匹上預先浸染其他材料，使色澤持久；同樣，天主先設法使靈魂徹底認真，然後傾注祂的恩寵。爲此，祂也沒有立刻派遣聖神，而是在第五十天。若有人問，爲何我們不在五旬節那季節施洗？我們可以回答：恩寵今日與那時相同；但藉着齋戒準備，心靈如今更爲提升。五旬節的季節也提供了一個合理的理由。那是什麼呢？我們的教父認爲聖洗聖事正是對邪惡貪慾的適當約束，並在歡樂時節教導節制的有力功課。\par

因此，好像我們正與基督一同宴飲，共享祂的筵席，讓我們不要行事隨意，而要度過齋戒、祈禱和極度節制的生活。因爲如果一個將要擔任某種現世職務的人，終其一生都準備自己，爲獲得某種尊嚴而花費錢財、耗費時間、忍受無盡勞苦；那麼我們這些如此漫不經心地接近天國，既在領受前不認真，領受後又疏忽的人，該當受何懲罰呢？不，這正是我們領受後疏忽的原因：我們在領受前未曾警醒。因此，許多人在領受後，立刻回到從前的嘔吐，變得更爲邪惡，爲自己招致更嚴厲的懲罰；當他們已從先前的罪惡中被解救出來時，卻更嚴重地激怒了審判者，因爲他們雖從如此重病中得救，卻仍未學會節制，反而應驗了基督對癱子所說的威脅：\Quote{看，你已痊愈了；不要再犯罪，免得遭遇更糟的事}\BibleRef{若 5:14}；以及祂對猶太人所預言的：\Quote{最後的狀況比先前更壞}\BibleRef{瑪 12:45}。因爲祂說，\Emph{如果}我沒有來，也沒有對他們說話，他們就沒有罪\BibleRef{若 15:22}，以此表明他們的忘恩負義將爲他們招致最壞的災禍；因此，在這些恩惠之後所犯之罪的罪責是雙倍、四倍的，因爲在我們受尊榮之後，我們卻表現得忘恩負義和邪惡。聖洗之池也絲毫不能幫助我們獲得較輕的懲罰。試想：一個人因謀殺、姦淫或其他罪行而犯了大罪；\Emph{這些}罪都藉着聖洗被赦免了。因爲沒有任何罪，沒有任何不敬，不屈服於這恩賜；因爲恩寵是神聖的。這人再次犯姦淫和謀殺；先前的姦淫確實被除去了，謀殺被赦免了，不再歸到他身上，\Quote{因爲天主的恩賜和召叫是決不會撤回的}\BibleRef{羅 11:29}；但對於聖洗之後所犯的罪，他將受與先前罪過一同被追究時同樣大的懲罰，甚至更重。因爲罪責不再單純相等，而是加倍和三倍。請看聖保祿的話，證明這些罪的刑罰更大：\Quote{人廢棄梅瑟法律，憑兩三個見證人，尚且不得憐恤而死；何況人踐踏了天主子，將那藉以成聖的盟約之血視爲不潔，又褻慢施恩的聖神，你們想，這人該受怎樣更重的懲罰呢？}\BibleRef{希 10:28}\par

或许我们现在已阻止了许多人领受圣洗圣事。但我们这样说并非出于此意，而是希望他们在领受后，能持守节制与极大的温和。有人说：\Quote{但我害怕。}你若害怕，就该领受并守护它。他说：\Quote{不，} \Quote{但这正是我不领受的原因——因为我害怕。}难道你不害怕就这样离去吗？他说：\Quote{天主是仁慈的。}那么，领受圣洗罢，因为祂是仁慈的，且乐于助人。但你，在需要认真对待的事上，却不提祂的仁慈；你只在乐意时才想到这一点。然而，那时才是该依靠天主仁慈的时候，当我们尽了自己的本分，我们才最确信能获得它。因为那将一切交托给天主，并在领洗后犯罪（作为人，这是可能的）而悔改的人，必获得仁慈；而那玩弄天主的仁慈、未分享那恩宠就离世的人，必受惩罚，无话可说。你说：\Quote{但若他领受了恩宠后离去呢？}他将空无一切善工而去。因为依我之见，不可能——是的，不可能——那因如此希望而轻慢圣洗的人能成就任何慷慨而良善的事。你为什么怀着这样的恐惧，寄望于未来的不确定？为何不将这恐惧转化为劳苦与热忱，而你将成为伟大而可钦佩的？哪个更好：恐惧还是劳苦？设想有人让你无所事事地待在一所摇摇欲坠的房子里，说：\Quote{等着腐朽的屋顶落在你头上：因为它也许会落下，也许不会；但若你宁愿它不落，那就去劳苦，住进更安全的房间：你宁可选择那伴随着恐惧的闲散状态，还是这带有信心的劳苦？}那么，现在就照样行事罢。因为不确定的未来就像一所腐朽的房子，随时可能倒塌；但这劳苦，尽管辛苦，却确保安全。\par

但愿我们不至于陷入如此大的困境，以致在领洗后犯罪。然而，即使发生这样的事，天主是仁慈的，祂赐给我们许多甚至在领洗后获得赦免的方法。但正如那些在领洗后犯罪的人比慕道者受更严厉的惩罚，同样，那些知道在悔改中有药方却不愿使用的人，将受更重的责罚。因为天主的仁慈越大，若我们不善用那仁慈，惩罚也越重。人哪，你说什么？当你充满如此严重的邪恶、已被定案时，你突然成了朋友，被提升至最高的尊荣，不是靠你自己的劳苦，而是靠天主的恩赐；你再次回到先前的不良行为中；尽管你理应受重罚，天主却没有转面不顾，反而赐下无数救恩的机会，使你仍可成为朋友：然而你仍不愿劳苦。从此以后，你还能得到什么宽恕？外邦人岂不有理由嘲笑你为无用的懒汉？他们说：\Quote{若你那教义有能力，那么这众多未入门的人是怎么回事？}若奥迹是卓越而可欲的，就让人不要在临终时才领受圣洗。因为那不是举行奥迹的时候，而是立遗嘱的时候；举行奥迹的时候是在心智健康、灵魂健全之时。因为，若一个人不愿在那种状态下立遗嘱，若他真那样做了，就会为日后的诉讼提供把柄（这就是立遗嘱者先写下这些话的原因：\Quote{活着，神智清醒，且健康，我如此处置我的财产：}），那么，一个不再能自主的人，怎能正确地进行神圣奥迹的预备呢？因为，若在世务上，世俗的法律不允许一个心智不完全健全的人立遗嘱，尽管他是在为自己的事务立法；那么，当你接受关于天国和那世界不可言喻的财富的教导时，若你很可能因疾病的剧痛而神志不清，又怎能清楚地学习一切呢？你何时将对基督说出那些话，在即将与此世离别、与祂同葬之时？因为我们必须用行动和言语向祂表明我们的善意。\BibleRef{罗 6:4}你现在所做的，就像一个人想在战争即将结束时入伍，或想在观众已离席时才脱衣进入竞技场。因为你被授予武器，不是要你立即离去，而是要你装备好，向敌人举起胜利的旗帜。愿没有人认为现在讨论这个话题不合时宜，因为现在不是四旬期。不，这正是使我烦恼的，因为你们在这种事上讲究特定时间。而那个太监，虽是野蛮人且在路上，甚至在公路上，也不寻求特定时间\BibleRef{宗 8:27}；狱警也是如此，虽然他身处一群囚犯之中，他所见的导师是受过鞭打、带着锁链的人，且他仍要监管他\BibleRef{宗 16:29}。但在这里，既不是在监狱里，也不是在路上，许多人却将圣洗推迟到最后一口气。\par

现在，如果你仍然质疑基督是天主，就离开教会：不要留在这里，甚至不要作为圣言的听取者和一个慕道者；但如果你确信这一点，并清楚知道这一真理，为什么迟延？为什么退缩和犹豫？你说，因为害怕犯罪。但你不害怕更糟的事吗？即带着如此沉重的负担离开现世？因为两者并不同等可原谅：没有领受摆在你面前的恩宠，与未能努力正直生活。如果有人质问你：你为什么不去领受它？你将如何回答？在另一种情况下，你可以以情欲的重担和德行的困难为托辞；但在这里却毫无此类借口。因为这里是恩宠，自由地赐予解放。但你害怕犯罪？让这话在你领受圣洗圣事后再说：那时再怀有这种恐惧，以便持守你已经获得的自由；而不是现在，以免阻碍你领受如此恩赐。然而，你现在在圣洗前谨慎，圣洗后却疏忽。但你在等待四旬期：为什么？那个时期有什么优势？不，宗徒们并不是在逾越节领受恩宠，而是在另一个时期；然后有三千（路加说）和五千人受了洗（\BibleRef{宗 2:41}；\BibleRef{宗 4:4}，及\BibleRef{宗 10}）；还有科尔乃略。因此，让我们不要等待固定的时间，免得因犹豫和拖延而空手离去，失去如此伟大的恩赐。你以为当我听说有人未经圣洗就被带走，同时我思考那生命不可忍受的惩罚、无情的审判时，我的痛苦如何？再者，我多么悲伤地看到其他人临终喘气，却连那时都未清醒。由此也发生与这个恩赐极不相称的情景。因为本应欢乐、跳舞、欢跃、戴花环，当另一个人受洗时；病人的妻子一听到医生嘱咐这事，就悲痛欲绝，仿佛遭遇了什么可怕的灾祸；她发出最大的哀号，整个屋子只听见哭号和哀恸，就像被判刑的罪犯被带往刑场一样。病人自己那时更加伤心；如果他病愈了，也如同遭受了什么大害一样烦恼。因为他没有为德行生活做好准备，所以对随后要进行的战斗毫无心思，一想到就退缩。你看见魔鬼的诡计、羞耻和嘲笑吗？让我们摆脱这种耻辱；让我们按基督所吩咐的生活。祂赐给我们圣洗圣事，不是让我们领受后就离开，而是要在以后的生活中结出果实。对那即将离去、颓丧的人，怎能说：结出果实？难道你没有听说：\BibleQuote{圣神的果实是爱、喜乐、平安？} \BibleRef{迦 5:22} 那么，为什么这里却恰恰相反？因为妻子站在那里悲恸，她本应喜乐；孩子们哭泣，他们本应一同欢喜；病人自己躺在黑暗里，被噪音和喧闹包围，他本应大庆祝；因想到留下子女为孤儿、妻子为寡妇、房屋荒凉而极度沮丧。这是接近奥迹的状态吗？回答我；这是接近祭台的状态吗？这种事能容忍吗？如果皇帝发送信件释放监牢里的囚犯，就会有喜乐和高兴：天主从天上派遣圣神，为赦免的不仅是金钱的欠债，而是整个罪孽，你们却都哀哭悲痛？哎，这是何等的不合适！更不用说有时水倒在死人身上，神圣奥迹被扔在地上。然而，这不能怪我们，而是那些如此悖逆的人。因此，我劝你们放下一切，立即转向圣洗，满怀热忱，以便我们在这现时表现出极大的热忱后，能为将来的事获得信心；为此我们得以达到，愿我们众人都能因我们的主耶稣基督的恩宠和仁慈而获得，愿光荣和权能归于祂，直到永远。阿们。\par

\SourceRecord{Translated by J. Walker, J. Sheppard and H. Browne, and revised by George B. Stevens.；Nicene and Post-Nicene Fathers, First Series；Vol. 11.；Edited by Philip Schaff.；Buffalo, NY: Christian Literature Publishing Co.,；1889.；<http://www.newadvent.org/fathers/210101.htm>.}

% structure: p0001 source=h1 target=section reason=page-title

\section{关于宗徒大事录的第二篇讲道}

\BibleRef{宗 1:6}\par

\BibleQuote{当他们聚在一起时，他们问他说：主，此时你要复兴以色列国吗？}\par

门徒们打算问任何事时，就一同来到祂面前，藉着人数众多迫使祂应允。他们深知，在祂先前所说的\BibleQuote{但那一日，没有人知道}\BibleRef{玛 24:36}中，祂只不过是不愿告诉他们；因此他们再次上前提出这个问题。他们若真满意那回答，就不会提出这个问题了。因为他们听说即将领受圣神，既然现在堪受教导，便渴望知道。他们也已准备好获得自由：因为他们无心面对危险；他们所愿的是重新自由呼吸；因为发生在他们身上的并非轻浮之事，而是极度的危险曾悬在他们头顶。他们未向祂提及任何关于圣神的事，提出了这个问题：\BibleQuote{主，此时你要复兴以色列国吗？} 他们不是问\Quote{何时？}而是问\BibleQuote{此时}。他们是如此渴望那日子。事实上，在我看来，他们对那个国度的性质尚无清晰概念，因为圣神尚未教导他们。他们不说\Quote{这些事何时发生？}却以更大敬意近前说：\BibleQuote{你要此时复兴这国吗？} 似乎它已经倾覆了。因为在那里他们仍然受感官对象影响，因为他们尚未超越前人；而在这里，他们从此对基督怀有崇高的认识。既然他们的心思被提升，主也以更高的方式对他们说话。因为祂不再对他们说\BibleQuote{那日子连人子也不知道}\BibleRef{谷 13:32}；而是说：\BibleQuote{圣父凭自己的权柄所定的时刻、日期，不是你们可以知道的。}\BibleRef{宗 1:7} 祂似乎在说：你们所问的超过了你们的能力。然而，即使现在，他们所学习的事远比这更大。为使你们看到确实如此，请看我要列举多少事。请问，有什么事比他们所学习的事更大呢？他们学到了有天主之子，且天主有位与自己同等尊严的儿子\BibleRef{若 5:17}；他们学到了将有复活\BibleRef{玛 17:9}；祂升天后坐在天主右边\BibleRef{路 22:69}；更奇妙的是，肉身坐在天上，受天使朝拜，且祂将再来\BibleRef{谷 16:19}；他们学到了审判时将要发生的事\BibleRef{玛 16:27}；学到了他们将要坐在那里审判以色列十二支派\BibleRef{路 21:27}；学到了犹太人将被弃绝，而外邦人将代替他们进入\BibleRef{玛 19:28}。请告诉我，哪一样更大？是知道一个人将要为王，还是知道什么时候？\BibleRef{路 21:24} 保禄学到了\BibleQuote{人不可说的事}\BibleRef{格后 12:4}；世界造成以前的事，他全都学到了。哪一样更难？是开始还是终结？显然是知道开始。这事梅瑟学得了，以及时间、多久以前，他还历数了年代。明智的撒罗满说：\Quote{我要提及从世界开始的事。} 而且他们知道时候近了：正如保禄所说：\BibleQuote{主已经近了，应当一无挂虑。}\BibleRef{斐 4:5} 这些事他们那时并不知道，然而祂提到了许多征兆\EditorialNote{玛窦福音24}。但是，正如祂刚刚所说的\BibleQuote{不多几日以后}，希望他们保持警醒，并没有公开宣布确切时刻，这里也是如此。然而，他们现在问祂的不是关于世界的终结，而是说：\BibleQuote{主，此时你要复兴以色列国吗？} 连这事祂也没有启示给他们。他们先前也问过这事（关于世界的终结）：正如那次祂回答时，引导他们离开以为救恩临近的想法，反而将他们投入危险之中；同样，这次也如此，只是更加温和。为使门徒不以为受了亏待，以为这些只是托词，请听祂说什么：祂立刻赐给他们那令他们欢喜的事：因为祂接着说：\BibleQuote{但圣神降临于你们身上，你们将获得能力，并要在耶路撒冷、全犹太和撒玛黎雅，直到地极，为我作证人。}\BibleRef{宗 1:8} 然后，为使他们不再追问，祂就被接升天。因此，正如先前祂以敬畏和说\BibleQuote{我不知道}来遮蔽他们的心智，这里祂也以被接升天来这样做。因为他们对此事极其热切，不会罢休；然而，他们不应该知道这事是非常必要的。请问，外邦人最不信什么？是世界将有终结，还是天主降生成人，且出于童贞女？但我羞于在这点上多言，仿佛这是件难事。再者，为使门徒不说：你为什么悬而不决？祂补充说：\BibleQuote{那圣父放在自己权柄中的。} 然而祂宣告圣父的权能与祂的权能是一体的：正如这话：\BibleQuote{因为圣父怎样复活死人并使他们活着，子也随意使谁活着。}\BibleRef{若 5:21} 如果在需要工作之处，祢以与圣父同等的权能行事；在需要知道之处，祢难道不以同一权能知道吗？然而，复活死人远大于知道日子。如果更大的事靠着权能，那么更小的事更是如此。\par

但正如我们看见一个小孩哭泣，固执地想要从我们这里得到某样对他无益的东西时，我们藏起那东西，向他展示我们的空手，说：\Quote{看，我们没有；}基督在这里对宗徒们也做了同样的事。但是，即使当我们向他展示[空手]时，小孩仍执拗地哭泣，意识到自己被欺骗了；然后我们离开他，离去，说：\Quote{某人叫我；}我们反而给他别的东西，以转移他的欲望，告诉他那东西比先前的好得多，然后匆匆走开；基督也是如此行事的。门徒们请求得到某样东西，祂说祂没有。起初祂吓唬了他们。然后他们再次请求现在就得到它：祂说祂没有；现在祂没有吓唬他们，但在展示了[空手]之后，祂这样做，并给了他们一个似乎合理的理由：\Quote{那日期}祂说，\BibleQuote{是父所定的。}什么？你竟不知道父的事！你认识祂，却不知道属于祂的事！而祢曾说：\BibleQuote{除了子，没有人认识父}\BibleRef{路 10:25}；并且\BibleQuote{圣神洞察一切，甚至天主的深奥}\BibleRef{格前 2:10}；而祢竟不知道这事！但他们不敢再问祂，怕祂又说：\BibleQuote{你们也是不明白的吗？}\BibleRef{玛 15:26}因为他们现在比以前更怕祂。\BibleQuote{但你们将领受圣神降临于你们身上的能力。}正如先前祂没有回答他们的问题（因为教师的职责不是教导门徒所选择的，而是教导他们需要学习的有益之事），同样，祂事先告诉他们这些，为使他们不致烦恼。事实上，他们仍然软弱。但为激励他们的信心，祂提升他们的灵魂，并隐藏了令人忧伤的事。由于祂很快就要离开他们，因此在这段话中祂没有说任何痛苦的事。但如何呢？祂将那些痛苦的事夸大为大事：几乎在说：\BibleQuote{不要怕：因为圣神降临于你们身上时，你们将领受能力；并要在耶路撒冷、全犹太和撒玛黎雅，直到地极，为我作证人。}因为祂曾说：\BibleQuote{外邦人的路，你们不要走，撒玛黎雅人的城，你们不要进}\BibleRef{玛 10:5}，祂在那里未说出的，在这里补充了：\BibleQuote{直到地极；}在说了这比一切更可怕的话之后，为使门徒不再质问祂，祂沉默了。\BibleQuote{说了这话，他们正观看时，就被接上升；有一朵云彩把祂接去，离开他们的眼界。}\BibleRef{宗 1:9}你看见他们确实传了道并完成了福音吗？因为祂赐给他们的恩赐是伟大的。祂说，在你们害怕的地方，即耶路撒冷，你们先在那里宣讲，然后直到地极。然后为证实祂所说的话，\BibleQuote{他们正观看时，就被接上升。}祂复活不是\BibleQuote{他们正观看时}，但祂被接上升是\BibleQuote{他们正观看时}。然而，因为在这里肉眼所见并不完全足够；因为复活时他们看见了结局，却未看见开始；而升天时他们看见了开始，却未看见结局：因为在复活中，看见开始是多余的，因为主亲自在场说了这些话，坟墓也清楚表明祂不在那里；但在升天中，他们需要由他人之口告知后续：既然他们的眼睛不足以向他们展示上面的高处，也不足以告诉祂是否真的升了天堂，还是只是表面上升了天堂，那么请看接下来发生的事。他们知道那是耶稣本人，因为祂曾与他们交谈，因为如果他们在远处看见，他们不可能凭视觉认出祂，但祂被接升天，是天使们亲自告诉他们的。注意如何安排，并非一切都由圣神完成，眼睛也各尽其分。但为什么\BibleQuote{有一朵云彩把祂接去}？这也是一个确切的记号，表明祂升入了天堂。不是火，像厄里亚的情形，也不是火马车，而是\BibleQuote{一朵云彩把祂接去}；这是天堂的象征，正如先知所说；\BibleQuote{祂以云彩为车辇}\BibleRef{咏 103:3}；这是指着圣父说的。因此他说\BibleQuote{在云上}；他是在说，这是天主能力的象征，因为没有别的能力被看见显现于云上。因为再听另一位先知说什么：\BibleQuote{上主乘驾轻云}\BibleRef{依 19:1}。因为就在他们非常专注地听祂说话，而且这是对一个非常有趣问题的回答，他们的心智完全被唤醒并十分清醒时，这件事发生了。在西乃山上，云彩也是因为祂：因为梅瑟也进入了黑暗，但那里的云彩不是因为梅瑟。祂不只是说\BibleQuote{我去}，免得他们再次悲伤，但祂说\BibleQuote{我派遣圣神}\BibleRef{若 16:5}；并且祂升入天堂，他们亲眼看见了。噢，他们被恩赐了何等景象！\BibleQuote{他们正凝视着天上}，记载说，\BibleQuote{祂上升时，忽然有两个人穿着白衣，站在他们面前，说：加里肋亚人，你们为什么站着望天呢？这位离开你们被接到天上的耶稣}——他们用了\BibleQuote{这位}作为指示词，说：\BibleQuote{这位离开你们被接到天上的耶稣，将来也要这样}——指示性地，\BibleQuote{这样}——\BibleQuote{你们看见祂怎样升了天，也要怎样降来。}（10-11节）再次，外貌令人振奋[\BibleQuote{穿着白衣}]。他们是天使，以人形显现。他们说：\BibleQuote{加里肋亚人}：他们通过说\BibleQuote{加里肋亚人}来表明自己为门徒所信任。因为这是含义：否则，告诉他们早已熟知的出身地有何必要？他们也通过外表吸引他们的注意，表明他们是来自天上。但为什么基督不亲自告诉他们这些事，而是由天使来说？祂事先已告诉他们一切；\BibleQuote{如果你们看见人子升到祂先前所在之处，又怎样呢？}\BibleRef{若 6:62}\par

况且天使并没有说\Quote{你们所看见被接去的}，而是说\InnerQuote{升入天上}：是\Quote{升天}这个词，而不是\Quote{被提}；\Quote{被接去}的说法属于肉身。为此他们又说：\InnerQuote{这位从你们中被接去的，将要这样来}，不是\InnerQuote{将被差遣}，而是\InnerQuote{将要来}。\BibleQuote{那升上的，也正是那降下的}\BibleRef{弗 4:10}。同样，\InnerQuote{一朵云彩接住了祂}：因为祂亲自登上了云彩。在这些说法中，有些适应门徒的理解，有些符合天主的尊威。现在，当他们观看时，他们的思想被提升：祂给了他们对第二次来临的性质不小的暗示。因为\InnerQuote{将要这样来}，意思是有肉身，这是他们渴望听到的；并且祂要\InnerQuote{这样}驾云再来审判。\InnerQuote{忽然，有两个人站在他们旁边。}为什么说是\InnerQuote{人}？因为他们完全装扮成人的样子，以免观看者承受不了。\InnerQuote{他们也说：}并且他们的话旨在安抚：\InnerQuote{为什么站着望天呢？}他们不再让他们在那里等候祂。这里，他们又说出更伟大的事，而将较小的事略去不谈。他们说\InnerQuote{祂将要这样来}，并且\InnerQuote{你们必须从天上等候祂。}至于其余的事，他们将他们从那景象引开，转向他们的话，免得他们因为看不见祂，就以为祂没有升天，甚至在他们说话时，祂就会在不察觉中临在。因为若他们先前说过：\BibleQuote{祢往哪里去？}\BibleRef{若 13:36}，现在更会这样说。\par

（复述）他们说：\InnerQuote{主啊，祢在这时候就要复兴以色列国吗？}他们如此深知祂的温和，以致在祂受难后，还问祂：\InnerQuote{祢要复兴吗？}然而祂先前曾对他们说：\InnerQuote{你们要听见战争和战争的风声，但末期还没有到}，耶路撒冷也不会被攻取。但现在他们问的是关于国度，而不是关于结局。此外，祂复活后并没有与他们长篇大论。于是他们提出这个问题，因为他们想，若这事实现，自己必大受尊荣。但祂（因为关于这个复兴，祂没有公开宣告它不会发生；他们何须学习这个？因此他们不再问：\InnerQuote{祢来临和世界终局的预兆是什么？}因为他们不敢说那话；而是问：\InnerQuote{祢要复兴以色列国吗？}因为他们以为有这样的国度），但祂，我说，在比喻中已经表明时间不近了，在这里他们问到时，祂也回答：\InnerQuote{你们将要领受圣神的能力，当圣神降临在你们身上时。}\InnerQuote{降临在你们身上}，不是\InnerQuote{被差遣}[为显示圣神同等的尊威。那么，你这反对圣神的人，怎敢称祂为受造物？]。\InnerQuote{你们要为我作证人。}祂暗示了升天。[祂说了这些话之后。] 这些他们先前听过，现在祂提醒他们。[\InnerQuote{祂被接去了。}] 已经表明，祂升上去了。[\InnerQuote{一朵云彩，等等。}] 圣经说：\BibleQuote{云彩和幽暗在祂的脚下}\BibleRef{咏 17:9}：因为\InnerQuote{一朵云彩接住了祂}这话正是宣告这一点：是说祂是天主。因为正如君王由御辇显示，所以御辇为祂被差派。[看，两个人，等等。] 好使他们不致发出哀叹，也不致如同厄里叟那样 \BibleRef{列下 2:12}，当他的老师被接去时，撕裂了外衣。他们说什么？\InnerQuote{这位从你们中被接去升天的耶稣，将要这样来。}并且，\BibleQuote{忽然有两个人站在他们旁边}\BibleRef{玛 18:16}。有充分的理由：因为\BibleQuote{凭两个证人的口，每句话都得成立}\BibleRef{申 17:6}；并且他们说了同样的话。又记载，他们穿着\InnerQuote{白衣。}就像他们先前在坟墓那里看见一位天使，甚至告诉了他们心里所想的；同样，这里也有一位天使传布祂的升天；虽然先知们早已屡次预言了这事，以及复活。\par

处处都有天使：如在诞生时，\BibleQuote{她所怀的，一位说，是由圣神来的}\BibleRef{玛 1:20}；又对玛利亚说：\BibleQuote{玛利亚，不要怕}\BibleRef{路 1:30}。在复活时：\BibleQuote{祂不在这里，祂已复活了，并在你们前走}\BibleRef{路 24:6}。\BibleQuote{你们来看！}\BibleRef{玛 28:6}。在第二次来临时也是如此。为了不使他们完全惊惶失措，因此加上：\BibleQuote{将要这样来}\BibleRef{玛 25:31}。他们稍微喘了一口气；如果祂真要再来，如果祂也要这样来，并且不是不可接近的！还有，那\InnerQuote{从他们中}被接去的说法，也不是随意添加的。关于复活，基督亲自作证（因为在一切事中，这事仅次于诞生，甚至超过诞生，最为奇妙：祂使自己复活）：因为祂说：\BibleQuote{你们拆毁这座圣殿，三天之内我要把它重建起来}\BibleRef{若 2:19}。因此，如果有人渴望看见基督；如果有人因未曾见过祂而悲伤；既闻此言，就当活出堪当的生命，他必定会见祂，且不会失望。因为基督将以更大的光荣来临，虽然\InnerQuote{这样}，以这种方式，带着肉身；并且看见祂从天降下，将是更加奇妙的事。但祂为何而来，他们没有补充。\par

\Quote{就这样来，等等。}这是复活的确证；因为如果祂带着身体被提升天，那么祂必定更带着身体复活了。那些不相信复活的人在哪里？请问，他们是谁？是外邦人，还是基督徒？因为我不知道。不，我清楚知道：他们是外邦人，他们也不相信创造之工。因为这两种否认是并存的：否认天主从无中创造万物，和否认祂使埋葬的身体复活。但是，为了羞于被人认为像那些\Quote{不知道天主的权能}\BibleRef{玛 22:29}的人，免得我们归咎于他们，他们便声称：我们这样说并非此意，而是因为身体没有必要。真正适时地可以说：\Quote{愚人说愚妄话。}\BibleRef{依 32:6}你们岂不羞于不承认天主能从无中创造吗？如果祂从已有的物质中创造，那祂与人又有何区别？但你问：恶从何而来？纵然你不知道从何而来，难道你因此就该在认知恶的事上引进另一种恶吗？由此产生了两个荒谬之处。因为如果你不承认天主从无中造了万有，那么你就更不知道恶从何而来；然后，你又引进另一种恶，即断言恶（\OriginalTerm{邪恶}{\GreekText{τὴν} \GreekText{κακίαν}}）是非受造的。现在考虑一下：当你想要寻找恶的根源时，你既不知道它，又添上另一种恶。追究恶的起源，但不要亵渎天主。我怎么亵渎了？他说。当你把恶说成拥有与天主同等的权能，即一种非受造的权能时。因为，请看保禄所说的：\Quote{其实，自从天主创世以来，祂那看不见的事，即祂永远的大能和神性，都可凭祂所造的万物，辨认洞察出来。}\BibleRef{罗 1:20}但魔鬼却要使两者都出自物质，以便没有任何东西使我们认识天主。请告诉我：哪一样更难：取本性为恶的东西（如果真有这样的东西；因为我按你们的原理说，因为本性没有恶），使之成为善，甚至成为善的合作者？还是从无中创造？哪一样更容易（我是指性质而言）：引入不存在的性质，还是取已有的性质并改变成它的对立面？在没有房屋的地方造成房屋，还是在房屋完全毁坏时使之重新存在？因为正如这是不可能的，那样也是不可能的：使一物变成它的对立面。告诉我，哪一样更难：制造香水，还是使污秽产生香水效果？说说看，哪一样更容易（既然我们使天主受我们的推理束缚；不，不是我们，而是你们）：形成眼睛，还是使瞎子在继续瞎的情况下看得见，甚至比看得见的人更敏锐？使瞎变成看见，聋变成听见？在我看来，另一方面似乎更容易。那么，你们承认天主能做更难的事，却不承认更容易的事？此外，他们还断言灵魂是出于天主的实体。你们看，这里有多少不敬和荒谬！首先，他们为要表明恶出自天主，便引进了比这更不敬的事，即恶与天主同等尊严，天主在它们之先并不存在，竟把这伟大的特权也赋予了它们！其次，他们断言恶是不可毁灭的：因为如果那非受造的可以毁灭，你们看这亵渎！于是便得出这样的结论：要么除了这些以外，没有任何东西出于天主；要么这些就是天主！第三，正如我以前说过的，他们在这点上自相矛盾，并为自己激起新的愤怒。第四，他们断言无序的物质具有如此\OriginalTerm{内在的}{\GreekText{ἐπιτηδειότητα}}能力。第五，他们断言恶是天主美善的原因，没有恶，美善就不是美善。第六，他们阻断了我们达到认识天主的途径。第七，他们把天主降格为人，甚至降到植物和木头中。因为如果我们的灵魂是出于天主的实体，而它转生到新身体的过程最终将它带入黄瓜、甜瓜和洋葱之中，那么天主的实体就会进入黄瓜之中！如果我们说，圣神在童贞玛利亚内形成了[我主身体的]圣殿，他们就嘲笑我们；如果我们说，祂居住在那属灵的圣殿里，他们又嘲笑；而他们自己却不羞于把天主的实体带到黄瓜、甜瓜、苍蝇、毛毛虫和驴子中，从而构想出一种新的偶像崇拜方式：因为不要像埃及人那样说：\Quote{洋葱是神；}而是：\Quote{神在洋葱里！}为什么你对天主进入身体的概念退缩？\InnerQuote{这令人震惊，}他说。那么，这更令人震惊。然而，实际上，这并不令人震惊——怎么可能呢？——这件事本身若是进入我们里面！\InnerQuote{但你们的概念确实令人震惊。}你们看见他们不敬的污秽了吗？——但他们为什么不希望身体复活？为什么他们说身体是恶的？那么请告诉我：你藉着什么认识天主？你藉着什么认识存在的事物？哲学家呢：他藉着什么是哲学家，如果身体在其中毫无作用？麻木感官，然后学习需要知道的事物！还有什么比灵魂更愚蠢呢，如果从一开始它的感官就麻木了？如果麻木一个部分，即大脑，就会使它完全受损；如果其余的都麻木了，它还有什么用？给我看一个没有身体的灵魂。你没有听见医生说：疾病的存在严重削弱灵魂？你要等到什么时候才上吊？身体是物质的吗？告诉我。\Quote{当然，是物质的。}那么你应该恨恶它。你为什么喂养它，为什么珍惜它？你应该摆脱这个监狱。但此外：\Quote{天主不能克服物质，除非祂\OriginalTerm{交织}{\GreekText{συμπλακἥ}}与物质纠缠：因为祂不能命令它（多么软弱！）直到祂与之结合，并\OriginalTerm{站立}{\GreekText{σταθἥ}}在其中站立。}然而，一个国王尚且可以发号施令成就一切；但天主却不能命令恶！简而言之，如果恶完全不参与任何善，它根本不能存在。因为恶不能存在，除非它抓住德性的某些属性；所以，如果它从前完全与善无涉，它早已消亡了：因为恶的情况正是如此。假设一个放荡的人，完全不约束自己，他能活十天吗？假设一个强盗，对任何人完全不顾良心，甚至对同伙也是如此，他能活下去吗？假设一个小偷，毫无羞耻，不知脸红，公开在公共场所偷窃。恶的本性不能存在，除非它至少获得一点善。因此，照这些人的说法，是天主使它们存在。假设一座恶人城，它能站立吗？假设他们不仅对善人，而且彼此之间都作恶。这样的城不可能站立。确实，\Quote{他们自称聪明，反而成为愚拙。}\BibleRef{罗 1:22}如果身体的本性是恶的，那么一切可见之物都是徒然和虚空的，无论是水、土地、太阳，还是空气；因为空气也是身体，虽然不坚固。那么，有必要说：\Quote{恶人向我讲说愚事。}\BibleRef{咏 118:85}但我们不要容忍他们，我们要堵住耳朵不听他们。因为确实有，确实有身体的复活。这在耶路撒冷的坟墓宣告了，这被鞭打时捆绑祂的柱子也宣告了。因为经上说：\Quote{我们曾与祂一同吃喝。}让我们相信复活，并做配得复活的事，好使我们获得将来要来的美善，藉着我们的主耶稣基督，祂与圣父及圣神，一同享有权能和荣耀，从现在直到永远，世世无穷。阿们。\par

\SourceRecord{Translated by J. Walker, J. Sheppard and H. Browne, and revised by George B. Stevens.；Nicene and Post-Nicene Fathers, First Series；Vol. 11.；Edited by Philip Schaff.；Buffalo, NY: Christian Literature Publishing Co.,；1889.；<http://www.newadvent.org/fathers/210102.htm>.}

% structure: p0001 source=h1 target=section reason=page-title

\section{论宗徒大事录第三篇讲道}

\BibleRef{宗 1:12}\par

\BibleQuote{他们就从名叫橄榄的山回了耶路撒冷，这山离耶路撒冷有安息日可走的路程。}\par

\BibleQuote{他们就回了，}据说：即他们听见之后。因为如果天使没有\OriginalTerm{延迟}{\GreekText{ὑπερέθετο}}将他们指向另一次来临，他们原是受不了的。在我看来，这些事也发生在安息日；因为他不会如此具体说明距离，说\BibleQuote{从名叫橄榄的山起，离耶路撒冷有安息日可走的路程，}除非他们那时正在安息日走了某一定距离。\BibleQuote{他们一进来，}经上说，\BibleQuote{就上了他们所住的一间楼房，}可见他们复活后仍留在耶路撒冷：\BibleQuote{有伯多禄、雅各伯、若望，}不再只提后者与他的兄弟，而是与伯多禄一起提到两人：\BibleQuote{并安德肋、斐理伯、多默、巴尔多禄茂、玛窦、阿耳斐的儿子雅各伯、热诚者西满、雅各伯的兄弟犹达。}\BibleRef{宗 1:13}他很好地提到了门徒们：因为既然一人出卖了基督，另一人不信，他就借此表明，除了第一个，他们都得以保全。\par

\BibleQuote{这些人都同心合意地同几个妇女一同祈祷。}\BibleRef{宗 1:14}因为这在诱惑中是一强力武器；他们对此已受训练。[\Quote{同心合意。}]好。\OriginalTerm{好}{\GreekText{καλώς}}。此外，当前的困境驱使他们如此：因为他们极其害怕犹太人。\BibleQuote{同几个妇女，}据说：因为经上说过她们跟随了他：\BibleQuote{并同耶稣的母亲玛利亚。}\BibleRef{路 23:55}那么，[那位门徒]怎么那时把她接到自己家里呢？\BibleRef{若 19:26}但那时主已将他们重新聚在一起，于是又回去了。\BibleQuote{同他的弟兄们。}\BibleRef{若 17:5}这些人以前也不信。\BibleQuote{在那些日子，}经上说，\BibleQuote{伯多禄在门徒中间站起来说：}\BibleRef{宗 1:15}他既热心，又受基督委托牧养羊群，且在尊荣上有优先，故此常率先发言。（\BibleQuote{一起的人数约有一百二十名。}\Quote{诸位仁人弟兄，}他说：\BibleQuote{这圣经必须应验，圣神从前藉着……所说的，}\BibleRef{宗 1:16}等等）为什么他不求基督另给一人代替犹达斯呢？这样更好。首先，他们忙于别事；其次，基督与他们同在的最大证明即是：正如祂在他们中间时选择了人，如今祂不在时也照样选择。这为他们并非小事，能得安慰。但请注意伯多禄如何凡事取得共识，决不独断。他这样说并非无意义。但请注意他如何就所发生的事安慰他们。事实上，所发生的事令他们惊惶不小。因为若如今有许多人推敲此事，那么当时他们该说些什么呢？\par

\BibleQuote{诸位仁人弟兄，}伯多禄说。因为若主称他们为弟兄，何况他呢。[\Quote{诸位仁人，}他说：]他们全体都在场。请看教会的尊荣，天使般的境界！那里没有分别，\BibleQuote{不分男也不分女。}但愿现今的教会也是如此！那里没有人心里装满世俗之事，没有人焦虑地想着家务。诱惑何等有益，苦难何等有利！\par

\Quote{这经}他说，\Quote{必须应验，就是圣神从前所说的。}他总是用预言来安慰他们。基督也是这样，在一切事上。同样，他在这里表明，并没有发生什么奇怪的事，而是早已预言过的。\Quote{这经必须应验}他说，\Quote{就是圣神藉达味的口预先所说的。}他没有说达味，而是说圣神藉着他。请看这位作者在书的一开头就教导了怎样的道理。你看见了吗？我在这部作品的开头说，这本书是圣神的政制，并不是没有原因的。\Quote{就是圣神藉达味的口预先所说的。}请注意他如何将达味\OriginalTerm{归为己用}{\GreekText{οἰκειοὕται}}；并且这对他们是有利的，因为这话是达味说的，而不是别的先知说的。\Quote{关于犹达斯}他说，\Quote{他曾是向导。}这里再次请注意这个人的哲学气质：他如何不轻蔑地提到他，也没有说\Quote{那个恶徒}\Quote{那个歹徒}，而是简单地陈述事实；他甚至没有说\Quote{出卖他的那个人}，而是尽力将罪责推给别人；他也没有严厉谴责这些人。\Quote{就是给那些捉拿耶稣的人作向导的}他说。此外，在他说明达味在哪里说了这话之前，他先讲述了犹达斯的情况，以便从眼前的事中获取对将来的事的保证，并表明这人已经得到了他应得的报应。\Quote{因为他本来被列在我们中间}他说，\Quote{并且得了这职务的一份。这人用不义的代价买了一块田地。}（第17-18节）他使他的言论转向道德，并暗地里提到了恶行的原因，因为这带有责备。而且他没有说犹太人，而是说\Quote{这人}买了它。因为软弱之人的心思，正如关注现在的事一样，不关注将来的事，所以他讲论了立即施加的惩罚。\Quote{他俯冲下来，肚腹崩裂}他说，\Quote{所有的肠子都流了出来。}他很好地详述了惩罚而非罪过。\Quote{并且这事传遍了所有住在耶路撒冷的人，以致那块田地在他们的乡谈中被称为阿刻耳达玛，就是血田的意思。}\BibleRef{宗 1:19}犹太人给它起这个名字，不是因为这个原因，而是因为犹达斯；但在这里，伯多禄使它具有这层含义，并且当他提出反对者作为证人时，无论是通过他们命名的事实，还是通过说\Quote{在他们的乡谈中}，他指的就是这个意思。\par

此后，他恰当地引用了先知的话，说：\BibleQuote{因为\BibleBook{圣咏}{集}上记载：愿他的居所荒凉，无人居住。}\BibleRef{宗 1:20}\BibleRef{咏 68:25} 这是指着那块田地和住所说的。\BibleQuote{至于他的监督职分让别人取得}，那是指他的职务，他的司祭职。所以他说，这不是我的建议，而是那位预先宣告这些事的天主的建议。因为他为了不显得自己承担一件大事，且是基督曾做过的事，他引用了先知作为证人。\BibleQuote{所以，从这些一直与我们为伴的人中，必须选立一位。}\BibleRef{宗 1:21} 为什么他也让他们参与这事？为了使这事不致成为纷争的对象，也不致引起他们之间的争论。因为如果宗徒们自己曾经这样做，那么那些人就更有可能这样做。他一直是避免这种情况的。因此，他在开头说：\BibleQuote{诸位弟兄，必须从你们当中选出。} 他让全体众人来做决定，从而既使被选者受人尊敬，又使自己远离对其他人的嫉妒。因为这样的场合总会引起大恶。既然必须有人被指派，他引用了先知作为证人；但从哪些人中呢？他说：\BibleQuote{从这些始终与我们为伴的人中。} 如果说必须由那些配得的人自己站出来，那就会侮辱其他人；但他却把事情交由时间长短来决定；因为他不是简单地说：\BibleQuote{这些与我们为伴的人}，而是说：\BibleQuote{从若翰施洗起，直到主耶稣从我们中被接去的日子为止，凡主耶稣在我们中间出入的整个时期，必须有一人与我们一起作祂复活的见证。}\BibleRef{宗 1:22} 使他们的\OriginalTerm{团体}{\GreekText{ὁ} \GreekText{χορὸς}}不致残缺不全。那么，为何不让伯多禄亲自选举呢？动机是什么？是为了不让他显得是施恩于人。此外，他当时尚未领受圣神。\BibleQuote{他们就推选了两个人：若瑟·巴尔撒巴，号称犹斯托的，和玛弟亚。}\BibleRef{宗 1:23} 不是他指定了他们，而是他提出了这样的动议，同时指出，这甚至不是他自己的想法，而是自古藉由预言而来的；所以他扮演的是解释者，而不是教导者的角色。\BibleQuote{若瑟·巴尔撒巴，号称犹斯托的。} 或许给出了两个名字，因为同名的也有别人，因为宗徒中也有几个相同的名字：如雅各伯和阿尔斐的儿子雅各伯；西满伯多禄和西满热诚者；犹达斯（雅各伯的兄弟）和加略人犹达斯。不过，这个称号可能是由于生活的转变而产生的，也很可能是由于道德品格的改变。\BibleQuote{他们推选了两个人，} 经上说，\BibleQuote{若瑟·巴尔撒巴，号称犹斯托的，和玛弟亚。他们就祈祷说：主啊，祢知道众人的心，求祢指示，这两个人中，祢拣选了哪一个，使他分担这职务和宗徒的职位，犹达斯因背信舍弃了，去了他自己的地方。}\BibleRef{宗 1:23-25} 他们提及犹达斯的罪过做得很好，从而表明他们要求的是见证人；不是增加人数，而是不让数目减少。\BibleQuote{他们就摇签}\BibleRef{宗 1:26}（因为圣神尚未赐下），\BibleQuote{签落在玛弟亚身上：他就被列入十一位宗徒之中。}\BibleRef{宗 1:26}\par

\BibleQuote{然后，他们从名叫橄榄山的山上回了耶路撒冷}（回顾），\BibleQuote{这山靠近耶路撒冷，约有安息日可走的路程。}这样就没有很长的路要走，以免成为他们因战栗恐惧而惊慌的原因。\BibleQuote{他们进了城，便上了楼。}他们不敢在城里露面。他们还上了楼，这样做也很好，因为这样就更不容易立即被逮捕。\BibleQuote{他们继续同心合意地祈祷。}你看到他们多么警醒吗？\BibleQuote{恒心祈祷，}和\BibleQuote{同心合意，}好像只有一个灵魂，持续在其中：这是两件值得称赞的事。\BibleQuote{他们居住的地方，}等等，直到\BibleQuote{耶稣的母亲玛利亚和祂的弟兄们。}现在若瑟可能已经去世了：因为不能设想，当弟兄们成为信徒时，若瑟却没有相信；实际上他比任何人都先信了。可以肯定的是，我们无处发现他把基督只当作人来看待。正如祂的母亲所说：\BibleQuote{你的父亲和我忧苦地找你。}\BibleRef{路 2:48} 在另一场合，\BibleQuote{你的母亲和你的弟兄们找你。}\BibleRef{玛 13:47}所以若瑟比所有人都先知道这事。基督对他们（弟兄们）说：\BibleQuote{世界不能恨你们，却恨我。}\BibleRef{若 7:7}\par

再请看雅各伯的谦和。他曾接受耶路撒冷的主教职，而在这里他什么也没说。也请注意其他宗徒的极大谦和，他们如何将宝座让给他，不再彼此争执。因为那个教会如同在天上：与俗世的事务无关；它的荣耀不在于墙壁，也不在于人数，而在于组成集会的那些人他们的热心。他们约有\BibleQuote{一百二十人，}经上说。那七十人或许是基督亲自选定的，以及其他更热心的门徒，如若瑟和玛弟亚。\BibleRef{宗 1:14}还有妇女，他说，有许多，跟随了祂。\BibleRef{谷 15:41}\BibleQuote{名字的数目一起。}同心合意地，他们时常如此。\par

\Quote{诸位仁人弟兄，}等。此处有提供教师的先见；此处是首位按立教师的人。他并没有说：\Quote{我们足够了。}他是如此远离一切虚荣，只专注于一件事。然而他拥有与他们全体相同的按立权柄。这些事之所以如此进行，全因这人高尚的精神，并且当时的管理职务并非荣耀之事，而是对被管理者之眷顾。这既不会使被选者骄傲（因为他们被召是去面对危险），也不会使未被选者因感到受辱而抱怨。但现在事情却不如此进行，反而截然相反。——因为请注意，他们共有一百二十人，而他只从全体中选出一人：他完全有权这样做，因为他被委派管理他们；因为基督曾对他说：\BibleQuote{你回头以后，要坚固你的弟兄。}\BibleRef{路 22:32}\par

\BibleQuote{因为他原是我们中被列为一数的}（\OriginalTerm{他拥有此事最初的权柄}{\GreekText{πρῶτος} \GreekText{τοῦ} \GreekText{πράγματος} \GreekText{αὐθεντεῖ}} 于 A.B.C 抄本中缺失），伯多禄说。为此理当推举另一人，代替他作见证。请看他是如何效法他的恩师，总是从圣经中论述，且尚未谈及基督的事，即基督曾多次亲自预言此事。他并未提及圣经中论及犹达斯背叛之处，例如：\BibleQuote{邪恶的口与欺诈的口已张开攻击我}\BibleRef{咏 108:1}，而是只提及他的惩罚；因为这对他们最有益处。这也再次显示了主的慈爱：\BibleQuote{因为他原是我们中被列为一数的}（\OriginalTerm{因为这对他们最有益处}{\GreekText{τοῦτο} \GreekText{γὰρ} \GreekText{αὐτοὺς} \GreekText{μάλιστα} \GreekText{ὠφέλει}}）。\par

\SourceRecord{Translated by J. Walker, J. Sheppard and H. Browne, and revised by George B. Stevens.；Nicene and Post-Nicene Fathers, First Series；Vol. 11.；Edited by Philip Schaff.；Buffalo, NY: Christian Literature Publishing Co.,；1889.；<http://www.newadvent.org/fathers/210103.htm>.}

% structure: p0001 source=h1 target=section reason=page-title

\section{\WorkTitle{论宗徒大事录第四讲}}

\BibleRef{宗 2:1-2}\par

\BibleQuote{五旬节日一到，众人都聚集在同一地方。忽然从天上来了一阵响声。}\par

你看到预像了吗？这五旬节是什么？是镰刀割禾、收成入仓的时候。如今看这实事：到了镰刀——也就是圣言——下镰之时，因为正如镰刀般锋利的圣神降临了下来。请听基督的话：\BibleQuote{你们举目观看，田地已经发白，可以收割了。}\BibleRef{若 4:35}又说：\BibleQuote{庄稼实在多，工人却少。}\BibleRef{玛 9:38}但作为这收获的初果，祂亲自取了〔我们的本性〕，并把它带到了高处。祂自己首先下了镰刀。因此，祂也称圣言为种子。经上说：\BibleQuote{到了五旬节日}\BibleRef{路 8:5}，即大致在五旬节期间。因为当前的事件也必须在庆节期间发生，好使那些目睹基督受难的人也目睹这些事。\BibleQuote{忽然从天上来了一阵响声。}\BibleRef{宗 2:2}为什么这件事没有以可感觉的表征发生？原因在此：如果在这种情况下，人们还说\Quote{他们被新酒灌满了}，否则他们又会说什么呢？而且不止是\BibleQuote{来了响声}，而且是\BibleQuote{从天上来}。突然性也惊动了他们，把众人聚集到那地方。\BibleQuote{好像暴风刮来}：这表示圣神的极其强烈。\BibleQuote{充满了整个屋子}：以致在场的人都信了，并且（\EditorialNote{希腊文\OriginalTerm{他们}{\GreekText{τούτους}}}）就这样显为配得。还不止此；更令人敬畏的是：经上说\BibleQuote{又有舌头出现，好像火焰，分开}\BibleRef{宗 2:3}。请注意，总是\Quote{好像}，这是对的，免得你对圣神有粗俗可感的观念。还有\Quote{好像一阵暴风}：所以并不是真的风。\BibleQuote{好像火焰}。因为当圣神要显示给若翰时，祂以鸽子的形状降在基督头上；但现在，当一大群人即将归化时，却是\BibleQuote{好像火焰。并且停留在他们每人身上}。意思是它留下并安息在他们身上。因为\Quote{停留}表示安稳和持久。\par

难道是降临在十二人身上？不是；而是降临在那一百二十人身上。因为伯多禄不会无缘无故引用先知的话说：\BibleQuote{到末日，天主说：我要将我的神倾注在一切有血肉的人身上；你们的儿子们和女儿们要说预言，你们的青年人要见异象，你们的老年人要做异梦。}\BibleRef{岳 2:28}\BibleQuote{他们都充满了圣神。}\BibleRef{宗 2:4}因为，为使效果不单单是恐惧，所以同时有\BibleQuote{以圣神和火。于是他们就开始说起别国的话来，正如圣神赐给他们发言。}\BibleRef{玛 3:11}他们领受的标记没有别的，只有这第一个；因为这对他们来说是崭新的，也不需要别的标记。作者说：\BibleQuote{它停在每人身上。}现在请注意，那没有被选上如同玛弟亚一样的人，再也没有理由悲伤了。他说：\BibleQuote{他们都充满了}，不仅领受了圣神的恩宠，而是\BibleQuote{充满了。于是他们就开始说起别国的话来，正如圣神赐给他们发言。}如果不是其余的人也分受了，就不会说\Emph{所有}，连宗徒们也在场。因为若非如此，既然上文曾特别提名提到宗徒们，他就不会现在把他们与其余的人合在一起。因为如果单是提到他们在场时，他就分开提及宗徒，那么在此假设的情况下，他更会这样做了。请注意，当一个人\Emph{恒心祈祷}，当一个人处于爱德之中时，圣神就临近了。这也使他们想起另一个景象：因为天主在荆棘中也是以火显现的。正如圣神赐给他们发言的能力，\OriginalTerm{发言}{\GreekText{ἀποφθέγγεσθαι}}\BibleRef{出 3:2}，因为他们所说的是\OriginalTerm{深奥言论}{\GreekText{ἀποφθέγματα}}，即深奥的话语。经上说：\BibleQuote{那时，有虔诚的犹太人，从天下各国来，住在耶路撒冷。}\BibleRef{宗 2:5}他们住在那里的这一事实是虔诚的标志：他们来自这么多民族，竟然离开了自己的国家、家乡和亲属，住在那里。因为经上说：有来自天下各国的虔诚犹太人住在耶路撒冷。当这事传开时，众人就聚集，都惊慌失措。\BibleRef{宗 2:6}因为这事发生在一间屋里，所以他们自然从外面聚集而来。众人\Emph{惊慌失措}，一片骚乱。他们惊讶不已；\BibleQuote{因为人人都听见他们用各人的语言说话。他们都惊讶，惊奇地说：看，这些说话的不都是加里肋亚人吗？}\BibleRef{宗 2:7}他们立刻转目注视宗徒们。\BibleQuote{怎么我们每人听见他们说我们生来所用的方言呢？帕提雅人、玛待人、厄蓝人，住在美索不达米亚、犹太、卡帕多细亚、本都、亚细亚、夫黎基雅、旁非里雅、埃及，及基勒乃附近利比亚一带的人，}请注意他们从东到西的分布：\BibleQuote{以及罗马来的旅客、犹太人和归依犹太教的人、克里特人和阿拉伯人，我们都听见他们用我们的方言讲论天主伟大的作为。}他们都惊讶不已，疑惑纷纷，彼此说：这是什么意思？还有些人讥笑说：他们喝醉了新酒。哦，极端的愚昧！哦，极端的恶毒！因为那时甚至还不是喝酒的季节，因为正是五旬节。这使得情况更糟：当那些人承认时——他们是犹太人、罗马人、归依者，甚至可能曾钉死祂的人——而这些人，在如此伟大的标记之后，却说：\BibleQuote{他们喝醉了新酒！}\par

但让我们回顾一下从一开始所说的内容。（重述。）他说：\BibleQuote{到了五旬节那天，}等等，\BibleQuote{它充满了，}\BibleQuote{那间屋子。}那阵风\OriginalTerm{风}{\GreekText{πνοὴ}}好似一池水。这象征了丰盛，正如火象征了猛烈。这种情况从未发生在先知们身上：因为对于未被灌醉的灵魂，这样的降临不会带来很大扰动；但是\Quote{当他们喝足了，}才真正像这里一样，但先知们的情况却不同。\BibleRef{则 3:3}有一卷书卷赐给了他，厄则克耳吃了将要说出的话。经上说：\BibleQuote{它在他口中成为甘甜如蜜。}（再者，天主的手触摸了另一位先知的口舌；但在这里是圣神亲自：\BibleRef{耶 1:9}，祂在荣耀上与圣父和圣子同等。）另一方面，厄则克耳又称之为\BibleQuote{哀歌、悲叹和灾祸。}\BibleRef{则 2:10}对他们而言，以书卷的形式出现是合适的，因为他们仍然需要比喻。那些先知只同一个民族打交道，而且是同自己的百姓；但这些人却要与整个世界打交道，与素不相识的人打交道。此外，厄里叟藉着一件外衣领受了恩宠\EditorialNote{列王纪下 13}；另有人藉着油，如达味\BibleRef{撒上 16:13}；梅瑟藉着火，如我们在荆棘丛中读到他的事。\BibleRef{出 3:2}但在目前的情况下却并非如此；因为火本身停在他们的头上。（但为什么火没有出现充满屋子？因为他们会害怕。）但故事表明，这里和那里是一样的。因为你们不应停留于此，即\BibleQuote{有分开的舌头显现给他们，}但要注意它们是\BibleQuote{火的。}这样的火能够点燃无穷的燃料。而且，说\Emph{分开的}很恰当，因为它们来自同一根源；这样你们可以知道，那是从护慰者那里发出的作为。\par

但请观察，那些人也是先被显示为堪当，然后才领受圣神，因为堪当。例如，达味：他在羊圈中所做的，与他在胜利和凯旋后所做的一样；这是为了显示他的信德是多么单纯和绝对。再看梅瑟轻视王权，舍弃一切，四十年后引领百姓\BibleRef{出 2:11}；撒慕尔在圣殿中服事\BibleRef{撒上 3:3}；厄里叟舍弃一切\BibleRef{列上 19:21}；厄则克耳又藉此后发生的事显明。你看，同样地，这些人也舍弃了他们所有的一切。他们也藉所受的苦难，认识了人性的软弱；他们明白了，他们所做的这些善工并非徒然。\BibleRef{撒上 11:6} 连撒乌耳也是先获得见证他是好人，然后才领受圣神。但没有人像这里一样领受。梅瑟是先知中最大的，然而当别人要领受圣神时，他自己却受了亏损。但这里并非如此；反而正如火能点燃它愿意的许多火焰，同样，这里圣神的丰盛得以显示，因为每人领受了一个圣神的泉源；正如主亲自预言的，凡信祂的人，要有\Quote{涌到永生的水泉}\BibleRef{若 4:14}。这理当如此。因为他们不是去与法郎辩论，而是与魔鬼摔跤。奇妙的是，当他们被派遣时，没有提出反对；他们没有说自己是\Quote{拙口笨舌的}\BibleRef{出 4:10}。因为梅瑟已教导了他们更好的。他们没有说太年轻\BibleRef{耶 1:6}。耶肋米亚已使他们有智慧。然而他们听到了许多可怕的事，远比古时那些人的事更大；但他们惧怕反对。因为他们是光明之子，是天上之事的仆役[\Quote{忽然从天上来了一阵响声}等等]。对古时那些人，没有人从天上来显现，因为他们还在追求地上的召叫；但现在，人已升到高处，圣神也大能地从高处降下。\Quote{好像暴风刮来}；借此显明，没有什么能抵挡他们，他们必将一切敌人像尘土一样吹散。\Quote{充满了整个房屋}。房屋也象征世界。\Quote{舌头如火焰显现}，等等，以及\Quote{众人聚集，都惊慌失措}。请观察他们的虔诚；他们不妄下判断，而是困惑不解；而那些鲁莽的人立刻断言：\Quote{这些人喝新酒醉了。}这些\Quote{从天下各国来的虔诚人}住在那里，是为了能依法在一年三次上圣殿。请注意，作者无意奉承他们。因为他没有说他们发表了什么意见；而是什么呢？\Quote{这声音一响，众人都来聚集，惊慌失措}。他们惊慌是理所当然的；因为他们以为事情要和他们作对，因为对基督所行的暴行。良心也搅动他们的灵魂，手上还沾着血，一切都令他们恐惧。\Quote{看，这些说话的不都是加里肋亚人吗？}因为这确实是公认的。[\Quote{我们怎么听见}]声音如此之大，使他们恐惧。[\Quote{各人用我们生来所用的言语}等等]因为它找到了当时聚集在那里的世界大部分的人。这鼓舞了宗徒们：因为他们不知道用帕提亚语说话是怎么回事，但现在从那些人的话中得知了。这里提到敌对他们的民族：克里特人、阿拉伯人、埃及人、波斯人；显明他们将要征服所有这一切。至于他们在那些地方，许多人是在那里被掳的；或者，法律的教义已在那些地方的外邦人中传播。因此，见证来自各方：来自本地人、外邦人、归依者。\Quote{我们听见他们用我们的语言，讲论天主的奇事。}不仅因为他们说话（用他们的语言），而且他们所说的事是奇妙的。他们困惑不解是理所当然的：因为从未有过这样的事。请观这些人坦率。他们惊奇困惑，说：\Quote{这是什么事？}但\Quote{另有些人讥诮说：\InnerQuote{这些人喝新酒醉了}}\BibleRef{若 8:48}，因此讥诮。多么厚颜无耻！有什么奇怪呢？因为连主自己赶鬼时，他们也说祂有鬼！因为就是这样；凡有厚颜无耻之处，它只有一个目的：不顾一切地说话，不管说什么；不是要人说出真实切题的话，而是要他不顾一切地说话。[\Quote{他们喝新酒醉了}]这自然是常事（不是吗？），那些身处如此危险之中，惧怕最坏之事，且如此沮丧的人，竟有勇气说出这样的话！请注意：由于这不像真的（因为在那个时辰他们不会喝太多），他们便将整个事情归咎于酒的性质，说他们\Quote{充满了}酒。\Quote{伯多禄却同十一位宗徒站起来，高声对他们说。}先前你看到他的远见，这里你看到他的勇敢。因为他们若惊奇诧异，那么他这个未受教育和无知的人，在如此众多的人群中能找到话语，岂不也是奇事？若一个人在朋友面前说话会紧张，那么在敌人和嗜血者面前岂不更紧张？他立刻用自己的声音表明他们没有醉，不是像占卜者那样神魂超拔；也表明他们不是被某种强制力量所逼。\Quote{与十一位宗徒}是什么意思？他们用一个共同的声音表达自己，他是众人的口。那十一位站在旁边，为他所说的话作证。\Quote{他高声}说，这是说他以极大的信心讲话，好让他们看到圣神的恩宠。那个连一个使女的问话都受不了的人，现在在众人中间，在充满杀气的环境中，如此自信地讲论，这本身就成了复活无可置疑的证据：在那些能够嘲笑和戏弄这样事情的人中间！这是何等的厚颜无耻！何等的亵渎！何等的无耻！因为圣神在哪里临在，祂就会用泥土造出金子。我请求你，现在看看伯多禄：仔细察看过那个胆怯、缺乏悟性的人，正如基督说：\Quote{连你们还是不明白吗？}\BibleRef{玛 15:16}这个人在那次奇妙的宣认之后被称为\Quote{撒殚}\BibleRef{玛 16:23}。也想想宗徒们的一致。他们自己把说话的任务让给他；因为不必所有人都说话。\Quote{他高声}向众人极其勇敢地讲话。成为属神的人就是这样的事！只要我们也将自己作成堪配天上恩宠的状态，一切就变得容易了。因为一个火一般的人落入干草中，不会受伤害，反而伤害别人：他不会受任何伤害，而攻击他的人却自取灭亡。因为这里的情况正是如此：如同一个携带干草的人攻击一个带火的人；同样，宗徒们以极大的勇气迎战这些敌人。\par

因为人数虽如此众多，对他们有什么损害呢？他们不是耗尽了自己的一切怒气吗？他们不是把痛苦转到了自己身上吗？在众人中，有谁曾像那些人那样，既充满恐惧又充满怒气呢？他们岂不焦虑、惊慌、战栗吗？请听他们所说的话：\BibleQuote{你们愿意把此人的血归到我们身上吗？}\BibleRef{宗 5:28}宗徒们岂非与贫穷和饥饿争战，与耻辱和恶名争战（因为他们被视为骗子），与讥讽、愤怒和嘲笑争战？——因为他们身上遇见了相反的事：有人嘲笑他们，有人惩罚他们——他们岂非成了整个城巿愤怒与戏弄的对象？暴露于派系与阴谋，火与剑，与野兽之中？战争岂非从四面八方以千万种形式围攻他们？而这一切事对他们心灵的影响，岂不比他们在梦中或图画中看见它们更甚？他们以赤身露体迎战所有全副武装的人，尽管所有人对他们都有任意权柄[对付他们的，是]：统治者的恐怖、武力的强压、城邑与坚固城墙；没有经验，没有口才，且处于完全普通人的状态，对阵耍弄妖术者、骗子、全体智者、修辞学家、在学园与逍遥学派走廊里腐朽的哲学家——他们与所有这些人大战一场。那个曾在湖边谋生的人，如此完全地制胜了他们，仿佛他毫不费力，甚至像与哑鱼角力一样：因为他所须智取的对头，确实比鱼更沉默，所以他就如此轻易地胜过了他们！而柏拉图，那个在他时代胡言乱语的人，如今沉默了，而这个人却在四处发出他的声音；不仅在他的同胞中，也在帕提亚人、玛代人、厄蓝人中间，在印度，在地球的各个角落，直到地极。如今，希腊自大在哪里？雅典的名声在哪里？哲学家的妄言在哪里？加里肋亚人、贝特赛达人，那个粗野的乡下人，已经战胜了他们全部。你们岂不羞愧——承认吧——一听到击败你们的那个人的家乡之名？但若你们也听到他的名字，知道他叫\OriginalTerm{刻法}{Cephas}，你们更会掩面。这，这正是你们的彻底失败；因为你们以此为耻辱，以口若悬河为赞美，以不善辞令为耻辱。你们没有走上应当选择的道路，反而离开了那又容易又平坦的王道，踏上了一条崎岖、陡峭、艰难的路。因此，你们没有获得天国。\par

那么，有人会问，基督为何没有影响柏拉图和毕达哥拉斯呢？因为伯多禄的心智远比他们的心智更具哲学性。他们实在是孩子，被虚荣心四处摇摆；而这个人却是一位哲学家，能够接受恩宠。如果你嘲笑这些话，并不奇怪；因为从前那些人也曾嘲笑，并说这些人被新酒灌满了。但后来，当他们遭受那些悲惨的灾难，其苦楚超越一切时；当他们看见自己的城市坍塌，火焰熊熊，城墙被夷为平地，以及那多种言语无法表达的疯狂恐怖时，他们便不再嘲笑了。那时你也会嘲笑，如果你有心去嘲笑，当地狱的时刻临近，当火焰为你的灵魂点燃之时。但我为何谈论未来呢？让我向你展示伯多禄是什么，柏拉图的哲学又是什么？现在让我们考察他们各自的生活习惯，看看他们各自的追求。一个将时间浪费在一套无用无益的教条上，并如他所说，进行哲学思考，为的是让我们学到，我们哲学家的灵魂变成了一只苍蝇。说得最真切，一只苍蝇！并非真的变成了苍蝇，但一只苍蝇必定占据了那曾居于柏拉图灵魂的住所；因为对于这样的思想，还有什么比苍蝇更相配呢！这个人充满讽刺和嫉妒之情，反对其他所有人，仿佛他的雄心就是既不从自己的头脑中也不从他人的观点中引入任何有用的东西。因此他从别人那里接受了灵魂转世，从自己那里产生了《共和国》，在其中他制定了那些充满粗鄙丑恶的法律。他说，让女人公有，让处女裸体，让她们在爱人面前摔跤，也让他们有共有的父亲，并让所生的孩子共有。但在我们这里，并非自然造成共有的父亲，而是伯多禄的哲学做到了这一点；至于那另一种哲学，它却取消了所有的父权。因为柏拉图的体系只倾向于使真正的父亲几乎不为人知，而引进了虚假的父亲。它使灵魂陷入了一种醉态和污秽的翻滚。他说，让所有人毫无畏惧地与女人交合。我不检查诗人的格言，是因为我可能被指控重提神话。然而我所说的神话比那些还要可笑得多。诗人哪里想出过如此怪诞的东西？但（不讨论他的其他格言），你怎么看待这些——当他给女性装备武器、头盔和胫甲，并说人类没有理由不同于犬类！因为他说，狗，雌雄都同样做公有的事，所以让女人做和男人同样的工作，让一切都颠倒过来。魔鬼总是通过他们来证明我们的种族不比野兽更高贵；事实上，有些人已发展到如此荒谬的地步（\OriginalTerm{虚妄}{\GreekText{κενοδοξίας}}），以至于断言无理性的受造物具有理性。且看他是以多少种不同的方式在那些人的头脑中肆虐！因为当他们的领袖人物断言我们的灵魂进入苍蝇、狗和野兽的身体时；后人对此感到羞愧，便堕入了另一种污秽，赋予野兽一切理性科学，并断言那些因我们而存在的受造物在各方面都比我们更尊贵！他们甚至归给它们预知和虔诚。他们说，乌鸦认识天主，渡鸦也同样，它们拥有预言的天赋，能预知未来；它们之间有正义、政体和法律。也许你不相信我告诉你的事。这不足为怪，因为你所受的是纯正教义的养育；因为如果有人以此种食物为生，他决不会相信有人能从吃粪中找到乐趣。按照柏拉图，他们的狗也会嫉妒。但当我们告诉他们这些事是神话且充满荒谬时，他们说：\Quote{你们没有进入（\OriginalTerm{在内}{\GreekText{ἐ}} \OriginalTerm{理解}{\GreekText{νοήσατε}}）更高深的意义。}不，我们没有进入你们这种超越的胡说，但愿我们永远不会：因为这当然需要一个极其深奥的头脑，来告诉我，所有这些不敬和混乱要成为什么。你们这些无知的人，是在用乌鸦的语言说话，如同孩子们（玩耍时）那样吗？因为你们实在是孩子，就像他们一样。但伯多禄从未想过说任何这些事：他发出了一种声音，如同黑暗中闪耀的大光，一种驱散了整个世界迷雾和黑暗的声音。再者，他的举止是多么\OriginalTerm{温和}{\GreekText{ἐπιεικὲς}}，多么体谅人；多么超越一切虚荣；他如何毫无自高自大地仰望天堂，甚至在复活死人时也是如此！但如果那些无知者中的任何一个（当然只是在幻想中）有能力做任何类似的事，难道他不会立刻为自己寻求祭台和圣殿，并想与神明平等吗？因为事实上，当没有这样的征兆出现时，他们却永远沉溺于这样的幻想。请问，他们的米涅瓦、阿波罗和朱诺是什么？它们是不同种类的魔鬼。他们中有一位国王，认为值得为被当作与神明平等而死。但这里的人并非如此：恰好相反。请听他们在跛子痊愈时的说话：\BibleQuote{以色列人哪，为什么因这事惊奇？为什么注视我们，好像是我们因自己的能力和虔诚使他行走呢？}\BibleRef{宗 3:12} \BibleQuote{我们也是和你们有同样性情的人。}\BibleRef{宗 14:14} 但那些人却充满了自高自大，充满了夸耀；一切都是为了人的尊荣，没有一丝为了真理和德行的纯洁之爱（\OriginalTerm{为了哲学}{\GreekText{φιλοσοφίας} \GreekText{ἔνεκεν}}）。因为凡是为荣耀而做的事，都是无价值的。因为即使一个人拥有一切，但如果他不能控制这个（欲望），他就失去了真正哲学的一切资格，他就被更暴虐、更可耻的情欲所束缚。轻视荣耀：这足以教导一切善，并将一切有害的情欲逐出灵魂。因此，我劝你们尽最大的努力，将这种情欲连根拔除；没有别的方法能使你得到天主的赞赏，并能吸引那永不睡眠的眼睛的慈悲注视。因此，让我们竭尽全力，以获得那属天影响的享受，从而既逃避当前邪恶的考验，又达到未来的福乐，赖我们的主耶稣基督的恩宠和慈爱，祂偕同圣父及圣神，永受荣耀、权能、尊崇，从现在直到永远，万世无疆。亚孟。\par

\SourceRecord{Translated by J. Walker, J. Sheppard and H. Browne, and revised by George B. Stevens.；Nicene and Post-Nicene Fathers, First Series；Vol. 11.；Edited by Philip Schaff.；Buffalo, NY: Christian Literature Publishing Co.,；1889.；<http://www.newadvent.org/fathers/210104.htm>.}

% structure: p0001 source=h1 target=section reason=page-title

\section{论宗徒大事录第五篇讲道}

\BibleRef{宗 2:14}\par

\Quote{你们犹太人和一切住在耶路撒冷的人，这事你们要知道，要听我的话。}\par

[\Quote{你们犹太人和一切住在耶路撒冷的人，}] 上文作者称他们为陌生人。在此，他转向那些嘲讽者发言，似乎在与那些人理论，实则在纠正这些人。因为天主安排\Quote{有些人嘲讽}，使他有一个辩护的起点，并藉此辩护进行教导。[\Quote{以及一切住在耶路撒冷的人。}] 似乎他们也认为住在耶路撒冷是极高的赞誉。\Quote{这事，}他说，\Quote{你们要知道，并要听我的话。}起初，他使他们更愿意听他。他说：\Quote{因为这些人不是像你们所想的那样喝醉了。}你注意到他辩护的温和吗？\BibleRef{宗 2:15} 尽管多数人站在他一边，他却温和地与那些人辩论；先消除恶意的揣测，再确立自己的辩护。因此，他没有说\Quote{像你们嘲笑}，或\Quote{像你们讥讽}，而是说\Quote{像你们所想}，想使人看出他们并非当真这样说，目前只责备他们无知而非恶意。\Quote{因为这些人不是像你们所想的那样喝醉了，因为现在只是白天第三时辰。}为何如此？难道第三时辰就不能喝醉吗？但他并不字面强调这一点，因为他们身上没有醉酒的样子；是那些人只出于嘲讽才这么说。由此我们学到：在非本质的问题上不必多费言辞。况且，接下来的话足以证明这一点：所以现在这篇讲论是对所有人共同的。但这正是由先知\Person{岳厄尔}所说的：在末日，上主天主说，我要将我的圣神倾注在一切有血肉的人身上。（16-17节。\BibleRef{岳 2:28}）至今尚未提及基督的名字或祂的应许，而是天父的应许。请看智慧：请看体谅的宽容：（\OriginalTerm{体谅}{\GreekText{συγκατάβασιν}}。）他没有立刻转而谈论与基督有关的事，即基督在受难后应许了这神恩；那确实会颠覆一切。然而，你会说，这足以证明祂的天主性。诚然，若相信（重点正是要被相信），则足够；但若不相信，就会导致他们被石头砸死。\Quote{我要将我的圣神倾注在一切有血肉的人身上。} 他甚至向他们提供了美好的希望，如果他们愿意接受。至此，他没有让人将这视为他自己和同伴的专利，否则会招人嫉妒；从而消除了所有妒意。\Quote{你们的儿子们要说预言。} 然而，他说，这成就、这殊荣不是属于你们的；恩赐已传到你们的子孙身上。他把自己和同伴称为他们的儿子，而把那些（他所训话的人）称为他和他们的父亲。\Quote{你们的青年要见异象，你们的老人要做梦；在那些日子里，我要将我的圣神倾注在我的仆人和婢女身上，他们要说预言。} 至此他表明自己和同伴蒙受了恩宠，因为他们\OriginalTerm{堪受}{\GreekText{καταξιωθέντας}}了圣神；他所训话的人则不然，因为他们钉死了主。同样，基督为平息他们的怒气，说：\Quote{你们的儿子是藉谁驱魔？}\BibleRef{玛 12:27} 他没有说\Quote{我的门徒}，因为那听起来像是奉承。同样，\Person{伯多禄}也没有说：\Quote{他们不是喝醉了，而是藉圣神发言。}而是投靠先知，在他的庇护下如此发言。至于醉酒指控，他用自己的陈述为自己辩解；而对于恩宠，他请先知作证。\Quote{我要将我的圣神倾注在一切有血肉的人身上。}[\Quote{你们的儿子们，}等等。] 有些人藉梦领受了恩宠，有些人则公开倾注。因为先知们确实藉异象看见并领受了启示。\par

接着他继续讲预言，其中也包含可怕的内容。\Quote{我要在天上显示奇事，在地上显示征兆}[\Quote{在下地}]。\BibleRef{宗 2:19} 这些话既说到将来的审判，也说到耶路撒冷的陷落。\Quote{血、火、和烟雾的蒸气。} 请看他是如何描述攻城的。\Quote{太阳要变为黑暗，月亮要变为血。}\BibleRef{宗 2:20} 这源于受害者的\OriginalTerm{内在感受}{\GreekText{διαθέσεως}}。据说，天上确实出现了许多这样的现象，如\Person{若瑟夫}所证实的。同时，宗徒使他们感到恐惧，提醒他们近期发生的黑暗，并引导他们期待未来的事。\Quote{在主伟大而显赫的日子来到之前。} 他意即：不要因为目前你们犯罪而免受惩罚就自信。因为这些都是那伟大而可怕之日的预兆。你看到他如何使他们的灵魂战栗消融，将他们的嘲笑转为求饶？如果这些是那日子的预兆，那么极端的危险即将来临。但接下来呢？他又让他们喘口气，补充说：\Quote{将来，凡呼号主名的人，必得救。}\BibleRef{罗 10:13} 这是关于基督说的，正如\Person{保禄}所肯定的，但\Person{伯多禄}还不敢透露这一点。\par

好，让我们再回顾一下刚才说的话。他处理得很巧妙，当人们嘲笑讥讽时，他站起来开始说：\Quote{这件事你们都当知道，也当听我的话。} 但他开头说：\Quote{犹太人们。} 他用 \OriginalTerm{犹太人}{\GreekText{᾿Ιουδαῖοι}} 这个说法，我认为他是要指那些住在犹太地的人。——如果你愿意，让我们比较一下福音中的那些话，好让你看看伯多禄发生了多么突然的变化。经上写着：\Quote{有一个使女出来对他说：你也是同纳匝肋人耶稣一起的。} 而他说：\Quote{我不认识那个人。} 再被问时，\Quote{他就开始诅咒发誓。} \BibleRef{玛 26:69} 但请看这里他的胆量和极大的言论自由。——他没有称赞那些说\Quote{我们听见他们用我们的语言讲论天主的奇事}的人；相反，他通过严厉对待那些人，使这些人更加热切，同时他的言论毫无谄媚之嫌。还应该注意，无论宗徒们怎样俯就听众（\OriginalTerm{俯就}{\GreekText{συγκατάβασις}}），他们的语言既无谄媚也无傲慢之态：这是很难做到的。\par

现在，这些事发生在\Quote{第三时辰}，并非没有原因。因为当人们不工作、不是吃饭的时候，当大白天，所有人都在市场的时候，这火的光明显现了。你看到他的言论充满自由吗？\Quote{当听我的话。} 他没有添加什么，而是说：\Quote{这事，}他说，\Quote{就是由岳厄尔先知所说的：在末日将要发生。} 事实上，他表明结局已经临近，而\Quote{在末日}这句话带有强调。[ \Quote{我要倾注，等等。} ] 然后，为了不使人以为这特权只限于儿子们，他补充说：\Quote{你们的老年人要梦见梦境。} 请注意顺序：先是儿子；正如达味所说：\Quote{你的儿子们代替了你的父亲们。} \BibleRef{咏 44:17} 又玛拉基亚说：\Quote{他必使父亲的心转向儿女。又在我的仆人和使女身上。} \BibleRef{拉 4:6} 这也是卓越的标记，因为我们成为了祂的仆人，因从罪恶中获得了释放。这恩赐是伟大的，因为恩宠也临到另一性别，不像古时只限于一两个人，如德波辣和胡耳达。他没有说这是圣神，也没有解释先知的话；他只是引进预言使其自行争战。此外，他还没有提到犹达斯；然而所有人已经知道他所遭受的报应和惩罚；因为没有比用预言与他们辩论更有力的；这甚至比事实更有力。因为当基督行奇迹时，他们常常反驳祂。但当基督引用先知的话说：\Quote{上主对我主说：你坐在我右边，}他们就沉默了，\Quote{没有人，}经上记载，\Quote{能回答祂一句话。} \BibleRef{咏 89:1} 而且祂自己也时常引用圣经；例如：\Quote{如果祂称那些承受天主圣言的人为神。} \BibleRef{若 10:35} 在许多地方都可以看到这一点。因此，这里伯多禄也说：\Quote{我要将我的神倾注在一切有血肉的人身上；}即也倾注在外邦人身上。但他尚未揭示这一点，也没有解释；事实上，不解释更好（正如这句隐晦的话：\Quote{我要在天上显奇事，}因其隐晦而使他们更加恐惧。）如果一开始就解释，反而会更冒犯人。此外，即使明白，他也略过，希望他们只当做如此。但后来，当他论及复活，并以其讲道铺平道路后，他立刻给他们解释（\Emph{以下} 39节）。因为美好的事物不足以吸引他们，[所以补充说：\Quote{我要显奇事，等等。} ] 然而这事从未实现。因为那时[在以前的审判中]无人逃脱，但如今在维司巴西安时代，信徒们逃脱了。这正是主所说的：\Quote{除非那些日子被缩短，凡有血肉的都不会得救。} —— [ \Quote{血、火和烟柱。} ] \BibleRef{玛 24:22} 最坏的先来：即居民被掳，然后城市被夷平烧毁。然后他停留在比喻上，将倾覆和掳掠的情景摆在听众眼前。\Quote{太阳要变为黑暗，月亮要变为血。} 月亮变为血是什么意思？它表示杀戮的过度。这话充满了无力的惊惶。（\Emph{以上} 32页。）\Quote{将来，凡呼求上主之名的人，都必得救。} 他说：\Quote{凡呼求的，} 无论他是司铎（但他尚未揭示含义），无论为奴的，无论自主的。因为在基督耶稣内，没有男也没有女，没有为奴的，也没有自主的。\BibleRef{迦 3:28} 这很好，因为这一切不过是影子。因为在君王的宫殿里，没有高贵或低贱之分，各人按其行为显现；在艺术中，各人由其作品显明；更何况在那智慧的学校（\OriginalTerm{哲学}{\GreekText{φιλοσοφια}}）里。\Quote{凡呼求的。} 呼求：不是随便的，因为经上记着：\Quote{不是凡向我说\InnerQuote{主啊，主啊}的人，} 而是怀着（\OriginalTerm{心态}{\GreekText{διαθέσεως}}）内在的真诚情感，带着超乎寻常的善良生活，带着应有的信心。然而，他使讲道变得轻松，通过引入有关信德的内容，以及有关惩罚的可怕内容。因为得救在于呼求。\par

请问，您说的是什么意思？您在十字架之后还为他们谈论救恩吗？请容忍我一下。天主的仁慈何其伟大。而这件事本身，不亚于复活，证明祂是天主，甚至不亚于祂的奇迹——即祂召唤这些人归向祂。因为超越一切的良善，尤其为天主所特有。因此祂也说：\Quote{\BibleQuote{除了惟独天主一位外，没有良善的。}}\BibleRef{路 18:19}只是我们不要把这良善当作疏忽的借口。因为祂也作为天主施行惩罚。事实上，此处所说的惩罚，正是由祂实现的，即使祂说过：\Quote{凡呼号主的名号的人，必将获救。}我指的是耶路撒冷的命运；那不堪忍受的惩罚：我将告诉你们其中一些细节，对我们在与马西昂派和许多其他异端的论战中有所裨益。因为，既然他们区分基督为良善的天主，而那个邪恶的天主〔旧约的〕为邪恶，让我们看看是谁成就了这些事。是那个邪恶的天主，为基督报仇？还是并非如此？那么，祂怎么会与祂不一致呢？但难道是良善的天主吗？不，但事实表明，圣父和圣子都做了这些事。圣父在许多地方如此行；例如，当祂在葡萄园的比喻中说：\Quote{\BibleQuote{他要凶恶地消灭那些邪恶的农夫}}\BibleRef{玛 21:41}；又在婚宴的比喻中，国王被说\Quote{派遣他的军队}\BibleRef{玛 22:7}；圣子也是如此，当祂说：\Quote{\BibleQuote{至于那些反对我作王的仇敌，把他们带到这里来，在我面前杀掉。}}\BibleRef{路 19:27} * * *。他们派人来说：我们不愿意你作我们的王。那么，你愿意听听实际发生的事吗？而且，基督自己也预言了未来的苦难，从未有过比这更可怕的事；从未有过比那时所行更残忍的事，我亲爱的！而且祂亲自预言了。因为你还想看到比这些更悲惨的事吗？* * * 用匕首刺他们！* * * 但我还要向你们讲述那令人震惊的妇人的事，那悲剧故事吗？* * *（约瑟夫，《犹太战争》VI.3.4）难道实际发生的事不是使一切苦难都黯然失色吗？但我还要向你们讲述饥荒和瘟疫吗？还可以讲述无数恐怖的事：自然法则被破坏，法律被无视，他们比野兽还要凶残。确实，这些苦难来自战争的命运；但因为天主，因为基督如此愿意。这些事实既适用于反对马西昂派，也适用于反对那些不相信有地狱的人：因为它们足以堵住他们的无耻之口。这些灾难不是比巴比伦的更严重吗？这些痛苦不是比当时的饥荒更悲惨吗？是的，因为\Quote{\BibleQuote{从未有过从创世以来的，也决不会再有。}}\BibleRef{玛 24:21}而这是基督自己的宣告。那么，你们以为，说基督赦免了他们的罪是什么意思？也许这似乎是一个普通的问题：但你们来解答吧。——在任何地方，甚至在小说中，都不可能像这里实际发生的那样。而且，如果是一个基督徒写了这段历史，这事可能会受到怀疑：但如果他是一个犹太人，一个犹太的热诚者，并且在福音之后，那么事实的意义怎么能不为众人所明瞭呢？因为你们会看到那个人，他如何处处总是赞美犹太人的事。——所以，有地狱存在，人啊！而且天主的良善。——是的，你听到这些恐怖的事是否战栗？但这些发生在今世的事，与来世将要发生的事相比，算不了什么。我再次被逼得显得严厉、不悦、冷酷。但我能做什么？我被派定做这事：正如一个严厉的教师被派定被学生憎恨一样，我们也是如此。因为，那些被国王授予某个职位的人，无论任务多么令人不悦，都去完成所委托的事；而我们，却因怕你们的责备，就要撇下所委托的任务不去完成，这岂不是怪事吗？另一个人有另一个工作。你们中许多人，其工作是显示仁慈，行仁道，对你们的恩人喜悦和满意。但对我们所行善的人，我们显得冷峻严厉，麻烦且不悦。因为我们行善，不是通过给予快乐，而是通过施加痛苦。医生也是如此：虽然他并非过分令人不悦，因为他的技艺带来的益处立时可见；我们的益处则在将来。同样，官员对混乱和叛乱者是可憎的；立法者对受他立法的人也是烦扰的。但邀请享乐的人并非如此，预备公开庆典和娱乐、使众人头戴花环的人并非如此：不，这些人是受欢迎的人，他们设宴款待整个城市，用各种表演；大量捐助，承担所有费用。因此，那些受招待的人，以欢迎和祝福的话语，以悬挂挂毯\OriginalTerm{挂毯}{\GreekText{παραπετάσματα}}、灯火辉煌、花环、树枝和华丽服装来回报这些享乐。然而，一看见医生，病人就悲伤沮丧；一看见官员，闹事者就收敛：那时没有骚乱，没有嬉戏，除非他也加入他们的行列。那么，让我们看看，哪些人对自己的城市服务最好：是那些提供这些庆典、宴席、昂贵娱乐和多种欢庆的人；还是那些制止这一切行为的人，他们带着枷锁、鞭子、刽子手、可怕的士兵和充满恐怖的声音；发号施令，使人垂头，用棍子驱散市场上的闲人。让我们看看，我说；这些是令人不悦的，那些是受欢迎的：让我们看看益处在哪里。\OriginalTerm{说吧}{\GreekText{λήλει}}那么，你们那些提供快乐的人得到了什么？一种冷淡的享受，持续到晚上，明天就消失了；放纵的欢乐，不雅的言语和放荡。这些又得到了什么？敬畏、清醒、克制的思想；理性的头脑，闲散的终结；对内心情欲的约束；在天主之后，抵御外敌的防御墙。正是通过这些，我们拥有各自的财产，但通过那些毁灭性的庆典，我们将其挥霍。确实，强盗没有入侵，但虚荣和快乐扮演了强盗的角色。每个人都看到强盗在他眼前拿走一切，却为此高兴！这是一种新的抢劫方式，诱使人们在被抢劫时高兴！在另一方面，根本没有这种事：但天主，作为共同的父，像一堵墙一样保护我们免受一切可见和不可见的〔掠夺者〕的侵害。因为祂说：\Quote{\BibleQuote{你们要当心，不要在人前施行你们的施舍。}}\BibleRef{玛 6:1}灵魂从一方面学习放纵，从另一方面学习逃避不义。因为不义不仅在于攫取不属于我们的更多财富，而且在于给肚腹超过所需的供养，把欢乐带入不当的界限，使之陷入疯狂的过度。从一方面，它学会清醒；从另一方面，学会不贞。因为不贞不仅在于与女人发生肉体的关系，甚至在于用不贞的眼睛看女人。从一方面，它学会谦逊；从另一方面，学会自负的妄自尊大。因为宗徒说：\Quote{\BibleQuote{凡事我都可行，但不全有益。}}\BibleRef{格前 6:12}从一方面，学会端庄；从另一方面，学会不雅。至于剧场里的事，我就不提了。但为了让你们看到，那甚至不是快乐，而是悲伤，请给我看一下，节日后仅一天，那些花钱设宴的人，以及那些观看表演而享乐的人：你们会看到他们全都十分沮丧，但最沮丧的是你们那位\OriginalTerm{那个人}{\GreekText{ἔ} \GreekText{κεἵνον}}花钱的有名人物。这是公平的：因为前一天，他使普通人快乐，普通人确实兴高采烈、享受，并为华丽的衣服而欢喜，但随后无法使用它，看到自己被剥夺了，就悲伤烦恼；并想成为大人物，看到自己的快乐与那人相比是多么渺小。因此，第二天，他们交换位置，现在他，大人物，得到了更大的沮丧份额。\par

如今，若在世俗事务上，娱乐伴随如此的不满，而令人不快的事却如此有益，那么在属神之事上更是如此。为什么没有人反对法律，反而人人都视之为共同的福祉？因为制定这些法令的并非来自他方的陌生人，也不是受法律约束者的敌人，而是公民自己、他们的保护者、他们的恩人；而制定法律这一举动本身，就是恩惠和善意的标志。然而，法律充满了惩罚和约束，没有不带刑罚和强制的法律。那么，那些解释这些法律的人被称为解救者、恩人和保护者，而我们谈论天主的法律时却被视为麻烦和烦扰，这难道不合理吗？当我们谈论地狱时，我们便提出那些法律；正如在世俗事务中，人们援引关于谋杀、拦路抢劫等法律，我们也援引刑法：这些法律不是人制定的，而是天主的独生子亲自制定的。祂说：\Quote{那不怜悯人的，必须受惩罚。}\BibleRef{玛 18:23}；因为比喻的含义正是如此。让那记恨的人付出最后的代价。让那无故发怒的人被投入火中。让那辱骂的人在地狱中得到应得的报应。如果你觉得这些法律陌生，不要惊讶。因为如果基督不是要制定新法律，祂为何而来？那些其他的法律对我们来说是明显的；我们知道杀人犯和奸淫者应当受惩罚。那么，如果我们只是被告知同样的事情，哪里还需要天上的导师呢？因此，祂不说：\Quote{让奸淫者受惩罚}，而是说：\Quote{凡用淫邪的目光观看的。}并且在何处、何时那人将受惩罚，祂在那里告诉我们。祂没有将祂的法律存放在精美的公共纪念碑上，也没有藏在某个看不见的地方；祂没有竖起铜柱并在上面刻字，而是为我们兴起了十二个灵魂，即宗徒们的灵魂，并借圣神将这篇铭文刻在他们的心中。我们向你们引用这篇。如果这曾向犹太人授权，使无人能借口无知，那么对我们更是如此。但若有人说：\Quote{我没有听见，因此没有罪过}，恰恰因此，他最该受罚。因为，如果没有导师，还能以此为由；但如果有导师，便不再可能。请看主如何对犹太人说话，剥夺了他们所有的借口：\BibleQuote{我若没有来，没有对他们说话，他们就没有罪；}\BibleRef{若 15:22}保禄又说：\BibleQuote{但我说：他们没有听见吗？诚然，他们的声音传遍天下。}\BibleRef{罗 10:18}因为当无人告知时，尚有借口；但当守望者坐在那里，以此为毕生事业时，便再无借口。相反，基督的意愿不仅是让我们只看这些写下的柱石，更是让我们自己成为柱石。但既然我们使自己不配承受这书写，至少让我们仰望那些柱石。正如柱石威胁他人，但本身不受惩罚，法律亦然，蒙福的宗徒们也是如此。请看：这柱石不仅立在一处，它的书写传遍普世。无论你到印度人中去，你都会听到这话；无论去西班牙，或到地极，无人未闻，除非他自己忽视。那么，不要被冒犯，而要留心所讲的话，使你能把握德行的行为，并在我们的主耶稣基督内获得永福，愿光荣、权能、尊威归于祂，偕同圣父及圣神，从今世直到永远。阿们。\par

\SourceRecord{Translated by J. Walker, J. Sheppard and H. Browne, and revised by George B. Stevens.；Nicene and Post-Nicene Fathers, First Series；Vol. 11.；Edited by Philip Schaff.；Buffalo, NY: Christian Literature Publishing Co.,；1889.；<http://www.newadvent.org/fathers/210105.htm>.}

% structure: p0001 source=h1 target=section reason=page-title

\section{宗徒大事录讲道 第六篇}

\BibleRef{宗 2:22}\par

\Quote{以色列人，请听我这些话。}\par

【\Quote{诸位以色列人}】：他使用这个称呼并非出于奉承；相反，因他曾严厉责备他们，现在稍加缓和，提醒他们想起伟大的先祖\Person{以色列}。在此他又以引言开始，免得他们激动，因为他即将明确向他们提及耶稣；因为在之前的内容中，当先知是论述主题时，他们没有理由激动；但耶稣的名字一开始就会引起反感。——他并没有说\Quote{照我吩咐的去做}，而是\Emph{听}，显得毫不强求。且看他如何避免谈论崇高的事，而从最卑微的开始：他说\Quote{耶稣}，然后立刻提到他所出的地方，是一个被轻视的地方：\Quote{纳匝肋的耶稣}。他至今尚未说任何关于他的伟大之处，甚至没有说他是一位先知：他说：\Quote{纳匝肋的耶稣，一位从天主在你们中间所证明的人。}请看他所说\Quote{他奉天主差遣}是何等大事？因为这是他和若翰以及宗徒们始终不渝要表明的要点。请听若翰说：\BibleQuote{那曾对我说：你看见圣神降下并住在他身上的，就是这一位。}\BibleRef{若 1:33}但基督本人也极其注重这一点：\BibleQuote{我不是凭我自己来的，是他派遣了我。}\BibleRef{若 7:28}在圣经各处，这似乎都是最刻意强调的要点。因此这位蒙福团体的神圣领袖，基督的挚爱者，善牧，受托掌管天国之钥的人，领受了属神智慧的人，首先以恐惧慑服了犹太人，并显示了赐予门徒们的何等伟大恩典，以及他们具有何种被相信的权利，然后才首次开始谈论他。想想看，在凶手中间说他已经复活，这是何等胆量！然而他并没有立刻说\Quote{他复活了}，而是说什么？他说：\Quote{他来自天主：这由那些征兆——}他还没有说是耶稣本人行的，而是说什么？\Quote{是天主藉着他在你们中间所行的。}他叫他们自己作证人。\Quote{一位从天主在你们中间所证明的人，藉着奇迹、奇事和征兆，就是天主藉着他在你们中间所行的，正如你们自己所知道的。}接着，当他谈到他们亵渎的暴行时，请看他如何竭力为他们开脱罪行：他说：\BibleQuote{他，按天主的定旨和预知被交出。}\BibleRef{宗 2:23}【然而又补充说：】\BibleQuote{你们藉着不法的手把他钉在十字架上杀了。}因为虽然这是预定的，但他们仍是杀人犯。\Quote{按天主的定旨和预知}——他几乎用了若瑟同样的言语，正如若瑟对他的兄弟们说：\BibleQuote{你们在路上不要彼此生气：是天主派遣我到这里来。}\BibleRef{创 45:5}这是天主的作为。\Quote{那么我们呢？}（可能会有人说）\Quote{我们做得好。}为了避免他们这样说，所以他补充说：\Quote{你们藉着不法的手把他钉在十字架上杀了。}这里他暗示了犹达斯；同时向他们表明，这并非出于他们的力量，若他自己不允许，也不会发生：是天主把他交出来。他把全部罪恶转移到犹达斯头上，他当时已离开他们；因为正是他用亲吻把他交在他们手中。或者说，\Quote{藉着不法的手}指的是士兵们：因为他不是说简单的\Quote{你们杀了他}，而是说\Quote{你们藉着不法的人做了这事。}并且请看他如何处处强调必须先宣认苦难。他说：\BibleQuote{天主使他复活了。}\BibleRef{宗 2:24}这是大事；请看他如何把它放在讲论的中间：因为先前的事都已承认；包括奇迹、征兆和杀害——他说：\BibleQuote{天主使他复活了，解除了死亡的痛苦，因为他不应被死亡拘禁。}他在这暗示了伟大而崇高的事。因为\Quote{他不应}这个说法本身就是一种指派。它表明死亡在拘禁他时经历了产痛般的痛苦，极为困苦：而旧约中的\Quote{死亡的痛苦}（或\Quote{产痛}）指的是危险和灾难：并且他复活后永不再死。因为\Quote{他不应被死亡拘禁}这句话意味着他的复活与众人不同。然而，在他们尚未领会其意时，他便引到达味——一个消除一切人类推理的权威。\Quote{因为达味指着他说的。}并且请看看，这见证又多么卑微。所以他开始引证时，从更卑微的事情开始：因此死亡并不在痛苦之列，他说：\BibleQuote{我常将主置于我眼前，因他在我右边，使我不致动摇。}\BibleRef{宗 2:25}以及\BibleQuote{你绝不让我的灵魂留在阴府。}然后，在结束先知引文后，他补充说：\Quote{诸位弟兄。}\BibleRef{宗 2:29}当他准备说重大之事时，便用这种开场白来激发和安抚他们。他说：\Quote{请允许我放胆对你们说有关圣祖达味的话。}这是何等谦卑，因为他并未伤害谁，听众也没有理由生气。他没有说：这话不是指着达味，而是指着基督说的。但从另一角度看：他用对蒙福达味的尊敬言辞使他们敬畏；将一件公认的事说成似乎大胆之言，因而请求他们原谅他说了这话。于是他的说法不是简单说\Quote{关于达味}，而是说\Quote{关于圣祖达味，他死了且被埋葬了}；他也没有说\Quote{他没有复活}，而是以另一种方式（虽说这也非什么大事）说\Quote{他的坟墓直到今日还在我们这里}——他说了同样的事。然后——即使如此他仍未提及基督，而是接着做什么？——他继续称颂达味：\BibleQuote{他既是先知，明知天主曾以誓词向他起誓。}\BibleRef{宗 2:30}他说这话，是为了让他们因达味所受的尊荣以及出于他的后裔，而接受关于基督复活的话，因为他们看到，若事实并非如此，那便是对预言的损害，并有损于\OriginalTerm{对他们的尊荣}{\GreekText{τἥς} \GreekText{εἰς} \GreekText{αὐτοὺς} \GreekText{τιμἥς}}。\Quote{他知道}，他说，\Quote{天主曾以誓词向他起誓}——他并没有简单说应许——\Quote{要从他腰间的后裔，按肉身兴起基督，使他坐在他的宝座上。}请看他又如何仅仅暗示了崇高的事。因为现在他用言语安抚了他们之后，便自信地补充说：先知\BibleQuote{论他的复活说：他的灵魂没有被留在阴府，他的肉身也没有见到腐朽。}\BibleRef{宗 2:31}这又是奇妙的：这表明他的复活不像其他人。因为死亡虽抓住了他，却未能当场完成它的工作。——至于罪，他隐晦而模糊地提及；对于惩罚，他没加任何话；但他们杀了他，这一点他说得明白；至于其余，他来到天主所给的记号。一旦证明被杀的他，是义人，是天主所爱的，那么，就算你闭口不提惩罚，犯罪的人也必定会更多地定自己的罪，远超过你能定他的罪。因此，他把一切归于圣父，是为了让他们接受所说的话；而\Quote{不应}这一断言，他从预言中引用。那么，让我们再次审视以上所说。\par

\BibleQuote{纳匝肋人耶稣，一位由天主向你们证明（是被派遣）的人。}\BibleRef{宗 2:22}：一位，因祂的作为，无可怀疑；相反地，祂是被显明的。尼苛德摩也这样说：\BibleQuote{没有人能行你所行的这些奇迹——借着奇迹、奇事和记号，就是天主藉着祂在你们中间所行的。}\BibleRef{若 3:2}：不是隐秘的。从他所面对的听众所熟知的事实出发，然后才论及隐秘的事。接着【在说：\BibleQuote{借着天主的旨意和预知}】\BibleRef{宗 2:23}，他显示并非因为他们有能力行这事，而是事件中有智慧和天主的安排，因为它是出于天主。他迅速掠过不愉快的部分，【加上：\BibleQuote{天主使祂复活了等等}】\BibleRef{宗 2:24}。因为向他们显示祂曾经死亡，总是非常重要的一点。他说，纵然你否认，\OriginalTerm{那些人}{\GreekText{ἐκεῖνοι}}（在场）将为此作证。【\BibleQuote{解除了死亡的痛苦。}】那使死亡痛苦的，更能使那些钉死祂的人忧伤；然而，这里并没有说祂有能力杀死你们。同时，我们也当如此学习持守。因为一个在痛苦中的人，如同生产的妇人，不能紧握所握之物，不主动而被动，并急于摆脱它。并且说得很好：\Quote{因为达味论及祂说}\BibleRef{宗 2:25}；免得你将这话归于先知。——【\Quote{所以，他既是先知，又知道等等}\BibleRef{宗 2:30-31}】你注意到他如今如何解释预言，而不让它毫无注释吗？祂如何\Quote{使祂坐在}达味的\Quote{宝座上}？因为按照圣神的国度是在天上。注意他如何连同复活，也借着祂再次复活的事实宣告了国度。他显示先知是受约束的：因为预言是关于祂的。为什么他说的不是\Quote{关于祂的国度}（那是大事），而是\Quote{关于祂的复活}？祂又如何使祂坐在他的（达味的）宝座上？因为祂也作为君王统治犹太人，是的，更统治那些钉死祂的人。\BibleQuote{因为祂的肉身没有见到腐朽。}\BibleRef{宗 2:31}这似乎不如复活，但其实是同一件事。\par

\Quote{这位耶稣}——请注意他如何不另称祂——\Quote{天主已使祂复活；我们都是此事的见证人。他既被高举到天主的右边}（33-34节）：再次退避到圣父那里，尽管仅提上文就足够了；但他知道这是一个多么重要的论点。此处他也暗示了升天，并表明基督在天堂：但他并未公开这样说。\Quote{祂既领受了}，他说，\Quote{圣神的恩许。}请注意，在他讲论的开端，他并未说耶稣亲自派遣了圣神，而是说圣父；但现在，既然他已经提到了祂的征兆和犹太人向祂所行的事，并讲述了祂的复活，他便大胆地引入关于这些事的言论，再次以两种感官引证他们自己为证人：[\Quote{祂已倾注了这，就是你们所见所闻的。}]至于复活，他不断提及，但对他们残暴的行为，他只说了一次。\Quote{祂既领受了圣神的恩许。}这又是伟大的。\Quote{恩许，}他说；因为是在祂受难之前[应许的]。请注意他现在如何将这一切归于祂[\Quote{祂倾注了这}]，暗中提出了一个重要的论点。因为如果是祂倾注了它，那么先知在上文所说的就是关于祂：在末日，我将把我的神倾注在我的仆人和我的婢女身上，并在天上的高处显奇事。\BibleRef{宗 2:17} 请注意他暗中包含的意思！但接着，因为这是一件伟大的事，他又用\Quote{祂从圣父领受了}这一说法掩饰了它。他谈到了应验的善事、征兆；他说祂是君王，这触及了他们；说赐予圣神的是祂。(\WorkTitle{修辞学} 1.3.)（因为，无论一个人说多少话，如果最终不能产生某种利益，他就等于白说。）正如若翰所说：\Quote{那一位}，他说，\Quote{要用圣神给你们授洗。}\BibleRef{玛 3:11} 这表明十字架不仅没有使祂变小，反而使祂更显光荣，因为天主从古时就应许了它给祂，但现在已赐予了它。或者[也许]是祂应许给我们的\Quote{恩许}。祂如此预知了将要发生的事，并且在复活后将它更丰盛地赐给了我们。并且，他说\Quote{倾注了它}，并不要求配得；不是简单地赐予，而是丰盛地赐予。这从何显现呢？从此以后，在提到祂赐予圣神之后，他也自信地谈论祂升天的事；不仅如此，还再次引证证人，并提醒他们关于基督曾经谈到的那一位。\BibleRef{玛 22:43} \Quote{因为达味并没有}，他说，升到天上。\BibleRef{宗 2:34} 在这里，他不再用卑微的言辞说话，而是因着所说的话而有了信心；他也不说\Quote{请允许我发言}之类的话：\Quote{而是他自己说：主对我主说：你坐在我右边，直到我使你的仇敌作你的脚凳。}如果祂是达味的主，那么他们更不会藐视祂。\Quote{你坐在我右边；}他把整个问题摆在这里；\Quote{直到我使你的仇敌作你的脚凳：}这里他也给他们带来了巨大的恐惧，正如开头他显示了祂对朋友做什么、对仇敌做什么。此外，关于制伏的行为，为避免引起不信，他将此归于圣父。既然他讲了这些伟大的事，他又将话题降回到卑微的事上。\Quote{因此，}他说，以色列全家应确切知道：也就是说，你们不要质问，也不要怀疑；接着又以命令的语气说：\Quote{天主已使祂成为主}——这是他根据达味说的——\Quote{和基督，}\BibleRef{宗 2:36}，这是根据圣咏说的：因为如果正确结论应该是\Quote{因此以色列全家应确切知道祂坐在天主的右边}，这本是一件伟大的事，他却克制不提，而是引入另一件更为卑微的事，并且用了\Quote{使祂成为}这个说法，即\Quote{指定}；所以这里丝毫不涉及（\OriginalTerm{本质实体}{\GreekText{οὐσίωσις}}）实体相通的问题，而是指所提到的事。\Quote{就是这位你们钉在十字架上的耶稣。}他以这句话收尾，做得很好，从而搅动他们的心。因为当他显示了这事多么伟大时，他便揭露了他们的胆大妄为，以显明其更甚，并使他们充满了恐惧。因为人们受恩惠的吸引不如受恐惧的惩戒。\par

但那些可敬而伟大的、蒙天主所爱的人，并不需要这些动机：有些人，例如\Person{保禄}，既不关心天国，也不关心地狱。因为这才是真正爱基督，这才是不要做佣工，也不要把这事看作买卖和交易，而是要真正有德，并为爱天主而行一切。\BibleRef{罗 9:3} 当我们欠了如此大的债，却甚至连像商人那样寻求天国都不肯时，这是何等令人流泪的事！祂应许我们如此伟大的事物，我们却连听祂都不配？还有什么比这种敌意更甚的呢？然而，他们却疯狂追求赚钱，即使面对敌人，即使面对奴隶，即使面对对他们最敌对的人，即使面对极恶的人，只要他们指望能藉他们赚钱，他们就会做一切事，阿谀奉承、卑躬屈膝，使自己成为奴隶，并把他们看得比任何人都更可敬，以便从中获利：因为对金钱的希望不容他们考虑任何这样的念头。但天国并不像金钱那样有力；不，它远不及金钱那样有力。因为应许者并非凡人：但比天国本身更大的，是我们从这样一位施予者领受它！然而现在的情况是，就好像一位国王，在赐予一万种其他恩惠之后，愿意使我们作他的继承人、并与他的儿子同为继承人，却被轻视；而某个强盗头子，对我们和我们的父母做了一万种恶事，自身充满一万种邪恶，且完全毁坏了我们的荣誉和福祉，却因施舍一分钱而受到我们的崇拜。天主应许天国，却被轻视；魔鬼帮助我们下地狱，却受到尊崇！这里，是天主；那里，是魔鬼。但让我们看看所吩咐的任务有何不同。因为假如没有这些考虑：假如不是这里天主、那里魔鬼；不是这里帮助进天国、那里下地狱：所吩咐的任务本身的性质就足以诱导我们服从前者。因为每一条吩咐了什么？一条吩咐的是光荣之事；另一条吩咐的是羞耻之事：一条吩咐那些会带来一万种灾祸和耻辱的事；另一条吩咐那些本身带有极大安息的事。因为看：一条说：\Quote{你们跟我学吧，因为我是良善心谦的，这样你们必得灵魂的安息。}\BibleRef{玛 11:29}；另一条却说：你要凶残、不温和、易怒、暴躁，比人更像是野兽。让我们看看哪一条更有益，哪一条，请问，更有利？\Quote{不要谈这个，}你说。……但你要考虑他是魔鬼：首先确实，如果那被显明：也需要承受劳苦，另一方面，胜利的奖赏会更大。因为吩咐容易之事的并非恩人（\OriginalTerm{关心者}{\GreekText{κηδεμὼν}}），而是吩咐对我们有益之事的人。因为父亲们也吩咐令人不快的事；但正因如此他们才是父亲：同样，主人对奴隶也是如此；但绑架者和\OriginalTerm{破坏者}{\GreekText{λυμεώνες}}则相反。然而，基督的命令伴随着快乐，从那句话可以明显看出。因为你认为易怒的人属于哪一类，而宽容温和的人属于哪一类？那一个人（\OriginalTerm{那人}{\GreekText{ἐκείνου}}）的灵魂难道不是似乎处在一种孤寂的退隐中，享受极大的宁静；而另一个人（\OriginalTerm{这人}{\GreekText{τούτου}}）的灵魂则像是市场、喧嚣和城市中心，那里有出门者的大声喧哗、骆驼、骡子、驴的噪音，人们向遇见的人高声喊叫以免被踩踏，还有银匠、铜匠，人们推推搡搡，有的被压倒，有的压倒别人？但前者（\OriginalTerm{这人}{\GreekText{τούτου}}）的灵魂就像某个山顶，空气清新，阳光纯净，泉水清澈涌流，繁花似锦，春天的草地和花园披上灌木和花卉的羽毛，波光粼粼；如果那里听到什么声音，那也是甜美的，能引起耳朵极大的愉悦。因为要么唱歌的鸟儿栖息在枝繁叶茂的树梢上，蝉、夜莺和燕子交织在一起，奏出一种和弦；要么微风轻轻摇动树叶，从松树和冷杉中吹出哨音，常常像天鹅的歌声；玫瑰、紫罗兰和其他花朵轻轻摇曳，暗起涟漪（\OriginalTerm{暗波}{\GreekText{κυανίζοντα}}），如同大海被轻柔的微波轻轻拂过。不，人们可以找到许多意象。这样，当人们看到玫瑰时，会想象在其中看到彩虹；在紫罗兰中，看到起伏的大海；在百合中，看到天空。但这样的人不仅通过景象和观看带来愉悦，而且对于那个望向草地的人的身体，它更使他焕发精神，自由呼吸，以至于他觉得自己在天上多于在地上。还有一种不同的声音，当水从山崖上凭借自身力量穿过沟壑，在卵石床上轻轻溅起，发出催眠的声响，从而以愉悦的感觉放松我们的身体，很快使我们的眼睛垂下柔软的倦意。你愉快地听了描述：也许它也使你爱上了独处。但比这独处更甜美的是忍耐的灵魂……因为我们提出这个比喻并不是为了描述一片草地，也不是为了炫耀言辞；而是为了让你看到忍耐的喜乐何其大，以及与一个忍耐的人交谈会远比常去这种地方更令人愉悦和受益，从而追求这样的人。因为当这样一个灵魂连一丝暴力也不发出，只有温和宜人的话语时，那么微风的那种温柔柔和便得到了对应；还有不带一切傲慢的恳求，形成与那些有翼歌者的相似——这岂不更好得多？因为身体并非被言语的柔风所扇动；不，它使我们灼热发烫的灵魂得到清凉。一个医生即使极其用心，也不能像一个有耐心的人那样，用他自己话语的气息，使一个暴躁易怒、怒火中烧的人迅速冷却。我为什么提到医生？甚至连烧红的铁浸入水中，也不会像暴躁的人遇到忍耐的灵魂那样快地失去热量。但是，如果唱歌的鸟偶然飞进市场，它们在那里就毫无价值，正如我们的教训落入易怒之人的灵魂时一样。确实，温柔比苦涩和乖戾更甜美。——然而，一个是天主的命令，另一个是魔鬼的。你是否看到，我并非无缘故地说，即使没有魔鬼或天主的事，所吩咐的事本身就足以使我们（\OriginalTerm{疏远}{\GreekText{ἀποστῆσαι}}）反感？因为一条既令自己愉快，也服务他人；另一条令自己不悦，也伤害他人。没有什么比一个愤怒的人更令人不快的，没有什么比这更讨厌、更可憎、更可鄙的，正如没有什么比一个不知愤怒为何物的人更令人愉快的。与野兽同住胜过与暴躁的人同住。因为野兽一旦被驯服，就遵守它的规律；但人，无论你驯服他多少次，又会变野，除非他自己养成某种（温和的）习惯。\par

因为正如晴朗明媚的日子与阴郁沉闷的冬天相对，愤怒之人的灵魂与温和之人的灵魂也是如此。不过，我们现在不要看对他人造成的恶果，而要关注那些影响到自身的后果：尽管给他人带来伤害对自己也绝非轻伤，但此刻，暂且考虑这一点。有哪个刽子手能用鞭子这样撕裂肋骨，有哪种\OriginalTerm{红热的柳叶刀}{\GreekText{ὀ} \GreekText{βελίσκοι}}这样刺穿身体，有哪种疯狂能像愤怒与暴怒那样使人丧失本性理智？我知道许多例子，有人因放纵愤怒而致病：最严重的发烧正是这些。但如果它们如此伤害身体，想想灵魂吧。因为你不可争辩说你没有看见这损害，反而要思考：如果那接纳恶毒激情的器官受害如此，那么产生它的那部分又该遭受怎样的损害！许多人失明了，许多人患上了最严重的疾病。然而，那勇敢忍受的人，将轻易承受一切。但是，尽管如此，这些正是魔鬼命令的麻烦任务，而他为我们这些任务所定的工价就是地狱。他既是魔鬼，又是我们救恩的仇敌，而我们对他的命令的服从甚于对基督的服从，而基督是救主、恩主、护慰者，祂所说的话既更甘饴、更可敬、更有益、更有助，对我们自己和我们的同伴都是最大的祝福。亲爱的，没有什么比愤怒更糟，没有什么比不合时宜的怒气更糟。它不会长久拖延；它是一种迅速、尖锐的激情。许多次，一句气话脱口而出，却需要一生来医治；一个行为做出，却毁了人的一生。因为最糟的是，在片刻之间，因一个行为、一句言语，它常常将我们从永善的拥有中驱逐出去，并使无数的辛劳化为乌有。因此，我恳求你们尽力遏制这头野兽。至此，我已经讲了关于温和与忿怒；如果有人着手处理其他对立面，例如贪婪与对荣耀的疯狂激情，与轻视财富和荣耀相对；无节制与节制相对；嫉妒与仁爱相对；并将它们一一排列对比，那么人就会知道这差别有多大。看，从所命令的这些事中清楚显示：一位主人是天主，另一位是魔鬼！那么，让我们遵行天主的命令，不要将自己投入无底的深渊；趁着还有时间，让我们洗净一切玷污灵魂的东西，好使我们能获得永恒的祝福，借着我们的主耶稣基督的恩宠和慈爱，愿光荣、权能、尊崇归于祂及圣父和圣神，从现在直到永远，永无穷尽。阿们。\par

\SourceRecord{Translated by J. Walker, J. Sheppard and H. Browne, and revised by George B. Stevens.；Nicene and Post-Nicene Fathers, First Series；Vol. 11.；Edited by Philip Schaff.；Buffalo, NY: Christian Literature Publishing Co.,；1889.；<http://www.newadvent.org/fathers/210106.htm>.}

% structure: p0001 source=h1 target=section reason=page-title

\section{宗徒大事录讲道 第七讲}

\BibleRef{宗 2:37}\par

\BibleQuote{他们听见这些话（英文译本作\InnerQuote{这事}），心中刺痛，便向伯多禄和其他宗徒说：\InnerQuote{诸位兄弟，我们该做什么？}}\par

你看见温和是多么伟大的事吗？它比任何激烈更能刺透我们的心，造成更深的创伤。因为如同在身体变得麻木的情况下，击打它的人不能如此有力地影响感觉，但若他先使之柔软，使之变得敏感，然后就能有效地刺入；同样，在此情形中，也必须先使之柔软。然而能使柔软的，不是愤怒，不是严厉的指责，不是人身攻击；而是温和。前者反而加剧麻木，唯有后者能去除它。因此，你若想责备任何有过失的人，就要以最大的温柔接近他。请看这里：他温和地提醒他们所犯的暴行，不加任何评论；他宣示天主的恩赐，继而言及为这事作证的恩宠，因而将他的言论延伸至更大的长度。他们因此敬畏伯多禄的温和，因为他向那些曾钉死他的主、并对他和同伴满心杀机的人说话，却以慈父和导师的身份向他们讲论。他们不仅被说服，甚至自责，意识到过去的行为。因为他没有给他们的愤怒留有余地，以蒙蔽他们的判断，而是藉着谦卑，如同驱散了他们愤怒的雾霭与黑暗，然后向他们指出他们所做的胆大妄为之举。事实正是如此：当我们说自己受了伤害时，对方便试图证明他们并未伤害；但当我们说我们未受伤害，反而做错了时，对方就会采取相反的态度。因此，你若想使你的敌人处于\OriginalTerm{在争辩中}{\GreekText{εἰς} \GreekText{ἀγώνα}}的过错中，要谨慎指责他；相反，要为他辩护\OriginalTerm{争辩}{\GreekText{ἀγώνισαι}}，他必定会发现自己有罪。人天生有一种对抗精神。伯多禄的行为正是如此。他没有严厉地指责他们；相反，他几乎竭尽全力为他们辩护。这正是他能深入他们灵魂的原因。你会问，他们被刺透的证明在哪里？在他们自己的话中；他们说什么？\Quote{诸位兄弟，我们该做什么？}他们曾称之为欺骗者的人，如今称为\Quote{兄弟}：并非因此自视与他们平等，而是为了获得他们的兄弟情谊和善意；此外，也因为宗徒们屈尊以此称呼他们。他们又说：\Quote{我们该做什么？}他们没有立刻说\Quote{那么我们悔改吧}，而是把自己交托给门徒。就像一个在沉船边缘的人，看见舵手，或在病中看见医生，就把一切交在他手中，凡事听从；同样，这些人也承认自己处于极端危险中，毫无得救的希望。他们没有说\Quote{我们如何得救？}而是说\Quote{我们该做什么？}这里，尽管问题是对所有人提出的，但回答的仍是伯多禄。他说：\BibleQuote{你们悔改，并要因耶稣基督的名受洗。}\BibleRef{宗 2:38}他还没有说\Quote{相信}，而是\BibleQuote{你们每个人都当受洗。}因为他们在洗礼中领受了这个。然后他谈到益处：\BibleQuote{为得罪之赦，并将领受圣神的恩赐。}如果你要领受一个恩赐，如果洗礼带来赦免，为何耽延？接着他使他的讲话更具说服力，补充说：\BibleQuote{因为这恩许是给你们。}\BibleRef{宗 2:39}因为他之前提到了一个恩许。他说：\BibleQuote{和给你们的子女。}当这些人成为祝福的继承者时，恩赐更大。他继续：\BibleQuote{并给一切远方的人，}如果给远方的人，更给近处的你们：\BibleQuote{凡主我们的天主所召叫的。}请注意他选择说\BibleQuote{给远方的人}的时机。那是在他发现他们被安抚并自责的时候。因为当灵魂宣判自己有罪时，它就不再能嫉妒。\BibleQuote{他还用许多别的话作证并劝勉说。}\BibleRef{宗 2:40}请注意写作者始终追求简练，多么不慕虚荣和炫耀。\BibleQuote{他作证并劝勉说。}这是教导的完美，包含了某种畏惧和某种爱。\BibleQuote{你们应救自己脱离这邪恶的世代。}他什么也没有说关于未来的事，全部是关于现在，而人们主要受现在影响；他表明福音也释放人脱离现今的邪恶。\BibleQuote{于是，接受他话的人就受了洗；那一天约增加了三千人。}\BibleRef{宗 2:41}你不觉得这比神迹更使宗徒们欢欣吗？\BibleQuote{他们专心并同心合意地坚持宗徒的训诲和团契。}\BibleRef{宗 2:42}这里有两种美德：恒心和同心。他说：\BibleQuote{在宗徒的训诲中，}因为他们再次教导他们；\BibleQuote{和团契，以及擘饼和祈祷。}一切都是共同的，一切都有恒心。\BibleQuote{每人都起了敬畏之心。}\BibleRef{宗 2:43}指那些相信的人。因为他们没有像对待普通人那样轻视宗徒，也没有仅注视可见之物。实际上，他们的思想被点燃成炽热。由于伯多禄之前说了很多，并宣告了恩许和将来的事，他们满心畏惧是自然的。神迹也见证了这些话：\BibleQuote{宗徒们行了许多奇事异兆。}正如基督的情形：先有征兆，然后教训，然后奇事；现在也是如此。\BibleQuote{所有信了的人都在一处，凡物公用。}\BibleRef{宗 2:44}想想这是多么大的进步！因为团契不仅在祈祷中，也不仅在教义中，而且也在社会生活\OriginalTerm{社会生活}{\GreekText{πολιτεία}}中。\BibleQuote{并且卖去了田产和财物，分给众人，按各人的需要。}看看他们心中产生了何等敬畏！他说：\BibleQuote{他们分给众人，}显示出明智的管理\OriginalTerm{明智的管理}{\GreekText{τὸ} \GreekText{οἰκονομικὸν}}：\BibleQuote{按各人的需要。}并非像希腊某些哲学家那样随意，其中一些放弃土地，另一些将大量金钱投入大海；但这并非轻视财富，而只是愚蠢和疯狂。因为魔鬼普遍致力于贬低天主的受造物，好像不可能善用财富似的。\BibleQuote{并且每天同心合意地，}\BibleRef{宗 2:46}他们享受教导的益处。想想这些犹太人在大小事上除了专心致志地待在圣殿外，不做别的事。因为他们变得更加热忱，对那个地方也增加了虔诚。因为宗徒们暂时没有把他们从这一目标上拉开，以免伤害他们。\BibleQuote{并且挨户擘饼，以欢乐和诚实的心用饭，赞美天主，并获得了全体百姓的好感。}\BibleRef{宗 2:47}在我看来，他提到\Quote{饼}，在这里表示禁食和艰苦的生活；因为他们\Quote{用饭}，不是享用美味。他说：\Quote{以欢乐的心。}你看见不是美味，而是食物而非享乐\OriginalTerm{食物而非享乐}{\GreekText{τροφἥς} \GreekText{οὐ} \GreekText{τρυφἥς}}带来了快乐吗？因为享受美味的人处于惩罚和痛苦中；但这些并非如此。你看见伯多禄的话也包含了这一点，即生活的规范吗？[\Quote{并且一心一意。}]因为哪里没有纯朴，哪里就没有喜乐。他们如何获得\Quote{全体百姓的好感}？由于他们的施舍。请不要看大司祭出于嫉妒和恶意起来反对他们的事实，反而要考虑\Quote{他们获得了百姓的好感}。——\BibleQuote{主天天在一起\OriginalTerm{在一起}{\GreekText{ἐπὶ} \GreekText{τὸ} \GreekText{αὐτό}} [一起]把得救的人加给教会。——所有信了的人都在一处。}再一次，同心，爱德，是一切美善的原因！\par

[\BibleQuote{他们听见这话}，等等。\BibleQuote{伯多禄便对他们说}，等等。] \BibleRef{宗 2:37} 先前所说的并不足够。因为那些话确实足以引他们进入信德；但这些话是要表明信徒应当做什么。他没有说：在十字架中，而是说：\BibleQuote{你们每人要因耶稣基督的名受洗。} \BibleRef{宗 2:38} 他不常叫他们记起十字架，免得显得谴责他们，而只是简单地说：\BibleQuote{悔改}，为什么呢？是要我们受罚吗？不：\BibleQuote{你们每人要因耶稣基督的名受洗，为得罪赦。} 然而世上官府的法则却截然不同；但在福音\OriginalTerm{宣讲}{\GreekText{κηρύγματος}}的情形中，犯人一旦认罪，就被拯救了！请看伯多禄并没有立刻略过这点，而是还指明了条件，并加上：\BibleQuote{你们将领受圣神的恩惠；}这个宣告因宗徒们自己也领受了那恩惠的事实而得到印证。[\Quote{因为这恩许}等等。] \BibleRef{宗 2:39} \Quote{这恩许}即圣神的恩赐。到此为止，他谈论了容易的部分，以及附带有大恩赐的部分；然后引导他们付诸实践：因为预先尝到了如此大的祝福，将成为他们热心的基础。[\Quote{他又用许多别的话证道}等等。] \BibleRef{宗 2:40} 然而，既然听者想知道这些进一步的话语的要点和实质，他就告诉我们：[\Quote{说：你们当救自己脱离这邪恶的世代。}] [\Quote{于是，那些欣然领受他话的人}等等。] \BibleRef{宗 2:41} 他们赞同所说的，尽管其中充满恐惧，并在表示同意后，立即前去受洗。\BibleQuote{他们继续}，经上记载，\BibleQuote{恒心于宗徒的训诲}（或称\InnerQuote{教导}）\BibleRef{宗 2:42}：因为他们不是一天，也不是两三天还在受教，作为已经归入另一种生活方式的人。[\Quote{他们同心合意恒心于宗徒的训诲}等等。] 这个表述不是\OriginalTerm{一同}{\GreekText{ὁμοῦ}}\Quote{一起}，而是\OriginalTerm{同心合意}{\GreekText{ὁμοθυμαδόν}}\Quote{同心合意}；（\Quote{每天}，他说[之后]，\Quote{他们同心合意继续在圣殿中}）即一心一意。在这里，他又一次以其简洁，没有叙述所给予的教导；因为如同幼小的孩童，宗徒们用属神的食物哺养他们。\BibleQuote{众人都感到惧怕} \BibleRef{宗 2:43}：显然，也包括那些不信的人；就是说，因看到如此巨大的变化一下子完成，此外也因那些奇迹。[\Quote{凡信了的人都聚在一起，并将一切东西公用}等等。] \BibleRef{宗 2:44} 他们突然间都成为了天使；他们所有人都恒心于祈祷和听道，他们看到属神的事物是\Emph{共有的}，那里没有一个人比其他人拥有更多，他们迅速聚\OriginalTerm{在一起}{\GreekText{ἐπὶ} \GreekText{τὸ} \GreekText{αὐτό}}，归向同一共通之物，甚至分施给众人。\BibleQuote{所有信了的人} \BibleRef{宗 2:44}，经上说，都是\OriginalTerm{在一起}{\GreekText{ἐπὶ} \GreekText{τὸ} \GreekText{αὐτό}}：若要明白这并非指他们在空间上聚在一起，请看下文[\Quote{并将一切东西公用}]。\BibleQuote{所有人}经上说：没有一个人例外。这是天使般的共融体，不称任何东西为自己的。立刻，万恶之根被剪除了。藉着他们的作为，他们显明了所听见的：这就是他所说的：\BibleQuote{你们当救自己脱离这邪恶的世代。}——\BibleQuote{每天同心合意继续在圣殿中。} \BibleRef{宗 2:46} 既然他们已成为三千人，他们现在就把他们带到外面：并且圣神所赋予的胆量是大的：他们每天上到圣殿如同一个神圣的地方，正如我们常看到伯多禄和若望这样做：因为此时他们丝毫不扰乱犹太人的规条。这尊荣也传到了那地方；在家中吃饭。在什么家中？在圣殿中。请看虔诚的增长。他们抛弃了财富，并欢喜快乐，有极大的喜乐，因为他们所领受的财富是更大的，无需劳苦（\OriginalTerm{无劳苦}{\GreekText{ἄπονα}}，抄本异文作\OriginalTerm{美善}{\GreekText{ἀγαθά}}）。没有人责备，没有人嫉妒，没有人怨恨；没有骄傲，没有轻视。他们如同孩童，确实自认为处于受教之中：如同新生婴儿，他们的性情就是如此。然而为何用这模糊的形象？如果你记得天主用地震震动我们城市时，众人是如何降服的。（\Emph{参下文}，第四十一讲第二节。）那些悔改者当时的情形就是这样。那里没有狡诈，没有邪恶：这就是恐惧和苦难的效果！那时没有人谈论\Quote{我的}和\Quote{你的}。于是喜乐在餐桌旁等待他们；没有人似乎吃自己的或别人的——我承认这似乎是个谜。他们也不认为弟兄们的财产与自己无关；它是主人的财产；他们也不认为任何东西是自己的，一切都是弟兄们的。穷人不知羞耻，富人不骄傲。这就是喜乐。后者认为自己是有义务且幸运的一方；其他人感到自己在此受到尊敬，他们紧密地联合在一起。因为，当真，当人们发放钱财时，容易有侮辱、骄傲、怨恨；因此宗徒说：\BibleQuote{不要怨恨，或出于不得已。}—— \BibleRef{格后 9:7} [\Quote{怀着喜乐和心地纯朴}等等。] 请看，他为他们作见证的事有多少！真诚的信德、正直的行为、恒心于听道、于祈祷、于纯朴、于喜乐。[\Quote{赞美天主。}] \BibleRef{宗 2:47} 有两件事可能使他们沮丧：他们的俭朴生活和财产的损失。然而因这两件事他们都欢喜。[\Quote{并且获得全民众的爱戴。}] 因为谁会不爱这样性格的人，如同共同的父亲？他们彼此不怀恶意；他们把一切都交托给天主的恩宠。[\Quote{与全民众。}] 那里没有恐惧；是的，尽管他们置身于危险之中。然而，藉着\Emph{纯朴}，他指称他们的全部德行，远远超越他们对财富的轻视、他们的节制和他们在祈祷中的恒心。因为如此他们也向天主献上纯洁的赞美：这就是赞美天主。但也要注意这里他们如何立即获得赏报。\BibleQuote{获得全民众的爱戴。} 他们吸引人，大受爱戴。因为谁不珍视和钦佩他们品性的纯朴？谁不愿和一个没有暗藏诡计的人相连？救恩不属于这些人，又属于谁呢？那些伟大的奇事不属于这些人，又属于谁呢？福音不是首先宣讲给牧羊人吗？还有若瑟，一个心地纯朴的人，甚至不让怀疑通奸的恐惧使他做错事？天主不是选择了粗人、那些质朴的人吗？因为经上记载：\BibleQuote{纯朴的灵魂是有福的。} \BibleRef{箴 11:25} 又说：\BibleQuote{行事纯朴的，步履稳妥。} \BibleRef{箴 10:9} \Quote{不错，}你会说，\Quote{但也需要谨慎。} 请问，纯朴不就是谨慎吗？因为当你不怀疑恶时，你也不会制造恶；当你没有烦恼时，你也不会记恨伤害。有没有人侮辱你？你不感到痛苦。有没有人辱骂你？你毫不受伤。他嫉妒你？你仍然没有伤害。纯朴是通往真正哲学的大道。没有比纯朴的人灵魂更美丽的。因为就外貌而言，一个忧郁、垂头丧气、矜持（\OriginalTerm{忧郁}{\GreekText{σύννους}}）的人，即使相貌英俊，也失去了许多美丽；而一个放松面容、微微微笑的人，则增加了他的美貌；同样，在灵魂方面，一个矜持的人，即使有万千优点，也使之丑化；而坦诚纯朴的人则正好相反。后一种人可以安全地交朋友，发生争执时也容易和解。与这样的人相处，无需警卫和岗哨，无需锁链和脚镣；但他自己有极大的自由，与他交往的人也是如此。但你会说，这样的人若落入恶人中间，会怎样？下令我们心地纯朴的天主，会伸出祂的手。谁比达味更无伪？谁比撒乌耳更邪恶？但谁胜利了？又如若瑟的例子：他难道不是以纯朴接近主母，而她以邪术接近他？请问，他又有何损失？再者，谁比亚伯尔更纯朴？谁比加音更恶毒？若瑟再次，他对兄弟难道不是毫无心计地相处？这不就是他显赫的原因吗，因为他毫无猜疑地说出话来，而他们却以恶意接受他的话？他毫无保留地一再述说他的梦；然后他又携带食物出发去他们那里；他没有谨慎；他把一切交托给天主：不仅如此，他们越把他当敌人看待，他就越发待他们如兄弟。天主有能力不让若瑟落入他们手中；但为使奇迹显明：尽管他们作恶多端，他却比他们更高；虽然打击临到他，却是来自别人，而非来自他自己。相反，恶人先打击自己，而且是惟独他自己。\BibleQuote{因为独生，}经上说，\BibleQuote{独自担当自己的苦难。} \BibleRef{箴 9:12} 他灵魂常充满沮丧，心思常纠结：无论必须听什么或说什么，他都带着抱怨和控诉行一切事。友谊与和谐离这样的人很远，很远：那里只有争斗、敌意和一切不快。这样的人甚至怀疑自己。对他们来说，连睡眠也不甜美，其他一切也是如此。如果他们也有妻子，看，他们成为敌人，与所有人作战：无尽的嫉妒，不断的恐惧！是了，恶人，\OriginalTerm{恶人}{\GreekText{πονηρὸς}}源自\OriginalTerm{劳苦}{\GreekText{πονεἵν}}，即\Quote{受苦}。确实，圣经常称\Quote{邪恶}为劳苦；例如：\BibleQuote{舌下尽是辛劳和劳苦；}又说：\BibleQuote{在他们中间尽是辛劳和劳苦。} \BibleRef{咏 10:7}\par

如今若有人诧异，那些先前属于末等人之列的人，如今为何如此不同，就让他明白：患难是原因——患难，这天上智慧的教师，这虔诚之母。当财富被消除时，邪恶也随之消失。你说：没错，这正是我想问的；但如今所有的邪恶从何而来？那三千人和五千人怎么立刻就有了选择德行的心意，并且同时成了基督的哲学家，而如今几乎找不到一个？他们那时为何如此和睦？是什么使他们果敢而活跃？是什么突然点燃了他们？原因是，他们怀着极大的虔诚亲近；荣誉不像现在那样被追求；他们将思想转向未来的事物，不指望任何眼前的东西。这是热切心灵的标志，敢于面对危险；这是他们对基督信仰的理解。我们持不同观点，我们在此世寻求安慰。结果是，当时候来到，我们甚至连这个也得不到。那些人问：\Quote{我们该做什么？} 我们却相反——\Quote{我们该做什么？} 该做的事，他们做了。我们，恰恰相反。那些人定自己的罪，对自己得救不抱希望。正是这一点使他们成为那样。他们知道自己领受了何等恩赐。但你如何能像他们一样，当你以相反的精神行事时？他们听了，立刻受了洗。他们没有说我们现在说的那些冷漠的话，也没有制造拖延（p. 47, note 3）；然而他们已听到所有的要求；但那句\Quote{从这世代中得救}的话使他们不再懈怠；他们欣然接受了劝勉；他们用行动证明了自己欣然接受，展示了他们是怎样的人。他们立刻进入赛场，脱下了外衣；而我们，我们进入赛场，却打算穿着外衣战斗。这就是为什么我们的对手如此轻松，因为我们被自己的动作缠住，不断跌倒。我们做的正像一个人要与一个狂怒的、喷火的人搏斗，看到那职业摔跤手满身尘土、肤色黝黑、赤身露体、因沙土和太阳而满身泥垢，流着汗、油和污物；而他自己，香气扑鼻，穿着丝绸衣服、金鞋、拖到脚跟的长袍、头上戴着金饰，就这样下到竞技场去与他扭打。这样的人不仅会被阻碍，而且一心想着不弄脏或撕破他的漂亮衣服，会在第一回合就摔倒，并且还会遭受他最害怕的事——他心爱之物的损坏。比赛的时刻到了，你却说：你穿上丝绸衣服了？这是锻炼的时候，赛跑的时刻，你却在装饰自己如同准备游行？不要看外表，要看内心。因为对这些事物的思想像强韧的绳索一样从四面八方束缚灵魂，使她无法让你举起手来与敌人争战，使你变得软弱无力。即使摆脱了所有这些束缚，一个人也会认为自己很幸运才能战胜那不洁之力。为此，基督也不允许仅仅舍弃财富就足够，但祂说什么？\Quote{变卖你所有的一切，施舍给穷人，然后来跟随我。}\BibleRef{谷 10:21} 现在，即使我们抛弃了财富，仍不安全，还需进一步的技巧和勤奋实践；那么，如果我们保留财富，就更不能成就大事，反而会成为旁观者和那恶者本身的笑柄。因为即使没有魔鬼，没有人与我们角力，也有千万条路从四面八方引导贪财者下地狱。那些问\OriginalTerm{魔鬼为什么被造}{\GreekText{διατί} \GreekText{ὁ} \GreekText{δ}. \GreekText{γέγονεν}}的人如今在哪里？看，这里魔鬼没有插手，一切都是我们自己做的。那些山上的人或许有权这样说，他们证明了他们的节制、对财富的轻蔑和对这一切的漠不关心之后，无限地优先选择了舍弃父亲、房屋、田地、妻子和儿女。然而，他们是最不说这话的；而那些在任何时候都不该说这话的人，却说了。那些确实是与魔鬼的角力；这些，魔鬼认为不值得介入。你会说：但这贪欲也是魔鬼灌输的。那么，逃避它，不要容纳它，人啊。假设你看见有人从楼上倒出污物，同时有一个人看见污物被倒出，却站在那里，让污物全落在自己头上；你不仅不怜悯他，反而生气，告诉他活该；每个人都对他喊：\Quote{别傻了！}并把过错更多地归咎于他让污物落在自己身上，而不是那倒出污物的人。但现在，你知道贪欲来自魔鬼；你知道这是千万邪恶的原因；你看见他像倒出污物一样倒出他那恶臭的想象；你却看不到你正用光着的头接受他的肮脏，而只需稍微转开一点就能完全避开？正如我们的人只要改变位置就能避开一样；所以，你拒绝接受这样的想象，抵挡欲望。我怎能做到这一点？你会问。如果你是一个外邦人，只盯着眼前的事物，也许这事相当困难，但即使外邦人也做到了；而你——一个期待天堂和天上福乐的人——却问\Quote{我怎能抵挡恶念？} 如果我说相反的话，你可能怀疑：如果我说贪恋财富，\Quote{我怎能贪恋财富？}你可能会回答\Quote{既然我看到这样的事？}告诉我，如果金银宝石摆在你面前，我说渴望铅，你岂不犹豫？因为你会说：我怎能这样？但如果我说不要渴望它，这就更容易理解了。我不惊叹那些轻视财富的人，而是惊叹那些不轻视财富的人。这是极其愚笨的灵魂的特点，与苍蝇和蚊子无异，在地上爬行，在污秽中打滚，没有崇高思想。你说什么？你注定要承受永生；你说：我怎能为了未来而轻视现在的生活？什么，这些事能相提并论吗？你要接受王袍；你说：我怎能轻视这些破布？你要被领进王宫；你说：我怎能轻视这破棚屋？的确，事事都是我们自己的错，我们不愿意让自己哪怕稍微被唤醒。因为愿意的人已经成功了，而且是带着极大的热情和容易做到的。但愿你能被我们的劝勉说服，也能成功，并成为成功者的效仿者，赖我们的主耶稣基督的恩宠和怜悯，愿荣耀、权能、尊崇归于祂及圣父与圣神，从今世直到永远，世世无尽。阿们。\par

\SourceRecord{Translated by J. Walker, J. Sheppard and H. Browne, and revised by George B. Stevens.；Nicene and Post-Nicene Fathers, First Series；Vol. 11.；Edited by Philip Schaff.；Buffalo, NY: Christian Literature Publishing Co.,；1889.；<http://www.newadvent.org/fathers/210107.htm>.}

% structure: p0001 source=h1 target=section reason=page-title

\section{宗徒大事录第八篇讲道}

\BibleRef{宗 3:1}\par

\BibleQuote{伯多禄和若望一同上圣殿去，祈祷的时辰，即第九时辰。}\par

处处可见这两位宗徒同心协力。\BibleQuote{西满伯多禄向他示意。} \BibleRef{若 13:24} 他俩也一同来到坟墓。\BibleRef{若 20:3} 至于若望，伯多禄曾对基督说：\BibleQuote{这人将来怎样？} \BibleRef{若 21:21} 至于其他奇迹，本书作者略过不提；但他提及那个引起众人骚动的奇迹。再看，他们不是特意去找那些人；他们如此淡泊名利，如此效法他们的师傅。他们为何上圣殿？他们仍在像犹太人那样生活吗？不，而是为了\OriginalTerm{权宜之用}{\GreekText{χρησίμως}}。又一个奇迹发生了，既巩固了归化者，又吸引了其余的人；而这个奇迹，正是他们行的一个标记。那人的病是天生，医药无效。他瘸了四十年\BibleRef{宗 4:20}，正如作者后文所说，期间无人治好他。最难治的病是先天性的。这是大不幸，甚至连供养自己必需的食物都不可能。那人因地点和病痛都引人注目。听作者如何叙述。\BibleQuote{有一个人，从母胎中就瘸了，每天被人抬来，放在那叫做丽门的殿门口，向进殿的人求施舍。} \BibleRef{宗 3:2} 他想得到施舍，却不知来者是谁。\BibleQuote{他看见伯多禄和若望要进殿，便求施舍。伯多禄和若望定睛看他，说：你看我们！}（第3-4节）即便如此，那人的思想仍未提高，仍然坚持乞讨。贫穷就是这样：遭拒绝时，它迫使人继续坚持。愿这让我们感到羞愧，我们在祈祷时却退缩。但请看伯多禄的温良：他说：\BibleQuote{你看我们。} 他们的神态本身，就显出了他们的品格。\BibleQuote{他就留意他们，期待从他们得些什么。伯多禄却说：银钱我没有，但我所有的，我给你。}（第5-6节）他没有说：我给你比金银好得多的东西；而是说什么？\BibleQuote{因纳匝肋耶稣基督的名字，起来行走罢。于是拉著他的右手，扶他起来。} \BibleRef{宗 3:7} 基督的方法也是如此：他常以言语治愈，有时用行动，有时也伸手，对于那些信心稍弱的人，使治愈不显得是自发。\BibleQuote{于是拉著他的右手，扶他起来。} 这一行动彰显了复活，因为它正是复活的图像。\BibleQuote{立时他的脚和踝骨就健壮了。他跳起来，站着行走。} \BibleRef{宗 3:8} 也许他是在试验自己，这样进一步检验，看所做的事是否徒劳。他的脚是软弱的，并非失去。有人说他甚至不知如何行走。\BibleQuote{同他们进了圣殿。} 真是奇妙。宗徒们没有催促他，但他自愿跟随，以此追随行动指明他的恩人。\BibleQuote{跳着赞美天主}；不是赞美他们，而是赞美藉他们行事的天主。那人是知恩的。\par

［ \BibleQuote{伯多禄和若望一同上到圣殿去，}等等。］ 你看他们如何恒心祈祷。 \BibleQuote{第九时辰：}他们在那里一同祈祷。 ［ \BibleQuote{有一个人，}等等。］ 那人当时正被人抬来。 ［ \BibleQuote{他们每天把他放在，}等等。］ （抬他的人把他抬来：） ［ \BibleQuote{在门口，}等等。］ 正当人们进入圣殿的时候。并且为使你不至于以为他们抬他是为了别的目的，而是为了使他能得到施舍，请听作者所说： \BibleQuote{好使他能向进圣殿的人求施舍。} \BibleRef{宗 3:1} 这也是为什么他也提到这些地方，以证实他所叙述的事。 \Quote{那是怎么回事，} 你可能会问，\Quote{他们为什么不把他带给基督？} 也许他们是某些不信的人，常常出没于圣殿，正如事实上当宗徒们行了如此伟大的奇迹后，他们也没有把他带到宗徒面前。 \BibleQuote{他求，} 经上记载，\BibleQuote{要得到施舍。} \BibleRef{宗 3:3} 他们的举止表明他们是虔诚而正义的人。 ［ \BibleQuote{伯多禄和若望注目前视，说道，}等等。］ （第4～5节。） 并且注意若望是如何处处沉默，而伯多禄也为他辩护； \BibleQuote{金银，} 他说，\BibleQuote{我都没有。} \BibleRef{宗 3:6} 他没有说：我在这里没有，像我们常说的那样，而是绝对地说：我没有。 \Quote{那么怎样？} 他可能会说，\Quote{你对我这个求告者不予理会吗？} 不是这样，而是：我所有的，你拿去吧。你是否注意到伯多禄何等谦逊，即使对受惠者也毫无炫耀？ ［ \BibleQuote{因这名，}等等。\BibleQuote{他握住他的手，}等等。］ \BibleRef{宗 3:7} 口和手做了一切。犹太人就是这类人；瘸腿的，本该求健康，却求钱财，匍匐在地：因为他们之所以聚集在圣殿就是为了谋财。那么伯多禄做了什么？他没有轻视他；他没有寻找某个富人；他没有说：如果奇迹不是行在某个\OriginalTerm{到那人}{\GreekText{εἰς} \GreekText{ἐκεἵνον}}身上，就算不得大事；他没有从他那里寻求什么荣誉，也没有在众人面前治好他；因为那人是在入口处，不在人群当中，即里面。但伯多禄没有追求这些；进入时也没有宣称此事；不，他是以他的举止吸引瘸腿的人来求。令人惊奇的是，他如此迅速地相信了。因为那些从长期疾病中获得释放的人，几乎不相信自己的眼睛。一旦治好，他就同宗徒们在一起，感谢天主。 \BibleQuote{他进入，} 经上说，\BibleQuote{同他们进了圣殿，走着跳着赞美天主。} \BibleRef{宗 3:8} 看他多么不安定，在喜悦的急切中，同时堵住了犹太人的口。并且，他跳跃是为了避免伪装的嫌疑；因为毕竟这是不可能欺骗的。因为若先前他完全不能行走，即使饥饿紧逼（如果他能够自己谋生，他本不会选择与抬他的人分享乞讨所得），那么现在更是如此。他又怎能为那些没有施舍给他的人假装呢？但那人即使在痊愈后仍是感恩的。如此，他的信德在两方面都显明出来，既因他的感恩，也因最近发生的事件。\par

他非常为众人所熟知，以致 \BibleQuote{他们认出了他。众百姓，} 经上说，\BibleQuote{看见他走路赞美天主；他们就\OriginalTerm{认出}{\GreekText{ἐπεγίνωσκον}}他就是那坐在圣殿美门口求施舍的人。} \BibleRef{宗 3:9} 说 \Quote{他们认出} 是很好的，因为由于所发生的事，他如今成了不识的人：因为我们用这词指难以辨认的对象。 ［ \BibleQuote{他们充满了惊奇诧异，因所发生在他身上的事。} ］ 必须相信，基督的名能赦免罪过，因为它甚至产生这样的效果。 （ \BibleQuote{他拉着伯多禄和若望的时候，众百姓全都跑到那叫撒罗满廊下，大为惊讶。} \BibleRef{宗 3:11} ）由于对宗徒们的好感和爱戴，那瘸子不愿离开他们；也许他公开感谢他们，赞美他们。 \BibleQuote{众百姓，} 经上说，\BibleQuote{都跑到他们跟前。伯多禄看见，就发言。} \BibleRef{宗 3:12} 又是他行事，并向百姓讲话。\par

在先前的场合，是舌音的现象引起了他们的注意，而现在是这奇迹；那时，他从他们的控告中找到了讲论的契机，而现在，是从他们的猜想中。那么，让我们来看看这篇演讲与先前的那篇有何不同，又有何相同。先前的那篇是在一间屋子里进行的，那时还没有人皈依，他们自己也还没有行过任何奇事；而这篇是在众人惊异、被治愈的人站在身旁时发表的，那时无人怀疑，不像另一场合有人曾说：\Quote{这些人充满了新酒。}\BibleRef{宗 12:13} 在那一次，他讲话时被所有宗徒围着；而这一次，只有若望与他同在，因为这时他已经胆大，变得更精力充沛了。这就是德性的本质：一旦开始，它就前进，永不停歇。也请注意，这奇迹发生在圣殿里，是天主上智的安排，以便其他人也能变得勇敢，而宗徒们不是在\OriginalTerm{隐秘之处}{\GreekText{εἰς} \GreekText{καταδύσεις}}、角落和暗中工作：虽然也不在圣殿内部人数众多之处。那么，请问，它是如何被相信的呢？被治愈的人自己宣告了这恩惠。他没有理由撒谎，也没有理由加入另一群人。那么，要么是因为地方宽敞，他在那里行奇迹，要么是因为那地方僻静。请注意结果：他们为一个目的上去，却成就了另一个目的。科尔乃略也是这样：他祈祷并禁食……* * * 但至此为止，他们总是称他为\Quote{纳匝肋人}。\Quote{因纳匝肋人耶稣基督之名，}伯多禄说，\Quote{行走。}因为在起初，所需要的是，他被相信。\par

我恳求你们，不要在故事的开头就放弃：若有人提到了某一特定的德性成就，然后又中断了一段时间，让我们重新开始。如果我们进入了\OriginalTerm{正确心境}{\GreekText{ἐν} \GreekText{ἕξει}}，我们将很快到达终点，很快达到顶峰。因为据说，热诚产生热诚，迟钝产生迟钝。一个已经完成了一点小改革的人，由此得到鼓励去接近更伟大的事，再从那里前进到超越它的事；正如火一样，它抓住的木材越多，就越猛烈，同样，热诚越点燃虔诚的思想，就越有效地武装起来对抗相反的事物。例如，在我们里面，像许多荆棘一样，安放着：发假誓、说谎、虚伪、诡诈、不诚实、辱骂、嘲笑、戏谑、猥亵、粗鄙；又另一类：贪财、掠夺、不义、诽谤、奸诈；又：邪恶的情欲、不洁、淫荡、私通、奸淫；又：嫉妒、争竞、愤怒、忿恨、怨恨、报复、亵渎，以及其他无数。如果我们首先在那些事上改革，不仅在这些事上成功了，而且通过它们，在随后的事上也成功了。因为理性那时获得了更大的力量去推翻那些别的恶习。例如，一个经常发誓的人，一旦根除了这种撒殚般的习惯，不仅赢得了这一点，而且在其他方面也会引入虔诚的习惯。因为我想，一个厌恶发誓的人，不会轻易同意去做任何其他邪恶的事；他会对已获得的德性心存敬畏。正如一个穿着漂亮长袍的人，会羞于在泥泞中打滚；这里也是如此。从此开始，他将学会不发怒，不打人，不侮辱。因为一旦他在小事上走上了正轨，整个事情就完成了。然而，常常发生这样的事：一个人已经改革了一次，然后又因疏忽而很容易地回到了旧罪中，以致情况不可挽回。例如，我们给自己定下不起誓的规矩；我们做得好，有三天甚至四天；之后被逼到绝境，我们把所有积累的收获都挥霍掉了；然后我们陷入懈怠和鲁莽。然而，放弃是不对的；必须再次热切地开始工作。因为据说，一个建造了房屋的人，然后看到自己的建筑被拆毁，会缺少再次建造的勇气。是的，尽管如此，也不应气馁，而应再次热切地开始工作。\par

那么让我们为自己订立每日的律法。目前先从比较容易的开始。让我们减去所有多余的路途，给我们的舌头加上辔头；谁也不要指着天主发誓。这里没有花费，没有疲劳，没有时间成本。只要愿意，一切就都成了。这是一个习惯问题。我恳求并乞求你们，让我们贡献出同样多的热忱。告诉我，如果我命令你们献出钱财，你们每个人不是都愿意按照自己的能力投入吗？如果你们看见我处于极度危险中，如果可能，你们不是会割下自己的肉来给我吗？我现在正处在危险中，而且极大的危险，即使我被关在地牢里，或遭受一万次鞭打，或在矿场中为囚犯，我也不能受更多的苦。那么就向我伸出援手吧。想想这危险有多大，我竟然未能改革这最微小的事情：我说\Quote{最微小}，指的是所需要的劳力。将来这样被质问时，我该说什么呢？\Quote{为什么你没有提出异议？为什么你没有命令？为什么你没有把律法摆在他们面前？为什么你没有制止那些不顺服的人？} 我光说我曾劝诫是不够的。会有人回答说：\Quote{你应该更严厉地责备；因为厄里也曾劝诫过。}\BibleRef{撒上 2:24} 但天主不许我把你们与厄里的儿子们相比。事实上，他确实劝诫了他们，说：\BibleQuote{不，我的儿子们，不要这样做；我听到关于你们的报告是邪恶的。}\BibleRef{撒上 2:24} 但随后圣经说，他没有劝诫他的儿子们；因为他没有严厉地或带着威胁劝诫他们。\BibleRef{撒上 3:13} 因为在犹太人的会堂里，律法如此有威力，教师所吩咐的无不遵行；而我们在这里却如此被轻视和拒绝，这难道不奇怪吗？我所关心的不是我自己的荣耀（我的荣耀是你们的好名声），而是为了你们的救恩。我们每天提高声音，在你们耳边呼喊。却没有人听。但我们仍然没有采取强硬措施。我怕在将来的审判之日，我们不得不为这过分和不合时宜的宽容而交账。\par

因此，我以响亮清晰的声音向众人宣告并作证，那些在这一过犯上臭名昭著的人，那些说出\BibleQuote{出自那邪恶者}\BibleRef{玛 5:37}（因为发誓就是这样）的话语的人，不得跨过教会的门槛。让这个月份作为你们在这件事上改正的期限。不要对我说\Quote{事务的紧迫迫使我发誓，否则人们不信我。}首先，戒除那些仅仅出于习惯的誓言。我知道许多人会嘲笑，但与其将来为之哭泣，不如现在被嘲笑。他们会嘲笑，那些人是疯狂的。请问，一个神志正常的人谁会嘲笑遵守诫命呢？但假设他们嘲笑；那么，他们嘲笑的不是我们，而是基督。你们听到这句话震惊了吧！我知道你们会。如果这条法律是我制定的，那么嘲笑就落在我身上；但如果立法者是另一位，那讥讽就转嫁到他身上。是的，基督曾被人吐唾沫，被掌击，被打脸。现在他也忍受着这些，这\OriginalTerm{并不奇怪}{\GreekText{οὐδὲν} \GreekText{ἀπεικὸς}}！为此，地狱已被预备；为此，不死的虫。看，我再次说并作证；谁愿意嘲笑就让他嘲笑，谁愿意讥讽就让他讥讽。我们正是为此而被设立的，要被人嘲笑和戏弄，要忍受一切。我们成了世上的\BibleQuote{垃圾}\BibleRef{格前 4:13}，正如真福保禄所说。如果有人拒绝遵守这一命令，那么那个人，我以我的话语，仿佛号角之声，禁止他踏上教会的门槛，无论他是王子，甚至是头戴王冠的人。要么把我从这个职位上免职，要么如果我要留下，就不要把我置于危险之中。我无法登上这宝座而不带来一些重大的改革。因为如果这是不可能的，那还不如站在下面。没有什么比一个对人民毫无益处的统治者更可怜的了。你们要努力，并注意这点，我恳求你们；让我们努力，确实会有更多的成果。守斋，恳求天主（我们也将与你们一同做），使这有害的习惯得以根除。成为世界的教师并不是什么了不起的事；到处都说在这城里确实没有一个人发誓，这并非小荣誉。如果此事实现，你们不仅将因自己的善工而获得报酬；实际上，我对你们是什么，你们也将成为世界是什么。确实，其他人也会效法你们；确实，你们将成为一个放在灯台上的蜡烛。\par

你会说：\Quote{就这样吗？}不，这不是全部，而是其他德行的开端。不发誓的人，必定会在其他方面达至虔诚，无论他愿不愿意，出于自尊和敬畏。但你可能会坚持说，大多数人不遵守，反而堕落。好吧，遵行天主旨意的一人，胜过万万千千的悖逆者。事实上，正因如此，一切都颠倒了，一切都混乱了——因为这就像戏院一样，我们追求数量，而不是精选的数目。众多人能有什么益处呢？你愿意知道，是圣人，而非数量，造就了群众吗？带领十万人去打仗，加上一位圣人，我们看看谁成就最多？农的儿子若苏厄出去打仗，独自成就了一切；其余的人毫无用处。亲爱的，你愿意看到，庞大的群众若不遵行天主的旨意，就与虚无无异吗？我确实希望、渴望，并且乐意被撕裂，以群众——是的，精选的群众——来装饰教会；但如果这不可能，那么少数人是精选的，就是我的愿望。你没有看见，拥有一颗宝石，胜过一万个铜钱吗？你没有看见，眼睛健全，胜过满身肥肉却失去视力吗？你没有看见，一只健康的羊，胜过一万只患瘟疫的羊吗？优秀的子女虽少，胜过许多患病的孩子；在天国里人少，在地狱里人多。我与群众有何相干？有何益处？没有。反而他们成了其他人的祸害。这就好比一个人可以在十个健康人或一万个病人中选择，却选择了后者加上那十个。那些行不出任何善事的人，只会让我们将来受惩罚，现在蒙羞。因为没有人会以我们人数多为荣；我们会因无用而受责备。事实上，当我们说\Quote{我们人数众多}时，人们总是告诉我们：\Quote{是啊，但都是坏的。}\par

看，我再警告，大声宣告：让没有人以为这是儿戏；我将排除和禁止悖逆者；只要我坐在这宝座上，我绝不放弃它的任何权利。如果有人把我从宝座上赶下来，那么我就不再负责；只要我负责，我就不能忽视他们；不是为了我自己的惩罚，而是为了你们的救恩。因为我极其渴望你们的救恩。为推进它，我忍受痛苦和烦恼。但你们要顺从，好使你们在此世和来世都能领受丰厚的赏报，并使我们共同获享永恒的福乐；借着天主独生子的恩宠和仁慈，愿光荣、权能和尊崇归于祂，及圣父和圣神，从今世直到永远。阿们。\par

\SourceRecord{Translated by J. Walker, J. Sheppard and H. Browne, and revised by George B. Stevens.；Nicene and Post-Nicene Fathers, First Series；Vol. 11.；Edited by Philip Schaff.；Buffalo, NY: Christian Literature Publishing Co.,；1889.；<http://www.newadvent.org/fathers/210108.htm>.}

% structure: p0001 source=h1 target=section reason=page-title

\section{宗徒大事录讲道辞 第九篇}

\BibleRef{宗 3:12}\par

\BibleQuote{伯多禄看见，就回答百姓说：以色列人哪，为什么因这事惊奇？为什么注视我们，好像因我们自己的能力和圣德使这人行走呢？}\par

篇讲道比上一篇更自由奔放。并非他在前一篇中有所畏惧，而是因为当时他所对之人是嘲弄者和讥诮者，不会忍受。因此，在先前那篇讲道的开头，他也通过开场白引起他们的注意：\BibleQuote{你们当知道这事，也请听我的话。}\BibleRef{宗 2:14} 但此处无需这种（\OriginalTerm{开场准备}{\GreekText{κατασκευἥς}}）。因为听众并非漠不关心。奇迹已激起他们全部关注，他们甚至充满恐惧和惊愕。因此，无需从那时开始，而应从另一个话题入手；事实上，通过拒绝来自他们的荣耀，他有力地安抚了他们。因为没有什么比不讲自己可荣耀的事、反而打消一切想要这样做的猜测，更有利、更能平息听众的了。事实上，他们因轻视荣耀而更增添了自己的荣耀，并表明刚才发生的事不是人的行为，而是天主的工程；他们的本分是与旁观者一同惊叹，而非接受惊叹。你是否看到他毫无野心，如何拒绝了那加给他的尊荣？古圣祖们也是如此；例如，达尼尔说：\BibleQuote{不是因我有什么智慧。}\BibleRef{达 2:30} 若瑟又说：\BibleQuote{解梦不是属于天主吗？}\BibleRef{创 11:8} 达味说：\BibleQuote{当狮子或熊来时，我奉主的名亲手撕裂它们。}\BibleRef{撒上 17:34} 同样，在此宗徒们说：\BibleQuote{你们为什么注视我们，好像因我们自己的能力和圣德使这人行走呢？}\BibleRef{宗 3:13} 不，甚至不是这样；因为我们不是因自己的功劳招致天主的作为。\BibleQuote{亚巴郎的天主、依撒格的天主、雅各伯的天主，我们祖先的天主。} 看他如何殷勤地将自己（\OriginalTerm{强行引入}{\GreekText{εἰσωθει}}）与古圣祖联系在一起，以免被人视为在传授新教义。在上一篇讲道中，他引用了先祖达味，在此则引用了亚巴郎等人。\BibleQuote{光荣了祂的仆人耶稣。} 又是一个谦卑的表达，如同开头的讲道一样。\par

但在此时，他开始详细说明暴行，并强调此事的可憎，不再像先前那样为其遮掩。他这样做，是希望能更有力地影响他们。因为他越证明他们负有责任，他的目的就越得以实现。\BibleQuote{光荣了祂的仆人耶稣，就是你们所交付的，并在比拉多面前否认的，那时比拉多已决定释放祂。} 控诉有两层：比拉多愿意释放祂，而你们却在祂愿意时不愿。\BibleQuote{但你们否认了那圣洁和公义者，反而要求把凶手赐给你们；你们杀害了生命的元首，天主却从死者中复活了祂，我们都是这事的见证人。}（第14-15节）他显示这行为的极大加重。\Quote{生命的元首，}他说。借这些话，他确立了复活的道理。\BibleQuote{天主从死者中复活了祂。}\BibleRef{宗 2:26} \Quote{这事从哪里可以看出来？}他不再引用先知，而是引用自己，因为现在他有权得到相信。之前，当他肯定祂复活时，他引用了达味的见证；现在，他说出之后，则引用了宗徒团体。\Quote{我们都是见证人，}他说。\par

\BibleQuote{因着对祂之名的信德，祂的名使你们所看见所认识的这人强壮了；那由祂而来的信德，在你们众人面前赐给了他完全的健康。}他\OriginalTerm{寻求讲述这事}{\GreekText{ζητὥν} \GreekText{τὸ} \GreekText{πρἅγμα} \GreekText{εὶπεἵν}}，立刻提出证据说：\BibleQuote{在你们众人面前，}他说，\BibleQuote{你们众人面前。}由于他曾严厉指责他们，指出他们所钉死的那位已经复活，他再次放宽，给予他们悔改的能力：\BibleQuote{如今，弟兄们，我知道你们是出于无知而行了这事，你们的首领也是如此。} \BibleRef{宗 3:17}这是一个开脱的理由。第二个是另一种。正如若瑟对他的弟兄们说：\BibleQuote{是天主打发我在你们以先来。} \BibleRef{创 45:5}他在前一次讲道中曾简要说的，即\BibleQuote{祂按天主的定旨和预知被交付，你们把祂取了去，}——这里他加以扩展：\BibleQuote{但天主先前藉众先知的口所预言的，祂的基督当受苦难的事，祂就这样应验了。} \BibleRef{宗 3:18}同时表明，若此事乃按天主的计划成就，就不是出于他们的作为。他暗示他们曾在十字架上辱骂祂的话，即\BibleQuote{让祂拯救祂，如果祂愿意祂；因为祂说过：我是天主子。如果祂信赖天主，就让祂现在从十字架上下来。} \BibleRef{玛 27:42}愚昧的人啊，这些是空话吗？这些事必须如此成就，先知们也为此作证。因此，如果祂没有下来，并非由于祂的软弱，而是由于大能。伯多禄以此作为对犹太人的辩护，希望他们接受他的话。\BibleQuote{祂就这样应验了，}他说。你们看到他现在如何将一切归于那源头？\BibleQuote{所以你们要悔改，}他说，\BibleQuote{并要归正。}他没有说\BibleQuote{从你们的罪中离弃，}而是说\BibleQuote{使你们的罪得以涂抹，}意思相同。然后他添加益处：\BibleQuote{这样，那安舒的时期就必从主面前来到。} \BibleRef{宗 3:19}这表明他们处于悲伤的境况，因许多战争而卑微。因为如同一个人被火焚烧，渴望安慰，这个表达正适用于此。现在看他如何进展。在他的首次讲道中，只略微提及复活和基督在天上坐着；但这里他也谈到祂可见的来临。祂必将为你们预定的耶稣基督派遣来，\BibleQuote{就是天必须（\Emph{即}必然必须）收留的那位，直到万物复兴的时候。} \BibleQuote{就是天主藉祂的圣先知的口所预言的，}他继续说，\BibleQuote{从创世以来。梅瑟确实对祖先说过：主、你们的天主，要从你们弟兄中给你们兴起一位像我一样的先知；你们当在一切事上听从祂。}之前，他提到了达味，这里他提到梅瑟。\BibleQuote{一切事，}他说，\BibleQuote{就是祂所说的。}但他不是说\Quote{是基督所说的，}而是\Quote{是天主藉众圣先知的口，从创世以来所说的。}（第20-21节）然后他转向可信的根据，说：\BibleQuote{主、你们的天主，要从你们弟兄中给你们兴起一位像我一样的先知；你们当在一切事上听从祂。}接着是惩罚的严重性：\BibleQuote{凡不听从那位先知的，必从民中全然灭绝。撒慕尔和其后所有发言的先知，也都预言过这些日子。}（第23-24节）他在这里所作的区分很好。因为每当他说到大事时，都引证古人的话。他找到了一段同时包含两重真理的经文；正如在另一篇讲道中他说：\BibleQuote{直到祂把仇敌放在祂脚下。} \BibleRef{宗 2:35}值得注意的事实是，两件事同时存在，即服从与不服从，以及惩罚。\BibleQuote{像我一样，}他说。那么你们为何惊慌？\BibleQuote{你们是先知之子} \BibleRef{宗 3:25}：因此，先知们是对你们说话，这一切事情都是为了你们而成就。因为他们认为由于自己的暴行，他们已经与天主隔绝（的确没有合理的理由，那位现在被钉死的，竟会接纳他们为自己的子民），他向证明，两者都符合预言。\BibleQuote{你们是，}他说，\BibleQuote{先知和盟约之子，就是天主与我们的祖先所立的盟约，对亚巴郎说：\InnerQuote{地上万族都要因你的后裔蒙受祝福。}}他继续说，\BibleQuote{天主先兴起\OriginalTerm{祂的圣子}{\GreekText{τὸν} \GreekText{Παῖδα}}，派遣了祂到你们这里来。} \Quote{固然也到别人那里，但先到你们这些钉死祂的人那里。} \BibleQuote{为祝福你们，}他添加说，\BibleQuote{叫你们各人回转，离开你们的邪恶。} \BibleRef{宗 3:26}\par

现在让我们更仔细地思考刚才所读的内容。（复习。）首先，他确立了奇迹是由他们行的这一点，说：\Quote{你们为什么惊奇？}他不让这个断言被怀疑：为了增加其分量，他预料到他们的判断。他说：\Quote{你们为什么这样注视我们，好像是我们用自己的能力或虔诚使这人行走似的？} \BibleRef{宗 3:12}如果这事困扰并迷惑了你们，要学习谁是行此者，不要惊讶。并且注意，每当提到天主，并说万物都来自祂时，他总是毫无畏惧地责备他们：正如上面他所说的：\Quote{一位在你们中间为天主所认可的人。} \BibleRef{宗 2:22}在每种场合，他都提醒他们所犯的暴行，以便复活的事实得以确立。但在这里，他还附加了别的；因为他不再说\Quote{纳匝肋的}，而是说：\Quote{我们祖先的天主光荣了祂的仆人耶稣。} \BibleRef{宗 3:13}也要注意他的谦卑。他没有责备他们，也没有立即说：\Quote{那么现在就相信吧：看，一个瘸了四十年的人，因耶稣基督的名被扶起来了。}他没有这样说，因为这会激起反对。相反，他开始称赞他们钦佩这事，并再次以他们的祖先称呼他们：\Quote{以色列人。}此外，他没有说治愈他的是耶稣；而是说：\Quote{我们祖先的天主光荣了}等等。但随后，恐怕他们会说：这怎么合理——天主会光荣一个罪犯？因此他提醒他们在比拉多面前的审判，表明如果他们考虑一下，祂并不是罪犯；否则比拉多不会愿意释放祂。他没有说\Quote{当比拉多希望时}，而是说\Quote{坚决要释放祂。} \Quote{但你们否认了那圣者}等等。（13-14节。）你们要求释放那个杀害他人的人；你们不愿意拥有那使被杀者复活的一位！并且，恐怕他们会再次问：既然以前祂没有提供帮助，天主现在又怎么会光荣祂呢？他提出先知作证，说这样是应当的。\Quote{但天主先前所预告的}等等。 \BibleRef{宗 3:18}然后，恐怕他们会以为天主的计划是他们的借口，他先责备了他们。此外，否认祂\Quote{当面}不是寻常的事；因为比拉多愿意释放祂。而你们不能否认这事，因为那个被优先询问的人是反对你们的见证。这同样也是深奥计划的一部分。这里显示出他们的无耻和厚颜：一个外邦人，初次见到祂的人，虽然没听到任何惊人之处，却释放了祂；而他们这些在祂的奇迹中长大的人，却做了完全相反的事！因为，正如他所说的：\Quote{当比拉多已决定释放祂时}，以免被认为他是出于偏爱而这样做，我们读到：他说，你们有一个惯例要释放一个囚犯：那么你们愿意我释放这个人给你们吗？ \BibleRef{玛 27:15}\Quote{但你们否认了那圣者和义人。} \BibleRef{谷 15:6}他没有说：\Quote{你们交付了}；而是处处说：\Quote{你们否认了。}因为，他们说：\Quote{我们没有王，只有凯撒。} \BibleRef{若 19:15}他不仅说：你们没有求释放无辜者，以及\Quote{你们否认了}祂，而且说：\Quote{你们杀了}祂。当他们刚硬时，他避免使用这样的语言；但当他们的心最被感动时，他才击打要害，因为他们现在能够感受了。正如当人醉酒时我们什么也不对他们说，但等他们清醒，从醉态中恢复时，我们才斥责他们；伯多禄也是如此：当人们能够理解他的话时，他才尖锐自己的舌头，对他们提出许多控告：那就是，天主光荣了的那一位，他们却交付了；比拉多本想无罪释放的那一位，他们当面否认了；他们把强盗看得比祂更优先。\par

再观察他如何隐蔽地谈及基督的权能，显示祂自己复活了：正如他在第一篇讲词中说过的，\Quote{因为祂不能被它控制}\BibleRef{宗 2:24}，同样，这里他说，\Quote{你们杀了生命之主}\BibleRef{宗 3:15}。由此可见祂所拥有的生命并非来自他人。罪恶的元首（或创始者）乃是首先带来罪恶的那位；凶杀的元首或创始者是首先发起凶杀的那位；同样，生命之主（或创始者）必须是从自身拥有生命的那位。\Quote{天主却从死者中复活了祂}，他继续说：说了这话之后，他又加了一句：\Quote{因祂的名号，这人的强壮是因信祂的名号，并由祂而来的信德赐给他在你们众人面前完全的健康。}【\OriginalTerm{经由祂的信德}{\GreekText{ἡ} \GreekText{δι}' \GreekText{αὐτοῦ} \GreekText{πίστις}}】然而，其实是\OriginalTerm{对祂的信德}{\GreekText{ἡ} \GreekText{εἰς} \GreekText{αὐτὸν} \GreekText{πίστις}}，即\Quote{对祂的信德}（作为对象）成就了一切。因为宗徒们没有说\Quote{经由这名}，而是说\Quote{因这名}，那人是在祂（\OriginalTerm{对祂}{\GreekText{εἰς} \GreekText{αὐτὸν}}）内相信的。但他们还不敢使用\Quote{对祂的信德}这种表达。为使\Quote{经由祂}的说法不至于太低微，请注意，在说了\Quote{因祂名号的信仰}之后，他又加上\Quote{祂的名号使他强壮}，然后才说：\Quote{是的，那经由祂的信德赐给了他这完全的健全。}请观察他如何暗示，在\OriginalTerm{且那个}{\GreekText{καὶ} \GreekText{ἐκεῖνο}}先前的表达中，\Quote{天主却从死者中复活了祂}，他不过是迁就他们低下的领悟力。因为那位祂的名号能使一个瘫子（如同死人一般）起来，祂自己复活并不需要别人的帮助。请注意他如何处处引用他们自己的见证。如此上文中他说：\Quote{正如你们自己所知道的}；以及\Quote{在你们中间}；这里又：\Quote{你们所看见和认识的：在你们众人面前}\BibleRef{宗 2:22}。虽然他们不知道那是\Quote{因祂的名号}，但他们知道那人原是瘸腿的，如今站立在那里健康无恙。行此事的人自己承认，不是凭自己的能力，而是凭基督的能力。如果这断言毫无根据，如果他们自己并未真正相信基督已复活，他们就不会寻求为一个死人而不是为自己建立荣耀，尤其是在众人注视之下。然后，当他们的心意惊惶时，他立刻用\Quote{弟兄}的称呼鼓励他们：\Quote{现今，弟兄们，我知道等等}。因为在先前的讲词中，他什么也没有预言，只讲到基督：\Quote{所以以色列全家当确实知道}；在这里他加上劝告。在那里他等百姓说话；在这里他知道他们已经成就了多少，并且当前的集会对他们有更好的态度。\Quote{你们是出于无知而做了}。然而上述情况不能归因于无知。选择强盗，拒绝被判无罪的那位，甚至想要消灭祂——如何能归为无知？尽管如此，他给他们自由否认，并改变对已发生之事的想法。你们知道自己杀了无辜者；但你们杀害\Quote{生命之主}这件事，大概你们并不知道。他不仅为这些人辩护，也为那恶谋的主要策划者辩护：\Quote{你们和你们的首领}；因为如果他转为控告，必定会引起反对。因为作恶的人，当你指责他犯的某种邪恶时，他为了开脱自己而变得更激烈。他不再说\Quote{你们钉死了}、\Quote{你们杀了}，而是说\Quote{你们做了这事}，引导他们寻求宽恕。如果那些首领出于无知而做，那么在场的这些人更是如此。\Quote{但天主先前借着众先知的口所预言的事}\BibleRef{宗 3:18}等。值得注意的是，无论第一篇还是第二篇讲词，他说相同的内容，即在前一篇：\Quote{天主按祂的定旨和预知}；在这一篇：\Quote{天主先前预示基督当受难}；两者都没有引用任何特定经文作为证据。事实上，每一段这样的经文都伴随着许多控告，并提到为他们存留的惩罚：\Quote{我将交付}，有人说，\Quote{恶人作为祂坟墓的报偿，富人作为祂死亡的报偿}\BibleRef{依 53:9}。又，* * * \Quote{那些事}，他说，\Quote{天主先前借着众先知的口所预言基督当受难的事，祂已这样应验了。}这显示了那\Quote{定旨}的伟大，因为\Emph{全部}都说到它，而非仅一位。不能说因为这事出于无知，它就无关天主的安排。请看天主的智慧多么伟大，当祂利用别人的恶行来实现那必须发生的事时。\Quote{祂已应验了}，他说：使他们不至于认为有什么缺乏；因为凡是基督必须受的，都已应验了。但不要因为先知说了这话，你们出于无知做了这事，就以为这足以开脱你们。不过，他并没有这样说，而是用更温和的语气说：\Quote{你们应悔改}\BibleRef{宗 3:19}。\Quote{为什么？因为要么是出于无知，要么是天主的安排。}\Quote{为叫你们的罪得以涂去。}我不是指在钉十字架时所犯的罪行；也许那些是出于无知；但为使你们其他的罪得以涂抹：仅仅是这个。\Quote{这样，安舒的时期必临到你们。}这里他暗示性地谈论复活。因为那些实在是安舒的时期，也是保禄所期望的，他说：\Quote{我们在这帐棚里叹息，背负重担}\BibleRef{格后 5:4}。然后为证明基督是那安舒日子的原因，他说：\Quote{祂也必差遣耶稣基督，就是预先为你们指定的}\BibleRef{宗 3:20}。他没有说\Quote{使你的罪得以涂去}，而是\Quote{你们的罪}；因为他也暗示那罪。\Quote{祂必差遣}。从何处？\Quote{就是天堂必须收留的}\BibleRef{宗 3:21}。仍然\Quote{必须收留}？为什么不直说\Quote{天堂收纳了祂}？这如同谈论古代的事；他说，这是天意安排，如此决定；对祂的永恒存在则只字未提。——\Quote{梅瑟曾对祖先说：主必要给你们兴起一位先知}；\Quote{凡祂对你们所说的一切，你们都要听从}；在说了\Quote{天主借着众先知的口所说的一切话}\BibleRef{宗 3:22}之后，现在他确实引出了基督自己。因为如果祂预言了许多事，并且必须听从祂，那么说先知们说了这些事也不算错。但此外，他愿意显示先知们确实预言了同样的事。如果有人仔细去查考，就会发现旧约中虽然隐晦，但确实说了这些事。\Quote{祂是特意设计的}，他说：在祂毫无新奇之处。这里他也用还有许多事要应验这一思想来警告他们。但如果是这样，他怎么又说\Quote{应验了}\BibleRef{宗 3:18}？那些必须\Quote{基督当受}的事已经应验；但必须发生的事尚未。\Quote{主天主要从你们弟兄中给你们兴起一位先知，像我一样。}这最能平息他们。你是否注意到低微与崇高并列——那位要升到天上的竟像梅瑟？但这也是一件大事。因为事实上祂并不简单地像梅瑟，如果说\Quote{凡不听从的必从民中灭绝}。还可以提到无数其他事情表明祂不像梅瑟；所以他所处理的是一段强有力的经文。\Quote{天主必从你们弟兄中给你们兴起一位}，梅瑟说，等等：因此梅瑟自己威胁那些不听从的人。\Quote{是的，并且所有先知}等等：这一切都是为了吸引人。\Quote{是的，并且所有先知}，宗徒说，\Quote{从撒慕尔起}。他避免一一列举，以免使讲词太长；但引述了梅瑟那决定性的见证之后，便略过其余。\Quote{你们}，他说，\Quote{是先知和那盟约之子，是与天主所立的盟约}\BibleRef{宗 3:25}。\Quote{盟约之子}，即继承人。为了不让他们以为这恩赐是出于伯多禄的情面，他表明自古这恩赐就是属于他们的，为使他们更相信这也是天主的旨意。\Quote{天主先给你们兴起祂的圣子耶稣，差遣祂来}\BibleRef{宗 3:26}。他没有简单说\Quote{给你们差遣了圣子}，而且是在复活之后，当祂被钉死之后。为使不让他们以为是他自己给了这恩赐，而非圣父，他说：\Quote{为祝福你们}。因为如果祂是你们的弟兄，并祝福你们，这事就是一个应许。\Quote{先给你们}。就是说，你们不但没有与这些祝福无分，而且祂还愿你们成为他人这些祝福的促进者和来源。因为你们不该有被弃的感觉。\Quote{兴起}：再次，复活。\Quote{叫你们各人回头}，他说，\Quote{离开罪恶}。这样祂祝福你们：不是笼统的。这是什么祝福？一个伟大的祝福。因为当然，使人离开恶行本身并不足以赦免它们。如果不足以赦免，又怎能赐予祝福？因为不能认为犯罪的人立刻就变成了有福的；他只是被从罪中释放。但这\Quote{像我一样}绝不适用。\Quote{你们要听从祂}，他说；不仅如此，他还加了一句：\Quote{凡不听从那位先知的，必从民中全然灭绝。}当他显示了他们犯了罪，并授予他们宽恕，应许福事之后，那时，的确那时，他说：\Quote{梅瑟也说了同样的事。}这是什么联系：\Quote{直到万物复兴的时候}，然后引出梅瑟，说基督所说的一切必将实现？然后，另一方面，作为称赞（所以为此你们也当听从）：\Quote{你们是先知和盟约之子}，即继承人。那么为什么你们对待本该属于你们的，好像是他人的？的确，你们行了当受谴责的事；但你们仍可获得宽恕。说了这些，他现在有理由说：\Quote{天主差遣了圣子耶稣来祝福你们。}他没有说\Quote{拯救你们}，而是说更大的事：被钉死的耶稣祝福了钉死祂的人。\par

那么，让我们也效法祂。让我们驱逐那谋杀和敌意的精神。不报复还不够（因为甚至在旧约时代，这一点就已显明），但让我们为伤害我们的人做一切，如同我们对亲密朋友一样，如同对我们自己一样。我们是祂的跟随者，我们是祂的门徒，祂在被钉十字架后，仍为那些杀害祂的人调度一切，并为此派遣了祂的宗徒。然而，我们时常遭受公义的惩罚；但那些人不仅行事不义，而且不敬；因为祂是他们的恩主，祂没有作恶，他们却把祂钉在十字架上。原因何在？为了他们的名声。但祂自己却使他们成为可敬的对象。\BibleQuote{经师和法利塞人坐在梅瑟的座位上，凡他们所吩咐你们的，你们都要遵守，但不要效法他们的行为。}\BibleRef{玛 23:2} 又在另一处说：\BibleQuote{去，把你的身体给司铎看。}\BibleRef{玛 8:4} 此外，当时祂本可以毁灭他们，却拯救了他们。那么，让我们效法祂，愿没有人是敌人，没有人是仇敌，除了魔鬼。\par

不发誓的习惯对达到这个目的贡献不小：我指的是不轻易发怒；而通过不轻易发怒，我们也将不会有敌人。砍掉一个人的誓言，你就剪断了他愤怒的翅膀，你窒息了他所有的激情。据说，发誓是愤怒的风。降下船帆；当没有风时，不需要船帆。那么，如果我们不吵闹，也不发誓，我们就切断了激情的筋腱。如果你对此怀疑，只需试验一下。把永远不发誓作为一条法律强加给那个暴躁的人，你就不必再向他宣讲节制。于是整个事情就完成了。因为即使你还没有发虚誓，但仅仅发誓，难道你不知道你使自己陷入了多么荒谬的后果吗——把自己捆绑在一种绝对的必然性上，如同用绳索一般，并使自己陷入各种权变，如同人们研究如何从无法逃脱的邪恶中拯救自己的灵魂，或者，如果做不到，则（由于那自我强加的必然性）被迫从此在烦恼、争吵中度过余生，并诅咒你的愤怒？但一切都是徒劳，毫无意义。威胁吧，\OriginalTerm{划定界限}{\GreekText{διόρισαι}}，做一切事，无论什么，都不要发誓；[这样]：你就有能力\OriginalTerm{撤销}{\GreekText{ἀναλὕσαι}}你所说和所做的，只要你愿意。因此今天我必须更温和地对你们说话。因为你们已经听了我的话，并且大部分改革已由你们完成，那么现在让我们看看，为什么引入发誓的做法，又为什么允许它存在。在向你们讲述它们的起源，以及它们何时被构想，如何，由谁构想时，我们将以此报告来回报你们的服从。因为适合的是，那已经正确实践的人，应该被教导这事的哲理，但尚未行义的人，不配知道这段历史。\par

在亚巴郎的时代，他们立了许多盟约，宰杀祭牲，奉献祭品，但那时还没有誓言。那么誓言是从何而来的呢？当邪恶增加，一切都混乱颠倒，当人们转向偶像崇拜时，那时，那时，当人们似乎不再值得信任时，他们呼求天主作证，仿佛为所说的话提供了充分的保证。事实上，誓言正是如此：它是一种在人的原则不可信赖时的担保。因此，在指控发誓者时，第一项指控就是——没有保证人，且是一个伟大的保证人，他就不可信任：因为如此极端的无信，以致他们不以人为保证人，而偏要天主！第二，同样的指控也针对接受誓言者：在立约的问题上，他必须将天主拖来作担保，并且除非得到天主，否则拒绝满足。哦，何等过分的愚昧，何等的傲慢行为！你这虫豸，尘土，灰烬，蒸气，竟将你的主拖来作保证人，并迫使他人也如此行！告诉我，如果你的仆人们彼此争吵，互相交换保证，而这位同仆宣称他本人除非得到他们的共同主人作为保证人，否则不会满意，难道他不会受到无数的鞭打，并明白主人另有用途，不应被用于此类事吗？为什么我提到同仆？因为即使他选择任何体面的人，那人难道不会认为这是一种侮辱吗？但你说，我不想这样做。那么，也不要迫使别人这样做：因为当只涉及人时，这事如此——如果你的对方说：\Quote{我以某人为我的保证人，}你不允许他。\Quote{那么，}你说，\Quote{我就要失去我所给的吗？}我不是说这个；而是说你允许他侮辱天主。因此，强迫他人发誓的人将受到比发誓者更重的不可逃避的惩罚；同样，当无人要求时，主动给人发誓也是如此。更糟的是，每个人都随时准备发誓，为了一个铜钱，为了琐碎小事，为了自己的不义。这一切都是在没有伪誓的情况下说的；但如果伪誓随之而来，那么强迫发誓的人和接受誓言的人都把一切颠倒了。\Quote{但有些事，}你会说，\Quote{是未知的。}那么请将这些考虑在内，不要疏忽行事；但如果你疏忽了，就自己承受损失作为惩罚。这样损失总比以另一种截然不同的方式损失更好。告诉我——你强迫人发誓，是出于什么期望？是希望他发伪誓吗？这完全是疯狂；审判将落在你头上；你宁失钱财，也不愿他丧亡。为何这样行事，既损害自己，又侮辱天主？这是野兽和不信之人的精神。但你这样行，是期望他不会发伪誓？那么没有誓言也信任他。\Quote{不，有许多人，}你回答，\Quote{在没有誓言时竟敢欺诈；但一旦发了誓，就会克制。}你自欺欺人，人哪。一个人一旦学会了偷盗和损害邻人，就常常敢践踏誓言；反之，如果他回避发誓，他更会回避不义。\Quote{但他不情愿地受影响了。}那么，他应得宽恕。\par

然而，我为何谈及这类誓言，却略过市集上的那些呢？因为关于后者，你无法提出任何借口。为了十个小钱，你在那里又发誓又背誓。事实上，因为雷电并未真的从天而降，因为万物并未倾覆，你竟将天主捆绑在你的羁绊中：为了一点点蔬菜、一双鞋、一点钱财，就呼求祂作证。这是什么意思？我们不要以为，因为我们未受惩罚，因此我们就无罪；这出于天主的仁慈，并非我们的功德。让你的誓言成为针对你自己的孩子、你自己的诅咒吧：你说：\Quote{否则让刽子手鞭打我的肋骨。}但你不敢。难道天主不如你的肋骨宝贵？祂不如你的头颅珍贵？你说：\Quote{否则让我瞎眼。}但你不说。基督如此宽待我们，甚至不许我们指着自己的头起誓；而我们却如此不珍惜天主的荣耀，在任何场合都要拖祂出来！你们不知道天主是什么，也不知道当以怎样的嘴唇呼求祂。当我们谈论某位杰出人物时，我们说：\Quote{先洗你的嘴，然后再提他。}然而，那个宝贵的名号，超乎万名之上的名号，那在普世奇妙的名号，那魔鬼听见就战栗的名号，我们却随意拉扯！唉！习惯的力量！因此那名号变得廉价。无疑，如果你要求某人在圣殿中发誓，你会觉得那个誓言是可怕的。然而，为何这样它就显得可怕呢？岂不是因为我们平时随意使用那个，而不使用这个吗？因为仅仅提名天主时，难道不该感到战栗吗？如今，在犹太人当中，祂的名号被如此尊敬，以至于写在牌子上，除了大司祭之外，没有人允许佩戴那字符。我们却像普通词汇一样滥用祂的名号。如果连称呼天主都不允许所有人，那么呼求祂作证，这是何等狂妄！不，这是何等疯狂！因为如果需要（而不是这样）丢弃你所有的一切，难道你不该乐意舍弃一切吗？看，我郑重宣告并作证：改革这些论坛上的誓言，这些多余的誓言，并把所有想要发誓的人带到我这里来。看，在这集会面前，我嘱咐那些被派去照料祈祷之所的人，我劝勉并命令他们：不许任何人随意发这样的誓；或者更确切地说，不许任何人以其他方式发誓，而是无论何人，都应被带到我这里来，因为即使在这些小事上，也不得不来找我们，就像对待小孩子一样。但愿没有这个必要！在某些事上仍需受教，这是羞耻。你竟敢触摸祭台，而你是一个未受洗的人？不，更糟的是，你受过洗，竟敢把手放在祭台上——这甚至不是所有圣职者都能触摸的——并借此发誓。你不会去把手放在你孩子的头上，却触摸祭台，而不战栗、不害怕？把这些人带到我这里来；我要审判，并让他们欢欢喜喜地离开，无论哪一方。随你选择，我立法规定：根本不可发誓。如果我们将一切都化为乌有，那还有什么得救的希望？难道这就是你的账单和契据的结局，以至于你牺牲自己的灵魂？你得到的利益有多大，竟能抵得上损失？他背誓了吗？你毁了他也毁了你自己。但他没有背誓吗？即便如此，你仍然毁了他（和你自己），因为强迫他违犯诫命。让我们从灵魂中驱除这个疾病；至少让我们把它赶出论坛，赶出店铺，赶出其他工作场所；我们的收益只会更大。不要以为通过违犯天主的法律就能确保世俗计划的成功。\Quote{但他不信我，}你说；事实上我有时听有些人这样说：\Quote{除非我发无数的誓，否则那人不会信我。}是的，这要怪你自己，因为你轻易发誓。若非如此，反而让所有人都清楚你不发誓，请相信我的话，你只需点头就会比那些发誓成千上万的人更易被相信。试想：你更相信谁？是我这个不发誓的人，还是那些发誓的人？\Quote{是的，}你说，\Quote{但你是领袖和主教。}那么，假如我向你证明，并非仅仅因为这个原因呢？请诚实地回答我，我恳求你；如果我习惯不停地发誓，我的职位能为我做什么？一点也不能。你看到不是因为此原因吗？你究竟得到了什么？回答我。保禄忍受过饥饿；那么你难道不也选择饥饿，而不愿违犯天主的任何一条诫命吗？你为什么如此不信？你在这里，准备为不发誓而做一切、忍受一切；难道祂不会赏报你吗？祂，每日供养发誓者和背誓者，当你服从祂时，会让你饿死吗？愿所有人都看到，在这教会集会的人中没有一个发誓者。也藉此让我们显明出来，而不仅凭我们的信经；让我们以此标记将我们与外邦人及所有人区分开来。让我们把它当作来自天上的印记，好让我们处处被看作君王自己的羊群。让我们凭口舌被人认出，首先，如同外邦人被他们的口舌认出；正如说希腊语的人与外邦人区别开来，我们也应如此被认出。回答我：那些被称为鹦鹉的鸟，如何知道它们是鹦鹉？岂不是因为它们像人一样说话？那么就让我们像宗徒一样说话；像天使一样说话。如果有人要你发誓，告诉他：\Quote{基督已发言，我不发誓。}这就足以为一切德行打开道路。它是通往宗教的门户，是通往虔诚智慧的大道；是一种训练学校。让我们遵守这些，以便我们也能获得将来的恩赐，借着我们主耶稣基督的恩宠和仁慈，愿光荣、权力和尊崇归于祂，与圣父及圣神，自今至永远，世世无尽。阿们。\par

\SourceRecord{Translated by J. Walker, J. Sheppard and H. Browne, and revised by George B. Stevens.；Nicene and Post-Nicene Fathers, First Series；Vol. 11.；Edited by Philip Schaff.；Buffalo, NY: Christian Literature Publishing Co.,；1889.；<http://www.newadvent.org/fathers/210109.htm>.}

% structure: p0001 source=h1 target=section reason=page-title

\section{宗徒大事录讲道第十篇}

\BibleRef{宗 4:1}\par

\BibleQuote{当他们向百姓说话时，司祭和圣殿警官来到他们那里。}\par

他们初次考验之后还没来得及喘口气，立刻就进入了另一些考验。看事件是如何安排的。首先，他们一同被嘲笑；这不是小考验。其次，他们进入危险之中。这两件事并非接连发生；而是在宗徒们先因两篇讲道赢得钦佩，并行了显著奇迹之后，他们壮起胆来，由于天主的安排，才进入了比武场。但我愿你们思考：那些在基督的事上需要找人把祂交出来的人，如今竟亲手逮捕宗徒，他们在基督被钉死后变得更加胆大妄为、更加厚颜无耻。实在，罪在萌生之际还带有某种羞耻感；但一旦完全成熟，就使犯罪者更不知羞耻。据说：\BibleQuote{和圣殿警官}。其目的又是要使人所行的事带有公开的犯罪性质，而不是当作个人行为来追究：事实上，他们一贯如此行事。\par

\BibleQuote{他们因他们教训百姓而烦恼。}\BibleRef{宗 4:2} 不仅仅因为他们教训，更因为他们宣告，不但基督自己从死者中复活了，而且我们藉着祂也要复活。\BibleQuote{因为他们教训百姓，并藉着耶稣宣讲死人复活。} 祂的复活如此有力，以致祂也成为他人复活的原因。他们就下手拿住他们，把他们押到第二天，因为天已经晚了。\BibleRef{宗 4:3} 何等无耻！他们不惧怕群众；因为圣殿警官也和他们在一起；他们手上还沾着前一个牺牲者的血。据说：\BibleQuote{因为天已经晚了。} 那些人这样做，是要挫他们的锐气，并看守他们；但这拖延反而使宗徒们更加勇敢。想想被捕的是谁。他们是宗徒的首领，如今成为其余宗徒的榜样，叫他们不再彼此寻求支持，也不再想要在一起。\BibleQuote{然而，许多听了话的人都信了；男人的数目约有五千。}\BibleRef{宗 4:4} 这是怎么回事？他们看见他们受尊敬吗？他们不是看见他们被捆绑吗？他们怎么信了呢？你看见明显的功效吗？就连那些已经相信的人，本可能变得软弱。但不再是那样了：因为伯多禄的讲道已经把种子深深埋在他们里面，抓住了他们的理智。因此，［他们的敌人］被激怒，因为他们不惧怕他们，不把眼前的困难当回事。因为，他们说，如果那位被钉死的能行这样的大事，使瘸子行走，我们也不怕这些人了。这又是天主的安排。因为现在相信的人比先前的更多。因此，他们当着众人的面捆绑宗徒，是要使众人也更害怕。但结果相反。他们不在百姓面前审问他们，而是私下审问，以免听者因他们的勇敢而获益。\par

\BibleQuote{次日，他們的首領、長老和經師，以及大司祭亞納斯、蓋法、若望、亞歷山大，並大司祭家族的人，都聚集在耶路撒冷。}（第5-6節。）因為現在，連同其他的邪惡（時代的），法律已不再被遵守。他們又用法庭的形式來處理此事，要藉著不義的判決來定他們的罪。\BibleQuote{他們把他們放在中間，問說：你們憑什麼能力，或憑誰的名，做了這事？}\BibleRef{宗 4:7}但他們對此心知肚明，因為他們是因\Quote{他們藉耶穌宣講復活}而惱怒，才逮捕了他們。那麼他們質問他們是什麼目的？他們期望在場的人數會使他們改變主意，以爲藉此可以挽回一切。請看他們說什麼：\BibleQuote{你們憑誰的名做了這事？那時，伯多祿充滿聖神，對他們說。}\BibleRef{宗 4:8}如今，請你們記起基督的話：當他們把你們交到會堂時，不要思慮如何或說什麼，因為那時說話的是你們父的聖神。\BibleRef{路 12:11}所以他們享有大能。那麼伯多祿說什麼？\BibleQuote{你們百姓的首領和以色列的長老。}請注意這人的基督徒智慧，多麼充滿信心：他沒有一句侮辱的話，而是恭敬地說：\BibleQuote{你們百姓的首領和以色列的長老，如果我們今天因對這瘸子所行的善事被查問。}他勇敢地對付他們；在講話開頭就揭露了他們，並提醒他們過去的事：他們是因一件慈善工作而傳喚他們。彷彿他在說：按理我們應該因這事而得冠冕，被宣稱為恩人；但既然\Quote{我們竟因對一個瘸子所行的善事而受審}，那人既非富人、非有權勢、非貴族——然而在這種情況下，誰還會嫉妒呢？這實在是一種十分有力的（\OriginalTerm{宣告}{\GreekText{ἀπαγγελία}}，另一抄本作\OriginalTerm{應許}{\GreekText{ἐπαγγελία}}）表述方式；他表明他們是在刺傷自己：——\BibleQuote{這個人怎樣痊愈：你們全體和全以色列百姓都該知道，是因納匝肋人耶穌基督的名。}——這是最使他們惱怒的。因為這正是基督曾告訴門徒的：\BibleQuote{你們在暗處所聽到的，要在屋頂上宣講出來。}\BibleRef{瑪 10:27} \BibleQuote{奉耶穌基督的名，}他說，\BibleQuote{納匝肋人耶穌，你們所釘死、天主從死者中復活的那一位，正是憑著祂，這個人站在你們面前痊愈了。}\BibleRef{宗 4:10}不要以為，他說我們隱瞞地域或死亡的方式。\BibleQuote{你們所釘死、天主從死者中復活的那一位，正是憑著祂，這個人站在你們面前痊愈了。}再次提到死亡，再次提到復活。\BibleQuote{這石頭，}他說，\BibleQuote{是你們匠人所棄而不用的，卻成了屋角的基石。}\BibleRef{宗 4:11}他也提醒他們一句足以驚嚇他們的話。因為經上說過：\BibleQuote{誰跌在這石頭上，必會跌碎；這石頭落在誰身上，必會把他壓碎。}\BibleRef{瑪 21:44}——除祂以外，沒有救恩，\BibleRef{宗 4:12}伯多祿說。你想，這些話會使他們受到多大的創傷！\BibleQuote{因為在天下人間，}他繼續說，\BibleQuote{沒有賜下別的名，使我們賴以得救的。}這裏他也說出了崇高的話。因為當目的不是為了成功辯論，而是僅僅顯示膽量時，他毫不留情；因為他並不害怕打擊得太深。他不只是簡單地說：\Quote{憑另一個，}而是說：\Quote{除祂以外，沒有救恩，}這就是說，祂能夠拯救我們。這樣他壓制了他們的威脅。\par

\BibleQuote{他们看见伯多禄和若望的胆量，又发现他们是没学问、无知的人，就惊讶，并认出了他们，知道他们是同耶稣一起的。}\BibleRef{宗 4:13}这两个没学问的人以他们的雄辩驳倒了他们和司祭长。因为不是他们在说话，而是圣神的恩宠。\BibleQuote{又看见那被治好的人同他们站在一起，就无话可驳。}\BibleRef{宗 4:14}那人的胆量何等大！即使在公议会中，他也没有离开他们。因为假如他们否认事实，他可在那里驳斥他们。\BibleQuote{但他们命他们从公议会出去，彼此商议说：我们当怎样处置这些人？}\BibleRef{宗 4:15}你看他们陷入何等困境，对人的畏惧又如何支配一切。正如基督的情形，他们（如俗语所说）既不能抹杀既成事实，也不能将其掩盖；他们越阻碍，信仰反而越发展；如今亦然。\Quote{我们该怎么办？}啊，何等的愚昧！竟以为那些尝过争战的人如今会畏惧；竟指望在起初证明无能后，在展现了如此雄辩的榜样后，还能做出什么新事！他们越想阻止，事情愈发在他们手中扩大。但他们说什么？\BibleQuote{因为确实，他们行了一件明显的奇迹，凡住在耶路撒冷的人都知道了，我们不能否认。但为叫这事不再在民间传播，我们当恐吓他们，叫他们今后不再以这名向任何人发言。他们遂叫了他们来，命令他们绝不可以耶稣的名发言或施教。}\BibleRef{宗 4:16}看这些人何其厚颜无耻，宗徒们又何其心胸伟大。\BibleQuote{伯多禄和若望回答说：在天主面前，听你们的而不听天主的，这是否合理，你们自己判断吧！因为我们不得不说我们所见所闻的事。他们于是又恐吓了一番，但因民众的缘故，找不到惩罚他们的方法，便释放了他们。}\BibleRef{宗 4:19}奇迹堵住了他们的口：他们甚至不容宗徒们把话说完，极其狂妄地打断了他们的话。\BibleQuote{众人因所发生的事都光荣天主。原来那被治好的人已有四十多岁了。}\BibleRef{宗 4:22}但让我们从头回顾一下所说的话。\par

\Quote{他们正向百姓说话等等，因为他们教导百姓，并藉着耶稣宣讲死人的复活，而感到烦恼。}（回顾，1至2节）如此，起初他们做一切是为了人的意见（或名声），但现在添加了另一个动机：即他们不应被认为有杀人罪，正如他们后来所说的：\BibleQuote{你们愿意把这人血归在我们身上吗？}\BibleRef{宗 5:28}啊，何等的愚昧！他们确信他已复活，并接受了这证据，竟指望那死亡不能拘捕的他，能被他们的阴谋掩盖！还有什么比这更愚昧！邪恶的本性就是如此：它对任何事都视而不见，却时时陷入慌乱。发现自己被压倒，他们感觉像被智胜的人：就像那些在某事上被抢先并受愚弄的人。然而他们到处宣称是天主使他复活；但他们说话却是\Quote{以耶稣的名}，表明耶稣已复活。\Quote{藉着耶稣，死人的复活}：因为他们自己也相信复活，一种冷漠幼稚的道理，但毕竟他们相信。单是这一点，难道不足以促使他们不对宗徒们做什么吗？——我是说，门徒们以那样的胆量行事时，难道这还不能阻止他们吗？说，犹太人哪，你为什么不信呢？你本应注意所行的奇迹和所说的话，而不是众人的恶意。\Quote{他们教导百姓。}因为因他们对基督所做的事，他们已在他们中声名狼藉；所以他们反而加重了自己的恶名。\BibleQuote{他们遂下手拿住他们，将他们收在监里，直到第二天，因为那时已是晚上。}\BibleRef{宗 4:3}然而，对于基督，他们却没有这样做；他们在半夜捉了他，立刻带去，毫不耽搁，因为他们非常害怕民众；但在这里对于宗徒们，他们却大胆了。他们不再把他们带到比拉多那里，因想起前事而羞愧脸红，免得他们也为那事受责罚。\par

\Quote{次日，他们的首领、长老和经师都聚集在耶路撒冷。}\BibleRef{宗 4:5} 再次在耶路撒冷：在那里人的血被倾倒；他们也不敬畏自己的城市；\Quote{以及亚纳斯和盖法，}等等。\BibleRef{宗 4:6} \Quote{亚纳斯，}经上说，\Quote{和盖法。}正是他的使女盘问伯多禄，而他无法忍受：就在他的屋里伯多禄否认了，那时另一个人被捆绑在那里：但现在，当他来到他们众人中间，看他如何说话！\Quote{你们凭什么名做这事？}你为什么不说它是什么，却把它隐藏起来？\Quote{你们凭什么名做这事？}\BibleRef{宗 4:7} 然而他断言，不是我们做的。\Quote{你们人民的首领，}等等。\BibleRef{宗 4:8} 注意他的智慧：他不直接说，\Quote{我们因耶稣的名做了这事，}而是怎样？\Quote{因祂的名，这人}——他不说，\Quote{被我们治好；}而是——\Quote{现在站在你们面前痊愈了。}又说，\Quote{若我们因向那软弱的人行了善事而被查问。}\BibleRef{宗 4:9} 他狠狠打击他们，说他们总是将此类行为定罪，挑剔向人行善的工作：并提醒他们以前的行为，即\Emph{他们}跑去杀人，不仅如此，还将行善定为犯罪。你是否也（在修辞上）注意到他们以何等尊严表达自己？甚至在用词上，他们因实践而变得老练，从此不再被击倒。\Quote{你们众人和全以色列百姓都应知道，}等等。\BibleRef{宗 4:10} 借此他表明，他们反而违背自己意志地宣讲基督；他们自己通过查问和审问来赞扬这教导。哦，极大的胆量——\Quote{你们钉死的！天主复活的！}——这更加大胆！不要以为我们隐藏了可羞耻的事。他几乎是嘲弄地说这话：不仅说，而且详细阐述。\Quote{这，}他说，\Quote{就是你们匠人弃而不用的石头；}然后他继续教导他们，补充说，\Quote{成了屋角的基石}\BibleRef{宗 4:11}；也就是说，这石头确实被认可了！他们现在因奇迹而有的胆量是巨大的。当需要教导时，看他们如何说话并引用许多预言；但当需要大胆发言时，他们只断然说话。因此，\Quote{除祂以外，}他说，\Quote{别无拯救，因为在天下人间，没有赐下别的名，使我们赖以得救。}\BibleRef{宗 4:12} 他说，这对所有人都是明显的，因为那名不是单单赐给我们的；他甚至以他们为证人。因为他们问，\Quote{你们凭什么名做这事？}\Quote{凭基督的名，}他说：没有别的名。你怎么问呢？这到处都是明显的。\Quote{因为天下人间没有赐下别的名，使我们赖以得救。}这是一个已经\OriginalTerm{弃绝}{\GreekText{κατεγνωκυίας}}了今生的灵魂所说的话。他极其坦率的态度在此证明，当他以卑微的言辞论及基督时，他这样做并非出于恐惧，而是出于\OriginalTerm{俯就}{\GreekText{συγκαταβαίνων}}的智慧：但现在时机合适，他不再说卑微的话：正是以此意图使他们惊慌。看哪，另一个奇迹，不亚于前一个。\Quote{他们看见伯多禄和若望的胆量，}等等。\Quote{又认出他们原是耶稣的门徒。}\BibleRef{宗 4:13} 圣史写下这段不是没有意义的；但在说\Quote{认出他们原是耶稣的门徒}时，他是指在祂受难时：因为那时只有这些人与祂在一起，而那时他们确实见过他们谦卑、沮丧：而这正是最让他们惊讶的：变化之大。事实上，亚纳斯和盖法及其同伴也在那里，这些人当时也曾站在祂旁边，而现在他们的胆量令他们惊异。\Quote{他们看见胆量。}因为不仅是他们的话语；他们的举止也显示了这一点：他们如此无畏地站在那里，在这样的罪案中受审，且面临着极端的危险！不仅通过言语，也通过姿态、眼神和声音，简而言之，通过他们的一切，都表明了面对百姓时的胆量。从他们所说的话，他们也许感到惊奇：\Quote{他们是没有学问的普通人：}因为一个人可能没有学问，却不是一个普通或平民；一个普通人，却不是没有学问。\Quote{他们看出，}经上说。从何而看？从他们说的话？伯多禄并没有长篇大论，但通过他的\OriginalTerm{方式与方法}{\GreekText{τῆς} \GreekText{ἀπαγγελίας} \GreekText{καὶ} \GreekText{τῆς} \GreekText{συνθήκης}}，他表明了自信。\Quote{他们认出他们原是耶稣的门徒。}这情况使他们相信，正是从耶稣他们学到了这些，并且他们以门徒的身份行事。但这些人声音并不亚于奇迹，奇迹本身就发出了声音：而这个标记本身堵住了他们的口。[\Quote{他们看见那人，}等等。]所以，他们会\OriginalTerm{断然要求}{\GreekText{ἐπέσκηψαν}}与他们决裂，如果那人不在他们中间。\Quote{我们不能否认它。}所以他们会否认它，如果事情不是这样：如果见证不是来自众人。\Quote{但恐怕这事越发传扬在民间。}\BibleRef{宗 4:17} 然而，这对所有人都是明显的！但邪恶的本性即是如此：它在各处蒙羞。\Quote{我们恐吓他们。}你说什么？恐吓？你们期望阻止宣讲？然而所有开头都是艰难而考验的。你们杀了主，却没有阻止它；现在，如果你们恐吓，你们就期望我们回转？监禁没有使我们屈从地说话，你们就能吗？\Quote{于是叫了他们来，吩咐他们，}等等（第18、19节）。放他们走对他们好得多。\Quote{伯多禄和若望回答说：你们自己斟酌：在天主面前，听从你们而不听从天主，是否合理。}当恐惧消退时（因为那命令等同于释放他们），那时宗徒们说话也更温和：他们远非单纯的好勇斗狠：\Quote{是否合理，}他说：\Quote{我们不能不说。在天主面前听从你们而不听从天主，是否合理。}\BibleRef{宗 4:20} 这里（在\Quote{天主}一词中）他们意指基督，因为是他命令了他们。他们再次证实了祂复活的事实。\Quote{因为我们不得不说我们所看见和听见的事。}这样我们就是有资格被相信的见证人。\Quote{于是他们更加恐吓他们。}\BibleRef{宗 4:21} 他们再次徒然恐吓。\Quote{于是他们释放了他们，因为找不到惩罚他们的理由，为了百姓，因为众人为所发生的事光荣天主。}于是百姓光荣天主，而这些人却试图毁灭他们：他们真是与天主作战的人！因此他们使他们更加引人注目和荣耀。\Quote{因为我的能力，}经上说，\Quote{在软弱中才显得完全。}\BibleRef{格后 12:9}\par

这些作为殉道者已经作了见证：他们置身于对抗众人的战斗中，说：\Quote{我们不能不说我们所看见所听见的事。}如果我们所说的是假的，你们就指责；如果是真的，为何阻拦？这就是哲学！那些人困惑，这些人喜乐；那些人蒙受极大的羞耻，这些人行事坦然无惧；那些人害怕，这些人有自信。试问，谁害怕？是那些说\Quote{免得它再在百姓中蔓延}的人，还是这些说\Quote{我们不能不说我们所看见所听见的事}的人？这些人有喜乐、言论自由、超越一切的欢愉；那些人却有沮丧、羞耻、恐惧，因为他们惧怕百姓。但这些并不怕那些人；相反，这些人说他们想说的，那些人却做不了他们想做的。谁在锁链和危险中？岂不正是这些吗？\par

让我们持守德行；不要让这些话仅止于喜乐和某种精神的提升。这不是剧场，不是为\OriginalTerm{歌手}{\GreekText{κιθαρῴδων}}、悲剧演员和\OriginalTerm{乐师}{\GreekText{κιθαριστῶν}}所设，那里的果实仅在于享受，而享受本身也随流逝的日子而逝。不，但愿仅是享受而已，而不是享受中还有祸害！但事实如此：每个人都将那里所见的许多东西带回家，像瘟疫感染一样粘在他身上：年轻者从那些撒旦的戏剧中随处采集一些歌曲片段，尽力记在脑海中，回家后哼唱；年长者，本应稳重而不至于轻浮，虽不如此，但却记住那里所说的一切——无论是讲道还是言辞本身；另一个人则记住一些污秽荒谬的小调。你从此处离开，什么也不带走。——我们立了一条法律——不，不是我们：天主禁止！因为经上说：\Quote{不要称地上的人为你们的师傅}\BibleRef{玛 23:8}；基督立了一条法律，禁止人发誓。现在，关于这条法律，做了什么？因为我不停地讲论它，如同宗徒所说：\Quote{免得我再来时，不得不严厉。}\BibleRef{格后 13:2}我问你们，你们将这事放在心上了吗？你们认真思考过吗？你们是认真对待，还是我们又要再次谈论同一主题？不，无论你们是否认真，我们都要再次谈论，好使你们认真思考，或者如果你们已经放在心上，可以更确定地遵行，并劝勉他人。那么，我请问，我们从哪里开始？从旧约开始？因为这也是我们的耻辱：我们本应超越法律，却连法律的规定也不遵守！因为我们本不应听到这类事情：这些是适应\OriginalTerm{犹太人的卑微}{\GreekText{τῆς} \GreekText{Ἰουδαϊκῆς} \GreekText{εὐτελείας}}之规条；\Emph{我们}本应听到那些完美的劝谕：\Quote{舍弃你的财产，勇敢地站立，为福音舍命，轻看世上所有美物，与今生无分；若有人亏待你，善待他；若有人欺骗你，祝福他；若有人辱骂你，尊敬他；超越一切。}（\Person{圣盎博罗削}，\WorkTitle{论职份}卷一第二章）这些和类似的话是我们本应听到的。但我们却在谈论关于发誓的事；我们的情况就如同一个人本应成为哲学家，却被从大师那里带走，让他去一个音节一个音节地拼读一样！试想，一个留着长须、手持杖杖、肩披斗篷的人，却去儿童学校，和他们做同样的功课，岂不是极其可笑？然而，我们的可笑之处甚至更大。因为犹太人的政体与我们的政体之间的差距，如同哲学与拼读课之间的差距吗？不，确实不是，而是如同天使与人的差距。试想，若有人从天上请来一位天使，并让他站在这里听我们的讲道，好像他必须遵行一样，岂不荒谬可笑？但如果像孩童一样学习这些事已是可笑，那么连这些事都不留意，该是何等的定罪，何等的羞耻！身为基督徒，却还要学习不可发誓！然而，我们忍受这个，免得陷入更大的可笑。\par

那么，今天我们就从旧约来对你们说话。它告诉我们什么？\BibleQuote{不可让你的口惯于起誓；也不可让你自己习惯于呼号圣者之名。}而且为什么呢？\BibleQuote{因为仆役常受责打，不会没有伤痕；同样，常起誓的人亦然。}\BibleRef{德 23:10}请看这位智慧人的明辨。他没有说\BibleQuote{不可让你的心惯于起誓}，而是\BibleQuote{你的口}；因为这完全是口舌之事，因此容易纠正。因为最终这成为不经意的习惯；例如，有许多人进入公共浴池时，一跨过门槛，便画十字\BibleRef{格前 11:7}（\OriginalTerm{画十字}{\GreekText{σφραγίζονται}}）。这是手已经习惯做的，无需任何人吩咐，出于习惯。同样，在点蜡烛时，常常心思专注于别的事，手却做出记号。同样，口也不经心思的同意，仅凭习惯说出那个词，整个事情就在舌头上。他说：\BibleQuote{也不要让你自己习惯于呼号圣者之名。因为仆役常受责打，不会没有伤痕；同样，常起誓的人亦然。}他这里不是指虚假的誓言，而是砍掉所有誓言，并且也给它们指定了惩罚。既然如此，发誓就是罪。因为灵魂实在是这样：满布所有这些溃疡、所有这些伤痕。但你们看不见吗？是的，这就是祸害所在；然而如果你们愿意，你们是能看见的，因为天主给了你们眼睛。先知就是用这样的眼睛看见的，他说道：\BibleQuote{我的创伤腐烂流脓，因了我的愚昧。}\BibleRef{咏 38:6}我们轻视了天主，我们憎恨了那个圣善之名，我们把基督踏在脚下，我们失去了所有敬畏，没有人以尊敬提到天主的圣名。然而，如果你爱一个人，甚至一听到他的名字就站起；但你却如此不断地呼求天主，却不以为然。为你的仇敌的好处呼求祂；为你自己灵魂的救恩呼求祂；那么祂就会临在，那时你会使祂喜悦；而现在你却激怒祂。要像斯德望那样呼求祂；他说：\BibleQuote{主，不要将这罪归于他们。}\BibleRef{宗 7:59}要像厄耳卡纳的妻子那样呼求祂，流着泪、抽泣着、祈祷着。\BibleRef{撒上 1:10}我不阻止你们，反而诚恳地劝你们这样做。要像梅瑟那样呼求祂，是的，他大声呼喊，为那些曾使他流放的人代祷。你若随意提及任何重要人物，便被视作侮辱；而你却在言谈中随意谈论天主，不分时宜？我不想阻止你们常将天主铭记于心：反之，这正是我所渴望和祈祷的，只要你们这样做是为了尊崇祂。若我们只在应该呼求的时候、为应该呼求的事呼求天主，这会为我们带来多大的益处！我还要问，为何宗徒时代行了如此多的奇迹，而我们的时代却没有？然而仍是同一天主，同一个圣名。但不，情况不同。因为那时他们只为我说过的那些目的呼求祂；而我们呼求祂不是为了这些，而是为了完全不同的目的。——如果有人不肯相信你，因此你起誓，你就对他说：\Quote{请相信我：}然而，如果你非要起誓，就指着你自己起誓吧。我说这话，并非要设立一条反对基督律法的法律；绝对不可！因为经上说：\BibleQuote{你们的\InnerQuote{是}就是\InnerQuote{是}，\InnerQuote{非}就是\InnerQuote{非}}\BibleRef{玛 5:37}：但这只是迁就你们目前的水平，好使我更容易引导你们实行这条诫命，使你们脱离这个暴虐的习惯。有多少在其他方面做得好的人，被这些行为毁灭了！要不要我告诉你们，为什么古人被允许起誓？（因为虚假的誓也不允许他们。）因为他们指着偶像起誓。但你们难道不羞愧于停留在那些因软弱而被引导向更好事物的法律上吗？的确，当我带领一个外邦人时，我不会立即对他下这条禁令，而是首先劝他认识基督；但若信徒，既已学习并听到了基督，却仍需要与外邦人同样的宽容，那么（他的基督徒身份）有什么用处，有什么益处呢？——但习惯很强，你无法摆脱它？那么，既然习惯的暴政这么大，就把它转移到另一个渠道吧。怎么做呢？你会问。我常说的话，现在也说了：要有许多的\OriginalTerm{监督者}{\GreekText{λογισταὶ}}，要有许多的审查者和检验者（\OriginalTerm{审查者}{\GreekText{ἐξετασταὶ}}，\OriginalTerm{检验者}{\GreekText{δοκιμασταί}}）。比如说，如果你偶然把外衣穿反了，你允许你的仆人纠正你的错误，并羞于向他学习，尽管这本身很可耻；而这里，当你的灵魂受到伤害时，你难道羞于接受他人的教导吗？你容忍仆人整理你的衣服、系你的鞋，却不愿容忍那愿意整理你灵魂的人吗？让你的仆人、你的孩子、你的妻子、你的朋友、你的亲属、你的邻人，在这件事上做你的老师吧。因为当一只野兽从四面八方被追捕时，它不可能逃脱；同样，一个人有如此多的人监视他、如此多的人责备他、从四面八方受到打击，他不可能不保持警惕。第一天他会觉得难以忍受，第二天、第三天也是如此；但之后就会变得容易，过了第四天，甚至就没什么可做的了。如果你怀疑我，就试试看；我恳求你考虑一下。犯错和改正都不是小事；两方面都有善恶的重大意义。愿这善事得以成就，借着我们的主耶稣基督的恩宠和慈爱，愿光荣、权能和尊崇归于祂，与圣父及圣神，自今至永远，世世无尽。阿们。\par

\SourceRecord{Translated by J. Walker, J. Sheppard and H. Browne, and revised by George B. Stevens.；Nicene and Post-Nicene Fathers, First Series；Vol. 11.；Edited by Philip Schaff.；Buffalo, NY: Christian Literature Publishing Co.,；1889.；<http://www.newadvent.org/fathers/210110.htm>.}

% structure: p0001 source=h1 target=section reason=page-title

\section{\WorkTitle{宗徒大事录讲道集第十一篇}}

\BibleRef{宗 4:23}\par

\Quote{他们被释放后，便回到自己人那里，将司祭长和长老所说的一切，都报告了。}\par

他们讲述这事，并非为了自己的光荣——他们怎能存这样的动机？——而是展示了其中所彰显的基督恩宠的证据。他们将敌人所说的一切都报告了；很可能省略了自己的部分：这使得听者更加勇敢。然后呢？他们再次逃向真正的援助，向那无敌的联盟，并且再次，\BibleQuote{同心合意。当他们听见了这事后，}经上说，\BibleQuote{他们就同心合意地高声向天主说：}\BibleRef{宗 4:24}并且极其恳切，因为这不是随便的祈祷。请注意他们的祈祷是如何恰如其分地构架的：当他们恳求指明谁适合宗徒职分时，他们说：\BibleQuote{主啊，祢知道众人的心，求祢指示。}\BibleRef{宗 1:24}因为那需要预知。但在此处，需要的是堵住敌人的口，他们就谈及主权的领域；因此他们这样开始：主，\BibleQuote{（\OriginalTerm{主人}{\GreekText{Δέσποτα}}）创造天地、海和其中万物的天主，祢曾藉着圣神，藉祢的仆人，我们的祖先达味的口说：\InnerQuote{外邦人为什么叛乱？万民为什么妄想？世上的君王起来，首领聚集在一起，反对主，反对祂的基督。}}\BibleRef{宗 4:24}这可以说是依据祂自己的盟约向天主申诉，他们这样提出预言；同时也安慰自己，认为敌人的一切妄想都是徒劳的。这就是他们所说的：使这些话实现，并显明他们\BibleQuote{妄想虚妄的事。}——事实上，他们接着说，\BibleQuote{在这城里，果然聚集了黑落德和般雀比拉多，同外邦人和以色列民众，反对祢的圣童耶稣（\OriginalTerm{孩童}{\GreekText{Παἵδα}}），是祢所傅油的，为要实行祢的手和祢的旨意所预定要成就的一切。现在，主，请看他们的恐吓。}\BibleRef{宗 4:27}请留意他们的\OriginalTerm{博大胸怀}{\GreekText{φιλοσοφίαν}}。这些并非咒诅的话。说\Quote{他们的恐吓}，他们并非特指某件具体被恐吓的事，而只是泛指恐吓这一事实，也许是视为可怕。事实上，作者叙述简洁。请注意，他们不说：\Quote{压碎他们，推倒他们；}而是说什么？\BibleQuote{求祢赐予祢的仆人们，使他们能放胆讲论祢的圣言。}我们也当这样学习祈祷。然而，当一个人落在那些蓄意杀害他并以此相威胁的人中间时，他会多么充满愤怒？多么充满敌意？但这些圣人们并非如此。\BibleQuote{伸出祢的手来医治，并且藉祢圣童耶稣的名字显征兆和奇迹。}\BibleRef{宗 4:30}如果那名字行了大事，胆量就会大增。\par

\Quote{他们祈祷后，他们聚集的地方震动起来。} \BibleRef{宗 4:31} 这是他们蒙垂听和祂眷顾的证明。 \Quote{他们都充满了圣神。} 是什么意思？\Quote{他们充满了？}意思是，他们被点燃了，恩赐在他们内燃烧。 \Quote{他们大胆地宣讲天主的圣言。相信了的人都是一心一灵。} \BibleRef{宗 4:32} 你看见他们与天主的恩宠一同贡献了自己的部分吗？因为处处都该留意，他们与天主的恩宠一同做了自己的部分。正如伯多禄上面说的：\Quote{我没有金银}；又说，\Quote{他们都在一处。} \BibleRef{宗 3:6} 但在此处，提到他们蒙垂听后，圣史接着也讲述他们所显示的德行。此外，他正要进入撒斐辣和阿纳尼雅的叙述，为显示那一对的恶劣行为，他先论述其余人的高贵行为。请问，是他们的爱德产生了贫穷，还是贫穷产生了爱德？依我之见，爱德产生了贫穷，然后贫穷拉紧了爱德的绳索。因为看他说的：\Quote{他们都一心一灵。} 你看，心和灵造就了\Quote{一起。} \Quote{他们中没有一个人说自己的财物是属于自己的，一切都归公用。宗徒们以大能力为复活作证（\OriginalTerm{偿付}{\GreekText{ἀπεδίδουν}}）。} \BibleRef{宗 4:33} 这个说法表示他们如同受托之人：他视之为债务或义务，也就是说，他们放胆地作证，即向所有人偿付了见证。 \Quote{众人都大蒙恩宠。他们中没有一个缺乏的。} \BibleRef{宗 4:34} 他们的感觉如同在父家的屋顶下，大家暂时平等共享。不能说，虽然他们供养了其余人，却仍感到供养所用的资财是自己的。不，可钦佩的情况是：他们先让渡了财产，然后供养其余人，为的是供养不是出于私有的资财，而是出于公共财产。 \Quote{凡拥有田地或房屋的都卖了，把卖得的价钱带来，放在宗徒们脚前，照每人所需分配。} \BibleRef{宗 4:35} 这是一个极大的尊敬标记，即\Quote{他们放在宗徒们脚前。若瑟，宗徒们给他起名叫巴尔纳伯（解说劝慰之子）。} \BibleRef{宗 4:36} 我不认为这人与玛弟亚的同伴是同一个人；因为那人也叫做犹斯托和[巴尔撒巴，但这人是若瑟和]\Quote{巴尔纳伯}[\Quote{劝慰之子}]。我想他因德行而得此名，因为他适合并胜任此职。 \Quote{他是肋未人，生于塞浦路斯。} 请注意作者处处如何指出法律的中断。但他怎么又是\Quote{塞浦路斯人？}因为那时他们甚至移居到其他国家，仍被称为肋未人。 \Quote{他有田地，卖了，把价钱带来，放在宗徒们脚前。}\par

现在让我们再看一遍前面所说的话。\Quote{他们被释放后，回到自己的人那里，将司祭长和长老们所说的一切都报告了。} \BibleRef{宗 4:23} 请看宗徒们谦逊的作风和宽广的心胸。他们并不四处夸耀，说：\Quote{我们怎样\OriginalTerm{对待}{\GreekText{ἀπεχρησάμεθα}}了司祭！} 他们也不贪求荣誉；正如我们读到，他们回到自己的人那里。请注意他们如何不将自己投入试探，而是在试探出现时勇敢地忍受它们。如果是别的门徒，也许会被群众的支持所鼓动，出言侮辱，说出许多刻薄的话来。但这些真正的哲人却不如此；他们一切都以温和与良善而行。\Quote{他们听了，就同心合意地高声向天主呼求。} \BibleRef{宗 4:24} 那呼喊出自喜悦和极大的激动。这样的祈祷确实有效，充满真正的智慧，为这样的目标，由这样的人，在这样的场合，以这样的方式献上；而其他所有的祈祷都是可憎和亵渎的。\Quote{主，祢是造天地、海和其中万物的天主。} \BibleRef{宗 4:24} 请注意他们如何不说闲话，不说老妇的无稽之谈，而是谈论祂的能力。正如基督亲自对犹太人说：\Quote{我若靠天主的圣神驱魔，} 看，圣父也藉着圣神说话。因为经上说什么？\Quote{主，天主，祢曾藉着圣神，通过祢的仆人、我们的祖先达味的口说：外邦人为什么愤怒？} \BibleRef{宗 4:25} 圣经惯于这样以多数来描述一位。\Quote{实在，主，黑落德和般雀比拉多，联合反对祢的圣子耶稣，祢所傅油的那一位，等等。} \BibleRef{宗 4:27} 请注意，即使在祈祷中，他们也详细描述受难，并将一切都归于天主。——也就是说，不是他们有能力做这事：而是祢成就了一切，祢是那允许的、追问的，却又使之实现的，祢是全能和智慧的，祢利用祢的仇敌为祢自己的喜悦服务。\BibleRef{宗 4:28} \Quote{因为祢的手所预定要成就的一切，} 等。这里他们谈论祂的无穷技巧、智慧和能力。因此，他们如同敌人聚集，怀着杀人的目的，并且反对自己，但他们却做了祢所愿意的事：正如经上所说：\Quote{因为祢的手和祢的旨意所预定要成就的事，} \Quote{祢的手}是什么意思？在我看来，他在这里似乎用能力和旨意表示同一件事，意思是对于祢，只要有愿意就够了：因为人并不是凭能力来决定。\Quote{凡是祢的手}等等，即凡是祢所预定的；要么是这个意思，要么是祂用手实现。\Quote{现在，主，求祢垂顾他们的恐吓。} \BibleRef{宗 4:29} 正如那时，他们\Quote{图谋虚妄的事，} 同样\Quote{现在，} 求祢使他们的图谋成为虚空：即不要让他们恐吓得逞。他们这样说，不是因为自己躲避艰难，而是为了宣讲的缘故。因为他们不说：\Quote{救我们脱离危险，} 而说什么？\Quote{赐给祢的仆人们，使他们能放胆讲论祢的圣言。} 祢既然成就了先前的计划，也求祢使这些得以实现。请注意，他们如何确认天主是他们胆量的根源；他们如何为天主而求一切，不为自己的荣耀或野心求任何事物。他们自己承诺不会丧胆；但他们祈求，\Quote{藉着伸出祢的手来医治，并且使神迹和奇事得以施行。} \BibleRef{宗 4:30} 因为没有这些，无论他们表现出多大的热忱，都将徒劳无功。天主应允了他们的祈祷，并藉着震动那地方来表明。因为经上说：\Quote{他们祈祷后，那地方就震动了。} \BibleRef{宗 4:31} 至于为什么发生这事，你们要从先知的话中听，他说：\Quote{祂注视大地，大地便震动。} \BibleRef{咏 103:32} 因为藉此祂表明祂临在于他们的祈祷中。另一位先知又说：\Quote{大地震动颤抖，因上主的临在。} \BibleRef{咏 17:7} 天主这样做，既为使这更可畏，也为带领他们走向勇敢的信赖。\Quote{他们都充满了圣神，并放胆讲论天主的圣言。} 他们更加勇敢了。因为这是（他们工作的）开始，他们曾恳求一个可感知的标记来使他们确信（\OriginalTerm{为使他们相信}{\GreekText{πρὸς} \GreekText{τὸ} \GreekText{πεισθῆναι} \GreekText{αὐτούς}}）——但此后我们再也没有发现类似的事情发生——因此他们所得的鼓励是巨大的。事实上，除了藉着奇迹的标记，他们无法证明祂已复活。所以他们寻求的不仅是自己的\OriginalTerm{保证}{\GreekText{ἀσφάλειαν}}，而且是不致蒙羞，能放胆宣讲。\Quote{那地方震动了，} 这使他们更加坚定不移。因为这有时是愤怒的标记，有时是恩宠和眷顾的标记，但在目前的情况下，是愤怒的标记。例如，在十字架上，大地震动；主亲自说：\Quote{那时将有饥荒、瘟疫和各地地震。} \BibleRef{玛 24:7} 但那时它所象征的愤怒是针对敌人的；至于门徒，则使他们充满圣神。请注意，就连宗徒们，在祈祷后，也\Quote{充满了圣神。} \Quote{众信徒的心意}等等。\BibleRef{宗 4:32} 你看到，这件事的德能有多大，因为在那个团体中也需要这（恩宠）。因为这是一切美善的基础，他现在第二次提到这一点，劝勉众人轻视财富：\Quote{他们中没有人说自己的东西是自己的，} \Quote{而是凡物公用。} 因为这不单是奇迹标记的结果，也是他们自己意愿的果效，从撒斐辣和阿纳尼雅的事例可以清楚看出。\Quote{宗徒们以大德能作证}等等。\BibleRef{宗 4:33} 宗徒们展示复活的见证，不是言语，而是以德能：正如保禄所说：\Quote{我的宣讲不是凭人的智慧所劝诱的言词，而是凭圣神和德能的显示。} 不但是\Quote{以德能，} 而且是\Quote{以大德能。} \BibleRef{格前 2:4} 经上说：\Quote{大恩宠降在他们众人身上，因为他们中没有一个人缺乏。} \BibleRef{宗 4:34} 这就是为什么恩宠（在他们众人身上），因为\Quote{没有一个缺乏的。} 也就是说，由于施予者的极大热忱，没有人缺乏。因为他们不全给一部分、又保留一部分；即使全部给出，也不当作是自己的财产。而且他们生活丰裕：他们消除了彼此间的一切不平等，建立了美好的秩序。\Quote{凡拥有田产房屋的，} 等等。他们极其恭敬地这样做：不敢直接交给他们手里，也不炫耀地献上，而是放在宗徒们脚前。他们让宗徒们作分施者，使他们成为所有者，从此一切费用都从公共财产中支出，而非私人财产。这也帮助他们抵御虚荣。如果现在这样做，我们，无论贫富，都会过得更快乐，而且穷人不会比富人更快乐。如果你们愿意，让我们现在用语言描述一番，至少从中得到这份乐趣，因为你们无意在行动上实行。因为无论如何，即使从当时发生的事实也可清楚看出，他们变卖财产，并没有陷入匮乏，反而使那些匮乏的人富足。但让我们现在用语言描绘这种状况，让所有人都变卖财产，放入公共基金——我是指用语言说；无论是富人或穷人，都不要激动。你们认为会收集多少钱金？就我而言，我推测——当然不可能精确地说——假设这里所有的人，男人和女人，都倒出他们所有的财产、田地、资产、房屋——因为我不提奴隶，因为那时没有这样的事情，但毫无疑问，曾是奴隶的都被释放了——那么也许能收集到十万磅重的黄金；不，或许两倍或三倍。你们想想；我们的城市按多少\Quote{ \OriginalTerm{轭}{juga} }估税？我们说有多少人口是基督徒？我们说十万，其余是希腊人和犹太人？那么能收集到多少千磅黄金啊！穷人的数目是多少？我想不会超过五万。那么每天养活这些人，该多么丰裕啊！即使食物是由公共供给，大家一同用餐，毕竟也不需要那么大的开销。但你们会问，钱花完后我们怎么办？你们以为钱会用完吗？天主的恩宠岂不更多一万倍？天主的恩宠岂不确实丰厚地倾注？难道我们不会使之成为地上的天堂吗？如果当人数只有三千和五千时，这样做就取得了如此辉煌的成功，且他们中没有一人抱怨贫穷，那么在如此庞大的群众中，这岂不更荣耀吗？甚至那些外教人，谁不会捐献呢？——但是，为表明分开生活是昂贵且导致贫穷的，假设一所房子里有十个孩子：妻子和丈夫，一个做羊毛活，另一个从户外工作带来收入；现在告诉我，他们怎样花费最多？是共同用餐、同住一所房子，还是分开居住？当然是分开居住。因为如果十个孩子必须分开住，他们就需要十个不同的房间、十张桌子、十个仆人，而收入也按比例增加开支。不正是因为这个原因，当有许多仆人时，他们都共用一张桌子，这样花费就不会那么大？因为事情就是这样，分裂总是导致减少，和睦一致带来增加。隐修院中的居民就像当时的信徒那样生活：他们中有谁曾饿死吗？有谁不曾得到一切丰盛的供给吗？现在，人们似乎更害怕这个，甚于坠入无底深渊。但如果我们实际尝试过，那么我们就会大胆地实行\OriginalTerm{这个计划}{\GreekText{τοῦ} \GreekText{πράγματος}}。你想，又将有何等的恩宠啊！因为如果在那时，除了三千和五千人之外没有信者，当时全世界都是敌人，他们无处寻求安慰，却仍勇敢地实行这个计划并取得如此成功；那么现在，因天主的恩宠，信者遍布天下，岂不更应如此？那时还有什么外邦人留下？就我而言，我想一个也不会有：我们必会如此吸引众人，将他们引领到我们这里来。但只要我们取得良好的进展，我信赖天主，连这一切也必会实现。只管照我说的去做，让我们按部就班地成功实现这些事情；如果天主赐予生命，我相信我们不久就能引导你们归向这种生活方式。\par

首先，关于那发誓的律法：实行它，使之坚固；让持守的人指出未持守的人，并追究责任，严厉斥责他。因为那（\Emph{参见上文}，第八篇讲道）指定的时间（\OriginalTerm{指定时间}{\GreekText{ἡ} \GreekText{προθεσμία}}）已经临近，我正在此事上进行查问，凡发现有罪者，我要放逐并排除。但愿天主不许我们中间有任何这样的人；反而但愿显明，所有的人都严格遵守了这属灵的口令。如同在战争中，通过口令区分友敌，现在也当如此；因为如今我们也确实在作战；好让我们知道谁是真正的弟兄。因为我们无论在本地或异乡，以此为标志是何等美善！这是何等武器，直攻魔鬼的头顶！不能发誓的口，很快就能在祈祷中与天主交往，并给魔鬼致命一击。不能发誓的口也必不会使用辱骂的言语。将这把火从你的舌头上赶出去，如同从房屋里赶出：把火拖出去。让你的舌头稍作休息：减轻这疮毒的恶性。是的，我恳求你们这样做，好让我能继续给你们另一个教诲：因为只要这件事没有做好，我就不敢转到别的事上。让这教训被完美地掌握，你们将意识到这成就，然后我会带你们认识其他律法，更确切地说，不是我自己，而是基督。在你们灵魂中植入这善事，渐渐地你们将成为天主的乐园，远胜于古时的乐园。你们中间没有蛇，没有致命的树，也没有诸如此类的东西。深深固定这习惯。如果这事成就，不仅你们在场的会受益，全世界的人也会受益；不仅他们，后世的人也会受益。因为一旦善习进入，并为大家所持守，就会传到久远的世代，任何情况都不能将其抹去。如果那在安息日拾柴的人被石头砸死——那么做比拾柴更可恶之事的，堆积罪孽的人，因为如此多的誓言，他将遭受什么？他还有什么不能忍受的？如果这事在你们手中顺利成就，你们将从天主那里得到巨大的帮助。如果我说：不要辱骂，你们会立即向我申辩你们的愤怒；如果我说：不要嫉妒，你们会提出别的借口。但在这种情况下，你们无话可说。因此，我从容易的规诫开始，这也是一切技艺中的惯例。因此，人通过先学习那些容易的，进而达到更高的责任。你们将看到这有多么容易，当靠天主的恩宠在这件事上成功之后，你们将接受另一条规诫。\par

请赐我能力，在外邦人和犹太人面前，尤其是在天主面前，大胆说话。是的，我以那爱情，以我为你们生产所受的产痛恳求你们。\Quote{我的孩子们}。我不再加下面的话：\Quote{我再次为你们受产痛}；也不说：\Quote{直到基督在你们内形成}。\BibleRef{迦 4:19}。因为我深信基督已在你们内形成。我将用其他的话对你们说：\Quote{我的弟兄们，我亲爱的和怀念的，我的喜乐和我的冠冕}。\BibleRef{斐 4:1}。请相信我，我不会再用别的话。如果此刻有千万顶镶嵌珠宝的王冠戴在我头上，它们也不能给我像你们在圣德上成长那样的喜乐；甚至我认为，君王本人也没有我对你们所怀的那种喜乐。让他战胜了他所交战的所有民族凯旋归来，让他除了本国王冠之外还赢得了许多其他冠冕，并接受其他王冠作为胜利的标志：我不认为他会为他的战利品而欢喜，像我因你们灵魂的进步而欢喜那样。因为我欣喜，仿佛我头上戴着千顶王冠；我理当欢喜。因为如果靠天主的恩宠你们成就了这个善习，你们将赢得一千场远比他更艰难的战役：通过摔跤和与恶毒的魔鬼、凶恶的邪灵战斗，用的是舌头，而不是刀剑，而是意志。因为请想，如果你们成功，将获得多少益处！你们已经根除了：第一，可恶的习惯；第二，恶念，一切邪恶的根源，即认为事情无关紧要、不会带来伤害的看法；第三，愤怒；第四，贪吝；因为这一切都是发誓的产物。不仅如此，你们由此将在通往所有其他美德的路上获得稳固的立足点。因为如同儿童学习字母时，他们不仅学习字母，而且借助它们逐渐学会阅读；你们也将如此。那恶念不会再欺骗你们，你们不会再说：这无关紧要；你们不再只是出于习惯说话，而是勇敢地对抗一切，以便在各方面成全那合乎天主的德行，你们可以收获永生的福分，借着祂独生子的恩宠和慈爱，愿光荣、权能和尊威归于祂，偕同圣父和圣神，自今而始，至永远，世世无尽。阿们。\par

\SourceRecord{Translated by J. Walker, J. Sheppard and H. Browne, and revised by George B. Stevens.；Nicene and Post-Nicene Fathers, First Series；Vol. 11.；Edited by Philip Schaff.；Buffalo, NY: Christian Literature Publishing Co.,；1889.；<http://www.newadvent.org/fathers/210111.htm>.}

% structure: p0001 source=h1 target=section reason=page-title

\section{宗徒大事录第十二篇讲道}

\BibleRef{宗 4:36-37}\par

\BibleQuote{有位若瑟，宗徒们称之为巴尔纳伯，解释是\InnerQuote{安慰之子}，是肋未人，生于塞浦路斯。他有田地，卖了以后，把银钱带来，放在宗徒们脚前。}\par

作者现在将要叙述阿纳尼雅和撒斐辣的事。为了显示那个人的罪是何等严重，他首先提到那行善事的人；这样，既然有如此多的群众都在行同样的事，有如此大的恩宠，如此大的奇迹，他却不为这些所教导，反被贪婪蒙蔽，招致毁灭在自己头上。\BibleQuote{他有田地——意思是这是他所有的一切——卖掉了，把银钱带来，放在宗徒们脚前。但是一个名叫阿纳尼雅的人，与他的妻子撒斐辣，卖了一块田地，从银钱中保留了一部分，他的妻子也同意这事，只带来一部分，放在宗徒们脚前。}\EditorialNote{见：宗 5:1-2}加重情节在于，这罪是共谋的，也没有别人看见所做的事。这个不幸的人怎么会想到犯这罪的？\BibleQuote{但伯多禄说：阿纳尼雅！为什么撒殚充满了你的心，使你欺骗圣神，并保留一部分田地的钱呢？}\BibleRef{宗 5:3}请同样观察，这里行了一个大奇迹，比前者大得多。\BibleQuote{当它还留着时}他说，\BibleQuote{难道不属于你吗？卖了以后，不是还在你权下吗？}\BibleRef{宗 5:4}这就是说，\Quote{有什么义务和强制吗？我们强迫你违反你的意愿吗？}\BibleQuote{你为什么心里起了这事？你不是欺骗人，而是欺骗天主。阿纳尼雅一听这话，就跌倒，断了气。}\BibleRef{宗 5:5}这个奇迹比瘸子的奇迹更大，在于死亡被施加，以及知道心里的念头，甚至知道私下所做的事。\BibleQuote{所有听见这事的人，都十分害怕。少年人就起来，把他包裹起来，抬出去埋葬了。大约过了三个小时，他的妻子不知道所发生的事，进来了。伯多禄回答她说：你告诉我，你们卖田地的钱就是这么多吗？}\BibleRef{宗 5:6}他愿意挽救那女人，因为男人是这罪的主谋；因此他给她机会为自己辩解，并给她悔改的机会，说：\BibleQuote{\InnerQuote{你告诉我，你们卖田地的钱就是这么多吗？}她说：\InnerQuote{是，就是这么多。}伯多禄便对她说：\InnerQuote{你们怎么约定试探圣神呢？你看，埋葬你丈夫的人的脚已到门口，他们也要把你抬出去。}她立刻跌倒在他脚前，断了气；少年们进来，见她死了，就抬出去，埋在她丈夫旁边。全教会和所有听见这些事的人，都十分害怕。}\BibleRef{宗 5:9}\par

在恐惧降临他们之后，他行了更多的奇迹；伯多禄和其他宗徒都一样；\BibleQuote{靠宗徒的手，在百姓中行了许多征兆和奇事；他们都同心合意地在撒罗满廊下。其余的人没有一个敢与他们联合，}即，与宗徒联合；\BibleQuote{但百姓却尊重他们，}即犹太百姓。如果\BibleQuote{没有一个人敢与他们联合}，宗徒们，\BibleQuote{却}有\BibleQuote{更多的人信了主，男女信徒的数目增加了很多，以致他们把病人抬到街上，放在床和褥子上，好让伯多禄经过的时候，至少他的影子能遮住他们。}\BibleRef{宗 5:12}因为伯多禄是那奇妙的一位，人们更注意他，既因为他公开的演讲——第一篇、第二篇和第三篇——也因为他所行的奇迹；因为他行了奇迹——第一次、第二次、第三次：因为当前的奇迹是双重的：第一，揭露内心的思想；第二，因他一句话就施以死亡。\BibleQuote{好让伯多禄经过的时候，至少他的影子能遮住他们}等等。这在基督的历史中尚未发生过；但请看，主所告诉他们的如今真正实现了，即\BibleQuote{凡信我的，我所做的事业，他也要做，并要做出比这些更大的事业。}\BibleRef{若 14:12}\BibleQuote{还有许多人从四周的城市来到耶路撒冷，带着病人和受污灵折磨的人，他们都得到了痊愈。}\BibleRef{宗 5:16}\par

现在，我要你们注意他们的整个生活是如何交织的。首先，因为基督被夺去而感到沮丧，然后因圣神降临而喜乐；再次，因嘲笑者而沮丧，然后因他们自己的辩护结果而喜乐。在这里我们又同时发现沮丧和喜乐。因他们变得引人注目，天主向他们启示而喜乐；因他们砍掉了自己的一些同伴而悲伤。再一次：成功中有喜乐，又因大司祭而悲伤。整件事中我们看到都是如此。古代圣人也同样如此。——但让我们重新回顾所说的话。\par

\Quote{他们变卖了}，经上写着，\Quote{把价银拿来，放在宗徒脚前}。\BibleRef{宗 4:34}看，我亲爱的弟兄们，他们不是让宗徒去卖，而是自己变卖，把价银交到他们手中。\Quote{但是有一个人，名叫阿纳尼雅}，等等。\BibleRef{宗 5:1}这段历史也极其有力地触动主教们。阿纳尼雅的妻子知道所做的事，因此伯多禄审问她。但也许有人会说，他对她处理得极其严厉。你说什么？什么严厉？如果因为捡柴而被石头砸死，那么因亵渎圣物更应如此；因为这笔钱已成为圣物。选择变卖自己的财物并分施，然后又撤回的人，是犯亵渎圣物之罪。如果从自己的财物中取回是亵渎，那么从别人的财物中取走更是如此。不要以为因为现在后果不同，罪行就会免于惩罚。你看见了吗？这就是对阿纳尼雅的控告：他把钱变成圣物后，又私自藏起来。你岂不能卖了田地后，任意使用价银吗？有人禁止你吗？你为什么在许诺之后又这样做？请看，在起初，魔鬼就发动了攻击；在如此多的征兆和奇事之中，这个人怎样心硬！在旧约里也曾发生过类似的事。加尔米的儿子贪图当灭之物；因为请看那里，对罪的惩罚是怎样的。亵渎圣物，亲爱的，是一项极其严重的罪行，是侮辱和轻蔑。宗徒说：我们没有强迫你卖，也没有强迫你卖后交出你的钱；你自愿做的，为什么从圣库里偷窃？\Quote{为什么}他说，\Quote{撒殚充满了你的心？}\BibleRef{宗 5:3}如果撒殚做了这事，为什么这人要为此负责？因为容许了魔鬼的影响，并被它充满。你会说，他们本应纠正他。但他不会接受纠正；因为见了如此多的事仍然无动于衷，那么任何其他事也必定无济于事；这件事不能简单地放过；像坏疽一样，必须割除，以免感染身体其余部分。实际上，这人自己也受益，因为未让他进一步作恶；其余的人也受益，因为他们变得更殷勤；否则就会相反。接着，伯多禄证明他有罪，表明这事没有瞒过他，然后宣布判决。但是，你为什么、出于什么目的做这事？你想保留它？你自始至终本该保留它，而从未声称要奉献。亲爱的，亵渎圣物是严重的。另一个人可能贪图不属于自己的东西；但保留自己的东西是你自由决定的。你为什么先把它圣化，然后又取走？你是出于极端的轻蔑做的。这行为不容宽恕，无可辩护。——因此，如果现在也有亵渎圣物的人，不要让他们成为绊脚石。如果那时就有这样的人，那么现在当邪恶众多时，更会有。但让我们\Quote{在众人面前斥责他们，使其余的人也有所惧怕。}\BibleRef{弟前 5:20}犹达斯是亵渎圣物的，但这并未成为门徒们的绊脚石。你看见从贪爱金钱生出多少邪恶？\Quote{所有听见这事的人，都甚惧怕。}\BibleRef{宗 5:5}那人受了惩罚，别人因此受益。并非无因。然而，先前也曾行过奇事：确实，但没有如此惧怕。所以那句话是真实的：\Quote{上主藉执行审判，被人认识。}\BibleRef{咏 9:16}同样的事曾在约柜的事上发生：乌匝受罚，其余的人惧怕。\BibleRef{撒下 6:7}但在那件事上，君王因惧怕而将约柜移开；但在这里，门徒们变得更加谨慎。[\Quote{大约隔了三个时辰，他的妻子不知道所发生的事，便进来了}，等等。]\BibleRef{宗 5:7}但请看，伯多禄没有派人去叫她，而是等她进来；也没有任何人敢于传出消息。这样的教师可畏，这样的门徒敬畏，这样的服从！\Quote{间隔三个时辰}——那妇人还没有听到，在场的人也无人报告，尽管时间足够让这事传开；但他们惧怕。福音史家甚至以惊叹的方式叙述这一情况，说：\Quote{她不知道所发生的事，就进来了。}\Quote{伯多禄回答她说，}等等。\BibleRef{宗 5:8}而她甚至可以从这一点看出伯多禄知道秘密。因为为什么不问别人，却问你？难道不是因为知道才问吗？但她如此心硬，不肯设法逃避罪责；她极其自信地回答，因为她以为只是在跟一个人说话。这罪的加重在于他们同心合意，如同事先约好。\Quote{你们为什么同心试探主的圣神？}他说，\Quote{埋葬你丈夫的人的脚已到门口。}\BibleRef{宗 5:9}他先让她认识罪，然后表明她将公正地遭受与丈夫同样的惩罚，因为她行了同样的恶：\Quote{他们必把你抬出去。她立刻仆倒在他脚前}，因为她正站在他旁边，\Quote{便断了气。}\BibleRef{宗 5:10}他们完全是自招惩罚！此后谁不敬畏？谁不怕宗徒？谁不惊奇？谁不恐惧？\Quote{他们全都同心合意地聚集在撒罗满廊下，}\BibleRef{宗 5:12}不再在屋里，而是占据了圣殿，在那里度日！他们不再防备触摸不洁之物，甚至毫无顾忌地处理死人。请看，他们对自己人严厉，对外人却不行使权力。\Quote{百姓，}他说，\Quote{都尊重他们。}\BibleRef{宗 5:13}既然他提到他们\Quote{在撒罗满廊下}，为使你不惊异群众如何容许这事，他告诉我们，他们甚至不敢接近他们：因为\Quote{没有一个人敢接近他们。}但是\Quote{信主的人越发增加，男女众多。甚至有人把病人抬到街上，放在床上或褥子上，希望伯多禄走过来时，至少他的影子能遮住一些人。}（14–15节）大信德，超过在基督身上所显示的。这是怎么回事？因为基督曾宣告：\Quote{因为我往父那里去，你们要作比这更大的事。}\BibleRef{若 14:12}这些事是百姓做的，而宗徒们留在那里，没有从一处移到另一处；也从别处把所有病人都用床和褥子抬来；从四面八方，新的惊奇归于他们：来自相信的人，来自被医治的人，来自受罚的人，来自他们向那些人（他们的对手）说话的大胆，来自信徒的美善行为：因为产生这效果的不单是奇迹。虽然宗徒们谦逊地将这一切归于这个原因，宣称他们因基督的名做这些事，但同时这些人的生活和高尚行为也帮助产生了效果。\Quote{信主的人越发增加，男女众多。}请注意，他现在不再告诉相信的人数；信仰如此快速进展，达到极其众多的人群，复活如此广泛地被传扬。因此百姓都尊重他们；但他们现在不再像以前那样轻易被藐视：因为在短时间内，只一次转机，渔夫和税吏就产生了如此的效果！大地成了天堂，因生活方式、讲话的胆量、奇事、以及其他一切；像天使一样，他们被惊异地看待；不关心嘲笑、威胁、危险；有同情心，乐于行善；一些人用金钱救助，一些人用言语，一些人用医治他们的身体和灵魂；\OriginalTerm{每一种治疗}{\GreekText{πἅν} \GreekText{εἶδος} \GreekText{ἰατρείας}}他们都施行了。\par

\Person{伯多禄}几乎为自己辩护，在即将施加惩罚之时，同时也给其他人一个教训。因为此行为似乎极其严厉，所以他在这件事上如此谨慎。关于那妇人，审判的过程也是可怕的。但请看，从亵圣中生出多少邪恶：贪婪、藐视天主、不敬；而且在这些事上，他也在会众面前为自己辩护，因为他没有立即施行惩罚，而是先揭露了罪过。无人呻吟，无人哀哭，众人都惊慌。因为随着他们的信德增长，奇事也增多，他们自己队伍中甚是恐惧：因为外部事物并不像我们自己人的行为那样\OriginalTerm{攻击}{\GreekText{πολεμεἵ}}我们的平安。如果我们紧密结合在一起，就没有什么战争是艰难的；但祸患在于分裂和瓦解。如今他们在公共场所走动：大胆地攻击市场，在敌人中间得胜，那句话就应验了：\BibleQuote{你要在仇敌中作王。}\BibleRef{咏 109:2}这乃是更大的奇事：他们被捕、下狱，竟然做出如此的事来！\par

如果那些说谎的人遭受了这样的事，那么作伪证的人将不会遭受什么呢？因为她简单地肯定说：\Quote{是的，就这么多}，你们看她遭受了什么。那么你们想一想；那些发誓和背誓的人，应当受到怎样的惩罚呢？今日甚至从《旧约》适当地出现，向你们显示伪誓的严重性。\Quote{曾有}，经上说，\Quote{一飞行的镰刀，宽十肘。}\BibleRef{匝 5:2} \Quote{飞行}表示追逐誓言的惩罚迅速来临；它长宽有许多肘，表示灾祸的力量和巨大；它从天上飞来，表明惩罚来自高处的审判座；它形状如镰刀，表明审判的不可避免性：因为正如镰刀，当它来到并钩住脖子时，不会空手而回，而是将头割下，同样，临到发誓者的惩罚是严厉的，必不止息直到完成工作。但若我们发誓却逃脱，不要自信；这反而是祸患。因为你们想什么？自从\Person{阿纳尼雅}和\Person{撒斐辣}以来，有多少人做了与他们相同的事？那么，你们说，为什么他们没有遭遇同样的命运？不是因为在他们那里被允许，而是因为他们被留待更大的惩罚。因为那些常犯罪而不受罚的人，比受罚的人更有理由恐惧和战栗。因为惩罚因他们目前的不受罚和天主的容忍而增加。那么，不要看我们不受罚这一点，而要省察我们是否犯了罪：若犯罪而不受罚，我们更有理由战栗。比方说，你若有一个奴隶，你只恐吓他，却不打他；他什么时候最恐惧，什么时候最想逃跑？岂不正是你只恐吓他的时候？因此我们彼此劝诫不要连续使用威吓，宁可藉恐怖搅动心灵，比打击更撕裂它。因为一边的惩罚是暂时的，另一边的却是永久的。如果没有人感到镰刀的击打，不要看这一点，而要让每人省察自己是否犯这样的罪。现在许多像洪水以前所行的事仍在行，却没有洪水降下：因为有地狱和惩罚被预告。许多人犯罪如同\Place{索多玛}人，却没有降下火雨：因为一条火河已预备。许多人作出与\Person{法郎}一样的事，却没有像\Person{法郎}那样结局，没有被淹死在\Place{红海}：因为等待他们的海是无底坑的海，在那里惩罚不伴有无知觉，没有窒息结束一切，而是在不断延长的折磨、焚烧、窒息中，他们被消灭在那里。许多人像以色列子民一样犯罪，却没有蛇吞吃他们：等待他们的是不死的虫。许多人像\Person{革哈齐}一样，却没有患癞病：因为代替癞病，他们将被斩成两段，与伪善者同列。许多人既发誓又背誓；但若他们真的逃脱了，我们不要自信：切齿等待他们。是的，在这里他们也要受许多严重的灾祸，虽然可能不是立即，而是在更多过犯之后，使惩罚更大：因为我们也常从小罪开始，然后因大罪失去一切。因此，当你看到有什么事发生在你身上时，要想起你的那一件特别的罪。\Person{雅各伯}的子孙是此例证。记住\Person{若瑟}的兄弟们：他们卖了他们的兄弟，甚至企图杀他；而且，就意愿而言，他们已经杀了他；他们欺骗并使老人悲伤；他们什么也没受。多年后，他们被带到极大的危险中，如今他们的罪被纪念。这件事的引入非常明智。听他们说什么：\Quote{我们实在是因我们的兄弟而有罪。} \BibleRef{创 42:21} 那么，你也照样，当有什么事发生时，说：我们实在有罪，因为我们没有服从基督；因为我们发了誓；我多次发誓，我假发誓，已落在我自己头上。你当认罪；因为他们也认了，并得救了。因为即使惩罚不立即跟随，又怎样呢？因为\Person{阿哈布}在\Person{纳波特}的事上犯罪后，也没有立即遭受他最终还是遭受了的惩罚。\BibleRef{列上 21:19} 这原因是什么？天主给你一段时期，在其中洗净自己；但如果你坚持，最后祂会降下惩罚。你们已看见说谎者的命运。想一想假发誓者的命运，想一想，然后停止。一个发誓者不可能不使自己背誓，无论他愿意与否；没有伪誓者能得救。一次假誓足以结束一切，为我们引来整个惩罚。那么让我们谨慎自守，好使我们逃脱这过犯应得的惩罚，并堪当承受天主的恩爱，藉祂独生子的恩宠和怜悯，愿光荣、权能和尊崇归于祂、圣父及圣神，自今至永远，世世无尽。阿们。\par

\SourceRecord{Translated by J. Walker, J. Sheppard and H. Browne, and revised by George B. Stevens.；Nicene and Post-Nicene Fathers, First Series；Vol. 11.；Edited by Philip Schaff.；Buffalo, NY: Christian Literature Publishing Co.,；1889.；<http://www.newadvent.org/fathers/210112.htm>.}

% structure: p0001 source=h1 target=section reason=page-title

\section{论宗徒大事录第十三篇讲道}

\BibleRef{宗 5:17-18}\par

\Quote{那时，大司祭和他同党的人（即撒杜塞派）充满愤怒，就下手拿住宗徒，把他们押在公共监牢里。}\par

\BibleQuote{起来以后}，即被唤醒，为所发生的事所激动，大司祭和他同党的人（即撒杜塞派）满怀愤恨，下手拿住宗徒：如今他们更猛烈地攻击他们：\BibleQuote{且把他们押在公共监狱里}。但他们没有立刻提审，因为他们以为他们又会软化下来。\BibleQuote{但是主的天使夜间开了狱门，领他们出来，说：你们去，站到圣殿里，把这一切生命的话讲给百姓听。} \BibleQuote{他们听了这话，天一亮就进入圣殿施教。} \BibleRef{宗 5:19} 这事既是为鼓励门徒，也是为其他的人得益和教训。请注意，目前发生的事与基督自己所行的事完全一样。即在祂的奇迹中，虽然祂不让人看到奇迹正在发生，却提供了使人能得知所行之事的方式：例如，在祂的复活中，祂不让他们看到祂如何复活；在变水为酒中，宾客们没有看到这事发生，因为他们已喝了很多，祂将辨认留给他人。同样在此情形中，他们没有看到他们被领出的情景，但能看到可以推断所发生之事的证据。而且天使是在夜间领他们出来的。为什么这样？因为这样他们比另一种方式更被相信：这样，人们甚至没有机会提出问题：他们不会以其他方式相信。古时在拿步高王的事上也是如此：他看见他们在火窑中赞美天主，那时他确实惊愕不已。\BibleRef{达 3:24} 然而这些司祭本该首先问：你们是怎么出来的？但相反，他们好像什么事都没发生，问道：\BibleQuote{我们不是严格禁止你们讲话吗？} \BibleRef{宗 5:28} 请注意，通过他人的报告，他们知道了所有情况：他们看见监狱仍然紧闭安全，守卫站在门前。这是双重的保障；正如在坟墓那里，既有封印，又有守夜的人。看他们是如何对抗天主！你说，所发生的这些事，难道是人的作为吗？门关闭时，是谁领他们出来的？看守站在门前，他们是怎么出来的？他们这样说话，一定是疯了或醉了。这里有一些人，监狱、锁链、关闭的门都不能将他们关住；而他们竟期望制服他们：这是何等幼稚的愚妄！他们的差役来承认所发生的事，好像故意要剥夺他们任何理由似的。你看到吗，奇迹不断，种类不同，有些由宗徒们行出，有些临到他们身上，且比前者更显著？\BibleQuote{他们听了这话，天一亮就进入圣殿施教。但大司祭和他同党的人来了，召集公议会，并以色列子民的全长老议会，差人到监狱去提他们。差役来到，在监狱里找不到他们，便回去报告说：我们看见监狱关得严严实实，守卫站在门外；但打开后，里面却空无一人。大司祭、圣殿警官和司祭长们听了这些话，不知这事将会如何发展。} \BibleRef{宗 5:21} 这安排得很好：消息没有立刻传到他们那里，而是他们先全然不知所措，当他们仔细思考并看出其中有天主的大能时，才得知整个情况。\BibleQuote{这时有人来报告说：你们关在监狱里的人，正站在圣殿里教训百姓。于是圣殿警官同差役前去，不用武力把他们带来，因为他们害怕群众，免得被石头砸死。}（25-26节）人们的愚妄啊！\BibleQuote{他们害怕}，他说，\BibleQuote{群众。}为什么？群众何时帮助过宗徒？他们本该害怕那位不断像带翅膀的飞鸟一样解救他们脱离他们权势的天主，却反而害怕群众！\BibleQuote{大司祭}，无耻、鲁莽、愚蠢，\BibleQuote{问他们说：我们不是严格命令你们，不可因这名施教吗？看，你们竟把你们的道理传遍了耶路撒冷，想要把这人的血归到我们身上。}（27-28节）那么（宗徒们说什么？）他们再次以温和的态度对他们说话；虽然他们本可以说：\Quote{你是谁，竟敢违抗天主？}但他们说什么？他们再次以劝诫和忠告的方式，并带着极大的温和回答。\BibleQuote{伯多禄和其余的宗徒回答说：听天主的命应胜过听人的命。} \BibleRef{宗 5:29} 何等崇高的胸怀！他也向他们表明他们是在对抗天主。因为他说：你们所杀死的，天主却复活了祂。\BibleQuote{我们祖先的天主复活了耶稣，就是你们挂在木头上所杀死的。天主以右手高举祂，使祂作元首和救主，为把悔改和赦罪赐给以色列人。}（30-31节）他们再次将一切归于父，以免祂似乎与父不相干。\BibleQuote{祂以右手}，祂说，\BibleQuote{高举了祂。}祂不仅肯定复活，也肯定高举。\BibleQuote{为把悔改赐给以色列人。}请注意这里如前所述的获益（对他们的）：请注意以道歉形式表达的教理完美。\BibleQuote{我们就是这些事的见证人。} \BibleRef{宗 5:32} 多么大胆的言论！他们可信的理由：\BibleQuote{还有圣神，就是天主赐给服从祂之人的圣神。}你是否注意到他们不仅引证圣神的见证？他们也没有说\BibleQuote{祂赐给了}我们，而是\BibleQuote{赐给了服从祂的人}：这样既显出自己的谦逊，又暗示恩赐的伟大，并告诉听众他们也有可能领受圣神。看，这些人既通过行为又通过言语受教，却仍不留意，为使他们受审是公正的。因为天主允许宗徒们受审，正是为此目的：使他们的仇敌受教，使众人学习，并使宗徒们被激励而大胆讲话。\BibleQuote{他们听了这话，内心被刺透。} \BibleRef{宗 5:33} 其他人（在先前场合）\BibleQuote{听了这话，觉得扎心}；这里他们如被\OriginalTerm{锯开}{\GreekText{διεπρίοντο}}，\BibleQuote{就想要杀害他们。} \BibleRef{宗 2:37}\par

但是现在有必要重新审视我们所读的。\Quote{但是，主的天使在夜间打开了监狱的门，领他们出来，说：\InnerQuote{去，站在圣殿里，把这一切生命的话讲给百姓听。}领他们出来。}（复述，第19-20节。）他领他们出来，并非为了让他们自己得到益处，而是说：\Quote{站，}他说，\Quote{在圣殿里讲给百姓听。}但如果卫兵赶他们出去，如同那些人认为的那样，他们就会逃跑，也就是说，如果他们被诱导出来的话；如果那些人把他们赶出去，他们就不会站在圣殿里，而是躲藏起来。没有人愚钝到看不出来这一点。\BibleQuote{我们不是严厉地命令过你们吗？}\BibleRef{宗 5:28} 嗯，如果他们答应服从你们，你们要他们交代是应该的；但如果就在那时他们告诉你们不会服从，你们有什么理由叫他们交代，他们又有什么辩护可言？\Quote{看，你们以你们的教训充满了耶路撒冷，并且想把这个人血债归于我们。}注意这些指控的矛盾和极大的愚蠢。他们现在想让人以为犹太人的性情是残忍的，好像他们做这些事不是为了真理，而是为了报复。因此，宗徒们也没有用挑衅（\OriginalTerm{鲁莽地}{\GreekText{θρασέως}}）来回答他们，因为他们是老师。然而，有哪一个有整个城市支持、又蒙受了如此大恩宠的人，不会说话并说些大话？但这些宗徒却不是这样：因为他们没有发怒；不，他们怜悯这些人，为他们哭泣，并思考如何使他们摆脱错误和愤怒。他们不再对他们说：\BibleQuote{你们判断吧！}\BibleRef{宗 4:19} 而是简单地肯定说：\Quote{天主所复活的那位，我们宣讲：这些事是出于天主的旨意进行的。}他们没有说：难道我们当时没有告诉你们，\BibleQuote{我们不得不说我们所见所闻的事吗？}\BibleRef{宗 4:20} 因为他们不争荣誉；而是重复那个故事——十字架、复活。他们没有说祂为何被钉十字架——那是为了我们；但他们暗示了这一点，但目前还没有公开，想暂时吓唬他们。然而，这算是什么修辞？完全没有，但到处仍然是受难、复活和升天，以及目的：\BibleQuote{我们祖先的天主复活了耶稣，等等。}\BibleRef{宗 5:30-31} 这些是多么不可能的断言！非常不可能，毫无疑问；但尽管如此，没有统治者，没有人民，能够反驳他们：但那些人哑口无言，这些人接受了教导。他说：\BibleQuote{我们是这些事的见证人。}\BibleRef{宗 5:32} 哪些事？关于祂应许了宽恕和悔改：因为复活现在确实已被承认。但祂赐予宽恕，我们既是见证人，\Quote{圣神也是，}祂不会降临，除非罪恶已被先赦免：所以这是一个无可辩驳的证明。\BibleQuote{他们听了这话，就心如刀割（痛彻心扉），并商议要杀害他们。}\BibleRef{宗 5:33} 你听到了罪恶的宽恕，啊，可悲的人，以及天主不要求惩罚，而你却想杀害他们？这是何等的邪恶！然而，他们要么应该证明他们在说谎，要么如果无法证明，就应该相信；但如果他们选择不相信，他们也不应该杀害他们。因为有什么是值得一死的？他们的醉意如此之深，甚至没有看到所发生的事。请注意，宗徒们每当提到罪行时，都会加上宽恕的提及；表明，虽然所做的应受死刑，但所赐予的却是作为恩人而给予他们的！还有其他方式能说服他们吗？\par

\Quote{于是大司祭站起来，}等等。作为享有盛誉的人，这些人（即宗徒们）即将跻身于先知之列。撒杜塞人最反对复活之说。但也许有人会说：既然宗徒们蒙受了如此恩赐，什么人不会变得伟大呢？但请想一想，在他们领受恩宠之前，他们是如何\BibleQuote{同心合意恒切祈祷}\BibleRef{宗 1:14}，并依靠从上而来的援助。而你呢，我亲爱的，盼望天国，却不忍受任何苦楚吗？你既领受了圣神，竟不遭受这些，也不遇见危险吗？但他们，从先前的危险中尚未喘息，又被带入新的危险。况且，没有骄傲，没有自大，这是何等的美善！以温和交谈，是何等的收获！因为他们所做的一切，并非完全是恩宠的直接作为，也有许多他们自己热忱的记号。恩宠的恩赐在他们身上闪耀，这是来自他们自己的勤勉。请看，例如，从起初，\Person{伯多禄}是多么谨慎、庄重而警觉；那些信的人如何舍弃财富，没有私有财产，恒心祈祷，显示他们同心合意，禁食度日。我请问，单靠恩宠，能做到这一切吗？因此，祂藉着他们自己的差役，将证据带到他们面前。正如在\Person{基督}的情况下，正是他们的差役说：\BibleQuote{从来没有像这人那样说话的。}\BibleRef{若 7:46}这些证据比复活更易令人相信。——也请注意首领们自己所表现的温和，以及他们如何退让。\Quote{大司祭问他们说}等等。\BibleRef{宗 5:27}这里他确实以温和的口气与他们理论，因为他害怕了；他确实更愿阻止他们，而非杀死他们，因为他不能这样做；并且为了激起他们众人，并向他们表明他们所面临的极大危险，\Quote{并且打算，}他对宗徒们说，\Quote{将这人血归到我们身上。}你仍然以为祂只是\Emph{人吗？}他想要使那禁令显得是为他们自身安全所必需的。但请注意（\Person{伯多禄}）所说的话：\BibleQuote{天主用右手高举了祂，立祂为元首和救主，为将悔改和赦罪赐给以色列。}\BibleRef{宗 5:31}这里他避免提到外邦人，以免给他们把柄攻击他。\Quote{他们想要}，经上说，\Quote{杀害他们}。\BibleRef{宗 5:33}再看这些人在困惑中，这些人在痛苦中；而那些人在安宁、喜乐和欢愉中。不只是说他们忧伤，而是\Quote{他们被刺透}（心肠）。这确实应验了那句谚语：\Quote{害人者自害}，正如我们在此例中所见。这些人被捆绑，被带到审判台前，而审判他们的人却陷于困苦和无助的困惑中。因为如同击打金刚石的人自己受到震击一样，这些人也是如此。但他们看到，不仅他们放胆宣讲没有被阻止，反而他们的宣讲越发增长，并且他们毫无畏惧地讲论，不给他们任何把柄。\par

让我们效法这些人，亲爱的诸位：让我们在一切危险中无所畏惧。对于敬畏天主的人，没有什么可畏惧的；一切可畏惧的都是为别人的。因为一个人若脱离了私欲偏情，视今世一切如影，请问，他还会从谁那里遭受什么可怕的事呢？他还需要畏惧谁呢？他还需要向谁恳求呢？让我们逃向这不可动摇的磐石。假如有人为我们建造一座城，围起城墙，将我们迁到无人居住之地，那里没有人打扰我们，给我们供应一切丰盛，不让我们因任何人而烦恼，他所给予我们的安全，也不及基督现在所赐予的。就算是一座铜城，如果你愿意，四面有城墙环绕，高耸坚固，没有敌人靠近；土地肥沃富饶，又有其他丰盛之物；居民温和良善，没有作恶者，没有强盗，没有窃贼，没有告密者，没有法庭，只有\OriginalTerm{协议}{\GreekText{συναλλάλματα}}；让我们住在这城里：即使如此，也无法安然生活。为什么？因为难免与仆人、妻子、儿女有分歧，成为许多烦恼的根源。但这里却没有这种事；因为这里没有任何事情使他们痛苦或引起不快。而且，更奇妙的是，那些被认为引起不快的事，反而成了所有喜乐和欢愉的来源。请问，有什么事情使他们烦恼？有什么可抱怨的？我们要举一个具体的例子与他们比较吗？好吧，假设有一个拥有执政官尊位的人，家财万贯，住在帝都，与人无争，只知享乐，别无他事，处于财富、荣誉和权力的巅峰；我们再拿一位\Person{彼得}与他对比，假设\Person{彼得}被锁链束缚，遭遇无数患难：我们就会发现，\Person{彼得}才是活得最快乐的人。因为当喜乐如此满溢，以至在锁链中仍能欢喜时，试想那喜乐是多么伟大！正如身居高位的人，无论发生什么祸患，都感觉不到，继续享受快乐：同样，这些人反而因这些患难更加喜乐。因为言语无法形容那些为基督受苦的人所获得的喜乐是多么巨大：他们在受苦中，而非在顺境中欢喜。凡爱基督的人，都知道我在说什么。———但谈到安全呢？请问，有谁，即使再富有，能避开如此多的危险，周游于如此众多的民族之间，只为改变他们的生活方式？因为他们所向披靡，仿佛奉了王命，甚至更轻易，因为从未有命令能像他们的话语那样有效。因为皇帝的法令强制人服从，而这些人却吸引人自愿、自发地跟随，并且满怀难以言表的感激。请问，有什么皇帝的法令能说服人舍弃全部财产和生命；蔑视家庭、祖国、亲属，甚至保全自己？然而，渔夫和制帐篷的人的声音却能成就此事。因此，他们既幸福，又比所有人更有能力和力量。\Quote{是的，}你说，\Quote{那些当然行过奇迹。}\EditorialNote{见前文第83页注释4}但我要问，那些相信的人——三千人和五千人——行了什么奇迹？然而，我们读到他们却在喜乐中度日。他们理应如此：因为一切烦恼的根源——拥有财富——被消除了。因为那，那，我说，一直是战争、争斗、悲伤、烦恼和一切罪恶的原因：它使生活充满劳苦和烦忧。实际上，你会发现富人比穷人有更多理由悲伤。如果有人觉得这不真实，那他们的想法不是来自事物的本质，而是来自自己的幻想。如果富人确实享受某种乐趣，这并不奇怪：因为即使是浑身痒痒的人，也有不少快感。因为富人与这些人完全相似，他们的心智也受到同样影响，从这一点可以清楚看出。他们的忧虑困扰着他们，而他们却为了短暂的快感选择沉浸其中；而那些摆脱这些情感的人，则健康而无烦恼。请问，哪一样更愉快，哪一样更安全？只需要操心一块面包和一套衣服，还是要操心一个大家族，包括奴隶和自由人，而不只是为自己操心？因为正如这个人为自己担心，你也为你所依赖的那些人担心。请问，为什么贫穷看似应当逃避？正如其他好东西在许多人看来也是应当鄙弃的。\Quote{是的，}你说，\Quote{但并非那些好东西应当鄙弃，而是它们难以获得。}那么，贫穷也是如此，不是应当鄙弃的，而是难以承受的：所以，如果一个人能忍受它，就没有理由鄙弃它。宗徒们为何不鄙弃它？为何许多人甚至选择它，非但不鄙弃，反而奔向它？因为真正应当鄙弃的东西，除了疯子，没有人会选择它。但如果是有哲学思想和高尚心灵的人投身于此，如同进入一个安全有益的退省之所，那么其他的人对它有不同的看法就不足为奇了。因为，事实上，在我看来，富人就像一座没有城墙的城，位于平原，四面招引攻击者；而贫穷则是一座坚固的堡垒，坚如铜墙，通往它的路是艰难的。\Quote{然而，}你说，\Quote{事实恰恰相反：因为正是这些人常常被拖进法庭，正是这些人被压倒、受虐待。}不：并不是穷人作为穷人，而是那些贫穷却想发财的人。但我说的不是他们，而是那些以安于贫穷为务的人。请问，为什么从来没有人把山中的弟兄们拖进法庭？然而，如果贫穷是受压迫的标志，那么那些人最应该被拖进去，因为他们比所有人都穷。为什么没有人把普通乞丐拖进法庭？因为他们已到了贫穷的极致。为什么没有人对他们施暴，没有人恶意控告他们？因为他们住在一个对此太过安全的堡垒中。有多少人认为贫穷和乞讨是一种难以抗拒的境况！那么，请问，乞讨是好事吗？\Quote{如果有安慰，}你说，\Quote{如果有人给，那就是好事：众所周知，这是一种无忧无虑、没有挫折的生活。}但我并非要称赞这个；断乎不可！我所劝诫的是不要追求财富。\par

请问，你愿称谁为有福？是那些亲近美德的人\OriginalTerm{亲近美德者}{\GreekText{ἐπιτηδείους} \GreekText{πρὸς} \GreekText{ἀρετήν}}，还是那些远离美德的人？当然，是那些亲近之人。那么你说，这两种人中，哪一种人更能学到有益之事，并在真智慧中发光？是前者还是后者？所有人都必看到是前者。若你怀疑，请这样验证。从集市上带一个可怜的乞丐来；让他是跛子、瘸腿、残废的；然后再带另一个人，容貌俊美，身体强壮，四肢充满活力与精力，富甲一方；让他出身显赫，拥有大权。然后我们将这两人带进哲学学校：请问，哪一个更有可能接受所教的内容？第一条训诫，一开始就说：\Quote{要谦卑和节制}（因为这是基督的命令）：哪一个人最能履行它，是这个还是那个？\BibleQuote{哀恸的人是有福的}\BibleRef{玛 5:4}：哪一个人最能接受这句话？\Quote{温良的人是有福的：}哪一个人最听从这话？\BibleQuote{心里洁净的人是有福的。饥渴慕义的人是有福的。为义受迫害的人是有福的}\BibleRef{玛 5:8,6,10}。哪一个人能轻易接受这些话？你若愿意，让我们将这些规则应用于他们所有人，看看它们如何适合。难道不是一个人全身发炎肿胀，而另一个人则始终保持谦卑和受控的举止吗？这很清楚。是的，在外邦人中也有这样的说法：\Quote{（我）是一个奴隶，名叫埃皮克提图，身体残疾，贫穷如伊鲁斯，却是不朽者的朋友。}因为请问，怎么能不然呢？富人的灵魂必然充满邪恶：愚妄、虚荣、无数情欲、愤怒和激情、贪婪、不义，等等。所以，甚至对于哲学而言，前者比后者更为适宜\OriginalTerm{适宜}{\GreekText{ἐπιτηδεία}}。无论如何，要试图查明哪种生活方式更令人愉悦：因为我看到这一点到处都在讨论，某人的生活是否更有乐趣。然而，即使在这方面，我们也不应怀疑；因为接近健康，也意味着享有更多的快乐。但请问，这两种人中，哪一个更适宜于\OriginalTerm{最适宜}{\GreekText{ἐπιτήδειος}}当前的事务，即我们必须完成的事情——我们的法律，我是说——是穷人还是富人？哪一个更容易起誓？是那个有孩子可激怒的人，那个与无数伙伴有契约的人，还是那个只关心讨一块面包或一件衣服的人？这个人甚至不需要起誓，如果他愿意，他总是不为世俗事务所累；不仅如此，人们常常看到，那训练自己绝不发誓的人，也会轻视财富；人们会在他的整个行为中看到，他的道路全部源于这一良好的习惯，导向温良、轻视财富、虔敬、灵魂的谦卑、心灵的痛悔。那么，我亲爱的，让我们不要懈怠，而要再次显出极大的热忱：那些已经成功的人，为要保守已取得的成功，以免轻易被退去的波浪抓住，或退潮将他们冲回；[那些尚未达到的人，为要重新崛起，并努力补足所缺的。同时，让那些已经成功的人，帮助那些未能做到同样的人]：伸出手，如同向深水中挣扎的人伸出援手，将他们接进不宣誓的港湾\OriginalTerm{不宣誓}{\GreekText{ἀνωμοσίας}}。因为完全不宣誓确实是一个安全的港湾：无论什么风暴袭击我们，我们都没有沉溺的危险：无论是愤怒、侮辱、激情，还是别的什么，灵魂都稳固地安息；是的，即使一个人不慎说出了某些不该说、最好不说的话，他也没有使自己处于任何必要性或法律之下。（\Emph{Supra}, Hom. ix. §5. ad. Pop. Ant. viii. §3.）请看黑落德因他的誓言做了什么：他斩了前驱的头。\BibleQuote{但因他的誓言，}经上说，\BibleQuote{和同席的人}\BibleRef{谷 6:26}，他斩了先知的头。请想，各支派因他们在本雅明支派事件中的誓言遭受了什么\BibleRef{民 21:5}；撒乌耳因他的誓言遭受了什么\BibleRef{撒上 14:24}。因为撒乌耳确实发了虚誓，但黑落德做了比背誓更坏的事，他犯了谋杀。若苏厄也是如此——你知道他在基贝红事件中因誓言而遭遇了什么。\EditorialNote{若苏厄书第9章}因为这起誓实在是撒殚的罗网。让我们挣断绳索；让我们使自己处于（不）容易起誓的状态；让我们摆脱一切纠缠，脱离撒殚的这个罗网。让我们敬畏主的命令；让我们定居于最好的习惯中：这样，我们不断进步，在完成了这一条和其他诫命之后，就能获得那为爱祂的人所应许的美物，借着我们的主耶稣基督的恩宠和慈爱，祂和圣父及圣神，一同享有荣耀、权能和尊崇，从现在直到永远，世世无尽。阿们。\par

\SourceRecord{Translated by J. Walker, J. Sheppard and H. Browne, and revised by George B. Stevens.；Nicene and Post-Nicene Fathers, First Series；Vol. 11.；Edited by Philip Schaff.；Buffalo, NY: Christian Literature Publishing Co.,；1889.；<http://www.newadvent.org/fathers/210113.htm>.}

% structure: p0001 source=h1 target=section reason=page-title

\section{关于\WorkTitle{宗徒大事录}的第十四篇讲道}

\BibleRef{宗 5:34}\par

\Quote{那时，在公议会中站起了一位法利塞人，名叫加玛里耳，是一位法学士，众百姓都尊敬他，他命人把那些人暂且带出去。}\par

这位加玛里耳是保禄的老师。人们或许会感到奇怪，他既然判断如此正确，又精通法律，为何还没有相信。但他不可能一直不信到底。事实上，从他这里所说的话可以明显看出。他\Quote{吩咐}，圣经说，\Quote{把那些人暂时带出去〔然后向他们说〕。}请注意他多么明智地组织他的言论，一开头就使他们感到恐惧。而且为了不被人怀疑偏袒他们，他说话时好像自己与他们意见一致，并不激烈，而是像对因激情而陶醉的人说话一样，他这样说道：\Quote{以色列人哪，你们要小心，对于这些人，你们打算怎样做。}\BibleRef{宗 5:35}他要说的是，不要轻率仓促地行事。\Quote{因为在那些日子以前，有个特乌达起来，自命不凡，约有四百人附从了他；他被杀后，所有服从他的人都四散消灭了。}\BibleRef{宗 5:36}他用实例教导他们谨慎；并且为鼓励起见，最后提到那个迷惑最多的人。在他举出例子之前，他说：\Quote{你们要小心；}但引证之后，他发表意见说：\Quote{不要管这些人。}因为，他说：\Quote{在户口登记的日子，有个加里肋亚人犹大起来，引走了许多百姓；他也灭亡了，所有服从他的人也都四散了。现在我告诉你们，不要管这些人，任凭他们吧！因为这计划或这事业若是出于人，必要消亡；但若是出于天主，你们就不能推翻他们。}\BibleRef{宗 5:37}（另作\Emph{它}）那么，他说，有什么能阻止你们被推翻呢？因为，他说（注意）：\Quote{恐怕你们反而成了与天主作对的人。}他既从事情不可能的角度，又从对你们无益的角度来劝阻他们。他没有说这些人被谁消灭，只说他们\Quote{四散了}，他们的联盟归于虚无。因为，他说，若是出于人，你们何必费心？但若是出于天主，你们再费心也不能胜过它。这个论证是无法反驳的。\Quote{他们被他说服了。}\BibleRef{宗 5:40}他们如何被说服？以致没有杀死他们，只是鞭打。因为圣经说：\Quote{他们叫来宗徒，鞭打了他们，又命令他们不可再因耶稣的名发言，然后释放了他们。}看哪，在行了那么伟大的事之后，他们竟被鞭打！而他们的教导反而更加扩展：因为在家里和在圣殿里，他们\Quote{离开公议会，心里欢喜，因被算是配为这名受辱。他们就每天在圣殿里，并挨家挨户，不断教导传扬耶稣基督。（41-42节）}在那些日子，门徒的数目加增的时候，说希腊话的希伯来人向希伯来人发怨言，因为在日常的供应上，他们的寡妇被忽略了。\BibleRef{宗 6:1}并非绝对就在那些紧接着的日子里；因为圣经的习惯是将即将发生的事说成是紧接着发生的。但我想他所说的\Quote{说希腊话的}是指那些讲希腊语的人（\Quote{向希伯来人}）：因为他们并不使用希腊语。看，又一个考验！注意，从一开始，内外就有争战！\Quote{那时，}圣经说，\Quote{十二宗徒召集众门徒来说：我们撇下天主的圣言，去伺候饮食，实在不合宜。}\BibleRef{宗 6:2}说得好：因为必须将必要的让给更必要的。但是请看，他们如何立即既关心这些（次要事务），又不忽略宣讲。\Quote{因为他们的寡妇被忽略：}因为那些（希伯来人）\OriginalTerm{被视为更重要的人}{\GreekText{αἰδεσιμώτεροι}}。\Quote{所以，弟兄们，当从你们中间选出七位有好名声、充满圣神和智慧的人，派他们管理这事。至于我们，我们要专务祈祷和宣讲圣言。这话使全体群众喜悦，他们就选了斯德望，一个充满信德和圣神的人}\BibleRef{宗 6:3}其余的人也同样充满信德；以免发生像犹大、阿纳尼雅和撒斐辣那样的事——\Quote{并选出了斐理伯、普洛曷洛、尼加诺尔、提摩、帕尔默纳和安提约基雅的进教者尼各老。他们把他们立在宗徒面前，宗徒们祈祷后，就给他们覆手。天主的圣言逐渐兴盛，门徒的数目在耶路撒冷大为增加，也有许多司祭服从了信德。}\BibleRef{宗 6:5}\par

但是，让我们再回顾一下刚才所说的话。\BibleQuote{以色列人，你们应小心谨慎对待这些人。}\BibleRef{宗 5:35}请看，我请你们注意，加玛里耳是如何温和地推理，他只用几句话对他们说话，并没有讲述古代的历史，虽然他本可以这样做，而是举出最近的实例，这些实例最有助于产生信服。为此，他自己暗示说：\BibleQuote{因为在这些日子以前}\BibleRef{宗 5:36}，意思是说，没有多少天以前。若是他立即说：\Quote{放这些人走吧}，那么他自己就会陷入嫌疑，他的话也不会那么有效；但在举例之后，他的话就有了自己的效力。他不仅举了一个例子，还举了第二个：\Quote{因为}，圣经说，\BibleQuote{凭两三个人的口}\BibleRef{玛 18:16}；虽然他本可以举出第三个例子。\BibleQuote{不要管这些人。}\BibleRef{宗 5:38}请看他的态度是多么温和，他的话不多而精炼，他提到那些（骗子）时也毫无激情：\Quote{凡听从他们的人，都四散了。}尽管这一切，他并没有亵渎基督。这些不信的人听到了他，这些犹太人听到了他。[\BibleQuote{如果这计划或这工作是从人来的，必要消灭。}]那么，既然它没有消灭，就不是从人来的。[\BibleQuote{但如果它是从天主来的，你们就不能推翻。}]\BibleRef{宗 5:39}他又用事情的不可能和不合理来制止他们，说：\BibleQuote{免得你们被发现甚至是在与天主交战。}\BibleRef{宗 5:39}他并没有说：如果基督是天主；而是工作（本身）表明这一点。他没有断定它是\Quote{从人来的}还是\Quote{从天主来的}，而是把证明留给将来。\BibleQuote{他们［被他］说服了。}\BibleRef{宗 5:40}那么，有人会问，为什么你们还鞭打他们？他的讲话是那样无可辩驳的公正，他们无法正视它；然而，他们满足了自己的敌意；并且他们期望用这种方式恐吓他们。他不在宗徒面前说这些话，这事本身使他比否则更能被听进去；然后，他话语的温和与所说的公正，帮助说服了他们。事实上，这个人几乎就是在宣讲福音。\Quote{你们被说服了}，有人会说，\Quote{你们没有力量推翻它。你们为什么不相信呢？}这就是连仇敌也作的见证。那里有四百人，那里有四千人；而这里是十二位发起者。不要让那不断增加的人数吓倒你们。\EditorialNote{参阅 2:41; 4:4}他本可以再举另一个例子，比如那个埃及人的，但他所说的已经足够了。他以一个令人震惊的话题结束他的讲话：\Quote{免得}等等。他没有断定，免得他显得是在为他们辩护；而是从事物的结局出发用三段论推理。他不敢断定它不是从人来的，也不敢断定它是从天主来的；因为如果他说它是从天主来的，他们会反驳他；但如果他说它是从人来的，他们又会立即采取行动。因此他让他们等待结局，说：\Quote{不要管。}但他们再次威胁，明知他们毫无用处，却按他们的方式行事。邪恶的本性就是如此：它甚至尝试不可能的事。——\Quote{此后犹达斯起来}等等。若瑟夫详细记述了这些事。（\WorkTitle{Ant.} xx. 8; ib. v. 2; xviii. 1. \WorkTitle{B. J.} ii. 8. 1.）他敢于断定那是从天主来的，这在后来从事件中得到了证实！这是多么伟大的大胆直言，多么伟大的不顾情面！他没有说：\Quote{但如果你们不推翻它，它就是从天主来的}，而是说：\Quote{如果它是从天主来的，它就不会被推翻。}\BibleQuote{他们同意了他的话。}\BibleRef{宗 5:40}他们尊敬这人的高尚品格。\BibleQuote{他们离开公议会，心里欢喜，因被算是配为这名受辱。}\BibleRef{宗 5:41}有什么奇迹能像这样令人惊奇？在古圣者中没有这样的记载：因为耶肋米亚的确为天主的话受过鞭打，他们威胁过厄里亚，还有其他人；但在这里，正是通过这件事，而不仅仅是他们的奇迹，这些人彰显了天主的能力。他并没有说他们不痛苦，而是说他们虽然痛苦却欢喜。这如何显明呢？从他们后来的勇敢：他们即使在被打之后，仍然坚持不懈地宣讲圣言。\BibleQuote{在圣殿里，并在各家中，他们不断地教训人，传耶稣基督为默西亚。}\BibleRef{宗 5:42}\Quote{在那时候}——当这些事件发生时，当有鞭打，有威胁，当门徒增多时——同时，经上说，\BibleQuote{起了怨言。}\BibleRef{宗 6:1}这是由于人多：因为要在众多人中保持严格秩序是不可能的。\Quote{起了怨言}等等，直到\BibleQuote{有许多司祭也服从了信仰。}——\Quote{对希伯来人起了怨言}——因为那一类似乎更尊贵——\BibleQuote{因为在每日的供应上，他们的寡妇被疏忽了。}\BibleRef{宗 6:1}那么，每日都有对寡妇的供应。请注意他如何称其为\Quote{服务}\OriginalTerm{服务}{\GreekText{διακονία}}，而非直接说施舍：以此同时赞扬了施予者和接受者。\Quote{被疏忽了}。这不是出于恶意，而是或许由于众人的疏忽。因此他公开提出此事，因为这并非小事。请注意，即使在开始，祸患不仅来自外部，也来自内部。因为你们不仅要注意到它被纠正了，还要注意它本身是一个大祸患。\BibleQuote{那时，十二位}等等。\BibleRef{宗 6:2}你们是否注意到外部事务如何接踵而至？他们不是凭自己的判断行事，而是向会众申辩。现在也应该这样做。\BibleQuote{我们放弃天主的圣言去管理饭食，是不合理的。}\BibleRef{宗 6:2}他首先向他们指出此事的不合理之处：即不可能以同样的注意力做这两件事；正如当他们要祝圣玛弟亚时，他们首先说明事情的必需性，即缺少一位，必须有十二位。同样，他们在这里也说明了必需性；他们并没有更早做，而是等到怨言发生后；另一方面，他们也没有让此事蔓延太远。看！他们把决定权留给他们：那些使众人喜悦的人，那些在众人中名誉良好的人，他们推荐出来：不是十二位，而是\BibleQuote{七位，充满圣神和智慧：在品行上有好名声的人}\BibleRef{宗 6:3}。当要推荐玛弟亚时，经上说：\BibleQuote{所以必须从这些与我们常年在一起的人中}\BibleRef{宗 1:21}；但这里不是这样：因为情况不同。他们现在也不用拈阄；他们本可以自己选择，如同受圣神感动；但他们仍然希望有百姓的见证。决定人数，祝圣他们，以及为此类事务，在于他们；但选择人则交给百姓，免得他们显得是出于偏爱；正如天主也留给梅瑟选择他所认识的人作长老。\BibleRef{户 11:16}\BibleQuote{和智慧。}因为在这种服务中确实需要许多智慧。不要因为这样的人没有受托宣讲圣言，就以为他不需要智慧：他需要，而且很多。\Quote{但我们}，他说，\BibleQuote{要专务祈祷，并专务服务圣言。}\BibleRef{宗 6:4}他们再次为自己申辩，从头到尾以此开始和结束。\BibleQuote{要专务}，他说。因为应该如此，不只是做做样子，而是经常做。\BibleQuote{这话}，我们被告知，\BibleQuote{使全体百姓都满意。}（第5-6节）这也与他们智慧相称。大家都赞同所说的话，因为这很有道理。经上说（又是百姓［\OriginalTerm{他们}{\GreekText{αὐτοί}}］选的）：\BibleQuote{他们选了斯德望，一个充满信德和圣神的人，以及斐理伯、普洛苛、尼加诺尔、提摩、帕尔默纳和来自安提约基雅的尼苛劳，一个归依犹太教的人；他们把他们带到宗徒跟前；祈祷后，就给他们覆手。}\BibleRef{宗 6:5}他们把他们从众人中分别出来，是百姓（\OriginalTerm{他们}{\GreekText{αὐτοί}}）把他们带来，而不是宗徒引领他们。请注意他如何避免一切多余的话：他没有说明是以什么方式做的，只说他们被祝圣（\OriginalTerm{被按立}{\GreekText{ἐχειροτονήθησαν}}）并伴随着祈祷：因为这就是\OriginalTerm{按手}{\GreekText{χειροτονία}}的意思，即\Quote{按手}或祝圣：人的手按在（那人）身上，但整个工作是天主的，如果那人恰当地被祝圣，是天主的手触摸被祝圣者的头。\BibleQuote{天主的圣言更加兴盛；门徒的数目大增。}\BibleRef{宗 6:7}他不是无缘无故说这话的：它表明施舍和良好秩序的德行有多大。当他要在后面\OriginalTerm{扩大}{\GreekText{αὔξειν}}关于斯德望的事时，他首先提出导致这事的原因。\BibleQuote{并且有许多司祭}，他说，\BibleQuote{服从了信仰。}因为他们看到他们的领袖和老师有这样的心意，他们就以事实来检验这事。——同样令人惊奇的是，群众在选择那些人时没有分裂，宗徒也没有被他们拒绝。但这些人担任什么等级，接受什么职务，这是我们需要的。是执事吗？而在教会中并非如此。是长老管理吗？但当时没有主教，只有宗徒。因此我认为清楚而明显的是，他们既不是执事也不是长老的称号：他们是为了这个特殊目的而被祝圣的。这事并没有简单地交给他们而不加仪式，而是宗徒为他们祈祷，使他们获得能力。但是，请注意，如果为此需要七个人，那么相应的金钱数目一定很大，相应的寡妇人数也一定很多。因此，祈祷不是随意做的，而是经过深思熟虑；并且这职务，如同宣讲一样，被有效地执行；因为他们所做的，主要是通过这些（他们的祈祷）成就的。这样，他们就能专注于属灵的事；同样，这些人也能自由地长途旅行；这样，这些人就受托圣言。但作者没有说这些，也没有赞扬他们，而是说他们\Quote{不合理}离开交给他们的工作。他们从梅瑟的榜样学到不要独自承担一切管理。\BibleRef{户 11:14}\BibleQuote{只要纪念穷人}\BibleRef{迦 2:10}他们是如何推荐这些人的呢？他们禁食。\BibleQuote{你们要从你们中选出七位}等等。\BibleRef{宗 6:3}不是简单地说是属灵的人，而是\BibleQuote{充满圣神和智慧}，而且需要有大灵性\OriginalTerm{灵性}{\GreekText{φιλοσοφίας}}来忍受寡妇的抱怨。如果管理施舍的人不偷窃，但却浪费一切，或者严厉易怒，又有什么益处呢？\BibleQuote{他们选了斯德望，一个充满信德和圣神的人。}\BibleRef{宗 6:5}在这方面，斐理伯也令人钦佩：因为关于他，作者说：\BibleQuote{我们进了传福音者斐理伯的家，他是七位之一；与他同住。}——\BibleRef{宗 21:8}你们是否注意到，事情的处理完全不同于人的方式？\BibleQuote{在耶路撒冷，门徒的数目更加增多。}\BibleRef{宗 6:7}在耶路撒冷人数增多。奇妙，在那里基督被杀死的地方，宣讲反而增长了！并且不仅没有像阿纳尼雅那样有人被冒犯，反而敬畏更大了：这些人被鞭打，那些人恐吓，那些人试探圣神，那些人发怨言。但我要你们注意，人数是在什么情况下增多的：在这些试炼之后，人数才增多，而不是之前。也要注意天主的仁慈何等伟大。那些司祭长们，就是那些愤恨、极端不满并喊叫说\Quote{他救了别人，却救不了自己}的人中，经上说：\BibleQuote{有许多人服从了信仰。}\BibleRef{玛 27:42}\par

因此，让我们也效法祂。祂接纳了他们，没有驱逐他们。同样，让我们回报那些仇敌，他们曾给我们造成无数祸害。凡我们所有的任何好处，都要分施给他们；在行善时不要忽略他们。因为如果我们应该以受苦来平息他们的愤怒，那么以行善来平息就更应如此：因为后者比前者更容易。因为以善报敌和甘心忍受比他所愿（施加）的更大的冤屈并不相同：从前者我们可以达到后者。这是基督门徒的尊贵。那些人将祂钉死，而祂正是为了向他们行善而来；他们鞭打了祂的门徒；而在这之后，祂仍使他们与祂的门徒同享尊荣，使他们同分祂的恩赐。我恳求你们，让我们做基督的效法者：在这方面，我们可以效法祂；这使人肖似天主；这超越人性。让我们紧握仁慈：她是那更高智慧的导师和教师。学会了怜悯受苦者的人，也将学会不记恨伤害；学会了这一点的人，将能甚至善待仇敌。让我们学习感受邻人所受的祸害，这样我们就能学会忍受他们施加的祸害。让我们问那虐待我们的人本人：难道他不谴责自己吗？难道他不乐意表现出更高尚的\OriginalTerm{哲学思考}{\GreekText{φιλοσοφεἵν}}吗？难道他不承认自己的行为只是冲动，是心胸狭窄、可悲的吗？难道他不愿成为受冤屈而沉默的人，而不是那些做错事而冲动发狂的人吗？他能不钦佩那忍耐受苦者而离去吗？不要以为这使人可鄙。再没有什么比傲慢和伤害他人的行为更使人可鄙；再没有什么比在傲慢和伤害下忍耐更使人可敬。因为前者是暴徒，后者是哲学家；前者不如人，后者等同天使。因为尽管他比行恶者低下，然而他若有意，也完全有能力报复。此外，前者被众人怜悯，后者被众人憎恨。那么怎样？前者将远胜于后者：因为人人都会把前者视为疯子，后者视为理智的人。他不能用恶意谈论他：是的，有人说，你担心，恐怕他并非（如你所描述的）那样。最好连心里也不要说恶言；其次，不要向别人说。那么，不要向天主祈祷反对这个人；你若听见有人讲他的坏话，要为他辩护：说，那是冲动说了这样的话，不是这个人；说，是愤怒，不是我的朋友；是他的疯狂，不是他的心。让我们这样看待每一件冒犯。不要等火燃起，而要趁它未到之前制止；不要激怒那野兽，更不要让它发怒：因为一旦火焰点燃，你便再也无法控制。因为那人叫你什么？\Quote{你这愚人和傻子。}那么，谁应得这个称呼？是被叫的人，还是叫的人？因为前者，无论多么聪明，反被人看作愚人；而后者，即使是个傻子，反被认为聪明，有哲学气质。你说，谁是傻子？是那诬告别人的人，还是那即使受此对待也毫不动心的人？因为如果真正的哲学标志是无论怎样受激都不动心，那么当没有人激怒你时却动怒——这是何等的愚蠢！我还不说，那些对邻人口出辱骂和诋毁的人将受何等严厉的惩罚。但是怎样？他叫你\Quote{下贱人和出身低微的人，可怜虫和出身卑贱的人}？这又使他反遭讥讽。因为对方将显得高尚可敬，而他实在是个可怜虫：因为把出身低微这等事作为羞辱来提起，实在是心胸狭窄；而对方将被视为一个伟大可敬的人，因为他对此类嘲弄毫不在意，就像被告知他身上有任何其他平常无关紧要的情况一样。但他若是叫你\Quote{奸夫}之类的话？对此你甚至可以发笑：因为良心不受谴责时，没有生气的理由。 * * 因为当一个人考虑到他揭露了多么恶劣和可耻的秘密时，仍然没有必要悲伤。他现在只是揭露了那些迟早每个人都会知道的事；同时，就他自己而言，他向所有人表明他不可信赖，因为他不知道如何遮掩邻人的过错：他羞辱自己更甚于羞辱对方；他堵塞了自己每一个港口；使自己未来的审判台变得可怕。因为被人厌恶的，不是（秘密被泄露的）那个人，而是那个在不该揭开帷幕的地方剥去衣服的人。但你，不要说出你知道的秘密，要沉默，如果你想赢得美名。因为这样你不仅会推翻已经说出的话并将其掩盖，还会带来另一个重大结果：你会阻止对自己的判决。有人讲你的坏话吗？你要说：\Quote{他若知道一切，就不会只说到这个程度。}——那么，你赞赏这话，并且欣然接受吗？是的，但你必须践行它。因为我们向你们讲述所有这些外邦道德家的格言，不是因为圣经中没有成百上千这样的说法，而是因为这些话更能使你们羞愧。实际上，圣经本身也习惯用这种诉诸羞耻感的方式；例如，当它说：\Quote{你们也要像外邦人一样行事。}\BibleRef{耶 35:3} 耶肋米亚先知把勒加布子孙带到众人面前，他们如何不肯违背他们父亲的命令。——米黎盎和她的同伴毁谤梅瑟，而他立即为他们求情免受惩罚；甚至不愿让人知道他的事得了申冤。\EditorialNote{户籍纪 12} 但我们却不这样：相反，这正是我们最想做的；让所有人都知道他们没有不受惩罚。我们要呼吸地上的气息到几时？——一方不能争斗。若你从两边拉开疯子，会更激怒他们；但你若从右边或左边拉开，就熄灭了冲动。打人者，若遇到不甘心挨打的人，会更加执着；但若遇到退让的人，他会更快泄气，他的拳力会落在自己身上。因为没有一个训练有素的拳击手能像那受冤屈却不还手的人一样，使对手的力量耗尽。因为对手只会在羞愧中离去，并被自己的良心和在场的所有人谴责。还有一句谚语说：\Quote{尊敬别人就是尊敬自己}；因此，羞辱别人就是羞辱自己。我再重复一遍，除非我们伤害自己，否则没有人能伤害我们；除非我使自己贫穷，否则没有人能使我贫穷。因为，来吧，让我们这样看。假设我有一副乞丐般的灵魂，而所有人都把他们的财产倾注给我，那又怎样？只要灵魂不变，一切都是枉然。假设我有一颗高贵的灵魂，而所有人都从我这里夺走财产，那又怎样？只要你不使灵魂变得乞丐般，就没有伤害。假设我的生活不洁，而所有人都说我相反的话，那又怎样？因为他们虽然这样说，心里却不如此判断我。再假设我的生活纯洁，而所有都说我相反的话，那又怎样？因为他们在自己的良心中会谴责自己：因为他们并不相信自己所说的。正如我们不应接受赞美，同样也不应接受指责。我为什么说这些话？只要愿意，没有人能谋害我们，或对我们提出任何恶名。因为现在，我问你？让他把我拖上法庭，让他提出烦扰的指控，让他甚至（若你愿意）夺走我的灵魂：那又怎样？为了片刻，无辜地遭受这些，有什么要紧呢？\Quote{可是，}你说，\Quote{这本身就是一件坏事。}但是，无辜地受苦，这本身就是一件好事。什么？你希望这苦难是应得的吗？让我再提一个智者的哲学故事。故事说，有一个人被处死了。智者的一位门徒对他说：\Quote{哀哉，他竟无辜受害！}智者转身对他说：\Quote{怎么？}他说，\Quote{你愿意他罪有应得地受害吗？}（\Person{苏格拉底}，见\Person{狄奥根尼·拉尔修}和\Person{色诺芬}\WorkTitle{回忆苏格拉底}）若翰不也是被冤杀的吗？那么，你更同情谁：是那些罪有应得而死的人，还是[他？难道你不认为他们可怜，而]你甚至钦佩他？那么，一个人从死亡本身获得了巨大收益，不但没有受害，又怎会受害呢？如果那人确实是不死的，而这件事使他成为可死的，那当然是受害；但若他是可死的，按自然规律本应在稍后死去，而他的仇敌只是加速了他的死亡，并使之与荣耀相伴，那又有何害？只要我们灵魂秩序良好，外界就没有伤害。但你不处于荣耀状态？那又怎样？财富的道理同样适用于荣耀：我若\OriginalTerm{慷慨}{\GreekText{μεγαλοπρεπής}}，便一无所需；我若贪图虚荣，所得越多，所需越多。这样，我反而最能成为杰出的人，获得更大的荣耀；那就是，如果我轻视荣耀。知道了这些，让我们感谢那白白赐给我们如此生命的那一位，并为祂的荣耀而追求它；因为荣耀属于祂，直到永远。阿们。\par

\SourceRecord{Translated by J. Walker, J. Sheppard and H. Browne, and revised by George B. Stevens.；Nicene and Post-Nicene Fathers, First Series；Vol. 11.；Edited by Philip Schaff.；Buffalo, NY: Christian Literature Publishing Co.,；1889.；<http://www.newadvent.org/fathers/210114.htm>.}

% structure: p0001 source=h1 target=section reason=page-title

\section{论\WorkTitle{宗徒大事录}第十五篇讲道}

\BibleRef{宗 6:8}\par

\BibleQuote{斯德望充满信德和能力，在百姓中行了大奇事和神迹。}\par

看，在七人之中有一个是何等卓越，赢得了头彩。\newline{}因为虽然祝圣对他和他们是共同的，但他却为自己吸引了更大的恩宠。\newline{}且要注意，此前他并未行任何（征兆和奇迹），只是在他变得为人所知之后；为表明仅有恩宠是不够的，还必须有祝圣；这样才有圣神的进一步降临。\newline{}因为他们若充满圣神，那是来自圣洗池的圣神。\newline{} \BibleQuote{然后有会堂中的几个人起来。}\BibleRef{宗 6:9} 他再次使用 \Quote{起来} 一词（\OriginalTerm{起来}{\GreekText{ἀνάστασιν}}，Hom. xiii. p. 81），以表达他们的恼怒和愤恨。\newline{}这里有一大群人。并且请注意控告形式的不同：因为加玛里耳已阻止他们以前一理由找茬，他们便提出另一指控。\newline{} \BibleQuote{并且，经上说，有称为（\OriginalTerm{被称为的}{\GreekText{τὥν} \GreekText{λεγομένων}}．Edd. \OriginalTerm{被称为的}{\GreekText{τἥς} \GreekText{λεγομένης}}）\InnerQuote{自由人}的会堂中的人，以及基勒乃人和亚历山大里亚人，并基里基雅和亚细亚的人，起来与斯德望辩论。他们敌不住他的智慧和藉他说话的神。于是他们唆使一些人说：我们听见他说亵渎梅瑟和天主的话。}\BibleRef{宗 6:9} \newline{}为坐实控告，这句话是 \Quote{他说反对天主，反对梅瑟。} \newline{}他们也为此目的而辩论，为要逼他说出什么。但他现在讲论得更加公开，或许谈到了神律的终止；或没有明说，但暗示了这一点：因为他若明说，就不需要唆使的人，也不需要假见证。\newline{}会堂是多样的：［即 \Quote{自由人的}］：基勒乃人的，即亚历山大里亚以外地区的人［\Quote{亚历山大里亚人的}等］。他们似乎也按不同的民族有自己的会堂；因为许多人留在那里，以免需要不断旅行。自由人或许是罗马的自由民。那里有许多外邦人居住，所以他们有自己的会堂，在那里宣读法律。\newline{} \Quote{与斯德望辩论。} 请注意他，不是自任为师，而是被迫如此。奇迹再次为他招来敌意；但当他辩论得胜时，竟然成了假见证！因为他们不愿杀他们无法容忍的人。\newline{} \Quote{他们不能抵抗等：于是唆使一些人。} 处处不在掌握之中，但藉着一个判决，为要也伤害他们的名誉：并抛开那些人（宗徒），攻击这些人（门徒），以为这样也能恐吓那些人。他们不说 \Quote{他说}，而是说 \Quote{他不住地说。他们煽动了百姓、长老和经师，一同来捉住他，解到公议会，并设立假见证说：这个人不住地说亵渎圣所和法律的话。}（12-13节）\Quote{不住地说}，他们说，好像他以此为业。\BibleQuote{因为我们听见他说：这个纳匝肋人耶稣要毁灭这地方，并要更改梅瑟所传授给我们的规例。}\BibleRef{宗 6:14} \Quote{耶稣}，他们说，\InnerQuote{纳匝肋人}，作为侮辱性称呼，\InnerQuote{将要毁灭这地方，并要更改规例。} 这也是他们对基督所说的话：\BibleQuote{你这拆毁圣殿的。}\BibleRef{玛 27:40} 因为他们对圣殿（他们确实选择了\OriginalTerm{迁移}{\GreekText{μετοικεἴν}}到它附近）和梅瑟的名号怀有极大的敬意。控诉是双重的。如果祂\Quote{更改规例}，祂也会引入其他规例代替：请注意这个控诉是何等尖刻，充满危险。\newline{} \BibleQuote{所有坐在公议会的人都注目看他，见他的面貌好像天使的面貌。}\BibleRef{宗 6:15} 所以即使是地位较低的人也可能发光。我问，这个人比宗徒们少什么？他不缺奇迹，他表现出了极大的胆量。——\BibleQuote{他们看见他的面貌，经上说，好像天使的面貌。}\BibleRef{出 34:30} 这就是他的恩宠，这就是梅瑟的荣耀。天主使他面容如此充满恩宠（\OriginalTerm{恩宠的}{\GreekText{ἐπίχαριν}}），因为他即将说话，从而凭他的面容立刻使他们敬畏。因为确实有一些面孔充满了属灵的恩宠，对爱他们的人是可爱的，对恨者和敌人是可畏的。经文也提到了他们为何容忍他发言的原因。——\newline{} \BibleQuote{果真有这些事吗？}\BibleRef{宗 7:1} 请注意，问话很温和，为要造成某种大恶。为此，斯德望也以温和的语气开始他的演说，说：\BibleQuote{诸位仁人弟兄、诸位父老，请听！荣耀的天主曾显现给我们的祖先亚巴郎，当他还在美索不达米亚，在迁居哈兰以前。}\BibleRef{宗 7:2} 一开始他就推翻了他们的狂妄，并使他所说的显出圣殿算不得什么，规例也算不得什么，而他们尚未察觉他的用意：同时也表明他们不能胜过宣讲；并且天主总从无能（\OriginalTerm{无能的}{\GreekText{ἀμηχάνων}}）的事物为自己安排有能（\OriginalTerm{有能的}{\GreekText{εὐμήχανα}}）的工具。请注意这些线索如何构成整篇演说之纹理：并且他们虽然始终享受了极致的恩惠，却仍以相反的行为回报了他们的恩主，而现在正在尝试不可能的事。\newline{} \Quote{荣耀的天主曾显现给我们的祖先亚巴郎，当他还在美索不达米亚，在来到哈兰以前。} 当时既没有圣殿，也没有祭献，但天主却赐予了亚巴郎一个神视，而他的祖先还是波斯人，又处在异乡。他在演说开头称祂为\Quote{荣耀的天主}极为恰当：因为祂使没有荣耀的人成为荣耀。\Quote{因为}（他说）\Quote{是祂使他们成为荣耀的，祂也会使我们如此。} 请注意他如何首先将他们从肉体的事物、从地方引开，因为地方正是问题所在。\Quote{荣耀的天主}，他说：又暗示，祂并不需要来自我们的荣耀、来自圣殿的荣耀：因为祂自己是荣耀的泉源。你们不要想，他说，这样来光荣祂。\Quote{并从你的亲族中出来。} 那么圣经怎么说，说亚巴郎的父亲愿意出去？由此我们得知，是由于亚巴郎的神视，他的父亲才被感动加入迁移。 \BibleRef{创 11:31} \BibleQuote{你要离开你的故乡和你的亲族，往我所要指给你的地方去。}\BibleRef{宗 7:3} 这显出这些人与亚巴郎的子孙相去多远，他多么服从。\Quote{并从你的亲族中出来。} \OriginalTerm{沉重的}{\GreekText{φορτικὰ}}反思，既是他承受了劳苦，而你们收割果实，又是你们所有祖先都处于困境。\newline{} \Quote{那时他离开了加色丁人的地方，住在哈兰；他父亲死了以后，天主又叫他迁移到你们现在所住的地方。在这里他没有给我留下产业，连脚掌那么大的一块也没有。}（4-5节）请看，他如何将他们的思想从（拥有）这地提升。因为如果祂说（祂会赐予：显然一切都是从他来的），而不是从他们自己来的。因为他来了，离开了亲族和故乡。那么为什么天主没有赐给他呢？这其实是另一块地的预象。\Quote{但天主应许把这地赐给他。} 你们可看到，他并非只是重新拾起话头？\Quote{没有赐给他}，他说，\Quote{却应许了；并要赐给他的后代，那时他还没有孩子。} 再次显出天主的作为：祂从不可能中成就一切。因为这里有一个在波斯的人，如此遥远，这个人天主说祂要使成为巴勒斯坦的主人。但让我们回顾先前所说的话。\par

请问，那恩宠如何绽放于\Person{斯德望}的面容之上？（复述。）作者在上文为他作证，说他\BibleQuote{充满信德}。见：\BibleRef{宗 6:8}。因为有可能拥有一种恩宠，并不显于治病的异能：因为圣神的恩宠以\OriginalTerm{如此的}{\GreekText{τοιὣσδε}}方式赐予一人，见：\BibleRef{格前 12:8}。但在这里，在我看来，经文说他容貌也令人喜爱：\BibleQuote{他们看见他的面容好像天使的面容。}\BibleQuote{充满信德和能力}\BibleRef{宗 6:15}，这也是\Person{巴尔纳伯}的特征——\BibleQuote{他是个好人，充满信德和圣神}\BibleRef{宗 11:24}。由此我们得知，真诚纯洁的人比众人更蒙救恩，而且这些人也更加可爱。\BibleQuote{于是他们唆使人说：我们听见他说亵渎的话。}见：\BibleRef{宗 6:11}。在宗徒们的事上，他们恼怒宗徒宣讲复活，且有许多人归向他们；但在此事上，他们恼怒人们得到疾病痊愈\BibleRef{宗 4:2}。他们本应为这些事感恩，却反以此为指责之由：何等疯狂！他们指望以言语胜过那些以作为胜过他们的人！这与他们在基督身上的行径如出一辙，总是迫使他们用言语辩解。因为他们羞于无故逮捕他们，毕竟无任何可指控之事。且注意，并非将被告带到审判前的人亲自作证指控他们，否则他们会被驳倒；而是雇请他人，以免显得是纯粹暴力行为。这与他们在基督身上的做法如出一辙。且观察宣讲的能力，尽管他们不仅受鞭打，还被石头砸，仍能得胜：不仅是他们身为凡人被拖上法庭，且四面受敌，而敌人自己也作证，他们不仅被打败，甚至\BibleQuote{不能抵抗}\BibleRef{宗 6:10}，尽管他们极其无耻：宣讲如此有力地击败了他们，尽管他们用荒谬的捏造（例如说祂有魔鬼——祂可是驱魔者！）。因为这场战斗不是人的，而是天主与人的战斗。有许多人联合起来；不仅耶路撒冷的人，还有其他人。见：\BibleRef{宗 6:9}。因为\BibleQuote{我们听见他}，他们说，\BibleQuote{说亵渎梅瑟和天主的话。}见：\BibleRef{宗 6:11}。无耻之徒们！你们行亵渎之事，却毫不在意。所以提到梅瑟——因为他们对天主的事并不在意：他们总是提及梅瑟：\BibleQuote{这位梅瑟，曾领我们出来。}见：\BibleRef{宗 7:30}。\BibleQuote{他们煽动百姓。}见：\BibleRef{宗 6:12}。民众反复无常！然而一个亵渎者怎能如此成功？一个亵渎者怎能行这样的奇迹？但不受管束的民众却使辩不过他们的人变得强大。这正是最令他们烦恼的事。\BibleQuote{我们听见他}，他们说，\BibleQuote{说亵渎梅瑟和天主的话}\BibleRef{宗 6:13}；又说：\BibleQuote{这人不住地说亵渎这圣地和法律的话}，并且加上了：\BibleQuote{传统}\BibleQuote{就是梅瑟传给我们的}\BibleRef{宗 6:14}；是梅瑟，而非天主。他们假设他有意推翻他们的\OriginalTerm{生活方式}{\GreekText{πολιτείας}}，便也指控他不敬。但为表明如此之人不可能说这样的话，且严厉地［\BibleQuote{于是}，经文说，\BibleQuote{所有在公议会里的人，注目看他，见他的面容好像天使的面容}\BibleRef{宗 6:15}］：他面容如此温和。因为，在那些未被诬告的案例中，圣经并未提及此类事情；但在此，既然全是诬告，天主合情合理地藉此人的容貌来纠正。因为宗徒们确实未被诬告，而是被禁止；但这人却遭诬告，因此他的容貌首先为他辩护。这甚至使\OriginalTerm{司铎}{priest}感到羞愧。\BibleQuote{他便说}等等\BibleRef{宗 7:1}。他在此表明，应许是在地方之前、割损之前、祭献之前、圣殿之前作出的，并且他们领受割损或法律并非出于功德，而是土地本身是唯独服从的赏报。而且，应许也并非在实行割损时实现。这些也是预像，同样，他奉天主命令离开家乡——并不违反法律（因为家乡就是天主引导的地方）：\BibleQuote{于是他出来}，经文说，\BibleQuote{离开加色丁人的地}\BibleRef{宗 7:4}——而且，若仔细查看，犹太人是波斯裔；并且，即使没有奇迹，也必须遵行天主的吩咐，无论后果如何艰难；因为圣祖遵从天主的命令，离开了父亲的坟墓和他所有的一切。但如果\Person{亚巴郎}的父亲未被允许与他同享迁移至巴勒斯坦的特权，因为他不配：那么儿女们（最终将被排除）更是如此，尽管他们可能已在路上走了很远。\BibleQuote{天主应许}，经文说，\BibleQuote{要将此地赐给他和他的后裔。}\BibleRef{宗 7:5}在此显示出天主的仁慈和\Person{亚巴郎}的信德何等伟大。因为\BibleQuote{当时他还没有孩子}这句话确实显示了他的服从和信德。\BibleQuote{应许赐给他和他的后裔。}然而事实却相反：他来了以后，连\BibleQuote{立足之地}都没有，也没有孩子；这些事都与他的信德相反。\par

这些事既已看见，让我们也照样，无论天主应许什么，都接受它，尽管事件可能相反。然而在我们身上，它们并不相反，反而十分合适。因为哪里有望，哪里就有相反的事发生时，它们就真是相反的；但在我们身上恰恰相反：因为祂已告诉我们，我们在此世会有磨难，但我们的安息在来世。我们为何混淆时期？为何颠倒一切？请问，你受痛苦，生活在贫困和沮丧中吗？不要烦恼：因为若你在那世界注定受痛苦，那才值得烦恼；至于这现世的痛苦，它是安息的原因。祂说：\BibleQuote{这病不是向着死亡的。}\BibleRef{若 11:4} 那种痛苦是惩罚；这种是管教和纠正。这现世生命是一场竞赛：若是如此，战斗就是我们现在的事业：这是战争和战役。在战争中，人不寻求安息，在战争中，人不寻求美食，人不为财富焦虑，那时他不关心妻子；他只关注一件事：如何战胜敌人。这也该是我们的关注：如果我们战胜，带着胜利归来，天主将赐予我们一切。让这是我们唯一的学习，如何战胜魔鬼：尽管毕竟不是我们自己的学习成就此事，而是天主的恩宠成就全部。让这是我们唯一的学习，如何吸引祂的恩宠，如何把那种帮助引向我们。\BibleQuote{如果天主偕同我们，谁能反对我们？}\BibleRef{罗 8:31} 让我们专务一事：使祂不成为我们的敌人，使祂不远离我们。\par

受苦难并不是坏事；坏事是犯罪。这是最严重的痛苦，即使我们过着奢华的生活——且不说来世，就在今生也是如此。思考良心如何被悔恨刺痛，这是否比任何酷刑更糟糕！我想探问那些生活在邪恶中的人（\OriginalTerm{在邪恶中}{\GreekText{ἐν} \GreekText{κακοῖς}}），他们是否从不反思自己的罪，是否不战栗，是否不恐惧和痛苦，是否不认为那些生活在节制中的人、山中的人、严格规则中的人（\OriginalTerm{在严格规则中}{\GreekText{τοὺς} \GreekText{ἐν} \GreekText{πολλῇ} \GreekText{φιλοσοφίᾳ}}）是有福的？\par

\SourceRecord{Translated by J. Walker, J. Sheppard and H. Browne, and revised by George B. Stevens.；Nicene and Post-Nicene Fathers, First Series；Vol. 11.；Edited by Philip Schaff.；Buffalo, NY: Christian Literature Publishing Co.,；1889.；<http://www.newadvent.org/fathers/210115.htm>.}

% structure: p0001 source=h1 target=section reason=page-title

\section{\WorkTitle{宗徒大事录第十六篇讲道}}

\BibleRef{宗 7:6-7}\par

\BibleQuote{天主这样说过：他的子孙要侨居外方，在那里要受人虐待四百年之久。以后，我要惩罚那奴役他们的民族，天主又说：以后，他们要出来，在这里事奉我。}\par

看，那预许已经赐下多少年了，预许的方式如何，且没有祭献，没有割损！\newline{}这里他表明，天主如何容许他们受苦，并非祂有什么可指控他们的。\newline{}\Quote{他们要奴役他们}等等。\newline{}但尽管如此，他们作这些事并非不受惩罚。\newline{}\Quote{他们所奴役的民族，我要审判，天主说。}\newline{}为了表明他们不可凭此判断谁是虔敬的（由于他们说：\Quote{他信赖天主，让祂解救他吧}）\BibleRef{玛 27:43}。——那位预许者，那位赐予土地者，先容许了灾祸。\newline{}同样，现在，虽然祂预许了天国，却仍容许我们在试探中受磨炼。如果这里的自由要到四百年后才实现，那么关于天国，又何足为奇呢？然而祂实现了它，时间的流逝并未使祂的话落空。再者，他们所遭受的也不是普通的奴役。此事不仅以惩罚那些人（压迫者）而告终；祂说，他们自己也要享有伟大的救恩。这里他也提醒他们所享有的恩惠。\newline{}\Quote{祂赐给他割损的盟约：于是，他生了依撒格。}这里他降格讨论较低俗的事。\Quote{他在第八天给他行了割损：依撒格（生了）雅各伯，雅各伯生了十二位圣祖。}\BibleRef{宗 7:8}。——这里他似乎现在暗示了预像。\newline{}\Quote{圣祖们因嫉妒，将若瑟卖到了埃及。}\BibleRef{宗 7:9}这里再次有基督的预像。虽然他们挑不出他什么毛病，而且他本是专程来给他们送粮食的，他们却这样虐待他。然而这里预许再次得到实现，虽然过了很长时间。\Quote{天主与他同在}——这也是为他们——\Quote{并救他脱离了一切苦难。}\BibleRef{宗 7:10}。\newline{}他表明他们在不知不觉中促成了预言的实现，而他们自己是原因，灾祸落到了他们自己头上。\Quote{使他在埃及王法郎眼前获得恩宠和智慧，恩宠}在外邦人眼中，施予他这奴隶、这俘虏：他的弟兄卖了他，而这（外邦人）却尊敬了他。\newline{}\Quote{那时埃及和迦南全地发生饥荒，大灾难，我们的祖先找不到食物。但雅各伯听说埃及有粮食，就先打发我们的祖先去。第二次，若瑟才被他的弟兄认出。}\BibleRef{宗 7:11}。\newline{}他们下去购买谷物，样样都得依赖他。那么他做什么？[ \Quote{他向他的弟兄表明了自己：} ]他的友善不止于此；他还将他们介绍给法郎，并带他们下了那地。\Quote{若瑟的亲属被报告给法郎。于是若瑟派人去请他父亲雅各伯和他所有的亲属，共七十五人。雅各伯就下到埃及，他和我们的祖先都死在那里，被运到舍根，安放在亚巴郎用银钱从舍根父哈摩尔的子孙那里买来的坟墓里。但天主向亚巴郎所发誓的预许时期临近时，人民在埃及繁衍增多，直到另一位不认识若瑟的君王兴起。}\BibleRef{宗 7:13}。\newline{}然后，又是新的失望（\OriginalTerm{无望}{\GreekText{ἀνελπιστία}}）：首先，饥荒，但他们度过了；其次，落入敌人手中；第三，被君王消灭。然后（表明）天主的万般巧妙（\OriginalTerm{上智的安排}{\GreekText{εὐμήχανον}}），经上说：\Quote{那时，}又说：\Quote{梅瑟出生了，且极其俊美。}\BibleRef{宗 7:20}\newline{}如果说前一件事——若瑟被他的弟兄出卖——是奇妙的，那么这里又有一件事更加奇妙：那君王竟抚养了将要推翻他统治的人，而他自己正是将要灭亡的人。\newline{}你看到自始至终，仿佛是死者复活的一种预演吗？但天主自己行事，与事物按\OriginalTerm{人的意图}{\GreekText{προαίρεσις}}发生，是不同的。因为这些事情固然是按人的意图发生的〔但复活本身是独立的〕。——经上说：\Quote{他很有能力，}又说：\Quote{在言语和行为上，}\BibleRef{宗 7:22}：这位本应死去的人。\newline{}然后他又表明他们对自己的恩人多么忘恩。因为，正如前一个例子中，他们因受伤害的若瑟而得救，这里他们又因另一个受伤害的人，即梅瑟而得救。\Quote{当他满四十岁时，}等等。虽然他们并没有实际杀死他，但意图上他们杀了他，正如前例中其他人所做的那样。在那里，他们从本乡卖到外乡；在这里，他们从外乡赶往外乡；前一个例子中，他是送食物的人；这里，他是给予善言的人；他是那人在天主之下欠他生命的人！注意这显示加玛里耳的那句话：\Quote{如果这是从天主来的，你们不能推翻它。}\BibleRef{宗 7:39}\newline{}看，被谋害的人最终成了谋害者的救恩：人民谋害自己，也被他人谋害，却仍得救！一场饥荒，没有消灭他们：不仅如此，他们竟借着他们本以为会（被他们）消灭的那人得救。一道御旨，没有消灭他们：反而在那位\Quote{认识}他们的人死后，他们的人数增长最多。他们想杀死自己的救主，却终究无力做到。你看到吗，魔鬼试图使天主的预许落空的手段，恰恰促进了预许的实现？\par

\BibleQuote{天主这样说，}等等。（复述，第6、7节。）这也适合在此处说：天主在方式和方法上是富有的，将我们从这里提升。因为首先这显示了天主资源的丰富，即在其逆转\OriginalTerm{逆转}{\GreekText{ἀποστροφῇ}}中，这个民族增加了，虽然被奴役，受虐待，并被企图消灭。这就是应许的伟大之处。因为如果它在本土增加，就不会如此奇妙。此外，他们在异国的时间也不短：而是四百年。由此我们学到了哲学忍耐\OriginalTerm{哲学忍耐}{\GreekText{φιλοσοφίαν}}的伟大教训：——他们待他们不像主人待奴隶，而是像敌人和暴君——并且他预言他们将被置于大自由中：因为这就是那句话的意思，\BibleQuote{他们必事奉（我），并且要回到这里}\OriginalTerm{返回这里}{\GreekText{ἐνταὕθα} \GreekText{ἐπανελεύσονται}}；而且不受惩罚。——并且要注意，他如何在似乎让步于割礼的同时，实际上并不允许它什么\BibleRef{宗 7:8}；因为应许在它之前，而它随后。——\BibleQuote{列祖}，他说，\BibleQuote{心生嫉妒}\BibleRef{宗 7:9}。在无伤大雅之处，他宽容\OriginalTerm{宽容}{\GreekText{χαρίζεται}}他们：因为他们也以此自夸。——并且他显示，圣人并非免于磨难，而是在他们的磨难中获得了帮助。并且这些人自己也有助于促成结果，他们想要缩短这些（苦难）：正如这些使\Person{若瑟}更荣耀：正如王对待\Person{梅瑟}，下令杀害婴孩：因为若他不下令，就不会这样：正如那个（希伯来人）将\Person{梅瑟}赶入流放，使他得以在那里见到异象，成为配得的人。同样，那个被卖为奴的人，他使他在那里为王，在那里他被认为是一个奴隶。同样，\Person{基督}也以他的死证明了他的权能：同样，他在那里为王，在那里他们出卖了他。\BibleQuote{赐他恩宠和智慧，}等等。\BibleRef{宗 7:10} 这不仅是为了尊荣，而且使他对自己权能有信心。\BibleQuote{立他为埃及和全家之主。}\BibleQuote{如今发生了饥荒，}等等。因为饥荒——他正在做这样的准备——\BibleQuote{共有七十五人，}他说，\BibleQuote{\Person{雅各伯}下到埃及，死了，他和我们的列祖，被运到\Place{舍根}，安放在\Person{亚巴郎}用银子从\Place{舍根}之父\Person{厄默尔}的儿子们那里买来的坟墓里。}\BibleRef{宗 7:11}。这显示，他们甚至不能自主拥有一块墓地。\BibleQuote{但应许的时候临近，就是\Person{亚巴郎}所誓许的，百姓在埃及增长增多，直到另一位不认识\Person{若瑟}的王兴起}（第17、18节）。注意，并不是在四百年间他使他们增多，而是（仅）在结局临近的时候。而且四百年已经过去，不，更多的时候，在埃及。但这就是其奇妙之处。\BibleQuote{他狡猾地对待我们的亲族，虐待我们的列祖，要他们丢弃婴孩，使他们不能存活。}\BibleRef{宗 7:19}\BibleQuote{狡猾地对待：}他暗示他们不愿公开消灭他们：\BibleQuote{要他们丢弃婴孩，}它说。\BibleQuote{那时\Person{梅瑟}出生，极其俊美。}\BibleRef{宗 7:20} 这就是奇妙之处，那将成为他们捍卫者的人，不是在这些事之后或之前，而是在风暴\OriginalTerm{风暴}{\GreekText{θυμῷ}}的中央出生。\BibleQuote{在家里被抚养了三个月。}但当人的帮助被绝望，他们抛弃了他时，天主的恩惠就显明了。\BibleQuote{他被抛弃时，\Person{法郎}的女儿捡起他，养为自己儿子。}\BibleRef{宗 7:21} 没有一句话提到圣殿，没有一句话提到祭献，而所有这些眷顾正在发生。他在蛮夷之家被抚养。\BibleQuote{\Person{梅瑟}学习了埃及人的一切智慧，说话行事都有能力。}\BibleRef{宗 7:22}\BibleQuote{受过训练，}在纪律和学识上。\BibleQuote{他满了四十岁。}\BibleRef{宗 7:23} 他在那里四十年，并未因割礼而被发现。注意，他们在安全中，如何忽略自己的利益，他和\Person{若瑟}都是，为了拯救他人：\BibleQuote{他满了四十岁，心中起意去看望他的弟兄以色列子民。看见一人受屈，就保护他，为受欺压的人伸冤，打死了那埃及人：他以为弟兄们必明白天主藉他的手拯救他们，但他们不明白。}\BibleRef{宗 7:23} ——看，到此时为止，他还没有得罪他们；他们听他讲这些事。而且\BibleQuote{他的面容}，我们读到，\BibleQuote{好像天使的面容}\BibleRef{宗 6:15}。——\BibleQuote{他以为，}等等。然而他的捍卫是通过行为表明的；这里需要什么智慧？但尽管如此，\BibleQuote{他们不明白。第二天，他们争闹时，他向他们显现，想要劝和他们，说：诸位，你们是弟兄，为什么彼此亏负？}\BibleRef{宗 7:26} 你注意他多么温和地对他们说话？他在别人身上显示愤怒，却在自己身上显示温和。\BibleQuote{但那欺负邻舍的人推开他，说：谁立你作首领和审判官呢？难道你要杀我，像昨天杀那埃及人一样吗？}注意：这正是他们对\Person{基督}说的话：\BibleQuote{谁立你作首领和审判官呢？}犹太人如此惯于在受恩惠时亏负（他们的恩人）！再次，注意那可怕的卑鄙：(\OriginalTerm{卑鄙}{\GreekText{μιαρίαν}} al. \OriginalTerm{邪恶}{\GreekText{μοχθηρίαν}}, Sav. marg.) \BibleQuote{像昨天杀那埃及人一样吗？\Person{梅瑟}听见这话就逃走了，寄居在\Place{米德扬}地，在那里生了两个儿子。}\BibleRef{宗 7:29} 但流亡并未消灭天主的眷顾计划，正如死亡（即\Person{基督}的死亡）也未消灭。\par

\Quote{过了四十年，在西乃山的旷野里，有一位天使在火焰中从荆棘里向他显现。} \BibleRef{宗 7:30} 你注意到时间并未阻止它吗？因为当他流亡、寄居异地，在国外度过了许多时日，甚至有了两个儿子，不再期望回去的时候，天使向他显现。他称天主子为天使，也如同称祂为人。显现于旷野，而不是圣殿。请看多少奇迹正在发生，却没有提及圣殿，没有提及祭献。而且这里也不仅是在旷野，而是在荆棘中。\Quote{梅瑟一见，就惊异这异象；他正上前观看时，有主的声音向他说：}\BibleRef{宗 7:31} 看！他也堪当听见那声音。\Quote{我是你祖先的天主，亚巴郎的天主，依撒格的天主，雅各伯的天主。}（32-33节）看！祂如何表明祂就是那一位，不是别的，正是\Quote{亚巴郎的天主，依撒格的天主，雅各伯的天主}——祂就是\Quote{大谋士的天使}。（\BibleRef{依 9:6}。七十子译本：\Quote{奇妙、谋士}，英文译本。）这里他表明天主在此彰显了何等伟大的慈爱。\Quote{梅瑟就战栗起来，不敢注视。主向他说：\InnerQuote{将你脚上的鞋脱下，因为你站的地方是圣地。}} 没有提及圣殿，这地方因基督的显现和作为而成为圣地。这比至圣所里的那地方要奇妙得多：因为那里从未说过天主以这种方式显现，梅瑟也没有这样战栗。然后祂的慈爱是何等伟大。\Quote{我看见了，我看见了我在埃及的百姓所受的苦难，听到了他们的哀号，就下来拯救他们。现在你来，我要打发你往埃及去。}\BibleRef{宗 7:34} 请看，他如何表明，天主用恩惠、惩罚和奇迹吸引他们归向祂，但他们仍然如故。他们学到了天主无所不在。\par

听到这些事，让我们在苦难中投奔祂。祂说：\Quote{他们的哀号，我已听见了：}不只是\Quote{因为他们遭遇的灾难。}但若有人问：既然如此，祂为什么容许他们在那里受恶待？首先，对每个义人来说，受苦是他得赏报的原因。其次，祂为什么磨难他们：是为了显示祂的权能，祂能（做一切），不仅如此，祂还要训练他们。请看事实：当他们身在旷野时，他们\InnerQuote{养肥了，长壮了，饱享口福，就踢踏了}\BibleRef{申 32:15}；舒适安逸永远是祸患。因此，从起初祂就对亚当说：\Quote{你必汗流满面，才有饭吃。}\BibleRef{创 3:19} 这也是为使他们经历诸多痛苦后进入安息，得以感谢天主。因为苦难是极大的善。请听先知的话：\Quote{为我的益处，你磨难了我。}\BibleRef{咏 118:71} 若苦难对伟大而奇妙的人是极大的善，那么对我们更是如此。你若愿意，让我们考察苦难本身的本质。设若有一个人极其欢喜、快活、放纵欢乐：还有什么比这更失态、更愚昧的呢？设若有一个人忧伤、沮丧：还有什么比这更真正富有哲思的呢？因为经上记载：\Quote{往哀悼之家去，比往欢笑之家更好。}\BibleRef{训 7:2} 然而，很可能你不喜欢这话，想要回避它。不过，让我们看看亚当在乐园里是什么样子，后来是什么样子；加音从前是什么样子，后来是什么样子。灵魂不能稳守其位，反而像被急流（\OriginalTerm{水流}{\GreekText{ῥεύματος}}，某些编者读作\OriginalTerm{风}{\GreekText{πνεύματος}}）托起浮起，因快乐而飘飘然，毫无定力：轻易许诺，轻率应承；思绪纷乱：不合时宜的笑，无缘无故的欢乐，空洞无物的闲谈。何必说别人？让我们拿某位圣人来看，看看他在快乐时是什么样子，在痛苦中又是什么样子。我们看看达味本人吧：当他快乐欢喜时，因他众多的战利品、胜利、冠冕、奢侈生活和自信，请看他说了什么、做了什么：\Quote{我曾在顺境中说：我永不会动摇。}\BibleRef{咏 29:6} 但当他陷入苦难时，请听他说的：\Quote{若祂对我说：我不喜悦你；看，我在这里，愿祂照祂眼中看为好的待我。}\BibleRef{撒下 15:26} 还有什么比这些话更富有哲思呢？他说：\Quote{凡天主所喜悦的，就这样成就吧。}他又对撒乌耳说：\Quote{若上主激动你来攻击我，愿你的祭品蒙悦纳。}\BibleRef{撒上 26:19} 那时，他身处苦难，甚至宽恕了他的仇敌；但后来，连朋友、甚至那些没有伤害过他的人，他都不宽恕。再者，雅各伯在苦难时说：\Quote{若上主赐给我食物吃、衣服穿。}\BibleRef{创 28:20} 同样，诺厄的儿子先前没有做任何类似的事；但当他不再为安全担忧时，你听他变得多么放纵。\BibleRef{创 9:22} 再看希则克雅：他在苦难时为了获得解救做了些什么？他穿上苦衣等类的东西；但他在安乐时，却因心高气傲而跌倒。\BibleRef{列下 19:20} 因为圣经说：\Quote{你吃饱了，喝足了，富足了，要谨慎。}\BibleRef{申 6:11} 因为富足的地位如同在悬崖边缘般危险。祂说：\Quote{你要谨慎。}当以色列子民受磨难时，他们人数更加增多；但当祂任由他们时，他们就全都败亡了。何必说古人的例子？在我们这个时代，让我们看看，如果你们愿意，难道不是这样吗：大多数人在顺境时变得自高自大、对所有人都充满敌意、易怒，而权力在他们手中；若权力被夺去，他们就变得温和、谦卑（如）常人，意识到自己的本性。因此圣经说：\Quote{骄傲环绕他们直到尽头；他们的邪恶仿佛从脂肪中流出。}\BibleRef{咏 72:6}\par

我这样说，是为了叫我们不要事事以享乐为目的。那么，保禄为什么说：\Quote{常常喜乐？}他不只是简单地说：\Quote{喜乐，}而是加上：\Quote{在主内。} \BibleRef{斐 4:4}这是最大的喜乐，就是宗徒所喜乐的；监牢、鞭打、迫害、恶名以及一切痛苦之事，都是这喜乐的根源、根本和机遇；因此它也导致幸福的结局。但世俗的喜乐相反，以甜蜜开始，以苦涩结束。我也不禁止在主内喜乐，反而诚恳地劝勉这样行。宗徒们受鞭打，却喜乐；被捆绑，却感恩；被石头砸，却宣讲。这正是我也愿意有的喜乐：它根本不源于肉体，而源于属灵之事。凡按世俗方式喜乐的人，不可能也按属神的方式喜乐：因为凡按世俗方式喜乐的人，他的喜乐在于财富、奢华、荣誉、权势、傲慢；但按天主旨意喜乐的人，他的喜乐在于为天主而受辱、贫穷、匮乏、禁食、谦卑。你看见了吗？两者的（喜乐）基础是多么对立！在此没有喜乐，也就是没有悲伤：在此没有悲伤，也就是没有快乐。事实上，这些（痛苦之事）才产生真正的喜乐，因为其他事物只有喜乐之名，却全然由痛苦组成。傲慢之人忍受何等痛苦！他在傲慢中如何\OriginalTerm{被截断}{\GreekText{διακόπτεται}}，为自己招来无数羞辱、众多仇恨、巨大敌意、极度的恶意和许多恶眼！无论是被比他大的人侮辱，他悲伤；还是他无法抵抗所有人，他苦恼。而谦卑的人生活充满享受：不期望从任何人得荣誉，若得荣誉则欣喜，若不得也不悲伤。他满足于被尊敬；但更重要的是，无人羞辱他。不求荣誉而仍得荣誉——其中的享受必定是巨大的。但在另一种情况下，恰恰相反：他寻求荣誉，却得不到。荣誉带来的快乐对于追求者与不追求者并不相同。追求者无论得到多少，都认为自己一无所获；不追求者，即使你给他极少，他也仿佛得到了一切。再者，生活富裕奢华的人有无数的事务，即使他的收入流入得如此容易，仿佛源自丰沛的泉源，他仍担心奢华生活带来的邪恶和未来的不确定性；但另一种人则总是处于安全与享受中，因为他已习惯饮食的匮乏。因为他并不因未能享用丰盛的筵席而如此哀叹，反而因不担心未来的不确定性而得意洋洋。但奢华生活带来的邪恶是何等多且大，无人不知；然而，现在必须提及它们。双重的战争：在身体和灵魂上；双重的风暴；双重的疾病；不仅于此，而且因为它们既不可治愈，又带来巨大祸患。节俭则不然：它带来双重的健康，双重的益处。我们读到：\BibleQuote{健康的睡眠在于节食。}\BibleRef{德 31:20}因为处处，保持适度令人愉悦，超过适度则不再令人愉悦。比方说：在小火花上放一大堆柴，你便再也看不到火焰闪烁，只有大量令人不快的烟雾。把一个超重的担子放在一个非常强壮高大的人身上，你会看到他连担子一起扑倒在地。船上装载过多的货物，你便注定了一场惨重的海难。此处正是如此。因为正如在超载的船上，水手、舵手、船头的人和乘客一片混乱，他们把上面的和下面的东西抛入海中；同样，这里的人也因上吐下泻，损害体质，毁灭自己。而最可耻的是，嘴被用来做下部的职务，那个肢体就成了更可耻的器官。但若嘴已如此可耻，试想灵魂中又是何等情形！因为那里确实全是迷雾、风暴、黑暗，思想因如此拥挤压榨而喧闹，灵魂自身因所受的虐待而呼喊：所有（部分和官能）彼此抱怨，恳求、哀求，让污秽从某处排出。排出之后，骚动仍未结束；然后发烧和疾病接踵而来。你说：\Quote{这怎么可能呢，人们怎能看见这些奢华享乐者，骑着马，处境良好？}你说：\Quote{这是什么胡言乱语，你告诉我们疾病？我才是患病的，我才是受折磨的，我才是令人恶心的，而我什么吃的都没有。}唉！听了这样的话，实在令人悲伤。但那些患痛风的人、被抬在担架上的人、被绷带缠绕的人，我问你，我们从哪一类人身上看到这些？的确，若不是他们认为这是侮辱，以为我的话语可耻，我早已指名道姓了。\Quote{但是，他们之中有些人的身体情况也不错。}因为他们不只沉溺于奢华，也从事劳动。否则，请指给我一个人，他什么也不做，只是养肥自己，躺着毫无痛苦，没有任何焦虑。因为即使不计其数的医生聚集，也无法将他从疾病中解救出来。这是违反本性的。因为我要对你进行一次医学讲解。进入腹中的物质，并非全部变成营养；因为即使在食物本身，也并非全部有营养，一部分在消化过程中变成粪便，一部分变成营养。如果消化过程运作完美，结果便是如此，每种成分各归其位；有益有用的部分进入合适的地方，而多余无用的部分自行分离并排出。但如果数量过大，那么连其中的营养部分也变得有害。为了让你更明白，我用比喻来说：在小麦中，一部分是精粉，一部分是粗粉，一部分是麸皮；如果磨能碾碎（放入的东西），它就把这些都分离出来；但如果放入太多，所有东西便混在一起。再者，葡萄酒如果经过适当的形成过程并在季节的适当影响下，那么起初全部混在一起，一会儿一部分沉为渣滓，一部分浮为泡沫，一部分留作享用者之乐，这部分是好的，不会轻易变化。但他们所谓的\Quote{营养}，既不是酒，也不是渣滓，而是一切混在一起。在河流中也可看到同样情形，当水流形成漩涡洪水时。此时我们看到鱼漂浮在水面上，死了，它们的眼睛首先被泥泞的黏液弄瞎：我们也是如此。因为当暴饮暴食如同暴雨一般浸透内部时，它使一切陷入混乱，使那此前健康、生活在纯洁元素中的\OriginalTerm{理性能力}{\GreekText{λογισμοὶ}}像死了一样漂浮在水面。既然通过这些例子，我们已经展示了危害有多大，让我们不要再为那些我们本应认为悲哀的事而羡慕那些人，也不要再为我们本应感到幸福的事而自怜，而要以满足的心迎接知足。难道你没有听见医生也告诉你，\Quote{匮乏是健康之母吗？} 但我要说，匮乏不仅是身体健康的母亲，也是灵魂健康的母亲。这些事情，保禄，那位真正的医生，也大声宣告，他说：\BibleQuote{只要有吃有穿，就应当知足。}\BibleRef{弟前 6:8}因此，让我们遵从他的命令，这样，在健康的状态下，我们便能履行我们应做的工作，借着我们的主基督耶稣，愿光荣、权能和尊崇归于圣父和圣神，自今至永远，世世无尽。阿们。\par

\SourceRecord{Translated by J. Walker, J. Sheppard and H. Browne, and revised by George B. Stevens.；Nicene and Post-Nicene Fathers, First Series；Vol. 11.；Edited by Philip Schaff.；Buffalo, NY: Christian Literature Publishing Co.,；1889.；<http://www.newadvent.org/fathers/210116.htm>.}

% structure: p0001 source=h1 target=section reason=page-title

\section{论宗徒大事录第十七讲道}

\BibleRef{宗 7:35}\par

\BibleQuote{这梅瑟，就是他们曾拒绝说\InnerQuote{谁立了你做我们的首领和审判者}的，天主却藉着在荆棘丛中向他显现的天使的手，派遣他做首领和解救者。}\par

这话非常适合当前的主题。他说：\Quote{这梅瑟}。\Quote{这}就是那曾面临丧命危险的人；就是那曾被他们藐视的人；\Quote{这}就是他们曾拒绝的人：\Quote{这}同一位，天主兴起他，派遣到他们那里去。\BibleQuote{他们曾拒绝说\InnerQuote{谁立了你做首领？}}正如他们自己（听众）说的：\BibleQuote{我们没有君王，只有凯撒。}\BibleRef{若 19:15}他也在此指出，当时所行的事，是由基督完成的。\BibleQuote{天主藉着天使的手派遣了同一位。}这位天使曾对他说：\BibleQuote{我是亚巴郎的天主。}他说：\Quote{这}同一位梅瑟——注意他如何指出他的名声——\Quote{这}同一位梅瑟，他说：\BibleQuote{领他们出来，先在埃及地，在红海，在旷野四十年间，行了奇事和神迹。这就是那梅瑟，曾对以色列子民说：\InnerQuote{主你的天主要从你们弟兄中给你们兴起一位先知，像我一样。}}（36-37节）如同我一样被藐视。同样，黑落德也想杀害他，他在埃及保全了性命，正如前一位一样，甚至当他还是婴孩时，就有人图谋毁灭他。\BibleQuote{这就是那在旷野的教会中，与那在西乃山上对他说话的天使，以及我们的祖先同在的；他领受了活泼的圣言传给我们。}\BibleRef{宗 7:38}再次没有提及圣殿，也没有提及祭献。经上说：\Quote{与天使一起}，\Quote{他领受了活泼的圣言，要传给祖先。}这表明他不仅行了奇迹，也颁布了法律，如同基督所做的一样。正如基督先施行奇迹，然后立法；梅瑟也是如此。但他们不听从他，即使在奇迹之后仍保持悖逆：他说：\Quote{那些人}，\Quote{我们的祖先不肯听从。}\BibleRef{宗 7:39}在那些四十年间所行的奇事之后。不仅如此，反而恰恰相反：\BibleQuote{反而推开他，心里转向埃及，对亚郎说：\InnerQuote{给我们造神像在我们前面引路，因为领我们从埃及地上来的这个梅瑟，我们不知道他遭遇了什么事。}那时，他们造了一个牛犊，向那偶像献祭，并因自己手所做的而欢庆。于是天主转身，任凭他们事奉天上的万象，正如先知书上所记载的：\InnerQuote{以色列家啊，你们在旷野四十年间，曾给我献过牺牲和祭品吗？你们抬着摩洛的帐幕，和你们的神朗番的星，就是你们所造为朝拜他们的像，因此我要把你们迁徙到巴比伦以外去。}}（40-43节）\Quote{任凭他们}这个说法，意思是\Quote{祂容许了}。\BibleQuote{我们的祖先在旷野里有见证的帐幕，正如那吩咐梅瑟的话，叫他按所看见的样式去制造。}\BibleRef{宗 7:44}即使有会幕，仍没有祭献。\BibleQuote{你们曾给我奉献过宰杀的牲畜和祭品吗？}\BibleRef{亚 5:25}虽有\Quote{见证的会幕}，却对他们毫无益处，他们反而被消灭了。但无论前后，奇迹都未能使他们获益。\Quote{这也正是我们后来的祖先带进来的。}你看见了吗？圣所就在天主所在的地方。正为此，他也说了\Quote{在旷野}，以比较不同的地方。然后是（赐予他们的）恩惠：我们的祖先后来在若苏厄的带领下，将它带进外邦人的产业中，就是天主从我们祖先面前驱逐出去的那些外邦人之地，直到达味的日子。达味在天主面前蒙恩，渴望为雅各伯的天主寻找一个居所。（45-46节）达味\Quote{渴望蒙恩}，却没有建造——他，那奇妙的、伟大的——而是那被弃绝的撒罗满建造了。经上说：\Quote{但是撒罗满}，\Quote{为祂建造了殿宇。然而至高者并不住人手所造的。}\BibleRef{宗 7:47}这一点其实已经由前面所说的表明了，但也由一位先知的声音表明：\BibleQuote{你们要为我建造什么样的殿宇？主天主说。正如先知所说：\InnerQuote{天是我的宝座，地是我的脚凳：你们要为我建造什么样的殿宇？主说：或是我安息的地方在哪里？这一切不都是我手所造的吗？}}\BibleRef{依 66:1}\par

不要惊奇，他说，如果那些蒙基督施恩的人拒绝祂的王国，因为梅瑟的情况也是如此。（回顾）。\Quote{祂将他们领出来}，并且不仅仅是一般地拯救他们，而是当他们在旷野时也如此。\Quote{奇事与征兆}，等等。\BibleRef{宗 7:35} 你注意到他们（斯德望的听众）也与那些古老的奇迹有关吗？\Quote{这就是那位梅瑟}，\BibleRef{宗 7:37} 他，与天主交谈的那位；他，从如此奇妙而奇异的情形中被救出的那位；他，行了如此伟大奇事、拥有如此大能的那位。[\Quote{他曾对以色列子民说：有一位先知}，等等。] 他表明，预言必须完全应验，并且梅瑟并不反对祂。\Quote{这就是那位在旷野的教会中，并对以色列子民说话的那一位。}\BibleRef{宗 7:38} 你注意到根源由此而来，并且\Quote{救恩出自犹太人}吗？\BibleRef{若 4:22} \Quote{与天使一起}，它说，\Quote{那位曾对他说话的天使。}\BibleRef{罗 11:16} 看，他再次断言是祂（基督）赐予了法律，既然梅瑟在旷野的教会中与\Quote{祂}同在。在这里他提醒他们一件伟大的奇事，就是在山上所发生的事：\Quote{谁接受了活生生的神谕以赐给我们。}每当谈及梅瑟，他都令人赞叹，在需要立法时也是如此。\Quote{活生生的神谕}（\OriginalTerm{活生生的神谕}{\GreekText{λόγια}}）这个词是什么意思？就是那些其结局\OriginalTerm{通过言语}{\GreekText{διὰ} \GreekText{λόγων}}所显明的；换句话说，他指的是预言。紧接着是控告，首先是对祖宗的控告，在接受了\Quote{活生生的神谕}之后，\Quote{我们祖先不愿听从。}\BibleRef{宗 7:39} 但关于那些，厄则克耳说它们不是\Quote{活的}；正如他所说：\Quote{我也给了你们不美好的法律。}\BibleRef{则 20:25} 他指的是那些，他说：\Quote{活的。但他们却推开他，心中转回埃及}——那地方他们曾呻吟，曾哭喊，曾呼求天主。\Quote{并对亚郎说：请给我们造神像，在我们前面引路。}\BibleRef{宗 7:40} 何等的愚昧！\Quote{造，}他们说，\Quote{使它们在我们前面引路。}往哪里？\Quote{往埃及。}看他们多么难以摆脱埃及的习俗！你说什么？什么，不等那领你们出来的人，反而逃避恩惠，否认恩主？留意他们是多么侮辱人：\Quote{至于这位梅瑟，}他们说：——\Quote{曾领我们出埃及地}，却绝口不提天主的名；相反，他们将一切都归于梅瑟。他们应当感谢天主的地方，却将梅瑟推到前面；在应当遵照法律行事的地方，他们不再顾及梅瑟。\Quote{我们不知道他遭遇了什么事。}而他曾告诉他们他将上去接受法律，他们竟没有耐心等候四十天。\Quote{请给我们造神像}——他们没有说\Quote{一位神}。——然而人们对此大可惊奇，他们竟然不知道。——\Quote{那时他们造了一个牛犊，向偶像献祭，并因自己双手的作为而欢庆}\BibleRef{宗 7:41}：为此他们本该掩面。既然你们不认识梅瑟，也不认识借着如此奇事显现的天主，那么你们不认识基督又何足为奇？但他们不仅不认识祂，还以另一种方式侮辱祂，即造偶像。\Quote{于是天主转身，任凭他们崇拜天象}\BibleRef{宗 7:42} 因此这些\Quote{习俗}起源于此，祭献也起源于此：是他们自己首先向偶像献祭！正因为如此，经上记载：\Quote{他们在曷勒布造了牛犊，并向偶像献祭：}可见，在这之前从未提及祭献之名，只有活生生的法规和\Quote{活生生的神谕。并欢庆}——这就是节庆的原因。\BibleRef{出 32:5} \Quote{正如先知书上所记载的}——且看他引用经文并非无目的，而是借此表明无需祭献；他说：\Quote{你们曾向我献过宰杀的牲畜和祭品吗？}——他强调这个词（向我？）。你不能说，是因为向我献祭，你们才转而向他们献祭：——\Quote{四十年之久}：并且也是\Quote{在旷野}，在那里祂最显著地显明自己是他们的保护者。\Quote{是的，你们抬起了摩洛赫的帐幕，和你们神蓝潘的星，就是你们所造为要敬拜它们的像。}祭献的原因！\Quote{我要把你们掳到巴比伦以外。}\BibleRef{宗 7:43} 连被掳掠也是对他们恶行的控诉！\Quote{但有一个帐幕，}你们说，\Quote{就是\InnerQuote{见证的帐幕}。}\BibleRef{宗 7:44} （是的，）这就是原因：他们应有天主作见证，仅此而已。\Quote{按照山上所显示给你的样式，}经上说，\Quote{。}所以在山上就有原型。而且这个帐幕，另外，\Quote{在旷野}中，是可移动的，并非固定在一地。他称之为\Quote{见证的帐幕：}即（为见证）奇迹和法令。这就是为什么它和那些（祖先）都没有圣殿。\Quote{正如曾对梅瑟说话的那位所指定的，要他按照他所看见的样式制造。}再者，给予样式的正是祂（基督）。\Quote{直到达味的日子}\BibleRef{宗 7:45}：那时还没有圣殿！然而外邦人也已被驱逐：这就是他提及此事的原因：\Quote{天主驱逐了他们，}他说，\Quote{在我们祖先面前。祂驱逐了他们，}他说：即使那时，也没有圣殿！如此多的奇迹，却没有提及圣殿！所以，虽然起初有一个帐幕，但任何地方都没有圣殿。\Quote{直到达味的日子，}他说：即使达味，也没有圣殿！\Quote{他寻求在天主面前蒙恩}\BibleRef{宗 7:46}：却没有建造——可见圣殿并非什么大事！\Quote{但撒罗满为祂建造了一座殿宇。}\BibleRef{宗 7:47} 他们以为撒罗满伟大，但他并不比他父亲好，甚至不如他父亲，这是明显的。\Quote{然而至高者并不住人手所造的殿宇，正如先知所说：天是我的宝座，地是我的脚凳。}（48-49节。）不仅这些，连那些也不堪与天主相配，因为它们是被造的，即受造物，是祂手的化工。看他如何一步步引导他们，表明连这些也不值得一提。预言又公开说：\Quote{你们要为我建造怎样的殿宇？}等等。\BibleRef{宗 7:50}\par

为什么此时他用\OriginalTerm{谴责的语气}{\GreekText{καταφορικὥς}}说话？他即将赴死时，说话何等大胆；因为我认为他确实知道情况如此。他说道：\Quote{你们这些执拗的人，}和\Quote{心耳未受割损的人。}这也是出自先知，并非他自己的话。\BibleQuote{你们时常反抗圣神；你们的祖先怎样，你们也怎样。}\BibleRef{宗 7:51} 当祂不愿有祭献时，你们却献祭；当祂愿意时，你们又不献祭；当祂不愿给你们诫命时，你们却强求；当你们得到时，又置之不理。再者，当圣殿存在时，你们敬拜偶像；当祂愿意在没有圣殿的情况下受敬拜时，你们却反其道而行。请注意，他说的不是\Quote{你们反抗天主}，而是\Quote{反抗圣神}：他丝毫不认为二者有何区别。还有更甚者：他说：\Quote{你们的祖先怎样，你们也怎样。}基督也曾这样（责备他们），因为他们总是以祖先自夸。他说：\Quote{哪一个先知没有受过你们祖先的迫害？他们杀了那些预言那位正义者来临的人。}他仍称祂为\Quote{那位正义者}，想要阻止他们：\Quote{如今你们竟成了那位的出卖者和凶手}——他向他们提出两项控诉——\Quote{你们接受了天使所安排的律法，却没有遵守。}\BibleRef{宗 7:52} 如何是\Quote{由天使所安排？}有人说（律法）由天使处置，或由那在荆棘丛中向他显现的天使交在他手中；难道祂是人吗？那位成就了那些奇事者，也成就了这些事，这不足为奇。\Quote{你们杀了那些宣讲祂的人}，更甚者，还杀了祂自己。他表明他们不顺服于天主、天使、先知、圣神和一切。正如圣经另一处所说：\Quote{上主，他们杀了你的先知，拆毁了你的祭台。}\BibleRef{列上 19:10} 于是，他们维护律法，说：\Quote{他亵渎了梅瑟。}因此，他表明真正亵渎的是他们，而且他们（的亵渎）不仅反对梅瑟，更是反对天主；表明\Quote{他们}从一开始就这样做；表明\Quote{他们}自己毁坏了他们的\Quote{惯例}，因此无需这些；表明在他们指控他，说他反对梅瑟时，他们自己却反对圣神；而且不仅是反对，还加上杀害；并且他们从起初就一直怀有敌意。你看见了吗？他表明他们既反对梅瑟和所有其他人，又不遵守律法。然而梅瑟曾说过：\Quote{上主必给你们兴起一位先知。}其他先知也预言了这位（基督）要来；先知又说：\Quote{你们要为我建造什么殿宇？}又说：\Quote{你们曾向我献过牺牲和祭献，那四十年？}\BibleRef{申 18:18}\par

如此就是背负十字架之人的放胆直言。那么，我们也当效法这胆量：虽非战争之时，却常是放胆直言之时。因为有人曾说：\Quote{我在君王前论及祢的证言，并不羞愧。} \BibleRef{咏 118:46} 倘若我们偶然置身于外邦人中，要如此堵住他们的口：不动怒，不严厉。（参阅：论格林多前书讲道集第四篇第六节；第三十三篇第四、五节；哥罗森书讲道集第十一篇第二节。）因为若我们带着怒气行事，那就不再像是（因自己的事业而自信之人的）胆量，而是激情；但若以温和待之，这才是真正的胆量。因为成功与失败不可能同时出现在同一件事上。胆量是成功，愤怒是失败。因此，我们若要拥有胆量，就必须远离愤怒，免得有人将我们的话归咎于愤怒。无论你的话多么公义，若你带着怒气说话，你就毁了一切：无论你多么放胆直言，多么公平地责备，或其他。看这人，当他向他们讲话时，多么没有激情！因为他没有辱骂他们：他只是提醒他们先知的话。因为，为显示那不是愤怒，就在他正遭受他们恶待的那一刻，他祈祷说：\Quote{不要将这罪归于他们。} 他绝非在愤怒中说出这些话；不，他是出于对他们的哀痛和悲伤而说的。的确，这就是为何它提到他的面容，说\Quote{他们看见他的面容好像天使的面容，}正是为使他们相信。那么，让我们远离愤怒。愤怒之处，圣神不住；愤怒者受诅咒。凡有愤怒发作之处，任何有益的事物都无法靠近。因为如同在海上风暴中，喧嚣大起，喧哗震耳，那时便不是学习智慧（\OriginalTerm{智慧思考}{\GreekText{φιλοσοφεῖν}}）的时机；愤怒之时亦然。如果灵魂要处于能够言说或接受任何智慧训练的状态，它必须先到达港湾。难道你们没有看见，当我们希望谈论重大事务时，我们寻找安静之地，那里一切都寂静、平静，以免我们被扰乱和不安吗？但若外来的噪音能扰乱我们，那么内在的扰扰更大。无论谁祈祷，若\Quote{带着愤怒和争论：}祈祷，都是徒劳；\BibleRef{弟前 2:8} 无论谁说话，只会使自己可笑；无论谁保持沉默，同样甚至会更糟；无论谁吃，即使在那时也会受伤害；无论谁喝，或无论谁不喝；无论谁坐，或站，或行走；无论谁睡觉：因为即使在睡梦中，这样的幻想也纠缠着他们。因为这样的人身上有什么不是令人不快的呢？眼睛难看，口歪眼斜，四肢扭曲肿胀，舌头污秽不饶人，心神错乱，姿态不雅：令人十分厌恶。看看附魔者的眼睛，以及醉汉和疯人的眼睛；他们彼此有何不同？难道不全是疯狂吗？虽然只是暂时的，那又如何？疯人也是暂时被附身：但有什么比这更糟糕呢？而他们并不以那个借口为耻；\Quote{我不知道（有人说）我说了什么。} 你怎么会不知道这个，你这理性的人，你拥有理性的恩赐，正是为了不让你做出无理性受造物的行为，就像一匹野马，被愤怒和激情冲昏头脑？事实上，这个借口本身就是有罪的。因为你本该知道你说过什么。你说：\Quote{是激情，}你说，\Quote{说了那些话，不是我。} 这怎么可能呢？因为激情本身没有力量，除非它从你那里得到力量。你还不如说：\Quote{是我的手造成了伤口，不是我。} 你以为，什么场合最需要愤怒？难道你不会说：战争与战斗？但即便如此，若是带着怒气行事，一切都被破坏和毁掉。因为所有人中，那些战斗的人最好不愤怒：所有人中，那些想要伤害他人（\OriginalTerm{伤害他人的人}{\GreekText{τοὺς} \GreekText{ὑβρίζοντας}}）的人最好不愤怒。那么又如何可能战斗呢？你会问。带着理性，带着自制（\OriginalTerm{温和自制}{\GreekText{ἐπιεικεία}}）。\par

\SourceRecord{Translated by J. Walker, J. Sheppard and H. Browne, and revised by George B. Stevens.；Nicene and Post-Nicene Fathers, First Series；Vol. 11.；Edited by Philip Schaff.；Buffalo, NY: Christian Literature Publishing Co.,；1889.；<http://www.newadvent.org/fathers/210117.htm>.}

% structure: p0001 source=h1 target=section reason=page-title

\section{宗徒大事录第十八篇讲道}

\BibleRef{宗 7:54}\par

他们听见这些话，就\BibleQuote{刺入心肠，并向他咬牙切齿}。\par

看，作恶者又一次陷入困境。正如犹太人困惑地说：\BibleQuote{我们该怎样对待这些人？}这些人也同样\BibleQuote{刺入心肠}。\BibleRef{宗 4:16}然而，他本有充分的理由愤怒，因为他没有行任何不义，却被当作罪犯对待，且被恶意诽谤。但诽谤者最终却自食其果。我常重复的那句谚语果真如此真实：\Quote{行恶必遭恶报。}然而他（在控告他们时）并未诉诸诽谤，而是证明了他所说的。我们确知，当我们因良心无愧的事而蒙受羞辱时，我们并不会因此受损。——\Quote{如果他们想要}你说，\Quote{杀他，他们怎么没有利用他所说的话作为借口来杀他呢？}他们想为自己的暴行找一个貌似正当的借口。\Quote{那么，羞辱他们难道不是正当的借口吗？}如果他们感到被羞辱，那并不是他的错，而是先知对他们的控告。此外，他们不愿让人以为他们是因为他所说的反对他们的话而杀他——正如他们在基督的事上所作的那样；不，而是因为不敬：而现在他所说的这番话正是虔敬的表现。因此，当他们试图在杀害他之外也损害他的名声时，\BibleQuote{他们被刺入心肠}。因为他们害怕，唯恐他反而成为更受尊敬的对象。因此，他们在基督的事上所作的，在这里也照样行。因为正如祂说：\BibleQuote{你们将要看见人子坐在天主的右边}\BibleRef{玛 26:64}，他们却称之为亵渎，\Quote{冲向他}；这里也是如此。在那里，他们\BibleQuote{撕裂衣服}；在这里，他们\BibleQuote{捂着耳朵}。但斯德望充满圣神，定睛望天，看见天主的荣耀，并耶稣站在天主的右边，就说：\BibleQuote{看，我看见天开了，人子站在天主的右边。他们就大声喊叫，捂着耳朵，齐心冲向他，把他推出城外，并拿石头砸死他。}\BibleRef{宗 7:55}然而，如果他说谎，他们本应认为他是疯癫，而放他走。——但他想要引导他们，\Quote{就说：看，}等等，因为他既然已经谈到基督的死亡，而未提及祂的复活，就愿意也加上这个道理。\BibleQuote{站在天主的右边。}祂以这种方式向他显现，为的是，即使如此，犹太人也能接受祂：因为由于（祂）坐在（天主右边）的观念令他们反感，所以暂时他只提出与祂复活有关的事。这也是为什么他的脸被显荣的原因。因为慈悲的天主，愿意利用他们的阴谋作为召叫他们归向自己的工具。看，行了许多奇迹！\BibleQuote{他们把他推出城外，并拿石头砸死他。}此处再次，\BibleQuote{城外}，甚至在死亡中，也是宣认与宣讲。\BibleRef{希 13:21}\BibleQuote{证人将衣服放在一个名叫扫禄的少年人脚前。他们拿石头砸斯德望，他呼求天主说：主耶稣，接收我的灵魂。}\BibleRef{宗 7:59}这是为了向他们表明他并未灭亡，并教训他们。\BibleQuote{他跪下，大声呼喊说：主，不要将这罪归到他们身上。}\BibleRef{宗 7:60}为了替自己辩白，并表明他先前的话也不是出于情绪，他说：\BibleQuote{主}\BibleQuote{不要将这罪归到他们身上}：也愿意以这种方式争取他们。因为要表明他宽恕了他们杀害他的愤怒和狂怒，并且他自己的灵魂脱离了一切情绪，这样他的话语才能被接纳。\par

\BibleQuote{扫禄赞同杀死他。}由此兴起一场迫害，且成为一场大迫害。\BibleQuote{那时，耶路撒冷的教会遭受大迫害；除了宗徒以外，众人都分散到犹太和撒玛黎雅各处。}\BibleRef{宗 8:1}请注意天主又一次允许试探兴起；注意并好好观察，这些事件如何由天主眷顾所安排。他们因奇迹而受钦佩；被鞭打后，他们并未因此受损；有人被派去管理寡妇的事；圣言增长；天主又一次允许一个巨大的阻碍兴起。一场非同寻常的迫害[\BibleQuote{众人都分散了}等等]；因为他们害怕他们的敌人，如今敌人变得更加大胆；同时，这也证明那些害怕、逃跑的也不过是人。因为，免得你事后说他们单靠恩宠成就事工，他们也曾受迫害，自己也变得更加恐惧，而他们的对手则更加大胆。\BibleQuote{众人都分散了}经上说，\BibleQuote{除了宗徒。}但这乃是天主的安排，好叫他们不再都留在耶路撒冷。\BibleQuote{有虔诚的人}经上说，\BibleQuote{把斯德望埋葬了，并为他大大哀悼。}\BibleRef{宗 8:2}如果他们是\BibleQuote{虔诚的}，为何\BibleQuote{为他大大哀悼？}他们尚未成全。这人温良可爱：这也表明他们是人——不仅是他们的恐惧，还有他们的悲伤和哀悼。谁不会流泪看到那位温良如羔羊的人被石头砸死，躺卧在地？圣史记录了他墓前相称的悼词，正在这一句话中：\BibleQuote{不要将这罪归到他们身上。}——\BibleQuote{又}\BibleQuote{为他大大哀悼。}——但让我们再回顾一下以上所说。\par

他（天使）显现的原因\EditorialNote{重述，7:54；8:2}；\Quote{但斯德望充满圣神，注目向天，看见天主的光荣，并看见耶稣站在天主右边。}当他说：\Quote{我看见天开了，他们就掩耳，齐心拥上前去。}\BibleRef{宗 7:56-57}这些事有何可责之处？\Quote{他们拥上前去，}向着这位行了如此奇事、言词超越众人、并能发出如此宣讲的人！好像他们得到了自己所渴望的，立刻发泄他们的怒气。\Quote{证人，}他说，\Quote{把衣服放在一个名叫扫禄的少年人脚前。}\BibleRef{宗 7:58}请注意他如何详细地提及关于保禄的事，为叫你看出那在他内运行的能力是出于天主。但经历过这一切之后，他不但不信，还千方百计地攻击祂：因此圣经说：\Quote{扫禄也赞同杀死他。}——这位有福之人不仅祈祷，而且恳切地祈祷：\Quote{跪下。}看他神圣的去世！主只许他的灵魂存留到那时。\Quote{说完这话，就长眠了。}\BibleRef{宗 7:60}——他们都散往犹太和撒玛黎雅各处。\BibleRef{宗 8:1}如今他们毫无顾虑地与撒玛黎雅人来往，虽然曾对他们说：\Quote{不要外邦人的路，}\Quote{也不要进入撒玛黎雅人的城。}\BibleRef{玛 10:5}\Quote{除了宗徒们，}经上说：他们这样行，一方面想赢得犹太人——但不愿离开耶路撒冷——另一方面借此激励别人勇敢。\par

\Quote{扫禄却摧残教会，进入各家，拉着男女，把他们投入监里。}\BibleRef{宗 8:3}他的疯狂极大：独自一人，甚至进入各家：因为他已准备为法律舍命。\Quote{拉着，}经上说：\Quote{男女：}请注意他的胆量、暴力和疯狂。凡落在他手里的，他都施以各种虐待：因为刚发生的谋杀使他更加胆大妄为。\Quote{因此那些散开的人，走遍各处，宣讲真道。斐理伯下到撒玛黎雅城，向他们宣讲基督。群众同心合意地听从斐理伯所说的话，听见并看见他所行的奇迹。因为有许多附了邪魔的人，大声喊叫着出来了；许多瘫痪和瘸子都得了痊愈。在那城里大有喜乐。}\BibleRef{宗 8:4}请看另一考验：西满的事件。\Quote{他自称，}经上说：\Quote{他是某位大人物。众人都从最小到最大的都听从他说：这人就是天主的大能。他们听从了他，因为很久以来他常用邪术迷惑他们。但当他们相信了斐理伯所宣讲的天主国和耶稣基督的名时，男女都受了洗。西满自己也相信了，受洗后常与斐理伯在一起，看见所行的奇迹和大能，就很惊奇。耶路撒冷的宗徒们听说撒玛黎雅领受了天主的话，就派伯多禄和若望到他们那里去。他们二人下来，就为他们祈祷，使他们领受圣神。}\BibleRef{宗 8:10}\Quote{因为那时圣神还没有降在任何人身上；他们只因主耶稣的名受了洗。于是他们给覆手，他们就领受了圣神。}为指明这是实情，且他们尚未领受行奇迹的圣神，请看结果：西满看见这事，就来请求这个能力。\Quote{西满看见因宗徒覆手，便赐下圣神，就给他们银钱说：也赐给我这个权能，使我无论给谁覆手，谁就领受圣神。}\par

\Quote{迫害，}你们说，\Quote{加重了。}诚然，但就在那时，它给从前（被敌对势力）附身的人带来了释放。因为它将奇迹像堡垒一样，栽在敌人的国境中心。——连斯德望的死都没有扑灭他们的怒气，反而使其更盛：它广泛地分散了教师们，以至于门徒队伍更大了。——\InnerQuote{并且有喜乐。}然而以前曾有\InnerQuote{极大的悲哀}：诚然；但再次注意这善——\InnerQuote{很久以来}这病就存在，但这个人带给了他们释放。——他怎么也给西满施洗呢？正如基督拣选了犹达斯一样。——\InnerQuote{看见这些征兆}因他所作的，因为其他的人没有领受（行奇迹的）能力，他不敢问这事。——那么，为什么他们没有把他打死，像对待阿纳尼雅和撒斐辣那样呢？因为在古时候，那（在安息日）拾柴的人也被处死，作为他人的警戒\BibleRef{户 15:32}。在其他情况下，没有任何人遭受同样的命运。同样在此时，\BibleQuote{伯多禄对他说：愿你的银子跟你一同丧亡，因为你想天主的恩赐可以用银钱买得。}——\BibleRef{宗 8:20}为什么这些人受洗时没有领受圣神呢？或者因为斐理伯将这尊荣留给了宗徒们；或者因为他没有这（分施的）恩赐；或者他是七位之一：这更可能是所说的。因此，我想这位斐理伯是宗徒之一。但注意：那些人没有出去；上智安排这些出去而那些人缺乏，因为圣神：因为他们领受了行奇迹的能力，却没有也分施圣神给他人：这是宗徒的特权。并且注意（他们怎么派遣）主要的人：不是其他任何人，而是伯多禄和若望。\BibleQuote{当西满，经上说，看见藉着宗徒们覆手，圣神被赐予后，}他不会说\InnerQuote{看见了，}除非有某种可感觉的显明。\BibleQuote{于是他们给他们覆手}等，正如保禄所做的，当时他们说方言。\BibleRef{宗 19:6}注意西满的可憎行为。\Quote{他献上钱，}用意何在？但他没有见伯多禄为钱而做这事。他这样行，不是出于无知；而是因为他要试探他们，因为他想得到控告他们的把柄。因此伯多禄也说：\BibleQuote{在这事上你没有分，也没有权，因为在天主面前你的心不正，因为你思想了等等。}\BibleRef{宗 8:21}他再次揭露他心里所想，因为西满以为能逃避察觉。\BibleQuote{所以你当悔改你这罪恶，并祈求天主，或者你心中的意念可得赦免。因为我看出你是在不义的枷锁中。西满回答说：请你们为我祈求主，使你们所说的没有一样临到我身上。}\BibleRef{宗 8:22}连这事他也只是形式上做的，像例行公事的话，而他本当像一个忏悔者哭泣哀恸。\InnerQuote{或许可得赦免。}并非表示他若哭泣就不得赦免，但这也是先知的方式，就是绝对地宣判，（\OriginalTerm{禁止}{\GreekText{ἀπαγορεύειν}}），而不说：\Quote{然而，你若这样行，你的罪必得赦免，}而是说刑罚无论如何必要实现。\par

（a）\Quote{因此，那些逃散的人四处奔走，宣讲着圣言。}但我愿你们钦佩，即使在患难之时，他们也未曾忽略宣讲。\Quote{听见并看见他所行的奇迹。}\BibleRef{宗 8:4}正如在梅瑟的事上（与术士相比）奇迹是明显的奇迹，这里也是如此。那里有魔法，所以这些征兆是显明的。（b）\Quote{因为有许多附魔的人，邪灵从他们身上出来了}\BibleRef{宗 8:7}；这是一个明显的奇迹——不像术士所行的：因为另一人（西满）很可能用符咒捆绑人——\Quote{还有许多，}经上说，\Quote{瘫痪的和瘸腿的都被治愈了。}这里没有欺诈：因为他们只需行走和工作。\Quote{众人都注意他，说：\InnerQuote{这人就是天主的大能}}\BibleRef{宗 8:10}这应验了基督所说的话：\Quote{将有假基督和假先知冒我的名而来。}——\BibleRef{玛 24:24}\Quote{他们注意他，因为他长久以来用邪术迷惑了他们。}\BibleRef{宗 8:11}（a）然而那里应该没有一个附魔者，因为他长久以来用邪术迷惑了他们：但如果有许多附魔者，许多瘫痪者，这些伪装就不是真理。但斐理伯在此也藉着言语赢得了他们，谈论关于基督的国。\BibleRef{宗 8:12}\Quote{西满，}经上说，他受了洗，常与斐理伯在一起\BibleRef{宗 8:13}：不是为信德的缘故，而是为能成为像他一样的人。（b）但他们为何不立即纠正他？他们满足于他自责。因为这也属于他们教导的工作（\OriginalTerm{}{\GreekText{τἥς} \GreekText{διδασκαλίας}}）。但当他无力抵抗时，他就假装，正如术士们所说，\Quote{这是天主的手指。}而确实，为免他被再次赶走，因此他\Quote{常与斐理伯在一起，}并不离开他。\Quote{那时，在耶路撒冷的宗徒，}等等。（13、14节）请看，通过斯德望的死，天主的眷顾成就了多少事！（a）\Quote{但他们，}经上说，\Quote{下来后，为他们祈祷，使他们领受圣神：因为圣神还没有降在他们任何人身上。于是他们给他们覆手，他们就领受了圣神。}\BibleRef{宗 8:15}你们看，这并非以平常方式成就，而是需要大能才能赐予圣神？因为获得罪的赦免与领受如此大能不是一回事。（b）他们渐渐领受这恩赐。这是一个双重的记号：既赐予那些人，又不赐予这人。然而，这人本该相反地祈求领受圣神，却因不关心此事，而求赐给别人的能力。而那些人并没有领受这赐予的能力：但这人竟想比斐理伯更显赫，而斐理伯是在门徒中！（a）\Quote{他拿钱给他们。}（18、19节）什么？他看见别人这样做吗？他看见斐理伯吗？他以为他们不知道他来时的念头吗？（b）\Quote{你的银子与你一同丧亡吧！}\BibleRef{宗 8:20}：因为你没有善用它。这不是咒诅的话，而是惩罚。\Quote{归你，}他说，归你吧：既然你如此。就如有人说，让你的银子与你的意图一同灭亡。你竟对天主的恩赐有如此卑劣的想法，以为它完全是人的事？并非如此。（a）因此伯多禄也称这事为恩赐：\Quote{你曾以为天主的恩赐可以用金钱买得。}你们看，他们如何时时远离金钱？\Quote{因为你的心在天主面前不正。}\BibleRef{宗 8:21}你们看，他如何出于恶意行事？然而，朴实是必需的。（b）因为若他出于朴实而行，他甚至会欢迎他的情愿之心。你们看，对大事有卑劣的想法是双重的罪？因此，他命令他两件事：\Quote{你当悔改，并祈祷，或许你心中的意念可得赦免。}\BibleRef{宗 8:22}你们看，他存有恶念？因此他说，\Quote{或许可得赦免：}因为他知道他是不可救药的。（a）\Quote{因为我看出你正处在苦胆的苦毒和罪恶的束缚中。}\BibleRef{宗 8:23}这是极其愤怒的话！但他没有以别样惩罚他：为使以后信德不至于被迫；使事情不至于显得残忍；使他能引入悔改的主题：或者，因为责备他、告诉他心里的事、使他承认自己被战胜\OriginalTerm{}{\GreekText{ὅ}}\OriginalTerm{}{\GreekText{τι} \GreekText{ἐ}\&}已足以改正。\par

\SourceRecord{Translated by J. Walker, J. Sheppard and H. Browne, and revised by George B. Stevens.；Nicene and Post-Nicene Fathers, First Series；Vol. 11.；Edited by Philip Schaff.；Buffalo, NY: Christian Literature Publishing Co.,；1889.；<http://www.newadvent.org/fathers/210118.htm>.}

% structure: p0001 source=h1 target=section reason=page-title

\section{论宗徒大事录第十九篇讲道}

\BibleRef{宗 8:26-27}\par

\BibleQuote{主的天使向斐理伯说：\Quote{起来，往南走，沿着从耶路撒冷下到迦萨的路，那路是旷野。}他就起来去了。}\par

依我看来，这位（斐理伯）是七位中的一位；因为从耶路撒冷他不会向南走，而会向北；但从撒玛黎雅则是\BibleQuote{往南。那路是旷野}：因此不必担心犹太人的攻击。他没有问\InnerQuote{为什么？}而是\BibleQuote{起来去了}。看，经上说：\BibleQuote{有一个厄提约丕雅人，是在厄提约丕雅女王甘达刻手下有权势的太监，总管她全部财宝，他曾上耶路撒冷去朝拜，正在回去，坐在车上，诵读依撒意亚先知。}（27-28节）这人对他的赞美很高：他住在厄提约丕雅，被这么多事务缠身，又没有节日，住在那个迷信的城市，却来到耶路撒冷朝拜。他也很勤学，甚至在车上坐着诵读。经上说：\BibleQuote{圣神对斐理伯说：\InnerQuote{你上前去，靠近这辆车。}斐理伯就跑过去，听见他诵读依撒意亚先知，便说：\InnerQuote{你所诵读的，你明白吗？}他说：\InnerQuote{若没有人指教我，怎能明白呢？}}\BibleRef{宗 8:29}再看他的虔诚：虽然他不明白，仍然诵读，诵读后加以查考。\BibleQuote{他请斐理伯上车与他同坐。他所诵读的经文是：祂如同羊被牵到宰杀之地，又像羔羊在剪毛者前缄默，同样祂不开口。在谦卑中，祂的审判被夺去，谁能述说祂的世代？因为祂的生命从地上被夺去。太监回答斐理伯说：\InnerQuote{请问，先知这话是指着谁说的？是指自己，还是指别人？}斐理伯就开口，从这段经文开始，向他宣讲耶稣。}\BibleRef{宗 8:32}注意这是如何天意安排的。先诵读而不明白；然后诵读的正是包含受难、复活和恩赐的经文。\BibleQuote{他们沿路前行，来到有水的地方，太监说：\InnerQuote{看，这里有水，有什么阻止我受洗呢？}}\BibleRef{宗 8:36}注意他的急切渴望，注意他的确切知识。\BibleQuote{他吩咐车停住，斐理伯和太监两人都下到水中，斐理伯给他施了洗。他们从水中上来时，主的神把斐理伯提去，太监不再看见他，就欢欢喜喜继续前行。}（38-39节）但为何主的神把他提去？这使所发生的事显得更加奇妙。当时太监并不认识他。因此这样做，是为了以后斐理伯成为他惊奇的对象。经上说：\BibleQuote{他欢欢喜喜继续前行。但斐理伯却在阿左托被发现；他经过各城，宣讲福音，直到凯撒勒雅。}\BibleRef{宗 8:40}这位（斐理伯）就是七位中的一位；因为事实上之后他在凯撒勒雅被发现。因此，圣神提去斐理伯是好的、适宜的；否则太监愿意与他同去，斐理伯因拒绝他的请求而使他伤心，因为时候尚未到来。（\Emph{a}）但同时这对他们也是一个鼓励的保证：他们也将胜过外邦人：因为（最初）信徒的高尚品格（\OriginalTerm{可信性}{\GreekText{τὸ} \GreekText{ἀξιόπιστον}}）足以感动他们。如果太监留在那里，又能找到什么过错？\EditorialNote{但他不认识他}：因为这就是为什么经上说\BibleQuote{他欢欢喜喜继续前行}；这样，如果他认识他，就不会如此高兴了。\par

\Quote{上主的天使，}等等。\BibleRef{宗 8:26}（\Emph{b}）请看天使辅助宣讲，而非亲自宣讲，而是召叫这些人（去工作）。但此事的奇妙性也由此显示：古时罕见、难以成就的事，如今轻易发生，且看何其频繁！（\Emph{c}）\Quote{一个太监，}经上说，\Quote{一位有权势的人，在甘达刻，厄提约丕雅女王手下。} \BibleRef{宗 8:27} 因为那里古时妇女掌权，这是他们中的法律。斐理伯还不知道他为了谁来到旷野：（\Emph{d}）但在车上，有什么能阻止他准确地了解这一切（详情）呢？\Quote{正在诵读依撒意亚先知。} \BibleRef{宗 8:28} 因为道路是旷野，此事并无炫耀。也请注意时间：在最酷热的时刻。（\Emph{e}）\Quote{圣神向他说。} \BibleRef{宗 8:29} 现在不是天使，而是圣神催促他。为何如此？那时，异象通过天使以更粗浅的形式发生，因为这适合那些更属肉体的人，但圣神是为更属神的人。祂如何向他说？当然，是默启他。为何天使不向那一位显现，带他到斐理伯那里？因为他可能不会被说服，反而会害怕。请看斐理伯的智慧：他没有指责他，也没有说\Quote{我对这些知道得很清楚}，没有奉承他，说\Quote{你诵读真是有福}。但注意他的言语，既无严厉也无谄媚，而是一个善良友好之人的言语。\Quote{你明白你所诵读的吗？} \BibleRef{宗 8:30} 因为需要他自己发问，自己渴望。他明确暗示他知道对方一无所知，并说\Quote{你明白你所诵读的吗？}同时向他表明其中蕴含着极大的宝藏。这件事也表明，这个太监不看人的\OriginalTerm{外貌}{\GreekText{σχἥμα}}，不说\Quote{你是谁？}不责骂，不自高自大，也不说自己知道。相反，他承认自己的无知，因此也得以学习。他向医生展示他的伤处：一眼看出那人既知道这事，又愿意教导。看他多么没有傲慢；外貌并不显赫。他如此渴望学习，留心听他的话；那句\Quote{寻求的，就必找到}的话\BibleRef{玛 7:8}在他身上应验了。\Quote{于是，}经上说，\Quote{他恳请斐理伯上来与他同坐。} \BibleRef{宗 8:31} 你注意到他的热切与渴望吗？但若有人说他应当等候斐理伯（开口），（答案是）他不知道是怎么回事：他丝毫不知对方要对他说什么，只以为将接受某种预言（教训）。而且，这样更显尊敬，因为他没有拉他上自己的车，而是恳请他。\Quote{斐理伯，}我们读到，\Quote{跑向他，听见他诵读；}单是跑这一举动，就表明他想说（什么）。\Quote{而他所诵读的经文，}经上说，\Quote{正是这一段：祂如同一只羊被牵去宰杀。} \BibleRef{宗 8:32} 这一情况也表明他拥有\OriginalTerm{高尚情操}{\GreekText{φιλοσοφίας}}，因为他手中有这位超越所有其他先知的先知。斐理伯没有随意地向他说教，而是安静地：不，在受到询问之前什么也不说。先前他曾祈求，现在也是如此，说\Quote{请问先知说这话是指谁呢？}他能知道先知们以不同的方式谈论不同的人，或他们以第三人称谈论自己——这个问题显出他非常深思熟虑。让我们，无论贫富，都被这个太监羞辱吧。然后，经上说，\Quote{他们来到一处水边，他说，看，这里有水。}\BibleRef{宗 8:36}他又主动请求，说\Quote{什么阻碍我受洗呢？}再看他的谦逊：他不说\Quote{为我施洗}，也不沉默；而是说出一种介于强烈渴望与敬畏之间的言语，说\Quote{什么阻碍我呢？}你注意到他拥有（信仰）完美的教理吗？因为先知确实包含了全部：降生成人、受难、复活、升天、将来的审判。如果他表现出极其热切的渴望，不要惊讶。你们这些尚未受洗的人，都该羞愧。\Quote{于是，}经上说，\Quote{他吩咐车停住。} \BibleRef{宗 8:38} 他说话，同时下令，还未听到（斐理伯的答复）。\Quote{当他们从水中上来时，主的神把斐理伯提去；} \BibleRef{宗 8:39} 为使这事件显示出于天主，使他不会以为这仅仅是出于人。\Quote{他就，}经上说，\Quote{欢欢喜喜地前行。} （页121，注2）这暗示，如果他知道了，他会悲伤：因为他的喜乐如此之大，得到了圣神的眷顾，他甚至不留意眼前的事——\Quote{但斐理伯在阿左托出现。} \BibleRef{宗 8:40} 斐理伯也大有收获：——他从前听说的关于先知的事，关于哈巴谷、关于厄则克耳及其他人的事，如今亲眼看见在自己身上成就。\BibleRef{达 14:36} 由此可见他走了很远的路，因为他\Quote{在阿左托出现}。（圣神）把他放在那里，他便从那里开始宣讲：\Quote{他经过各地，在所有的城中宣讲，直到凯撒勒雅。}\par

\BibleQuote{扫禄仍然向主的门徒口出恐吓和凶杀之气，他去见大司祭，请求给他发信到大马士革的各会堂，如果找到属于这种道路的人，不论男女，都把他捆绑带到耶路撒冷。} \EditorialNote{见：宗9:1-2} 他恰当地提到保禄的热心，并显示就在他热心之际被吸引。\BibleQuote{仍然口出恐吓和凶杀之气，} 而且还不满足于杀害斯德望，还不满足于迫害教会和分散教会。看，基督所说的话应验了：\BibleQuote{凡杀害你们的，还以为是恭敬天主。} \BibleRef{若 16:2} 他这样做，不像犹太人——绝对不像！因为他出于热心，甚至前往外邦城市；而犹太人连耶路撒冷的人也不关心，他们只求尊荣。但为什么他去大马士革？那是一座大城，一座王城，他担心那里先被占领。注意他强烈的渴望和热情，以及他如何严格按照法律行事：他不去找总督，而是去找 \BibleQuote{司祭。他若找到属于这种道路的人，} 因为信徒是这样称呼的，大概因为他们走上了通往天堂的正路。为什么他没有得到授权就地惩罚他们，而是把他们带到耶路撒冷？他在这里做这些事更有权威。且看他将自己置于何等危险之中。他不怕自己受到伤害，但他还带了别人同去，\BibleQuote{他若找到} 经上说，\BibleQuote{任何属于这种道路的人，不论男女}——哦，何等的残忍！——\BibleQuote{就捆绑带回。} 他这次旅行，想向所有人表明他的做法：他们是如此不热心于这件事。注意他之前也把人投入监狱。所以其他人没有成功，但这个人因他热切的头脑而成功。\BibleQuote{当他行路临近大马士革时，忽然从天上有光四面照着他，他就跌倒在地，听见有声音对他说：扫禄，扫禄，你为什么迫害我？} （3-4节）。为什么不在耶路撒冷？为什么不在大马士革？为免不同的人以不同方式叙述此事，唯独他这位为此而去的可信的叙述者，\OriginalTerm{可信的}{\GreekText{ἀξιόπιστος}}，他为此而去。事实上，他既在台阶上的演讲中，也在阿格黎帕面前申辩时说了这事。\BibleQuote{跌倒在地}：\EditorialNote{见：宗22:6；26:12} 因为过强的光容易使人震惊，因为眼睛有其限度；据说过强的声音也会使人耳聋昏迷（如发作），\OriginalTerm{震惊的}{\GreekText{ἀποπλἥγας}}。但只使他眼瞎，并以恐惧熄灭他的激情，使他能听清所说的话。\BibleQuote{扫禄，扫禄，} 他说，\BibleQuote{你为什么迫害我？} 他什么也不告诉他：没有说\Quote{你要相信}，或任何类似的话；而是与他理论，几乎是说：\Quote{你从我的受了什么或大或小的伤害，以致你做这些事？} \BibleQuote{他说：主，你是谁？} \BibleRef{宗 9:5} 这样首先承认自己是他的仆人。\BibleQuote{主说：我是你所迫害的耶稣。} 不要以为你是在同人作战。与他同行的人听见保禄的声音，却看不见回答他的人——因为主允许他们听到不甚重要的事。若他们听到另一个声音，就不会相信；但见保禄在回答某人，他们惊奇。\BibleQuote{起来，进城去，必有人告诉你你该做什么。} \BibleRef{宗 9:6} 注意，他并没有立即把一切都说出来，而是先软化他的心。同样，他第二次召叫门徒。\BibleQuote{必有人告诉你} 等：他给他美好的希望并暗示他将重见光明。\BibleQuote{与他同行的人站着说不出话来，听见声音，却看不见人。扫禄从地上起来，睁开眼却什么也看不见，他们牵着他的手，领他进了大马士革} （7-8节）——魔鬼的战利品（\OriginalTerm{他的财物}{\GreekText{τὰ} \GreekText{σκεύη} \GreekText{αὐτοὕ}}），\BibleQuote{他的财物} \BibleRef{玛 13:29}，如同从某个被攻陷的大城，是的，某个被攻陷的都会。奇妙的是，敌人和仇敌自己当着众人领他进城！\BibleQuote{三天之久，他因眼瞎不吃不喝。} \BibleRef{宗 9:9} 有什么能与此相比？为安慰斯德望事件带来的沮丧，保禄的归化带来了鼓励：虽然那悲伤因斯德望的结局得到安慰，却也得了这一进一步的安慰；此外，撒玛黎雅众村庄的归化也带来了极大的安慰。——但为何不在起初就发生，而是在这些事之后？为要显示基督确实复活了。这个疯狂的攻击者，这个不信基督死而复活的人，迫害他的门徒，若没有复活的巨大能力，怎能使他相信？就算其他宗徒倾向基督，但这个人呢？为何不在复活后立即发生？为要更清楚地显示他的敌意是公开的战争。这个疯狂到流血、把人投入监狱的人，突然就相信了！他从未与基督相处还不够，他还以激烈敌意迫害信徒：他使任何人都无法比他更狂暴。但当他失明时，他看到了基督主权和仁慈的证明：于是他回答：\BibleQuote{主，你要我做什么？} 以免有人说他伪善，这个渴求鲜血、去见司祭、投身如此危险、迫害甚至外地的人——在此情况下他现在承认基督的主权。为什么那光照他不发生在城内，而发生在城外？许多人不相信，因为即使在耶路撒冷，当人们听到天上的声音时，他们说 \BibleQuote{打雷了} \BibleRef{若 12:29}；但这个人足以权威地报告自己亲历的事。他被领进城，虽无捆绑，但他们牵引他，而他原想牵引别人。\BibleQuote{他不吃不喝：} 他自责过去，他认罪、祈祷、恳求天主。若有人说这是强迫的结果，我们回答：同样的事情发生在厄里玛斯身上，为什么他没有改变？（十三章，\WorkTitle{de Laud. Pauli Hom.} 第四篇§1，第二卷491页）。还有什么比复活时的地震、士兵的报告、其他奇迹、看见复活的主更具有强迫性？但这些并不强迫，而是教导，\OriginalTerm{不是强迫性的，而是教导性的}{\GreekText{οὐκ} \GreekText{ἀναγκαστικὰ} \GreekText{ἀλλὰ} \GreekText{διδακτικά}}。为什么犹太人听到这些事却不相信？他说话的真实性显而易见：因为若不是这事发生，他不会改变，所以所有人都该相信。他不逊于宣讲复活的那些人，因他突然的归化而更可信。他没有与任何信徒交往；他是在大马士革归化的，或者更准确地说，是在到达大马士革之前发生的。我问犹太人：保禄是如何归化的？他看见这么多神迹，却没有归化；他的老师加玛里耳归化了，而他仍不归化，\Emph{上文}第87页注1。谁说服了他——不仅说服，还立即激起如此热切的热情？为什么他竟愿为基督下地狱？事实的真相显而易见。\par

但是，正如我说过的，目前让我们（想到太监时）感到羞愧，无论是他的圣洗还是他的阅读。你们是否注意到他处于极大的权威地位，他拥有财富，甚至在旅途中也不允许自己休息？他在家中休闲时该是怎样的，这个在旅行中也不休息的人？他在夜间该是怎样的？你们这些身居高位的人，请听：效法他的谦卑（\Emph{de Lazaro, Conc.} iii. §3, t. i. p. 748. c），他的虔诚。虽然即将回家，他并没有对自己说：\Quote{我要回国，在那里领受圣洗}——这是大多数人所说的那些冷漠之词！他不需要征兆，不需要奇迹：仅从先知就相信了（\Emph{b}）。但为什么（这样安排）他不在上耶路撒冷之前见到（\Person{斐理伯}），而是在去了之后？他不应该在迫害中见到宗徒。因为他仍然软弱，先知并不容易（但先知还是）教导了他。因为即使现在，如果你们中有人愿意钻研先知书，他就不需要奇迹。而且，如果你们愿意，我们来看这段预言本身：\BibleQuote{他如同绵羊被牵去宰杀，又像羔羊在剪毛者前缄默，他不开口。他受屈辱时，审判被夺去；谁能述说他的世代？因为他的生命从地上被夺去。}\BibleRef{依 53:7-8}[现在他得知]很可能他听说他被钉死，\Quote{他的生命从地上被夺去}以及其余的：\Quote{他没有犯罪，口中也没有诡诈}，他也能拯救别人，以及他是谁，他的世代不可言说。很可能他看见了那里的裂石，以及幔子如何裂开，黑暗如何降临等等；\Person{斐理伯}只是从先知书中取题提到了这些事。阅读圣经实在是一件伟大的事！\Person{梅瑟}所说的应验了：\Quote{坐、卧、起、行，要记念上主你的天主。}\BibleRef{申 6:7}因为道路，尤其是孤寂的道路，给我们反思的机会，无人打扰。既有这个人在道路上，也有\Person{保禄}在道路上：然而后者无人牵引，唯独\Person{基督}。这对宗徒来说是一项太大的工作：此外，更甚的是，宗徒们在耶路撒冷，而\Place{大马士革}没有权威人士，他仍然从那处归化回来：但\Place{大马士革}的人知道他并非从耶路撒冷归化而来，因为他带着信件，为要捆绑信徒。如同一位完美医师，当热病达到顶点时，\Person{基督}帮助了他：因为需要在他的狂乱中制服他。因为那时最可能被降服，并谴责自己犯下了可怕的鲁莽行为（\Emph{a}）。为此\Person{保禄}哀叹自己，说：\Quote{然而，我蒙受了怜悯，正是为叫耶稣基督在我这罪魁身上显示他的永恒忍耐。}\BibleRef{弟前 1:13}实在，我们有理由钦佩这个太监。他没有看见\Person{基督}，没有看见奇迹：他看见耶路撒冷\OriginalTerm{仍然屹立}{\GreekText{συνεστὥτα}}，他的灵魂\OriginalTerm{殷切}{\GreekText{μεμεριμνημένη}}。他相信了\Person{斐理伯}。他为何如此行事？他的灵魂是殷切的。然而，那个（在十字架上的）强盗见过奇迹：贤士们见过星辰：但此人什么也没有。如此，认真阅读圣经是多么伟大的事！那么\Person{保禄}呢？难道他没有研究法律吗？但在我看来，他是特别被保留的，因为那目的我已经预先提及了，因为\Person{基督}愿意从各方面吸引犹太人归向自己。因为如果他们神志清醒，没有什么比这更对他们有益；因为这事比奇迹和一切更能吸引他们：而另一方面，没有什么比这事更容易使迟钝的人跌倒。看，在宗徒之后，我们也有天主施行奇迹。他们控告宗徒（在这些奇迹之后），把他们投入监狱：看，随后天主行奇迹。例如，带他们出监狱是天主的奇迹：带\Person{斐理伯}去是天主的奇迹：带\Person{保禄}回来也是天主的奇迹。—— 注意\Person{保禄}如何被尊荣，太监如何被尊荣。那里，\Person{基督}显现，可能是因为他的顽固，以及若非如此\Person{阿纳尼雅}就不会被说服。沉浸于这些奇事，让我们证明自己配得上。但现今许多人，甚至在他们来到教会时，不知道所读的是什么；而太监甚至在公共场合（\OriginalTerm{在公共场所}{\GreekText{ἐπ᾿} \GreekText{ἀγορᾶς}}）和坐在车中时，也致力于阅读圣经。但你们不是这样：没有人拿起圣经；相反，什么都有，就是没有圣经。\par

请告诉我，圣经是为什么的？就你所能做的，一切都已败坏。教会是为什么的？把圣经捆起来吧：或许审判不会如此，惩罚也不会如此；若有人把它们埋在粪堆里，以免听见它们，他所做的侮辱也不及你现在所做的。因为请告诉我，那里的侮辱是什么？那人把它们埋了。而这里的侮辱是什么？我们不听它们。请问，当一个人沉默不语，无人回应时，与当他说话却无人理会时，哪一样更受侮辱？所以目前的侮辱更大，当天主发言而你却不愿听时：轻蔑更大。我们读到，古时他们向先知们说：\BibleQuote{不要向我们说话} \BibleRef{依 30:10}，但你们做得更坏，说：说吧，我们不去做。因为那里他们转离他们，甚至不让他们说话，觉得从声音本身得到某种敬畏和约束；而你们，由于轻蔑至极，连这点也不做。相信我，若你们用手捂住我们的口来阻止我们，侮辱也不会像现在这样大。请说，哪一种表现出更大的轻蔑：是那听了但以动作阻止的人，还是那连听都不肯的人？请说——若我们视之为一个受侮辱的案例——假设一人阻止那侮辱他的人，并捂住他的口，因为被侮辱所伤；而另一人毫不在意，甚至假装没听见：哪一人表现出最大的轻蔑？难道你们不说是后者吗？因为前者显示出他感到被击中；后者几乎堵住了天主的嘴。你们对所听的战栗了吗？为什么？天主借以说话的口，就是天主的口。正如我们的口是我们灵魂的口，虽然灵魂没有口，同样先知的口是天主的口。听着，战栗吧。那里，执事（对全体会众）高声喊道：\Quote{让我们留意诵读。}那是整个教会的共同声音，他所发出的声音，然而无人留意。之后，诵读员开始说：\Quote{\Person{依撒意亚}的预言，}仍然无人留意，尽管预言毫无人的成分。此后，他又说：\Quote{主这样说，}仍然无人留意。此后是惩罚和报复，仍然无人留意。但常见的借口是什么？\Quote{老是同样的东西。}这最是毁灭你的。即使你知道了这些事，你也绝不该转离：因为在剧院里，不也是总演同样的东西，而你并不觉得厌烦吗？你竟敢说\Quote{同样的东西}，而你对先知名讳一无所知？你难道不羞愧地说，这就是你为何不听，因为\Quote{老是同样的东西}，而你对所读的那些名讳并不知晓，虽然你总在听同样的东西？你自己已承认说了同样的事。如果我这样说作为责备你的理由，你就得求助于一个完全不同的借口，而不是这个正是你所指责的。——难道你不劝诫你的儿子？现在若他说：\Quote{总是同样的东西！}难道你不认为这是一种侮辱？谈论\Quote{同样的东西}本应在我们既知道又行出来之后。或者更确切地说，即使那时，诵读也不是多余的。有谁像\Person{弟茂德}那样？告诉我那个：然而\Person{保禄}对他说：\BibleQuote{你要注重诵读、劝勉。} \BibleRef{弟前 4:13}因为不可能，我说不可能，耗尽圣经的心意。它是一口无底的井。训道者说：\BibleQuote{我说，我已有了智慧，智慧却离我而去。}—— \BibleRef{训 7:24} 我该向你们证明这些事不是\Quote{同样的}吗？你们以为有多少人论述过福音？但所有人都以新颖清新的方式论述。因为人越沉浸其中，就越有洞察，越看清纯洁的光。请看，我即将讲多少事：——请说，什么是叙述？什么是预言？什么是比喻？什么是预象？什么是寓意？什么是象征？什么是福音？只需回答我这个简单的问题：为什么它们被称为福音，\Quote{喜讯}？然而你们常听说喜讯不该有悲伤：但这\Quote{喜讯}却充满悲伤。它说：\BibleQuote{他们的火永不熄灭，他们的虫不死：} \BibleRef{谷 9:44} 它说：\BibleQuote{要划分他的份，与伪善者，与那些被腰斩的人：然后祂要说，我不认识你们，你们这些作恶的人，离开我去吧！} \BibleRef{玛 24:51} 确实，我们不自欺，我们以为在用你们的母语告诉你们。\par

\SourceRecord{Translated by J. Walker, J. Sheppard and H. Browne, and revised by George B. Stevens.；Nicene and Post-Nicene Fathers, First Series；Vol. 11.；Edited by Philip Schaff.；Buffalo, NY: Christian Literature Publishing Co.,；1889.；<http://www.newadvent.org/fathers/210119.htm>.}

% structure: p0001 source=h1 target=section reason=page-title

\section{论宗徒大事录第二十篇讲道}

\BibleRef{宗 9:10,12}\par

\BibleQuote{在达马士革有个门徒，名叫阿纳尼雅，主在异象中对他说：阿纳尼雅。他说：主，我在这里。主对他说：起来，往那条名叫直街的街上去，在犹大家里寻找一个名叫扫禄的塔尔索人；看，他正在祈祷，并在异象中看见一个名叫阿纳尼雅的人进来，按手在他身上，使他复明。}\par

为什么祂没有召请任何有权威和重要的人，也没有让这样的人来教导保禄呢？这是因为他不应被人引导，而只应由基督亲自引导；事实上，这个人没有教他什么，只是给他施了洗；因为一受洗（\OriginalTerm{光照}{\GreekText{φωτισθείς}}），他就该凭自己的热心和极大的热诚，为自己汲取圣神的恩宠。阿纳尼雅并不是什么非常杰出的人，这是显然的。因为经上说：\BibleQuote{主在异象中对他说，阿纳尼雅回答说：主，我听见许多人说这个人在耶路撒冷对你的圣徒行了多少恶事。}\BibleRef{宗 9:13} 因为如果他向主提出异议，那么若主派遣一位天使，他更会如此了。这就是为什么在先前的事例中，斐理伯也没有被告知是怎么回事；他看见天使，然后圣神就命他走近那辆车。但请注意主如何解除他的恐惧：祂说：\Quote{他是瞎的，} 祂说，\Quote{他正在祈祷，你还害怕吗？} 同样，梅瑟也害怕：所以这些话表明他害怕，退缩，并非不信。他说，我听见许多人说这个人。你说什么？天主说话，你却犹豫。他们那时还不明白基督的大能。\BibleQuote{并且他在这里有从大司祭得来的权柄，可以捆绑一切呼求你名的人。}\BibleRef{宗 9:14} 这事怎样知道的呢？很可能他们在恐惧中，仔细查问了。他说这话，并不是以为基督不知道事实，而是说：\Quote{既然如此，} 他说，\Quote{这些事怎么能成呢？} 正如福音中那些人说的：\BibleQuote{谁能得救？}\BibleRef{谷 10:26} 这样作，是为使保禄相信那来见他的人：\Quote{他在异象中看见了：} 它预先指示了他：\Quote{他祈祷，} 主说：不要怕。请注意，祂没有向他说所成就的事，教训我们不要谈论自己的成就。而且，虽然祂见他害怕，但为了这一切，祂没有说这事。\Quote{你不会不被相信：} \Quote{他看见了，} 祂说，\Quote{在异象中一个人，名叫阿纳尼雅：} 这就是为什么这是 \Quote{在异象中}，就是因为他瞎了。甚至连事之奇异也没有占据门徒的心，他如此害怕。但请注意：保禄瞎了眼，祂就是这样使他复明。\BibleQuote{去吧，因为他是我所拣选的器皿，为把我的名字带到外邦人、君王和以色列子民面前：因为我要指示他，为我的名字必须受多少苦。}\BibleRef{宗 9:15-16} \BibleQuote{不单，} 祂说，他将会是一个信徒，而且还会是一个老师，并会显示极大的胆量：\BibleQuote{在外邦人和君王面前}——教义将如此传播！——正如祂以前者使他惊讶，祂也以后者（使他更惊讶）。\BibleQuote{阿纳尼雅就去了，进了那家，按手在他身上，说：扫禄弟兄}——他立刻用这个名字称呼他如同朋友——\BibleQuote{耶稣，就是你在来的路上所显现给你的}（其实基督没有告诉他这事，是他从圣神得知的）\BibleQuote{打发我来，叫你能看见，并充满圣神。}\BibleRef{宗 9:17} 他说这话时，就按手在他身上。\BibleQuote{立刻有像鳞片的东西从他的眼睛上掉下来。}\BibleRef{宗 9:18} 有人说这是他瞎眼的记号。为什么祂没有完全弄瞎他的眼睛呢？这更奇妙，他睁着眼却看不见：\BibleRef{宗 9:8} 这正是他在法律方面的情况，直到耶稣的名放在他身上。\BibleQuote{他立刻看见了，就起来，受了洗。吃了食物，就恢复了力量。}\BibleRef{宗 9:19} 所以他软弱了，既因旅途，也因恐惧；既因饥饿，也因心灵沮丧。因此，为了加深他的沮丧，主使他瞎眼直到阿纳尼雅来到；并且，免得他想那瞎眼只是幻觉，这就是鳞片的缘故。他不需要别的教导：所发生的事成了他的教导。\BibleQuote{他同达马士革的门徒住了些日子。立刻在各会堂里宣讲耶稣，说他是天主子。}\BibleRef{宗 9:20} 看，他立刻就在会堂里做了老师。他不以这改变为耻，也不害怕，而从前他引以为荣的事，他现在却摧毁它们。甚至从他初上台时，这就是一个致人死命、准备好流血的人：你看见多么明显的标记吗？正是这事，他使众人害怕：因为他们说：他到这里来也就是为了这事。\BibleQuote{但凡听见的人都惊奇，说：这不是在耶路撒冷毁灭呼求这名的人吗？他到这里来，也正是为了要捆绑他们，带到司祭长那里吗？但扫禄越发强壮，驳倒了住在大马士革的犹太人，证明这位的确是基督。}\BibleRef{宗 9:21-22} 作为一个精通法律的人，他堵住了他们的口，不容他们说话。他们以为除掉了斯德望，就摆脱了这类问题的争论，却找到了另一个比斯德望更激烈的人。\par

（事略。）但让我们看看有关阿纳尼雅的事。主没有对他说：你去与他交谈，并教导他。因为，当祂说：\Quote{他正在祈祷，并看见一个人按手在他身上}（第11-12节）时，祂并未说服他，更不用说祂说这话了。所以，为使他不致不信你，祂说：\Quote{他在异象中看见了。}请注意，在先前的例子中，斐理伯也并非立即被告知一切。祂说：\Quote{不要惧怕，因为这个人是我所拣选的器皿。}\BibleRef{宗 9:15}祂完全解除了他的恐惧，如果情况是这个人将为我们的缘故如此热心，甚至要受许多苦。他也恰当地被称为\Quote{器皿}（或工具）——因为理性表明邪恶不是一种物理性质：祂说：\Quote{被选之器}（或：被拣选的工具）；因为我们选择那被认可的。不要有人以为阿纳尼雅因不信所说的话而发言，以为基督被欺骗了：远非如此！但惊慌战栗的他，甚至没有留心听所说的话，只因听到了保禄的名字。此外，主没有告诉他祂已使他失明：一提到这个名字，恐惧已占据了他的灵魂：他说：\Quote{看，}他说，\Quote{祢把我出卖给谁了：\InnerQuote{他到这里来，正是为了捆绑所有呼求祢名的人。}我怕他把我带到耶路撒冷：祢为什么把我投进狮子的口中？}他说话时仍是惊恐的；好让我们从各方面认识这人刚毅的\OriginalTerm{美德}{\GreekText{ἀρετήν}}。因为这些话若由犹太人说，并不奇怪；但这些人（信徒们）如此恐惧，正是天主大能的最有力证明。恐惧被显示出来，而恐惧后的服从更加伟大。因为确实需要力量。既然祂说\Quote{被选之器}，使你不致以为天主将做一切，祂又加上：\Quote{在外邦人、君王和以色列子民前，宣扬我的名。}阿纳尼雅听到了他最渴望的事——他将站在犹太人面前：这尤其给了他勇气。因为祂说：\Quote{我要指示他，为我的名必须受多么大的苦。}同时，这话也以令阿纳尼雅羞愧的方式说：如果那个如此疯狂的人将要受一切苦，你竟甚至不愿给他付洗！他说：\Quote{好吧，让他继续瞎眼}（这就是他说这话的原因）；\Quote{他是瞎的：祢为什么吩咐我开他的眼，使他再次捆绑（人）？}不要惧怕将来：因为他重见光明，将不是用来反对你，而是为了你（这是针对\Quote{使他复明}这句说的\BibleRef{宗 9:12}）：因为他不但不会伤害你，反而\Quote{要受许多苦}。奇妙的是，他将先知道\Quote{他要受多么大的苦}，然后才迎战危险。——\Quote{扫禄兄弟，主耶稣}——他没有说\Quote{那位使你瞎眼的}，而是说\Quote{在路上向你显现的那位，打发我来，叫你看见}\BibleRef{宗 9:17}：也请看这人如何不夸耀，就如伯多禄在瘸子的事上所说：\Quote{你们为什么注视我们，好像因我们自己的能力或虔诚使他行走呢？}\BibleRef{宗 3:12}同样，他也说：\Quote{那位向你显现的耶稣。}\Emph{(b)} 或者，是为了使对方相信：他没有说\Quote{被钉十字架的那位}、\Quote{天主子}、\Quote{行奇迹的那位}，而是说什么？\Quote{那位向你显现的}：根据对方所知道的说。正如基督也没有多说，没有说\Quote{我是耶稣，被钉死的，复活的}，而是说：\Quote{你所迫害的。}阿纳尼雅没有说\Quote{被迫害的}，免得他好像对他\OriginalTerm{}{\GreekText{ἐπενθουσί}\& 139v. 20, 21.}\Quote{他立刻在各会堂宣讲耶稣}——因为他在信德上是精确的——\Quote{说祂是天主子。凡听见的人都惊奇，}等等，因为确实难以置信。\Quote{但扫禄越发强壮，}等等。因此\Quote{过了多日}这事发生了：即犹太人\Quote{商议要杀掉他。扫禄知道了他们的计谋。}\EditorialNote{宗 22:22-24}这是什么意思？很可能有一段时间他不愿离开那里，虽然也许有许多人恳求他；但当他知道后，他就允许了他的门徒们：因为他立刻就有了门徒。\par

\BibleQuote{然后，门徒们，}等等。\BibleRef{宗 9:25} 论及此事，他说：\BibleQuote{阿勒特王的总督把守了大马士革人的城，想要捉拿我。}\BibleRef{格后 11:32} 但请注意这位作者在此处，他并未雄心勃勃地讲述故事，以显示保禄是何等重要的人物，说\Quote{因为他们煽动了王}等等，而只是说：\BibleQuote{于是门徒们在夜间用筐子把他从城墙上缒下去。}因为他们只让他一人出去，无人陪伴。他们这样做很好；结果是他得以在耶路撒冷出现在宗徒们面前。他们打发他出去，似乎必须藉逃跑来保全自己；但他却反其道而行之——他跳入了那些对他疯狂的人中间。这就是如火般的热情，这就是真正的热忱！从那天起，他明白了宗徒们所听到的一切诫命：\BibleQuote{谁若不背起自己的十字架跟随我}\BibleRef{玛 10:38} 正是他比别人晚来这一事实使他更加热心：因为\BibleQuote{那多得赦免的}\BibleRef{路 7:47}爱得更多，所以他来得越晚，爱得越深；* * * 而且既行了万千恶事，他认为自己无论如何做都不足以遮掩以前的作为。\BibleQuote{证明}\BibleRef{宗 9:22}，它说：即以温和教导。并且请注意，他们没有对他说：\Quote{你就是那毁灭者，为何改变了？}因为他们羞愧；但他们对自己说了。因为他会对他们说，这正是应教训你们的事，事实上他在阿格黎帕王面前正是这样辩护的。让我们效法这人，让我们将灵魂捧在手中，准备好面对一切危险。——（他从大马士革逃走）这并非怯懦：他保全自己是为了宣讲。如果他怯懦，他就不会去耶路撒冷，不会立即开始教导：他会稍微减弱他的热忱，因为他从斯德望的结局受了教训。他并不怯懦，但他也是 \OriginalTerm{有智慧的}{\GreekText{οἰκονομικός}}（在管理自己方面）。因此他认为为福音而死并非大事，除非这样做有极大益处：他甚至不愿看见基督——他最渴望看见的那位——只要他管理人的工作尚未完成。\BibleRef{斐 1:23}。基督徒的灵魂应当如此。从他初次出现的一开始，保禄的品性就显明了：甚至在此以前，甚至在他所做\BibleQuote{不按知识}\BibleRef{罗 10:2}的事中，他不是由于人的推理而行动如他所做的。因为如果过了那么久以后，他仍满足于不离去，那么在他贸易航行的开始，当他刚刚离开港口时更是如此！基督将许多事留给（普通的）人的智慧去做，使我们知道（祂的门徒）是人，并非处处都要藉恩宠完成：否则他们就会成为不动的木头；但在许多事上他们自己处理事情。——这并不亚于殉道——为了众多人的救恩而不逃避任何苦难。没有什么如此令天主喜悦。我将再次重复我常说的话；我重复，因为我极其渴望如此：正如基督也做了同样的事，当论及宽恕时：\BibleQuote{当你们祈祷时，若你们对某人有什么不满，就宽恕吧。}\BibleRef{谷 11:25}；又对伯多禄说：\BibleQuote{我不对你说，直到七次，而是到七十个七次。}\BibleRef{玛 18:22}。而祂自己也确实宽恕了对祂的过犯。我们也当如此，因为我们知道这正是基督徒的目标，我们不断谈论此事。没有什么比一个不关心他人救恩的基督徒更冷漠了。你不能在此以贫穷为借口：因为那投入两个小钱的寡妇要控告你。\BibleRef{路 21:1}。伯多禄说：\BibleQuote{金银我都没有。}\BibleRef{宗 3:6}。保禄如此贫穷，常常挨饿，缺乏必需的食物。你不能以出身卑微为借口：因为他们也是出身低微、父母低微的人。你不能以缺乏教育为借口：因为他们也是\BibleQuote{没有学问的人}\BibleRef{宗 4:13}。所以即使你是奴隶，甚至是逃亡的奴隶，你也能尽你的本分：因为敖尼息摩就是如此；然而看看保禄怎样称呼他，把他提升到何等尊荣的地位：他说：\BibleQuote{使他在我被囚时与我在一起}\BibleRef{费 1:13}。你不能以体弱为借口：因为弟茂德就是如此，常患疾病；因为宗徒说：\BibleQuote{为了你的胃和你的屡次患病，稍微用点酒吧。}\BibleRef{弟前 5:23}。每个人都能有益于邻人，如果他愿意尽自己的本分。你们岂不见那不结果的树，它们多么强壮、美丽，又高大、光滑、参天？但如果有一个花园，我们宁愿要石榴树或多结果实的橄榄树：因为前者只悦目，而无益处，它们的益处很小。那些只顾自己利益的人就是这样：不，甚至不是这样，因为这些人只适合焚烧：而那些树既可用于建造又可保护其中的人。那些童女也是如此，她们的确贞洁、端正、谦逊，但对谁都没有益处\BibleRef{玛 25:1}，因此她们被烧了。那些没有喂养基督的人就是这样。因为请注意，那些人中没有谁被指控犯有特定的罪，如奸淫或伪证；简言之，除了对他人无益之外，没有任何罪。那埋藏自己塔冷通的人就是这样，他固然显示了无可指责的生活，但对他人没有益处\BibleRef{玛 25:25}。这样的人怎能是基督徒呢？你说，如果酵母与面粉混合后没有将整个面团变成自己的性质，那还能叫酵母吗？同样，如果香水没有向靠近它的人散发香气，我们能称它为香水吗？不要说：\Quote{我不可能引导他人（成为基督徒）}——因为如果你是一个基督徒，就不可能不发生这样的事。因为正如事物的本性不可否认，这里也是如此：这事是基督徒本性的一部分。不要侮辱天主。说太阳不能发光，就是侮辱祂；说基督徒不能行善，就是侮辱天主，称祂为说谎者。因为太阳不发热、不发光，比基督徒不发光更容易；光变成黑暗，比这更容易。不要对我说这是不可能的：相反，那是不可能的。不要侮辱天主。如果我们把自己的事务整顿好，其他事必然会自然必然地随之而来。基督徒的光不可能隐藏；如此显眼的灯不可能被掩盖。我们不要疏忽。因为，正如美德的好处既达于我们自己，也达于受益的人；同样，邪恶带来双重的损失，既达于我们自己，也达于受害的人。假如（如果你愿意）有一个普通人，从某人那里遭受了无数祸患，没有人替他出头，然而那个人仍然回报善意；还有什么教导比这更有力？什么言语、什么劝勉能与之相比？什么愤怒不能因而熄灭和软化？因此，知道了这些事，让我们持守美德，知道若不在此生行这些善工，就不可能得救，好使我们也能获得未来的美好事物，借着我们的主耶稣基督的恩宠和慈爱，愿祂与圣父及圣神，一同享受光荣、威能、尊崇，从今世直到永远。阿们。\par

\SourceRecord{Translated by J. Walker, J. Sheppard and H. Browne, and revised by George B. Stevens.；Nicene and Post-Nicene Fathers, First Series；Vol. 11.；Edited by Philip Schaff.；Buffalo, NY: Christian Literature Publishing Co.,；1889.；<http://www.newadvent.org/fathers/210120.htm>.}

% structure: p0001 source=h1 target=section reason=page-title

\section{宗徒大事录讲道集 第二十一讲}

\BibleRef{宗 9:26-27}\par

\BibleQuote{当扫禄来到耶路撒冷，设法同门徒交结；众人都怕他，不信他是门徒。但巴尔纳伯接待了他，领他去见宗徒，并给他们讲述他在路上怎样看见了主。}\par

在此处，人可能大为困惑，难以理解何以在\BibleBook{迦拉达}{书}中，保禄说\BibleQuote{我没有上耶路撒冷}，而是\BibleQuote{去了阿剌伯}和\BibleQuote{到了大马士革}，并且\BibleQuote{过了三年，我才上耶路撒冷}，\BibleQuote{去见伯多禄}\BibleRef{迦 1:17}，（\OriginalTerm{访查}{\GreekText{ἱστορἥσαι}} Cat.），而在此处作者却说相反的情况。（那里，保禄说：）\Quote{我没有见过任何一位宗徒；但这里却说（巴尔纳伯）带他到宗徒那里。}——那么，要么（保禄）的意思是，我上去并不是为了向他们请教或依附他们（\OriginalTerm{请教}{\GreekText{ἀναθέσθαι}}）——因为他说了什么？\BibleQuote{我没有请教自己，也没有上耶路撒冷去见那些先我作宗徒的人：}要么，就是他在大马士革遇害是在他从阿剌伯回来之后；要么，又或者，他访问耶路撒冷是在他从阿剌伯回来之后。诚然，他并非自愿去见宗徒，而是\BibleQuote{尝试加入门徒中间}——作为一位教师，而非门徒——\BibleQuote{我没有去，}他说，\BibleQuote{为的是去见那些先我作宗徒的人：我确实没有从他们那里学到什么。}或者，他并非指这一次访问，而是略过了它，因此顺序是：\Quote{我去了阿剌伯，然后来到大马士革，然后到耶路撒冷，然后到叙利亚：}要么，又或者，他先上耶路撒冷，然后被派往大马士革，然后去阿剌伯，然后又去大马士革，然后到凯撒勒雅。另外，那\BibleQuote{过了十四年}的访问，很可能就是他带着（施舍给）弟兄们与巴尔纳伯一起上去的那一次：\BibleRef{迦 2:1}要么，他指的是另一个场合。\BibleRef{宗 11:30}因为历史家为求简洁，常常省略事件，压缩时间。请注意作者是多么不事张扬，他甚至没有叙述\BibleRef{宗 22:17}那个神视，而是略过了它。\BibleQuote{他尝试，}经上说，\BibleQuote{加入门徒中间。他们却都怕他。}这再次显明了保禄性格的热诚：不仅从阿纳尼雅的口中，以及那些在那里惊讶于他的人，也从耶路撒冷的人：\BibleQuote{他们不相信他是门徒：}因为那确实超乎人之常情。他不再是一头野兽，而是一个温柔和善的人！请注意，他没有去见宗徒，这是他的克制，而是去见门徒，如同为门徒一般。他不被信任。\BibleQuote{但巴尔纳伯}——\Quote{安慰之子}是他的称号，正因如此他也易于接近那个人：因为\BibleQuote{他是一个好人}\BibleRef{宗 11:24}，极其良善，这从眼前的事例以及关于若望（马尔谷）的事上都得到证明——\BibleQuote{他带了他来，领他到宗徒那里，并给他们述说了他怎样在路上看见了主。}\BibleRef{宗 15:39}很可能在大马士革他也听到了关于他的一切：因此\Emph{他}并不害怕，而其它人却害怕，因为他是一个目光慑人的人。\BibleQuote{怎样，}经上说，\BibleQuote{他在路上看见了主，主怎样向他说话，以及他在大马士革怎样因主的名放胆讲论。以后他就在耶路撒冷同他们出入，并且因耶稣的名放胆讲论}\BibleRef{宗 9:28}：这些事证实了先前的事，并且他用自己的行动证实了关于他所讲的话。\BibleQuote{他说话并和说希腊话的犹太人辩论。}\BibleRef{宗 9:29}所以门徒们怕他，宗徒们也不信任他；因此他消除了他们的恐惧。\Quote{和说希腊话的犹太人：}他是指那些说希腊话的犹太人：他这样做非常明智；因为其它那些深谙希伯来语的人根本不愿意见他。\BibleQuote{他们，}经上说，\BibleQuote{却设法要杀他：}这是他精力旺盛、取得胜利的记号，也是他们对所发生的事极其恼怒的证明。于是，怕后果与斯德望的情形一样，他们就打发他到凯撒勒雅去。因为经上说：\BibleQuote{弟兄们知道了这事，就带他下到凯撒勒雅，打发他往塔尔索去}\BibleRef{宗 9:30}，同时去传道，也似乎在更安全的地方，因为是他本乡。但请注意，我恳求你，事情绝非完全靠（奇迹性的）恩宠；相反，天主在许多事上让他们凭自己的智慧和人的方式自行处理；这样就堵住了懒散的借口：因为如果在保禄身上是这样，在他们身上就更该如此。然后，经上说：\BibleQuote{教会全犹太、加里肋亚和撒玛黎雅各地平安无事，他们而被建立，在敬畏主中行走，并因圣神的安慰而充满。}\BibleRef{宗 9:31}他即将叙述伯多禄下去（耶路撒冷），因此免得你将此归之于恐惧，他先说了这话。因为当有迫害时，他在耶路撒冷，但当教会到处都安全时，他反而离开了耶路撒冷。请看他是多么热心和积极！因为他没有想，既然有和平，就不需要他的临在。保禄离开了，和平就来了：既没有战争也没有动乱。那些（宗徒）他们最尊重，因为他们常常站在他们一边，并且受众人钦佩；而他人，他们则鄙视，并且对他更加残暴。请看，多大的战争，立刻就有和平！请看那战争产生了什么效果。它驱散了和平的缔造者。在撒玛黎雅，西满蒙羞；在犹太，发生了撒丕辣的事件。并非因为有了和平，事情就松懈了，而是这种和平也需要劝勉。\BibleQuote{当伯多禄巡行各处时，也到了住在里达的圣徒那里。}\BibleRef{宗 9:32}如同军队的统帅，他四处巡视，检阅队伍，看哪部分紧密，哪部分有序，哪部分需要他的临在。请看，他如何时常在前头行走。当要选宗徒时，他走在最前；当要对犹太人讲这些话不是醉了时，他走在最前；当要治好瘸子时，当要发表演讲时，他都在其他人之先；当要对官长说话时，他是那个人；当阿纳尼雅时，他\EditorialNote{宗 1:15；2:15；3:4-12；4:8；5:3-15}；当藉影子施行医治时，仍然是他。请看：哪里有危险，他在那里，哪里需要管理；但当一切平静时，他们都共同行事，他不要求比（其他人）更大的尊荣。当需要行奇迹时，他挺身而出，在这里又是他劳苦工作。\BibleQuote{在那里他遇见一个人，名叫艾尼雅，瘫痪病躺在床上已有八年。伯多禄对他说：艾尼雅，耶稣基督治好你了。起来，收拾你的床。他立刻起来了。}\BibleRef{宗 9:33-34}为什么他不等那个人表示信心，问他是否愿意被治好？首先，这奇迹对许多人都有劝勉作用：那么请听收获多大。\BibleQuote{所有住在里达和沙龙的人看见了他，就归向了主。}\BibleRef{宗 9:35}因为那个人是显赫的。\BibleQuote{起来，收拾你的床：}他很好地给出了奇迹的证明：因为他们不仅使人脱离疾病，而且在赐予健康的同时也赐予力量。此外，那时他们还没有能力证明他们的权能，所以不能合理地要求那人表明信心，正如在瘸子的事上他们也没有要求那样。\BibleRef{宗 3:6}正如基督在开始行奇迹时不要求信心，这些人也不要求。因为在耶路撒冷，确实是合理的，当事人的信心先显明：\BibleQuote{他们把病人抬到街上，若伯多禄经过，至少有他的影子可以遮在一些人身上}\BibleRef{宗 5:15}；因为那里行了许多奇迹；但在这里这是第一次发生。因为在奇迹中，有些是为了吸引（他人来信）而行的；有些是为了安慰那些已相信的人。\BibleQuote{在约培有一个女门徒，名叫塔彼达，翻译出来叫多尔卡：这妇人充满善行和施舍，她经常这样做。在那些日子里，她患病而死；他们洗了她，停放在楼上。因为里达靠近约培，门徒们听说伯多禄在那里，就打发两个人到他那里，求他不要延迟到他们那里去。}\BibleRef{宗 9:36-38}为什么他们等到她死了？为什么之前不请伯多禄（\OriginalTerm{打扰}{\GreekText{ἐσκύλη}}？他们是如此正直（\OriginalTerm{有道}{\GreekText{φιλοσοφοὕντες}}），认为不该为这类事\OriginalTerm{打扰}{\GreekText{σκύλλειν}}门徒们，使他们离开传道的事；正因此，经文提到那地方近，因为他们求这事是出于意外，并非在常规之中。\BibleQuote{不要延迟到他们那里去：}因为她是一个门徒。伯多禄就起身跟他们同去。到了，就领他上楼。\BibleRef{宗 9:39}他们没有恳求，而是让他来决定是否赐她生命（\OriginalTerm{救恩}{\GreekText{σωτηρίαν}}。）请看这里给人施舍有多么令人振奋的鼓励！\BibleQuote{所有的寡妇，}经上说，\BibleQuote{都站在他旁边哭，拿出多尔卡与她同住时所做的里衣外衣。}伯多禄进了房间，仿佛镇定自若，但请看后来获得了多大的成就！记述者并非无意义地告诉我们这妇人的名字，而是为了显示她的名字（\OriginalTerm{名如其人}{\GreekText{φερώνυμος} \GreekText{ἦν}}）与她的品格相符；她像羚羊一样活跃和警醒。因为在许多情形下，取名中带有天主的眷顾，正如我们常常对你们所说的。\BibleQuote{她充满，}经上说，\BibleQuote{善行：}不仅施舍，而是\BibleQuote{善行，}首先，然后尤其是这一善行。\BibleQuote{这些，}经上说，\BibleQuote{是多尔卡与她同住时所做的。}何等谦逊！不像我们这样；而是她们都在一起，她与她们一起做这些事并工作。\BibleQuote{伯多禄叫众人都出去，跪下祈祷，然后转向遗体说：塔彼达，起来。她睁开眼睛，见了伯多禄，就坐起来。}\BibleRef{宗 9:40}为什么他叫众人都出去？为不被他们的哭声扰乱和干扰。\BibleQuote{跪下祈祷。}请注意他祈祷的专注。\BibleQuote{他给她手。}\BibleRef{宗 9:41}基督对雅依洛的女儿也是这样做的：\BibleQuote{（福音记者说）拉着她的手。}请注意分别，先是生命，然后力量进入她，一个藉言语，一个藉他的手——\BibleQuote{他给她手，扶她起来，然后叫圣徒和寡妇们来，把活着的她交给他们：}对一些人是安慰，因为他们收回了自己的姐妹，也因为他们看见了奇迹，对另一些人是仁慈的照顾（\OriginalTerm{保护}{\GreekText{προστασίαν}}）。\BibleQuote{这事传遍了整个约培，许多人信了主。此后，他在约培住了多日，和一位硝皮匠西满同住。}\BibleRef{宗 9:42}请注意伯多禄的不张扬的态度，请注意他的节制，他没有与这位贵妇或以另一个显赫的人同住，而是与一个硝皮匠同住：他的一切行动都引导人谦卑，既不使低贱者羞愧，也不使高贵者傲慢！\Quote{多日；}因为那些因奇迹而相信的人需要他的教训。——让我们再回顾一下已经说过的话。\par

\BibleQuote{想}，经上说，\BibleQuote{与门徒结交}。\BibleRef{宗 9:26} 他没有厚颜地来到他们面前，而是带着谦逊的态度。那时，他们都因着伟大的德行被称为\BibleQuote{门徒}，因为门徒的样式清楚可见。\BibleQuote{但他们都害怕他}。请看他们怎样害怕危险，惊慌仍如何在他们心中高悬。\BibleQuote{但巴尔纳伯}等等。\BibleRef{宗 9:27} ——在我看来，巴尔纳伯从前是他的朋友——\BibleQuote{并述说了}等等：请注意保禄自己对此只字未提；若不是迫不得已，他也不会向别人提出此事。\BibleQuote{保禄就在耶路撒冷同他们来往，也大胆地讲道}（宗9:28-29）。这给了他们众人信心。\BibleQuote{但他们却想杀害他；弟兄们知道了}等等。\BibleRef{宗 9:30} 你看到吗？无论在那里（大马士革）还是在这里，其余的人都照顾他，为他提供离开的方法，我们哪里看到他此时受到了来自天主的（直接超自然）援助？所以他的性格活力得到了彰显。\BibleQuote{领他下到凯撒勒雅，打发他往塔尔索去}。所以，我猜想，他没有继续走陆路，而是航行了剩下的路程。这次（离开）是天主上智的安排，使他也能在那里宣讲；同样，针对他的阴谋也是天主上智安排的，他的耶路撒冷之行也是，关于他的故事不再被怀疑。因为在那里他\BibleQuote{大胆地讲道}，经上说，\BibleQuote{奉主耶稣的名；他并同希腊化犹太人交谈辩论}；又，\BibleQuote{他在他们中间来往出入。——教会既在全犹太、加里肋亚和撒玛黎雅得了平安}——即，她增长了；与自己和好，那是真正的平安：因为外来的战争不会伤害他们——\BibleQuote{建立，怀着敬畏主的心行事，并充满圣神的安慰}。圣神既借着奇迹和作为安慰他们，又独立于这些之外临到每个人身上。\BibleQuote{此后，伯多禄向他说：艾乃阿}等等。\BibleRef{宗 9:32} 但在讲论、劝诫之前，他对瘸子本人说：\BibleQuote{耶稣基督治好你}。这句话他无论如何都相信，于是痊愈了。请看他是多么不居功：因为他没有说\Quote{因这名}，而是作为一个记号，他叙述了奇迹本身，并作其宣讲者。\BibleQuote{凡住在里达和沙龙的人，看见了他，就都归向了主。——在约培有个女门徒}等等（宗9:35-36）。处处看到奇迹发生。但我们应相信它们，如同现在亲见一样。经上不是简单说塔彼达死了，而是说她死了，她一直在患病中。（然而）他们没有叫伯多禄直到她死了；然后\BibleQuote{他们派人去告诉他，请他不要迟延来到他们那里}。请看，他们派人去叫他，由他人代请。他来了：他并不认为被两个人召唤是有失尊重；因为经上说：\BibleQuote{他们派了两个人到他那里}。——亲爱的，痛苦是件大事，它凝聚我们的灵魂。那里没有一点哀哭，也没有丧恸。看事情是多么洁净！\BibleQuote{他们洗净了她}，经上说，\BibleQuote{把她停放在楼上}：也就是说，他们为尸体做了一切（妥当的事）。然后伯多禄来了，\BibleQuote{跪下祈祷；转身向遗体说：塔彼达，起来}。\BibleRef{宗 9:40} 他们不是同样轻易地行所有奇迹。但这对他们有益：因为天主不仅为别人的救恩着想，也为他们自己的着想。那位凭影子就治愈这么多人的，怎么现在要先做这么多事？也有些情况是求者的信德合作。这是他复活的第一位死者。请看他是如何像唤醒她出睡一样：首先她睁开眼：然后看见（伯多禄）就坐起来：接着从他的手得了力量。\BibleQuote{这事传遍了整个约培，许多人信了主}。\BibleRef{宗 9:42} 注意收获，注意果实，这不是为了炫耀。的确，这就是为什么他打发众人出去，在此事上也效法他的师傅。\par

因为哪里有眼泪——更确切地说，哪里有奇迹，那里就不该有眼泪；不是在那里庆祝这样的奥迹。听，我恳求你们：虽然现在没有发生类似的事，但在我们死者的情形中，同样有一个伟大的奥迹正在庆祝。请问，如果我们坐在一起，皇帝派人来邀请我们中的一位进宫，我们还能哭泣哀悼吗？天使在场，受命从天堂而来，从那而来，被君王亲自派遣来召唤他们的同仆，他们说，你在哭泣吗？难道你不知道正在发生的是何等奥迹，何等可畏，何等可怖，实在值得颂歌和赞美吗？你愿意学习，好让你知道，现在不是流泪的时候吗？因为这是天主上智的一个非常伟大的奥迹。灵魂仿佛离开她的居所，出发，急速前往她的主那里，而你却哀悼？那么，你倒应该在孩子出生时这样做：因为实际上这也是一种出生，而且是更好的一种。因为在这里，她出发到一个非常不同的光明中，如同从监狱中被释放，如同从竞赛中退出。\Quote{是的，}你们说，\Quote{说这些固然很好，但只是对于那些我们已经确定其救恩的人。}那么，人啊，你有什么毛病，连在这种情况下，你也不这样看待呢？请问，你能在小孩身上找到什么可定罪的？你为什么为他哀悼？新受洗的又有什么？因为他也被带入同样的状况：你为什么为他哀悼？因为如同太阳升起，清澈明亮，同样，灵魂带着纯洁的良心离开身体，欢欣地闪耀。皇帝驾临城市时的景象（\OriginalTerm{进入城市}{\GreekText{ἐπιβαίνοντα} \GreekText{πόλεως}}）不能与之相比，那种敬畏的寂静不能与之相比，就像灵魂离开身体后，在天使的陪伴下离去。想想那时灵魂会是怎样！何等惊讶，何等惊奇，何等喜悦！你为什么哀悼？回答我。——但你们只在罪人的情况下才这样做？但愿如此，我就不会禁止你们的哀悼了；但愿这是目的！这样的哀悼是宗徒的，是效法主的榜样；因为连耶稣也曾为耶路撒冷哭泣。但愿你们的哀悼能以此标准来分辨。但当你们说出那些想要召回（死者）的话，并谈到你们长久以来的亲密和他的恩惠时，你们只是为了这个哀悼（并非因为他是个罪人），你们只是假意这样说。哀悼吧，为罪人哭泣，我也会放声流泪；我比你更甚，因为他作为罪人应受的惩罚更大：我也会为此目的而哀恸。但不仅你一个人必须为这样的人哀悼；整个城市都必须如此，以及所有在路上遇见你的人，如同人们为那些被带去处死的人哀哭。因为这才是真正的死亡，一种邪恶的死亡，罪人的死亡。但在你们那里，一切都完全颠倒了。这样的哀悼标志着崇高的心灵，并传达许多教诲；另一种则标志着灵魂的渺小。如果我们都用这种哀悼来哀悼，我们就能在这些人还活着的时候纠正他们。因为，如果由你决定使用药物来防止肉体死亡，你会使用它们；同样，如果\Emph{这种}死亡是你们哀悼的死亡，你们就会防止它的发生，无论是在你们自己身上还是在他身上。然而现在我们的行为完全是个谜：我们本有能力阻止它的到来，却任由它发生，等它来了再为此哀悼。他们确实值得哀悼（当我们思考时），当他们站在基督的审判台前时，他们那时将听到什么话，将遭受什么！这些人活着毫无目的：不，不是毫无目的，而是成了邪恶的目的！对他们也可以恰当地说，\Quote{那人没有生下来为他倒好。} \BibleRef{谷 14:21}请问，花了这么多时间损害自己，有什么益处呢？如果只是白白浪费，我问你，这难道还不够惩罚吗？如果一个做了二十年雇工的人发现他所有的劳动都是徒劳，他岂不要哭泣哀悼，认为自己是最悲惨的人吗？看，这里有一个人，失去了整个一生的劳动：他没有一天为自己而活，而是为了奢侈、放荡、贪婪、罪恶、魔鬼。那么，请问，我们难道不该为这个人哀哭吗？我们难道不该尽力把他从危险中拯救出来吗？因为，是的，如果愿意，我们有可能减轻他的惩罚，如果我们为他不断祈祷，如果我们为他施舍。无论他多么不配，天主都会因我们的恳求而让步。因为如果保禄曾怜悯一个（本不应受怜悯的人），并且为了他人的缘故饶恕了一个（他本不会饶恕的人），那么我们更应该这样做。借着你的财物，借着你的，借着你的一切方式，帮助他：倒上油，不，倒上水。他没有自己的施舍可以展示吗？至少让他拥有他亲属的施舍。他自己没有做过施舍？至少让他拥有那些为他做的施舍，这样他的妻子能在那天自信地为他求情，因为她为他付出了赎金。他需要承担的罪越多，就越需要施舍，不仅因为这一点，还因为现在的施舍没有同样的功效，而是小得多：因为自己去做和别人替他做是不一样的；因此，既然功效较小，让我们通过数量使它成为最大的。我们不要忙于墓碑，不要忙于纪念物。这才是最大的纪念：让寡妇们围着他站立。告诉她们他的名字：吩咐她们都为他祈祷，为他恳求：这能说服天主：即使不是那人自己做的，但为了他的缘故，另一个人成了施舍的施行者。这甚至也属于天主的仁慈：\Quote{围着他站立的寡妇们在哭泣}知道如何救助，不是从现在的死亡中，而是从将来的死亡中。许多人甚至因别人为他们所做的施舍而受益：因为即使他们没有获得完全的（拯救），至少他们从中得到了一些安慰。如果不是这样，孩子们是如何获得救恩的？然而在那里，孩子们自己没有任何贡献，而是他们的父母做了一切；而且常常有妇女得到她们的孩子，尽管孩子自己没有任何贡献。天主赐予我们救恩的方式很多，只让我们不要疏忽。\par

\Quote{那么，如果一个人贫穷呢？}你说。我再次说，施舍的伟大并不取决于给予的数量，而是取决于心意。只要你给予的不少于你的能力，你就已经付清了。那么，你说，如果他孤苦伶仃，是外乡人，没有人照顾他呢？我请问你，为什么他没有人呢？正是在这件事上，你遭受了应得的报应，就是你没有这样一个有德行的朋友。这是如此安排的，为的是，即使我们自己没有德行，我们也可以努力拥有有德行的伴侣和朋友——无论是妻子、儿子还是朋友——因为通过他们我们可以获得一些好处，虽然微小，但确实是好处。如果你不把娶富有的妻子作为主要目标，而是娶一个虔诚的妻子，有一个虔诚的女儿，你就将获得这种安慰；如果你努力使你的儿子不是富有而是虔诚，你也将获得这种安慰。如果你以此为目标，那么你自己也会变得像他们一样。这也是德行的一部分：选择这样的朋友、这样的妻子和儿女。为逝者所做的祭品不是徒然的，祈祷不是徒然的，施舍也不是徒然的：这一切都是圣神所安排的，希望我们彼此受益。看：他得益处，你得益处；因为他，你轻视了财富，被激励去做一些慷慨的事：你使他获得救恩，而他为你提供了施舍的机会。不要怀疑他会因此得到益处。执事呼喊：\Quote{为那些在基督内安眠的人，以及为那些为他们举行纪念的人。}这不是执事发出的声音，而是圣神：我指的是恩赐。你说什么？祭献正在进行，一切已妥善安排：天使在场，总领天使在场，天主子也在那里：大家都如此敬畏地站着，在一片肃静中，那些站在旁边的人大声呼喊：你以为所做的一切都是徒然的吗？那么其余的一切不也是徒然的吗？包括为教会、为司铎、为全体所做的祭品？绝对不可！一切都是凭信德做的。你以为为殉道者所做的祭品如何？在那个时刻所作的呼求，即使他们是殉道者，但甚至\Quote{为殉道者？}在主面前被提名是极大的尊荣，当那个纪念举行时，可畏的祭献，不可言喻的奥秘。因为正如当皇帝就座时，是恳求者实现他愿望的时刻，但当他起身时，无论他说什么，都是徒然的；同样，在举行奥秘的时候，对所有人来说，被提名为最大的尊荣。因为看：那时宣告了可畏的奥秘：天主为世界舍弃了自己；伴随这个奥秘，他适时地提醒天主那些犯罪的人。因为正如当皇帝的胜利庆典进行时，那些在胜利中有份的人得到称赞，同时那些被捆绑的人因这个场合而获得释放；但当庆典过去后，那时没有获得这种恩惠的人就不再得到了：同样在这里：这是一个胜利庆典的时刻。因为经上说：\BibleQuote{你们每次吃这饼，就是宣告主的死亡。}那么，我们不要漫不经心地接近，也不要以为这些事情是随意进行的。但我们在另一种意义上提及殉道者，这是为了确证主没有死；这是死亡已受致命一击的记号，死亡本身已死。知道这些，让我们为逝者设想安慰，而不是眼泪，而不是哀悼，而不是坟墓，而是我们的施舍、祈祷、祭品，为使他们和我们都能获得所应许的福乐，因祂的独生子、我们的主耶稣基督的恩宠和慈爱，愿光荣、权能、尊崇偕同圣父和圣神归于祂，自今而始，迄于无穷，永世无疆。阿们。\par

\SourceRecord{Translated by J. Walker, J. Sheppard and H. Browne, and revised by George B. Stevens.；Nicene and Post-Nicene Fathers, First Series；Vol. 11.；Edited by Philip Schaff.；Buffalo, NY: Christian Literature Publishing Co.,；1889.；<http://www.newadvent.org/fathers/210121.htm>.}

% structure: p0001 source=h1 target=section reason=page-title

\section{论宗徒大事录第二十二篇讲道}

\BibleRef{宗 10:1-4}\par

\BibleQuote{在\OriginalTerm{凯撒勒雅}{Cæsarea}有一个人，名叫\OriginalTerm{科尔乃略}{Cornelius}，是意大利营的百夫长。他是个虔诚的人，合家都敬畏天主，多多周济百姓，常常向天主祈祷。约在第九时辰，他在异象中清楚看见一位天使进来，向他说：科尔乃略！他注目看天使，就害怕，说：主，有什么事？天使向他说：你的祈祷和你的周济已上达天主面前，蒙纪念。}\par

这人不是犹太人，也不是属于法律的人，但他早已预尝了我们的生活方式。请注意，到目前为止，有两位身居高位的人接受了信仰：加萨的太监和这个人；以及对这两个人所费的心力。但不要以为这是因为他们的高位：绝对不可！而是因为他们的虔诚。因为圣经提及他们显赫的地位，是为了显示他们虔诚的伟大；因为一个人身处财富和权力之位却能如此，就更显奇妙。称赞前者的是：他进行如此长途的旅行，且并非节期所要求，在路上、坐在车上时阅读，恳求斐理伯，以及无数其他优点；称赞后者的是：他施舍、祈祷，是正义的人，且身居如此要职。作者如此详尽地描述此人，是为了没有人会说圣经记载的是虚假的：他说：\Quote{科尔乃略，} \Quote{是意大利营的百夫长。} \BibleRef{宗 10:1} \Quote{营}，\OriginalTerm{营}{\GreekText{σπεἵρα}}，就是我们今天所说的\Quote{大队}。\Quote{他是一个虔诚的人，}他说，\Quote{且全家敬畏天主} \BibleRef{宗 10:2}：好叫你不要以为是因为他的高位才发生这些事。——当保禄要归化时，没有天使，而是主自己；主没有打发他去某个大人物，而是去一个非常平凡的人；但在这里，相反地，祂把首席宗徒带到（这些外邦人）那里，而没有打发他们到他那里：在此祂俯就他们的软弱，知道这样的人需要如何被对待。事实上，在许多场合我们发现基督亲自急忙（到这样的人那里），因为他们更为软弱。或者（可能是因为）（科尔乃略）自己不能离开家。但这里再次对施舍给予高度赞扬，正如先前通过塔彼达所给的。\Quote{一个虔诚的人，}它说，\Quote{且全家敬畏天主。}让我们听听这个，我们中间谁忽略了自己家中的人，而这个人却关心他的兵士。\Quote{又多多施舍，}它说，\Quote{给百姓。}他的教义和生活都是正直的。\Quote{有一天，约在第九时辰，他在异像中清楚看见天主的一位天使进来，向他说：科尔乃略。} \BibleRef{宗 10:3} 为什么他看见天使？这也是为了伯多禄的确信，或者更确切地说，不是为了他，而是为了其他人，那些软弱的人。\Quote{在第九时辰，}当他摆脱了忧虑、安静下来时，当他专心祈祷和痛悔时。\Quote{他注目看天使，便害怕起来。} \BibleRef{宗 10:4} 请注意天使所说的话并非立即说出，而是先激动、提升他的心神。看见时，有恐惧，但恐惧适中，恰足以引起他的注意。然后话语又解除了他的恐惧。恐惧激动了他；赞赏减轻了恐惧中不愉快的一面。\Quote{你的祈祷，}他说，\Quote{和你的施舍已上达天主面前，获得纪念。现在你要打发人往约培去，请一位名叫西满、号伯多禄的人来。} \BibleRef{宗 10:5} 免得他们找错人，他不仅用姓氏指定此人，还用地方。\Quote{这人，}他说，\Quote{住在一位名叫西满的皮匠家里，他的房子在海边。} \BibleRef{宗 10:6} 你看见宗徒们如何为了爱清静和安宁，喜欢城市的偏僻角落吗？\Quote{住在一位名叫西满的皮匠家里：}那么，假如偶然另有其人呢？看，还有另一标记：他住在海边。三个标记不可能（在其他地方）同时成立。他没有告诉他原因，免得削弱他热切的渴望，而是让他急切盼望将要听到的话。\Quote{向科尔乃略说话的天使走了以后，他便叫了两个家人，和常服侍他的一个虔诚兵士，把一切事都讲给他们听，然后打发他们往约培去。}\BibleRef{宗 10:7-8}你看见作者这样说并非没有用意吗？（这说明）那些\Quote{常服侍他的}也是如此之人。\Quote{他把一切事都讲给他们听：}请注意此人的谦逊：他没有说：\Quote{把伯多禄给我叫来}，而是为了引他来，把一切事都讲明了——这是上智的安排——因为他不愿利用自己地位的权柄把伯多禄召来；因此\Quote{他把事情讲明了}；这是此人的节制：而且对一个住在皮匠家里的人，也不要抱太高想法。\Quote{第二天，他们上路，将近那城的时候} \BibleRef{宗 10:9} ——请注意圣神如何连接时间：不早不晚，恰在此刻祂让此事发生——\Quote{伯多禄约在第六时辰上屋顶祈祷：}即私下、安静地，如同在楼上。\Quote{他就觉得饿了，想要吃东西；但他们还在预备的时候，\OriginalTerm{神魂超拔}{\GreekText{ἔκστασις}}来到他身上。} \BibleRef{宗 10:10} 这句话是什么意思，\Quote{神魂超拔？}更确切地说，有一种\OriginalTerm{属灵的观见}{\GreekText{θεωρία}}呈现给他：灵魂，可以说，\OriginalTerm{被引出体外}{\GreekText{ἐξέστη}}。\Quote{他看见天开了，一块好像大布的东西，四角系着，从天降下，放到地上。里面有地上各种四足兽、野兽、爬虫和天空的飞鸟。有声音向他说：伯多禄，起来，宰了吃！伯多禄却说：主，绝对不可！因为我从未吃过任何凡俗或不洁之物。第二次又有声音向他说：天主洁净的，你不可称为凡俗！这事一连三次，那布就立刻收回天上去了。} \BibleRef{宗 10:11} 这是什么？它是全世界的象征。那个人是未受割损的；并且——因为他与犹太人毫无共同之处——他们全会指责他是违法者：\Quote{你进了未受割损的人的家里，且与他们一同吃饭了！} \BibleRef{宗 11:3} 这对他们来说完全是冒犯的事；那么请看上智如何安排。他自己也说：\Quote{我从未吃过：}并非他自己害怕——愿此事远离我们！——而是圣神如此安排，为使他能对那些指责他的人回答说，他曾反对过：因为他们遵守法律是完全必要的。他将被派往外邦人那里；因此，为了不让这些（犹太人）也指责他，请看有多少安排（天主的眷顾）。因为，为避免被认为只是想像，\Quote{这事一连三次。我说，}他说，\Quote{主，绝对不可！因为我从未吃过任何凡俗或不洁之物。——声音又向他说：天主洁净的，你不可称为凡俗！} \EditorialNote{参见 11:8 及 10:14} 这话似乎是对他说的，但全部是对犹太人说的。因为如果教师被责备，这些人更甚。那么，大地，这就是布所象征的，其中的野兽是外邦人，命令\Quote{起来，宰了吃}表示他必须也到他们那里去；并且这事行了三次，表示圣洗。\Quote{天主洁净的，}它说，\Quote{你不可称为凡俗。}好大的胆量！为什么他反对？为了没有人说天主在试探他，如同在亚巴郎的事上，这就是为什么他说：\Quote{主，绝对不可！}等等，并非反驳——正如祂对斐理伯说：\Quote{你们有多少饼？}不是为了知道，而是为试探他，或\Quote{考验他。}然而，同一位主曾（在法律中）论述过洁净和不洁之物。但在那块布中也有\Quote{地上各种四足兽：}洁净的与不洁的在一起。尽管这样，他仍不知道这是什么意思。\Quote{伯多禄正在猜疑，这异像究竟有什么意思，忽然科尔乃略派来的人，已经查问到西满的家，站在门前，喊着问：有没有一位号称伯多禄的西满在这里寄宿？但伯多禄，}它说，\Quote{正在猜疑之间}\BibleRef{宗 10:17-18}，那些人适时来到，解除了他的猜疑：正如（主）先让若瑟内心不安，然后才派天使：因为当灵魂先困惑时，就更容易接受解释。他的困惑既不是（发生时的）长久的，也不是在此之前，而恰好在他们\Quote{问是否他寄宿在那里时。伯多禄正沉思异像，圣神向他说：看，有三个人找你。起来，下去，不要疑虑，跟他们一同去，因为是我打发他们来的。}（参照前文142页和145页注7；\BibleRef{宗 10:19-20}）这又是为伯多禄在门徒面前的一个辩护：他曾疑惑，并被教导不要疑虑。\Quote{因为，}祂说，\Quote{是我打发他们来的。}圣神的权威何其大！天主所做的，圣神也被说成是做。天使不是这样，而是先说了\Quote{你的祈祷和你的施舍已上达天主面前，获得纪念}，以表明他从那里被派来，然后才接着说：\Quote{现在你要打发人}等等：圣神却不是这样，而是说：\Quote{因为是我打发他们来的。于是伯多禄下来，到科尔乃略派来的那些人那里，说：看，我就是你们所找的人。你们为什么事而来？他们说：科尔乃略百夫长，一个正义的人，敬畏天主，为全犹太民族所称赞，他从一位圣天使获得指示，要请你到他家里去，听你的话。}\BibleRef{宗 10:21-22}他们赞扬他，为说服他天使确实向他显现了。\Quote{于是伯多禄请他们进去，}（b）好使他们不受伤害，\Quote{招待他们住下：}从此以后，他毫无顾忌地与他们一同吃饭。\Quote{第二天，伯多禄起身和他们一同去了，还有几个凯撒勒雅的弟兄陪着他们。又过了一天，他们进了凯撒勒雅。}\BibleRef{宗 10:23-24}那人是个重要人物，而且当时他在一个重要的城市里。\par

（a）但我们再来审视一下已经说过的话。\BibleQuote{在凯撒勒雅有一个人，等等。}（重述，第1-2节）注意，外邦人的开端是从什么样的人开始的——是从一个\BibleQuote{虔诚的人}，一个经其行为证明为相称的人。因为如果情况如此，他们仍感到冒犯，倘若情况不是这样，后果将会多么严重！但你要注意这确证的伟大之处。\Emph{(c)}为此，一切都是（以这种方式）完成的，事情是从犹太开始的。\Emph{(d)}\BibleQuote{他明明在异象中看见，}等等。\BibleRef{宗 10:3}。天使不是在睡梦中向他显现，而是在他清醒时，白天，\BibleQuote{约在第九时辰，他看见天主的一位天使进来，向他说：科尔乃略。他注目看天使，就害怕起来。}他如此专注于自己的事。意指，他看见天使是因为天使以声音呼唤他；因为若天使没有呼唤他，他就不会看见他：他完全沉浸在自己所做的事中。但天使向他说：\BibleQuote{你的祈祷和施舍已上升到天主面前，成为纪念。现在你派人往约培去，请一位叫西满、人称伯多禄的人来。}\BibleRef{宗 10:5}至此，他暗示派人去请他会带来好的结果，但如何好，他并未透露。同样，伯多禄也没有讲述全部，而是处处只作部分叙述，为使听者专注于所说的内容。\BibleQuote{你派人去请西满来。}同样，天使只呼唤斐理伯。\BibleQuote{他们上路前行，将近城的时候，}\BibleRef{宗 10:9}：为的是伯多禄不至于长期困惑。\BibleQuote{伯多禄上了房顶，}等等。注意，甚至连饥饿也没有迫使他靠近那块布。有声音说：\BibleQuote{伯多禄，起来，}\BibleQuote{宰了吃吧。}\BibleRef{宗 10:13}很可能他看见异象时是跪着的。——在我看来，这也象征福音（或\Quote{宣讲}）。所发生之事出于天主（环境已表明，即），他看见它从上面降下，且他处于神魂超拔之中；声音从那里传来，三次承认那里的动物是不洁的，它从那里降下，又被收回那里（这一切），都是洁净（赋予它们）的强力标记。——但为什么要这样做？是为了以后那些他将要向他们讲述此事的人。因为曾对他自己说：\BibleQuote{外邦人的路，你们不要走。}\BibleRef{玛 10:5} \* \* 因为倘若保禄既需要（受）割礼，又需要（献）祭，那么当时，在宣讲之初，当人们尚且软弱时，（更需要某种确证）。\BibleRef{宗 16:3}——还要注意，他如何没有立刻接待他们。因为经上说，他们\BibleQuote{喊叫，问是否有一位叫西满、又称伯多禄的在这里寄宿。}\BibleRef{宗 10:18}因为房子看起来很简陋，他们在下面问，向邻居打听。\BibleQuote{伯多禄还在思索时，圣神向他说：起来，下去，不要疑惑，尽管跟他们去，因为是我差他们来的。}（第19-20节）他并没有说：为此异象向你显现；而是说：\BibleQuote{是我差他们来的。}于是伯多禄下去\BibleRef{宗 10:21}——这就是必须服从圣神的方式，不问理由。因为对一切确证而言，被祂告知做这个，信这个就足够了：不再需要别的——\BibleQuote{于是伯多禄下去，向那两个人说：看，我就是你们所要找的人。你们来的原因是什么？}他看见一个士兵，看见一个人：他并不是害怕，相反，他先承认自己就是他们所找的人，然后才问（他们来的）原因；为免有人以为他问原因是因为想隐藏自己：（他问原因）是为了，如果事情紧急，他可以立即跟他们走，如果不是，则可接待他们为客。\BibleQuote{他们便说，等等，进入他的家。}\BibleRef{宗 10:22}这是他曾吩咐他们的。不要以为他这样做是出于轻视：他派人去并非出于轻视，而是奉命行事。\BibleQuote{科尔乃略等待着他们，并召集了他的亲戚和密友。}\BibleRef{宗 10:24}他的亲戚和朋友们聚集到他那里是合宜的。但他们既在场，就会从他那里听到（所发生的事）。\par

请看哀矜的德行多么伟大，无论是在先前的言论中，还是在这里！在那里，它从暂时的死亡中拯救；在这里，它从永恒的死亡中拯救，并打开了天堂之门。为使\Person{科乃略}归信，付出了如此大的努力：天使被派遣，圣神运行，宗徒之长被召到他那里，如此异象被显示，简言之，没有一件事被忽略。那时有多少百夫长、千夫长和君王，却没有一个得到这个人所得的！所有身居军职、侍立君王身边的人，请听。经上说：\BibleQuote{一个义人，敬畏天主，虔诚}（2,22节）；而最重要的是，他全家如此。不像我们（谁）：我们为了让仆人们害怕我们，什么都做，却不是为了让他们虔诚。对家人也是如此，\*\*。这个人却不是这样；他\BibleQuote{全家敬畏天主}\BibleRef{宗 10:2}，因为他如同与他一起的人和所有其他（在他指挥下的）人的共同父亲。但请注意（士兵）自己说的话。因为，害怕\*\*，他还加上：\BibleQuote{全国人都称赞他}。即使他未受割礼，那又怎样？但那些人称赞他。没有什么比哀矜更伟大：这实践的德行是巨大的，当哀矜从纯洁的宝库中涌出时；因为如果它从不义的宝库中涌出，就如同喷出泥浆的泉源；但如果从正当的所得中涌出，它就像乐园中清澈纯净的溪流，悦目，悦触，轻盈凉爽，在正午的炎热中施予。哀矜正是如此。在这泉源旁，不是白杨、松树或柏树，而是其他比这些更好的树木，高耸挺拔：与天主的友谊，人的赞美，对天主的荣耀，来自众人的善意；罪的涂抹，巨大的胆量，对财富的轻视。这是滋养爱德之树的泉源：因为没有什么比慈悲更滋养爱德；它使枝条高高扬起。这泉源比乐园中的那泉源更好\BibleRef{创 2:10}；不是分为四道，而是直抵天堂本身；它生出了那条\BibleQuote{涌向永生}的河流\BibleRef{若 4:14}。让死亡落在这泉源上，它便如火花被泉源熄灭：无论它滴落在哪里，都会产生如此巨大的祝福。这水甚至像熄灭火花一样熄灭那火河：这水如此窒息那虫子，没有别的能做到\BibleRef{谷 9:44}。拥有这水的人，将不会咬牙切齿。让这水的一滴落在锁链上，便能解开；只要它触及火把，便能全部熄灭。——泉源不会时而涌流，时而干涸——否则它就不再是泉源——而是不断涌出：所以让我们的泉源更加丰沛地涌出慈悲（哀矜）的溪流。这使接受者喜悦：这就是哀矜，不仅要涌出丰沛的溪流，而且要涌出永久的溪流。如果你愿天主像从泉源般倾泻祂的慈悲给你，你也要有泉源。然而，两者无法相比：因为如果你打开这泉源的出口，天主泉源的口是如此之大，以至超越一切深渊。天主不过是在寻找我们这边的机会，便从祂的宝库中倾泻祝福。当祂花费、当祂慷慨施予时，祂便富有，祂便丰裕。那泉源的口宽阔：其水纯净清澈。如果你不堵塞这里的泉源，那泉源也不会被堵塞。——不要让不结果子的树立在它旁边，以免浪费它的水汽。你有财富吗？不要在那里种白杨：因为奢侈就是如此：它消耗很多，却显不出什么，反而毁坏果实。不要种松树——这是衣着上的放荡，只见其美，毫无用处——也不要种枞树，或其他这类只消耗却无用的树。让它密植嫩枝：将一切结果子的，种在穷人手中，种你愿意的一切。没有比这土地更肥沃的了。虽然手伸出的范围很小，但所栽种的树却直冲天际，屹立不倒。这就是栽种。因为栽种在地上的树，即使不是现在，也终将在百年后朽坏。你种许多树，却不能享用其果实，在你能享用之前，死亡便临到你。这棵树会在你死后给你果实。——如果你栽种，不要栽种在饕餮的胃口里，免得果实最终进入粪池；而要栽种在饥饿的肚腹中，使果实升上天堂。安慰穷人的困乏灵魂，免得你压抑自己宽敞的灵魂。——你岂不见，浇水过多的植物根部腐烂，而适度浇水则生长茂盛？同样，不要过分浇灌你自己的肚腹，免得树根腐烂：要浇灌那干渴的，使它结果。如果你适度浇水，太阳不会晒枯；但如果过量，太阳便会晒枯：这就是太阳的本性。凡事过度都有害；因此让我们断绝过度，使我们也能得到我们所求的。——据说，泉源常在最高的地方涌出。让我们的灵魂高升，我们的哀矜便会如急流涌出：高升的灵魂不能不慈悲，慈悲的灵魂也不能不高升。因为轻视财富的人，高于万恶之根。——泉源常在偏僻处被发现：让我们的灵魂从人群中退隐，哀矜便会从我们身上涌出。泉源越是被清理，流得越丰沛：同样，我们越花费，一切美善就越是增长。——拥有泉源的人，无所畏惧：因此我们也不要害怕。因为这泉源对我们有益，可供饮用、灌溉、建造，一切用处。没有比这饮水更好的了：它不会醉人。拥有这样的泉源，胜过拥有涌出黄金的泉源。比一切产金之地更好的，是产生这黄金的灵魂。因为它提升我们，不是进入地上的宫殿，而是进入天上的宫殿。黄金成为天主教会装饰。这黄金铸成了\BibleQuote{圣神的剑}\BibleRef{弗 6:17}，就是斩杀巨龙的剑。从这泉源中，产生了君王头上的宝石。那么，让我们不要忽视如此巨大的财富，而要慷慨地奉献我们的哀矜，使我们堪当获得天主的慈悲，借着祂独生子的恩宠和温柔怜悯；愿光荣、权能、尊崇归于祂，偕同圣父及圣神，从今世直到永远。亚孟。\par

\SourceRecord{Translated by J. Walker, J. Sheppard and H. Browne, and revised by George B. Stevens.；Nicene and Post-Nicene Fathers, First Series；Vol. 11.；Edited by Philip Schaff.；Buffalo, NY: Christian Literature Publishing Co.,；1889.；<http://www.newadvent.org/fathers/210122.htm>.}

% structure: p0001 source=h1 target=section reason=page-title

\section{论宗徒大事录讲道词 第23篇}

\BibleRef{宗 10:23-24}\par

\BibleQuote{于是请他们进去，款待了他们。次日，伯多禄起身与他们一同去了，还有几个约培的弟兄陪伴着他。又次日，他们进了凯撒勒雅。科尔乃略等候他们，并且已经召集了他的亲戚和挚友。}\par

\Quote{他}叫他们进来，招待他们住宿。好，首先他以友善的态度对待这些人，在他们旅途劳顿之后，使他们感到自在；\Quote{第二天}，便与他们一同出发。也有一些人陪着他：这也是上智的安排，为叫他们日后在伯多禄需要为自己辩护时，能作见证。\Quote{科尔乃略等着他们，并已召集了他的亲属和密友。}这是朋友的本分，这是虔诚人的本分：面对如此福分，他关心让他的近亲都能分享一切。当然，他的\Quote{密友}，就是那些他素来完全信任的人；因着如此重大的利害关系，他不敢将此事托付给别人。依我之见，正是科尔乃略本人把朋友和亲属都引向了更好的心志。\Quote{当伯多禄进来时，科尔乃略迎上前去，俯伏在他脚前，朝拜了他。}\BibleRef{宗 10:25}这既是为教训众人，也是为感谢天主，并显示自己的谦卑：由此显明，虽然他受了命令，但他本人却怀有极大的虔诚。那么伯多禄做了什么？\Quote{伯多禄却扶起他来，说：你起来，我自己也是一个人。}\BibleRef{宗 10:26}你是否注意到，在一切之先，宗徒们教导他们这一点：不要对他们存有高看？\Quote{伯多禄和他交谈着进去，见有许多人聚集。便对他们说：你们知道，犹太人是不允许与外邦人结交或来往的；但天主指示了我，不要称任何人为凡俗或不洁的。}(27-28节)你看，他立刻谈到天主的仁慈，并指出天主向他们显示了大恩。你也注意到，他说出大事的同时，也谦卑地说话。因为他没有说：我们这些不屑与任何人结交的人，来到了你们这里；而是说什么？\Quote{你们知道}——是天主命令了这事——\Quote{与外邦人结交或来往是违反法律的。}然后他接着说：\Quote{天主却指示了我}——他说这话，为叫无人将感谢归于他——\Quote{不要称任何人}——免得显得是讨好他——\Quote{不要称任何人}，他说，\Quote{为凡俗或不洁的。}\BibleRef{宗 10:29}\Quote{因此}——为叫他们不认为这事是他违犯了法律，也不叫科尔乃略以为因他身居要职而屈从，而是叫他们将一切归于天主——\Quote{因此我一被请，就毫无异议地来了：}(虽然)不仅结交，甚至来往都是不被允许的。\Quote{所以我请问：你们为什么请我来？}伯多禄已从士兵那里听到了整个事情，但他希望他们先承认，并使他们顺服于信德。那么科尔乃略做什么？他没有说：难道士兵没有告诉你们吗？但再看他是如何谦卑地说话。因为他说：\Quote{从第四天起，我禁食直到此时；在第九时辰，我在家中祈祷，忽然有一个人站在我面前，穿着光亮的衣服，说：科尔乃略，你的祈祷蒙垂听，你的施舍已在天主面前成为记念。在第九时辰，}他说，\Quote{我正在祈祷。}(30-31节)依我看来，这个人还为自己设定了更严格生活的固定时间，并在某些日子。这就是为什么他说：\Quote{从第四天。}请看祈祷是多么伟大的事！当他进德日深时，天使就显现给他。\Quote{从第四天：}即周中第四天；不是\Quote{四天前。}因为\Quote{第二天伯多禄同他们出发，又过了一天，他们进了凯撒勒雅：}这是一天；所派之人来到约培的那一天也是一天；第三天天使显现：所以从科尔乃略祈祷那天以后，共有两天。\Quote{忽然有一个人站在我面前，穿着光亮的衣服：}他没有说天使，他是多么谦虚：\Quote{说：科尔乃略，你的祈祷蒙垂听，你的施舍在天主面前已蒙纪念。所以你当派人往约培去，请那称为伯多禄的西满来；他住在海边一个皮匠西满的家里；他来时，必会对你说话。于是我立刻打发人去请你；你来了真好。现在我们都在天主面前，要听天主所命令你的一切话。}\BibleRef{宗 10:31}\Emph{(b)}请看怎样的信德，怎样的虔诚！他知道伯多禄所说的\Quote{天主指示了我}不是人的话。于是那人说：我们在场，要听天主所命令你的一切话。\Emph{(a)}因此，伯多禄问\Quote{你们为什么请我来？}正是为了让他说出这些话。\Emph{(d)}\Quote{伯多禄就开口说：我真看出天主是不偏待人的；但在各邦国中，那敬畏祂而行义的人，都为祂所悦纳。}(34-35节)就是说，无论是未受割礼的或受割礼的。\Emph{(c)}保禄也声明这事说：\Quote{因为天主不偏待人。}\BibleRef{罗 2:11}\Emph{(e)}那么怎样呢？有人会问：在波斯的那个人也为祂所悦纳吗？如果他是配得的，在这方面他就是悦纳的，即被赐予得以被带到信德面前（\OriginalTerm{被赋予信仰}{\GreekText{τῷ} \GreekText{καταξιωθῆναι} \GreekText{τῆς} \GreekText{πίστεως}}）。祂没有忽略埃塞俄比亚的太监。\Quote{那么那些虔诚却被忽略的人，该怎么说呢？}并不存在任何被忽略的人。但他所说的是这意思：天主不拒绝任何人。\Quote{在各邦国中，那敬畏祂而行义的：}(行义)他指一切德行。请注意，他如何消除他内心的骄傲。为叫犹太人似乎不处于被弃的状态，他加上：\Quote{祂藉耶稣基督——祂是万有的主——传给以色列子民和平的福音。}\BibleRef{宗 10:36}他说这话也是为在场的犹太人，为叫他们也信服；这也是他促使科尔乃略说话的原因。\Quote{祂}，他说，\Quote{是万有的主。}但请注意他在开头说：\Quote{祂传给以色列子民的话}；他给他们优先权。然后他引证这些外邦人自己为证人：\Quote{你们知道}，他说，\Quote{这事传遍了全犹太，从加里肋亚开始}——然后他也以此证实——\Quote{在若翰宣传洗礼之后}\BibleRef{宗 10:37}——\Quote{就是纳匝肋的耶稣，天主怎样以圣神和德能傅了祂。}\BibleRef{宗 10:38}他不是说你们认识耶稣，因为他们不认识祂，而是说祂所行的事：\Quote{祂周游各处，施恩行善，治好所有受魔鬼压制的人；}由此表明许多失去知觉或瘫痪的肢体是魔鬼的工作，是魔鬼对身体造成的扭曲；正如基督所说。\Quote{因为天主与祂同在。}又是谦卑的言辞。\Quote{我们都是祂在犹太地区和耶路撒冷所行一切事的见证人}\BibleRef{宗 10:39}：\Quote{我们}，他说，和你们。然后是受难，以及他们不信的理由：\Quote{他们竟把祂悬在木头上，杀了。第三天，天主使祂复活，并显现出来；不是显现给众百姓，而是显现给天主预先拣选的见证人，就是我们这些在祂从死者中复活后，与祂同吃同喝的人。}(40-41节)这是复活的证明。\Quote{祂吩咐我们向百姓宣讲，并证明祂就是天主所立的、审判活人死人的那一位。}\BibleRef{宗 10:42}这是伟大的。然后他引用先知们的见证：\Quote{一切先知都为祂作证：凡信祂的人，赖祂的名字，都必获得罪赦。}\BibleRef{宗 10:43}这是对将要发生之事的证明；这就是他在这里引用先知的原因。\par

但是，让我们重新审视与\Person{科尔乃略}有关的事情。（复述。）经上记载，他派人到\Place{约培}去请\Person{伯多禄}。\Quote{他正在等候他，}等等；请看他是多么深信\Person{伯多禄}一定会来：（\Emph{b}）\Quote{就俯伏在他脚前，敬拜他。}（第24、25节。）（\Emph{a}）请注意，在各方面都显明他是多么配得！同样，那里的太监请求\Person{斐理伯}上车同坐\BibleRef{宗 8:31}，虽然不知道他是谁，除了先知所给的引介（\OriginalTerm{宣告}{\GreekText{ἐπαγγελίας}}）之外，别无其他。但在这里，\Person{科尔乃略}俯伏在他的脚前。（\Emph{c}）\Quote{起来，我自己也是一个人。}\BibleRef{宗 10:26}请注意他的话处处毫无谄媚，充满谦逊。\Quote{与他交谈，他进去了。}（\Emph{a}）\BibleRef{宗 10:27}谈论什么？我想是说这些话：\Quote{我自己也是一个人。}（\Emph{e}）你注意到（\Person{伯多禄}）谦和的气质吗？他自己也表明他的到来是天主的作为：\Quote{你们知道，犹太人是不可以的，}等等。\BibleRef{宗 10:28}为什么他没有提到那麻布？请注意\Person{伯多禄}毫无虚荣：但他确实提到他是天主派遣的；至于他被派遣的方式，他目前没有说；当需要出现时，鉴于他已经说过：\Quote{你们知道，犹太人是不可以同外邦人往来或接近的，}他只是补充说：\Quote{但天主已经指示我，}等等。这里毫无虚荣之处。\Quote{你们众人，}他说，\Quote{都知道。}他让他们的知识为他作保。但\Person{科尔乃略}说：\Quote{我们现在都在天主面前，要听主所吩咐你的一切话。}\BibleRef{宗 10:33}不是在人面前，而是\Quote{在天主面前。}这就是应当听从天主仆人的方式。你看到他那觉醒的心灵吗？你看到他多么配得这一切吗？\Quote{\Person{伯多禄}，}经上记载，\Quote{就开口说：我真正明白了，天主是不偏待人的。}\BibleRef{宗 10:34}他说这话也是为了向在场的犹太人表白自己。因为即将要把圣言交给这些（外邦人），他先以此作为辩护。那么，难道祂以前是\Quote{偏待人}的吗？绝不！因为以前也同样如此：\Quote{凡敬畏祂、行义的人，}如他所说，\Quote{都为祂所悦纳。}正如\Person{保禄}所说的：\Quote{原来，凡没有法律的异民，顺着本性去行法律上的事，}\BibleRef{罗 2:14}\Quote{敬畏天主、行义的人，}他兼顾了教义和生活方式：\Quote{为祂所悦纳；}因为，若祂没有忽略贤士、埃塞俄比亚人、强盗、妓女，更何况那些行义且愿意的人，祂无论如何也不会忽略。\Quote{那么，你对此怎么说，有一些\OriginalTerm{温和的人}{\GreekText{ἐπιεικεῖς}}，性情柔和，却仍不愿相信？}（见上文，第149页，注2。）看，你自己说出了原因：他们不愿意。但他这里所说的温和人不是这种，而是\Quote{行义的人：}即，在所有方面都品德高尚、无可指责的人，当他拥有应有的敬畏天主之心时。但一个人是否如此，只有天主知道。请看这人如何被悦纳：请看他一听到就立即信服。\Quote{是的，现在也一样，}你说，\Quote{每个人都会信服，无论他是谁。}但现在的神迹远比那些更大、更奇妙。——接着，\Person{伯多禄}开始教导，并为犹太人保留了出身特权。\Quote{祂向以色列子民所传的话，}他说，\Quote{宣讲和平，}\BibleRef{宗 10:36}不带来审判。祂也被差遣到犹太人那里：然而尽管如此，祂并没有宽免他们。\Quote{藉\Person{耶稣基督}宣讲和平。祂是万有的主。}他首先谈论祂是主，并以极其崇高的言辞，因为他要面对一个超乎寻常崇高且热切接纳一切的心灵。然后他证明祂如何是万有的主，从祂所成就的事：\Quote{全犹太境内。因为你们知道，}他说，\Quote{那话传遍了全犹太：}而奇妙的部分是，\Quote{从\Place{加里肋亚}开始：\Person{若翰}宣讲洗礼以后。}\BibleRef{宗 10:37}他先说祂的成功，然后又说关于祂：\Quote{\Place{纳匝肋}人\Person{耶稣}。}哎呀，这出生地是何等绊脚石！\Quote{天主怎样以圣神和德能傅了祂。}\BibleRef{宗 10:38}接着又证明——这如何显明？——从祂所行的善事。\Quote{祂周流四方，行善，医治所有被魔鬼压制的人：}以及胜过魔鬼所显示的大能；原因是，\Quote{因为天主与祂同在。}因此犹太人也这样说：\Quote{我们知道祢是由天主而来的导师，因为除非天主与他同在，谁也不能行这些神迹。}\BibleRef{若 3:2}然后，当他表明祂是由天主派遣后，他接着谈到祂被杀害：使你不会想象什么荒谬之事。你看到他们如何毫不隐瞒十字架，甚至还连同其他情形一起说明了方式吗？\Quote{他们，}经上记载，\Quote{把祂挂在木头上杀了。并且使祂，}又补充说，\Quote{显现出来，不是给众百姓，而是给天主预定见证的人，就是我们。}然而，是（\Person{基督}）祂自己拣选了他们；但这也归于天主。\Quote{这预定，}他说，\Quote{就是我们，在祂从死者中复活后，与祂一同吃喝的人。}（第39、41节。）请看他从哪里取得复活的保证。为什么复活后祂没有行神迹，只吃喝呢？因为复活本身就是一个大奇迹，而其中最大的迹象就是吃喝。\Quote{祂命令我们，}他说——以一种令人警觉的方式——使他们无法以无知为借口：并且他没有说，\Quote{祂是天主子，}而是说，最令他们警觉的，\Quote{祂就是天主所立的，审判活人死人的那一位。}\BibleRef{宗 10:42}\Quote{所有先知都为祂作证，}等等。\BibleRef{宗 10:43}当他以恐惧激动他们后，便引入赦免，不是出于他自己，而是出于先知。可怕的话出于他，温和的话出于先知。\par

所有你们这些领受了这赦免的人，\newline{}所有你们这些得以蒙恩抵达信德的人，\newline{}我恳求你们，学习这恩赐的伟大，并努力不要对你们的恩主无礼。\newline{}因为我们获得了赦免，不是为了让我们变得更坏，而是为了使我们远为更好、更卓越。\newline{}谁也不要说天主是我们作恶的原因，因为祂没有惩罚，没有报复。\newline{}如果（如人们所说）一个统治者抓住了杀人犯，却释放了他，那么，难道他不被视为后来所犯谋杀的原因吗？\newline{}看哪，我们如何使天主暴露于恶人的口舌之下。\newline{}因为他们什么不说，什么不出口呢？他们说：\Quote{（天主）自己}，他们说，\Quote{允许了他们；因为祂本该按他们应得的惩罚他们，而不是尊敬他们，加冕他们，接纳他们到最优先的特权中，而是要惩罚他们，报复他们：但祂，非但不这样做，反而尊敬他们，使他们成了现在这个样子。}\newline{}我恳求并祈求你们，千万不要让任何人就我们而言说出这样的话。\newline{}纵然被埋葬一万次，也比天主因我们而被如此议论好得多。\newline{}我们读到，犹太人曾对（基督）自己说：\BibleQuote{你拆毁圣殿，三日内重建起来，从十字架上下来吧！}\BibleRef{玛 27:40}:并且又说：\BibleQuote{你如果是天主子：}但是，这里的责难比那些更严重，那就是因我们祂被称为邪恶的导师！\newline{}让我们使人们说截然相反的话，通过使我们的言行配得上那召唤我们的祂，并（以配得的方式）接近义子的圣洗。\newline{}因为圣洗（\OriginalTerm{光照（指圣洗）}{\GreekText{φωτίσματου}}）的能力确实伟大：它使领受这恩赐的人完全不同于从前；\newline{}它不让人们仅仅是常人（仅此而已）。\newline{}你使那外邦人（\GreekText{τὸν}\par

\SourceRecord{Translated by J. Walker, J. Sheppard and H. Browne, and revised by George B. Stevens.；Nicene and Post-Nicene Fathers, First Series；Vol. 11.；Edited by Philip Schaff.；Buffalo, NY: Christian Literature Publishing Co.,；1889.；<http://www.newadvent.org/fathers/210123.htm>.}

% structure: p0001 source=h1 target=section reason=page-title

\section{宗徒大事录讲道集第24讲}

\BibleRef{宗 10:44, 46}\par

\BibleQuote{当伯多禄还在说话的时候，圣神降在所有听道的人身上。那些受割损而信的人，和伯多禄一同来的，都惊讶，因为圣神的恩赐也倾注在外邦人身上。他们听见他们说方言，并称扬天主。}\par

请注意天主眷顾的安排。祂没有让讲道结束，也没有让圣洗在伯多禄的命令下举行，而是当祂显明了他们的心态多么令人钦佩，开始了教导的工作，并且他们相信圣洗肯定是罪的赦免之后，圣神立即降临在他们身上。这是天主如此安排，以给伯多禄提供一个强有力的成义理由。而且不只是圣神降临在他们身上，而且\BibleQuote{他们说了方言}：这正是令那些聚集的人惊讶的事。他们完全不喜欢这事，因此全是出于天主；至于伯多禄，几乎可以说，他只是在场（与他们一起）学习一个功课：他们必须着手处理外邦人，并且他们自己就是必须做这事的人。因为尽管发生了所有这些大事，在凯撒勒雅和耶路撒冷仍然对此有质疑，如果这些（迹象）没有随着事情的进展逐步出现，那将会怎样？因此，这件事被推向某种极端。伯多禄抓住他的优势，看他以此提出的理由：\BibleQuote{谁能在禁止用水，不让这些人受洗呢？他们和我们一样领受了圣神。}\BibleRef{宗 10:47}请注意他引出的问题；他一直努力引出这一点。他是如此（完全）抱有这种想法！他问：\BibleQuote{禁止用水？}这几乎可以说是一个人\OriginalTerm{乘胜追击}{\GreekText{ἐπεμβαίνοντος}}的语言，对抗那些会禁止、会说这事不应发生的人。他说，整件事已经完成，是事情中最本质的部分，即我们受洗的圣洗。\BibleQuote{他命令以耶稣基督之名给他们施洗。}\BibleRef{宗 10:48}在他为自己辩白之后，然后，而不是之前，他命令给他们施洗：用事实本身教导他们。犹太人是如此厌恶这事！因此他先为自己辩白，尽管事实本身在大声呼喊，然后才下命令。\BibleQuote{于是他们求他}——他们这样求是理所当然的——\BibleQuote{逗留几天}：从此他勇敢地逗留了几天。\par

\BibleQuote{宗徒和在犹太的弟兄听说外邦人也领受了天主的圣言。及至伯多禄上了耶路撒冷，那些受过割礼的人与他争辩，说：\InnerQuote{你进了未受割礼之人的家，且同他们吃了饭。}} \BibleRef{宗 11:1} 在如此伟大的事以后，\Quote{那些受过割礼的人争辩}，并不是宗徒；绝对不可！意思是他们大感不快。且看他们指控什么。他们不说：\Quote{你为何宣讲？}却说：\Quote{你为何同他们吃饭？}但伯多禄并没有止于这个冷淡的异议——实属冷淡——而是立足于\OriginalTerm{站立}{\GreekText{ἵσταται}}那伟大的论据：如果他们已被赐予圣神本身，谁还能拒绝给他们施行圣洗呢？但是，为何在撒玛黎雅人的事例中没有发生这事，反而在他们受洗前后都毫无争议，他们并未因之而恼怒，反倒一听到这话就为此目的派遣宗徒？\BibleRef{宗 8:14} 不错，但在当前事例中，他们所抱怨的也不是这事；因为他们知道是出于天主的恩宠：他们所说的是：\Quote{你为什么同他们吃饭？}此外，撒玛黎雅人与外邦人的差别并不那么大。况且，这是如此安排（作为天主计划的一部分），使他如此受指控：为的是他们可以学习；因为伯多禄若没有某种理由，就不会叙述那异象。但请看他是如何摆脱一切得意和虚荣。经上说：\BibleQuote{伯多禄便按次序从头述说这事，说：\InnerQuote{我在约培城祈祷的时候}}他没有说为何祈祷，也没有说在什么情况：\BibleQuote{在出神中我看见一个异象，有一件东西像一块大布，从天上缒下，由四角兜着，一直来到我面前。我定睛观看，见地上的四足兽、野兽、爬虫和天上的飞禽。我又听见有声音向我说：\InnerQuote{伯多禄，起来，宰了吃。}}（4-5节，6-7节）这本身足以说服我——我见了那麻布；但更有声音加添。\BibleQuote{我却说：\InnerQuote{主，万万不可！因为凡俗和不洁之物从来没有进过我的口。}} \BibleRef{宗 11:8} 你留意吗？\Quote{我尽了我的本分，}他说，\Quote{我说我从未吃过任何凡俗或不洁之物：}这是针对他们所说的\Quote{你进了他们的家，且同他们吃了饭。}但他对科尔乃略并没有提起这事，因为没有必要对他提及。\BibleQuote{然而声音第二次从天上回答我说：\InnerQuote{天主所洁净的，你不可称为凡俗。}这样一连三次，一切又都收回天上去了。}（9-10节）要紧的是那些（在凯撒勒雅发生的事）；但通过这些他为他们铺路。请注意他是如何为自己辩证（用理由），而没有运用他作为老师的权威。因为他表达得越温和，就越使他们易于接受。\Quote{从来没有，}他说，\Quote{凡俗或不洁之物进过我的口。——而且，看——这也是他辩护的一部分——就有三个人来到我所在的房子，是从凯撒勒雅差到我这里来的。圣神吩咐我同他们去，一点不要疑惑。}（11-12节）你留意吗？立法的权柄属于圣神！\Quote{这些弟兄也与我同去——}再谦卑不过了，当他请弟兄们作证时！——\Quote{这六个人，我们进了那人的家。他告诉我们，他看见一位天使站在他家里，说：\InnerQuote{你打发人到约培去，请那称呼伯多禄的西满来；他有话对你说，叫你和你的全家得救。}}（13-14节）他没有提及天使对科尔乃略所说的话：你的祈祷和你的赒济已上达天主面前，作为纪念，为免引起反感；但他说什么？\Quote{他有话对你说，叫你和你的全家得救。}这添加得很有理由。他也对那人的\OriginalTerm{温和}{\GreekText{ἐπιεικές}}只字未提。\Quote{圣神，}他可能会说，\Quote{派遣了我，天主命令了，一方面通过天使召唤我，另一方面勉励我，并解除我对这些事的疑惑，我还有什么可做的呢？}然而他并未提及这些事，却把重点放在最后发生的事上，这本身就无可辩驳的论证。\Quote{我一开讲，}等等。\BibleRef{宗 11:15} 那么为什么单单这事没有发生？天主出于\OriginalTerm{充盛}{\GreekText{ἐκ} \GreekText{περιουσίας}}行这事，为要显明开始也不是来自宗徒。但若他主动出发，没有这些事发生，他们就会大大受伤害：所以从一开始他就使他们对他怀有好感*：对他们说：\Quote{他们领受了圣神，和我们一样。}且不满足于此，他还提醒他们主的言语：\BibleQuote{我又想起主的话，他说：\InnerQuote{若翰用水施洗，但你们要受圣神的洗。}} \BibleRef{宗 11:16} 他的意思是，没有发生新事，只是主预言过的。\Quote{但无需施洗吗？}（参看第158页）但洗礼已经完成了。他没有说：\Quote{我命他们受了洗}；而是说什么？\BibleQuote{既然天主赐给他们同样的恩赐，正如赐给我们这些信主耶稣基督的人一样，我是谁，能阻挡天主呢？} \BibleRef{宗 11:17} 他显明自己什么也没做：因为我们所得的，他说，那些人也领受了。为更有效地堵住他们的口，他说\Quote{同样的恩赐。}你察觉他如何不让他们少得吗？当他们相信时，他说，天主赐给他们同样的恩赐，正如赐给我们信主的人，并亲自洁净他们。他没有说\Quote{给你们}，而是\Emph{给我们}。当\Emph{我们}称他们为同享者（与我们同享）时，你为何感到委屈？\BibleQuote{他们听见这话，就不再言语，只光荣天主说：\InnerQuote{那么天主也赐给外邦人悔改而得生命了。}} \BibleRef{宗 11:18} 你留意吗？这一切都来自伯多禄的讲话，他巧妙地叙述了事实？他们光荣天主，因为天主也赐给了他们\OriginalTerm{自己}{\GreekText{καὶ} \GreekText{αὐτοῖς}}悔改；他们被这些话谦卑了。从此为外邦人敞开了信仰之门。但若你乐意，我们再来回顾一下所说的话。\par

\BibleQuote{\Person{伯多禄}还在说话的时候}等等。（概述）他没有说\Person{伯多禄}惊奇，而是说\BibleQuote{那些受割礼的人}：因为他知道将要发生什么。然而他们本应为此惊奇：他们自己如何相信了。当他们听说他们相信了时，并不惊奇，但当天主赐给他们圣神时。然后\BibleQuote{\Person{伯多禄}回答说}等等。\BibleRef{宗 10:47}因此他说：\BibleQuote{天主指示了我，不可把任何人称为俗或不洁的。}\BibleRef{宗 10:28}他从起初就知道这事，并且事先计划了他的讲论（着眼于此事）。外邦人？此后还有什么外邦人？真理既已来临，他们不再是外邦人了。他说，若在圣洗圣事之前他们领受了圣神，这并不稀奇；在我们自己身上也发生了同样的事。\Person{伯多禄}表明，他们受洗的方式也不像其余的人，而是以一种好得多的方式。这就是事情以这种方式发生的原因，好让他们无话可说，并且即使这样，也认为他们与自己相等。\BibleQuote{他们便请求他，}经上说，\BibleQuote{住下几天。}\BibleRef{宗 10:48}\BibleQuote{宗徒和弟兄们等等。那些受割礼的人与他争辩。}\EditorialNote{宗 11:1-2}你注意到他们对他并不友善吗？说\BibleQuote{你进了未受割礼的人那里，且同他们一起吃饭。}\BibleRef{宗 11:3}你注意到他们对法律的热情吗？不是\Person{伯多禄}的权威使他们羞愧，不是已发生的奇迹，也不是所取得的成就——这是何等的事，外邦人\BibleQuote{领受了圣言}——而是他们为那些琐碎之事争论。因为即使那些（奇迹）都没有发生，难道这成功本身还不够吗？但\Person{伯多禄}并不这样为自己辩护：因为他是有智慧的，或者不如说，不是他的智慧，而是圣神在说话。并且通过他的辩护内容，他表明在任何一点上他都非主动，而在每一点上都是天主，他把全部责任都归给天主。\Quote{那神魂超拔，}他说——是他使我陷入其中，因为\BibleQuote{我在\Place{约培}}等等：那器皿——是他显示的；我反对了；他又说话，即使那时我也没听；圣神命令我去，即使那时，虽然我去了，我也没有跑；我告诉他们说天主派遣了我，在此之后，即使那时我也没有施洗，而是天主完成了全部。天主给他们施洗，不是我。他并没有说，那时再加添水难道不对吗？而是暗示什么也不缺，\BibleQuote{我是谁，竟能阻止天主？}这是何等的辩护啊！因为他不是说，既知道这些事，就沉默吧；而是什么？他抵挡他们的攻击，对他们的控告，他申辩说——\Quote{我是谁，能阻止天主？}我不可能阻止——这的确是有力的陈述，足以使他们羞愧。因此他们终于害怕了，\BibleQuote{他们便安静下来，并光荣了天主。}\par

同样，我们也应该为邻人所遇到的好事光荣天主，只是不要像其他新领洗者那样，当他们看到别人领受圣洗后立即离世，就感到受辱。即使所有人都得救，光荣天主也是对的；至于你，如果你愿意，你已领受了比他们更大的恩赐：我不是就圣洗而言，因为圣洗的恩赐对你和他是一样的，而是就你获得了一段赢得殊荣的时间而言。另一个人穿上了白衣，却未被允许在游行中展示它；而天主却给了你充分的机会，以正确的方式使用你的武器，从而验证它们。另一个人离开，只获得了他信德的赏报；而你站在赛场上，既能因你的善工获得丰厚的赏报，又能让自己比他更光荣，如同太阳比最小的星辰，如同将军，甚至皇帝本人比最低等的士兵。那么，责怪你自己吧，或者更确切地说，不是责怪，而是改正：因为仅仅责怪自己是不够的；你还有能力重新角力。你跌倒了吗？你受了重伤吗？站起来，恢复自己：你仍在比赛中，\OriginalTerm{赛场}{\GreekText{θέατρον}}尚未散场。你没有看见多少在摔跤中跌倒的人后来重新投入战斗吗？只是不要故意跌倒。你认为他因离世而幸福吗？不如认为你自己幸福。他解除了罪债吗？但你，若你愿意，不仅会洗去你的罪恶，还会获得善工的成就，这是他所不能的。我们有能力恢复自己。悔改的\OriginalTerm{药物}{\GreekText{φάρμακα}}何其伟大：不要有人对自己绝望。真正该绝望的人，是那对自己绝望的人；那人不再有救恩，也不再有任何希望。可怕的不在于陷入罪恶的深渊，而在于跌倒后躺卧不起；可怕的不在于陷入这样的境地，而在于轻视它，这是不虔敬的。那本该使你认真的事，请问，竟使你鲁莽吗？受了这么多伤，你又后退了吗？灵魂没有不可治愈的创伤；身体有许多，但灵魂没有：然而对于身体的创伤我们不停努力医治，对于灵魂的创伤我们却漫不经心。你没有看见那（十字架上的）强盗在多么短的时间内成就了（他的救恩）吗？你没有看见殉道者在多么短的时间内完成了整个事业吗？\Quote{但如今没有殉道的机会了。} 确实，但还有竞赛可进行，正如我常对你们说的，如果我们有这样的心意。\Quote{因为凡愿意，}宗徒说，\Quote{在基督耶稣内虔诚生活的，都必遭受迫害。}\BibleRef{弟后 3:12} 虔诚生活的人总是遭受迫害，如果不是来自人，至少来自恶神，那是更严重的迫害。是的，正是由于安逸和舒适，那些不警醒的人首先会遭受这种迫害。或者你认为安逸生活是一种微不足道的迫害吗？这比一切都更严重，这比迫害更糟。因为如同流动的恶疾，安逸使灵魂\OriginalTerm{松弛}{\GreekText{χαυνοἵ}}；如同夏季和冬季，迫害和安逸也是如此。但为向你们证明这是更糟的迫害，请听：它在灵魂中引入睡眠，过度的呵欠和困倦，在各方面激起情欲，武装骄傲，武装享乐，武装愤怒、嫉妒、虚荣、妒忌。但在迫害时期，这些都不能扰乱；而恐惧进入，猛烈地挥鞭，就像人对待吠叫的狗一样，不让这些情欲中的任何一个试图发出声音。谁能在迫害时期沉溺于虚荣？谁能在享乐中生活？没有一人。但是有许多战栗和恐惧，造成极大的平静，使港口安静下来，使灵魂充满敬畏。我从我们的祖先那里听说（但愿在我们这个时代天主不许它发生，因为我们奉命不要寻求试探），在从前的迫害中，可以看到真正的基督徒。没有人关心钱财，没有人关心妻子，没有人关心孩子、家庭或国家；所有人唯一的关心是保全自己的性命（或灵魂）。他们躲藏起来，有些在坟墓和墓穴里，有些在旷野：是的，就连娇嫩柔弱的妇女也在那里，与持续的饥饿搏斗。\par

那么，请想一想：一个躺在棺材旁\OriginalTerm{在棺材旁}{\GreekText{παρὰ} \GreekText{λάρνακι}}、等候婢女送饭、战战兢兢唯恐被捕、惊恐如卧火窑中的妇人，心中可曾有过一丝对奢华精致生活的向往？她甚至可知道世上曾有精致生活这回事，有没有衣饰、装饰品这些东西\OriginalTerm{什么装饰品这回事}{\GreekText{ὅτι} \GreekText{κόσμος} \GreekText{ὅλως} \GreekText{ἐστίν}}？你们看见吗？现在是迫害，我们的情欲如同野兽从四面扑来。现在是严峻的迫害，不但在这一点上，尤其在于它甚至不被认为是迫害。因为这种迫害还有一个害处：明明是战争，却被以为是和平，以致我们根本不武装自己，甚至不奋起；无人恐惧，无人战栗。但如果你们不相信我，请问问外邦人、迫害者：什么时候基督徒的行为更严谨？什么时候他们更受考验？那时他们在人数上虽少，却在德行上富足。请说，多量的干草有何益处，倘若可得的是宝石？数量不在于数目总和，而在于经得起考验的价值。厄里亚只有一人，然而全世界的价值抵不上他一人。世界虽由万民组成，但若连那一个人都比不上，便算不得万民。\Quote{敬畏上主的人，胜过千万不义的人}，因为那千万人尚未达到那一个人的境界。\Quote{不要羡慕不虔敬者的儿女众多，}\BibleRef{德 16:1}他们带给天主的亵渎，比不是基督徒的人更大。我要众多何用？不过是给火添加燃料。在身体上也能看到这点：与健康相配的适量食物，胜过带来损伤的肥牛犊。前者是食物，后者是疾病。在战争中也可见到：十个精勇的健儿，胜过一万个生手。后者非但无益，反而阻碍工作者。在船上也是如此：两个有经验的水手，胜过极多无经验的，后者会使船沉没。我说这话，并非以恶眼看你们的人数，而是希望你们都成为经得起考验的人，不倚靠人数。下地狱的人更多，但天国虽少，却大于地狱。以色列百姓多如海沙，但一人便拯救了他们。梅瑟仅一人，却比他们都更有效；若苏厄一人，能胜过六十万人。我们不要只研究（百姓）如何众多，而要研究他们如何出众；当这有了成果，那另外的也会随之而来。没有人一开始就建造宽敞的房子，而是先把它造得坚固牢靠，然后才宽敞；没有人打地基是为了被人讥笑。让我们先追求这个，然后才追求那个。有了这个，那个也会容易；没有这个，那个就算有了也无益。因为如果教会内有人能发光，人数也很快就会增多；但若没有这样的人，人数再多也无用。你们想，我们城中得救的人\OriginalTerm{得救的人}{\GreekText{τοὺς} \GreekText{σωζομένους}}可能有多少？我下面要说的话令人不快，但我仍要说：在这数万民众中，恐怕找不到一百个得救的人；甚至对这一点我也怀疑。请想想，青年中有何等邪恶，老年中有何等懈怠！无人视照顾自己的男孩为己任，无人因看见长者的榜样而受感动去效法。榜样已败坏，因此青年也不在品行上令人钦佩。不要对我说：\Quote{我们人数众多。}这是冷淡\OriginalTerm{冷淡的}{\GreekText{ψυχρῶν}}而又麻木之人的话。在人的事上，这样说或许有点道理；但在天主的事上，若说祂需要我们，那是绝对不能允许的。让我告诉你们，即使在前面的事上，这也是麻木\OriginalTerm{冷淡的}{\GreekText{ψυχρόν}}的言论。请听：一个人有很多仆人，如果他们都是败坏的，那将是何等的痛苦！没有仆人的人，其困苦无非是无人伺候；但如果仆人坏，其害处还包括他自身败坏，损害更大（仆人越多损害越大）。因为自己动手侍奉，远胜于与别人争斗、进行持续不断的战争。我说这些话，是希望没有人因教会的众多人数而赞赏，而是我们努力使群众经得起考验；使每个人都为自身的一份责任而热心——不仅为朋友、亲属、邻居，正如我常说的，也要吸引陌生人。例如，祈祷正在进行，他们（跪着）躺着，所有青年冷漠\OriginalTerm{冷淡的}{\GreekText{ψυχροὶ}}、漫不经心，连老人也是如此。他们不是青年，而是令人厌烦的污秽；咯咯笑、放声大笑、说话——因为我甚至听见这些事——并且互相嘲弄，跪着躺着。而你，年轻人或年长者，站在那里：如果你看见他们（这样行为），要斥责他们；如果有人不肯收敛，要更严厉地责备他；叫执事来，威胁他，尽你所能去做。如果他敢对你做什么，你必会得到众人的帮助。因为谁会那样不合理性，看见你为这样的事斥责，而他们受斥责，却不站到你一边？你离开时，必从祈祷中领受你的赏报。——在主人的家中，我们称那些见不得任何家具杂乱无章之仆人为最忠心的。请回答我：在家若看见银器被扔到门外，虽不是你的事，你会捡起来带进屋；若看见衣服被丢出原处，虽非你负责，即使你与负责者不睦，但出于对主人的善意，你不会把它放好吗？同样，在此情形下，这些也是家具的一部分。如果你看见它们散乱无序，就整理它们；向我汇报，我不怕麻烦；通知我，让我知道犯错的人。我不可能看见一切；请原谅我。看，何等邪恶弥漫全世界！我说我们如同干草（如翻腾之海）岂是无缘无故？我不是说那些人（年轻人）这样行为（我抱怨的是），进来这里的人如此昏睡麻木，以致连这种不当行为也不纠正。\par

我又看见别人在祈祷进行时站着谈话；那些更一致的人（这样做）不但祈祷时如此，甚至当司铎在降福时也是如此。啊，可怕！何时才有救恩？何时我们才能安抚天主？——士兵们去娱乐，你会看到他们全都随着舞蹈合拍，没有任何疏忽，正如刺绣和绘画中，从构图中每一部分的有序安排，立即产生极大的和谐与协调，这里也是如此：我们有一面盾牌，一个头，我们大家（共同）：若任何一点因疏忽而混乱，整体便混乱损坏，多人的良好秩序被一部的混乱破坏。而且，想到确实可怕，你们来到这里，不是为了娱乐，不是为了跳舞，却站着毫无秩序。你们不知道你们是站在天使中间吗？你们与他们一起咏唱，与他们一起唱圣歌，而你们却站着发笑？岂不奇妙，雷劈不仅没有降在那些人（如此行为者）身上，也没有降在我们身上？因为这样的行为确实该遭雷劈。皇帝在场，在检阅军队：而你们，即使在他眼前，却站着发笑，并且容忍看到别人发笑？我们还要责骂多久？抱怨多久？难道这样的人不该被视为害虫和讨厌的东西，视为被遗弃、无用的败类，充满无数祸患，被赶出教会吗？这些\OriginalTerm{在可畏的时刻}{\GreekText{ἑ} \GreekText{ν} \GreekText{ὣρᾳ} \GreekText{φρίκης}}发笑的人，何时才会停止发笑？那些在降福的瞬间说话的人，何时才会停止他们的轻浮？他们在在场的人面前没有羞耻感吗？他们不怕天主吗？我们自己的闲散念头还不够吗？我们在祈祷中四处游荡还不够吗？竟然还要插入笑声和欢闹？这里所做的事是戏院的娱乐吗？是的，但我想，是戏院造成了这一切：正是戏院使得你们大多数人如此拒绝被我们约束和改造。我们在这里建立的，在那里被推倒：不仅如此，听者自己也不免沾染其他污秽：所以情况就像一个想清洁一个地方，而上方却有喷出泥浆的泉源；无论你清除多少，更多的泥浆流进来。这里也是如此。因为当我们清洁人们的时候，他们带着污秽从戏院来到这里，又去那里，吸收更多的污秽，好像他们活着只是为了给我们添麻烦，然后回到我们这里，满载污秽，在他们的举止、动作、言语、笑声、闲散中。然后我们再次开始重新铲除，好像我们这样做只是为了清洁地送走他们后，再看到他们被污秽堵塞。因此，我郑重向你们这些健全的成员声明，这将成为你们的审判和定罪，我从今以后将你们交付于天主，若有人看见某人行为不轨，若有人看见某人说话，尤其是在礼仪的那部分，却不告发他，不使他回转。这样做比祈祷更好。放下你的祈祷，斥责他，这样你既能有益于他，自己也能受益，这样我们大家才能得救，并获享天国，因我们的主耶稣基督的恩宠和慈爱，愿光荣、权能、尊威归于祂，偕同圣父及圣神，自今至永远，及世之世。阿们。\par

\SourceRecord{Translated by J. Walker, J. Sheppard and H. Browne, and revised by George B. Stevens.；Nicene and Post-Nicene Fathers, First Series；Vol. 11.；Edited by Philip Schaff.；Buffalo, NY: Christian Literature Publishing Co.,；1889.；<http://www.newadvent.org/fathers/210124.htm>.}

% structure: p0001 source=h1 target=section reason=page-title

\section{\WorkTitle{论宗徒大事录讲道集 第二十五篇}}

\BibleRef{宗 11:19}\par

\BibleQuote{那时，因司提反的事遭遇患难而四散的人，直走到腓尼基、塞浦路斯和安提约基雅，只向犹太人宣讲圣言，不向别人宣讲。}\par

迫害成了不小的益处，因为\BibleQuote{天主使一切协助那些爱祂的人获得益处}。\BibleRef{罗 8:28} 如果他们曾特意研究如何最好地建立教会，他们也不会做出别的——他们分散了教师们。请注意宣讲扩展到了哪些地方。经文说：\BibleQuote{他们旅行直到腓尼基、塞浦路斯和安提约基雅，但只给犹太人宣讲圣言。}你看到天主眷顾何等智慧的目的，在科尔乃略的事上行了如此多事吗？这既为基督辩护，又控告犹太人。当斯德望被杀，当保禄两次遇险，当宗徒们被鞭打，那时外邦人接受了圣言，那时撒玛黎雅人也接受了。保禄也宣告：\BibleQuote{天主的圣言本来必须先讲给你们听，但因你们把它推开，断定自己不配得永生，我们转向外邦人。}\BibleRef{宗 13:46} 于是他们四处巡游，也向外邦人宣讲。\BibleQuote{但他们中有些是塞浦路斯和基勒乃的人，他们到了安提约基雅，也向希腊人讲论，宣讲主耶稣：}\BibleRef{宗 11:20} 因为很可能他们此时能说希腊话，安提约基雅也有这样的人。经文说：\BibleQuote{主的手与他们同在}，就是说，他们行了奇迹；\BibleQuote{于是有许多人相信而归向主。}\BibleRef{宗 11:21} 你看到为什么现在也需要留意奇迹（即）使他们相信吗？\BibleQuote{这事传到了耶路撒冷教会的耳中，他们就打发巴尔纳伯到安提约基雅去。}\BibleRef{宗 11:22} 为什么当这样一座城市接受圣言时，他们自己不来？因为犹太人。但他们打发巴尔纳伯。然而，即使如此，保禄去了那里，也是天主眷顾治理中不小的部分。他们对他反感，从而福音之声、天上的号角不被困在耶路撒冷，这既是自然的，也是明智的安排。你看到基督如何总把他们的恶意向转化为必需之事，并使教会受益吗？祂利用他们对这人的仇恨来建立教会。但请看这位圣人——我是指巴尔纳伯——他如何不顾自己的利益，急忙赶往塔尔索。\BibleQuote{他来了，看见天主的恩宠，非常喜乐，劝勉众人，要心存坚定，紧靠主。因为他是个好人，充满圣神和信德，于是许多人归向了主。}（23–24节）他是个非常仁慈、心地单纯且体谅的人（\OriginalTerm{体谅的}{\GreekText{συγγνωμονικ}{\GreekText{ο}}\GreekText{ς}}）。\BibleQuote{于是巴尔纳伯出发前往塔尔索，去寻找扫禄。}\BibleRef{宗 11:25} 他来到这位竞技搏斗者、将军（能统率军队）、单打冠军、狮子面前——我不知用什么词才好，无论怎样说——猎犬、杀狮者、力壮如牛、明灯、足以充满世界的口舌。\BibleQuote{找到了他，就带他到安提约基雅。}\BibleRef{宗 11:26} 确实，这就是为什么在那里他们被称为基督徒，因为保禄在那里住了这么久！\BibleQuote{他们整整一年与教会聚在一起，教导了许多人。门徒在安提约基雅最先被称为基督徒。}对那座城市而言，这是不小的赞美！这足以使它堪与众人匹敌，因为它如此长久地受益于那口舌，它最先、在众人之先，因此也正是在那里，人最先被配得上那个名号。你看到保禄给那座城市带来的益处吗？那名号像一面旗帜（\OriginalTerm{标准}{\GreekText{σημε}\GreekText{ἵον}}）将它高举到了何等高度？那里三千人、五千人相信，那里如此庞大的群众，却没有发生过这类事，他们被称为\Quote{属于那道的}；在这里他们被称为基督徒。\BibleQuote{在那几天，有先知从耶路撒冷下到安提约基雅。}\BibleRef{宗 11:27} 需要在那里也栽种施舍的果实。请看，如何\OriginalTerm{必然地}{\GreekText{ἀναγκα}\GreekText{ίως}}没有一位显要人物成为他们的教师。他们得到教师是塞浦路斯人、基勒乃人以及保禄——尽管他确实超过了（宗徒）他们自己——因为保禄也有阿纳尼雅和巴尔纳伯做老师。但在这里，这是必然的情况。\BibleQuote{其中一位名叫阿加波的站起来，藉着圣神指明天下将有大饥荒，这饥荒在克劳狄皇帝时果然发生了。}\BibleRef{宗 11:28} 经文说：\BibleQuote{藉着圣神}。为使人们不以为饥荒来临是因为基督信仰的传入、魔鬼的离去，圣神便预言了它。但这并不稀奇，因为基督实际已预言过。这不是原因，否则从一开始就应是如此；而是因为宗徒们所受的恶行——天主已长久忍耐他们；但当他们迫害宗徒时，一场大饥荒随之而来，向犹太人预示即将来临的灾祸。\Quote{如果是因他们而起，它（指饥荒）在存在时无论如何应该停止。外邦人做了什么坏事，应该分享灾祸？他们倒应被看作蒙赞许（\OriginalTerm{受赞许}{\GreekText{ε}\GreekText{ὐδοκιμ}\GreekText{ἥσαι}}），因为他们正尽本分，杀戮、惩罚、报仇、处处迫害。也要注意饥荒发生的时候：正是在外邦人从此加入教会的时候。但如果如你所说，是因为（犹太人做的）恶事，那么这些（外邦人）应该被豁免。}怎么会这样？基督预见到这反对，说：\BibleQuote{你们将有苦难。}\BibleRef{若 16:33} 正如你如果要说，他们也不该被鞭打。\BibleQuote{于是门徒决定，各按自己的力量，送救济品给住在犹太的弟兄们。}\BibleRef{宗 11:29} 请注意，饥荒如何成为他们得救的方法、施舍的机会、许多祝福的先兆。有人可以说，对你也是如此，如果你愿意，但你不愿意。但它的预言，为使他们预先准备好施舍。\BibleQuote{给住在犹太的弟兄们}；因为他们正遭受巨大困苦，但之前他们并没有遭受饥荒。\BibleQuote{他们也这样行了，托巴尔纳伯和扫禄送到长老们那里。}\BibleRef{宗 11:30} 你看见吗？他们一相信就结出果实，不仅为自己，也为远处的人？于是巴尔纳伯和扫禄被派去办理这事。关于这机会（\OriginalTerm{在此处}{\GreekText{Εντα}\GreekText{ὕθα}}），他（对迦拉达人）说：\BibleQuote{雅各伯、刻法和若望向我与巴尔纳伯伸出右手相交，只是（他们愿意）我们该纪念穷人。}\BibleRef{迦 2:9} 雅各伯那时还在世。\par

\Quote{如今那些因迫害而四散的人}等等。（重述。）你注意到吗？即使在患难中，他们没有像我们那样陷入哀叹和眼泪，反而投身于伟大而美好的工作？\Quote{一直走到腓尼基、塞浦路斯和安提约基雅}，在那里更安全地宣讲圣言。\Quote{其中有些是塞浦路斯和基勒乃人}等等。\BibleRef{宗 11:20} 他们并没有说：\Quote{我们这些基勒乃人和塞浦路斯人，竟要去攻取这座辉煌的大城！}而是依靠天主的恩宠，投身于教导的工作，而这些外邦人自己也并不轻视向他们学习任何事。请注意，天主怎样借着微小的事物成就一切：请注意宣讲怎样扩展：请注意在耶路撒冷的人，他们同样关心众人，把整个世界看作一家。\Quote{他们听说撒玛黎雅接受了圣言，} \BibleRef{宗 8:14} 就派遣宗徒到撒玛黎雅去；他们听说在安提约基雅发生的事，就派遣巴尔纳伯到安提约基雅去；他们还再次派遣了这些先知。因为路途遥远，而且当时宗徒们不宜离开那里，免得被人认为他们是逃亡者，从自己的人民那里逃离。但那时，几乎恰好是他们离开耶路撒冷的时候，当时犹太人的状况已显得无可救药，战争迫在眉睫，他们必定要灭亡：判决已成定局。因为直到保禄前往罗马，宗徒们一直留在耶路撒冷。但他们离开，并非因为害怕战争——怎么会呢？——因为他们所去的地方，正是那些将带来战争的人所在之处；而且战争是在宗徒们去世之后才爆发的。关于他们，圣经说：\Quote{忿怒已降临到他们身上直到终结。} \BibleRef{得前 2:16} 人越是微不足道，恩宠就越显赫，借着微小的事物成就伟大的结果。——\Quote{他劝勉他们要决心依靠主，因为他是一个良善的人。}（23–24节。）藉着\Quote{良善的人}，我认为他指的是一个仁慈、真诚的人，原文 \OriginalTerm{仁慈}{\GreekText{χρηστόν}}，极切渴望邻人得救的人——\Quote{因为他是一个良善的人，充满圣神和信德。要决心依靠主}（这是所说的）：带着赞美和颂扬。\Quote{有许多人归附了主：}因为这座城市如同肥沃的土地，接受了圣言，结出了许多果实。\Quote{于是巴尔纳伯往塔尔索去，}等等。\BibleRef{宗 11:25} 但他为什么要把保禄从塔尔索带到这里呢？不是没有充足的理由；因为这里有好的希望，有更大的城市，有众多的百姓。请看恩宠怎样运作一切，而不是保禄：事情借着微小的事物开始。当它变得困难时，宗徒们接手。为什么他们之前没有见到巴尔纳伯？因为他们忙于耶路撒冷的事（原文 \OriginalTerm{忙于}{\GreekText{ἠσχόληντο}}）。他们又向犹太人辩白，说外邦人正在接受（原文 \OriginalTerm{接受}{\GreekText{προσελάμβανε}}）圣言，即使没有享受到如此大的关注。将要有一个争论：因此科尔乃略的事先期解决了。那时，他们确实说：\Quote{我们往外邦人那里去，他们往受割礼的人那里去。} \BibleRef{迦 2:9} 请注意，从此以后，正是饥荒的压力引入了外邦人的共融，即通过施舍。因为他们接受了从他们那里送来的献仪。\par

\Quote{如今那些四散的人}等等。\BibleRef{宗 11:19} 不像我们这些在患难中哭泣流泪度日的人，但他们更加无所畏惧地度过时日，因为远离了那些阻碍他们的人，并置身于不惧怕犹太人的人中间：这也帮助了他们。他们来到塞浦路斯，那里有大海相隔，更大的无忧无虑：所以他们不把对人的恐惧放在心上，但他们（仍然）将律法的尊重置于首位：\Quote{他们只向犹太人说话。但在安提约基雅有些塞浦路斯和基勒乃人：}这些人，比所有其他人，最不关心犹太人：\Quote{他们向希腊人说话，宣讲主耶稣。} \BibleRef{宗 11:20} 大概是因为他们不会希伯来语，所以他们称他们为希腊人。\Quote{当}巴尔纳伯，圣经说，\Quote{来了并看到了天主的恩宠，}——不是人的努力——\Quote{他劝勉他们依靠主} \BibleRef{宗 11:23}：并且借此他使更多人归化。\Quote{有许多人归附了主。} 为什么他们不写信给保禄，却派巴尔纳伯去呢？他们还不认识这个人的德行：但这是天主的安排，使巴尔纳伯前来。既然那里有一大群人，又无人阻拦，信德自然增长，尤其是因为他们没有经历任何考验。保禄也宣讲，不再被迫逃亡。这也是很好的安排：不是他们谈论饥荒，而是先知。安提约基雅的人也没有因宗徒们没有来而不快，反而满足于他们的教师：他们都如此热衷圣言。他们没有等到饥荒来临，而是在此之前就送去了：\Quote{照各人所能的。} 请注意，在宗徒们中，其他人被委托这种信任，但这里是保禄和巴尔纳伯。因为这是上智不小的安排（原文 \OriginalTerm{安排}{\GreekText{οἰκονομία}}）。此外，这是开始，他们不应该被冒犯。\par

\Quote{每人按各自的能力，他们遣送了。}但现在无人这样做，尽管有比那更严重的饥荒。因为情况不同：所有人都共同承受灾难，而其余人富足时，穷人却在挨饿。这话也表明施予者本是穷人，因为经上说\Quote{每人按各自的能力。}这是双重的饥荒，正如富裕是双倍的：严重的饥荒，不是缺少听主的话，而是缺少靠施舍获得滋养。那时，犹大的穷人得享益处，安提约基雅的那些捐钱的人也如此；是的，后者更甚。但现在，我们和穷人都在挨饿：他们缺乏必要的生计，我们过着奢侈的生活，缺少天主的仁慈。但这是一种食物，没有比它更必要的。这不是一种会带来饱胀之恶的食物；不是一种大部分最终进入排泄物的食物（\OriginalTerm{废物}{\GreekText{ἀφεδρώνα}}）。没有什么比被这食物滋养的灵魂更美丽、更健康的了：它远高于一切疾病、瘟疫、消化不良和失调；没有人能胜过它（\OriginalTerm{胜过}{\GreekText{ἑλεῖν}}），但正如一个人的身体由金刚石造成，铁或任何东西都无法伤害它，同样当灵魂因施舍而坚固时，没有什么能胜过它。因为请问，什么能破坏它？贫穷吗？不可能，因为它藏在皇家宝库中。盗贼和闯入者吗？不，那是无人能穿透的墙。虫子吗？不，这财宝也远远超出了这种祸害。嫉妒和恶眼吗？不，这些也不能胜过。诬告和阴谋吗？不，也不能，因为这财宝像避难所一样安全。然而，若我似乎只指出施舍之益处的这些方面（免于这些恶），而不提其反面，那将是耻辱。因为事实上，它不只不受敌意侵害；它还从受益者那里获得大量祝福。正如残忍无情的人不仅得罪了的人为敌，而且未受伤害的人也分担悲伤并同声控诉；同样，行大善的人不仅得到受益者的称赞，而且未受影响的人也赞扬他们。再者，它不受恶人、强盗和闯入者的攻击——这到底是所有好处，还是这好处——它不仅不受减损，反而增长并增多？有什么比拿步高更可耻、更污秽、更不义呢？那人是不虔敬的；在无数征兆和神迹之后，他拒绝\OriginalTerm{醒悟}{\GreekText{ἀνενεγκεῖν}}，却把天主的仆人投入火窑；但做了这些事之后，他敬拜了。那么先知说了什么？\BibleQuote{因此，}他说，\BibleQuote{大王，愿我的劝言蒙你悦纳：以施舍\OriginalTerm{赎}{\GreekText{λύτρωσαι}}你的罪过，以怜悯穷人赎你的不义；或许你的过犯可得赦免。} \BibleRef{达 3:27}这样说时，他并非怀疑，而是满怀信心，但想要使他更加畏惧，并使他更迫切地做这些事。因为如果他说得毫无疑义，王就会更加怠惰：正如我们身上那样，当我们最劝某人（我们想说服的人）时，他们对我说，\Quote{劝勉这样的人，}却不加，\Quote{他必定会听，}而只说，\Quote{他或许会听：}因为留下疑问，恐惧就更大，更加催迫他。这就是先知不把事情确定给他的原因。你怎么说？如此大的不敬也能得到赦免？是的。没有施舍不能洁净的罪，没有施舍不能熄灭的罪；一切罪都在施舍之下：它是适合每个伤口的药。有什么比税吏更坏呢？他职业的\OriginalTerm{基础}{\GreekText{ὑπόθεσις}}完全是不义：然而匝凯洗净了这一切（罪）。注意基督自己如何显明这事，通过小心拥有一个钱囊并携带放入其中的捐献。保禄也说，\BibleQuote{只要我们记得穷人} \BibleRef{迦 2:10}；并且圣经处处有很多关于这事的话。\BibleQuote{赎价}，它说，\BibleQuote{人的生命的赎价是他自己的财富} \BibleRef{箴 13:8}；这是有道理的：因为（基督）说，\BibleQuote{你若愿意成全，就去变卖你所有的，施舍给穷人，然后来跟随我。} \BibleRef{玛 19:21}这可以说是成全的一部分。但施舍不仅可以用钱财，也可以用行动。例如：一个人可以友善地站在（\OriginalTerm{站旁}{\GreekText{προστῆναι}}）另一个人身边（帮助并保卫他），可以伸手相助；行为所给予的帮助（\OriginalTerm{保护}{\GreekText{προστασία}}）往往比钱财更有益。让我们实行各种施舍。你能用钱财施舍吗？不要懈怠。你能用善行施舍吗？不要说，因为我没钱，这就算不了什么。这是非常重要的一点：看做如同你捐了黄金。你能用关心（\OriginalTerm{照料}{\GreekText{θεραπεία}}）施舍吗？也这样做。例如，如果你是医生，拿出你的医术：因为这也是大事。你能用建议施舍吗？这（服务）比一切都大；这（施舍）比一切都好，或者更甚，因它的收益更大。因为这样做，你消除的不是饥饿，而是惨死。\EditorialNote{ch. 3:6; 6:4}\newline{}宗徒们大量拥有这样的施舍：因此他们把钱财的分发交给后继者，自己却展示出言语所显的（仁慈）。或者，你认为对一个失丧、被弃、处于最危险境地、被\OriginalTerm{热病}{\GreekText{πύρωσις}}烧着的灵魂，能够除去它的疾病，是微小的施舍吗？例如，你看见一个人被爱钱所困？怜悯他。他濒临窒息？熄灭他的火。\Quote{若他不肯听劝呢？}你要尽本分，不要懈怠。你看见他受捆绑？——因为财富确实是捆绑。\BibleRef{玛 25:35}到他那里去，探望他，安慰他，设法解开他的捆绑。若他拒绝，他自担责任。你看见他赤身露体、又是陌生人？——因为他确实赤身露体，对天国是陌生人。带他到你的客栈，给他穿上德行的外衣，把天上的城给他。\Quote{若我自己赤身露体呢？}你说。先也穿上自己：如果你知道自己赤身露体，那你一定知道自己需要穿上；如果你知道这是什么样的赤身露体。如今有多少妇女穿着丝绸衣服，却确实赤身露体，没有德行的外衣！让她们的丈夫给这些妇女穿上衣服。\Quote{但她们不肯接受那些衣服；她们宁愿要这些。}那也先做这事：诱导她们渴望那些衣服；给她们看她们是赤身的；对她们讲论将来的审判；回答我，在那里我们需要什么衣服？但若你容忍我，我也要给你看这赤身露体。赤身的人，在寒冷时，畏缩战栗，弯腰抱臂站着；但在夏季炎热时却不是这样。那么，如果我向你证明，你们的男女富人，穿得越多，就越赤身露体，请你不要见怪。我再问，当我们提起地狱之火和那里的痛苦时，他们难道不比那些赤身的人更战栗吗？他们不是痛苦地呻吟并自责吗？什么？当他们来到这个人或那个人面前，对他说，为我祈祷，难道不说与那些（赤身者）同样的话吗？\par

确实，在我们所能言说的一切之后，赤裸尚不明显；但在那里，它将昭然若揭。如何，以何种方式？当这些丝绸衣裳和宝石朽坏之后，唯有德行与罪恶的衣裳将显明众人：那时穷人将披上极荣光，而富人却赤裸可耻地被拖去受罚。有什么比那身披紫袍的富人更赤裸（版本作\Quote{更精致}）？有什么比拉匝禄更贫穷？那时，他们中谁说了乞丐的话？谁是富足的？试想，如果有人用大量挂毯装饰房屋，自己却赤身坐在里面，那有何益处？这些妇女的情况正是如此。诚然，灵魂的房屋——我指身体——她们以大量衣裳环绕；但屋内的主妇却赤身而坐。请借给我灵魂的眼睛，我便让你看见灵魂的赤裸。因为灵魂的衣裳是什么？当然是德行。它的赤裸是什么？恶习。正如若有人剥去任何体面之人的衣服，那人便会羞愧，缩身躲避；同样，灵魂若我们愿意看见，那没有这些衣裳的灵魂，也会羞愧脸红。你想，此刻有多少妇人感到羞愧，宁愿沉入深渊，仿佛寻找某种幔子或屏障，以便听不到这些话？但那些良心无亏的，却兴奋、喜乐、欢愉，并用所说的话欢欣地装饰自己\OriginalTerm{装饰自己}{\GreekText{ἐγκαλλωπίζονται}}。请听关于那位真福特克拉的故事：为了能看见保禄，她甚至献出了黄金；而你们却连一分钱也不愿付出，为能看见基督。你们羡慕她的行为，却不效法她。难道你们没有听见：\BibleQuote{怜悯人的人是有福的，因为他们要蒙受怜悯？}\BibleRef{玛 5:7}你们昂贵的衣裳有何益处？我们还要张大嘴巴凝视这装束多久？让我们穿上基督的荣耀：让我们用那美丽装扮自己，好使我们在这里受称赞，在那里获得永恒的福乐，因我们的主耶稣基督的恩宠和慈爱，愿光荣、权能、尊崇归于祂，并与圣父及圣神，从今世直到永远，世世无尽。阿们。\par

\SourceRecord{Translated by J. Walker, J. Sheppard and H. Browne, and revised by George B. Stevens.；Nicene and Post-Nicene Fathers, First Series；Vol. 11.；Edited by Philip Schaff.；Buffalo, NY: Christian Literature Publishing Co.,；1889.；<http://www.newadvent.org/fathers/210125.htm>.}

% structure: p0001 source=h1 target=section reason=page-title

\section{论宗徒大事录第二十六篇讲道}

\BibleRef{宗 12:1-2}\par

\Quote{那时，\Person{黑落德}王下手磨难教会中的一些人，用剑杀了\Person{若望}的兄弟\Person{雅各伯}。他见这事中悦了犹太人，便又去捉拿\Person{伯多禄}。那时正是无酵节的日子。}\par

\BibleQuote{那时}，自然指的是紧接着的时候：因为这是圣经的惯例。他恰当地说黑落德\BibleQuote{王}（做了这事）：这不是基督时代的那个。看，另一种试炼——请注意我在开头所说的，事情如何交织在一起，在整个历史的脉络中，安息与动荡如何交替——现在不是犹太人，也不是公议会，而是王。权力越大，战争越激烈，这更是为了讨好犹太人。\BibleQuote{并且，}它说，\BibleQuote{他用剑杀了若望的兄弟雅各伯：}（抓他）是随意而无选择的。但若有人提出疑问，为何天主允许这事，我们要说，这是为了这些（犹太人）本身的益处：首先，以此说服他们，即使被杀（宗徒们）仍然得胜，正如斯德望的情形；其次，给他们机会，在满足了他们的愤怒之后，从疯狂中清醒过来；第三，向他们表明，这事是出于祂的许可。\BibleQuote{当他看到，}它说，这事使犹太人喜悦，他继续逮捕伯多禄。哦，极度的邪恶！他这样毫无计划、毫无理由地杀人来讨好他们，是为了谁呢？\BibleQuote{正是无酵饼的日子。}又是犹太人那种无聊的精确：杀人不禁止，却在这样的时刻做这种事！\BibleQuote{他逮捕了他，把他关在监里，交给四班士兵看守。}\BibleRef{宗 12:4}这既是出于愤怒，也是出于恐惧。\BibleQuote{他杀了，}它说，\BibleQuote{若望的兄弟雅各伯用剑。}你注意到他们的勇气了吗？因为，免得有人说他们不畏死亡，是因为确信天主会拯救他们，所以祂允许一些人被杀，而且是一些首领，斯德望和雅各伯，以此说服那些杀害者自己，即使这些事情也不能使他们退缩和阻碍。\BibleQuote{因此伯多禄被囚在监里；教会却不住地为他向天主祈祷。}\BibleRef{宗 12:5}因为现在生死攸关：一个人的被杀使他们害怕，另一个人的下狱也使他们恐惧。\BibleQuote{当黑落德将要提他出来的时候，那一夜伯多禄睡在两个士兵中间，被两条铁链捆着；看守的人也在门前看守监狱。忽然，主的天使站在他旁边，有光照耀在监狱里；天使拍伯多禄的肋旁，唤醒他说：快快起来。那铁链就从他手上脱落下来。}（6-7节）那一夜他解救了他。\BibleQuote{有光照耀在监狱里，}使他不会以为是幻象：没有人看见那光，只有他看见。因为，即使这样，他还以为那是幻象，因为事情太出人意料；如果没有这光，他更会这样想：他如此准备好了去死。他在那里等待了许多天而没有得救，导致了这一点。那么，你说，为什么祂不让他落入黑落德手中，然后再解救他？因为那样会使人们惊奇，而这是可信的；他们甚至不会被认为是血肉之躯。至于斯德望的情形，祂岂没有做什么？祂岂没有使他的脸如同天使的脸？但在这里，祂又有什么没有做呢？\BibleQuote{天使对他说：束上腰带，穿上鞋。}\BibleRef{宗 12:8}这里又表明，这不是出于诡计：因为一个匆忙想要逃出（监狱）的人，不会这样仔细地穿上鞋、束上腰带。\BibleQuote{他那样做了。天使对他说：披上外衣，跟着我来。他就出去跟着走，不知道天使所作的是真的，以为见了异象。过了第一层和第二层守卫，来到通往城里的铁门，那门自动给他们开了。}（9-10节）看，第二个奇迹。\BibleQuote{他们出去，走过一条街，天使就立刻离开了他。伯多禄清醒过来，说：我现在真知道主差遣了祂的天使，救我脱离黑落德的手和犹太人民所期待的一切。}（10-11节）天使离开时，伯多禄才明白：\BibleQuote{现在我明白了，}他说，不是在那时。但为什么这样，伯多禄为什么没有察觉所发生的事，虽然他曾经历过类似的解救，那时所有人都被释放？\BibleRef{宗 5:18}（主）要让他一下子得到喜乐，让他先得自由，然后才意识到所发生的事。铁链从他手上脱落的事实，也强有力地证明他没有逃跑。\BibleQuote{他想了想，就往那称呼马尔谷的若望的母亲玛利亚家去，在那里有好些人聚集祈祷。}\BibleRef{宗 12:12}注意伯多禄没有立即退走，而是先给朋友们带来好消息。\BibleQuote{伯多禄敲外门时，有一个使女名叫洛达，前来探听。她认出是伯多禄的声音，就欢喜得顾不得开门，}——连使女们都记下，她们多么虔诚——\BibleQuote{跑进去告诉众人，伯多禄站在门外。}\BibleRef{宗 12:13}但尽管这样，他们仍摇头（不信）：他们对她说：你疯了。她却坚持说确实如此。他们说：那是他的天使。\BibleQuote{伯多禄不住地敲门；他们开了门，看见他，就甚惊奇。伯多禄摆手，叫他们不要作声，就告诉他们主怎样领他出监。又说：你们把这事告诉雅各伯和众弟兄。于是出去，往别处去了。}（16-17节）现在让我们回顾叙述的顺序。\par

（摘要。）\Quote{那时，}经上说，\Quote{黑落德王下手磨难教会中的几个人。}\BibleRef{宗 12:1}他好像野兽一般，不加分辨、不加考虑地攻击所有人。这正是基督所说的：\Quote{你们必要喝我的杯，并且你们也要受我所受的洗礼。}\BibleRef{谷 10:39}(b)\Quote{他杀了若望的兄弟雅各伯。}\BibleRef{宗 12:2}因为还有另一位雅各伯，即主的兄弟；所以为区别起见，经上说\Quote{若望的兄弟}。你是否注意，事情的关键在于这三个人，尤其是伯多禄和雅各伯？(a)为何他没有立刻杀死伯多禄？经上提到原因：\Quote{那正是无酵节的日子。}他宁愿藉他的死来\OriginalTerm{炫耀}{\GreekText{ἐκπομπεὕσαι}}。\Quote{他见犹太人喜欢。}\BibleRef{宗 12:3}至于他们，因加玛里耳的劝告，已避免流血；而且甚至没有捏造控告；而是藉他人之手达到同样目的。(c)这（加玛里耳的劝告）尤其是他们的定罪：因为宣讲已显明不再是人的事。\Quote{他又进一步要用死刑对待伯多禄。}\EditorialNote{参照：宗 5:8}这确实应验了：\Quote{我们被当作待宰的羊。}\BibleRef{咏 43:13}\Quote{可见，}经上说，\Quote{令犹太人喜欢。}\BibleRef{罗 8:36}喜欢的事，流血、不义的流血、邪恶、不虔！他纵容他们\OriginalTerm{愚昧}{\GreekText{ἀτόποις}}的欲望：因为他本应做相反的事，制止他们的怒火，却使他们更急切，仿佛他是刽子手，而不是医治他们病态心灵的医生。（而）尽管他有无数的警告，关于他的祖父和父亲黑落德：前者因杀害婴儿遭受了最大的灾难，后者因杀死若翰而为自己招来一场严重的战争。但他以为……他怕伯多禄因雅各伯被杀而离开；想要将他安全看管，就把他下在监里：\Quote{把他交给四班兵士看守。}\BibleRef{宗 12:4}看守越严，显示越奇妙。\Quote{因此伯多禄被囚在监里。}\BibleRef{宗 12:5}但这反而对伯多禄更好，他因此更受考验，显出自己的勇敢。经上说，\Quote{有人恳切祈祷。}这是孝爱的祈祷：他们是为慈父祈求，一位温和的慈父。经上说\Quote{有人}，\Quote{恳切祈祷。}请听他们对老师的态度。没有党派，没有慌乱：他们转向祈祷，转向那真正无敌的联盟，以此作为避难所。他们不会说：\Quote{什么？我这卑微的受造物，为他祈祷？}因为他们出于爱而行，并不考虑这些。且注意，正是在节日期间，（敌人）给他们带来这些考验，使他们的价值更受认可。\Quote{及至黑落德}等。\BibleRef{宗 12:6}看伯多禄在睡觉，没有忧愁或恐惧！就在隔天要被提出来那天晚上，他安然入睡，把一切都交托给天主。\Quote{睡在两个兵士中间，被两条锁链绑着。}\BibleRef{伯前 5:7}注意，看守何等严密！\Quote{说：起来。}\BibleRef{宗 12:7}卫兵与他一同沉睡，因此对发生的事一无所知。\Quote{有光照耀。}这光为什么？为使伯多禄既看见又听见，不以为一切都是幻象。那命令，\Quote{快快起来。}叫他不要懈怠。他也打了他；因为他睡得很沉。(a)\Quote{起来，}他说，\Quote{快些：}这不是为了催逼他（\OriginalTerm{催逼}{\GreekText{θορυβοὕντος}}），而是劝他不要迟延。(c)\Quote{并且}立刻\Quote{锁链从他手上脱落。}(b)怎么？请回答我：异端者在哪里？——让他们回答。\Quote{天使对他说}等。\BibleRef{宗 12:8}藉此也使他相信这不是幻象：为此他命他束上腰带，穿上鞋，使他摆脱睡意，知道这是真实的。(a)(e)\Quote{他不知天使所作的是真的，以为见了异象。}\BibleRef{宗 12:9}(e)这也难怪，因为所发生的事\OriginalTerm{过于伟大}{\GreekText{ὑπερβολήν}}。你是否注意，一个神迹\OriginalTerm{过于伟大}{\GreekText{ὑπερβολὴ} \GreekText{σημείου}}会怎样？如何使观看者\OriginalTerm{惊愕}{\GreekText{ἐκπλήττει}}？如何使人难以置信？因为如果伯多禄\Quote{以为见了异象}，尽管他已经束了腰带穿了鞋，别人又当如何？\Quote{并且，}经上说，\Quote{过了第一层和第二层守卫，他们来到铁门，那门自动给他们开了。}\BibleRef{宗 12:10}然而发生在（狱）内的事更为奇妙；但这件却更合乎人的常情。\Quote{出来以后，他们走过一条街，天使立刻离开了他。}\BibleRef{宗 12:11}因为没有阻碍时，天使才离开。因为伯多禄不会继续\OriginalTerm{前行}{\GreekText{προἥλθεν}}，因为有这么多阻碍。\Quote{他清醒过来：}因为确实，这是一件令人惊愕的事（\OriginalTerm{惊愕}{\GreekText{ἔκπληξις}}）。\Quote{如今，}他说，\Quote{我知道}——如今，不是当时在狱中时——\Quote{主派遣了祂的天使，救我脱离黑落德的手和犹太人民的一切期待。他考虑之后，}\BibleRef{宗 12:12}经上说：即他身在何处，或他不可就此离去，而要报答他的恩主：\Quote{他来到若望的母亲玛利亚家。}这位若望是谁？大概是经常与他们在一起的那位：因此他加上了他的\OriginalTerm{别号}{\GreekText{τὸ} \GreekText{παράσημον}}，\Quote{又名马尔谷。}但请注意，在夜间\Quote{祈祷}，他们获得了多少益处：患难是何等美好；它使他们何等警醒！你看见从斯德望的死所得的利益有多大吗？你看见从这次监禁所得的益处有多大吗？因为天主并非通过报复那些恶待他们的人来显示福音的伟大：而是在作恶者身上，不伤害那些人，他显明患难本身是何等强大的事物，好使我们绝不寻求脱离它们，也不寻求报复我们的冤屈。且看，此后连使女也与他们平等。经上说，\Quote{因欢喜，她没有开门。}（13-14节）这也做得好，使他们不致因立刻见到他而惊愕，也使他们不信，而心灵受到操练。\Quote{却跑进去}等等，正如我们惯常所做的，她急于亲自成为报喜讯的人，因为这确实是喜讯。\Quote{他们对她说：你疯了！她却坚持说确是如此。他们便说：那是他的天使。}\BibleRef{宗 12:15}这是一个真理：每人都有天使。天使要做什么？他们是在（夜间）那时候猜到的。但当他\Quote{继续敲门，他们开了门，看见他，就惊愕了。他却摆手}（16-17节），叫他们安静，听他所遭遇的一切。他现在成为门徒们更亲切渴望的对象，不仅因为他得救，更因他突然进来又立刻离去。现在，他的朋友们清楚地知道了一切；外人也知道了，如果他们愿意，但他们不愿意。在基督身上也发生了同样的事。他说：\Quote{把这些事告诉雅各伯和弟兄们。}何等毫无虚荣！他并没有说：把这些事传扬给各处的人，而是\Quote{告诉弟兄们。然后他退到别处去了。}因为他没有试探天主，也没有把自己投入试探：因为他们奉命这样做时才这样做。\Quote{你们去，}曾说过，\Quote{在圣殿里向百姓讲论。}\BibleRef{宗 5:20}但天使在这里并没有这样说；相反，他默默地带他出来，在夜间领他出去，给了他自由离开的许可——这样做的目地也是为使我们知道，许多事情是按人的方式由上智安排的——使他不再陷入危险。——为了使他们不能说：\Quote{那是他的天使，}在他走后，他们先这样说，然后看见他自己推翻了他们的想法。如果是天使，他会敲门，不会退到别处去。白天所发生的事使他们确信。\par

\BibleQuote{于是\Person{伯多禄}被拘留在监狱里}等\BibleRef{宗 12:5} 他们自由自在，正在祈祷：他却被捆锁，睡着了。\BibleQuote{他并不知道这事是真的。}\BibleRef{宗 12:9} 如果他以为所发生的事是真的，他必会惊奇，不会记住（所有细节）；但如今，他似乎在做梦，毫无扰动。\BibleQuote{当他们经过第一和第二道岗哨时}——你看守卫多么严密——\BibleQuote{他们来到了铁门。}\BibleRef{宗 12:10} \BibleQuote{如今我知道主派遣了他的天使。}\BibleRef{宗 12:11} 为什么这事不是由他们自己完成的？（我回答：）主也以此尊荣他们，藉着天使的职务拯救他们。那么，为什么在\Person{保禄}的事上并非如此呢？那里有充分的理由，因为狱警要归化，而在这里，只需释放宗徒。\BibleRef{宗 16:25} 天主以各种方式安排一切。在那里也同样美好，\Person{保禄}唱圣咏，而这里\Person{伯多禄}睡着了。\BibleQuote{他思考后，来到\Person{玛利亚}的家}等\BibleRef{宗 12:12} 那么，我们不要隐藏天主的奇迹，而为了我们的益处，让我们努力将这些宣扬出去，以建树他人。因为他为选择被捆锁而受钦佩，同样更值得钦佩的是，他没有退却，直到将所有事报告给他的朋友。\BibleQuote{他说：告诉\Person{雅各伯}和弟兄们。}\BibleRef{宗 12:17} 好使他们喜乐，不再忧虑。藉着这些人，那些人得知，而非那些人藉着他得知：他对卑微者有如此关怀！——\par

诚然，没有什么比\OriginalTerm{适度的苦难}{\GreekText{συμμέτρου}}更好的了。你想他们当时的心绪是何等充满喜乐！如今那些整夜沉睡的妇女在哪里？那些在床榻上连身也不翻的男子在哪里？你看到那警醒的灵魂吗？与妇女、孩童、女仆一同向天主唱圣咏，因苦难变得比天空更纯洁。但如今，我们若见到一点危险，就退缩。没有比那教会更辉煌的了。让我们效法他们，让我们以他们为榜样。黑夜不是为此而造，让我们整夜沉睡、无所事事。对此，手艺匠、搬运夫和商人（对此）作证，天主的教会深夜中起身。你也起身，观看星辰的合唱、深沉的寂静、深邃的安息；怀着敬畏默想你主人家庭的\OriginalTerm{安排}{\GreekText{οἰκονομίαν}}。那时你的灵魂更加纯洁：更轻盈、更精细、自由翱翔：黑暗本身、深沉的寂静，足以引导你痛悔。你若再看那布满星辰的天空，仿佛千万只眼睛；你若想到那些在白天呼喊、欢笑、雀跃、蹦跳、行不义、攫取、威胁、施行无数不义的人群，都如同死人一样躺卧，你必谴责人一切的任性与固执。睡眠\OriginalTerm{侵入并击败了}{\GreekText{ἡ} \GreekText{λεγξεν}}本性：它是死亡的肖像，是万物终局的肖像。你若（凭窗）俯身街道，甚至听不到一点声响；你若探视屋内，你会看到一切都仿佛躺在坟墓中。这一切足以唤醒灵魂，引导它默想万物的终局。\par

事实上，我的讲论是为男人和女人的。请屈膝，发出呻吟，恳求你的主施以仁慈：祂更在夜间被祈祷所感动，那时你把休息的时刻变成哀哭的时刻。记住那位君王所说的话：\BibleQuote{我因呻吟而疲倦：每夜我要洗我的床，用眼泪灌溉我的卧榻。}\BibleRef{咏 6:6} 无论你如何娇生惯养，你也不比他更娇嫩；无论你如何富有，你也不比达味更富有。同样，这位圣咏作者又说：\BibleQuote{半夜里我起来感谢你，因你正义的判决。}\BibleRef{咏 118:62} 那时，没有虚荣侵入你：怎么会呢，当众人都睡着，不看你的时候？那时，也没有懒惰或困倦侵袭你：怎么会呢，当你的灵魂因如此伟大的事物而被激动？这样的守夜之后，带来甜美的睡眠和奇妙的启示。你也该这样做，男人，不只是女人。让房屋成为一座教会，由男人和女人组成。因为不要以为因为你是唯一的男人，或她是唯一的女人，这便是什么障碍。\Quote{因为哪里有两个，}祂说，\Quote{因我的名聚在一起，我就在他们中间。}\BibleRef{玛 18:20} 哪里基督在中间，那里就有大群众。哪里基督在，那里就必须有天使，也必须有总领天使和其他众能者。那么你并不孤单，因为你有祂，祂是万有的主。再听先知也说：\BibleQuote{遵行主旨意的一人，胜过一万个违犯者。}\BibleRef{德 16:3} 没有比一群不义之人更软弱的，也没有比一个按天主法律生活的人更坚强的。如果你有儿女，也叫醒他们，让你的房屋整夜成为一座教会：但如果他们幼小，不能忍受守夜，让他们停留到第一次或第二次祈祷，然后送他们去休息：只激动你自己，使自己养成习惯。没有比那收纳如此祈祷的宝藏更好的。听先知说：\BibleQuote{我在床上想起了你，在清晨黎明中思念你。}\BibleRef{咏 62:7} 但你会说：我白天劳苦很多，不能。这只是借口和托辞。因为你无论劳苦多少，都不会像铁匠那样辛劳，他从高处放下如此沉重的锤子，落在（金属飞溅的）火花上，全身吸入烟雾：然而在这工作上他花了大部分夜晚。你也知道女人们，如果我们需要到乡下去，或出去守夜，她们整夜守望。那么，你也要有一个属灵的铁匠铺，在那里塑造的不是锅或壶，而是你自己的灵魂，这比铜匠或金匠所能塑造的好得多。你的灵魂，因罪而衰老，你把它投入告解的熔炉中：让锤子从高处落下：那就是对你\OriginalTerm{言语的谴责}{\GreekText{τὥν} \GreekText{ῥημάτων} \GreekText{τὴν} \GreekText{κατάγνωσιν}}：点燃圣神的火。你有一个远远更强大的手艺（比他们的）。你正在塑造的不是金器，而是灵魂，它比所有金子更珍贵，正如铁匠锤打出他的器皿。因为你正在工作的不是物质器皿，而是使你的灵魂摆脱属于此世的一切幻想。让一盏灯放在你身边，不是我们点的那盏，而是先知所有的，当他说：\BibleQuote{你的言语是我脚前的明灯。}\BibleRef{咏 118:105} 用祈祷把你的灵魂烧成赤热：当你看到它够热时，把它拉出来，塑造成你想要的任何形状。相信我，没有火能如此有效地烧去锈迹，如同夜间祈祷除去我们罪恶的锈迹。让守夜者，如果没有别人，使我们羞愧。他们，根据人的法律，在寒冷中巡逻，大声喊叫，行走在\OriginalTerm{窄巷}{\GreekText{στενωπὥν}}和小巷中，常常被雨淋透，冻得僵硬，为了你和你的安全，以及保护你的财产。他如此关心你的财产，而你甚至不为你的灵魂操心。然而，我不叫你像他一样在露天巡逻，也不大声喊叫、喊破喉咙：而是在你的密室或你的卧室里，屈膝，恳求你的主。为什么基督自己整夜在山上？难道不是为作我们的榜样？那时，植物在夜间呼吸，我是说：而灵魂吸入露水甚至比它们更多。太阳在白天晒干的东西，夜间又变凉。比一切露水更清新的是，夜间的眼泪降在我们的情欲上，以及灵魂的一切热和烧上，不让它受任何这样的影响。但如果它没有享受那露水的益处，就会在白日被烧干。但愿不是如此！反而，愿我们众人，得到更新，享受天主的仁慈，从罪恶的负担中释放，因我们的主耶稣基督的恩宠和仁慈，祂偕同圣父和圣神，共享光荣、权能、荣耀，从今世直到永远。阿们。\par

\SourceRecord{Translated by J. Walker, J. Sheppard and H. Browne, and revised by George B. Stevens.；Nicene and Post-Nicene Fathers, First Series；Vol. 11.；Edited by Philip Schaff.；Buffalo, NY: Christian Literature Publishing Co.,；1889.；<http://www.newadvent.org/fathers/210126.htm>.}

% structure: p0001 source=h1 target=section reason=page-title

\section{\WorkTitle{宗徒大事录第二十七篇讲道}}

\BibleRef{宗 12:18-19}\par

\BibleQuote{天一亮，在士兵中起了不小的骚动，不知道伯多禄发生了什么事。黑落德派人搜索他，没有找到，便审讯卫兵，下令把他们处死。然后他从犹太下去，到凯撒勒雅，住在那里。}\par

有些人可能感到困惑，为什么天主会静观其勇士被处死，而现在又因伯多禄的缘故处死那些士兵；然而，天主在（解救）伯多禄之后，本也可以解救他们。但审判的时刻尚未到来，无法按各人的行为报应。此外，并非伯多禄将他们交到黑落德手中。因为最使他恼怒的是被愚弄；正如他祖父被贤士欺骗时那样，使他痛彻心扉的正是被（逃避并）戏弄。\BibleQuote{审讯之后，}经上记着，\BibleQuote{他下令把他们处死。}\BibleRef{玛 2:16}然而，他曾从他们那里听说——因为他审讯了他们——锁链被留下，他取了凉鞋，而且直到那夜他还在他们中间。\BibleQuote{审讯之后：}但他们隐瞒了什么？为什么他们自己也不逃走？\BibleQuote{他下令把他们处死：}然而他本应惊讶，本应为此惊奇。结果，因这些人的死，（这事）向众人显明了：他的邪恶暴露无遗，并且（显明了）这奇迹（是）出于天主。\BibleQuote{他从犹太下去，到凯撒勒雅，住在那里。黑落德对提洛和漆冬人非常恼怒；但他们同心合意到他那里，先说服了国王的内侍布拉斯特，便请求和解，因为他们的地区是靠国王的地区供养。在约定的日子，黑落德穿上王袍，坐在审判台上，向他们讲话。民众呼喊说：\InnerQuote{这是神的声音，不是人的声音。}立刻，主的天使打击了他，因为他没有归光荣于天主；他被虫子咬死，断了气。}\BibleRef{宗 12:20} * * 但请看（作者）在此并没有隐瞒这些事。他为什么提及这段历史？请问，黑落德对提洛人和漆冬人发怒，这与福音有何相干？这并非小事，甚至这一点，正义如何立刻惩罚了他；虽然不是因伯多禄的事，而是因他的傲慢言论。然而，有人会说，如果那些人呼喊，与他何干？因为他接受了欢呼，因为他自认为配受敬拜。通过他，那些盲目奉承他的人（另译：\OriginalTerm{献媚者}{\GreekText{οἱ} \GreekText{κολακεύοντες}}）受到最大的教训。请再留意，双方都应受惩罚，但这人受了惩罚。因为这不是审判的时候，但祂惩罚了罪责最重的那人，让其他人从这人的遭遇中受益。\BibleQuote{天主的圣言，}经上记着，\BibleQuote{日渐繁衍，}即因此，\BibleQuote{而增多。}\BibleRef{宗 12:24}你注意到天主的眷顾安排了吗？\BibleQuote{巴尔纳伯和扫禄完成了他们的职务，便从耶路撒冷回来，带着那称为马尔谷的若望同去。}\BibleRef{宗 12:25}\BibleQuote{在安提约基雅的教会中，有一些先知和教师，就是巴尔纳伯、号称尼格尔的西默盎、基勒乃的路求、与分封侯黑落德一同长大的玛纳恒，以及扫禄。}\BibleRef{宗 13:1}他仍将巴尔纳伯列在前面：因为保禄尚未出名，尚未行过任何奇迹。\BibleQuote{他们敬拜主并禁食的时候，圣神说：'你们为我选出巴尔纳伯和扫禄，去作我召叫他们从事的工作。'他们禁食祈祷后，就给他们覆手，派他们走了。}（第2-3节）'敬拜'是什么意思？传道。\BibleQuote{你们为我选出，}它说，\BibleQuote{巴尔纳伯和扫禄。}\BibleQuote{你们为我选出}是什么意思？为工作，为宗徒职务。请再看，他是被怎样的人所祝圣的（\OriginalTerm{更赤裸的}{\GreekText{γυμνοτέρα}} . Cat. \OriginalTerm{更庄严的}{\GreekText{σεμνοτέρα}}, \Quote{更可敬畏的}）。是由基勒乃的路求和玛纳恒，或者更确切地说，是由圣神。这些人越卑微，恩宠就越明显。从此他被祝圣为宗徒，以便有权威地宣讲。那么，他自己怎么说：\BibleQuote{不是从人，也不是藉着人？}\BibleRef{迦 1:1}因为不是人召唤或带领他过来：这就是为什么他说：\BibleQuote{不是从人，也不是藉着人，}即他不是被这人（所派），而是被圣神所派。因此（作者）接着写道：\BibleQuote{他们既被圣神派遣，就下到色娄基雅，又从那里乘船往塞浦路斯去。}\BibleRef{宗 13:4}但让我们再看一遍刚才所说的。\par

（复述。）\BibleQuote{及至天亮，}等等。\BibleRef{宗 12:18} 因为如果天使也把兵士与\Person{伯多禄}一同带出来，人们就会以为是逃跑。那么，你可能会问，为什么没有用别的方式处置呢？何害之有？如果我们看见那些无辜受苦的人并未受害，我们就不会提出这些问题。因为你为什么不对\Person{雅各伯}说同样的话呢？为什么（天主）没有救他？\BibleQuote{兵士们起了不小的骚动。} 所以（显然）他们什么也没有察觉（所发生的事）。看，我替他们辩护。锁链在那里，看守在里面，监狱关闭，没有一处墙壁被拆毁，一切都说明同一件事：那人被带走了；你为什么定他们的罪？如果他们想放他走，他们早就会这样做，或者会和他一起出去。\Quote{但他给他们钱？}\BibleRef{宗 3:6} 而他连给穷人钱都没有，又怎会有钱给他们呢？而且锁链既没有被折断，也没有被解开。他本应看出，这事出于天主，而非人的作为。\BibleQuote{他从\Place{犹太}下到\Place{凯撒勒雅}，住在那里。\Person{黑落德}对\Place{提洛}和\Place{漆冬}的人极其恼怒，}等等。\BibleRef{宗 12:19} 他现在要提到一件史事：这是为什么他加上这些名字的原因，为要表明他如何在一切事上忠于真理。经上说：\BibleQuote{他们买通了国王的内侍\Person{布拉斯特}，就向他求和，因为他们的地区靠王家的地区供养。}（第20、21节）大概是因为发生了饥荒。\BibleQuote{到了约定的日子，}等等。（\Person{若瑟夫}\WorkTitle{犹太古史}第十九卷）\Person{若瑟夫}也这样说，说他患了慢性病。当时一般人不清楚这件事，但宗徒记述了下来；然而他们的无知反而有益，因为人们把（\Person{阿格黎帕}）所遭遇的事归咎于他杀害\Person{雅各伯}和那些兵士。请注意，当他杀害宗徒时，并未发生这类事；但当他杀害这些人时，却发生了；事实上他不知道对此说什么；因着困惑和羞愧，\BibleQuote{他从\Place{犹太}下到\Place{凯撒勒雅}。}我想他退离（\Place{耶路撒冷}）也是为了叫那些人（\Place{提洛}和\Place{漆冬}的人）来道歉：因为他恼怒那些人，却向这些人献殷勤。看，这人多么虚荣：他本意是要施恩于他们，却发表了一篇演讲。但\Person{若瑟夫}说，他还穿了一件用银子做的华丽袍子。注意那些奉承者是什么样的人，以及宗徒们表现出多么高的精神：全民族如此讨好的人，他们却藐视他。\BibleRef{宗 12:24} 但再次看到，他们获得了极大的安慰，以及从他受到的惩罚中产生的无数益处。但如果这个人，因为有人向他说：\BibleQuote{这是神的声音，不是人的声音}\BibleRef{宗 12:22}——尽管他自己什么也没说——就遭受了这样的事，那么基督若祂自己不是天主（就该受更重的惩罚），因为祂常说：\BibleQuote{我这些话不是我的}\BibleRef{若 14:10}，以及\BibleQuote{天使服侍我}之类的话。但那个人以可耻悲惨的方式结束了生命，从此不再出现。也注意他，甚至轻易被\Person{布拉斯特}说服，像个可怜虫，一会儿恼怒，一会儿又平静，在所有场合都是民众的奴隶，身上没有自由独立的东西。但也注意圣神的权威：\BibleQuote{他们敬拜主并禁食的时候，圣神说：你们为我分别出\Person{巴尔纳伯}和\Person{扫禄}。}\BibleRef{宗 13:2} 哪个受造物如果不在同一权威下，敢于说这样的话？\BibleQuote{分别出来，}等等。这样做是为了使他们不聚在一起。圣神看见他们有更大的能力，并且能足够应付许多事。祂如何对他们说话？大概是藉着先知：因此作者事先说明，那里也有先知。他们禁食并敬拜：为叫你们知道需要极大的清醒。他在\Place{安提约基雅}被祝圣，在那里讲道。为什么祂不说：为主分别，而说\BibleQuote{为我？}这表明祂具有同一权威和能力。\BibleQuote{他们禁食完了，}等等。你们看见禁食是多么伟大的事吗？\BibleQuote{他们既被圣神派遣：}这表明一切是圣神做的。\par

斋戒是一件伟大的善事，确实伟大：它不受任何限制。当需要授予圣职时，他们就斋戒；就在他们斋戒时，圣神向他们发言。我只吩咐这一点：（我说）不是斋戒，而是戒绝奢华。让我们寻求滋养的肉食，而非毁灭我们的东西；寻求作为食物的肉食，而非疾病（灵魂和身体的疾病）的缘由；寻求带来安慰的食物，而非充满不适的奢华：其一是奢华，其二是祸害；其一是快乐，其二是痛苦；其一是合乎本性的，其二是违反本性的。比方说，如果有人给你喝毒芹汁，难道不是违反本性吗？如果有人给你木头和石头，你难道会不拒绝吗？当然，因为它们是违反本性的。好吧，奢华也是如此。因为正如一座城市，在敌人入侵时遭受围困和骚乱，喧嚣极大；灵魂在酒和奢华的侵袭下也是如此。\BibleQuote{谁有祸？谁有骚乱？谁有不适和胡言乱语？不就是那些长久沉溺于酒的人吗？谁的眼睛布满血丝？}\BibleRef{箴 23:29} 但是，无论我们说什么，我们将无法挽回那些沉溺于奢华的人，除非我们引入一种相反的情感与之对抗。首先，让我们向妇女们进言。没有什么比一个沉溺于奢华的女人更丑陋，没有什么比一个酗酒的女人更丑陋。她肤色的光彩褪去了；眼睛中平静温和的神情变得混浊，就像云层遮住了阳光。这是一种粗俗的（\OriginalTerm{卑贱}{\GreekText{ἀνελεύθερον}}）、奴隶般的、完全低贱的习惯。一个女人呼出酒酸臭气时多么令人厌恶：她打嗝，呼出腐烂肉食的（\OriginalTerm{腐臭气味}{\GreekText{χυμόν}}）气味；她自己沉重不堪，无法保持直立；脸上泛着不自然的红色；不停地打哈欠，眼前一切都朦胧如雾！但戒绝奢华的女人并非如此：不，（这种节制使她看起来）成为一个更美丽、更端庄的（\OriginalTerm{更端庄}{\GreekText{σωφρονεστέρα}}）女人。因为即使对身体，灵魂的沉静也会赋予它自身的美。不要以为美的印象仅来自身体的特征。给我一个漂亮的女孩，但她\OriginalTerm{暴躁}{\GreekText{τεταραγμένην}}、多话、爱骂人、酗酒、挥霍无度，（告诉我）她不是比任何丑女人更难看吗？但如果她羞怯，如果她沉默，如果她学会脸红，如果她谦逊地（\OriginalTerm{适度地}{\GreekText{συμμέτρως}}）说话，如果她找到时间斋戒；她的美丽将加倍，她的青春会更加焕发，她的容貌更加动人，充满\OriginalTerm{节制和端庄}{\GreekText{σωφροσύνης} \GreekText{καὶ} \GreekText{κοσμιότητος}}。现在，我们来说说男人吗？还有什么比一个醉酒的男人更丑陋呢？他成为仆人的笑柄，敌人的笑柄，朋友的怜悯；应受无尽的谴责：更像野兽而非人类；因为吞食大量食物是豹子、狮子和熊的本性。它们这样做不足为奇，因为这些生物没有理性的灵魂。然而，即使它们，如果被食物塞满超过所需，超过本性所定的量，也会因此毁坏全身；何况我们呢？因此，天主将我们的胃收缩成小小的空间；因此，祂规定了少量的食粮，好教导我们关心灵魂。\par

让我们思考我们的构造，我们会发现我们内里只有一小部分有这种功能——因为我们的口和舌是为歌唱赞美诗，喉咙是为发音——因此，自然的需求约束了我们，使我们甚至不由自主地不会陷入许多\OriginalTerm{麻烦}{\GreekText{πραγματείαν}}（以这种方式）。因为，如果奢侈生活没有它的痛苦，没有疾病和虚弱，那或许还可以容忍：但事实是，祂藉着自然的限制约束了你，即使你想过度，你也无法如此。你所追求的是快乐吗，亲爱的？你将从节制中找到它。是健康吗？你也会这样获得。是心安吗？这也是。是自由吗？是身体的活力和良好状态吗？是头脑的清醒和警觉吗？（所有这一切你都会找到）；所以所有美好的事物都在那里，而在另一边则是这些的反面：不适、失调、疾病、窘迫——\OriginalTerm{吝啬}{\GreekText{ἀνελευθερία}}。那么你会问，我们大家都如此热切地追求它，这是怎么回事呢？这是出于疾病。因为，请说，是什么让病人渴望那伤害他的东西？这种渴望本身不就是他疾病的一部分吗？为什么跛足的人不能直立行走？这件事，难道是因为他懒惰，不愿意去找医生吗？因为有些事物中，它们所带来的快乐是暂时的，而惩罚却是持久的；另一些则恰恰相反，忍受是暂时的，快乐却是永恒的。因此，一个意志不够坚定和刚强，不愿为了将来而轻视眼前甜蜜的人，很快就会被打败。请说，厄撒乌是怎么被打败的？他怎么会将目前的快乐置于未来的尊荣之上？由于缺乏坚定和刚强的品格。\BibleRef{创 25:33}而你说，这缺点本身来自哪里？来自我们自己；从这一点就可以看出。当我们有心时，我们就会振作起来，变得能够忍受。可以肯定的是，如果有时必要落到我们身上，甚至常常只是出于竞争精神，我们就会清楚地看到什么对我们有益。因此，当你将要放纵奢侈时，想想快乐是何等短暂，想想损失——因为把钱花在伤害自己上确实是一种损失——疾病、虚弱：并鄙视奢侈。我还要列举多少人因放纵而遭受了恶果？诺厄醉了酒，赤身露体，看这带来了什么恶果。\BibleRef{创 9:20}厄撒乌因贪吃放弃了长子的名分，并起了杀弟之心。以色列子民\Quote{坐下吃喝，起来玩耍。}\BibleRef{出 32:6}因此圣经说：\Quote{当你吃了喝了，记念上主你的天主。}\BibleRef{申 6:12}因为他们陷入奢侈，如同跌下悬崖。他说：\Quote{那好宴乐的寡妇，虽活着也是死的。}\BibleRef{弟前 5:6}又说：\Quote{蒙爱的肥胖了，粗壮了，踢跳起来。}\BibleRef{申 32:15}再者，宗徒说：\Quote{不要为肉体安排，去放纵私欲。}\BibleRef{罗 13:14}我不是立法要求禁食，因为确实没有人会听；但我是在消除美味，为了你们的益处而砍断奢侈：因为奢侈如同冬天的急流，冲垮一切；没有什么能阻止它的进程：它把人逐出王国；有什么\OriginalTerm{益处}{\GreekText{τί} \GreekText{τὸ} \GreekText{πλέον}}呢？你愿意享受（真正的）奢侈吗？施舍给穷人，邀请基督，这样即使在撤席之后，你仍能享受这种奢侈。因为现在，你确实没有它，这并不奇怪：但那时你会有。你愿意品尝（真正的）奢侈吗？喂养你的灵魂，给她她习惯的食物，不要让她因饥饿而死。——现在是战争的时候，是竞赛的时候：而你却坐着享乐？你难道没看见那些执掌权杖的人，在出征时是如何节俭度日的吗？\Quote{我们不是对抗血肉。}\BibleRef{弗 6:12}而你在即将摔跤时却养肥自己？仇敌正在咬牙切齿，而你就放纵欢愉，专心于餐桌？我知道我说这些话是徒然的，但并非对所有人都是徒然。\Quote{有耳听的，就听吧！}\BibleRef{路 8:8}基督因饥饿而消瘦，而你却因贪食而\OriginalTerm{浪费}{\GreekText{διασπᾅς}}自己的生命？\OriginalTerm{两种无度}{\GreekText{Δύο} \GreekText{ἀμετρίαι}}。因为奢侈不产生什么恶呢？它自相矛盾：所以我不知它如何得名：但正如那真正是耻辱的被称为荣耀，那真正是贫穷的被称为财富，同样，奢侈之名被给予那实质上是令人作呕的东西。我们把自己当作屠宰场吗，这样养肥自己？为什么要为虫子准备丰盛的储藏室？为什么要增加它们的\OriginalTerm{体液}{\GreekText{ἰχώρας}}？为什么要储存汗水和臭味的来源？为什么要使自己对于任何事情都无用？你希望眼睛强壮吗？希望身体强健吗？因为在琴弦中，粗糙而未精炼的，不适合发出乐音，但那些被好好刮过、拉得紧、振谐完美的，才能奏出和谐。你为什么把灵魂活活埋葬？为什么把周围的墙加厚？为什么要增加烟雾和云彩，如同雾气从四面八方蒸腾而上？如果没有别的，就让摔跤手教导你：身体越瘦越强壮，并且（然后）灵魂也更活跃。事实上，这就像战车御者和马匹。但你看，就像那些放纵奢侈、使自己肥胖的人一样，肥胖的马匹笨拙，给御者带来许多麻烦。一个人可能认为自己\OriginalTerm{还算幸运}{\GreekText{ἀγαπητὸν}}，即使有一匹顺从缰绳、四肢健全的马，能够赢得奖品：但当御者被迫拖着马走，马摔倒时，即使他再怎么鞭策，也不能让它站起来，即使他自己多么有技巧，他也将失去胜利。那么，让我们不要容忍因肉体而伤害灵魂，而要让灵魂自己更敏锐，让她的翅膀轻盈，让她的束缚更松弛：让我们用言论、用节俭来喂养她，（喂养）肉体只到足以健康、强壮、喜悦而不痛苦的程度：这样妥善地处理我们的事务，我们才能得以抓住最高的德行，并因我们的主耶稣基督的恩宠和慈爱，与圣父及圣神一同，获得荣耀、权能、尊崇，从现在直到永远，世世无尽。阿们。\par

\SourceRecord{Translated by J. Walker, J. Sheppard and H. Browne, and revised by George B. Stevens.；Nicene and Post-Nicene Fathers, First Series；Vol. 11.；Edited by Philip Schaff.；Buffalo, NY: Christian Literature Publishing Co.,；1889.；<http://www.newadvent.org/fathers/210127.htm>.}

% structure: p0001 source=h1 target=section reason=page-title

\section{论宗徒大事录第二十八讲}

宗徒大事录第十三章四、五节\par

\BibleQuote{他们既被圣神派遣，就下到塞琉西亚，从那里坐船往塞浦路斯。到了撒拉米，就在犹太人的会堂中宣讲天主的圣言：他们也有若望作助手。}\par

他们一被祝圣就出发，急忙赶往塞浦路斯，因为那里没有蓄谋陷害他们的诡计，而且圣言也已经播在那里。在安提约基雅有足够的教师，腓尼基也靠近巴勒斯坦，但塞浦路斯却不是这样。然而，你们不当追问原因和理由，既然是圣神引领他们的行动：因为他们不仅是由圣神祝圣，也是由祂派遣的。\newline{} \BibleQuote{他们来到撒拉米，就在犹太人的会堂中宣讲天主的圣言。} 你看，他们如何刻意先向他们宣讲圣言，不使他们更加争辩？前面提到的人 \BibleQuote{只向犹太人讲道} \BibleRef{宗 11:19}，所以他们这里也前往会堂。\newline{} \BibleQuote{他们走遍全岛，直到帕佛，遇见一个术士，是个假先知，犹太人，名叫巴耶稣；这人常与总督色尔爵·保禄在一起，色尔爵·保禄是个明智的人，他叫了巴尔纳伯和扫罗来，想听天主的圣言。但术士厄里玛——他的名字就是这样解释的——反对他们，想要使总督背弃信德。} \BibleRef{宗 13:6} 又是一个犹太术士，如同\Person{西满}一样。请注意这个人，当他们向别人宣讲时，他并不太在意，但一当他们接近总督，就出来反对。论到总督，令人惊奇的是，虽然他受了那人的邪术影响，却仍愿意听宗徒们讲道。撒玛黎雅人也是这样：从比较（\OriginalTerm{竞争}{\GreekText{συγκρίσεως}}）中可以看出胜利，邪术被击败了。处处可见，虚荣和贪权是（丰富的）罪恶之源！\newline{} \BibleQuote{可是扫罗，又名保禄，}——\BibleRef{宗 13:9} 在这里他的名字同时被改变，正如\Person{伯多禄}的情形一样——\BibleQuote{充满了圣神，注视着他，说：你这充满各种诡诈和各种奸诈的人，魔鬼的儿子，} \BibleRef{宗 13:10} 并且请注意，这不是辱骂，而是控诉：因为应该这样谴责横蛮无耻的人。\BibleQuote{一切正义的敌人，} 他揭露了那人内心的想法，即假借拯救之名而实则毁灭总督。\BibleQuote{你还不停止，} 他说，\BibleQuote{扰乱主的正道吗？} （他说这话）既自信（\OriginalTerm{自信地}{\GreekText{αξιοπίστως}}），与你作战的不是我们，你也不是攻击我们，而是你 \BibleQuote{扰乱主的正道，} 并且称赞（这些道，他补充说）\BibleQuote{正直的} 道。\newline{} \BibleQuote{你看，现在主的手加在你身上，你要变成瞎子。} \BibleRef{宗 13:11} 这是他本人曾借此归化的标记，他愿用此来使这人悔改。同样，那句话，\BibleQuote{暂时，} 不是作为惩罚，而是为了他的悔改：若是惩罚，他就会使他永久瞎眼，但现在却不是这样，而是 \BibleQuote{暂时，}（并且）为赢得总督。因为，他既受邪术影响，用这惩罚来教训他是好的（也教训术士），正如（埃及的）术士曾被疮教训一样。\BibleRef{出 9:11}\newline{} \BibleQuote{立刻有雾和黑暗落在他身上，他就四处找人牵他的手。总督看见所发生的事，就信了，对主的道理感到惊奇。} \BibleRef{宗 13:12} 但请注意，他们并没有在那里停留，如同（可能被诱惑那样）现在总督成了信徒，也不因受人奉承尊敬而软弱，而是立刻继续工作，出发前往对面海岸的地区。\BibleQuote{保禄和他的同伴从帕佛起航，来到旁非里雅的佩尔革；若望却离开他们，回了耶路撒冷。当他们离开佩尔革，来到丕息狄雅的安提约基雅，在安息日进了会堂坐下。} （第十三、十四节）\newline{} 他们再次进入会堂，以犹太人的身份，以免被当作敌人而赶走：这样他们顺利完成了整个工作。\BibleQuote{在诵读法律和先知之后，会堂长派人到他们跟前说：诸位仁人弟兄，若你们有什么劝勉百姓的话，请说罢。} \BibleRef{宗 13:15} 从这一点，我们得知保禄行事的经过，正如以上所说的，我们也不少了解伯多禄。但让我们回顾一下所说的话。\par

（概述。）\BibleQuote{他们来到撒拉米时，}即塞浦路斯的首府，\BibleQuote{他们宣讲天主的圣言。}\BibleRef{宗 13:5}他们在安提约基雅已住了一年：理应也要到此处（塞浦路斯）来，而非永远留在他们所在之处（塞浦路斯的归化者）：需要更伟大的教师。也看他们如何不在色娄基雅停留，知道（那里的人）可能从邻近城市（安提约基雅）已获益良多；但他们赶去处理更迫切的事务。当他们来到该岛的首府时，他们热切地要\OriginalTerm{纠正}{\GreekText{διορθωσαι}}总督。但（作者）说\BibleQuote{他同总督在一起，总督是个明智之人}\BibleRef{宗 13:7}并非奉承，你可以从事实得知；因为他不需要许多论辩，而且他自己愿意听他们。他也提到了名字。 * * * 注意，他如何对术士一言不发，直到术士给了他机会；但他们只\BibleQuote{宣讲了主的圣言。}既然（厄吕马）看见其余的人都听从他们，他只看重这一件事，便是总督不可被争取过来。为什么（保禄）没有行别的奇迹？因为没有一件能与这件事相比，即俘虏敌人。注意，他先控告，然后惩罚他。他指出此人遭受苦难是多么公正，通过他说：\BibleQuote{哦，充满一切欺骗}\BibleRef{宗 13:10}：（\Quote{充满一切，}）他说：什么都不缺，满盈满溢；他说得好，一切的\Quote{欺骗，}因为此人扮演伪君子的角色。——\BibleQuote{魔鬼的儿子，}因为他做魔鬼的工作：\BibleQuote{一切正义的仇敌，}因为（他们所宣讲的）就是全部的正义（尽管同时如此）：我想他用这些话责备他的生活方式。他的话并非出于愤怒，为表明这一点，作者先说了\Quote{充满圣神，}即祂的运作。\BibleQuote{现在你看，主的手临到你。}\BibleRef{宗 13:11}这并非报复，而是医治；因为他好像说：\Quote{不是我在做，而是天主的手。}注意多么谦逊！没有\BibleQuote{光，}如同保禄的情形，\BibleQuote{四面照射着他。}\BibleRef{宗 9:3}\BibleQuote{你要成为瞎子，}他说，\BibleQuote{暂时不见日光，}使他有机会悔改：因为我们从未发现他们希望藉更严厉（行使权威）而引人注目，即使是对敌人如此；对于自己团体内的人（他们用了严厉），且理由充分，但对于外人则不然；这样（信德的服从）才不至于显得是强迫和恐惧的结果。他瞎眼的证据，就是他\Quote{寻找人用手领他。}\BibleRef{宗 5:1}总督看见瞎眼被施加，\BibleQuote{当他看见所发生的事，就相信了：}而且他们两人不止是相信了这个，且\BibleQuote{对主的道理感到惊奇}\BibleRef{宗 13:12}：他看见这些不是空话，也不是诡计。注意他多么喜欢从老师那里接受教导，尽管他身居如此高位。而（保禄）没有对术士说：\Quote{你还不停止败坏}总督吗？若望离开他们回去的原因是什么？因为\Quote{若望，}经上说，\Quote{离开他们，回了耶路撒冷}\BibleRef{宗 13:13}：（他这样做）是因为他们要走更远的路；而且他是他们的助手，至于危险，却是他们（而非他）去承担。——再者，当他们来到培尔革时，他们匆忙经过其他城市，因为他们急于赶到首府安提约基雅。注意历史学家是多么简洁。\Quote{他们进了会堂，}他说，\Quote{在安息日那天}（14,15节）：以便预先为圣言预备道路。他们不先说话，而是在被邀请时才说；因为作为陌生人，他们被邀请这样做。如果他们不曾等待，就不会有讲道。这里我们第一次看到保禄宣讲。注意他的谨慎：在圣言已播种的地方，他继续前行；但在没有人（宣讲）的地方，他停留；正如他自己所写：\BibleQuote{是的，我这样竭力传福音，不在基督已被提起的地方。}\BibleRef{罗 15:20}这也是极大的勇气。确实，从一开始，一个奇妙的人！被钉十字架，准备好面对一切遭遇（\OriginalTerm{准备战斗}{\GreekText{παρατεταγμένος}}），他知道自己获得了多大的恩宠，并以同等的热忱回应。他对若望没有生气：因为这不是他的风格；但他坚持工作，他不畏惧，他不惊慌，即使被围困在众人之中。注意上智的安排是多么明智，保禄不应在耶路撒冷宣讲：单是听说他已成了信徒，这本身对他们就足够了；若他宣讲，他们绝不能忍受，他们对他的仇恨如此之深：所以他远走他乡，到不认识他的地方。但这样做很好，\BibleQuote{他们进了会堂，在安息日那天}当所有人都聚集时。\BibleQuote{读了法律和先知以后，会堂长派人到他们跟前说：诸位仁人兄弟，如果你们有什么劝勉人民的话，请说罢。}\BibleRef{宗 13:15}看他们如何不情愿地这样做，但此后就不再如此了。如果你们真的希望这样，就更需要劝勉了。\par

他首先证实了那个术士，显示了他是什么样的人；而他就是这样的人，标记显示了出来：\BibleQuote{你将失明，不见太阳}，这是他灵魂失明的标记：\BibleQuote{暂时} \BibleRef{宗 13:11}，他说，为要引导他悔改。但是，嗐，那种爱权之心！嗐，那种虚荣的贪欲！它如何颠覆并毁灭一切；使人起来反对自己的利益，反对彼此的救恩；使他们真是盲目、黑暗，甚至不得不寻求别人来牵手引导！但愿他们这样做，但愿他们寻求哪怕只是有人来牵手引导！但不然，他们不再忍受这个，他们把整个事情揽在自己手里。（这个恶习）不让任何人看见：如同迷雾和浓密的黑暗，它笼罩在他们身上，不让任何人看透它。我们有什么借口可以提出？我们这些人，为了一种邪恶的情感而克服另一种邪恶的情感（\Emph{见前第176页}），却不为了敬畏天主！例如，许多既淫荡又贪婪的人，因着他们的吝啬而给情欲戴上缰绳；反之，另一些这样的人为了享乐而蔑视财富。同样，那些两者兼具的人，因着虚荣的贪欲而克服了这两者，毫不吝惜地挥霍金钱，并实行节欲却毫无正当目的；还有另一些人，极其虚荣，却蔑视那种邪恶的情感，为了他们的恋情或为了钱财而忍受许多卑劣的耻辱；还有另一些人，为了满足他们的愤怒，选择承受无尽的损失，毫不关心这些，只求能遂行己意。然而，情欲能对我们做什么，敬畏天主却无法实现！我为什么提情欲？羞耻在人前能对我们做什么，敬畏天主却没有力量实现！我们做了许多正确或错误的事，出于在人前的羞耻感；但我们不畏惧天主。有多少人因顾及人的意见而羞愧，从而扔掉了金钱！有多少人错误地将其视为荣誉，只献身于朋友的服务，却有害于己！有多少人出于对友谊的尊重而羞耻地做了无数错事！既然情欲及对人意见的顾及都能驱使我们去做正确或错误的事，那么说\Quote{我们不能}就是徒劳的：我们能够，只要我们有心；而且我们应当有心。为什么你不能克服对光荣的爱，而别人却能克服，他们与你拥有同样的灵魂、同样的身体，具有同样的形体，过着同样的生活？思念天主，思念来自天上的光荣：拿现今的事物与之比较，你就会迅速远离这世俗的光荣。如果你确实渴望光荣，那就去渴慕那真正是光荣的。那带来羞辱的光荣算什么光荣？那迫使一个人渴望来自较低者的尊荣并需要这种尊荣的光荣算什么光荣？真正的荣誉是获得比自己伟大者的尊重。如果你确实贪恋光荣，你不如贪恋来自天主的光荣。如果你贪恋那光荣而轻视这世界的光荣，你就会看到这世界的光荣是多么卑贱。但如果你看不见那光荣，你也就不能看见这世界的光荣是多么污秽、多么可笑。因为正如那些被某个邪恶、丑陋无比的女人所迷惑的人，只要还爱着她，就看不见她的丑陋，因为情欲使他们的判断蒙上黑暗：此处也是如此。只要我们被这种情欲所控制，我们就不能察觉它是什么东西。那么，我们怎样才能摆脱它呢？想想那些为了光荣而花费了无数金钱的人，如今却毫无益处；想想那些死者，他们得到了什么光荣，而这光荣如今已无处可留，一切都消亡归于乌有：你要思想它只是一个名字，内中并无实质。请问，什么是光荣？你给我下一个定义。你会说：\Quote{为众人所钦佩。} 是公正地，还是不公正地？如果不公正，这就不是钦佩，而是\OriginalTerm{控诉}{\GreekText{κατηγορία}}、奉承和\OriginalTerm{毁谤}{\GreekText{διαβολή}}。但如果你说是公正地，那是不可能的：因为在民众中没有正确的判断；那些迎合他们情欲的人，才是他们所钦佩的人。如果你愿意（看看这证据），请注意那些把财产献给娼妓、车手、舞者的人。但你会说，我们指的不是这些人，而是那些公正正直、能行伟大高贵善举的人。但愿他们愿意，他们很快就能行善；但事实上，他们并不这样做。请问，现在谁称赞那公正正直的人？不，恰恰相反。还有什么比一个正义的人在做任何这样的善事时，却要从大众那里寻求光荣——如同一个技艺精湛的工匠在画皇帝的肖像时，却希望得到无知的赞美——更荒谬的吗？此外，一个人如果寻求从人而来的尊荣，他很快就会停止德性所命令的行为。如果他总是张着嘴等着他们的赞美，他就会做他们所希望的事，而不是他自己所希望的事。那么，我劝你做什么呢？你必须只仰望天主，只仰望从祂而来的赞美，实行一切祂所喜悦的事，追求（与祂同在的）善事，不要为任何来自人的事张着嘴等待；因为这败坏斋戒、祈祷和施舍，使我们的所有善行都归于无效。为了避免我们落到这地步，让我们逃避这种情欲。让我们只注视一件事：从天主而来的赞美，被祂接纳，从我们共同的主帅那里得到赞许；这样，我们就能在此生度德性生涯之后，与那些爱祂的人一同获得所应许的福分，因着我们的主耶稣基督的恩宠和仁慈，愿光荣、威能、尊威归于祂——与圣父和圣神——自今至永远，世世无尽。阿们。\par

\SourceRecord{Translated by J. Walker, J. Sheppard and H. Browne, and revised by George B. Stevens.；Nicene and Post-Nicene Fathers, First Series；Vol. 11.；Edited by Philip Schaff.；Buffalo, NY: Christian Literature Publishing Co.,；1889.；<http://www.newadvent.org/fathers/210128.htm>.}

% structure: p0001 source=h1 target=section reason=page-title

\section{\WorkTitle{宗徒大事录}讲道集第二十九篇}

\BibleRef{宗 13:16-17}\par

\BibleQuote{\Person{保禄}就站起来，打了一个手势，说：诸位以色列人和敬畏天主的人，请听！以色列人的天主，拣选了我们的祖先，当这百姓寄居\Place{埃及}时，就举扬了他们，以大能的手臂从那里领他们出来。}\par

请看巴尔纳伯让位给保禄——这难道不是理所当然的吗？——就是给他从塔尔索带来的那位；正如我们发现若望在一切场合都让位于伯多禄一样。然而巴尔纳伯比保禄更受人尊重：确实如此，但他们只着眼于共同的益处。\BibleQuote{那时保禄站起来}，圣经说——这是犹太人的习惯——\BibleQuote{并用手示意}。请看他是如何事先为自己的讲演铺路的：他首先赞扬了他们，并表明了对他们的深厚敬意，用\BibleQuote{敬畏天主的人}这样的话，才开始他的讲论。他没有说\Quote{你们这些归化者}，因为那是一个不利的称谓。\BibleQuote{这百姓的天主拣选了我们的祖先，并抬举这百姓}——请看，他称天主本身特别地是\Emph{他们的}天主，而祂是众人共同的天主；并显示出从一开始祂的恩惠是多么浩大，正如斯德望所做的那样。他们这样做是为了教导他们，现在天主也按照同样的惯例行事，派遣了自己的圣子；\BibleRef{路 20:13}正如（基督）自己在葡萄园的比喻中所做的——\BibleQuote{这百姓，}他说，\BibleQuote{当它寄居在埃及地时，祂抬举了它}——然而实际情况正好相反：确实如此，但是他们的人数增多了；而且，奇迹是因他们而行的：\BibleQuote{并用大能的手将他们领了出来。}关于这些在埃及所行的事（奇事），先知们不断地提及。请注意，他是如何略过他们灾难的时期，并未提及他们的过犯，而只提到天主的仁慈，将这些留给他们自己去反思。\BibleQuote{在旷野中约四十年之久，祂容忍了他们的习性。}\BibleRef{宗 13:18}然后是定居。\BibleQuote{祂在客纳罕地消灭了七个民族之后，将他们的土地分给他们为业。}\BibleRef{宗 13:19}时间很长，四百五十年。\BibleQuote{此后，祂赐给他们民长，约有四百五十年，直到先知撒慕尔。}\BibleRef{宗 13:20}他在这里表明天主在各时期对他们的安排是不同的。\BibleQuote{后来他们要求一位君王：}（仍然）没有一句关于他们忘恩负义的话，而是始终在讲论天主的仁慈。\BibleQuote{天主就把基士的儿子撒乌耳赐给他们，他是本雅明支派的人，作王四十年。}\BibleRef{宗 13:21}\BibleQuote{撤去他以后，又给他们兴起达味为君王，并为他作证说：我找到了叶瑟的儿子达味，他是一个合乎我心意的人，他要承行我的一切旨意。从这人的后裔中，天主照应许给以色列兴起了一位救主，耶稣。}\BibleRef{宗 13:22-23}基督出自达味，这可不是小事。然后若望为此作证：\BibleQuote{若望在祂来临以前，先向以色列众民宣讲了悔改的洗礼。若望将要完成他的职分时，说：你们以为我是谁？我不是祂。但是，看哪，在我以后要来一位，我连解祂脚上的鞋带也不配。}\BibleRef{宗 13:24-25}若望不仅作证（事实），而且（在这样的时候）当人们要将荣耀归给他时，他拒绝了：因为不去追求一个无人想给你的荣誉是一回事；而当众人都准备给予时，不仅拒绝它，而且以如此谦卑的态度拒绝，则是另一回事。\BibleQuote{诸位仁人弟兄，亚巴郎的子孙和你们中间敬畏天主的人，这救恩的话是传给我们的。因为耶路撒冷的居民和他们的首领，由于不认识祂，也不明白每安息日所诵读的先知的话，就应验了这些话，而定祂的罪。虽然查不出什么该死的罪，还是要求比拉多把祂处死。}\BibleRef{宗 13:26}我们在一切场合都发现他们着力表明这一点，即这恩惠是特别属于他们的，好使他们不致因认为自己与祂无关而逃离（基督），因为他们曾把祂钉在十字架上。\BibleQuote{因为不认识祂，}他说：因此这罪是出于无知。看他是如何温柔地为那些人（钉十字架者）辩护的。不仅于此：他还补充说，这样的事是必须发生的。怎么会这样呢？\BibleQuote{他们定了祂的罪，就应验了先知的话。}接着又是引用圣经。\BibleQuote{他们既应验了所有关于祂所记载的，就把祂从木架上取下来，放在坟墓里。但天主却使祂从死者中复活了。祂多日显现给那些从加里肋亚同祂上耶路撒冷的人，这些人如今在百姓前是祂的见证人——}\BibleRef{宗 13:29}证明祂复活了。\BibleQuote{我们现今也给你们传报喜讯：就是那对祖先的应许，天主已为我们这些做子孙的成就了，叫耶稣复活了，正如诗篇第二篇所记载的：\InnerQuote{祢是我的儿子，我今日生了祢。}论到天主叫祂从死者中复活，不再归于朽坏，曾这样说：\InnerQuote{我要将达味的可靠恩惠赐给你们。}因此又在另一篇说：\InnerQuote{祢必不让祢的圣者见到朽坏。}达味固然在自己的世代遵行了天主的旨意，就长眠了，归到他的祖先那里，并见了朽坏；但天主所复活的那位，却没有见朽坏。所以，诸位仁人弟兄，你们要知道：通过这一位，赦罪之道传报给你们；在一切事上，凡相信的，都因祂而成义，这些事是你们凭梅瑟的法律不能成义的。}\BibleRef{宗 13:32-39}请注意保禄在此的讲论是多么激昂：我们从未见伯多禄说过这样的话。然后他又加上了令人畏惧的话：\BibleQuote{所以，你们要小心，免得先知书上所说的临到你们：\InnerQuote{藐视的人哪！你们要惊奇，要灭亡！因为在你们的日子，我行事，是你们无论如何也不会相信的，即使有人告诉你们。}}\BibleRef{宗 13:40-41}\par

(\Emph{a}) 留意他如何轮番从当下之事和先知书中编织讲论的线索。于是，\Quote{按应许，从（这人的）后裔中}— \BibleRef{宗 13:23} : (\Emph{c}) 达味之名是他们所珍视的；那么，他的儿子作他们的王，岂不正是（一件可欲之事）吗？— (\Emph{b}) 然后他引荐若翰；接着又引先知，他论道：\Quote{他们定罪时便应验了}，又论道：\Quote{凡经上所写的}；然后宗徒们作复活的见证；接着是达味作证。因为若单取旧约的证明，它们似乎并不像这样有力；同样，后者离开前者也是如此；故此他使它们互相证实。\Quote{诸位仁人弟兄}等。\BibleRef{宗 13:26} 因为他们被恐惧所控制，如同已将他杀死，良心使他们成了外人（即宗徒们），于是他不像对待杀基督者那样对他们讲话，也不像将一种不属于他们而唯独属于自己的好处交到他们手中。(\Emph{d})\Quote{因为凡住在耶路撒冷的人和他们的首领}：意思就是说，不是你们，而是他们；并且，他又为那些人辩护：\Quote{因为他们不认识他，也不明白每安息日所宣读的先知的话，在定罪他的时候，便将它们应验了。}这是一个重大的指控：他们不断听闻却未留心。但这不足为奇：因为前面关于埃及和旷野所说的，已足以表明他们的忘恩负义。且看这位宗徒，也如同被圣神亲自感动的那样，不断宣讲苦难、埋葬。(\Emph{g})\Quote{将他从木架上取下来。}请注意，他们多么强调这一点。他谈到他的死法。此外，他们突出地提及比拉多，以便通过提及（判他罪的）法庭来证明苦难的（事实），但同时为了更严厉地控告那些（钉死他的人），因为他将交付给了一个外人。而且他没有说\Quote{他们曾投诉（反对他）}(\OriginalTerm{}{\GreekText{ἐνέτυχον}}，另作 \OriginalTerm{}{\GreekText{ἐντυγχάνει}})，而是说\Quote{他们虽找不到死因，却求将他杀死。}(\Emph{e})\Quote{他显现了}，他说，\Quote{多日给那些从加里肋亚与他一同上耶路撒冷的人。}\BibleRef{罗 11:2} 代替**他说，\Quote{这些人就是他向百姓的见证人}，即\Quote{那些从加里肋亚与他一同上耶路撒冷的人。}然后他引达味和依撒意亚作证。\Quote{那忠实的（慈爱）}，那永久的（慈爱），那些永不朽坏的。(\Emph{h}) 保禄极其爱他们。请注意，他不详述祖先的忘恩负义，而是将他们所当惧怕的摆在他们面前。诚然，斯德望这样做很有理由，因为他即将被处死，并非在教导他们；他向他们显示，法律甚至现在即将被废除：\EditorialNote{宗 7} 但保禄却不是这样；他只是威胁并使他们恐惧。(\Emph{f}) 他不在这些事上久留，因为他视这话语当然被相信；也不详述他们惩罚的重大，不攻击他们锺爱的事物——通过显示法律即将被抛弃——而是专注于对他们有益的事（告诉他们），他们若顺从，将有大福；若悖逆，将有大祸。\par

但让我们重新审视刚才所说的。\Quote{以色列人哪}等。\BibleRef{宗 13:16} 那么，他说道，祖先领受了应许；你们领受了现实。(\Emph{j}) 且注意，他从未提及他们的善行，只提及天主的恩惠：\Quote{他拣选；提升；容忍他们的行为}：这些对他们来说并非可夸之事：\Quote{他们求，他赐予。}但达味他却称赞（且仅称赞他），因为基督要从他而来。\Quote{我找到了叶瑟的儿子达味，一个合我心意的人，他将遵行我的一切旨意。}\BibleRef{宗 13:22} (\Emph{i}) 也请注意：他提及达味时是带着称赞的：\Quote{达味在承行天主的旨意以后}，正如伯多禄——因为那时福音尚在开端——提及他时说道：\Quote{请允许我放胆对圣祖达味讲论。}\BibleRef{宗 2:29} 而且，他没有说\Quote{死了}，而是说\Quote{归到他列祖那里。}(\Emph{k})\Quote{从这人的后裔中}等。\Quote{当若翰}，他说，\Quote{在他进入以前先宣讲了}——他称\Emph{进入}意指降生成人——\Quote{给全以色列民宣讲悔改的洗礼。}\BibleRef{宗 13:23} 同样，若望在写他的福音时，也经常求助于他，因为他的名字在天下各处深受思念。并且请注意，他并没有凭自己说\Quote{从这人的后裔中}等，而是引用了若翰的见证。\par

\BibleQuote{诸位仁人弟兄，亚巴郎的子孙}——他也以他们的父亲称呼他们——\BibleQuote{这救恩之道是传给你们的}。\BibleRef{宗 13:26} 这里的措辞，\BibleQuote{传给你们}，并不是指传给（你们）犹太人，而是赋予他们权利，使他们与那些胆敢谋杀的割绝。他所加添的也清楚显明了这一点。\BibleQuote{因为，}他说，\BibleQuote{耶路撒冷的居民，因为不认识祂，}\BibleRef{宗 13:27} 你将会说，他们怎能无知，若翰曾告诉他们？这有什么奇怪，既然他们如此，而先知们不断向他们高呼？接着是另一项指控：\BibleQuote{虽然找不出死罪，}在这件事上无知毫无关系。因为让我们假设，他们不认祂为基督：为何他们也杀死祂？并且\BibleQuote{他们要求比拉多把祂处死}，他说。\BibleRef{宗 13:28} \BibleQuote{他们完成了所有关于祂的记载之后，}\BibleRef{宗 13:29} 注意他如何强调这件事是天主计划的整体。看，他们说什么来说服人？说祂被钉十字架？为什么，还有什么比这更不可信的呢？说祂被埋葬——在那些应许祂成为救恩的人手中？那位被埋葬的赦罪，甚至比法律还有能力？并且（注意），他并没有说，你们不能，而是\BibleQuote{你们不能从梅瑟法律中成义。}\BibleRef{宗 13:39} \BibleQuote{凡信的，}他说道：无论他是谁。因为那些（规条）毫无用处，除非有某种益处。这就是为什么他后来引入赦免：并显示它更伟大，当事情（否则）不可能时，它仍然成就。\BibleQuote{谁是他的见证人，}他说，\BibleQuote{向百姓，}——那杀害祂的百姓。若不是被一种神圣力量坚固，他们绝不会这样做：因为他们绝不会向那些嗜血的人，即杀害祂的人，作这样的见证。但是，\BibleQuote{天主已叫耶稣复活了：今天，}他说，\BibleQuote{我生了你。}\BibleRef{宗 13:33} 是的，其余的当然随之而来。为什么他没有引用一些经文来说服他们罪过得赦免是因祂？因为对他们来说，重要的首先是证明祂复活了；这一点被承认后，其他都是毋庸置疑的。\BibleQuote{通过这人，}不，通过祂，\BibleQuote{是赦罪。}\BibleRef{宗 13:38} 此外，他渴望引导他们盼望这件大事。那么，祂的死不是被舍弃，而是应验预言。——至于其他的，他提醒他们历史事实，他们因无知遭受无数痛苦。并在结尾暗示这一点，说：\BibleQuote{看，你们轻蔑的人，且观看。}注意，这很严厉，他缩短了。他说，不要让你们遭遇那为别人所说的：\BibleQuote{我要行一件工作，你们决不相信，即使有人告诉你们。}\BibleRef{宗 13:41} 不要惊奇它似乎难以置信：这件事从一开始就预言了——（不会被相信的）。\BibleQuote{看，你们轻蔑的人，}指那些不信复活的人。\par

这也可以有理由对我们说：\Quote{看，你们这些轻蔑的人。}因为教会确实处于非常糟糕的境地，尽管你们认为她的事务平安无事。其祸害在于，当我们受这么多罪恶折磨时，我们甚至不知道自己有任何罪恶。\Quote{你说什么？我们拥有我们的教堂、我们的教会财产和一切，礼仪举行，会众每天来教堂。}不错，但不应该从这些事情来判断教会的状况。那么从什么判断呢？是否有虔诚，我们每天是否带着益处回家，是否收获一些果实，无论多少，是否我们不仅仅出于习惯和形式上的履行责任（\OriginalTerm{形式上的履行}{\GreekText{ἀφοσιούμενοι}}）而做。有谁因每天参加礼仪整整一个月而变成了更好的人？问题的关键就在这里：否则，那似乎表明（教会）繁荣兴旺之事，实际上却表明了一种糟糕的状况，当这一切都做了，却没有产生任何结果。但愿天主（仅此而已），没有产生任何结果；但事实上，就目前情况而言，结果甚至更糟。你们从礼仪中得到了什么果实？如果你们确实从中获得了益处，你们早就应该过着真正智慧的生活（\OriginalTerm{真正的智慧}{\GreekText{τῆς} \GreekText{φιλοσοφίας}}），有那么多先知每周两次向你们讲论，那么多宗徒和福音作者，都在阐述救恩的教义，并非常精确地将那能正确塑造品性的事摆在你们面前。士兵通过操练，他的战术更加完善；摔跤手通过经常去健身房，摔跤技术更加熟练；医生通过跟随老师，变得更加准确，懂得更多，学到更多；而你呢——你得到了什么？我不是对仅仅加入教会一年的人说话，而是对那些从幼年就参加礼仪的人说话。你以为虔诚就是经常去教堂\OriginalTerm{参加聚会}{\GreekText{παραβάλλειν} \GreekText{τῇ} \GreekText{συνάξει}}吗？这毫无意义，除非我们为自己收获一些果实：如果我们从教堂的聚会中没有为自己收集（\OriginalTerm{聚集}{\GreekText{συνάγωμεν}}）些什么，那么待在家里反而更好。因为我们的祖先为我们建造了教堂，不仅仅是为了把我们从私人住所聚集起来，彼此相见：因为这在市场、浴场和公共游行中也能做到——而是为了把学习者和教师聚集在一起，使一方因另一方而变得更好。在我们这里，这一切都变成了纯粹的习惯性常规和形式上的履行职责：一件我们习以为常的事；仅此而已。复活节来临，然后是极大的骚动、极大的喧闹和拥挤——我宁愿不称他们为人，因为他们的行为不是常人所为。复活节过去，骚动平息，但随之而来的安静再次结不出善果。\Quote{守夜和圣歌咏唱。}——从这些中得到了什么？不，情况更糟。许多人这样做仅仅是出于虚荣。试想，看到这一切就像（大量的水）倒进有漏洞的木桶，这一定会让我多么痛心！但你一定会对我说，我们知道圣经。那又怎么样？如果你用行为实践圣经，那才是收获，那才是益处。教会是一个染缸：如果你们一次又一次地离开这里而没有染上任何颜色，那么不断来这里有什么用？哎呀，祸害反而更大。你们中谁曾在他从祖先那里继承的习俗上添加过什么？例如：某人有纪念他母亲、或妻子、或孩子的习惯；不管我们是否告诉他，他都这样做，是受习惯和良心的驱使。你问，这使你反感吗？千万不要：相反，我全心喜悦：只是，我希望他也从我们的讲论中获得一些果实，并且习惯所产生的效果，在我们（你们教师）身上也能产生同样的效果——即养成另一种习惯。否则，如果你们依旧停留于同一状态，如果教会的礼仪对你们没有益处，我为什么白费力气、说无用的废话呢？不，你们会说，我们祈祷。那又怎么样？\Quote{不是凡向我说\InnerQuote{主啊，主啊}的人，就能进入天国；而是那承行我在天之父旨意的人。} \BibleRef{玛 7:21}许多次，我决定保持沉默，因为看到你们从我的话中得不到益处；或者也许有一些益处，但因为我贪得无厌和强烈的渴望，我和那些疯狂追求财富的人一样受到影响。因为正如他们，无论得到多少，都认为一无所有；同样，我因热切渴望你们的救恩，直到见你们有了良好进步，才认为有所成就，因为我极其渴望你们能达到顶峰。我希望情形如此，是我的急切有错，而非你们的懒惰；但我担心我的猜测太正确了。因为你们必须确信，如果在这漫长的时间里产生了任何益处，我们现在早就应该停止讲论了。在那种情况下，你们就不需要言语了，因为在已经讲过的那些话中，对你们已经说得足够了，而且你们自己也能纠正别人了。但事实是，我们仍然有必要向你们讲论，这只能说明你们的事情并未处于完美状态。那么，我们要实现什么呢？因为不能仅仅指责。我恳求并劝告你们，不要以为进入教堂就够了，而是也要离开这里，带走一些东西，一些医治你们自己疾病的药物；如果不是从我们，至少从圣经中，你们有适合每个人的疗法。例如，有人易怒吗？让他留心读圣经，他肯定会在历史或劝诫中找到（疗法）。在劝诫中，当它说：\Quote{怒气的控制是他的毁灭} \BibleRef{德 1:22}；和\Quote{暴躁的人不体面} \BibleRef{箴 11:25}；以及诸如此类；还有\Quote{多言的人必不亨通} \BibleRef{咏 139:11}；基督又说：\Quote{凡向弟兄发怒而无故的} \BibleRef{玛 5:22}；先知又说：\Quote{你们动怒，却不要犯罪} \BibleRef{咏 4:4}；以及\Quote{他们的怒气可咒，因为猛烈} \BibleRef{创 49:7}。在历史中，如你们听到法老充满大怒，以及亚述人。同样，有人被贪财所俘虏吗？让他听，\Quote{没有比贪财的人更邪恶的了，因为这种人甚至出卖自己的灵魂} \BibleRef{德 9:9}；以及基督说：\Quote{你们不能事奉天主而又事奉玛门} \BibleRef{玛 6:24}；宗徒说：\Quote{贪财是万恶之源} \BibleRef{弟前 6:10}；先知说：\Quote{钱财流入，不要倾心其上} \BibleRef{咏 61:10}；以及其他许多类似的言语。从历史中你们听到革哈齐、犹达斯、众经师，以及\Quote{礼物蒙蔽智者的眼目} \BibleRef{申 16:19}。另有人骄傲吗？让他听，\Quote{天主抵挡骄傲的人} \BibleRef{雅 4:6}；以及\Quote{骄傲是罪恶的开端} \BibleRef{德 10:14}；和\Quote{凡心里高傲的，在上主面前为不洁} \BibleRef{箴 16:5}。在历史中，还有魔鬼和其余一切。总之，既然不可能一一叙述，就让每个人从神圣圣经中为自己的创伤选择疗法。\par

所以，即使不能一下子洗净全部，至少也要洗净一部分——今天一部分，明天一部分，最终全部洗净。关于悔改、告解、施舍、正义、节制以及所有其他事情，你也会发现许多例子。因为 \Quote{这一切，}宗徒说，\Quote{都是为规劝我们而写的。}\BibleRef{格前 10:11}既然圣经的一切论述都是为了规劝我们，我们就应当如何地专注聆听。为何要徒然欺骗自己？我担心我们也应验了那句：\Quote{我们的岁月在虚空中消逝，我们的年岁匆匆而过。}\BibleRef{咏 77:33}有谁因听我们宣讲而放弃了戏院？有谁放弃了贪财？有谁变得更乐于施舍？我愿知道这些，不是为了虚荣，而是为了看到我劳苦的果实清晰可见，从而获得更大的热忱。但如今的情形，当我看到尽管教义的雨水一波接一波倾注在你们身上，我们的庄稼却仍停留于同样的尺寸，植株并未长高丝毫，我怎能再着手工作？转眼间，扬场的时候到了，（且）祂拿着簸箕。我怕那全是糠秕；我怕我们全被投进火炉。夏天已过，冬天已至；我们坐着，无论老少，都被自己的邪恶情欲所俘虏。不要对我说：我没有犯奸淫。因为即使你不是奸淫者，却是贪财者，那又有什么益处呢？对落入网罗的麻雀而言，它不是每处都被紧紧捉住，而只是一只脚被夹住，那又有什么关系？它终究是只失丧的鸟；它在网中，即使翅膀自由，只要脚被夹住，便无济于事。同样，你不是被奸淫捉住，而是被爱财捉住；但你仍然被捉住了；关键不在于你\Emph{如何}被捉住，而在于你\Emph{被捉住}。不要让年轻人说：我不爱钱财。但或许你是奸淫者？那又有什么益处呢？事实上，所有情欲不可能在同一人生阶段同时攻击我们；它们被划分并区分开来，这也是天主的仁慈，免得它们同时袭击我们，变得不可战胜，使我们的搏斗更加困难。这是多么可悲的懒散——连逐一战胜情欲都做不到，反而在生命的每个时期都被击败，却为那些（暂且不提）不是因我们真心努力、而仅仅因人生阶段而蛰伏的情欲而自夸！看看那些战车御者，难道你没看到他们对自己在训练、劳动、饮食及其他一切方面是何等谨慎严格，以免从战车上摔下被（缰绳）拖行？——看艺术是多么美妙。常常一个强壮的男人都制伏不了一匹马；但一个学会了技艺的小男孩却常常能将双马制服，轻松地牵着它们去任何地方。据说在印度，一头巨大的大象会屈服于一个十五岁的少年，被他极其轻松地驾驭。我说这一切是为了什么？为了说明，若通过练习和训练，我们能\OriginalTerm{勒住}{\GreekText{ἄγχομεν}}甚至大象和野马，那么里面的情欲就更不在话下了。为何我们一生中（在每次交锋中）不断失败？我们从未练习过这门艺术；从未在闲暇时、在无比赛之际，与我们自己讨论什么对我们有益。我们从不在比赛开始前出现在战车位置上。因此我们在那里出了丑。我曾多次说过：让我们在试炼之前，先在自己的家人身上练习？在家里，我们常被自己的\OriginalTerm{仆人}{\GreekText{παῖδας}}激怒，让我们在那里平息愤怒，以便在与朋友交往时能够轻易控制它。同样，对于其他所有情欲，如果我们事先练习过，就不会在争战本身中出丑。但现在，对于其他事情——艺术和竞技场的技艺——我们有工具、练习和训练，但对于德行却没有。农夫不敢去修剪葡萄树，除非他先练习过栽种技术；舵手不敢坐在舵旁，除非他先熟练练习过；而我们，在所有方面都未经练习，却想赢得头奖！最好是保持沉默，最好在能够\OriginalTerm{安抚}{\GreekText{κατεπᾴδειν}}我们里面的野兽之前，不在言语和行为上与任何人交往。我所说的野兽，因为愤怒和色欲向我们来战时，难道不比任何野兽的攻击更可怕吗？当心，不要带着这些野兽闯入市场（\OriginalTerm{不要闯入市场}{\GreekText{Μὴ} \GreekText{ἐμβάλῃς} \GreekText{εἰς} \GreekText{ἀγοράν}}），直到你把嚼子牢牢套在它们的嘴上，驯服它们并使它们温顺。那些在市场牵着驯服狮子的人，你没看到他们从中获得多大的收益、多大的赞美吗？因为他们在无理性的野兽身上产生了如此的驯服——但是，如果狮子突然狂性大发，它如何吓得市场里所有人四散奔逃，那么牵狮子的人自己也会陷入危险，如果造成他人的死亡，那也是他的作为。因此，你也要先驯服你的狮子，然后才牵着它走，不是为了收钱，而是为了获得无可比拟的收益。因为没有一样东西能比得上温和，它对拥有者和受者都是极其有益的。让我们追求这一点，以便行走在德行之道上，尽一切努力跑完赛程，最终得享永恒的美善，赖我们的主耶稣基督的恩宠和仁慈，愿光荣、权能、尊崇归于祂及圣父和圣神，从现在直到永远，世世无穷。阿们。\par

\SourceRecord{Translated by J. Walker, J. Sheppard and H. Browne, and revised by George B. Stevens.；Nicene and Post-Nicene Fathers, First Series；Vol. 11.；Edited by Philip Schaff.；Buffalo, NY: Christian Literature Publishing Co.,；1889.；<http://www.newadvent.org/fathers/210129.htm>.}

% structure: p0001 source=h1 target=section reason=page-title

\section{宗徒大事录第三十篇讲道}

\BibleRef{宗 13:42}\par

\Quote{当他们出去时（领受本：\InnerQuote{从犹太人的会堂出来}），他们（外邦人）恳求这些话能在下一个安息日对他们宣讲。}\par

你们看出保禄的智慧了吗？他不仅当时获得了钦佩，还在他们心中激起了再次聆听的渴望，同时他在讲话中像\OriginalTerm{撒下了一些种子}{\GreekText{εἰπών} \GreekText{τινα} \GreekText{σπέρματα}}，并没有解决（提出的问题）或将主题推向结论，他的计划是引起他们的兴趣，赢得他们的好感，而不是通过一下子把所有东西都倾注给（那些初听他话的人）而使他们变得冷漠和漠不关心。他告诉他们事实：\BibleQuote{藉着这个人，罪过的赦免已向你们宣告}，但他没有说明如何成就。\BibleQuote{散了会，许多犹太人和虔敬的归化者跟随着保禄和巴尔纳伯}——从此以后他把保禄放在前面——\BibleQuote{二人对他们讲话，劝他们继续留在天主的恩宠中。}\BibleRef{宗 13:43} 你们看到那热忱有多大吗？经上说，他们\BibleQuote{跟随着}他们。为什么不立刻给他们施洗？不是适当的时候：需要说服他们，使他们坚定地留在其中。\BibleQuote{到了下一个安息日，几乎全城的人都聚集来听天主的话。}\BibleRef{宗 13:44} \BibleQuote{但犹太人看见这众多的人，就满心嫉妒，反驳保禄所说的话，并且毁谤。}\BibleRef{宗 13:45} 看，恶毒在伤害别人时也伤害了自己；这使宗徒们更加显眼——那些人提出的反对。起初是他们自己恳求讲道（如今却反对他们）：经上说，\BibleQuote{反驳}，\BibleQuote{并毁谤}。何等的轻率！\BibleQuote{那时，保禄和巴尔纳伯却放胆说：天主的道本来必须先传给你们，但因你们拒绝它，断定自己不配得永生，看，我们就转向外邦人。}\BibleRef{宗 13:46} 你们看见他们如何因争执而更加扩展了宣讲，（宗徒们）如何更加投身于外邦人，藉此为自己辩护，在他们自己人（在耶路撒冷）面前洗脱一切责备了吗？（\Emph{c}）看，他们因\BibleQuote{嫉妒}而造成了伟大而非他们所预料的事：他们使宗徒们放胆发言，并走向外邦人！因此他说：\BibleQuote{保禄和巴尔纳伯放胆说道}。他们将要向外邦人出发：但注意那勇气是带着节制的：因为如果伯多禄曾为自己辩护，这些人更需要一个辩解，因为没有人召唤他们去那里。\BibleRef{宗 11:4} 但通过说\BibleQuote{先给你们}，他表明向他们传道也是他的本分；而通过说\BibleQuote{必要}，他表明也必要传给他们。\BibleQuote{但既然你们拒绝它}——他没有说\BibleQuote{祸哉你们}和\BibleQuote{你们受了惩罚}，而是\BibleQuote{我们转向外邦人}。这勇气充满了极大的温柔！（\Emph{a}）他也没有说\BibleQuote{你们不配}，而是\BibleQuote{断定自己不配。看，我们转向外邦人。因为主这样吩咐我们说：我已立你作外邦人的光，使你成为救恩，直到地极。}\BibleRef{宗 13:47} 为使外邦人听到这些话不至于受伤，好像情形是，如果犹太人热切，他们自己就不会得到福分，因此他引用了预言，说：\BibleQuote{外邦人的光}，和\BibleQuote{成为救恩直到地极。外邦人听见}\BibleRef{宗 13:48}——这既使他们更加欢欣，因为他们看到情形是，那些人本来有权先听，而他们自己却享受了福分，同时也更刺痛了那些人——\BibleQuote{外邦人}，经上说，\BibleQuote{听见}这话，\BibleQuote{就欢喜，赞美主的道；凡预定得永生的人都信了}，即：分别为天主的人。注意他如何显出这益处的迅速：\BibleQuote{主的道传遍了那整个地区}\BibleRef{宗 13:49}，这里用了\OriginalTerm{被传播}{\GreekText{διεφέρετο}}，而不是\OriginalTerm{被运送}{\GreekText{διεκομίζετο}}，即\Quote{被携带或传递（通过该地区）}。（\Emph{d}）\BibleQuote{但犹太人煽动虔诚尊贵的妇女和城中的首领，迫害保禄和巴尔纳伯，将他们赶出境外。}\BibleRef{宗 13:50} \BibleQuote{虔诚的妇女}，（\Emph{b}）即归化妇女。他们不止于\BibleQuote{嫉妒}，还加上了行动。（\Emph{e}）你们看见他们因反对宣讲而成就了什么？他们使这些人（\BibleQuote{尊贵的妇女}）陷入何等的不名誉？\BibleQuote{二人却对着他们跺下脚上的尘土，往\Place{依科尼雍}去了。}\BibleRef{宗 13:51} 现在他们使用了基督所吩咐的那可怕的记号：\BibleQuote{若不接待你们，就把脚上的尘土跺下去}\BibleRef{玛 10:14}；但他们这样做并非出于轻率，而是因为他们被那些人赶逐。这对门徒并没有伤害；相反，他们更加坚持这道：\BibleQuote{门徒们满心喜乐，并被圣神充满}\BibleRef{宗 13:52}，因为老师的苦难并不削弱他的勇气，反而使门徒更加勇敢。\par

\Quote{他们二人一同进了依科尼雍的犹太人会堂。} \BibleRef{宗 14:1} 他们再次进入会堂。看他们远没有变得更加胆怯！说了\Quote{我们转向外邦人}后，他们（通过进入会堂）仍然极大地加强了自己的辩白（对犹太弟兄）。\Quote{以致}经上说，\Quote{有大群犹太人和希腊人都信了。}因为很可能他们也向希腊人讲论。\Quote{但那些不信的犹太人煽动外邦人，使他们的心恶意对待弟兄。} \BibleRef{宗 14:2} 现在他们（连同自己）也开始煽动外邦人，因为他们自己不够。那么宗徒们为什么不离开那里？他们没有被迫离开，只是受到攻击。\Quote{二人住了多日，靠着主放胆讲论，主借着他们的手施行神迹奇事，为祂恩宠的道作证。} \BibleRef{宗 14:3} 这使他们放胆；或者更确切地说，他们的放胆实在是由他们自己真诚的善意所引发——因此他们长期没有行神迹——而听者的归化是神迹的效果，尽管他们的放胆也贡献了一些。\Quote{城里的群众分裂了：有的附同犹太人，有的附同宗徒。} \BibleRef{宗 14:4} 这种分裂非同小可。这正是主所说的：\Quote{我来不是带和平，而是带刀剑。} \BibleRef{玛 10:34} \Quote{当外邦人和犹太人连同他们的官长一齐起来，要凌辱他们，用石头砸他们时，他们知道了，就逃往吕考尼雅的吕斯特辣和德尔贝两城，以及周围地区，在那里传福音。} \BibleRef{宗 14:5} 又如，仿佛他们有意在宣讲增长后扩展宣讲，他们再次派他们出去。看，在一切场合，迫害成就大善，挫败迫害者，使被迫害者声名显赫。因为到了吕斯特辣，他行了一个大神迹，使跛子起来。\Quote{在吕斯特辣坐着一个双脚无力的人，从母胎就是瘸子，从来没有走过。他听保禄讲论；保禄定睛看他，见他有信心可得痊愈，就大声说}——为什么大声？为叫群众相信——\Quote{你起来，两脚站直。}（第8-9节）但请注意，经上说，他留意听保禄所讲的话。你注意到此人思想的高超（\OriginalTerm{哲学}{\GreekText{φιλοσοφίαν}}）了吗？他因听讲的恳切而完全没有被他的跛足所阻碍（\OriginalTerm{受阻}{\GreekText{παρεβλάβη}}）。\Quote{保禄注视他，看出}经上说，\Quote{他有信心可得痊愈。}他心中早已有此意向。然而在其他人的情形中，恰恰相反：因为他们先是身体得医治，然后才被着手治疗灵魂，但这个人却不是这样。在我看来，保禄看到了他的灵魂。\Quote{他就跳起来，}经上说，\Quote{而且行走。} \BibleRef{宗 14:10} 他的跳跃是痊愈的明证。\Quote{众人看见保禄所做的事，就大声用吕考尼雅话\InnerQuote{神取了人形降到我们这里了}，他们称巴尔纳伯为朱庇特，称保禄为墨丘利，因为他是主要发言人。于是城门前朱庇特的司祭牵着牛，拿着花环，要与群众一同献祭。} \BibleRef{宗 14:11} 但这意图尚未显现，因为他们用本地话说：\Quote{神取了人形下降到我们这里了。}因此宗徒那时没有对他们说什么。但当他们看见花环时，他们就出来，撕裂衣服，\Quote{巴尔纳伯和保禄二位宗徒听见了，就撕裂衣服，冲进群众中间，喊着\InnerQuote{诸位，为什么做这些事？我们也是与你们有同样性情的人。}}（第14-15节）看，他们如何在一切场合清除了对光荣的贪恋，不但不贪求，甚至当光荣被献上时也拒绝：正如伯多禄也说过：\Quote{你们为什么注视我们，好像是我们凭自己的能力和圣德使他行走呢？} \BibleRef{宗 3:12} 所以这些宗徒也说了同样的话。若瑟关于梦也说过：\Quote{解梦不是出于天主吗？} \BibleRef{创 60:8} 达尼尔同样：\Quote{这奥秘启示给我，并非因我智慧超于众人。} \BibleRef{达 2:30} 保禄处处说这话，如他说：\Quote{这些事谁能当得起？并不是我们凭自己能够承担什么，好像出于自己，而是我们的能力出于天主。} \BibleRef{格后 2:16} 但让我们再回顾一下所说的话。\par

（摘要。）\BibleQuote{他们出去以后}等等。\BibleRef{宗 13:42}。不但群众被他们吸引，而且怎样？他们恳求将同样的话再对他们讲，并用行动表明他们的热忱。\BibleQuote{散会以后}等等。\BibleRef{宗 13:43}。请看宗徒们处处劝勉，不但接纳人、讨好他们，而且\BibleQuote{对他们讲论}，经上说，\BibleQuote{劝他们常在天主的恩宠中。但犹太人}等等。\BibleRef{宗 13:45}。为什么他们以前没有反驳？你们是否注意到，他们总是受激情驱使？他们不但反驳，而且还诽谤。的确，恶意无所不为。但请看这何等的胆量！\BibleQuote{必须，}他说，\BibleQuote{先把天主的圣言讲给你们听，但你们弃绝它}——\BibleRef{宗 13:46}。这并非冒犯（虽然）事实上他们对待先知也是如此：\BibleQuote{不要对我们说话，}他们说，\BibleQuote{用言语}——\BibleRef{依 30:10}：\BibleQuote{但你们弃绝它}——他说的是\Quote{它}而不是\Quote{我们}：因为你们的冒犯不是针对我们。为了不让人以为这是他们虔敬的表现（他说的）\BibleQuote{你们断定自己不配，}因此他先说\BibleQuote{你们弃绝它，}然后说\BibleQuote{我们就转向外邦人。}这说法充满温柔。他没有说\Quote{我们放弃你们}，而是说\Quote{这样我们可能还能再转向这里}：这也不是你们冒犯的结果，\BibleQuote{因为主曾这样命令我们。}——\BibleRef{宗 13:47} \Quote{那么你们为什么没有这样做？}外邦人确实需要听，而且这并不在你们之前：你们自己的行为造成了\Quote{在你们之前。}\BibleQuote{因为主曾这样命令我们：我已立你作外邦人的光明，使你成为救恩，}即直到救恩的知识，不仅是外邦人，也是所有的人，\BibleQuote{直到地极——凡预定得永生的人} \BibleRef{宗 13:48}。这也证明他们接纳这些外邦人是符合天主心意的。但\Quote{预定，}不是出于必然：\BibleQuote{祂所预知的人，}宗徒说，\BibleQuote{祂也预定。} \BibleRef{罗 8:29} \BibleQuote{主的话}等等。\BibleRef{宗 13:49} 他们的道理不再仅在城内传播，也传到（整个）地区。因为当外邦人听见后，他们不久也归信了。\BibleQuote{但犹太人唆使虔诚的妇女，挑起迫害}——注意甚至妇女所做的事，也是他们主使的——\BibleQuote{将他们，}经上说，\BibleQuote{赶出境外} \BibleRef{宗 13:50}，不仅是城市。然后，更可怕的是，\BibleQuote{他们对着他们跺掉脚上的尘土，来到\Place{依科尼雍}。但门徒，经上说，都充满喜乐和圣神。}（第51-52节）教师遭受迫害，而门徒却喜乐。\par

\BibleQuote{他们这样讲论，以致一大群人}等等。\BibleRef{宗 14:1} 你们注意福音的本性，它有多大的德能？\BibleQuote{使他们的心受恶影响，}经上说，\BibleQuote{敌对弟兄们：} \BibleRef{宗 14:2} 即诽谤宗徒，对他们提出无数控告：（这些人，本是单纯，他们）\Quote{使受恶影响，}使他们行恶。再看，他处处将一切归于天主。\BibleQuote{住了很久，}他说，\BibleQuote{他们靠着主放胆讲论，主为祂恩宠的言语作证。} \BibleRef{宗 14:3} 不要以为这（\Quote{作证}）有何贬损（主的神圣威严）：\BibleQuote{祂在般雀比拉多面前作证，}经上说。\BibleRef{弟前 6:13} 然后那胆量——\BibleQuote{并且藉他们的手施行征兆和奇事。}他这样说，如同关于自己的民族。\BibleQuote{城里的群众}等等。（第4-5节）因此他们没有等待，而是看到攻击他们的意图，就逃走了，从未激起他们的愤怒，\BibleQuote{往\Place{吕考尼雅}的城市，即\Place{吕斯特辣}、\Place{德尔贝}及周围地区去。} \BibleRef{宗 14:6} 他们逃到乡间，不仅城市。——请看外邦人的淳朴和犹太人的恶毒。他们以实际行动表明自己配听：仅仅因奇迹就如此尊敬他们。一类人尊敬他们如神，另一类人迫害他们如恶徒：而且（那些人）不仅不因宣讲而反感，反而说什么？\Quote{神取了人形，降临到我们这里了；}但犹太人却反感。\BibleQuote{他们称\Person{巴尔纳伯}为则乌斯，称\Person{保禄}为赫尔默斯。}（第11-12节）我想\Person{巴尔纳伯}也是仪表堂堂的人。这是一种新的考验，来自过度的热情，而且非同小可：但由此也显示了宗徒的德行，以及他们如何在一切事上将一切归于天主。\par

让我们效法他们：让我们认为没有什么是属于我们自己的，即使信德本身也不是我们的，而更是天主的。他说：\Quote{因为你们得救是靠着恩宠，借着信德，而且这也不是出于你们自己，而是天主的恩赐。}\BibleRef{弗 2:8} 因此，我们不要自视过高，也不要自高自大，因为我们不过是人，是尘土和灰烬，是烟雾和影子。请问，你为什么自视过高？你施舍了财物，散尽了自己的家产？那又怎样？想想，如果天主没有使你富有呢？想想那些贫穷的人，或者更确切地说，想想有多少人不仅舍了财产，还舍了身体，并且在无数牺牲之后，仍觉得自己是可悲的！你是为自己而舍，基督（不是为自己，而是）为你而舍：你不过是还了债，基督并不欠你什么。—— 看看未来之不确定，要\Quote{不可心高气傲，反要战战兢兢}\BibleRef{罗 11:20}；不要因自夸而减损你的美德。你想做真正伟大的事吗？绝不要让你心中产生自己伟大的想法。但你是贞女？那（福音中）的贞女也是如此，但因她们的残忍和无情，她们并未从贞洁中获得任何益处。\BibleRef{玛 25:12} 没有比谦卑更好的了：谦卑是一切美德的母亲、根源、乳母、基础和纽带：没有谦卑，我们便是可憎的、可咒的、污秽的。比如说——让一个人复活死人、医治瘸子、洁净癞病人，却自满骄傲，那么没有比这更可咒、更不虔敬、更可憎的了。不要认为任何事是出于你自己。你有口才和教导的恩宠吗？不要因此认为你比别人多什么。正因如此，你更应当谦卑，因为你得到了更丰厚的恩赐。因为那多得赦免的，爱得也更多\BibleRef{路 7:47}：如果是这样，那么你也应当谦卑，因为天主越过其他人，眷顾了你。因此你应当害怕：因为这常常成为你毁灭的原因，如果你不谨慎的话。你为何自视过高？因为你用言语教导？但这很容易，用口舌讲论哲学：用你的生活教导我：那是最好的教导。你说应当节制，并为此高谈阔论，滔滔不绝地施展口才？但有人会对你说：\Quote{那位用行为教导我的人比你好}——因为那些用言语讲述的教训，不如用事实教导的更能深入人心：如果你没有行为，你不仅没有用言语使他受益，反而更伤害了他——\Quote{你最好是沉默。}为什么？\Quote{因为你向我提出的事是不可能的：因为我想，如果你对这事说了这么多，却做不到，那我更是情有可原了。}为此先知说：\Quote{但天主对恶人说：你为什么传述我的法令？}\BibleRef{咏 59:16} 因为这是一件更坏的事，当一个人言语教导很好，却用行为反驳教导时。这成了教会中许多恶事的起因。所以请宽恕我，我恳求你们，让我在这邪恶的\OriginalTerm{情感}{\GreekText{πάθει}}上停留久一点。许多人费尽心力，以便能站在公众面前长篇大论：如果他们得到群众的掌声，他们便以为得到了（天国）本身；但如果他们讲完后一片沉默，那么他们心中的沮丧比地狱本身还可怕！这使教会天翻地覆，因为\Emph{你们}不愿听一篇能引导你们痛悔的讲道，而愿听一篇从字句结构和声调上能使你们欢喜的讲道，仿佛在听歌者和琴师（\OriginalTerm{歌者和琴师}{\GreekText{κιθαρῳδὥν} \GreekText{καὶ} \GreekText{κιθαριστὣν}}，\Emph{参见上文}第68页）一样：而\Emph{我们}也扮演了一个颠倒、可怜的角色，被你们的欲望牵引，而我们本应根除它们。这就像是一个父亲有一个虚弱、好哭的孩子（已经比应当更娇嫩了），尽管孩子如此软弱，他仍给他蛋糕和冷饮，以及任何能让孩子高兴的东西，却不考虑什么对他有益；然后当医生责备他时，他为自己辩护说：\Quote{我能怎么办？我不忍心看孩子哭。}你这可怜、悲惨的人，你这背叛者！因为我不能称这样的人为父亲：与其让这短暂的快乐成为持久忧伤的根源，还不如让你使孩子痛苦一阵子，使他永久恢复健康，这要好多少？我们的情况也是如此，当我们徒劳地忙于优美的措辞和句子的结构与和谐，以便使人喜悦，而非使人获益时；我们以被人钦佩为目标，而非教导；以使人快乐，而非刺伤心灵；以被人鼓掌并带着赞美离开，而非纠正人的品行！请相信我，我说的是真心话——当我在讲道中听到有人为我鼓掌时，那一刻我确实感到人性的喜悦（因为为什么我不承认真相？）：我很高兴，并沉溺于这种愉悦的感觉；但当我回家后，想到那些鼓掌的人并没有从我的讲道中得到益处，反而他们所应得的益处，因鼓掌称赞而失去了，我便痛苦、呻吟、哭泣，感到我似乎白讲了。我对自己说：\Quote{我的劳苦给我带来什么益处呢？因为听者不愿从他们从我们这里听到的话中获益。}不但如此，我常想制定一条规则，禁止一切鼓掌，并劝你们以安静和合宜的秩序聆听。但请容忍我，我恳求你们，并听从我，如果你们觉得合适，让我们现在就确立这条规则：不许任何听者在任何人讲道中间鼓掌，如果他想赞美，就让他默默赞美：没有人阻止他；让他全部的心思和热切愿望都放在领受所讲的话上。—— 那声音又是什么意思？我正在为这件事制定规则，而你们甚至不能忍受听我说！—— 这条规定会有许多好处：它将是一种哲学的操练。甚至连外教哲学家——我们听说他们讲论，却从未发现他们的话伴随嘈杂的掌声：我们听说宗徒们公开讲道，但从未有记载说在讲道中听者以高声的赞许打断讲道者。这对我们将是一大益处。但让我们确立这条规则：让我们都安静地听，并完整地说出（我们要讲的话）。因为如果确实我们离开时保留了所听见的，我所坚持的是，即使这样，赞美也是无益的——但为了避免过多地详细论述（这一点）；请没有人指责我粗鲁——但既然它毫无益处，甚至是有害的，就让我们解除这一障碍，停止这些跳跃，削减灵魂的这些欢跃。基督曾在山上公开讲道，但没有人说话，直到他讲完。我并不剥夺那些想得到掌声的人；相反，我使他们更受钦佩。听者安静地听，之后凭记忆随时随地赞美，无论是在家还是在外，这远比失去一切、空手回家、不拥有他所赞美的东西要好得多。否则听者岂不变得可笑？不仅如此，他还会被视为谄媚者，他的赞美无异于讽刺，因为他声称教师讲得好，却说不出来他讲了什么。这完全像阿谀奉承。因为当一个人听了歌者和演员时，如果他不能以同样的方式唱出那旋律，那并不奇怪；但既然这不是歌曲或声音的表演，而是思想和智慧反思（\OriginalTerm{哲学}{\GreekText{φιλοσοφίας}}）的目的和含义，并且每个人都可以轻易地讲述和报告所说的话，那么，如果有人不能说出他赞美讲道者的缘由，他怎能不受到指责呢？没有什么比安静和秩序更相称于教会了。喧嚣属于剧院、浴室、公共游行和市场；但当教导的道理，而且是这样伟大的道理时，应当有宁静、安静、沉静的反思，以及非常平安的港湾（\OriginalTerm{哲学和宁静的港湾}{\GreekText{φιλοσοφία} \GreekText{καὶ} \GreekText{πολὺς} \GreekText{ὁ} \GreekText{λιμήν}}）。我请求并恳求这些事：因为我正在寻找能使我帮助你们灵魂受益的方法。我认为这不是小方法：它不仅会造福你们，也会造福我们。这样，我们就不会\OriginalTerm{被骄傲冲昏头脑}{\GreekText{ἐκτραχηλίζεσθαί}}，不会受诱惑去爱赞美和荣誉，不会去讲那些使人喜悦的事，而是讲那些有益的事：这样，我们就会把全部时间和勤勉放在思想的内容和意义上，而不是放在文章技巧和辞藻优美上。到画家的画室里去，你会看见那里多么安静。这里也应当如此：因为我们在这里也是用美德的色彩绘制肖像——君王的肖像（每一幅都是），没有一幅是普通人的——怎么？又鼓掌了？这改革不容易，只因长期的习惯。——画笔是舌头，艺术家是圣神。请问，在举行奥迹时，有什么噪音吗？有什么扰乱吗？当我们施洗时（\OriginalTerm{受洗}{\GreekText{βαπτιζώμεθα}}），当我们做其他一切事时？不是万物都（似乎）以安静和沉默装饰吗？整个天面上散布着这种（宁静的）魅力。——正因为如此，我们甚至在外邦人中被诽谤，好像我们做一切事都是为了炫耀和夸张。但如果这被制止了，对首席的喜爱也会熄灭。如果有人爱慕赞美，他可以在听讲之后、一切结束时得到赞美，这足够了。是的，我恳求你们，让我们确立这条规则，使我们能按天主的旨意行事，并因祂的唯一圣子、我们的主耶稣基督的恩宠和慈悲，堪当获得祂的怜悯——愿光荣、权能、尊崇归于圣父及圣神，同祂一起，自今至永远，世世无尽。阿们。\par

\SourceRecord{Translated by J. Walker, J. Sheppard and H. Browne, and revised by George B. Stevens.；Nicene and Post-Nicene Fathers, First Series；Vol. 11.；Edited by Philip Schaff.；Buffalo, NY: Christian Literature Publishing Co.,；1889.；<http://www.newadvent.org/fathers/210130.htm>.}

% structure: p0001 source=h1 target=section reason=page-title

\section{论宗徒大事录第三十一篇讲道}

\BibleRef{宗 14:14-15}\par

\Quote{\Person{巴尔纳伯}和\Person{保禄}宗徒听说这事，就撕裂衣裳，冲进群众中间，喊着说：\InnerQuote{诸位，你们为什么做这些事？我们也是和你们有同样性情的人，我们向你们传福音，叫你们离开这些虚妄，归向永生的天主，他创造了天、地、海和其中的万物。}}\par

请注意宗徒们所做这一切的激烈程度：\Quote{撕裂衣服，跑进去，大声喊叫}，全都是出于灵魂的强烈情感，对所发生的事感到厌恶。因为这是一个悲伤，确实是一个无法安慰的悲伤，他们竟被认为神明，并引入了偶像崇拜，而这正是他们来要摧毁的！这也是魔鬼的伎俩——但他没有得逞。但他们说什么呢？\Quote{我们也是和你们一样性情的人。}一开始他们就推翻了邪恶。他们不是说简单的\Quote{人啊}，而是\Quote{像你们一样}。然后，为了不显得他们尊敬那些神明，请听他们补充的话：\Quote{向你们宣讲，叫你们离开这些虚无，归向永生的天主，祂创造了天、海和其中的万物。}请注意他们如何没有一处提到看不见的事物。（\Emph{b}）因为他们已经学会，一个人不应如此努力说一些与天主相称的话，而应说对听者有益的话。（\Emph{a}）那么怎样？如果祂是万物的创造者，为何祂不以祂的眷顾来照顾这些事呢？——\BibleQuote{祂在过去的时代，容许所有外邦人各行其道} \BibleRef{宗 14:16}但为何祂容许他们，他没有说，因为目前他专注于眼前重要的事，没有一处提及基督的名字。请注意，他不想加重对他们的指控，而是希望他们自己将一切归于天主。\BibleQuote{但他并没有让自己不留下见证，因祂行了善事，从天上给你们降雨，赐予收获的时节，使你们的心充满食物和喜乐。} \BibleRef{宗 14:17}（\Emph{c}）请看他如何隐藏地将指控放入\Quote{因祂行了善事}等等。然而，如果天主这样做了，祂就不可能\Quote{任由他们}；相反，他们应该受罚，因为他们享有如此大的恩惠，却没有承认祂，甚至不认祂为养育者。\Quote{从天上}，他说，\Quote{赐给你们雨水}。同样，达味也说\BibleQuote{从他们的五谷、酒和油的果实，他们得以丰盈} \BibleRef{咏 4:7}，在许多地方提到创造时，他提出这些恩惠；耶肋米亚先提到创造，然后提到通过雨水显示的眷顾，所以宗徒在这里的论述如同从那些圣经中学到的。\Quote{充满}，他说，\Quote{食物和喜乐。} \BibleRef{耶 5:24}食物是以大方的慷慨\OriginalTerm{大方的慷慨}{\GreekText{φιλοτιμίας}}赐下的，不仅仅是为了节俭的饱足，也不因需要而吝啬。\Quote{说了这些话后，他们好不容易才阻止了群众} \BibleRef{宗 14:18}——事实上正是这件事使他们获得了最大的钦佩——\Quote{不向他们献祭。}你是否注意到，这正是他们的目的，要终止那种疯狂？\Quote{但来了}，它说，\Quote{一些从安提约基雅和依科尼雍来的犹太人} \BibleRef{宗 14:19}。——真是魔鬼的儿女，不只在自己的城里，也在城外，他们做这些事，并且如同宗徒们热心建立宣讲一样，他们也同样致力于终结宣讲！——\Quote{他们说服了群众，用石头砸了保禄，把他拖到城外。}（\Emph{e}）那么，外邦人把他们看作神，但这些（犹太人）却\Quote{拖}他，\Quote{出城，以为他死了。他们说服了群众}——因为不可能所有人都如此尊敬他们。正是在他们受到这种敬重的城市里，他们又如此可怕地受到虐待。这也使旁观者受益。\Quote{免得有人}，他说，\Quote{把我看作超过他所看见我的，或从他那里听到的任何事。}（第20节）——\Quote{然而，当门徒们围着他时，他站起来，进了城。}（\Emph{d}）这里应验了那句话：\BibleQuote{我的恩宠够你用，因为我的能力在软弱中得以成全。} \BibleRef{格后 12:9}这比使瘸子站起来更伟大！（\Emph{f}）\Quote{进了城。}你看到这热忱吗？你看到他多么炽热，多么燃烧！他再次进入那城本身：这证明，如果他偶尔退去，那是因为他撒下了圣言，并且不应该激起他们的愤怒。（\Emph{h}）然后他们走遍了所有他们曾陷入危险的城市。\Quote{第二天}，它说，\Quote{他与巴尔纳伯出发去了德尔贝。当他们向那城宣讲了福音，并教导了许多人后，他们又回到吕斯特辣、依科尼雍和安提约基雅，坚固门徒们的灵魂，劝勉他们持守信德，并且我们必须经过许多苦难才能进入天主的国。}（第21-22节）这就是他们所说的，也是他们所显示的。但这是特意这样做的，不仅宗徒们，连门徒们也是如此，使他们从一开始就学会宣讲的力量，以及他们自己也必须忍受同样的事，使他们能高尚地站立，不只呆望着奇迹，而更准备好面对考验。因此，宗徒自己也说：\BibleQuote{你们在我身上所见所闻的同样斗争。} \BibleRef{斐 1:30}迫害接踵而至：战争、争斗、用石头砸。（\Emph{g}）这些事，不亚于奇迹，既使他们更显赫，也为他们预备了更大的喜乐。圣经没有一处说他们因行了奇迹而欢喜地回去，而是（说他们欢喜），因为\BibleQuote{他们算配为这名受辱。} \BibleRef{宗 5:41}这是基督教导他们的，说：\BibleQuote{不要因魔鬼服从你们而欢喜。} \BibleRef{路 10:20}因为真正的、纯粹的喜乐就是为基督的缘故受苦。（\Emph{i}）\Quote{并且必须经过许多苦难：}这是什么鼓励\OriginalTerm{鼓励}{\GreekText{προτροπή}}？他们怎么说服他们的，一开始就告诉他们苦难？然后还有另一个安慰。\BibleQuote{他们在各教会为他们选立了长老，并祈祷禁食后，就把他们交托给他们所信的主。} \BibleRef{宗 14:23}你看到保禄的热忱吗？——然后另一个安慰：\Quote{把他们交托}，它说，\Quote{给主。他们走遍了丕息狄雅，来到旁非里雅。他们在培尔革宣讲了圣言后，就下到阿塔里雅（第24-25节）：（\Emph{l}）从那里航行到安提约基雅，他们就是从那里被托付于天主的恩宠，去完成他们所成就的工作。} \BibleRef{宗 14:26}为什么他们回到安提约基雅？为了报告那里发生的事。此外，还有天主眷顾的伟大目的：因为需要他们从此以后放胆向外邦人宣讲。所以他们来报告这些事，好让他们能知道；并且上智安排，正好那时有人来了，禁止与外邦人来往，为使他们从耶路撒冷得到极大的鼓励，然后放胆离开。此外，这也表明他们的性情中没有任意妄为：因为他们来，同时显示了他们的勇敢，即没有（耶路撒冷）那些人的权威，他们已向外邦人宣讲；也显示了他们的顺服，即他们请求他们来裁决问题：因为他们没有被变得傲慢\OriginalTerm{被变得傲慢}{\GreekText{ἀπενοήθησαν}}，尽管取得了如此大的成功。\Quote{从那里}，它说，\Quote{他们被托付于天主的恩宠，去完成他们所成就的工作。}并且圣神还说：\BibleQuote{你们为我选出巴尔纳伯和扫禄，去做我召叫他们去的工作。} \BibleRef{宗 13:2}\Quote{他们到了，便聚集教会，述说天主藉他们所行的一切，以及祂如何为外邦人开了信德之门。他们在那里与门徒们住了不少日子。}（第27-28节）因为那城很大，需要教师。——但让我们再回顾一下所说的。\par

（概要。）\newline{}\newline{}\Quote{宗徒们一听说}等等。\newline{}\newline{}\BibleRef{宗 14:14}。首先，他们以目睹制止了他们，撕裂自己的衣服。农的儿子若苏厄在百姓战败时也曾这样做。那么，不要以为这行动与他们不相称：因为当时情绪如此激昂，否则他们无法克制，无法扑灭那烈焰（\OriginalTerm{火焰}{\GreekText{πύραν}}）。因此，当需要做一件合宜的事时，我们不要回避。因为即使做了这一切，他们也几乎无法说服他们；如果他们没有这样做，后果将会如何？如果他们不这样做，就会被认为是在炫耀\OriginalTerm{谦卑}{\GreekText{ταπεινοθρονεἵν}}，并且更加渴望得到荣誉。请注意他们的言语，责备时多么温和，既充满敬佩又含有责备。最妨碍他们的，正是这句话：\Quote{宣讲你们远离这些虚妄的事，归向永生的天主。} \BibleRef{宗 14:15}。他们说：我们确实是人，但比这些人更伟大：因为这些人都是死的。请注意，他们不仅驳斥（虚假的），而且教导（真实的），对看不见的事却只字不提——他们说：\Quote{他造了}天、地、海和其中的万物。\Quote{他在过去的世代}等等。（16–17节。）他甚至连（运行的）年岁也称为见证。\Quote{那时，有些犹太人从安提约基雅和依科尼雍来，}等等。\BibleRef{宗 14:19}。啊，犹太人的疯狂！在一个如此尊敬宗徒的民族中，他们竟敢前来，用石头砸保禄。\Quote{他们把他拖出城外，}害怕那些人——\Quote{以为他死了。}（k）\Quote{但门徒们一围上来，}等等。\Quote{他进了城。} \BibleRef{宗 14:20}。为了使门徒的精神不致消沉，因为他们曾被视为神的人却遭受了这样的待遇，他们便进去与他们交谈。\Quote{第二天，}等等。请注意，他先前往德尔贝，然后回到吕斯特辣、依科尼雍和安提约基雅，\BibleRef{宗 14:21}，当他们情绪激动时先退让，等平息后再重新攻击他们。你注意到吗，他们所行的一切并非全靠（超性）恩宠，而是靠他们自己的勤勉？经文说：\Quote{坚固}门徒的心：\OriginalTerm{坚固}{\GreekText{ἐπιστηρίζοντες}}，\Quote{进一步坚定；}所以他们已坚定，但他们进一步加增。\Quote{又说：我们必须}等等。\BibleRef{宗 14:22}：他们预言了这一点，好让他们不致跌倒。\Quote{他们在各教会为他们选立了长老，}等等。又是伴随着禁食的祝圣：又是禁食，那净化我们灵魂的禁食。（m）\Quote{他们祈祷后，}经文说，\Quote{禁食，就把他们交托于主。} \BibleRef{宗 14:23}：他们也教导他们在试炼中禁食。为什么不在塞浦路斯和撒玛黎雅选立长老？因为后者靠近耶路撒冷，前者靠近安提约基雅，而且那地的话语强有力；而在那些地区，他们需要更多的安慰，尤其是外邦人，需要很多教导。（n）\Quote{他们一来到，}等等。\BibleRef{宗 14:27}。他们来了，教导他们应当听从圣神的派遣。他们没有说自己做了什么，而是\Quote{天主借他们所做的一切。}（o）在我看来，他们指的是他们的试炼。他们来到这里并非徒然，也不是为了休息，而是受圣神引导的安排，为使向外邦人的宣讲坚定确立。（p）请注意保禄的热忱。他不问向外邦人宣讲是否合适，而是立即宣讲：因此他说：\Quote{我并没有与血肉之人商议。} \BibleRef{迦 1:16}\par

这确实是一件伟大之事，一个伟大的、慷慨的灵魂（像他那样）！自那时起有多少人信了，但没有一个人像他那样闪耀！我们需要的是热忱、极大的热情、一个准备好面对死亡的灵魂。否则，若不被钉在十字架上，就不可能到达天国。我们不要自欺。因为在战争中，若过着安逸、经商、叫卖和闲散的生活，就不可能安全脱身，更何况在这场战争中。或者你们不认为这是一场比所有其他战争更糟糕的战争吗？（\Emph{Infra}, p. 204, note 1.）\Quote{因为我们战斗不是}他说，\Quote{对抗血和肉。} \BibleRef{弗 6:12} 因为即使在吃饭、走路和沐浴时，仇敌也与我们一起，他知道没有休战的时间，除了睡眠而已；不但如此，甚至在那时他也进行战争，将不洁的思想注入我们心中，藉着梦境使我们放荡。我们不警醒，不振奋自己，不注视我们面对的众多敌军，不反思这本身就是最大的不幸——尽管被如此多的战争包围，我们却如同和平时期一样过着安逸的生活。相信我，现在可能必须承受比保禄所受的更糟的苦难。那些仇敌用石头打伤了他：有一种用言语造成的伤害，甚至比石头更糟。那么我们必须做什么？和他所做的一样：他不恨那些向他扔石头的人，但在他们把他拖出城外之后，他又进入他们的城市，成为那些对他做了如此不公之事的人的恩人。如果你也忍受那严厉侮辱你、对你不公的人，那么你也已经被石头打了。不要说：\Quote{我没有对他做过任何伤害。} 保禄做了什么伤害，以致他应该被石头打？他宣讲一个天国，他使人们脱离错误，带领他们归向天主：这些恩惠，配得冠冕，配得传报者的呼声，配得一千件好事——而不是石头。然而（远非怨恨），他做了完全相反的事。因为这正是光荣的胜利。\Quote{他们拖着他，} \BibleRef{宗 14:19} 经上说。这些（恶人）也常常拖拽（别人）：但你不要生气；相反，你要温和地宣讲圣言。有人侮辱了你吗？保持沉默，如果你能祝福，那么你也宣讲了圣言，给了温和的教训，给了温良的教训。我知道许多人受的伤不如言语造成的打击那样痛：确实，一个伤身体，另一个伤灵魂。但让我们不要感到疼痛，或者不如说，感到疼痛时让我们忍受。你没看见拳击手吗？他们的头被打得青肿，用牙齿咬住嘴唇，如此温和地忍受疼痛？无需咬牙切齿，无需咬嘴唇。记住你的师傅，藉着记念，你立刻就应用了补救。记住保禄：思考，受击打的人，你征服了，而那击打者，被击败了；藉此你治愈了全部。只需一时的权衡，你就成就了一切：不要匆忙，甚至不要移动，你就熄灭了整个（火）。为基督受苦有一种极大的说服口才：你宣讲的不是信德之言，而是\OriginalTerm{忍耐的智慧}{\GreekText{φιλοσοφίας}}。但你会说，他越看见我的温和，就越攻击我。那么你为此痛苦，是因为他更增加了你的赏报吗？\Quote{但这是途径，}你说，\Quote{使他变得无法忍受。} 这仅仅是你小心眼的借口：相反，另一种才是使他变得无法忍受的途径，即你报复自己。如果天主知道，因着不报复的容忍，不义之人会变得无法忍受，祂就不会这样做；相反，祂会说：报复你自己；但祂知道，除此之外的途径更可能做好事。你不要制定与天主相反的法律：照祂吩咐你的去做。你并不比造我们的祂更仁慈。祂曾说：\Quote{忍受被亏待：}你说：\Quote{我以冤报冤，使他不会变得无法忍受。} 那么你对他比天主对他还有更多的关心吗？这样的话只是冲动和坏脾气，是傲慢，是制定与天主法律相反的法律。因为即使那人（因我们的容忍）受到伤害，我们岂不也有责任服从？当天主命令任何事时，我们不要制定相反的法律。\Quote{柔和的回答，}我们读到，平息愤怒 \BibleRef{箴 16:1} : 不是反对的回答。如果它使你受益，也使他受益；但如果它伤害了你——你是要纠正他的人——那么它岂不更严重地伤害他？\Quote{医生，医治你自己。} 有人说了你坏话吗？你称赞他。有人辱骂你吗？你赞美他。有人图谋害你吗？你向他行善。以相反的事回报他，如果你至少关心他的救恩，而不想报复你自己的痛苦。然而，你会说，虽然他曾多次从我这里遇到耐心，但他变得更坏了。这不是你的事，而是他的事。你愿意知道天主受了什么冤屈吗？他们推倒祂的祭坛，杀了祂的先知 \BibleRef{列上 19:10}，然而祂忍受了一切。祂不能从上头发出雷电吗？不，当祂派了祂的先知，他们杀了他们，然后祂派了祂的圣子 \BibleRef{玛 21:37}，当他们行了更大的不敬时，祂反而给了他们更大的恩惠。你也是如此，如果你见一个人愤怒，就更加让步：因为这疯狂更需要\OriginalTerm{安慰}{\GreekText{παραμθίας}}。他越严厉地辱骂你，他就越需要你的温良：正如强风吹时，需要顺势屈服，发怒的人也是如此。当野兽最凶猛时，我们都逃跑；所以我们也应当从发怒的人那里逃跑。不要认为这是给他的荣誉：因为当我们从野兽和疯子那里转向一边时，岂不是给了他们荣誉吗？绝不是，这是羞辱和轻蔑：或者不如说不是羞辱和轻蔑，而是怜悯和仁慈。你没看见水手吗？当狂风猛烈吹时，他们降下船帆，使船不会沉没？当马匹脱离驾驶者时，他只是把它们引到开阔的平原，而不拉紧它们，以免自愿耗尽力气？你也要这样做。愤怒是火，是快速的火焰，需要燃料：不要给火提供食物，你就很快熄灭了邪恶。愤怒本身没有力量；必须有别人来喂养它。对你来说没有借口。他被疯狂占据，不知道自己在做什么；但当你看见他是什么样子，却陷入同样的罪恶，没有被这景象唤醒理智，你还有什么借口？如果来赴宴，你看见在宴席刚开始就有人醉酒、行为不雅，那么那个看见他之后自己喝醉的人，岂不更不可原谅？这里正是这样。我们是否以为说\Quote{我不是先开始的}是借口？这恰恰是对我们不利的，因为即使看见别人处在那种状态，也没有使我们清醒。这就像有人说：\Quote{我没有先杀他。} 正是这件事使你该受惩罚，因为即使有这样的景象警告，你也没有约束自己。如果你看见醉汉正在呕吐、干呕、迸发，他的眼睛紧张，用他的污秽充满桌子，每个人都赶紧躲开，然后你自己陷入同样的状态，你岂不更可恨？发怒的人就像他：比呕吐的人更甚，他的血管扩张，眼睛通红，内脏被折磨；他说出的话比那食物更污秽；他所吐出的都是生涩的，没有经过适当消化，因为愤怒不允许如此。但正如在那情况下，过量的\OriginalTerm{液汁}{\GreekText{χυμων}}在胃里喧闹，常常将一切内容物吐出；同样在这里，过量的热气在灵魂中引起骚动，使他不能隐藏那些本应不说的话，而是不论适合说不适合说，他都一概说出，不是使听者蒙羞，而是使自己蒙羞。因此，我们躲开呕吐的人，同样也躲开发怒的人。让我们把尘土撒在他们的呕吐物上：怎么做？保持沉默：让狗来吃掉呕吐物。我知道你听这话感到厌恶：但愿你看到这些事情发生时也有同样的厌恶，而不是喜欢这事。谩骂的人比转回自己呕吐物的狗更污秽。因为如果他吐了一次便不再吐，他就不会像那只狗；但如果他再次吐同样的东西，显然是因为他又吃了同样的东西。那么还有什么比这样的人更可憎？还有什么比那张咀嚼那样食物的嘴更污秽？然而这是自然的工作，但另一件不是，或者不如说两者都是违反自然的。怎么？因为无缘无故地谩骂不符合自然，而是违反自然：他说话不像人，而一部分像野兽，一部分像疯子。正如身体的疾病违反自然，这也一样。为表明它违反自然，如果他继续如此，他会逐渐灭亡；但如果他继续在自然中，他不会灭亡。我宁愿与一个吃泥土的人同桌，也不愿与一个说这样的话的人同桌。你们没看见猪吃粪吗？这些人也是如此。因为还有什么比谩骂之人所说的话更恶臭？他们专门说无益、不洁的话，而是无论什么下流、什么不雅，他们都着意去做去说；更糟的是，他们以为羞辱了别人，而实际上他们是在羞辱自己。因为他们羞辱自己是很明显的。因为，暂且不谈那些说谎的人（在他们的谩骂中），假设某个臭名昭著的娼妓，或者甚至舞台上的某个其他（堕落之人），让那个人与另一个人打架：然后让后者向前者提起她或他是什么人，前者则以同样的指责回敬后者：哪个人被言语伤害最大？前者只是被叫成实际上她或他是什么，这与另一个不同：所以第一个在羞耻方面没有增加什么（比以前多），而另一个则大大增加了耻辱。但再次，假设有一些隐藏的行为（修订文本 \OriginalTerm{已做之事}{\GreekText{εἰργασμένα}} \Quote{已经做了的事}），而且只有谩骂的人知道它们：然后，他保持沉默直到现在，让他公开\OriginalTerm{炫耀}{\GreekText{ἐκπομπευέτω}}那指责：即便如此，他自己比另一个更蒙羞。怎么？通过使自己成为邪恶的传报者，因此要么为自己招致不涉及任何此类事情的指控，要么获得不值得信任的性格。你会看见所有人立刻指责他：\Quote{如果他真的知道一桩谋杀发生，他应该全部揭露：}因此他们厌恶他，甚至不把他当人，他们恨他，他们说他是野兽，凶猛残忍；而另一个他们宽恕得多。因为我们并不那么恨那些有伤的人，而是恨那些强迫人揭开和展示伤口的人。因此，那个人不仅羞辱了别人，也羞辱了自己、听者和共同的人性：他伤害了听者，没有好处。因此保禄说：\Quote{若有什么有益于建造的话，好给听者带来恩宠。} \BibleRef{弗 4:29} 让我们获得只说好事的舌头，使我们可爱可亲。但事实上，一切都达到了这样的邪恶程度，以至于许多人夸耀那些他们本该掩面的事情。因为许多人的威胁是这样的：\Quote{你受不了我的舌头，}他们说。这些话，只配一个女人、一个堕落的老醉妇、一个在广场上被拖行的人、一个老鸨。没有什么比这些话更可耻，更不男子汉，更女性化，就是把力量放在舌头上，因为能谩骂而自以为了不起，就像游行中的家伙，像小丑、寄生虫和谄媚者。他们是猪而不是人，以此自豪。而你本应（更早）埋葬自己，如果别人给你这样的评价，你应该厌恶地回避这指控，认为它可憎和不男子汉，相反你使自己成为（你自己的）\OriginalTerm{侮辱}{\GreekText{ὕβρεων}}的传报者。但你不能伤害你谩骂的人。所以我恳求你们，考虑到邪恶已到了如此高度，许多人以此夸耀，让我们清醒，让我们挽救那些疯狂的人，让我们把这些商议从城里除去，让我们的舌头优雅，让我们除去一切恶言，使我们从罪恶中洁净，能为我们得来自天上的善意，并蒙天主的仁慈，藉着祂独生子的恩宠和慈悲，与祂一起，同圣父及圣神，永受光荣、权能、尊崇，自今至永远，世世无尽。阿们。\par

\SourceRecord{Translated by J. Walker, J. Sheppard and H. Browne, and revised by George B. Stevens.；Nicene and Post-Nicene Fathers, First Series；Vol. 11.；Edited by Philip Schaff.；Buffalo, NY: Christian Literature Publishing Co.,；1889.；<http://www.newadvent.org/fathers/210131.htm>.}

% structure: p0001 source=h1 target=section reason=page-title

\section{宗徒大事录第三十二篇讲道}

\BibleRef{宗 15:1}\par

\BibleQuote{有几个人从犹太下来，教训弟兄们说：\InnerQuote{你们若不按梅瑟的规条受割礼，就不能得救。}}\par

请注意，在关于外邦人的正确进展的每一步中，开端是如何作为必要之事被引入的。在此之前，（伯多禄）受人责难时为自己辩护，并以道歉的语气说了他所说的一切，这使他的话被人接受；然后，犹太人转身离开，于是（保禄）来到了外邦人那里。现在，看到另一种极端出现，（宗徒）便颁布了法律。因为很可能由于他们是受天主教导的，便不加区别地向所有人讲论，这激起了（已相信的）犹太人的嫉妒。他们不仅谈论割礼，而且说：你们甚至不能得救。然而情况恰恰相反：接受割礼，他们反而不能得救。你们看到，考验如何紧密地接踵而来，从内部，从外部？同样安排得当的是，这事发生在保禄在场时，以便他能回答他们。\BibleQuote{那时，保禄和巴尔纳伯与他们发生了不小的纷争和辩论，他们就决定，保禄和巴尔纳伯并其中的几个人，为这问题上耶路撒冷去见宗徒和长老。}\BibleRef{宗 15:2}保禄并没有说：什么？难道我不配在这么多奇迹之后被人相信吗？但他为了他们而顺从。\BibleQuote{教会送他们上路，他们经过腓尼基和撒玛黎雅，叙述外邦人的归化，使众弟兄大为欢喜。}\BibleRef{宗 15:3}并且要注意，结果是所有撒玛黎雅人也都知道了外邦人的遭遇，他们欢喜。\BibleQuote{他们到了耶路撒冷，教会、宗徒和长老接待了他们，他们述说了天主与他们所做的一切事。}\BibleRef{宗 15:4}看，这是何等上智的安排！\BibleQuote{但有几个信了的法利塞党人起来说：必须给他们行割礼，并命令他们遵守梅瑟的法律。宗徒和长老聚集商讨这事。经过许多辩论，伯多禄起来对他们说：诸位仁人弟兄，你们知道，早在最初，天主就在你们中间拣选了我，叫外邦人从我口中得听福音之道而相信。}\BibleRef{宗 15:5}请注意伯多禄从一开始就置身事外（\OriginalTerm{分离}{\GreekText{κεχωρισμένον}}），甚至到这时还在犹化。但他却说：\BibleQuote{你们知道。}\EditorialNote{参阅 宗 10:45; 11:2}也许那些从前因科尔乃略的事责难他，并与他一同进去的人也在场：因此他提出他们作证人。他说：\BibleQuote{从起初，}\BibleQuote{就在你们中间拣选了。}\Quote{在你们中间}是什么意思？要么是在巴勒斯坦，要么是你们在场。\BibleQuote{藉我的口。}请注意他如何显示是天主藉他说话，而非人言。\BibleQuote{那洞察人心的天主，给他们作了证：}他把他们引向属灵的见证：\BibleQuote{赐给他们圣神，如同赐给我们一样。}\BibleRef{宗 15:8}他处处将外邦人置于完全平等的地位。\BibleQuote{并借着信德洁净了他们的心，不分我们和他们之间的区别。}\BibleRef{宗 15:9}他说，他们单从信德就获得了同样的恩赐。这也是对那些反对者的一种教训；这甚至能教导他们：只需要信德，而不需要行为或割礼。因为他们说这一切，不仅是作为对外邦人的辩护，也是要教训（犹太信徒）放弃法律。然而，目前还没有说出这一点。\BibleQuote{那么，现在你们为什么试探天主，把重轭放在门徒的颈上呢？}\BibleRef{宗 15:10}\Quote{试探天主}是什么意思？好像他没有能力藉信德拯救人一样。因此，引入法律是由于缺乏信德。然后他表明他们自己也没有从中获益，并把全部（重点）转向攻击法律，而不是攻击他们，以此截断对他们的指控：\BibleQuote{这轭无论是我们的祖先或是我们都不能负。但我们相信，我们得救是因主耶稣的恩宠，和他们一样。}\BibleRef{宗 15:11}这些话多么有力！保禄在《罗马书》中详述的，伯多禄也在这里说了同样的话。\BibleQuote{因为如果亚巴郎是因行为成义，他就有可夸耀的，但在天主前却没有。}\BibleRef{罗 4:2}你们是否看出这一切更多是为教训他们，而非为外邦人辩护？然而，如果他没有说话的理由而说出这话，就会被人怀疑；机会来了，他就抓住它，无所畏惧地说了出来。请看，在所有场合，敌人的计谋如何为他们所利用。如果那些人没有挑起问题，这些事就不会被说出来，以下的事也不会发生。\par

（复述。）\newline{}（\Emph{b}）但让我们更仔细地看看所说的。\BibleQuote{有几个人}等等。在耶路撒冷，当时外邦人中并没有信徒：但在安提约基雅当然有。因此，有些人仍受这种喜爱权柄的疾病之苦，希望将外邦人拉拢到他们一边。而保禄，尽管他也精通法律，却没有受此影响。\BibleQuote{因此，保禄和巴尔纳伯与他们发生了不小的争论}等等。见：\BibleRef{宗 15:2} 但他从那里回来后，教义也变得更加精确。因为如果耶路撒冷的人没有颁布这样的规定，那么这些人更（无权这样做）。\BibleQuote{他们被人送行，}等等，\BibleQuote{他们使弟兄们大感喜乐。}见：\BibleRef{宗 15:3} 你注意到吗？凡是不贪爱权柄的人，都因他们的信德而喜乐。他们的叙述并非出于野心，也不是为了炫耀，而是为了证明向外邦人的宣讲。见：\BibleRef{宗 15:4} 因此，他们没有提及在犹太人中所发生的事。\BibleQuote{但有几个信了的法利塞派起来，}等等。见：\BibleRef{宗 15:5} （\Emph{a}）但即使他们想将外邦人拉到自己一边，他们得知宗徒们也不能忽视它。\BibleQuote{宗徒和长老们，}等等。见：\BibleRef{宗 15:6} \BibleQuote{在我们中，}他说，\BibleQuote{天主拣选了；}和\BibleQuote{从早先的日子：}他说是早就如此，不是现在。这一点也不小——在犹太人相信、而不是（从福音）转离的时候。\BibleQuote{在我们中；}一个从地点而来的论证：\BibleQuote{从早先的日子，}从时间而来。而那个说法，\BibleQuote{拣选：}正如在他们自己的情形中，他说的不是（这样）愿意，而是\BibleQuote{拣选了外邦人藉我的口听到福音的话而相信。}这从何证明？从圣神。然后他表明，给予他们的见证不仅是恩宠，也是他们的德行。\BibleQuote{那洞察人心的天主为他们作了证}见：\BibleRef{宗 15:8}；赐予他们不少于（我们），因为他说，\BibleQuote{在我们和他们之间不作区别}见：\BibleRef{宗 15:9} 那么，人心是处处必须看重的。而且这话说得很恰当：\BibleQuote{洞察人心的天主为他们作了证：}正如在先前的情况中，\BibleQuote{主，你洞察所有人的心。}见：\BibleRef{宗 1:24} 为了说明这是含义，请看他加上的话：\BibleQuote{在我们和他们之间不作区别。}当他提到为他们所作的证言时，他说出了那句伟大的话，就是保禄所说的：\BibleQuote{割损算不得什么，不割损也算不得什么。}见：\BibleRef{格前 7:19} \BibleQuote{为使两者在祂内成为一体。}见：\BibleRef{弗 2:5} 所有这些的种子都存在于伯多禄的讲论中。他不是说（在）受割损的人和他们之间，而是\BibleQuote{在我们之间}，即宗徒们，\BibleQuote{和在他们之间。}然后，为了使\BibleQuote{不作区别}这一表达不显得过分，在信德之后，他说——\BibleQuote{藉着信德洁净了他们的心}见：\BibleRef{宗 15:10}——他首先彻底洁净了他们。然后他表明，法律并非邪恶，而是他们自己软弱。——\BibleQuote{但我们相信，藉着主耶稣的恩宠，我们将得救，甚至如同他们一样。}见：\BibleRef{宗 15:11} 注意他是如何以一个可怕的考虑来结束的。他没有向他们讲述先知们，而是从当前的事情，他们自己就是见证。当然（先知们）自己后来也添加了他们的证言（\Emph{infra} 第15节），并通过已经发生的事使论证更强。并且请注意，他首先允许在教会中提出这个问题，然后才发言。\BibleQuote{不作区别}——他不是说受割损的人，而是\BibleQuote{我们和他们，}即外邦人：因为这（逐步进展）一点点变得更强烈。\BibleQuote{那么，你们为什么试探天主？}祂已成为外邦人的天主：因为这就是试探：* * * 祂是否能在法律之后仍然拯救。看他做了什么。他表明他们处于危险中。因为如果法律不能做到的，信德有能力做到，\BibleQuote{我们相信藉着主耶稣的恩宠，我们将得救，甚至如同他们一样}见：\BibleRef{迦 2:16}：但信德失落，看，他们自己（处于）毁灭中。而且他没有说：为什么你们不信？这更严厉，而是说：\BibleQuote{试探天主，}而且事实已经得到了证明。\par

（\Emph{c}）法利塞人的这种厚颜无耻，即使在信德之后，他们仍坚持法律，且不愿服从宗徒们。但看这些人，他们说话多么温和，没有权威的口吻：这样的话是亲切的，更易于铭刻在心中。注意，这绝不是言辞的展示，而是由事实、由圣神的证明。然而，尽管有如此证据，他们仍然说话温柔。并且注意，他们不是来控告安提约基雅的人，而是\BibleQuote{述说天主藉他们所行的一切事。}见：\BibleRef{宗 15:4} 但这些人又抓住机会（来实现他们自己的目的），\BibleQuote{但有几个人起来，}等等。见：\BibleRef{宗 15:1} 他们在爱权柄上如此费力：保禄和巴尔纳伯受责备，并不是在宗徒们的知情下。但他们仍然没有提出这些指控：但一旦他们证明了事实，那么（宗徒们）就用更强硬的词句写信。\par

因为温良处处是巨大的善：我说的是温良，不是愚昧的冷漠；温良，不是谄媚：因为这两者之间有天壤之别。没有什么使\Person{保禄}烦乱，没有什么使\Person{伯多禄}不安。当你拥有令人信服的证据时，为何发怒，使它们失效呢？一个发脾气的人永远无法说服人。昨天我们也谈论了愤怒；但没有理由今天我们不再谈论；或许紧接第一次劝诫之后的第二次劝诫会有些效果。因为实际上，一剂药虽然能治愈伤口，但除非不断更换，否则会破坏一切。且不要以为我们不断谈论同样的事是在谴责你们：因为如果我们谴责你们，我们就不会谈论；但如今，我们怀着希望你们获益良多的心说这些话。但愿我们真的不断谈论同样的事：但愿我们的讲论没有别的主题，除了我们如何克服自己的情欲。因为，皇帝们生活在奢华和如此巨大的尊荣中，无论是在坐席还是其他任何时候，所谈论的也无非是如何战胜敌人——因此他们每天集会，任命将军和士兵，征收税赋和贡品；而且，在所有国务中，推动性的原因就是这两个：战胜攻击他们的人，以及让臣民安居乐业——我们却对这样的主题毫无兴趣，甚至做梦也不曾去谈论它们：而如何买地、购买奴隶、增加财产，这些却是我们每天可以谈论而不厌倦的：至于我们自身内真正属于我们的事物，我们既不想谈论，也不愿接受劝告，更不愿容忍别人谈论它们？这难道不是完全不合情理吗？但请回答我，你们谈论什么？关于晚餐？那是厨子的话题。关于金钱？那是小贩和商人的话题。关于建筑？那是木匠和建筑者的话题。关于土地？那是农民的话题。但对我们而言，没有其他适当的行业，只有这一个：如何为灵魂积攒财富。那么，不要让这番谈论令你们厌烦。为何没有人指责医生总是谈论医术，也没有人指责其他手艺之人谈论他们的特殊技艺？如果确实已经达成了对情欲的掌控，以至于不需要提醒我们，那么我们就可以被合情合理地指责为野心和炫耀——或者更确切地说，即使那时也不会。因为即使获得了健康，仍然需要谈论，以免旧病复发而得不到纠正：正如事实上医生不仅对病人，也对健康者谈论，他们还有关于这个主题的书籍，一部分是关于如何摆脱疾病，另一部分是关于如何保持健康。因此，即使我们健康，我们也不可放弃，而必须做一切来保持健康。当我们生病时，有双重必要接受劝告：第一，好使我们摆脱疾病；第二，好使摆脱之后不再陷入其中。那么，我们现在是按照治疗病人的方法在谈论，而非按照健康者的保养规则。\par

那麼，如何能根除這邪惡的情慾？如何制伏（\OriginalTerm{绊倒}{\GreekText{ὑποσκελίσειε}}）這劇烈的熱症？讓我們看看它從何而生，並除去其原因。它通常從何而起？來自傲慢和極大的驕矜。那麼，我們除去這原因，疾病也隨之除去。但什麼是傲慢？它從何而生？因為我們或許需要追溯一個更高的起源。但無論理性可能指出什麼途徑，讓我們採取之，好使我們深入禍害的底部，並將其連根拔起。那麼，傲慢從何而來？來自我們不審視自己的事，反而忙於探究土地的本質（雖然我們不是農夫）、黃金的本質（雖然我們不是商人）、以及衣著等一切；而對我們自身和我們的本性，卻從不觀看。你會說：誰不認識自己的本性？許多人——或許所有人，除了少數人；你若願意，我將證明這點。因為，請告訴我：人是什麼？如果有人被問及，他能直接回答這些問題嗎：人與禽獸的區別何在？他與天上居民的相似之處何在？人能成為什麼？因為如同其他任何材料一樣，在這情況中也是如此：人是原料，但從中可以造出天使或野獸。這難道不是一個奇怪的話嗎？然而你常在聖經中聽到這話。因為關於某些人，經上說：\BibleQuote{他是主的使者}\BibleRef{拉 2:7}；又說：\BibleQuote{從他的嘴唇，他們要尋求判斷}\BibleRef{拉 3:1}；又說：\Quote{我派遣我的使者在你面前}；但關於某些人，說：\BibleQuote{蛇，毒蛇的種類}\BibleRef{瑪 12:34}。所以，一切都取決於運用。我為何說天使？人能成為神，並成為天主的子女。因為我們讀到：\BibleQuote{我曾說：你們是神，你們都是至高者的子女}\BibleRef{詠 81:6}。而更大的事是，成為神、成為天使、成為天主的子女的能力已交在他自己手中。是的，正是如此：人能成為天使的製造者。這話也許使你震驚？然而請聽基督說：\BibleQuote{在復活的時候，他們不娶也不嫁，卻像天使一樣}\BibleRef{瑪 22:30}。又說：\BibleQuote{能領受的，就讓他領受}\BibleRef{瑪 19:12}。總之，是德行使人成為天使；而這在我們的能力範圍內；所以我們能製造天使，雖非在本性上，卻確實在意志上。因為若缺少德行，即便是天使的本性也無益；魔鬼便是證明，他一度是天使；但若有德行，為人本性也無損失；若翰便是證明，他是一個人，還有厄里亞升了天，以及所有將要往那裡去的人。因為這些人雖有身體，卻不妨礙他們居住在天上；而那些雖無身體，卻不能留在天上。因此，誰也不要因自己的本性而憂傷或煩惱，彷彿它妨礙了他，而要因自己的意志。魔鬼從無形體成為獅子；因為看，它說：\BibleQuote{我們的仇敵如同咆哮的獅子，巡遊尋找可吞吃的人}\BibleRef{伯前 5:8}；我們從有形體成為天使。就如同一個人找到某種珍貴材料，但因不是工匠而輕視它，對他便是大損失，無論他見到的是珍珠、珍珠貝、或其他任何東西；同樣，我們若不認識自己的本性，便會大大輕視它；但若知道它是什麼，便會表現極大的熱忱，並獲得最大的利益。因為從這本性製成君王的袍子，從這本性建成君王的房屋，從這本性形成君王的肢體：一切都是王家的。因此，我們不要為自己的利益而錯用我們的本性。祂使我們\BibleQuote{稍微遜於天使}\BibleRef{詠 8:5}——我是說，因死亡之故；但即使是那點微小，我們現已恢復。所以，沒有什麼能阻止我們接近天使，只要我們願意。那麼，讓我們願意，讓我們願意，並徹底鍛鍊自己，將光榮歸於聖父、聖子及聖神，從今時直到永遠，世世無窮。阿們。\par

\SourceRecord{Translated by J. Walker, J. Sheppard and H. Browne, and revised by George B. Stevens.；Nicene and Post-Nicene Fathers, First Series；Vol. 11.；Edited by Philip Schaff.；Buffalo, NY: Christian Literature Publishing Co.,；1889.；<http://www.newadvent.org/fathers/210132.htm>.}

% structure: p0001 source=h1 target=section reason=page-title

\section{论\WorkTitle{宗徒大事录}第三十三篇讲道}

\BibleRef{宗 15:13,15}\par

\Quote{他们沉默以后，雅各伯回答说：\InnerQuote{诸位弟兄，请听我说：\Person{西默盎}述说了天主当初怎样眷顾外邦人，从他们中选取一个属于他名下的民族。先知的话也与此相符。}}\par

此雅各伯如他们所说，是主教，因此他最后发言，于此应验了那句话：\BibleQuote{凭两三个见证的口，句句都可成立。}\BibleRef{申 17:6} 但也要观察他所显示的谨慎，从先知书中，无论新与旧，使他的论证成立。因为他没有像伯多禄和保禄那样宣布自己的事迹。而且，这样的安排是明智的，将（行动）部分指派给他们，因为他们不打算固定驻留在耶路撒冷；而这里的（雅各伯）执行教师的角色，对所发生之事毫无责任，但却不与他们意见分歧。 (\Emph{b}) \Quote{诸位弟兄，}他说，\Quote{请听我说。} 这人之节制何等伟大。他的言辞也更完整，因为它确实为讨论之事画上句号。 (\Emph{a}) \Quote{西默盎}他说，\Quote{述说了：}（即）在\BibleBook{路加}{福音}中，他预言说：\BibleQuote{就是你在万民面前所准备的，为启示外邦人的光明，为你百姓以色列的荣耀。}\BibleRef{路 2:25} (\Emph{c}) \Quote{天主起初眷顾外邦人，从他们中选取一个百姓归于自己的名下。}\BibleRef{路 2:25} 然后，既然那（见证）虽从时间上确实显明，但因不古老而缺乏权威，因此他也举出古老的预言，说：\BibleQuote{此后我要回来，重新修建达味已倒塌的帐幕；我要重建其废墟，并要把它竖立起来。}\BibleRef{宗 15:16} 什么？耶路撒冷被重建了吗？它不反而被拆毁了吗？他称那（从巴比伦回归后）所发生的事是什么样的重建？\BibleQuote{叫剩余的人，就是凡称我名的外邦人，都寻求主。}\BibleRef{宗 15:17} 然后，使他的话语有权威的是——\BibleQuote{行这一切事的主说：}以及，因为这不是新鲜事，而是从起初就计划好的，\BibleQuote{天主从永远就知道他的一切作为。}\BibleRef{宗 15:18} 然后再次他的权柄 \OriginalTerm{再次，他的权柄}{\GreekText{καὶ} \GreekText{τὸ} \GreekText{ἀξίωμα} \GreekText{πάλιν}}（作为主教）：\BibleQuote{因此我意见是，不要搅扰那归向天主的外邦人；但要写信给他们，叫他们禁戒偶像的污秽和奸淫，并勒死的牲畜和血。因为自古以来，在各城里都有宣讲梅瑟的人，每逢安息日在会堂里诵读他。}\BibleRef{宗 15:19} 既然他们已听到法律，有充分理由他吩咐这些来自法律的事，以免他显得使之无效。而且（注意）他如何没有让他们从法律中听到这些事，而是从他本人，说：我不是从法律听到这些，而是怎样？\Quote{我们已经判定。} 然后，法令是共同制定的。\BibleQuote{那时，宗徒和长老同全教会定意，从他们中间拣选人}——你注意到他们不仅制定这些事，而且不止于此？——\BibleQuote{差他们到安提约基雅同保禄和巴尔纳伯去；就是犹达号称巴尔撒巴，和息拉，弟兄们中的首领，并藉他们写信如下。}\BibleRef{宗 15:22} 并且注意，为使法令更加可信，他们派遣自己的人，以免有人怀疑保禄及其同伴。\BibleQuote{宗徒和长老并弟兄们，问安提约基雅、叙利亚和基里基雅的外邦众弟兄的安。}\BibleRef{宗 15:23} 且看他们以何等克制，避免对那些人严厉的责骂，写下书信。\BibleQuote{我们听说有几个人从我们这里出去，用言语搅扰你们，倾覆你们的心，说你们必须受割礼，守法律，其实我们并没有吩咐他们。}\BibleRef{宗 15:24} 这对那些人的鲁莽已是足够的指控，也配得上宗徒们的节制，因为他们没有多说其他。然后，为表明他们并非独断专行，而是众人一致同意，经过深思写下这些——\BibleQuote{我们同心合意，定意拣选几个人} \BibleRef{宗 15:25}——然后，为免使差遣那些人显得是贬低保禄和巴尔纳伯，请看对他们所加的赞美——\BibleQuote{同我们亲爱的巴尔纳伯和保禄，他们是为主耶稣基督的名不顾性命的。我们就差了犹达和息拉，他们也要亲口告诉你们同样的事。因为圣神和我们定意} ——这不是人的作为，它说——\BibleQuote{不将更大的重担加给你们}——再次称法律为重担：然后为这些禁令辩解——\BibleQuote{唯有这几件事是不可少的} \BibleRef{宗 15:26}：\Quote{就是禁戒祭偶像的物和血，并勒死的牲畜和奸淫；你们若能保守自己不沾染这些，就作得好了。}\BibleRef{宗 15:29} 因为这些事新约并未吩咐；我们找不出基督曾谈论这些事；但这些事出自法律。\Quote{勒死的牲畜，}它说，\Quote{和血。} 此处禁止杀人。\BibleRef{创 9:5} \BibleQuote{于是他们受差遣，下安提约基雅，聚集了众人，把书信交给他们。众人读了，因这安慰的话就欢喜了。}\BibleRef{宗 15:30-31} 然后那些人（弟兄们）也劝勉了他们：使他们坚固，因为保禄曾被他们争辩对待，所以平安地离开他们。\BibleQuote{犹达和息拉也是先知，就用许多话劝勉弟兄们，坚固他们。住了些日子，弟兄们打发他们平平安安地回到宗徒那里去。}\BibleRef{宗 15:32-33} 从此不再有派别和争斗，而是保禄从那时起教导。\par

（重述。）\BibleQuote{于是全体群众都默不作声}等等。\BibleRef{宗 15:12}在教会中没有傲慢。伯多禄之后，保禄发言，无人制止他；雅各伯耐心等待，没有跳起来（为接续发言）。秩序井然（会程）。若望在此没有说一句话，其他宗徒也没有说话，都保持沉默，因为雅各伯被赋予了首席权，他们也毫不介意。他们的灵魂如此洁净，不受爱光荣的玷污。\BibleQuote{他们沉默之后，雅各伯回答说}等等。\BibleRef{宗 15:13}（\Emph{b}）伯多禄说得更强烈，但雅各伯在此却更温和：因为身居高位者应当如此，把不悦的话留给别人说，而自己则显得比较温和。（\Emph{a}）但\Quote{天主如何\OriginalTerm{首先}{\GreekText{πρὥτον}}眷顾}是什么意思？\BibleRef{宗 15:14}（意思是）\OriginalTerm{从起初}{\GreekText{ἐξ} \GreekText{ἀρχἥς}}。（\Emph{c}）另外，他很好地说道：\Quote{西默盎\OriginalTerm{解释}{\GreekText{ἐ} \GreekText{ξηγήσατο}}}（或诠释），暗示他也在表达别人的意见。\Quote{与此一致}等等。请注意他如何表明这是一项古老的教义。\Quote{从外邦人中}，他说，\Quote{取出一族归于祂的名下。}\BibleRef{宗 15:15}不单是拣选，而是\Quote{为祂的名}，即为了祂的光荣。祂的名不因\OriginalTerm{先取}{\GreekText{προλήψει}}外邦人而蒙羞，反而成为更大的光荣。——这里甚至暗示了一件大事：这些人是在众人之前被选中的。\BibleQuote{此后我要回来，重建达味已倒塌的帐幕。}\BibleRef{宗 15:16}但如果有人仔细考察这件事，达味的王国如今确实屹立不倒，他的后裔在各地为王。如果那里无人顺服，建筑物和城市又有何用？如果所有人都愿意献出自己的灵魂，城市的毁灭又造成什么损害？那比达味更光荣的一位已经来临：如今世界各地都在歌颂祂。这事已经实现：如果是这样，那么这事也必须实现：\BibleQuote{我要重新建造它的废墟，并把它竖起：}目的是什么？\BibleQuote{使剩余的人寻求主，并所有在外邦中，凡呼求我名的人。}\BibleRef{宗 15:17}如果城市复兴是为了这个目的（即因他们中将要来的那一位），那就表明城市建造的原因是外邦人的蒙召。谁是\Quote{剩余的人？}是那时剩下的人。\Quote{并所有在外邦中，凡呼求我名的人：}但请注意，他如何保持适当顺序，将他们放在第二位。\Quote{说这话的是主，祂做这些事。}不仅是\Quote{说}（只是），而是\Quote{做。}那么，这原是天主的工作。——\Quote{但问题不在于此（即伯多禄更清楚地说过，他们必须受割礼吗？）。那么你为什么滔滔不绝地讲这些事情？}因为反对者所断言的不是信德不能使人得救，而是必须遵守法律。伯多禄对此进行了有力的辩护：但此事比任何事都更困扰听众，因此他再度\OriginalTerm{纠正}{\GreekText{θεραπεύει}}这一点。请注意，他引入了必须作为规则的规定，即不必遵守法律；但雅各伯说出较温和的部分，即古老的真理，并详述那没有写下的东西，以便用公认的道理安抚他们的心思后，再适时地引入这点。\Quote{因此，}他说，\Quote{我的判断是，不要为难那些从外邦人归向天主的人}\BibleRef{宗 15:19}，即不要颠覆：因为如果天主召叫了他们，而这些规条颠覆他们，我们就是在对抗天主。他又说：\Quote{那些从外邦人中} \Quote{归向的人。}他带着权威说得好：\Quote{\Emph{我的}判断是。但我们要写信给他们，叫他们禁戒偶像的污染和奸淫}——（\Emph{b}）然而他们在教导时常强调这些点——但为了显得也尊重法律（他提及），他也提到这些，但说话（不过）不是出自梅瑟，而是出自宗徒；为了增加诫命数量，他把一条分为两条说：\Quote{并禁戒勒死的牲畜和血。}\BibleRef{宗 15:20}因为这些虽然与身体有关，但必须遵守，因为（这些事）造成了极大祸害，\Quote{因为梅瑟自古以来在各城}等等。\BibleRef{宗 15:21}这事最能使他们安静（\OriginalTerm{安抚}{\GreekText{ἀ} \GreekText{νέπαυσεν}}）。（\Emph{a}）为此我断定这样做很好（即\Quote{写信给他们}）。那么，我们为何不也给犹太人相同的命令呢？梅瑟向他们讲论。看，对他们的软弱是何等迁就！在无碍之处，他立梅瑟为导师，并在不阻碍任何事的情况下，允许犹太人听梅瑟的话，甚至同时引导外邦人离开梅瑟。何等智慧！他看似尊崇梅瑟，立他为自家子民的权威，而正借此把外邦人从他引开！\Quote{每逢安息日在会堂里诵读。}那么，他们为何不从他那里学习（该学的东西），例如**？由于这些人的执拗。他表明甚至这些人（犹太人）也无需遵守更多（这些必要之事）。我们若没有写信给他们，并不是因为他们有义务遵守更多，只是因为他们有人告诉他们。他没有说\Quote{不得冒犯}，也没有说\Quote{不要使他们回头}，那是保禄对迦拉达人说的，而是\Quote{不要为难他们：}他表明这\OriginalTerm{要点}{\GreekText{κατόρθωμα}}若执行，不过是徒增困扰。就这样，他结束了整个事件；他看似通过采纳这些规则来维护法律，实际上只取这些，解除了法律。（\Emph{c}）在辩论中也有上智的安排，使辩论之后，教义更加稳固。\Quote{于是宗徒们决意从他们中选派几个人}等等，不是普通人，而是\Quote{首领；写了}（信）\Quote{叫他们带去，信上说：致安提约基雅}，经上说，\Quote{和叙利亚并基里基雅。}（第22、23节）那里是病源的发源地。请注意，他们没有说任何\OriginalTerm{更严厉}{\GreekText{φορτικώτερον}}的话反对那些人，只关注一件事，即消除已造成的损害。因为这甚至会使那里的煽动者认错。他们没有说\Quote{诱惑者}、\Quote{祸患}之类的话；虽然需要时，保禄会这样做，例如他说：\Quote{你这充满各种诡诈的人}\BibleRef{宗 13:10}；但在此，论点已获胜，没有必要。请注意，他们没有说\Quote{我们中的某些人命令你们遵守法律}，而是\Quote{用言语扰乱你们，颠覆你们的心神}——没有比这个词\OriginalTerm{更恰当}{\GreekText{κυριώτερον}}的了：没有（其他发言者）这样描述那些人的所作所为。\Quote{心神}，他说，已经坚定建立，这些人却像在谈论建筑时那样，\Quote{再次拆下}（\OriginalTerm{拆毁}{\GreekText{ἀ} \GreekText{νασχευάζοντες}}），把他们从根基上\OriginalTerm{转移}{\GreekText{μετατιθέντες}}。\Quote{对他们}他说，\Quote{我们并没有给他们这样的命令。因此我们同心合意，决意选派几个人到你们那里，同我们亲爱的巴尔纳伯和保禄一起，他们都是为我们的主耶稣基督的名而冒过生命危险的人。}（第25、26节）如果说是\Quote{亲爱的}，他们就不会轻视他们；如果说\Quote{冒过生命危险}，他们自己就有权被相信。\Quote{我们已派了}，经上说，\Quote{犹达和息拉，他们也必亲口报告同样的事。}\BibleRef{宗 15:27}因为仅仅有书信在那里是不够的，免得他们说保禄和巴尔纳伯隐晦（真意），说一件事代替另一件事。对保禄的称赞堵住了他们的口。这就是为什么保禄不独来，巴尔纳伯也不独行，而是还有教会中的其他人同来；使他不被怀疑，因为他正是倡导那教义的人；也不是只从耶路撒冷派人。这表明他们有权被相信。\Quote{因为决意}他们说，\Quote{对圣神和我们}\BibleRef{宗 15:28}：并不是使自己与（圣神）相等——他们不至于如此疯狂。但为何这样写呢？为何他们加上\Quote{和我们}，而只说\Quote{圣神}就已经足够？一是\Quote{圣神}，叫他们不以为是出于人；另一是\Quote{和我们}，叫他们知道他们自己也接纳（外邦人），尽管他们自己仍守割礼。他们要对那些仍然软弱并害怕他们的人说话：这就是为何加上这一点的原因。这也表明他们说话不是出于迁就，也不是因怜悯他们或认为他们软弱，而是相反；因为师者也怀着极大的敬意。\Quote{不再加给你们什么重担}——他们总称之为重担——又说，\Quote{除了这些必要的事：}因为那是多余的重担。看，一封简短的书信，其中没有多余之物，既无\OriginalTerm{说服技巧}{\GreekText{κατασκευὰς}}，也无推理，只是简单的命令：因为那是圣神的立法。\BibleQuote{他们既被派遣，就下到安提约基雅，聚集了众人，把书信递给他们。}\BibleRef{宗 15:30}书信之后，犹达和息拉也亲自用言语劝勉他们\BibleRef{宗 15:31}：因为这也必要，使保禄和巴尔纳伯完全脱离嫌疑。\Quote{他们二人也是先知，}经上说，\Quote{用许多话}劝勉弟兄们。这里显示保禄和巴尔纳伯有权被相信。因为保禄本也可以这样做，但理当由这些人来做。他们住了些日子，就平安地打发他们去了。\BibleRef{宗 15:33}\par

不再有派系。我想，正是在此时，他们握了右手，正如他本人所说：\Quote{他们把共融的右手给了我和巴尔纳伯。}\BibleRef{迦 2:9}那里他说：\Quote{他们并没有加给我什么。}\BibleRef{迦 2:6}因为他们证实了他的看法：他们称赞并钦佩它。——这表明，即使从人的推理也能看出这一点，更不用说仅从圣神了，他们犯了一个不易纠正的罪。因为这样的事并不需要圣神。——这表明其余的事并非必要，而是多余的，既然这些事是必要的。它说：\Quote{你们若保守自己不受这些，就必行得好。}这表明他们一无所缺，但这就足够了。因为这件事本也可以不经书信而完成，但为了有一条成文的法律（他们寄出了这封信）；再者，为了使宗徒们遵守法律，他们也向那些人（讲了同样的事），他们便这样做了，\Quote{并坚固了他们，又住了一些时候，才平安地被释放离开。}\par

那么，我们不要因异端而跌倒。看，在这传道的最初阶段，有多少绊脚石：我不说那些由外人而来的，因为这些算不了什么；而是那内在的绊脚石。例如，先是阿纳尼雅，然后是\Quote{抱怨，}接着是术士西满；后来他们因科尔乃略的事控告伯多禄，接着是饥荒，最后是这最大的邪恶。因为任何好事发生，难免有坏事随之而来。我们不要因有人跌倒就烦乱，而要为此感谢天主，因为这使我们更受考验。不仅磨难，连诱惑也使我们更加卓越。人若在没有被引诱离开的情况下坚持真理，算不上伟大的真理爱好者；但当许多人拉他离开时仍坚守真理，这才显出他的坚定。那么，绊脚石出现是为了这个吗？我不是说天主是这些的根源，绝对不可！我的意思是，甚至从他们的邪恶中，天主也为我们生出益处：祂从未希望它们发生。祂说：\Quote{求祢赐给他们，}祂说，\Quote{使他们合而为一}\BibleRef{若 17:21}。但既然绊脚石来了，它们对这些人无害，反而有益：就像迫害者不情愿地将殉道者拖向殉道，从而有益于他们，然而他们并非受天主驱使；同样在这里。我们不要只看到有人跌倒，这本身就是教义卓越的证明——许多人刺激和模仿它：因为如果教义不好，就不会这样。我现在要显明这点，并多方使它清楚。香料中，芬芳的香料常被人掺假和伪造，例如荜澄茄叶。因为它们稀有且必需，所以才会有冒牌货。没人会费心去伪造普通物品。纯洁的生活有许多假冒者，没人会去假冒邪恶的生活，但会去假冒隐修士的生活。——那么，我们对外教人该说什么？有一个外教人来，说：\Quote{我希望成为基督徒，但我不知该加入谁：你们中间有许多争斗和党派，许多混乱：我该选择哪个教义？}我们如何回答他？他说：\Quote{你们每个人都断言：\InnerQuote{\Emph{我}讲的是真理。}}（\Emph{b}）当然，这有利于我们。如果我们叫你用论证说服，你可能会困惑；但如果我们叫你相信圣经，而圣经是简单真实的，决定对你很容易。如果有人符合圣经，他就是基督徒；如果有人反对圣经，他就远离这个准则。（\Emph{a}）\Quote{但我该相信谁呢？我连圣经一点也不懂。其他的人也为自己声称同样的事。那么（\Emph{c}）如果另一个人来说圣经有这个，而你说它有不同，并且你们各按己意拉扯经文的意义？}那么我问你，难道你没有理解力，没有判断力吗？他说：\Quote{我怎么能判断呢？我甚至不知如何判断你们的教义。我想做个学习者，而你们立刻使我成了老师。}如果他这样说，你，你对我们该如何回答？我们如何说服他？让我们问问，这一切是否只是借口和托词。让我们问他，他是否已经\OriginalTerm{决定}{\GreekText{κατ}\GreekText{εγνωκε}}外教人是错的。他肯定会肯定这个事实，因为当然，如果他没这样决定，他就不会来询问我们的事：让我们问他决定的基础，因为他肯定不是草率决定。显然他会说：\Quote{因为（他们的神）是受造物，不是非受造的天主。}好。那么如果他在其他\OriginalTerm{党派}{\GreekText{α}\GreekText{ιρ}\GreekText{εσεις}}中发现这点，而在我们中间发现相反，我们还需要什么论证？我们大家都承认基督是天主。但让我们看谁在反对（这真理），谁没有。现在我们肯定祂是天主，说祂相称于天主的事：祂有权能，祂不是奴隶，祂是自由的，祂自己做主；而另一个人说相反的。我再问：如果你想学医，* * *？然而他们中间有许多（不同的）学说。因为如果你毫不犹豫地接受别人告诉你的，这不是人的作为；但如果你有判断力和感觉，你定能知道什么是好的。我们肯定子是天主，我们\OriginalTerm{验证}{\GreekText{επαληθε}\GreekText{υομεν}}我们所肯定的；但他们虽然肯定，却实际不承认。——但提到一个更明显的事：那些人有一些人，他们由此得名，公开指明异端创始人本身的名字，每个异端也都如此；对我们，没有人给我们一个名字，而是信仰本身。然而，你这话只是借口和托词。你回答我：如果你要买一件斗篷，虽然你不懂编织技术，你怎不说这样的话：\Quote{我不知道怎么买，他们骗我}，反而尽力去学？同样，买任何东西都是如此；但在这里你说这些话？照此，你将不接受任何东西。因为假若有一个人没有任何（宗教）学说，他如你所说关于基督徒的话：\Quote{有这么多的人，他们有不同学说；这个是外教人，那个是犹太人，另一个是基督徒：不需要接受任何学说，因为彼此不一；但我是一个学习者，不想做判断}——但如果你已经让步到宣告\OriginalTerm{反对}{\GreekText{καταγιν}\GreekText{ωσκειν}}一个学说，这个借口对你不再成立。因为就如你能拒绝假的，同样，来到此处，你也能证明什么是有益的。因为那没有宣告反对任何学说的人，很容易说这话；但那已经宣告反对某学说的人，虽未选择任何学说，继续以同样方式，便能看见他该做什么。那么，我们不要找借口和托词，一切都会容易。因为为了向你显示这一切只是借口，回答我：你知道你该做什么，不该做什么吗？那么你为什么不做该做的？做那事，并凭正直的理智寻求天主，祂定会启示你。它说：\Quote{天主是不偏待人的，但在各邦国中，那敬畏祂并履行正义的，都蒙祂悦纳。}（\EditorialNote{宗 10:34-35}）没有偏见的人不可能不被说服。因为就如有一条规则，一切事物都应照它摆正，那就不需要多考虑，而容易发现\OriginalTerm{量度有误的人}{\GreekText{τ}\GreekText{ον} \GreekText{παραμετρο}\GreekText{υντα} \GreekText{λαβ}\GreekText{ειν}}；同样在这里。\Quote{那么他们为何不一眼看出？}许多事情是原因：既有先入之见，也有人为的\OriginalTerm{原因}{\GreekText{α}\GreekText{ιτ}\GreekText{ιαι}}。你说，其他的人对我们说同样的话。怎么？难道我们与教会分离？我们有我们的异端创始人？我们随人名而得名——比如他们中有马尔基昂、摩尼、亚略作其创始人和领袖？而如果我们也从某个人那里接受一个别名，我们不是取那些作成某些异端的创始人，而是那些在我们之上统领、管理教会的人。我们没有\Quote{地上的师傅}（求主宽恕），我们有\Quote{一位在天上的师傅}\BibleRef{玛 23:9}。他说：\Quote{那些人，}他说，\Quote{也说同样的话。}但是，凌驾于他们之上的名字指责他们，塞住他们的口。——怎么有许多外教人，没有一人问这些问题；在哲学家中间也有这些（分歧），但那些持守正确\OriginalTerm{党派}{\GreekText{α}\GreekText{ιρεσιν}}的人没有一个因此受阻碍？——当（其他人）提出这些问题时，（那些信徒）为什么不说：\Quote{这些和那些都是犹太人：我们该信谁？}但他们按应当的方式相信了。那么，我们也服从天主的法律，按祂的圣意做一切，以便我们现世善度后，能获得为爱祂的人所预许的美善，赖我们的主耶稣基督的恩宠和仁慈，祂与圣父及圣神一同，享受光荣、权能、尊威，从现在直到永远，世世无尽。阿们。\par

\SourceRecord{Translated by J. Walker, J. Sheppard and H. Browne, and revised by George B. Stevens.；Nicene and Post-Nicene Fathers, First Series；Vol. 11.；Edited by Philip Schaff.；Buffalo, NY: Christian Literature Publishing Co.,；1889.；<http://www.newadvent.org/fathers/210133.htm>.}

% structure: p0001 source=h1 target=section reason=page-title

\section{讲道第三十四篇：\WorkTitle{宗徒大事录}}

\BibleRef{宗 15:35}\par

\BibleQuote{保禄和巴尔纳伯继续留在安提约基雅，与许多人一同教导和宣讲主的圣言。}\par

请再观察他们的谦卑，他们如何容许他人也参与宣讲。\BibleQuote{几天后，保禄对巴尔纳伯说：让我们回到我们曾宣讲主圣言的各城，去探望弟兄们，看看他们怎么样了。巴尔纳伯决定也带着那名叫马尔谷的若望同去。但保禄认为不妥（\OriginalTerm{不妥}{\GreekText{ἡ} \GreekText{ξίου}} \EditorialNote{见注释3，第213页}）带他同去，因为他从前在旁非里雅离开了他们，没有同去工作。于是二人之间起了激烈的争论（或激怒），以致彼此分离。}\BibleRef{宗 15:36} 路加早已向我们描述了宗徒们的性格：一位较为温和宽厚，另一位则较为严肃刚毅。因为恩赐是多样化的——（我说的是恩赐），因为显然这是一种恩赐——但一位适合一种性格，另一位适合另一种性格，如果互换，则有害无益。（\Emph{b}）在先知们身上我们也发现这一点：不同的心态，不同的性格：例如，厄里亚严肃，梅瑟温和。在此保禄较为激烈。且请观察，尽管如此，他是多么温和。\BibleQuote{认为不妥，}它说，\BibleQuote{不要带他同去，因为他从前从旁非里雅离开了他们。}（\Emph{a}）似乎确实有一场\OriginalTerm{争论}{\GreekText{παροξυσμός}}，但实际上整件事是天主眷顾的计划，好使每人得到自己适当的位置：并且他们不该平起平坐，而应一人领导，另一人服从。\BibleQuote{于是巴尔纳伯带了马尔谷，乘船往塞浦路斯去；保禄却选了息拉，蒙弟兄们交托于天主的恩宠后，就出发了。他走遍了叙利亚和基里基雅，坚固各教会。}\BibleRef{宗 15:39} 这也是上智的安排。因为塞浦路斯人没有显露出像安提约基雅和其他地方那样的事，那些人需要较温和的性格，而这些则需要保禄这样的性格。\Quote{那么，}你说，\Quote{哪一位做得对？带走他的，还是离开他的？}***（\Emph{c}）因为正如一位将军不会选择让一个低贱的人常作他的行李搬运工，宗徒也是如此。这纠正了另一位，并教训了（马尔谷）本人。\Quote{那么巴尔纳伯做错了？}你说。\Quote{在一件如此小的事上竟产生如此大的恶，这如何不是\OriginalTerm{不适当}{\GreekText{ἄτοπον}}呢？}首先，并没有产生恶，因为二人各自身负整个民族之需，分开彼此，反倒是一桩大善事。况且，他们本不愿轻易分离。但我请你佩服史家，他连这一点也不隐瞒。\Quote{但无论如何，}你说，\Quote{如果他们必须分离，让他们毫无激愤地分离。}不，但即使没有别的，请观察这一点：在此也显明了人性在福音宣讲中的参与。因为如果类似的事也必须在基督所做的事上显明，那么在此更当如此。此外，当每人为了如此的目的而争论，且有公正的理由时，这争论不能说为恶。我承认，如果激愤是在谋求私利、为自己争荣誉，这固然应受责备；但如果两人都愿意教导训诲，一人取此道，另一人取彼道，又有何可指责？因为在许多事上，他们按人性判断行事；他们不是木石。且观察保禄如何指责（马尔谷），并说明理由。因他极其谦卑，敬重巴尔纳伯，因他曾与自己同工完成如此大的事业，且与他同在；但他并未因此敬重到忽略必要之事。现在谁的意见最好，不是我们可以断定的：但至少我们可以这样说，这是天主眷顾的伟大安排，如果这些人蒙恩得第二次探访，而那些人甚至一次也未被探访。\par

\Emph{a} \Quote{教导并宣讲主的圣言。} \BibleRef{宗 15:35} 他们不仅仅在安提约基雅停留，还进行教导。他们\Quote{教导}什么，又\Quote{宣讲}（传福音）什么？他们一面（教导）那些已经相信的人，一面（传福音）给那些尚未相信的人。\Quote{过了些日子，}等等。\BibleRef{宗 15:36} 因为有无数的阻碍，所以他们的在场是必要的。\Emph{d} \Quote{他们怎样了？}他说。这一点他原先不知道：自然如此。请看他是何等警醒、关切，不愿闲坐，尽管他历经无数危险。你注意到了吗？他来安提约基雅并非出于怯懦？他正像一个医生对待病人那样行动。他藉着说\Quote{我们在那里宣讲了圣言。巴尔纳伯决定，}等等，\BibleRef{宗 15:37}（于是）巴尔纳伯\Quote{就离开了，没有与他同去。}\Emph{b} 应当注意的，不是他们在意见上有分歧，而是他们彼此迁就（看到），这样分离是一件更大的好事：此事也以此为借口。那么怎样？他们怀着敌意离开吗？天主不许！事实上，此后你在书信中看到巴尔纳伯从保禄那里得到许多称赞。经文说，发生了\Quote{激烈的争论，}并不是敌意也不是争吵。争论足以使他们分开。\Quote{巴尔纳伯带了马尔谷，}等等。而且有理由：因为每个人以为有益的事，不会因与另一人的共融而放弃。不仅如此，在我看来，这次分离是\OriginalTerm{明智地}{\GreekText{κατὰ} \GreekText{σύνεσιν}}发生的，他们彼此说：\Quote{既然我不愿意，你愿意，因此我们不可争闹，让我们分派地方吧。}所以他们实际上这样做，完全彼此让步：因为巴尔纳伯希望保禄的计划成立，因此离开；另一方面，保禄希望另一人的计划成立，因此他离开。但愿我们也做出这样的分离，为出去宣讲。这是一个何等奇妙的人，极其伟大！对于马尔谷，这场争论极为有益。因为保禄所激发的敬畏使他皈依，而巴尔纳伯的仁慈使他未被抛弃：所以他们确实争论，但益处归于同一目的。因为，看见保禄选择离开他，他会极其敬畏，并自责；看见巴尔纳伯如此站在他一边，他会极其爱他；如此，门徒被师傅们的争论纠正了：他远未因此被绊倒：因为如果他们是为自己的名誉这样做，他可能会被绊倒；但如果是为了他的救恩，而且他们为同一目标争论，表明那曾尊重他的人* * * 已妥善决定，这有什么\OriginalTerm{不合理}{\GreekText{ἄτοπον}}呢？\par

\Emph{e} \Quote{但是保禄，}经上说，\Quote{离开了，选择了息拉，被交托于天主的恩宠。}这是什么？经文说他们祈祷了：他们恳求了天主。请看，在任何情况下，弟兄们的祈祷能成就大事。如今他陆路旅行，甚至希望通过他的旅行使那些\OriginalTerm{看见他的人}{\GreekText{τοὺς} \GreekText{ὁρὥντας}}受益。因为当赶时间时，他们乘船，但如今却不是这样。\Emph{c} \Quote{他经过叙利亚和基里基雅，坚固各教会。然后他来到德尔贝和吕斯特辣。}\BibleRef{宗 15:41} 注意保禄的智慧：他不去其他城市，直到他先探访了那些已接受圣言的地方。因为随意乱跑是愚昧。让我们也这样做：让我们首先教导第一批人，这样他们就不会成为后来者的阻碍。\par

\BibleQuote{看，在那里有一个门徒，名叫弟茂德，是一个信主的犹太妇人的儿子，但他父亲是希腊人。他在吕斯特辣和依科尼雍的弟兄们那里有很好的名声。保禄愿意他与他一同出去，就带了他，因那地方的犹太人而给他行了割损礼，因为他们都知道他父亲是希腊人。} \BibleRef{宗 16:1} 保禄的智慧实在令人惊奇！他为了割损礼经历了那么多的争斗，调动一切以达到此目的，不达目的不罢休，如今这道命令已确认，他却给门徒行了割损礼。他不仅不禁止别人，而且亲自做了这事。 (\Emph{b}) 经上说：\BibleQuote{他愿意他与他一同出去。}令人惊奇的是，他甚至接纳了他。经上说：\BibleQuote{因为那地方的犹太人，}因为他们不愿从未受割损的人那里听到道。 (\Emph{a}) 没有比这更明智的了。因此他凡事都看有益处：他做事从不凭自己的\OriginalTerm{偏好}{\GreekText{προλήψει}}。 (\Emph{c}) 那么，结果呢？看那成功：他行了割损，为要除去割损；因为他宣讲了宗徒们的命令。\BibleQuote{他们经过各城，把耶路撒冷宗徒和长老所定的法规交给他们遵守。於是各教会信德坚固，人数天天增加。}\BibleRef{宗 16:4,5}你看到争战，并藉着争战建造吗？不是被他人攻击，而是他们自己做相反的事，就这样建造了教会！他们颁布一条不割损的命令，而他行了割损！经上说：\BibleQuote{各教会}在信德中\BibleQuote{坚固}，并且人数\BibleQuote{天天增加}。然后他没有继续与这些人同住，因为他是来探望他们的。而是怎样？他继续前行。\BibleQuote{他们走遍了夫黎基雅和迦拉达地区，被圣神禁止在亚细亚讲道，} \BibleRef{宗 16:6} 离开了夫黎基雅和迦拉达，他们急忙进入内地。因为经上说：\BibleQuote{他们到了米息雅，试图往彼提尼雅去，但圣神不允许他们。} \BibleRef{宗 16:7} 为何他们被禁止，他没有说，但他说他们被\Quote{禁止}，教导我们服从而不问问题，并表明他们做了许多人性的事。经上说：\Quote{圣神不允许他们；他们经过米息雅，下到特洛阿。} \BibleRef{宗 16:8} \BibleQuote{夜间有异象显现给保禄：有一个马其顿人站着求他说：请你过到马其顿来帮助我们。} \BibleRef{宗 16:9} 为什么是异象，而不是圣神？因为圣神禁止了另一种方式。祂也希望用这种方式引导他们：因为对圣人们祂也在梦中显现，起初（保禄）自己也看见一个异象：\BibleQuote{一个人进来，按手在他身上。} \BibleRef{宗 9:12} 同样，基督也以这种方式向他显现，说：\Quote{你必须站在凯撒面前。} 为此祂也把他引到那里，使宣讲得以扩展。这就是为什么他被禁止在其他城市久留，基督催促他前行。因为这些城市的人长久享受若望的好处，也许并不十分需要他（保禄），但他必须去那里。现在他渡海出去。\BibleQuote{他见了这异象，我们随即设法往马其顿去，推断是主召我们去给他们传福音。} \BibleRef{宗 16:10} 然后作者也提到地点，如叙述历史一般，表明他在何处停留（即）在大城市，而路过其余地方。\BibleQuote{我们从特洛阿开船，直航撒摩辣刻，第二天到了乃阿颇里；从那里到斐理伯，那是马其顿这一带的首府，是罗马的驻防城。}\BibleRef{宗 16:11,12}对一座城来说，成为驻防城是极高的荣誉。\Quote{我们在那城里住了几天。}但让我们再审视一遍所说的话。\par

（复述。）\Quote{过了些日子，保禄说}等等。\BibleRef{宗 15:36}他向巴尔纳伯说明他们必须外出，说：\Quote{我们去探望那些我们曾宣讲主道的城市。}\Quote{但保禄恳求}等等。\BibleRef{宗 15:38}其实他无需恳求，因为他马上就要提出指控。这种事在天主与人之间也会发生：人请求，天主发怒。例如，当祂说：\Quote{如果她父亲吐唾沫在她脸上}\BibleRef{户 12:14}，又说：\Quote{你且由我，我要在怒气中消灭这人民}\BibleRef{出 32:32}，以及撒慕尔为撒乌耳哀悼时\BibleRef{撒上 15:35}。因为藉着两者，都成就了极大的善事。此处也是如此：一方发怒，另一方却不如此。同样的事也发生在与我们相关的事情上。那激烈的争论是合情合理的，为要使马尔谷受到教训，并且这事看起来不像演戏。因为不可认为那位命令\Quote{不可让太阳在你们含怒时落下}\BibleRef{弗 4:26}的人，会因这样的事而发怒；也不可认为那位在一切场合都让步的人，会在这里不让步，他如此热爱保禄，以前曾在塔尔索寻找他，带他到宗徒们面前，与他共同承担施舍之事，并共同处理有关法令的事务。他们却这样分开，为要教导和成全那些需要从他们那里接受教导的人。他确实责备别人，但劝勉众人向善。正如他在别处所说：\Quote{但你们，行善不要厌倦。}\BibleRef{得后 3:13}我们在平常的做法中也是如此。在我看来，这里似乎也有其他人对保禄不悦。于是他把他们也叫到一边，做了一切事，并劝诫和训导。和好能成就大事，爱德也能成就大事。即使你为一件大事祈求，即使你不配，你也会因你心中的意向而被垂听：不要害怕。\par

\Quote{他走遍}那些城市，经上说：\Quote{看，那里有一个门徒，名叫弟茂德，在吕斯特辣和依科尼雍的弟兄们中都有好名声。}（第41节；16:1。）弟茂德的恩宠极大。当巴尔纳伯\OriginalTerm{离开}{\GreekText{ἀπέστη}}时，他找到了另一个同等的人。关于他，保禄说：\Quote{我怀念你的眼泪和你那无伪的信德，这信德首先存在你外祖母罗依和你母亲欧尼刻心中。}\BibleRef{弟后 1:5}他的父亲一直是外邦人，因此弟茂德没有受割损。（\Emph{a}）注意法律已被打破。若非如此，我想他是在福音宣讲之后出生的，但也许并非如此。（\Emph{c}）他即将使他成为主教，而他不宜未受割损。（\Emph{e}）这不是小事，因为过了这么久它会得罪人：（\Emph{b}）\Quote{从小}，他说，\Quote{你就认识圣经。}\BibleRef{弟后 3:15}（\Emph{d}）\Quote{他们走遍各城，将那些规条交给他们遵守。}\BibleRef{宗 16:4}因为直到那时，外邦人无需遵守任何此类规条。废止的开始在于外邦人不守这些，却未受损害，在信德上也没有任何不足；不久后，他们自愿放弃了法律。（\Emph{f}）既然他要宣讲，为避免给犹太人双重打击，他给弟茂德行了割损。但他只有一半（犹太血统），父亲是希腊人；然而，因为那是在外邦人的事业中赢得的一个重大胜利，他并不在乎这点：因为圣言必须传播；所以他亲手给他行了割损。\Quote{于是各教会信德越发稳固。}你注意到在此处如何从反方向（他的目标）产生了极大的好处吗？\Quote{数目天天增加。}\BibleRef{宗 16:5}你观察到割损不仅无害，甚至大有裨益吗？\Quote{夜间保禄看见一个异象。}\BibleRef{宗 16:9}这次不是藉着天使，像对斐理伯、对科尔乃略那样，而是如何？藉着异象现在向他显示：以更人性的方式，而不是像先前（即6、7节）以更神圣的方式。因为哪里更容易顺从，就以更人性的方式成就；但哪里需要大力，就以更神圣的方式。既然他仅仅是被催促去宣讲，为此通过梦境向他显示；但要他不宣讲，他难以忍受；为此圣神向他启示。当时对伯多禄也是如此：\Quote{起来，下去。}\BibleRef{宗 10:20}当然，圣神不会去做那原本容易的事；但（这里）甚至一个梦就足够了。对若瑟，因他易于顺从，显现也在梦中，但对其他人则在清醒的异象中。\BibleRef{玛 1:20}这样对科尔乃略，对保禄自己也是如此。\Quote{看，有一个马其顿人}等等，不只是命令，而是\Quote{恳求}，而且是来自那些需要（灵性）医治的人本身。\EditorialNote{见：宗 10:3; 9:3}\Quote{我们便断定}，经上说，\Quote{是主召叫了我们。}\BibleRef{宗 16:10}也就是说，推断，既因保禄看到而别人没有，也因曾\Quote{被圣神禁止}，也因他们在边界上；从这一切他们推断。\Quote{于是从特洛阿开船，直航}等等。\BibleRef{宗 16:11}也就是说，就连航行也证明了这一点：因为没有延迟。它成了马其顿的根基。圣神并非总是在\Quote{激烈争论}中工作；但这（圣言的）如此迅速进展是事情超乎人性的标记。然而并没有说巴尔纳伯被激怒，而是\Quote{他们之间起了激烈的争论。}\BibleRef{宗 15:39}如果一方没有被激怒，另一方也没有。\par

知道了这事，我们不要只是\OriginalTerm{挑选}{\GreekText{ἐκλέγωμεν}}这些事，而要从它们学习并受教：因为它们不是无缘无故写下的。不熟悉圣经是巨大的恶：我们从本该得益的事物中反而招致恶。同样，具有治疗功效的药物，常常因使用者的无知而带来毁坏和毁灭；用于保护的武器，如果不知道如何佩戴，它们自己就成了死亡的原因。但原因是我们追求一切，唯独不求对我们自己有益的事。就房屋而言，我们寻求对它有益的事，不忍心看它因年久失修而朽坏、倾斜或受风暴损坏；但对于我们的灵魂，我们却毫不在意；不但如此，即使我们看到它的地基朽坏，或结构及屋顶损坏，我们也不在乎。再者，如果我们拥有牲畜，我们寻求对它们有益的事：我们请来喂马的和马医，以及所有其他；我们留意它们的住所，并嘱咐受托照管它们的人，不可随意或粗心地驱赶它们，不可在不适当的时间夜间带它们出去，也不可卖掉它们的饲料；我们为牲畜的利益制定了许多规定；但对灵魂却毫不顾及。但我为什么要提对我们有用的牲畜呢？有许多人养着小鸟（或\Quote{麻雀}），它们除了供人消遣外一无所用，但甚至对它们也有许多规定，没有一样被忽略或无序，我们为一切操心，唯独不为自身操心。这样我们就使自己比一切都更无价值。如果有人辱骂我们是\Quote{狗}，我们就恼怒；但我们自己侮辱自己，不是用言语，而是用行为，并且连对灵魂的关心都不如对狗的关心，我们却不以为意。你们看到一切都是多么充满黑暗吗？多少人关心他们的狗，不让它们吃得过饱，好让它们因饥渴而变得敏锐、适合打猎；但对自己却毫不关心避免奢侈：他们教导牲畜操练哲学，却让自己沉沦到牲畜的野蛮中。这真是一件怪事。\Quote{你那些懂哲学的牲畜在哪里？}有啊；或者，你说，当一条狗被饥饿折磨，在追捕并抓获猎物后，却禁食不吃；虽然它看到食物摆在面前，饥饿催逼着它，却仍然等待主人，难道你不认为这是哲学吗？你们应当羞愧：教导你们的肚腹也如此有哲理。你们没有借口。既然你们能在无理性的、既不能说话也不能听理性道理的本性中植入如此有哲理的自我控制，难道不能更容易地在自己身上植入它吗？因为这显然是人为关心的结果，而不是天性：否则所有的狗都应有这种习惯。那么你们就变成像狗一样吧。因为是你们逼我从那里取例证：本来应该从天上事物中取例，但因为我说起那些时，你们说\Quote{那些太高深了}，所以我就不谈天上事物；再者，我提到保禄时，你们说\Quote{他是宗徒}，所以我也不提保禄；如果我提到一个人，你们说\Quote{那人能做到}，所以我连人也不提，而是提一个牲畜，而且是一个本无此习惯的牲畜，免得你们说它天性如此，而不是出于选择（事实如此）；更奇妙的是，这选择并非自发的，而是你们关心的结果。这牲畜不理会它在追捕猎物中所遭受的疲劳和损耗，也不理会它凭自己的辛劳捕获猎物；而是抛开这一切，遵守主人的命令，表现出自己胜过食欲的渴望。\Quote{不错，因为它希望得到称赞，希望得到更大的食物。}那么你对自己说：狗因着对未来快乐的希望而轻视眼前的快乐；而你却不愿因对未来美善之事的希望而轻视眼前的；但狗确实知道，如果它在错误的时间、违背主人的意愿吃了那食物，它不但会被剥夺那食物，连分给它的也得不到，反而挨打而不是得食；而你们甚至不能明白这一点，狗凭习惯学到的东西，你们从理性也不能获得。让我们效仿狗吧。据说鹰和雕也是如此：狗对待野兔和鹿的方式，它们对待鸟类的方式也一样；它们也是从人那里学来的哲学。这些事实足以定我们的罪，足以证明我们有罪。再提一件事：—那些善于驯马的人，会牵来野性、凶猛、又踢又咬的马，短时间内便能训练它们，即使训师不在场，骑起来也是一种享受，它们的步伐如此井然有序；但灵魂的步伐可能全都乱套，却无人关心：它蹦跳、踢蹬，骑手像孩子一样被拖在地上，样子极其狼狈，却无人给灵魂戴上马衔、绑腿和嚼子，也不让精于骑术的骑手— 我是指基督—骑上她。因此一切都颠倒了。因为当你们既教导狗克制肚腹的欲望，又驯服狮子的怒火和马的不驯，还教导鸟说清楚话时，把理性的成就植入无理性的本性，同时把无理性生物的激情输入有理性的本性，这是多么矛盾？我们毫无借口，毫无。所有在克服激情上成功的人都会控告我们，无论是信徒还是非信徒：因为非信徒也曾如此成功；不仅如此，野兽和狗也如此，不仅仅人；我们也会控告自己，因为我们只要愿意就能成功，但当我们懒惰时，就被拖走。事实上，许多过着极其邪恶生活的人，常常在他们愿意的时候就改变了自己。但原因如我所说，我们四处寻求对别的事物有益的事，而不是对自己有益的事。如果你建造一栋华丽的房屋，你知道什么对房屋有益，而不是对自己有益；如果你穿上一件美丽的衣裳，你知道什么对身体有益，而不是对自己有益；如果你得到一匹好马，也是如此。没有人以灵魂如何美丽为目标；然而，当灵魂美丽时，那些东西就都不需要了；反之，如果灵魂不美丽，它们也无益。因为就像新娘一样，即使有挂着金绣帷幔的房间，有最美丽女子的歌队，有玫瑰和花环，有英俊的新郎，以及女仆和女性朋友，周围的人都英俊，但如果新娘本人充满缺陷，那么这一切都毫无益处；另一方面，如果新娘美丽，那么即使没有那些，也不会有什么损失，反而恰恰相反；因为对于丑陋的新娘，那些会使她显得更加丑陋，而对于美丽的新娘，则会使她显得更加美丽：同样，灵魂美丽时，不仅不需要那些附属物，反而它们会遮蔽她的美丽。因为我们看到哲学家闪耀光芒，不是在富有时，而是在贫穷时。因为在前者的情况下，许多人会将其归功于他的财富，认为他并非超越财富；但当他和贫穷为伴，却依然光芒四射，绝不容忍自己被强迫做任何卑劣之事，那时就无人与他分享哲学的冠冕。那么，如果我们想要富有，就让我们使灵魂美丽吧。如果你们的骡子确实洁白、肥壮、状态良好，但被它们拉着的你们却瘦弱、病态、丑陋，那有什么益处？如果你们的地毯确实柔软、美丽，满是精美的刺绣和艺术，但灵魂却衣衫褴褛，甚至赤裸污秽，那有什么收获？如果马步伐优美，与其说迈步不如说跳舞，而骑手连同他的歌队，装饰着比新娘更华丽的饰物，却比跛子更扭曲，手脚不如醉汉和疯子，那有什么益处？告诉我，如果有人给你一匹美丽的马，却扭曲你的身体，那有什么益处？现在你灵魂扭曲，却毫不在意？我恳求你们，让我们最终关心我们自己吧。不要让我们使自己比一切都更无价值。如果有人用言语侮辱我们，我们就恼怒和烦恼；但我们用行为侮辱自己，却毫不考虑。让我们，虽已迟延，终于醒悟，使我们能因着对灵魂多加关心、抓住德行，而获得永恒的美善之物，借着我们的主耶稣基督的恩宠和仁慈，愿光荣、权能、尊崇归于祂，与圣父及圣神，从今世直到永远，永无穷尽。阿们。\par

\SourceRecord{Translated by J. Walker, J. Sheppard and H. Browne, and revised by George B. Stevens.；Nicene and Post-Nicene Fathers, First Series；Vol. 11.；Edited by Philip Schaff.；Buffalo, NY: Christian Literature Publishing Co.,；1889.；<http://www.newadvent.org/fathers/210134.htm>.}

% structure: p0001 source=h1 target=section reason=page-title

\section{宗徒大事录讲道词 第三十五篇}

\BibleRef{宗 16:13-14}\par

\BibleQuote{安息日，我们出城到了河边，那里有一个祷告的地方（\EditorialNote{金口若望注：\InnerQuote{被认为是可能的}}），我们就坐下，对聚集在那里的妇女讲道。有一个名叫\Person{里狄雅}的妇人，是\Place{提雅提辣}城卖紫布的，她敬畏天主，听我们讲道；主开明了她的心，使她留意保禄所讲的话。}\par

再看看保禄犹太化。经上说：\Quote{那时以为}，从时间和地点，\Quote{祈祷将在城外，河边举行。}因为不应认为他们只在有会堂的地方祈祷；他们也在会堂外祈祷，但为此目的，他们指定一个地方，作为犹太人他们更属肉性——而且，\BibleQuote{在安息日}，此时很多人可能会聚集。\BibleQuote{我们坐下，向聚集在那里的妇女讲话。}再次注意免于一切骄傲。\BibleQuote{有一个妇女：}一个地位卑微的妇女，从其职业也可看出：但注意（在她身上）一个具有崇高思想（\OriginalTerm{有崇高思想者}{\GreekText{φιλόσοφον}}）的妇女。首先，天主召叫她的这一事实为她作证：\BibleQuote{当她受洗时，}经上说，\BibleQuote{她和她的一家}——注意他如何说服了他们全体——\BibleQuote{她恳求我们说：如果你们认为我对主是忠信的，请进到我家里住下。她强留了我们。}\BibleRef{宗 16:15}然后看她的智慧，她如何\OriginalTerm{恳求}{\GreekText{δυσωπεἵ}}，宗徒们，她的话多么充满谦卑，多么充满智慧。她说：\Quote{如果你们认为我忠信。}没有什么比这更有说服力。谁不会被这些话软化呢？她没有请求（或\Quote{要求}）没有单纯地恳求：而是留给他们决定，（却）极其强迫他们：\BibleQuote{她强留了我们，}经上说，借着那些话。又以另一种方式：看她又如何立刻结果实，并视为巨大收益。\Quote{如果你们认为我，}意思是，你们认为我显然是，由于你们交给我如此（神圣的）奥迹（即圣事，见第225页注3）：她之前不敢邀请他们。但保禄和同行者为什么不愿而需要被强留？要么是为了使她更热切渴望，要么是因为主说过：\BibleQuote{查问谁是配得的，就住在那里。}\BibleRef{路 10:8}（并非他们不愿），而是他们有意为之。——经上说：\BibleQuote{当我们去祈祷时，一个拥有占卜之神的使女迎着我们，她以占卜使主人们大得利益：她跟着保禄和我们喊叫说：\InnerQuote{这些人是至高者天主的仆人，向我们宣讲救恩之道。}}（第16-17节）可能什么原因使得恶魔说了这些话，而保禄禁止了它？二者都行事，一个恶毒，一个智慧：事实上恶魔想使自己变得可信。因为如果保禄接受了他的见证，被他接受就会欺骗许多信徒：因此他忍受说出对自己不利的话，为要建立对自己有利的话：所以恶魔自己也使用\OriginalTerm{迁就}{\GreekText{συγκαταβάσει}}以达到毁灭。起初，保禄不愿接受，而是鄙视它，不愿一下投入奇迹；但当它继续这样做，\OriginalTerm{并指出他们的工作}{\GreekText{καὶ} \GreekText{τὸ} \GreekText{ἔργον} \GreekText{ἐδείκνυ}}\BibleQuote{向我们宣讲救恩之道}时，他就命令它出来。因为经上说：\BibleQuote{保禄心烦，转身对那神灵说：\InnerQuote{我因耶稣基督的名命令你从她身上出来。}那时刻他就出来了。}（\Emph{a}）\BibleQuote{她的主人们看见他们得利的指望已经没了，就抓住保禄和息拉。}（第18-19节）（\Emph{d}）于是保禄做了所有，既行奇迹又教训，但危险中息拉也有分。为什么经上说\BibleQuote{但保禄心烦？}意思是，他看穿了恶魔的恶意，如他所说：\BibleQuote{因为我们不拘是他的意图。}\BibleRef{格后 2:11}（\Emph{b}）\BibleQuote{当她的主人们看见他们得利的指望已经没了。}到处金钱是邪恶的原因。哦，外邦人的残忍！他们希望这女孩继续被附魔，以便借她赚钱。经上说：\BibleQuote{他们抓住保禄和息拉，拖到市场上到官长面前，带到官吏那里说：\InnerQuote{这些人，是犹太人，极其扰乱我们的城！}}\BibleRef{宗 16:20}做了什么呢？那么为什么之前不拖他们来呢？\Quote{是犹太人：}这名字名声不好。\BibleQuote{传授我们罗马人不可接受也不可遵守的习俗。}\BibleRef{宗 16:21}他们\OriginalTerm{以叛国罪控告}{\GreekText{εἰς} \GreekText{καθοσίωσιν} \GreekText{ἤγαγον}}。（\Emph{e}）为什么他们不说，因为他们驱魔，就是犯了向天主的不敬？因为这对他们是一种失败：但相反，他们诉诸叛国罪（\OriginalTerm{以叛国罪}{\GreekText{ἐπὶ} \GreekText{καθοσίωσιν}}）：正如犹太人所说：\BibleQuote{我们除了凯撒没有王：凡自称为王的，就是反对凯撒。}\BibleRef{若 19:14}（\Emph{c}）\BibleQuote{群众一同起来攻击他们：官长撕掉他们的衣服，下令鞭打他们。}\BibleRef{宗 16:22}哦，不合理的行为！他们没有盘问，没有让他们说话。而且，既然发生了这样的奇迹，你们本应敬拜他们，把他们视为救主和恩人。因为如果你们想要钱，找到了如此大的财富，为什么不跑去得着它？这使你们更出名，拥有驱魔的能力比服从他们更强。看，甚至奇迹，但贪财更强大。（\Emph{f}）\BibleQuote{当他们重重鞭打他们后，将他们下在监里}——他们的愤怒很大——\BibleQuote{嘱咐狱警严紧看守}\BibleRef{宗 16:23}：\BibleQuote{他领了这样的命令，就把他们下在内监，两脚上了木狗。}\BibleRef{宗 16:24}注意，他又把他们下在\Quote{内}监：这也是出于上智的安排，因为将有伟大的奇迹。\par

(摘要。)\Quote{城外。}\BibleRef{宗 16:13}这地方便于听道，远离纷扰和危险。（乙）\Quote{在安息日。}由于没有工作，他们更留意听讲。（甲）\Quote{有一个女人，名叫里狄雅，她卖紫布。}\BibleRef{宗 16:14}请看，历史的作者并不羞于提及（归化者的）职业；（丙）况且斐理伯也不是一个大城。我们既得知这些事，也就不要以任何人为羞。伯多禄与一个硝皮匠同住。\BibleRef{宗 9:43}（保禄）与一个卖紫布的外邦妇女同住。骄傲在哪里？\Quote{主开导了［她的］心。}因此我们需要天主，开导人心；但天主开导那些愿意的心；因为可以看见有刚硬的心。\Quote{她就注意保禄所讲的话。}开导是天主的工作，注意是她的事；所以既有天主的作为，也有人作为。她领受了圣洗，\BibleRef{宗 16:15}并如此恳切地请求接纳宗徒们；甚至超过亚巴郎所用的。她不说别的凭据，只提那使她得救的凭据：\BibleRef{创 18:3}她不说：\Quote{如果你们看我为}一个伟大、虔诚的女人；而是说什么？\Quote{忠信于主}；若对主，何况对你们呢？\Quote{如果你们看我为}：如果你们不怀疑。她不说：\Quote{与我同住}，而是说：\Quote{请到我家里来住}：带着极大的恳切（她说这话）。真是一个忠信的女人！——\Quote{有一个丫头，附有占卜之神。}\BibleRef{宗 16:16}请问，这是什么魔鬼？就是他们所谓的\OriginalTerm{派松}{Python}神，因地方得名。你看见阿波罗也是个魔鬼吗？(魔鬼)想把他们引入诱惑；所以为激怒他们，\Quote{她跟在保禄和我们后面，喊说：这些人是至高者天主的仆人，给我们宣传救恩的道路。}\BibleRef{宗 16:17}你这该死的、可憎的！既然你知道\Quote{他们宣传}\Quote{祂的救恩的道路}，为什么你不自由地出来？但这魔鬼正做了西满所想的事，西满说：\Quote{请把这种能力也给我，使我无论给谁覆手，谁就能领受圣神。}\BibleRef{宗 8:19}因为魔鬼看见他们出名，就在这里也装假；它以为若宣讲同样的事，就能被允许留在身体里。但基督说：\Quote{我却不接受从人来的证据}\BibleRef{若 5:34}——指若翰——更何况从魔鬼呢？\Quote{赞美在罪人口中是不适宜的}\BibleRef{德 15:9}，何况从魔鬼呢？因为他们宣讲的不是出于人，而是出于圣神。因为他们行事并非出于夸耀。\Quote{保禄觉得厌烦}等等。他们以为用喊叫和喧嚷能恐吓（官长们），说：\Quote{这些人扰乱我们的城。}\BibleRef{宗 16:18}你说什么？你相信魔鬼吗？为什么这里又不信呢？魔鬼说：\Quote{他们是至高者天主的仆人}；你们说：\Quote{他们大大扰乱我们的城}；魔鬼说：\Quote{他们给我们宣传救恩的道路}；你们说：\Quote{他们教导我们不该接受的规矩。}\BibleRef{宗 16:21}请注意，他们甚至不听从魔鬼，只关注一样事：他们的贪财。但请看他们（保禄和息拉），如何不回答也不为自己辩护；（乙）\Quote{因为我什么时候软弱，什么时候就有能力了。我的恩宠为你足够了，因为我的能力在软弱中才完全。}\BibleRef{格后 12:9}这样，他们也因温良而受钦佩。（甲）\Quote{官长}等等，\Quote{命令狱警谨慎看守他们}\BibleRef{宗 16:22}：为使他们成为更大奇迹的媒介。（丙）看守越严，奇迹越大。官长们这样做，大概是想平息骚乱；因为他们看见群众催逼，想立刻抑制他们的激情，所以施加鞭打；同时他们也想要查问此事，所以把他们下在监里，并命令\Quote{谨慎看守}。据说，\Quote{［狱警］把他们两脚上了木狗}\BibleRef{宗 16:24}，即\OriginalTerm{那木具}{\GreekText{το} \GreekText{ξύλον}}，我们可以说是\OriginalTerm{刑枷}{\GreekText{νέρβον}}。\par

这些事岂不令人流泪！(想想)他们所受的苦，而我们（生活）在奢华之中，我们在戏院里，我们沉沦淹没（在放荡的生活中），总是寻求无益的消遣，不肯为基督受苦，甚至连言语上的苦也不受，连谈论也不肯。我恳求你们，让我们常常想起这些事：他们受了什么，忍受了什么，如何毫无畏惧，如何毫无怨言。他们做天主的工作，却受这些苦！他们并没有说：为什么我们宣讲这个，而天主却不帮助我们？但即使撇开真理，这件事本身对他们也是有益的；它使他们更刚强、更坚强、更无畏。\Quote{磨难生忍耐。}\BibleRef{罗 5:4}那么，我们不要寻求松散放荡的生活。因为就如在前一种情况中有双重好处：受苦者变得坚强，且赏报很大；同样在另一种情况中有双重恶果：这样的人变得更软弱，而且这软弱无益，只会带来邪恶。因为再没有比一个终生闲散奢华的人更无价值的了。如俗语所说，未经考验的人也是未经认可的；不仅在竞技场上，在一切事上也是如此。懈怠是无用的东西；在奢华本身中，没有比过奢华生活更符合目的的了：因为它使人因饱腻而厌烦，以至于既没有享用佳肴的乐趣，也没有休息的乐趣，一切都变得乏味，归于无有。\par

那么，让我们不要追求这个。因为如果我们思考一下，谁的生活更愉快，是劳苦工作的人，还是生活在奢侈中的人，我们会发现前者更愉快。因为首先，身体的感官既不清晰也不健全，而是迟钝（\OriginalTerm{迟钝}{\GreekText{χαῦναι}}）和萎靡的；当这些感官不正常时，即便健康也显然毫无乐趣可言。哪一匹马更有用，是养尊处优的还是受过锻炼的？哪一艘船更适用，是航行中的还是闲置不动的？哪一种水最好，是流动的还是停滞的？哪一种铁最好，是经常使用的还是不工作的？难道不是一块像银子一样闪闪发光，而另一块却生满铁锈、无用甚至失去一些自身的物质吗？同样的事情也发生在灵魂上，因为懒惰的后果：一种铁锈蔓延开来，腐蚀它的光采和一切。那么，如何擦去这铁锈呢？用磨难的磨刀石：这样就能使灵魂变得有用且适于一切。否则，请问，它如何能斩断情欲，当它的刃口已经卷曲（\OriginalTerm{弯曲}{\GreekText{ἀνακλώσης}}）并像铅一样弯曲时？它如何能伤害魔鬼？——那么，这样的人对他人而言，除了是可憎的景观外还能是什么——一个养得肥胖、像海豹一样拖着身体的人？我不是指那些天生如此体质的人，而是指那些因奢侈生活而把身体变成这样状况的人，那些天生瘦削的人。太阳升起，向四面八方射出明亮的光线，叫醒每个人去做工：农夫拿着铲子出去，铁匠拿着锤子，每个工匠拿着各自的工具，你会发现每个人都在摆弄自己合适的工具；女人要么拿着纺锤，要么拿着织布梭；而那个人，像猪一样，天刚亮就去填饱肚子，寻找如何享受盛宴的办法。然而，只有为屠宰而养的畜生才会从早吃到晚；甚至畜生中，那些负重能劳作的，也在夜间就出去干活。但这个人，从床上起来，当（中午的）太阳充满市场，人们已经厌倦了各自的工作时，这个人才起床，伸着懒腰，简直就像一头正在长膘的猪，在黑暗中浪费了白天最美妙的时光。然后他在床上坐很久，常常因为昨晚的酗酒甚至无法起身，又浪费了更多时间在这种（无精打采）中，然后开始打扮自己，出门，成为一种不雅的景象，毫无人性可言，完全是一头具有人形的野兽：他的眼睛因酒而流泪，* * * 而可怜的灵魂，就像瘸子一样，无法起身，拖着那大象般的大块肉体。然后他来到（各个）地方坐下，说和做那些事情，以至于他睡着比醒着更好。如果碰巧有坏消息传来，他表现得比任何女孩都软弱；如果是好消息，他表现得比任何孩子都愚蠢；脸上永远打着哈欠。他是所有想作恶者的目标，即便不是对所有人，至少对所有邪恶的情欲而言也是如此；愤怒容易激动这种人，还有情欲、嫉妒和所有其他情欲。所有的人都奉承他、讨好他，使他的灵魂比原来更软弱；而且每天他继续下去，加重他的疾病。如果他碰巧陷入任何事业的困难，他就成了尘土和灰烬，他的丝绸衣服对他毫无帮助。我们不是说这一切毫无目的，而是为了教导你们，你们中没有人应该游手好闲、漫无目的地生活。因为闲散和奢侈无益于工作、名誉或享受。因为谁会不谴责这样的人呢？家庭、朋友、亲戚（会说），他确实是大地的一个负担。这样的人来到世上毫无目的；或者不如说，不是毫无目的，而是对自己有害，走向自己的毁灭，并伤害他人。但说这更愉快——让我们看看这一点；因为这是问题所在。那么，有什么比一个无事可做的人（的状况）更不愉快呢？有什么比总是张着嘴、打着哈欠坐在市场上看着过往行人更悲惨、更可怜的呢？因为灵魂，由于其本性总是在运动中，无法忍受静止。天主使它成为一个行动的存在：工作是其本性；闲散违反其本性。因为让我们不要从那些有病的人来判断这些事，而是让我们用事实来检验事情本身。没有什么比闲暇和无事可做更有害的了：因此天主确实给我们施加了工作的必要性：因为闲散伤害一切。即使对于身体的肢体，不活动也是一种伤害。眼睛，如果不做它的工作，嘴巴，肚子，以及每一处可以提到的肢体，都会陷入最糟糕的疾病状态；但没有一个比灵魂更严重。但是，正如不活动是一种恶，那么活动在那些应当避免的事情上也是一种恶。因为就像牙齿一样，如果不吃东西，就会受到伤害；如果吃不对的东西，就会使它们不适并刺痛：这里也一样；如果灵魂不活动，或者在不正确的事情上活动，就会失去它应有的力量。那么，让我们避开这两者：不活动，以及比不活动更糟的活动。那可能是什么呢？贪婪、愤怒、嫉妒和其他情欲。关于这些，让我们以不活动为目标，以便我们能够获得应许给我们的美好事物，藉着我们的主耶稣基督的恩宠和怜悯，愿光荣、权能、尊崇归于祂，并与圣父及圣神，从今世直到永远。阿们。\par

\SourceRecord{Translated by J. Walker, J. Sheppard and H. Browne, and revised by George B. Stevens.；Nicene and Post-Nicene Fathers, First Series；Vol. 11.；Edited by Philip Schaff.；Buffalo, NY: Christian Literature Publishing Co.,；1889.；<http://www.newadvent.org/fathers/210135.htm>.}

% structure: p0001 source=h1 target=section reason=page-title

\section{论宗徒大事录讲道集 第三十六篇}

\BibleRef{宗 16:25-26}\par

\BibleQuote{午夜时分，保禄和息拉祈祷，并赞美天主；囚犯们都听他们。忽然，发生大地震，监牢的地基都震动了，立刻所有的门都开了，每个人的锁链都松开了。}\par

有什么能比得上这些灵魂呢？这些人受了鞭打，挨了许多鞭伤，遭了虐待，性命垂危，被投进内监，双脚上了木枷；尽管如此，他们并没有让自己入睡，而是整夜守夜。你们可曾注意到磨难是何等祝福？但我们，躺在柔软的床上，无所畏惧，整夜沉睡。但大概他们之所以守夜，正是因为他们处于这种境况。睡眠的专制不能压倒他们，疼痛的剧痛不能使他们屈服，对邪恶的恐惧不能使他们陷入无望的沮丧：不，正是这些事情使他们警醒：他们甚至充满了极大的喜乐。\BibleQuote{午夜时分，}经上说，\BibleQuote{囚犯们都听他们。}这真是奇异而惊人！\BibleQuote{忽然，发生大地震，监牢的地基都震动了，立刻所有的门都开了，每个人的锁链都松开了。狱吏醒来，看见监门大开，拔出剑来，想要自杀，以为囚犯已经逃走了。}\BibleRef{宗 16:27}地震是为了唤醒狱吏，门开了是为了让他对发生的事情感到惊奇：但这些事情囚犯们没有看见：否则他们都会逃走：但狱吏正要自杀，以为囚犯逃走了。\BibleQuote{保禄大声喊叫说：不要伤害自己，因为我们都在这里。}\BibleRef{宗 16:28}(\Emph{b})\BibleQuote{然后他叫了灯来，跳进去，战战兢兢，俯伏在保禄和息拉面前；又领他们出来，说：诸位先生，我应当做什么才能得救？}(29,30节)你们看到了吗，惊奇是如何压倒了他？(\Emph{a})他更惊奇于保禄的仁慈；他惊叹于保禄的勇敢胆量，当他有能力逃跑时却没有逃跑，还阻止他自杀。(\Emph{c})\BibleQuote{他们却说：相信主耶稣基督，你和你一家就必得救。他们就把主的话讲给他和他全家所有的人听。}(31,35节)于是立刻证明了他们对他的仁慈。\BibleQuote{就在那个时辰，他领他们去，洗了他们的鞭伤；他和所有属他的人立刻都受了洗。}\BibleRef{宗 16:33}他洗了他们，自己也受了洗，他和他的全家。\BibleQuote{他领他们进了自己的家，摆上饭食，和全家一起相信了天主，就充满了喜乐。天亮以后，官长派了差役来说：放那些人走吧。}(34,35节)很可能官长已经知道了发生的事，不敢亲自释放他们。\BibleQuote{狱吏把这话告诉保禄说：官长派人来释放你们，现在你们可以出去，平安地走吧。保禄却对他们说：他们公开地鞭打我们这些罗马人，又没有定案，就把我们下在监里，现在却私下赶我们出去吗？不行，让他们亲自来领我们出去吧！差役把这些话禀报官长，官长听说他们是罗马人，就害怕了。于是来劝他们，领他们出来，请他们离开那城。他们出了监，进了里狄雅的家，见了弟兄们，劝慰了他们，便离开了。}\BibleRef{宗 16:36}即使官长宣布释放，保禄也不出去，但为了里狄雅和其他人的缘故，他使他们恐惧：这样他们就不会被认为是凭自己的请求出来的，并且可以使其他人处于勇敢的姿态。指控是双重的：他们\Quote{是罗马人}，又\Quote{未被定罪}，就公开将他们下在监里。你们看到，在许多事情上，他们采取了人的措施。\par

（重述）\BibleQuote{在半夜}等等。\BibleRef{宗 16:25} 诸位亲爱的，让我们将那一夜与我们这些夜晚相比，这些夜晚充满了宴乐、醉酒和放荡，充满了如同死亡的睡眠，以及比睡眠更糟的守夜。因为有些人在毫无知觉中沉睡，而另一些人醒着却怀有可怜又可悲的意图，图谋欺诈，焦虑地想着钱财，研究如何报复那些亏负他们的人，默想仇恨，数算白天说过的辱骂之言：就这样，他们煽起忿怒的余烬，做出不可容忍的事。请注意伯多禄如何睡眠。\BibleRef{宗 12:6} 在那里，他睡着是明智的安排，因为天使来到他那里，必须不让任何人看见所发生的事；而在这里，保禄醒着也是良好安排，为使狱吏不能自杀。\BibleQuote{忽然发生了大地震。}\BibleRef{宗 16:26} 为什么没有行别的奇迹呢？因为，在所有奇迹中，这件事足以使他归化，因为他自身正处于危险中：使我们屈服的，并非奇迹本身，而是那些涉及我们自身获救的事。为使地震不显得是偶然发生的，有这件同时发生的事给它作证：\BibleQuote{门都开了，众人的锁链也解开了。}而且这事发生在夜间，因为宗徒们行事不是为了炫耀，而是为了人的得救。\BibleQuote{狱吏醒来，}等等。\BibleRef{宗 16:27} 这个狱吏并非恶人，他把他们\BibleQuote{下在内监里}，\BibleRef{宗 16:24} 是因为\BibleQuote{领受了这样的命令}，并非出于自己。那人心中一片慌乱。\BibleQuote{我应当做什么才能得救？}他问。为什么不在以前问呢？保禄大声喊叫，直到他看见，并且抢先对他说：\BibleQuote{我们都在这里。}经上说：\BibleQuote{他要了灯，就跳进去，俯伏在}囚犯\BibleQuote{的脚前，}这狱吏说：\BibleQuote{诸位先生，我应当做什么才能得救？}\BibleRef{宗 16:28} 为什么呢？他们说了什么？请注意，他发现自己安全后，并不以为万事大吉；他对这奇能感到敬畏。\par

你注意到过去的事和这里的事吗？那里有一个女孩从邪魔中被释放，他们就把宗徒们下在监里，因为她把他们从邪魔中释放了。这里，他们只展示了门大开，这却打开了他的心门，解开了两种锁链；那囚犯点燃了真光，因为他心中的光在照耀。\BibleQuote{他跳进来，俯伏在他们面前，}他不问：这是怎么回事？而是直接说：\BibleQuote{我应当做什么才能得救？}保禄回答什么？\BibleQuote{信主耶稣基督，你和你的全家必要得救。}\BibleRef{宗 16:31} 因这句话最赢得人心：全家也能得救。\BibleQuote{他们把主的话讲给他和他全家听。当夜在那时刻，他带他们去，洗了他们的鞭伤，}等等（32-33节），他洗了他们，自己也受了洗：他洗了他们的鞭伤，自己从罪中被洗净；他喂养他们，也被喂养。\BibleQuote{且欢喜，}经上说：虽然只有话语和美好的希望：他和他全家都信了天主。\BibleRef{宗 16:34} 这就是他信了的标记——他被解除了所有束缚。有什么比一个狱吏更坏、更残忍、更野蛮呢？他却以极大的尊荣款待他们。他欢乐不是因为自己安全，而是因为信了天主。（\Emph{a}）宗徒说\BibleQuote{信主}，因此作者在此说\BibleQuote{信了}。——（\Emph{d}）经上说\BibleQuote{所以现在，离开，平安地去罢}\BibleRef{宗 16:36}，就是说，安全地，不怕任何人。（\Emph{b}）\BibleQuote{保禄却对他们说}\BibleRef{宗 16:37}：他不愿显得像一个被定罪或做错事的人而获得自由，因此他说\BibleQuote{他们公开地未经审判就打了我们}等等——免得这成为他们的恩惠。（\Emph{e}）此外，他们希望狱吏自己脱离危险，免得他以后为此被追究。他们不说\BibleQuote{打了我们}这些行奇迹的人，因为官长们连这些也不顾；但最能震骇他们的话是\BibleQuote{未经审判，又是罗马人}。（\Emph{c}）请看恩宠如何多样地管理事情：伯多禄怎样出来，保禄怎样出来，虽然两人都是宗徒。经上说\BibleQuote{他们害怕}\BibleRef{宗 16:38}，因为那些人是罗马人，并非因他们不义地把他们下在监里，\BibleQuote{就求他们离开那城}\BibleRef{宗 16:39}：恳求他们作为恩惠。他们去了里狄雅的家，坚固了她，然后离去。因为不该让他们的女主人处于忧愁焦虑中。但他们出去，并非顺从官长的请求，而是急忙去传道：那城已因这奇迹得到充分益处，所以他们不应再留在那里。因为行奇迹者不在时，奇迹显得更大，它本身更响亮地宣告着：狱吏的信德本身就是一种声音。有什么事比这更伟大？他被戴上锁链，却释放别人；自己受束缚，却解开双重束缚；那束缚他的人，他因被束缚而释放。这些确实是（超性）恩宠的工作。\par

(\Emph{f}) 我们应常记住这位狱卒，而非那奇迹：一位囚徒（宗徒）如何说服了他的狱卒。异教徒说什么？他们说：\Quote{说服了什么样的人呢——一个鄙陋、可怜、无知的家伙，满脑子邪恶，别无他物，极易被引向任何事物？因为这些，他们说，就是那些信了的人：一个制革匠、一个卖紫布的、一个太监、奴隶和妇女。} 这就是他们所说的。那么，当我们举出有地位的人，百夫长、总督、从那时至今的统治者、皇帝们时，他们还能说什么？但对我来说，我要说另一件更伟大的事：让我们看看这些微不足道的人本身。你说：\Quote{奇在哪里？} 我说，这正是奇事。因为，若有人被说服接受普通事物，那不足为奇；但若论及复活、天国、哲学式的自制生活，且对卑微之人讲论这些而说服了他们，那就比说服智者更奇妙。因为当所说之事不带任何危险时，反对者或许能合理地推诿于对方的无知；但当宣讲者——如你所愿对奴隶说：\Quote{你若被我说服，便身陷险境，你将与所有人为敌，你必须死，你必须遭受无数苦难}——却仍能说服那人灵魂时，这里便不能再谈无知了。的确，若教义包含悦人之事，人们可合理地说此；但若连哲学家都不愿学习的事，奴隶却学会了，那奇事更大。如果你愿意，让我们把那位制革匠叫到面前，看看伯多禄与他谈论了些什么；或者，如果你愿意，就看看这位狱卒。保禄对他说了什么？你说：\Quote{基督复活了，死人要复活，有国度；而他毫不费力地说服了他，一个易被引向任何事物的人。} 如何呢？他没有谈及生活方式吗？他必须节制，必须超越金钱，不可冷酷无情，必须与人分享自己的财物？不能说，在这些事上被说服也是由于心智的缺乏；不，被引向这一切需要伟大的灵魂。姑且承认，就教义而言，他们因无知而更容易接受；但接受如此有德、克己的生活规范，怎能归因于理解力的缺陷？因此，一个人越不理解，却仍被说服接受哲学家们无法说服同行者的事，那奇事就越大——当妇女和奴隶被说服接受这些真理，并以行动证明时，柏拉图们和其他所有人都未能说服任何人。为何我说\Quote{任何人？} 不如说：连他们自己都不能：相反，柏拉图说服（其门徒）金钱不可轻视，因为他自己获得了如此丰厚的财产、金戒指和酒杯；而大众的荣誉不可轻视，苏格拉底本人也证明了这一点，尽管他对此可能讲得天花乱坠：因为他所做的一切，都关乎名声。若你熟悉他的言论，我可以在此长篇大论，并指出其中有多少\OriginalTerm{讽刺}{\GreekText{εἰρωνείαν}}——如果我们至少相信他的门徒对他的描述——以及他所有的著作都以虚荣为基础。但撇开他们，让我们把话题转向我们自己。因为除了已说的，还有这一点：人们是在自己的危险中被说服的。因此，你不要无耻，让我们想想那一夜、那足枷、那赞美的诗歌。让我们也这样做，我们就会为自己打开——不是监狱，而是——天堂。如果我们祈祷，我们甚至能打开天堂。厄里亚曾借祈祷关闭并打开了天堂。\BibleRef{雅 5:17} 天堂也有监狱。祂说：\Quote{凡你在地上所捆绑的，在天上也要被捆绑。} \BibleRef{玛 16:19} 让我们在夜间祈祷，我们就能解开这些束缚。因为祈祷能解开罪恶，让那位寡妇说服我们，让那位朋友说服我们——他在夜深人静时坚持敲门；\BibleRef{路 11:5} 让科尔乃略说服我们，因为经上说：\Quote{你的祈祷和你的施舍已上达天主面前。} \BibleRef{宗 10:4} 让保禄说服我们，他说：\Quote{那真正孤寡的寡妇，寄望于天主，昼夜不断恳求祈祷。} \BibleRef{弟前 5:5} 若他对寡妇、软弱的妇女这样说，更当对男人说。我之前已对你们讲过此事，现在再重复：让我们在夜间警醒自己；即使你不多做祈祷，只要警醒地做一个，就足够了，我不求更多；倘若不能在午夜，至少在黎明初现。表明夜晚不仅属于肉体，也属于灵魂；不要让它虚度，而要以此回报你的主；不，这益处本身会回到你身上。试问，若我们陷入任何困境，我们不去求谁？若我们很快得偿所愿，便又舒畅。你能有一个朋友去向他请求，他乐意接受并因你的请求而感激你，这是何等恩惠？何等恩惠——不必四处找人请求，而能找到一位甘愿者？无需通过他人代为求情？还有什么比这更伟大？因为这里有一位，当我们不向他人祈求而只向祂祈求时，祂便做得最多：正如一位真诚的朋友，当我们请别人代他向祂求情时，祂最会抱怨我们不信任祂的友谊。让我们也如此行。你会问：\Quote{但我若得罪了祂，怎么办？} 停止得罪，哭泣，然后靠近祂，你就能迅速为过去的罪赢得祂的宽恕。只要说：我得罪了；从灵魂深处真诚地说，一切都被赦免。你并不如祂那样渴望你的罪被赦免。证明你不那么渴望的证据是，你无意守夜，也无意慷慨施舍金钱；但祂，为赦免我们的罪，连祂的独生真子、共享祂宝座的那位也不吝惜。你看见祂是如何更渴望赦免你的罪（比你更渴望被赦免）吗？那么，我们不要懒惰，不要再拖延。祂是仁慈良善的；只让我们给祂一个机会。\par

（甚至）这一点（祂也寻求），只是为使我们不至于成为无用的，因为即便没有这个，祂本可以释放我们脱离它们；但就像我们（抱着同样的目的）为我们的仆人筹谋安排许多事情一样，祂在我们的救恩上也是如此。\Quote{让我们以感恩来迎接祂的慈颜。} \BibleRef{咏 94:2}，因为祂是良善仁慈的。但如果你不呼求祂，\Emph{祂}会做什么呢？你不愿意说：\Quote{宽恕吧}；你不会从心里说，只从口里说。什么是真正地呼求？是怀着内心的决心，带着恳切，以真诚的心意；就像人们对香料所说的：\Quote{这是纯正的，毫无虚假。}，这里也一样。真正呼求祂的人，真正向祂祈祷的人，不断专注于此事，直到获得所求才罢休；但那只\OriginalTerm{形式化地履行义务者}{\GreekText{ἀφοσιούμενος}}，并且连这也只是出于履行规矩，并非真正呼求。无论你是谁，不要只说：\Quote{我是罪人。}，还要努力摆脱这种状态；不要只说这个，还要悲伤。如果你悲伤，你就是认真的；如果你不认真，你就不悲伤；如果你不悲伤，你就是在敷衍。一个只说\Quote{我有病}，却不尽力摆脱疾病的人，算什么人？祈祷是一把强大的武器。\Quote{如果你们，}主说，\Quote{知道把好东西给你们的儿女，何况你们的父呢？} \BibleRef{路 11:13}那么你为什么不愿意到祂面前？祂爱你，祂比一切更有能力。祂既愿意，又能够，有什么阻碍呢？没有。但是，就我们而言，让我们怀着信德靠近，靠近，献上祂所渴望的礼物：忘怀冤屈、仁爱、温良。即使你是罪人，你也可以大胆地向祂祈求罪的赦免，如果你能证明你自己已经做到这一点；但即使你是义人，却不具备这种忘怀伤害的德行，你也无益。一个宽恕了邻舍的人，不可能得不到完全的宽恕；因为天主比我们更慈悲。你说什么？如果你说：\Quote{我受了冤屈，我抑制了我的愤怒，我因你的诫命承受了怒火的冲击，难道你不宽恕吗？}祂一定会宽恕：这是众人皆知的。因此，让我们清除灵魂中一切怨恨。这足以使我们蒙垂听；让我们警醒、恒切地祈祷，好能享受祂丰厚的仁慈，堪当承受所许诺的恩赐，因我们的主耶稣基督的恩宠和仁慈。愿光荣、权能、尊威归于祂，偕同圣父及圣神，自今而始，至于永远，世世无穷。阿们。\par

\SourceRecord{Translated by J. Walker, J. Sheppard and H. Browne, and revised by George B. Stevens.；Nicene and Post-Nicene Fathers, First Series；Vol. 11.；Edited by Philip Schaff.；Buffalo, NY: Christian Literature Publishing Co.,；1889.；<http://www.newadvent.org/fathers/210136.htm>.}

% structure: p0001 source=h1 target=section reason=page-title

\section{《宗徒大事錄》講道集第三十七篇}

《宗徒大事錄》第十七章 1、2、3\par

\BibleQuote{他們經過安非波里和阿波羅尼亞，來到得撒洛尼，在那裏有猶太人的會堂。保祿照常例進去，一連三個安息日，根據聖經與他們辯論，闡明並指出基督必須受苦，並從死者中復活，又說：\InnerQuote{我所傳給你們的這位耶穌，就是基督。}}\par

他們又急忙越過小城，趕往較大的城市，因為從那裏，聖言將如泉水般流入鄰近的的城市。\BibleQuote{保祿照常例，進了猶太人的會堂。}雖然他曾說：\BibleQuote{我們轉向外邦人}\BibleRef{宗 13:46}，他並沒有放棄這些人：他對他們懷有如此深切的渴望。因為請聽他說：\BibleQuote{弟兄們，我心裏所渴望的，向天主所求的，是為以色列人得救}\BibleRef{羅 10:1}以及：\BibleQuote{為我的弟兄，我本人寧願受詛咒，與基督隔絕。}\BibleRef{羅 9:3}他這樣做，是因為天主的許諾和光榮；也是為了不給外邦人留下絆腳石。經上說：\BibleQuote{他根據聖經，一連三個安息日與他們辯論，開示並提出基督必須受苦。}你要注意，他在所有事之前首先宣講苦難：他們如此不以苦難為恥，因為知道這是得救的原因。\BibleQuote{他們中有些人信了，便與保祿和息拉結交；虔敬的希臘人中有極大一群，尊貴的婦女也不少。}\BibleRef{宗 17:4}作者只提到講論的要點，他不願重複，也並非每次都報告講道詞。但不信的猶太人（最好的版本省略\Quote{不信的}）滿心嫉妒，糾集了一些市井無賴，聚眾騷亂，擾亂全城，衝進雅松的家，要將他們帶到民眾那裏。找不到他們，就把雅松和一些弟兄拉到地方官那裏，喊叫說：\BibleQuote{這些攪亂天下的人也到這裏來了，雅松接待了他們。這些人違背凱撒的法令，說有另一位王，就是耶穌。}\BibleRef{宗 17:5}哦！這是怎樣的控告！他們又給他們加上叛國的罪名，\BibleQuote{說有另一位王，耶穌。}（第8-9節）雅松堪稱可欽佩的人，他置自身於危險，卻使他們脫離危險。\BibleQuote{弟兄們當夜立即送保祿和息拉往貝洛雅去；他們到了那裏，就進了猶太人的會堂。}經上說：\BibleQuote{這些人比得撒洛尼人更開明，即\OriginalTerm{更溫和}{\GreekText{ἐπιεικέστεροι}}（在行為上）：因為他們以全備的熱誠接受聖言，且不是輕率地，而是以不帶激情的嚴格態度，天天考查聖經，看這些事是否屬實。}（第10-11節）\BibleQuote{因此他們中有許多人信了，也有不少希臘尊貴的婦女和男子。但得撒洛尼的猶太人得知保祿在貝洛雅傳講天主的聖言，也到那裏去煽動群眾。於是弟兄們立即送保祿往海邊去；息拉和弟茂德仍留在那裏。}\BibleRef{宗 17:12}請看，他有時退讓，有時又前進，在許多事情上按人性的考慮採取措施。\BibleQuote{護送保祿的人帶他到了雅典；他們領了命令，叫息拉和弟茂德趕快到他那裏去，就離開了。}\BibleRef{宗 17:15}現在，讓我們再回顧所說的話。\par

(摘要)\BibleQuote{三个安息日}，经上说，正是他们空闲的时候，\BibleQuote{他同他们辩论，讲解圣经}\BibleRef{宗 17:2}：因为基督也常如此行：正如我们在许多场合看到祂根据圣经辩论，而非总用奇迹（劝人）。因为对此，他们确实怀着敌意，称他们为骗子与江湖术士；但藉圣经的推理说服人者，则不会受此指责。此外，我们常看到（保禄）仅仅凭教训的力量就说服了人：在安提约基雅，\BibleQuote{全城的人都聚集了}\BibleRef{宗 13:44}，此事实在伟大，因为本身非小事，实乃一个大奇迹。并且为使他们不致认为这一切全凭自己力量，而是天主所准许，便有两件事发生，即：\BibleQuote{有些人信服了}等等。(\Emph{c})\BibleQuote{还有许多虔诚的希腊人及不少显要妇女}；但另一些人相反：\BibleQuote{犹太人满心嫉妒}等等。(4、5节)(\Emph{b})并且，从被召本身是出于天主预定这件事上，(\Emph{a})他们既不图自身光荣以为胜利属己，也不因负全责而恐惧。但为什么他说：\BibleQuote{我们应往外邦人那里去，他们往受割礼的人那里去}\BibleRef{迦 2:9}，却又向犹太人讲论？\Emph{(\OriginalTerm{\GreekText{α}}{\GreekText{α}})}祂这样做是额外的。\Emph{(\OriginalTerm{\GreekText{β}}{\GreekText{β}})}因为祂也做了其他超出本分的事。例如，基督规定他们应\BibleQuote{靠福音生活}\BibleRef{格前 9:14}，但我们的宗徒却没有这样做；基督未派遣他付洗，他却付了洗。看他是如何与众人平等的。伯多禄往受割礼的人那里，他往外邦人那里，且占大部分。\Emph{(\OriginalTerm{\GreekText{α}}{\GreekText{α}})}既然他必须向犹太人讲论，如何又说：\BibleQuote{那曾在伯多禄内运行以向受割礼者行事的，也同样在我内运行向外邦人行事}\BibleRef{迦 2:8}？正如那些宗徒也与外邦人来往，虽被派往受割礼者，我们的宗徒也是如此。他的工作大部分固然在外邦人，却也没有忽略犹太人，免得他们以为彼此隔绝。你或许要问，他为何首先进入会堂，仿佛这是他主要目的？不错，但他借犹太人说服外邦人，并借他向犹太人所讲论的事。他知道这对异邦人最为合适，最有利于信服。因此他说：我既然是\BibleQuote{外邦人的宗徒}\BibleRef{罗 11:13}。他的书信也都反对犹太人。——\BibleQuote{基督必须受难}，他说，\BibleRef{宗 17:3}。若祂必须受苦，无疑也必须复活：因为前者远比后者奇妙。若祂将无罪的人交于死亡，就更应使祂复活。但不信的犹太人招了一些无赖之徒，耸动全城\BibleRef{宗 17:5}，以致外邦人人数更多。犹太人自觉不足以掀起骚乱；因无合理借口，便常用骚乱并招揽无赖实现目的。\BibleQuote{找不着他们，}经上说，\BibleQuote{就拖了雅松和几个弟兄。}\BibleRef{宗 17:6} 何等的暴虐！无故将他们从家中拖出。\BibleQuote{这些人，}他们说，\BibleQuote{都违背凯撒的谕令}\BibleRef{宗 17:7}：因为他们所说的话没有违背谕令，也没有在市内滋事，他们便将另一罪状加给他们：\BibleQuote{说另有一个王，叫作耶稣。他们惊扰了百姓}等等。\BibleRef{宗 17:8} 你们既见祂死了，还怕什么？\Emph{(\OriginalTerm{\GreekText{β}}{\GreekText{β}})}\BibleQuote{他们取了保状}等等。\BibleRef{宗 17:9} 看，雅松如何交出保状送走保禄，甘愿冒生命危险。\Emph{(\OriginalTerm{\GreekText{α}}{\GreekText{α}})}\BibleQuote{弟兄们}等等。\BibleRef{宗 17:10} 看，每次迫害都扩展了宣讲。\BibleQuote{这些人，}经上说，\BibleQuote{比得撒洛尼的人更为高贵}\BibleRef{宗 17:11}：即他们不做卑鄙之事，有些人信了，其他人（未信的）也没有那样做。\Emph{(\OriginalTerm{\GreekText{β}}{\GreekText{β}})}\BibleQuote{天天，}经上说，\BibleQuote{考查圣经，看这些事是否如此：}并非出于一时的冲动或热诚。\BibleQuote{更为高贵，}经上说：即德性方面。\Emph{(\OriginalTerm{\GreekText{α}}{\GreekText{α}})}\BibleQuote{所以他们中间有许多人}等等。\BibleRef{宗 17:12} 这里又有希腊人。\Emph{(\OriginalTerm{\GreekText{β}}{\GreekText{β}})}\BibleQuote{但得撒洛尼的犹太人}等等。\BibleRef{宗 17:13}，因为那里有无赖之徒。然而那城市更大。但在大城市百姓更坏并不奇怪：当然，大城市常有恶人，骚乱的缘由也多。犹如身体中，病势因物质与燃料越烈，此处亦然。\Emph{(\OriginalTerm{\GreekText{α}}{\GreekText{α}})}但请看我求你，他们的逃亡是天主上智的安排，并非由于怯懦：否则他们早该停止宣讲，也不致激怒那些人。但从这次（逃亡）产生了两件事：那些（犹太人的）怒气被熄灭，宣讲也得以扩展。但以适合他们骚乱行为的话，他说：\BibleQuote{煽动群众。}\Emph{(\OriginalTerm{\GreekText{β}}{\GreekText{β}})}正像在依科尼雍所做的那样——使他们除了自毁之外，还要加上毁灭别人的罪。\EditorialNote{第十四章2、19节} 这就是保禄论述他们的话：\BibleQuote{阻止我们向外邦人讲道，为常充满他们的罪，因为义怒已至极限临到他们身上。}\BibleRef{得前 2:16} 为何他不留下？因为若\EditorialNote{在吕斯特辣，第十四章19、21节}那里，他被石头打死后仍停留许久，此处更应如此。为何？（主）不愿他们常行神迹；因为受迫害而不用神迹得胜，其本身就是神迹，并不亚于行神迹。因此，正如现今祂不用神迹得胜，那时祂也常愿如此得胜。所以宗徒们并不追求神迹：正如他亲自说：\BibleQuote{我们宣讲被钉的基督}\BibleRef{格前 1:23}——向那些渴求神迹、渴求智慧的人，我们给予那连神迹后也不能说服的，我们却说服了！所以这是一个大奇迹。看，当宣讲扩展时，他们并不急于追求神迹。因为从那时起，信众本身应成为其余人的大奇迹。然而，他们退却与前进，行了这些事。\Emph{(\OriginalTerm{\GreekText{α}}{\GreekText{α}})}\BibleQuote{弟兄们立刻，}经上说，\BibleQuote{送走了保禄。}\BibleRef{宗 17:14} 此时他们单独送保禄：因为他们为他担心，怕他受害，毕竟他是首脑和中心，非他人可比。\Emph{(\OriginalTerm{\GreekText{β}}{\GreekText{β}})}\BibleQuote{他们送走了他，}经上说，\BibleQuote{犹如往海边去：}使敌人不易捉拿他。因为目前他们自己无能为力；与他一起则有所成就。目前，经上说，他们愿救他。\Emph{(\OriginalTerm{\GreekText{α}}{\GreekText{α}})}绝非（超性）恩宠在所有场合都一律工作；相反，它让他们按人的判断采取措施，（只）激动他们，唤醒他们睡眠，使他们劳苦。所以，你看，它只安全地带领他们到斐理伯，此后不再。\BibleQuote{他们接受了命令，}经上说，\BibleQuote{叫息拉和弟茂德赶快到他那里去，就出发了。}\BibleRef{宗 17:15} 虽他是保禄，仍需要他们。天主催促他们去马其顿是有充分理由的，因为那里还有希腊光明在前。\BibleRef{宗 16:9}\par

请看其余门徒对首领所表现的热忱：不像我们现在这样，彼此分离，分作大小：我们有些人傲慢，另一些人嫉妒；凡嫉妒的原因，是因为我们骄傲，不肯忍受与他们同等。身体和谐的原因，是因为没有骄傲；而没有骄傲，是因为肢体必须互相需要，头也需要脚。天主这样安排了我们，尽管如此，我们仍不肯接受：即使没有这些，我们之间本该有爱。你们没有听见那些外人如何指责我们说：\Quote{需要才交朋友吗？}平信徒需要我们，我们又是为他们而存在。因为没有受教的人，教师或统治者就不会存在，也无法执行其职务——那是不可能的。正如土地需要农夫，农夫也需要土地，这里也是如此。教师若没有一个他所教导的学生，他得什么赏报呢？受教者若没有受到最好的教导，又得什么赏报呢？因此，我们彼此互相需要：治理者与被治理者，首领与服从者，因为统治者是为众人而存在的。因为没有人单靠自己就能做什么，无论是需要\OriginalTerm{任命}{\GreekText{χειροτονἥσαι}}，或是审察人的谋略和意见，但借着集会与多数，他们变得更有尊荣。例如，穷人需要施予者，施予者又需要接受者。他说：\BibleQuote{彼此相顾，激发爱德，勉励行善。}\BibleRef{希 10:24}为此，全教会的集会更有能力：每个人单独不能做的，当他与其他人结合时，就能做。因此，在这里为世界、为教会、从地极到地极、为和平、为患难中的人所献上的祈祷，是最必需的。\Person{保禄}也显示这一点，他说：\BibleQuote{藉着许多人的代祷，为我们所得的恩赐，使许多人因我们而感谢。}\BibleRef{格后 1:11}；意思是，祂将恩惠施予许多人。他也常常请求他们的祈祷。请看天主论及尼尼微人所说的话：\BibleQuote{何况这座城中有十二万多居民，我岂能不爱惜呢？}\BibleRef{纳 4:11}因为如果祂说：\BibleQuote{那里有两个或三个人因我的名聚在一起}\BibleRef{玛 18:20}，他们大有能力，何况是许多人呢？然而，即使你一个人，也能有所成就，只是不能同样地。为何你只是一个人？为何你不使许多人？为何你不成为爱的缔造者？为何你不\OriginalTerm{创造}{\GreekText{κατασκευάζεις}}友谊？你缺少了美德的首要卓越。因为正如人们共同作恶更激怒天主，同样，人们同心行善更令祂喜悦。祂说：\BibleQuote{不可随从多数去作恶。}\BibleRef{出 23:2}\BibleQuote{他们都离弃了正道，一同变为无用}\BibleRef{罗 3:12}，他们好像是在邪恶中合唱的人。你当为自己结交朋友，胜过家仆和一切。如果缔造和平的人是天主的子女，那么缔造友谊的人更当如何？\BibleRef{玛 5:9}如果仅仅劝人和好的人被称为天主的子女，那么使已和好的人成为朋友的人，将不配得什么？让我们投身于这项事业：使彼此为敌的人成为朋友，使那些并非敌对但也不是朋友的人团结起来，而首先，使我们自己团结起来。因为正如一个人在自己的家里与妻子有仇恨和分歧，他在劝和别人时便没有权威，反而会被说：\BibleQuote{医生，医治你自己吧}\BibleRef{路 4:23}，同样，人在这种情况也会被这样说。那么，我们内在的仇恨是什么？是灵魂对身体的仇恨，是恶习对美德的仇恨。让我们终止这仇恨，除去这战争，然后，在平安中，我们也要以极大的放胆向他人说话，我们的良心不控告我们。愤怒与温和交战，吝啬与轻视交战，嫉妒与心地善良交战。让我们终结这战争，推翻这些敌人，树立这些战利品，在我们自己的城内建立和平。我们内在有一座城和一个政治体，有许多公民和异乡人：但让我们驱逐异乡人，免得我们的子民被毁灭。不要让任何异域或虚假的教义进入，也不要有肉欲。我们岂不见，若有敌人被发现在城里，就被当作间谍审判？因此，让我们不仅驱逐异乡人，也要驱逐敌人。如果我们看见一个，就把他交给统治者（即良心\OriginalTerm{良心}{\GreekText{τᾥ} \GreekText{νᾥ}}），那是确实的异乡人、蛮族人，尽管披着公民的服饰。因为我们内在有许多这样的意念，它们本质上是敌人，却披着羊皮。就如波斯人，当他们脱下头巾、束腰外衣和蛮族鞋子，穿上我们寻常的服装，并剃净头发，用我们的语言交谈时，他们在外表下隐藏着战争：但一旦施加酷刑（\OriginalTerm{酷刑}{\GreekText{βασάνους}}或\Quote{考验}），你就揭露出隐藏之物。同样，在这里，通过反复用酷刑来审查（或\Quote{考验}）这样的意念，你很快会发现它的精神是陌生人的。但为了也以实例向你们显示魔鬼派来窥探我们内里之事的间谍种类，来，让我们剥去其中一个，在法庭上严格审问它：如果你们愿意，让我们提出一些保禄所揭露的。他说：\BibleQuote{这些事在自创的敬礼、谦卑和苦待身体上，确有智慧的外表，但在满足肉欲上，毫无价值。}\BibleRef{哥 2:23}魔鬼想引入犹太教：如果它以本来面目出现，就不会达到目的。因此，注意他是如何做的。他说：\Quote{你必须苦待身体，这是（真正的）哲学：不接纳食物，而是防范它们：这就是谦卑。}现今在我们的时代，关于异端者，他又想使我们贬低于受造物。请看他是如何装饰他的欺骗的。如果他说\Quote{敬拜受造物}，他就会被识破；但他说什么？他说：\Quote{天主（即圣子和圣神）是受造物。}但让我们为审判官的决定而揭露宗徒著作的含义：那里让我们带他来；他们自己会承认这宣讲和语言。许多人牟取不正当的利益，\Quote{为能有所施舍给穷人}，这不正当的利益：这同样是邪恶的意念。但让我们剥去它的外衣，定它的罪，以免被它欺骗；而是借着逃脱魔鬼的一切诡计，并严格持守纯正的教义，使我们能够平安地度过今生，并获得所应许的美物。这全是借着我们的主\Person{耶稣基督}的恩宠和慈爱；愿光荣、权能、尊崇归于祂偕同圣父及圣神，从今世直到永远。阿们。\par

\SourceRecord{Translated by J. Walker, J. Sheppard and H. Browne, and revised by George B. Stevens.；Nicene and Post-Nicene Fathers, First Series；Vol. 11.；Edited by Philip Schaff.；Buffalo, NY: Christian Literature Publishing Co.,；1889.；<http://www.newadvent.org/fathers/210137.htm>.}

% structure: p0001 source=h1 target=section reason=page-title

\section{论 \WorkTitle{宗徒大事录} 第三十八篇讲道}

\BibleRef{宗 17:16-17}\par

\BibleQuote{当\Person{保禄}在\Place{雅典}等候他们时，看见满城都是偶像，他的心神就在内被激动。因此他在会堂里同犹太人和虔诚的人辩论，也每天在市场上同所遇见的人辩论。}\par

请注意，他在犹太人当中比在外邦人当中遭遇了更大的考验。在雅典，他并未经历这类事；事情只到被嘲笑为止：然而他也劝化了一些人；而在犹太人中间，他却历经许多危险；他们对他的敌意如此之大。——\Quote{他的灵}，经上说，\Quote{在他里面激动，因为他看见满城都是偶像。} 在其他任何地方都看不到如此多的崇拜对象。但又一次，\Quote{他在会堂里与犹太人辩论，每日在市场上与遇到的人辩论。然后有几位斯多噶派和伊壁鸠鲁派的哲学家与他相遇。} \BibleRef{宗 17:18} 奇怪的是，哲学家们没有嘲笑他，尽管他那样说话。\Quote{有的说：这胡言乱语的要说什么？} 傲慢地，立刻——这离哲学甚远。\Quote{另有的说：他似乎是个宣传外邦鬼神的人，} 从讲道来看，因为他没有傲慢。他们不明白，也不理解他所讲的主题——他们怎能理解呢？他们中有人断言天主是形体，有人断言快乐是（真正的）幸福。\Quote{宣传外邦鬼神，因为他向他们宣讲耶稣和复活：} 事实上他们竟以为\Quote{\OriginalTerm{阿纳斯塔西斯}{Anastasis}（即复活）}是一位神祇，因为他们也习惯了崇拜女神。\Quote{于是他们带了他，把他带到亚略巴古} \BibleRef{宗 17:19} — 不是为了惩罚，而是为了学习 — \Quote{到亚略巴古}，那里是审判谋杀案件的地方。因此，请注意，他们怀着学习的希望（问他），说：\Quote{我们可以知道你所讲的这新道理是什么吗？因为你把一些奇怪的事传到我们耳中} \BibleRef{宗 17:20}；到处都以新奇为罪名：\Quote{因此我们愿意知道这些事是什么意思。} 那是一座爱说话的城市，他们的城市。\Quote{原来所有的雅典人和侨居在那里的外方人，除了谈论或听闻一些新事物外，没有别的事可做。于是保禄站在亚略巴古当中，说：雅典人啊，我看你们在各方面}（21–22节）——他用的是赞美之词：这词似乎没有什么冒犯之意——\OriginalTerm{更为虔诚}{\GreekText{δεισιδαιμονεστέρους}}，即\OriginalTerm{虔敬}{\GreekText{εὐλαβεστέρους}}，\Quote{更富有宗教热忱。因为我经过各处，观看你们所敬拜的，发现一座祭台，上面刻着：\InnerQuote{献给一位未知之神。}那么，你们所不认识而敬拜的，我现在向你们宣告。} \BibleRef{宗 17:23} \Quote{上面刻着：献给一位未知之神。} 也就是说，雅典人，因为多次他们也从外邦接受神祇——例如，密涅瓦神庙、潘神以及来自不同地方的其他神祇——他们担心可能还有其他一些他们未知的神祇，但在别处受到敬拜，为了更加稳妥，果然也给那位神立了一座祭台：因为那神是未知的，就刻着\Quote{献给一位未知之神}。那么，他告诉他们，这位神就是基督；或者更确切地说，是万有的天主。\Quote{我向你们宣告祂。} 请注意他如何表明他们早已接受了祂，并说：\Quote{这不是奇怪的事，} \Quote{我向你们介绍的并非新奇之物。} 一直以来，他们所说的正是：\Quote{你所讲的这新道理是什么？因为你把奇怪的事传到我们耳中。} 因此他立刻消除他们的疑虑，然后说：\Quote{创造世界及其中的万有的天主，祂是天地的主}——为使他们不致以为祂是众神之一，他立即纠正他们的这一点；并加添：\Quote{并不住人手所造的殿宇} \BibleRef{宗 17:24}，\Quote{也不被人手服侍，好像祂需要什么}——你是否注意到，他怎样一点一点地引入哲学？怎样嘲笑异教的谬误？\Quote{因为祂将生命、气息和万有赐予众人；并且从一人造了所有民族，使他们住在地上。} 这是天主特有的。那么，请看这些是否也可以归于圣子。他说：\Quote{既是天地的主}——他们视为属于天主的。他宣称创造是天主的工作，人类也是。\Quote{预先指定了他们存留的时间，和他们居住的疆界}（25–26节），\Quote{为叫他们寻求主，或者可以摸索而找到祂，其实祂离我们每个人并不远：因为我们生活、行动、存在都在祂内，正如你们自己的诗人也说过：\InnerQuote{我们也是祂的子孙。}}（27–28节）这是诗人阿拉托斯说的。请注意他如何从他们自己的作为和言论中引出论证。\Quote{既然我们是天主的子孙，就不该认为天主的神性仿佛金银或石头，受人工艺雕刻而成。} \BibleRef{宗 17:29} 然而正因如此，我们理当这样认为？绝非如此：因为我们当然不像（那些东西），我们的灵魂也不是。\Quote{和人的想象。} 怎么会呢？* * 但有人可能说：\Quote{我们并不这样想。} 但他是对众人讲话，而非对哲学。那么他们怎么对祂有如此不相称的想法呢？他再次归咎于他们的无知，说：\Quote{如今，天主对无知的时代暂且放过。} 他以恐惧激动他们的心思，然后加上这句话：然而他说：\Quote{但如今祂吩咐各处的人都要悔改。} \BibleRef{宗 17:30} \Quote{因为祂已定了一个日子，要藉着祂所设立的那人，按正义审判天下；祂已叫祂从死人中复活，向众人提供了凭据。} \BibleRef{宗 17:31} 但让我们重看前面所说的。\par

(摘要。)(\Emph{b})\BibleQuote{\Person{保禄}等候的时候，}等等。\BibleRef{宗 17:16}。这是上智的安排：他违背自己的意愿留在那里，等候那些人。(\Emph{a})\Quote{他的心神，}它说，\Quote{在他内}\OriginalTerm{激动}{\GreekText{παρωξύνετο}}。这不表示愤怒或恼怒；正如别处所说：\BibleQuote{他们之间发生了\OriginalTerm{争执}{\GreekText{παροξυσμὸς}}。}\BibleRef{宗 15:30}(\Emph{c})那么\OriginalTerm{激动}{\GreekText{παρωξύνετο}}是什么？被激发：因这恩宠远非愤怒和恼怒。他无法忍受，反而憔悴。\BibleQuote{所以他在会堂里辩论，}等等。\BibleRef{宗 17:17}注意他再次与犹太人辩论。所谓\Quote{虔诚的人}他指归化者。因为犹太人在基督来临之前分散各处（现代文本作\Quote{自从}），法律可说已渐渐瓦解，但同时（分散的犹太人）教导人宗教。但他们毫无成就，只是成了自己灾祸的见证人。(\Emph{e})\BibleQuote{有些哲学家，}等等。\BibleRef{宗 17:18}他们为何愿意与他交谈？当他们看见别人辩论，而那人（在交锋中）有声誉时。注意他们立刻以傲慢无礼说：\BibleQuote{有人说：这个多言者要说什么？因为属血气的人不领受圣神的事。}\BibleRef{格前 2:14}另一些人说：他似乎是传讲奇怪鬼神的。\OriginalTerm{鬼神}{\GreekText{δαιμονίων}}，因为他们如此称呼他们的神。\BibleQuote{他们拿住他，带他到}等等。\BibleRef{宗 17:19}(\Emph{a})雅典人不再享有自己的法律，已臣服于罗马人。(\Emph{g})（那么）他们为何拉他到\Place{亚略巴古}？为要威吓他——那是他们审理流血案件的地方。\Quote{我们能知道你所讲的新道理吗？因为你把奇怪的事传到我们耳中；我们愿意知道这些事是什么意思。原来所有的雅典人和侨居的外邦人，没有别的事比谈论或听些新鲜事更感兴趣。}（20–21节）这里值得注意的是：虽然他们只忙于谈论和听闻，却认为那些事奇怪——他们从未听过的事。\BibleQuote{于是\Person{保禄}站在\Place{亚略巴古}当中说：诸位雅典人，我看出你们在各方面更敬畏鬼神。}\BibleRef{宗 17:22}(\Emph{f})因为城市充满鬼神（\OriginalTerm{鬼神}{\GreekText{δαιμόνων}}，另作\OriginalTerm{偶像}{\GreekText{εἰδώλων}}）：(\Emph{h})这就是为什么他说\OriginalTerm{更敬畏鬼神}{\GreekText{δεισιδαιμονεστέρους}}。因为我巡行时，观看你们所敬拜的对象——他没有简单说\OriginalTerm{鬼神}{\GreekText{τοὺς} \GreekText{δαίμονας}}，而是为他的讲论铺路：\BibleQuote{我看见一座祭坛，}等等。\BibleRef{宗 17:23}这就是为什么他说：\Quote{我看出你们更敬畏鬼神}，即因那祭坛。他说：\BibleQuote{天主创造了世界。}\BibleRef{宗 17:24}他说了一句话，推翻了所有哲学家的学说。因为伊壁鸠鲁派认为万物偶然形成，由原子（聚合）；斯多亚派认为它是物体和火（\OriginalTerm{火化}{\GreekText{ἐκπύρωσιν}}）。\Quote{世界和其中所有的一切。}你注意到简洁，以及在简洁中的清晰吗？注意那些他们对之感到奇怪的事：天主创造了世界！现在连最普通的人都知道的事，雅典人和雅典的智者们却不知道。\Quote{祂是天地的主宰：}因为若祂创造了它们，显然祂是主。注意他宣称神性的标志——创造。而圣子也拥有这属性。\par

因为先知们处处肯定这一点，即创造是天主的特权。并非如那些人肯定有另一位是创造者却不是主，他们假定物质是未被创造的。现在他暗中肯定并确立自己的主张，同时推翻他们的学说。\Quote{不住人手所造的殿宇。}因为祂的确住在殿宇里，但非住在这种殿宇，而是住在人的灵魂里。他推翻了物质性的崇拜。那么怎样？祂不是住在耶路撒冷的圣殿里么？实在不是，祂只是在那里工作。\Quote{也不受人手的供奉。}\BibleRef{宗 17:25}那么在犹太人中间，祂又如何被人手供奉？并非用手，而是用理智。\Quote{似乎祂有所需要：}因为就连那些（崇拜行为）祂也不是这样寻求，\Quote{仿佛有所需要。我岂要吃公牛的肉，喝山羊的血？}\BibleRef{咏 49:13}这还不够——他肯定了一无所缺——因为这虽然是天主性的，但还需加添另一个属性。\Quote{因为祂将生命、气息和万有赐予众人。}天主性的两个证明：祂自己一无所缺，并供应万有给众人。让柏拉图和他所有关于天主的哲学，以及伊壁鸠鲁所有的一切都拿出来：与这相比，统统不过是琐屑！\Quote{祂赐予，}他说，\Quote{生命与气息。}看，他也使祂成为灵魂的创造者，而非灵魂的生育者。再看他是如何推翻关于物质的学说的。\Quote{祂从一人造出万族，使他们住在全地之上。}\BibleRef{宗 17:26}这些事比前者更好：这是对原子和物质的何等控诉，表明创造不是局部的工作，人的灵魂也不是局部的。但那些人所说的，并非成为创造者。——但祂是由理智和理解所崇拜的。——\Quote{是祂所赐予}等等。祂不是那些\OriginalTerm{局部神祇}{\GreekText{μερικοὶ} \GreekText{δαίμονες}}。\Quote{万有。}是\Quote{祂，}他说。——人又是如何形成的。——首先他表明\Quote{祂不住在}等等，然后宣称祂\Quote{不被供奉如同有所需要。}如果祂是天主，祂便造了万有；但若祂没有造，祂便不是天主。那些没有造天地的神，愿它们灭亡。他引入了更伟大的教义，虽然尚未提及伟大的教义；但他像对孩童一样向他们讲述。而这些比那些更为伟大。创造、主权、一无所缺、万善之源——这些他都已宣告。但祂是如何被崇拜的？你说。现在还不是适当的时候。有什么能与这崇高相比？这也是奇妙的——由一位造出如此众多的人：而且，在造了之后，祂亲自维持它们存在(\OriginalTerm{维持}{\GreekText{συγκρατεἵ}})，\Quote{赐予生命、气息和万有。}(\Emph{b})\Quote{祂预先定准了时期和居住的疆界，好叫他们寻求天主，或者可以摸索到祂而找到祂。}\BibleRef{宗 17:27}(\Emph{a})这意思或是：祂没有强迫他们周游世界寻求天主，而是按照他们居住的疆界；(\Emph{c})或是指：祂限定了他们寻求天主，但并非限定这事（去作）持续进行，而是限定了某些指定的时期（他们应当寻求）；现在显出：他们未经寻求却已找到；因为他们虽然寻求却没有找到，他表明天主如今已如此显明，好像\OriginalTerm{可触知地}{\GreekText{ψηλαφώμενος}}在他们中间。(\Emph{e})\Quote{其实祂离我们每个人都不远，}他说，\Quote{却离众人很近。}再看到天主的德能（或：\Quote{什么是天主}）。他说什么？祂不但赐予\Quote{生命、气息和万有，}而且，作为万有的总纲和实质，祂藉着赐予我们这些使我们能够找到并领悟祂的事物，引领我们认识祂自己。但我们并不愿意找到祂，尽管祂近在咫尺。\Quote{其实祂离我们每个人都不远。}请看！祂离所有人、全世界的每个人都很近！还有什么比这更伟大？看他如何清楚摆脱了那些\OriginalTerm{局部神祇}{\GreekText{τοὺς} \GreekText{μερικούς}}！我说什么\Quote{遥远？}祂如此之近，以至于没有祂我们无法生存：\Quote{因为我们生活、行动、存在都在祂内。}\BibleRef{宗 17:28}\Quote{在祂内：}用身体上的比喻来说，就像无法对弥漫在我们四周的空气无知，而空气\Quote{离我们每个人都不远，}不，其实是在我们内。(\Emph{d})因为并非在一处有天堂，在另一处却没有，也非在某些时间有，另一些时间却没有。所以，在每个\Quote{时间}和每个\Quote{疆界}都能找到祂。祂安排事务，使人们不为地方和时间所阻碍。因为，即使没有别的，单单这一点本身也是他们的帮助——天堂无处不在，它存于所有时间。(\Emph{f})看他如何（宣告）祂的眷顾和祂的支持能力(\OriginalTerm{支持}{\GreekText{συγκράτησιν}})；万有由祂而来，（由祂而来）它们的\OriginalTerm{运作}{\GreekText{τὸ} \GreekText{ἐνεργεἵν}}，（由祂而来的保存）使它们不消亡。他并没有说\Quote{藉着祂，}而是说了更亲密的\Quote{在祂内。}——那位诗人所说的话比不上这个：\Quote{因为我们是祂的子孙。}然而，他是指朱庇特，而保罗却取其指创造者，并非指与他相同的神，绝无此意！而是指那真正属于天主的谓语：正如他提及祭台是指向祂，而非指向他们所崇拜的神一样。正如说：\Quote{因为某些言语和行动是指向这位（真天主的），但你们不知道它们是指向祂。}试问，\Quote{献给未识之神}这话是适切地说到谁？是创造者，还是邪神？显然是创造者：因为祂他们并不认识，但他们认识另一位。再者，万物都充满（临在）——是天主的还是朱庇特的？朱庇特是一个可鄙的人，一个可憎的骗子！但保罗说这话并非与他同义，绝无此意！而是用完全不同的含义。因他说我们是天主的子孙，即天主的所属，仿佛是祂最亲近的邻人。\par

因为恐怕当他说：\BibleQuote{是天主的子孙} \BibleRef{宗 17:29}, 他们又会说：你带些新奇的事到我们耳中，他就引用了诗人。他不说：\BibleQuote{你们不应该把神性比作金或银,} 你们这些被诅咒和可憎的：而是以更卑微的方式说：\BibleQuote{不应该.} 因为（他说）什么？天主高于这个？不，他也没有这样说：但暂时是这样——\Quote{我们不应该认为神性像这样,} 因为没有比这更与人对立的了。\Quote{但我们不肯定神性像这样，因为谁会这样说呢？} 请注意他是如何引出无形体（的天主本性）的，当他说：\Quote{在他内,} 等，因为当心灵猜疑物体时，同时也暗示了距离的概念。（讲话）对群众他说：\BibleQuote{我们不应该认为神性像金、银或石头，是人手工艺的塑造,} 因为如果我们在灵魂方面不像那些（物质），那么天主（更不像）。至此，他使他们远离这种概念。但是，他也会说，神性也不屈从于任何其他人类概念。因为如果艺术或思想所发现的——这就是他为何这样说：\BibleQuote{属于人的技艺或想象}——如果那样，那么，人类艺术或思想所发现的就是天主，那么甚至在石头中也有天主的本质。——那么，如果我们\Quote{在他内生活,} 怎么我们找不到他呢？指责是双重的，既因为他们没有找到他，也因为他们找到了如此的东西。（人类）的理解本身完全不可信赖。——但当他通过表明他们无可推诿而激动了他们的灵魂时，看他说什么：\BibleQuote{天主对无知的时代不予追究，但现在命令各处的人都要悔改。} \BibleRef{宗 17:30} 那么如何？这些人中没有一个要受惩罚？愿意悔改的人中没有一个。他说的是这些人，不是那些已逝的人，而是他所命令悔改的人。他不追究你们，他会说。他不说：\OriginalTerm{不加注意}{\GreekText{παρεἵδεν}}；也不说：允许：而是说：你们是出于无知。\Quote{忽略,} 即不要求像对待应受惩罚的人那样予以惩罚。你们是出于无知。他也不说：你们故意作恶；而是通过上述的话表明了这一点。——\BibleQuote{各处的人都要悔改:} 他再次暗示整个世界。注意他如何使他们脱离零散的神祇！\BibleQuote{因为他已定了一个日子，要借着他所指定的人，按正义审判天下，为此他已给众人一个凭据，就是使他从死者中复活。} \BibleRef{宗 17:31} 请注意他如何再次宣告受难。再次注意那恐惧：因为审判是真实的，从他使他复活来看是清楚的：因为那是为此提出的证据。他所说的一切都是真实的，从他复活的事实来看是清楚的。因为他确实给了这个\BibleQuote{给众人的凭据,} 即他从死者中复活：这事（即审判），从此也是确定的。\par

\Quote{各处的人都必须悔改，因为祂已指定了日子，在那日祂要审判世界。}请看他是如何也将祂当作审判者引入的：祂，既眷顾世界，又仁慈、宽恕、有权能、有智慧，总而言之，具有创造者的一切属性。\Quote{给众人提供了保证}，即祂在（耶稣）从死者中复活的事上提供了证据。那么，让我们悔改吧，因为我们必被审判。如果基督没有复活，我们就不会被审判；但如果祂复活了，我们无疑将被审判。经上说：\Quote{正是为此，祂死了，为使祂成为死人和活人的主。}\BibleRef{罗 14:9}\Quote{因为我们众人都要站在基督的审判台前，好使每人按他所行的领受报应。}\BibleRef{格后 5:10}不要以为这些只是空话。看哪！他也引入了所有人复活的主题；因为世界无法以其他方式被审判。而那句\Quote{因为祂已使祂从死者中复活}，是指身体而言；因为那身体是死了的，是已经堕落了的。在希腊人中，他们关于创造的观念，同样关于审判的观念，都是儿童的幻想，是醉汉的胡言乱语。但我们这些准确知道这些事的人，应做一些有实际意义的事：让我们成为天主的友人。我们与祂为敌要到何时？我们讨厌祂要到何时？\Quote{绝对不可！}你会说：\Quote{你为什么说这样的话？}如果你们不做你们所做的事，我倒不希望说我现在所说的话；但事实上，当行为的明显证据如此高声呐喊时，现在保持沉默又有什么用呢？那么，我们该如何爱祂呢？我已告诉你们千百种方式，千百次了；但现在我也要再说一遍。我似乎发现了一个非常伟大而奇妙的方式。即，在向祂承认我们普遍欠下的恩情——没有人能表达（我的意思是），祂为我们每人亲自所做的——之后，我们也要想想这些，因为它们有巨大的力量：让我们各人与自己计算，并仔细查考，如同在一本书里记下天主的恩惠；例如，若曾陷入危险而从敌人手中逃脱；若曾在不合时宜的时刻出门旅行而脱险；若曾与恶人相遇而胜过他们；或曾患病而被众人放弃后又康复：因为这大大有助于我们依附天主。因为如果那位\Person{摩尔德开}，当他为国王所作的事被记起时，他竟获得这么大的荣耀\BibleRef{艾 6:2}，那么，我们若回忆并仔细查考这两点：我们犯了什么罪得罪天主，祂又为我们行了什么善，我们就会既感恩，又将我们所有的一切都自由地献给祂。但没有人关心这些事：关于我们的罪，我们只说我们是罪人，却不具体查问它们；同样，关于天主的恩惠，我们说天主行了善，却不具体查问在哪里、有多少、在什么时候。但从现在起，让我们非常精确地计算。因为若有人能回忆起那些很久以前发生的事，让他准确地全部计算，如同寻找大宝藏的人。这对我们也有益处，使我们不致绝望。因为当我们看到祂常常保护我们，我们就不会绝望，也不会以为被抛弃；反而会视之为祂眷顾我们的强有力保证，当我们想到，虽然我们犯了罪，却没有受罚，反而享受祂的保护。现在让我告诉你们一个我从某人那里听说的案例，关于一个孩子，有一次他和母亲在乡下，还不满十五岁。那时正有一股坏空气，结果两人都发烧，因为正是秋季。母亲设法先赶进了城，但男孩，当当地的医生命令他，在体内发烧的情况下漱口时，他拒绝了，自以为有道理，认为如果不吃任何东西就能更好地灭火，因此，在他不合时宜的反抗精神下，像孩子一样，他什么都不吃。但当他进了城，他的舌头瘫痪了，长时间不能说话，因此无法清晰地发音；然而，他确实能读书，并跟随老师很长时间，但仅此而已，没有任何进步。因此他所有（生活）的希望都破灭了，他的母亲充满了悲伤：尽管医生提出了许多方案，许多人也这样做，但没有人能治好他，直到仁慈的天主解开了他舌头的结\BibleRef{谷 7:35}，然后他康复了，恢复了以前的流利和清晰。他的母亲也讲述，当他还是很小的孩子时，他鼻子里长了个东西，他们称之为息肉：那时医生们也放弃了他，他的父亲咒骂他（因为父亲还活着），甚至他的母亲也祈祷他死；一切都充满痛苦。但他突然咳嗽，由于粘液的积累，用呼吸的力量将\OriginalTerm{那活物}{\GreekText{τὸ} \GreekText{θηρίον}}从鼻孔中排出，所有危险都消除了。但这场灾祸平息后，一种辛辣粘稠的眼液形成了厚厚的\OriginalTerm{眼垢}{\GreekText{τὰς} \GreekText{λήμας}}，像一层皮肤覆盖在瞳孔上，更糟的是，它威胁要失明，人人都说会是这个结局。但靠着天主的恩宠，他也很快摆脱了这种疾病。以上是我从别人那里听来的：现在我告诉你们我自己知道的一件事。有一次，在我们城中出现了暴君嫌疑——那时我还年轻——所有士兵都被派在城外警戒，正巧，他们严格搜查巫术和魔法书籍。而写那本书的人，把书\OriginalTerm{未装订的}{\GreekText{ἀκατασκέυαστον}}扔进了河里，并被捕了，当被要求交出书时，他交不出来，被捆绑着在城中各处游行；然而，当证据指向他，他受了惩罚之后，恰巧我正想去殉道者教堂，和另一个人一起从河边花园返回。他看到书浮在水面上，起初以为是块亚麻布，但靠近后发现是本书，于是他下去捡了起来。我却要求分赃，并对此发笑。但他说，让我们看看这到底是什么东西。于是他翻开一部分书页，发现内容涉及魔法。就在这时，一个士兵正好路过：* * * 然后他从里面拿出书，走开了。我俩吓得僵住了。因为谁会相信我们从河里捡到这本书的故事呢？那时所有人，甚至连没有嫌疑的人，都受到严密监视。我们不敢扔掉它，怕被人看见，撕碎它也同样危险。天主赐给我们办法，我们扔掉了它，最终暂时脱离了极大的危险。如果我要全部列举，我可以数出无数案例。我提到这些也是为了你们，所以，如果任何人还有其他案例，即使不像这些，让他也时常记在心里：例如，若曾有一块石头扔向你，刚要击中你却没有击中，你要永远记住这件事：这些事在我们心中产生对天主的大爱。因为如果想起那些曾救过我们的人，我们因无法回报他们而深感惭愧，那么对于天主更应该如此。这在其他方面也有用。当我们不想过分悲伤时，让我们说：\Quote{如果我们从主手中接受了福，难道不该接受祸吗？}\BibleRef{约 2:10}当\Person{保禄}告诉他们他从哪里被救出来时，\BibleRef{弟后 4:17}，原因是为了提醒他们。也要看\Person{雅各伯}如何将所有这一切记在心里：因此他也说：\Quote{那从我幼年以来救赎我的天使}\BibleRef{创 48:16}；不仅是他救赎了他，而且是如何以及为何目的。相应地，看他如何也特别回忆起他所领受的恩惠。他说：\Quote{我带着我的杖}\BibleRef{创 32:10}\Quote{渡过了约但河}。犹太人还常常记念他们祖先所经历的事，心中思索在埃及所发生的事。那么，我们更当牢记发生在我们身上的特殊恩惠，我们多少次陷入危险和灾祸，若不是天主用手扶着我们，我们早就灭亡了：我说，让我们所有人，天天思考这些事情并讲述它们，共同向天主感恩，永不停止光荣祂，这样我们就能因内心的感恩而获得丰厚的赏报，藉着祂唯一圣子的恩宠和怜悯，祂与圣父及圣神同享光荣、权能、尊威，从今世直到永远。阿们。\par

\SourceRecord{Translated by J. Walker, J. Sheppard and H. Browne, and revised by George B. Stevens.；Nicene and Post-Nicene Fathers, First Series；Vol. 11.；Edited by Philip Schaff.；Buffalo, NY: Christian Literature Publishing Co.,；1889.；<http://www.newadvent.org/fathers/210138.htm>.}

% structure: p0001 source=h1 target=section reason=page-title

\section{论宗徒大事录 第三十九篇讲道}

见：\BibleRef{宗 17:32-34; 18:1}\par

\BibleQuote{他们一听见死人复活的事，有的讥讽，有的说：\InnerQuote{关于这事，我们还要再听你讲。}于是保禄从他们当中离去。然而有些人依附他而相信了，其中有亚略巴古的狄约尼削，还有一个名叫达玛黎的妇人，以及其他同他们一起的人。此后，保禄离开雅典，来到格林多。}\par

是什么原因使保禄说服了（一些人甚至说）他们还会再听他讲，又没有危险，却如此急于离开雅典？大概他知道自己不会给他们带来什么大好处；再者，他也是因圣神而到了格林多。（\Emph{b}）因为雅典人虽然喜欢听新奇的事，却并不留意（他）；因为这并非他们的追求，而只是总要有话可说，这使他们疏远了他。但如果这是他们的习惯，他们又为何控告他，说\BibleQuote{他似乎是讲论外邦神明的}？见：\BibleRef{宗 17:18}。然而，这些事他们根本不知道如何理解。不过，他确实归化了亚略巴古的狄约尼削和另外一些人。因为那些关注（正当）生活的人，很快就接受了圣言；但其他人则不然。在保禄看来，播下教义种子已经足够了。（\Emph{a}）于是，如前所述，他被圣神领到格林多，他要在那里居住。（\Emph{c}）\BibleQuote{遇见一个犹太人，名叫阿桂拉，本都生，因喀劳狄曾命令所有犹太人都离开罗马，他刚从意大利来}——因为他大半生都在那里度过——\BibleQuote{和他的妻子普黎史拉。}见：\BibleRef{宗 18:2}。虽然对犹太人的战争是在尼禄在位时完成的，但从喀劳狄那时起，它就已经在远处酝酿了，这样，即使只是这样，他们也可能会醒悟过来；现在他们像共同的祸害一样被赶出罗马。因此天主上智安排保禄被作为囚犯带到那里，这样他就不会作为犹太人被驱逐，而是因受兵士看管而甚至可以在那里受到保护。（找到他们后，）\BibleQuote{他投奔了他们，因为同业，便留在他们那里作工；原来他们是制帐幕为业的。}见：\BibleRef{宗 18:3}。看，他找到了与他们同住一处的何等正当的理由！因为在这里，尤其在别的地方，他必须不收钱，就如他自己说的：\BibleQuote{为使他们所夸耀的，在我们身上也找不到。}见：\BibleRef{格后 11:12}。因此，上智安排他住在那里。\BibleQuote{每逢安息日，他在会堂里辩论，劝化了犹太人和希腊人。及至息拉和弟茂德从马其顿来到后，保禄就被圣言所困，向犹太人证明耶稣就是基督。}（第4、5节）\BibleQuote{及至犹太人反对并亵渎时}——即他们试图压倒他（\OriginalTerm{攻击}{\GreekText{ἐπηρέαζον}}），攻击他——于是保禄做什么？他与他们分离，而且以非常严厉的方式；虽然他现在没有说：\BibleQuote{天主的话本该先传给你们}，却暗地里向他们暗示了这一点：\BibleQuote{他们既反对并亵渎，保禄就抖着衣服说：\InnerQuote{你们的血归到你们自己的头上；我是洁净的。从今以后，我要到外邦人那里去。}}见：\BibleRef{宗 18:6}\BibleQuote{他就从那里离开，进了一个名叫犹斯托的人的家里，这人敬畏天主，他的家紧挨着会堂。}你看，他虽然又说\BibleQuote{从今以后——}，却并未忽略他们；所以他这样说是为了激发他们，随后就来到犹斯托家，其房屋与会堂相邻，这样即使从这邻近本身，他们也可能会起嫉妒。\BibleQuote{会堂长克黎斯颇和他全家都信了主。}这件事实在足以使他们转变。\BibleQuote{有许多格林多人听了，就相信并领了洗。夜间，主藉异象对保禄说：\InnerQuote{不要害怕，只管讲，不要缄默，因为我与你同在，没有人会下手害你，因为在这城里有许多属我的人民。}}见：\BibleRef{宗 18:8}。你看，祂用多少理由说服他，又把最能使他说服的理由放在最后：\BibleQuote{在这城里有许多属我的人民。}那么，你会问，他们又怎会向他下手呢？然而，作者告诉我们，他们没有得逞，只是把他带到总督面前。\BibleQuote{他就在那里住了一年六个月，在他们当中教导天主的话。当加里雍作阿哈雅总督时，犹太人齐心起来攻击保禄，把他带到审判台前。}（第11、12节）你注意到这些人为何总是设法将他们所控告的罪行公开化吗？请看这里：（\Emph{b}）\BibleQuote{说：\InnerQuote{这个人劝人违反法律敬拜天主。}保禄刚要开口，加里雍就对犹太人说：\InnerQuote{如果是犯法或邪恶的罪行，犹太人，我理当容忍你们。但如果是关于言论、名字和你们法律的问题，你们自己看吧！我不愿做这些事的判官。}就把他们赶出了审判台。}见：\BibleRef{宗 18:13}。这加里雍在我看来是个明白人。（\Emph{a}）你看，当这些人说\BibleQuote{他劝人违反法律敬拜天主}时，他\BibleQuote{对这些事一概不理}；留意他如何回答他们：\BibleQuote{如果真是}任何涉及城市的事，\BibleQuote{犯法或邪恶的罪行}等等。（\Emph{c}）\BibleQuote{于是众犹太人捉住会堂长索斯特乃，在审判台前打了他；加里雍对这些事一概不理。}见：\BibleRef{宗 18:17}。但他并未将他们打他视为对自己的侮辱。犹太人是如此暴躁。但让我们重新审视一下前面所说的话。\par

（摘要。）\Quote{他们听见了}，\BibleRef{宗 17:32}多么伟大高尚的道理，他们却不在意，反而讥笑复活的事！\Quote{因为属血气的人，}经上说，\Quote{不能领受圣神的事。}\BibleRef{格前 2:14}\Quote{于是，}经上说，\Quote{保禄出去了。}\BibleRef{宗 17:33}怎样出去的？有些人被说服了；另一些人讥笑他。\Quote{但有几个人，}经上说，\Quote{贴近他，信了，其中有阿勒约帕哥的狄约尼削和另一些人。}\BibleRef{宗 17:34}\Quote{此后，}等等。\Quote{遇见了一个犹太人，名叫阿桂拉，本都籍，最近从意大利来，因为喀劳狄曾命所有犹太人离开罗马；保禄来到他们那里，因是同业，就住在他们那里工作；因为他们以制造帐幕为业。}\BibleRef{宗 18:1}这个阿桂拉是本都人* * *。请注意，他不是在耶路撒冷或附近（危机中）急忙前来，而是在更远的地方。保禄同他住在一起，不以同住为耻，反而正为此住在他那里，因为他有一个合适的住所，这住所对保禄来说比任何王宫都合适。亲爱的，你听到他的职业时，不要发笑。因为（对他来说）这就像运动员的角力学校比精致的地毯更有用，战士的铁剑比金剑更有用一样。\Quote{他工作，}虽然他在宣讲。我们该羞愧，我们虽然没有宣讲之责，却过着懒散的生活。\Quote{每安息日在会堂里辩论，说服犹太人和希腊人}\BibleRef{宗 18:4}；但\Quote{当他们反对并亵渎时}，他退出了，借此希望能更好地吸引他们。因为为什么他离开那房子，来住在会堂附近？难道不是为了这个吗？他不是因为看到这里有危险。所以，保禄向他们作证——不是教导，而是作证——\Quote{抖了抖衣服}，不仅用言语，也用行动恐吓他们，\Quote{对他们说：你们自己的血归在你们自己头上}\BibleRef{宗 18:6}：他说得更加激烈，因为已经说服了许多人。\Quote{我，}他说，\Quote{是洁净的。}那么，如果我们忽视委托给我们的人，我们也要为他们的血负责。\Quote{从今以后，我要到外邦人那里去。}所以，当他说\Quote{从今以后，不要再有人麻烦我}\BibleRef{迦 6:17}时，他也是为了恐吓。因为惩罚本身并不那么可怕，而是这刺痛了他们。\Quote{离开那里，他来到一个名叫犹斯托的人家里，这人恭敬天主，他的房子紧挨着会堂}\BibleRef{宗 18:7}，就住在那里，借此希望说服他们是认真的（\OriginalTerm{他急切地转向外邦人}{\GreekText{πρὸς} \GreekText{τὰ} \GreekText{ἐθνη} \GreekText{ἠπείγετο}}）要往外邦人那里去。因此，注意当他一这样做，会堂长立刻归依了，还有许多其他人。\Quote{克利斯颇，会堂长，同他全家信了主；许多格林多人听了，也信了，并领了洗。}——\BibleRef{宗 18:8}\Quote{同他全家：}注意那时代的归依者是全家一起行动的。这个克利斯颇就是他写信时提到的：\Quote{我只给克利斯颇和加约付了洗}\BibleRef{格前 1:14}。我认为（同一个人）被称为索斯特乃——显然是信徒，以至于挨打，并且常与保禄在一起。\Quote{主夜间对保禄说，}等等。现在，即使是人数（\Quote{许多人}）说服了他，但基督宣称他们是自己的（更打动他）。但祂还说：\Quote{不要怕：}因为现在危险更大了，既因更多人信了，也因会堂长信了。这足以激起他的情绪。并不是责备他害怕；而是使他不要受任何伤害；\Quote{我与你同在，没有人会下手害你。}（第9、10节。）因为祂并不总是允许他们受苦，免得他们变得太软弱。因为没有什么比人们的不信和反对（真理）更让保禄痛心：这比危险更糟。因此现在（基督）向他显现。\Quote{他住了一年零六个月，}等等。\BibleRef{宗 18:11}一年零六个月之后，他们向他下手。\Quote{当加里雍作阿哈雅的总督时，}等等（第12、13节），因为他们不能再利用自己的法律了。（\Emph{c}）请注意他多么谨慎：他没有直接说\Quote{我不在乎}，而是说：\Quote{如果，}他说，\Quote{真的是为非作歹或邪恶淫乱的事，犹太人哪，按理我该容忍你们；但如果问题是关于道理、言语和你们的法律，你们自己看吧，因为我不愿意做这些事的审判官。}（第14、15节。）（\Emph{g}）他教导他们，这些不是需要司法判决的事，他们做一切事都无序。他没有说\Quote{这不是我的职责}，而是说\Quote{我不愿意}，免得他们再麻烦他。同样，比拉多在基督的案件中说：\Quote{你们自己带他去，按你们的法律审判他吧。}\BibleRef{若 18:31}但他们就像喝醉和疯狂的人。（\Emph{d}）\Quote{他把他们赶出了审判台}\BibleRef{宗 18:16}——他有效地关闭了法庭对他们不利。\Quote{然后众人}（犹太人）\Quote{抓住会堂长索斯特乃，在审判台前打他。加里雍对这些事一概不理。}\BibleRef{宗 18:17}。（\Emph{a}）这件事，比其他任何事，更促使他们（行暴）——他们确信总督甚至不会屈尊（注意）。（\Emph{e}）这是一次辉煌的胜利。哦，他们所遭受的羞耻！（\Emph{b}）因为从辩论中获胜是一回事，而那些人得知他对这事毫不关心又是另一回事。（\Emph{f}）\Quote{加里雍对这些事一概不理：}然而整个事件都是对他的侮辱！但，仿佛他们得到了授权（他们就这样做了）。为什么他（索斯特乃），虽然也有权柄，却没有打（他们）？但他们（原来）被训练：这样审判官就知道哪一方更合理。这对在场的人不是小益——这些人的合理和那些人的放肆。（\Emph{h}）他被打，什么也没说。\par

我们也该效法此人：对那殴打我们的人，我们要以温良、以沉默、以忍耐还击。更痛苦的是这些创伤，更重的是这打击，更沉的是这伤害。因为为表明身体受击并非痛苦，而是心中受击才痛苦，我们常常打人，但若是出于友谊，他们甚至喜欢；但若以傲慢方式打一个无关的人，你便使他极其痛苦，因为你触动了他的心。所以我们要击打他们的心。但温良比暴烈造成更大的打击，来，让我们尽可能用言语证明。确实，可靠的证明是靠行动和经验；但若你愿意，让我们也用言语来探究，尽管我们已多次这样做过。在侮辱中，没有什么比旁观者的意见更使我们痛苦；因为公开受辱与私下受辱不同，同样的侮辱，若在僻静无人见证或知晓的地方受之，我们甚至能轻易忍受。可见，使我们沮丧的不是侮辱本身，而是当着众人受辱；因为若有人在众人前尊敬我们，私下却侮辱我们，我们反倒会感激他。所以痛苦不在于侮辱的性质，而在于旁观者的意见：免得自己显得可鄙。那么，如果这意见偏向我们呢？那个企图羞辱我们的人自己不是更受羞辱吗，当人们意见偏向我们时？请问，旁观者鄙视谁？是那侮辱人的，还是那受辱却保持沉默的？情欲的确暗示他们鄙视受辱者；但让我们现在在不受激动时审视这事，以免到时候被冲昏头脑。请问，我们大家谴责谁？显然是那侮辱人的人；若他是个下等人，我们会说他甚至疯了；若是个平等人，说他愚蠢；若是个上等人，我们仍不会赞同他。因为，我问，哪个值得称赞？是那激动、被情欲风暴翻腾、像野兽般狂怒、如此违背共同本性的人，还是那生活在平静中、在安宁港湾、在德行平稳中的人？一个不像天使，另一个甚至不像人吗？因为一个连自己的恶都不能忍受，而另一个还忍受他人的恶：这里，人连自己都不能容忍；那里，他还容忍别人：一个处于沉船危险，另一个安全航行，船只顺风漂行：因为他没有让情欲的突风抓住帆并倾覆理智之舟；而是温柔甜蜜的气息，即宽容的气息，吹拂其上，使其非常平静地驶入智慧平稳的港湾。犹如当船有沉没危险时，水手们不知丢弃何物，他们拿到手里的是自己的还是他人的财产，他们不加区别地抛下所有货物，无论珍贵的与否：但当风暴平息后，他们清点所丢弃的一切，便流泪，因损失所抛之物而感受不到平静：同样，当情欲猛吹，风暴兴起时，人们抛出言语，不知如何有秩序或得体地使用；但当情欲平息后，回想他们曾说出怎样的言语，他们便思量损失，因记得那些使自己丢脸、遭受最严重损失——不是金钱损失，而是关于节制与温良品德的损失——的言语，而感受不到宁静。愤怒是黑暗。圣经说：\Quote{愚妄人}，圣经说，\Quote{心里说：没有天主。} \BibleRef{咏 12:1} 或许也适合对愤怒的人说同样的话：愤怒的人心里说：没有天主。因为圣经说：\Quote{因他怒气填胸，必不寻求} \BibleRef{咏 9:4}。因为无论什么虔诚思想进入，（愤怒）都推开并赶出一切，横冲直撞。（\Emph{b}）当你被告知，你所辱骂的人没有说一句刻薄话，难道你为此不比你加给他的痛苦更痛吗？（\Emph{a}）若你内心里不比你辱骂的人更痛，就继续辱骂吧；（但）即使没有人向你追责，你良心的审判私下抓住你，便要给你千次鞭打，（当你想到）你如何向一个如此温良、谦卑、忍耐的人倾泻了一连串辱骂。我们总是说这些事，但未看到它们体现在行动上。你，一个人，侮辱你的同伴？你，一个仆人，侮辱你的同仆？但我为何对此惊奇，当许多人甚至侮辱天主？让这成为你受辱时的安慰。你受辱？天主也受辱。你被辱骂？天主也被辱骂。你被藐视？我们的主也同样。在这些事上，祂与我们分享，但在相反的事上却不然。因为祂从未不义地侮辱过人：天主不准！祂从未辱骂，从未行不义。所以，我们才是与祂分享的人，而不是你们。因为受辱时忍耐是天主的事；单单辱骂是魔鬼的事。看这两面。\Quote{你附有魔鬼} \BibleRef{若 7:20}，有人对基督说；祂从大司祭的仆人那里挨了一记耳光。那些不义地侮辱人的人，与这些同类。因为如果伯多禄因一句话甚至被称为\Quote{撒殚} \BibleRef{玛 16:23}，那么这些人行犹太人的事时，更应被称为\Quote{魔鬼的子女} \BibleRef{若 8:44}，因为他们行魔鬼的事。你侮辱；请问你是谁（竟这样做）？不，不如说，你侮辱的原因正是你什么也不是：没有一个人会侮辱。所以争吵中所说的\Quote{你是谁？}应当反过来：\Quote{侮辱吧，因为你什么也不是。}而实际说法是：\Quote{你是谁，竟侮辱人？}\Quote{比你更好的人}是回答。然而事实恰恰相反：但因我们问法错误，所以他们回答错误：所以错在我们。因为我们似乎认为侮辱是伟人的事，所以我们问：\Quote{你是谁，竟侮辱人？}所以他们这样回答。\par

然而，相反，我们应当说：\Quote{你辱骂人吗？继续辱骂吧：因为你什么都不是：}至于对那些不辱骂的人，应当这样说：\Quote{你是何人，竟不辱骂？——你已经超越了人性。}这才是高贵，这才是宽宏，不说任何不宽厚的话，尽管某人可能配得上你对他说这样的话。请告诉我，有多少人不配被处死？然而，法官并不亲自做这事，而是审问他们；连这事他也不是亲自做。但如果法官在审判时不得亲自与罪犯交谈，而是一切通过第三人介入，那么我们就更有义务不侮辱与我们地位相等的人；因为我们从他们那里得到的一切好处，与其说是羞辱了他们，不如说是让我们认识到我们羞辱了自己。那么，对于恶人，这就是为什么我们不可侮辱（甚至他们也不可）；对于善人，则另有原因，因为他们不配被辱；还有第三个原因，因为辱骂是不对的。但实际情况如何，请看结果：被辱者是个人，辱人者也是个人，旁观者也是人。那么该怎么办？难道要让野兽介入他们之间来调解吗？因为只剩下这个办法了。因为当作恶者和以作恶为乐者都是人时，调解者的角色就留给了野兽：正如当主人在屋里争吵时，除了仆人去调解他们之外别无他法——即使结果并非如此，但事情的性质要求这样——这里也是如此。——你辱骂人吗？你尽可如此，因为你甚至不是人。傲慢被认为是高贵之物，似乎属于伟人；但它更属于奴隶；而说好话则属于自由人。因为行恶是那些人的事，而忍受恶则是这些人的事。——就如某个奴隶偷了主人的财物，某个老巫婆——辱骂者就是这样的人。又像一个可憎的贼和逃犯，蓄意潜入，四处张望，想要偷东西：这个人也一样，像他那样留心观察四面八方，研究如何抛出一些（辱骂）。或者我们可以用另一种例子来形容他。就如一个人从屋里偷出肮脏的器皿，在众人面前拿出来，被偷的东西并不如此羞辱被偷的人，而是羞辱贼自己：同样，这人在众人面前说出他的话，并不因这些话羞辱别人，而是羞辱自己，因为他宣泄这些言语，污染了自己的舌头和心灵。因为当我们与坏人争吵时，就好像为了打一个浸在腐烂污秽中的人，而把自己的手插进脏物中玷污自己一样。因此，反思这些事，让我们逃避由此产生的祸害，保持清洁的舌头，这样远离一切辱骂，我们便能严谨地度过现世生命，并赖我们的主耶稣基督的恩宠与慈悲，获得那应许给爱祂之人的美好事物；愿荣耀、权能、尊崇归于祂，与圣父及圣神一起，自今而后，直至永远，世世无尽。阿们。\par

\SourceRecord{Translated by J. Walker, J. Sheppard and H. Browne, and revised by George B. Stevens.；Nicene and Post-Nicene Fathers, First Series；Vol. 11.；Edited by Philip Schaff.；Buffalo, NY: Christian Literature Publishing Co.,；1889.；<http://www.newadvent.org/fathers/210139.htm>.}

% structure: p0001 source=h1 target=section reason=page-title

\section{论宗徒大事录第四十篇讲道}

\BibleRef{宗 18:18}\par

\BibleQuote{保禄在那里又住了些日子，然后辞别弟兄们，从那里乘船往\Place{叙利亚}去，同行的有\Person{普黎史拉}和\Person{阿桂拉}；他在\Place{耕格勒}剃了头，因为他有誓愿。}\par

看看，法律是如何在瓦解；看看，他们是怎样被良心所约束。这是犹太人的习俗，按照誓愿剃头。但是，当时也应该有祭献\BibleRef{宗 21:26}，这里却没有。——\BibleQuote{尚且住了些日子：}这是\Person{索斯特乃}挨打之后。因为他必须再逗留一段时间，就这些事安慰他们。\BibleQuote{他乘船往\Place{叙利亚}去。}他为什么又要到\Place{叙利亚}去？正是在那里，\BibleQuote{门徒们首先被称为基督徒}\BibleRef{宗 11:26}；在那里，他曾\BibleQuote{被交托于天主的恩宠}\BibleRef{宗 14:26}；在那里，他成就了有关教理的这些事。\BibleQuote{同行的有\Person{普黎史拉}}——看，还有一位妇女——\BibleQuote{和\Person{阿桂拉}。}但是他把他们留在\Place{厄弗所}。这样做是很有理由的，就是要他们教导人。因为他们与他同住了那么久，学到了许多东西；然而他目前并没有让他们脱离犹太人的习俗。\BibleQuote{他来到\Place{厄弗所}，把他们留在那里；他自己却进了会堂，与犹太人辩论。他们请求他多住些日子，他没有同意；却辞别了他们，说：我务必到\Place{耶路撒冷}去，赶上那里的节庆。}因此，他无法前往\Place{亚细亚}，因为被催迫去处理紧急之事。所以你看到，他虽然被请求留下，但因为不能答应，急于离开，就\BibleQuote{辞别了他们。}然而，他并没有就这样离开他们，而是许下（会回来）：\BibleQuote{但是我若愿意，就必再回到你们这里来。于是他从\Place{厄弗所}开船。}\BibleRef{宗 18:19}\BibleQuote{他在\Place{凯撒勒雅}登岸，上了\Place{耶路撒冷}，问候了教会，然后下到\Place{安提约基雅}。在那里住了些日子后，又出发，按次序走遍\Place{迦拉达}和\Place{夫黎基雅}地方，坚固所有门徒。}\BibleRef{宗 18:22}他再次来到他以前访问过的地方。\BibleQuote{有一个犹太人，名叫\Person{阿颇罗}，生在\Place{亚历山大里亚}，是个有口才的人，对圣经很有研究，来到\Place{厄弗所}。}\BibleRef{宗 18:24}看，连有学问的人现在也急切起来，门徒从此往各处去。你注意到讲道的扩展了吗？\BibleQuote{这人已经在主的道上受了教导，并且在圣神内心神火热，讲话和教导主的事都很热切，虽然他只知道\Person{若翰}的洗礼。他开始在会堂里放胆讲论；\Person{阿桂拉}和\Person{普黎史拉}听了，就把他接来，更详尽地给他讲解了天主的道。}\BibleRef{宗 18:25}如果这人只知道\Person{若翰}的洗礼，他怎么能够\BibleQuote{心神火热}呢？因为圣神并不是那样赐下的。如果那些在他之后的人需要\Person{基督}的洗礼，那么他更需要了。那么该说什么呢？作者把这两件事连在一起，并不是没有意义的。在我看来，他是那与宗徒们一同受洗的一百二十人之一；否则，就是发生在\Person{科尔乃略}身上的事，也发生在这人身上。但是他并没有领受洗礼。那么，\BibleQuote{他们更详尽地给他讲解了}这句话，在我看来是指他也应该受洗。因为那十二位对准确的事一无所知，甚至连与\Person{耶稣}有关的事也不知道。而且他很可能是确实受了洗。但是如果这些\Person{若翰}的门徒，在那次洗礼之后又领受了洗礼，这对门徒们也是必要的吗？那么为什么需要水呢？这些人与他大不相同，他们甚至不知道是否有圣神。\BibleQuote{他心神火热，然后，在圣神内，只知道\Person{若翰}的洗礼：但这些人更详尽地给他讲解了。}\BibleRef{宗 18:27}他也愿意去\Place{阿哈雅}，这些人也鼓励他，并给了他推荐信。\BibleQuote{他到了那里，藉着恩宠大大帮助了那些信了的人；因为他有力地驳倒了犹太人，公开地用圣经证明\Person{耶稣}就是\Person{基督}。}\BibleRef{宗 18:28}\BibleQuote{当\Person{阿颇罗}在\Place{格林多}的时候，\Person{保禄}经过了上边地区（这是指我们读过的关于\Place{凯撒勒雅}和其他地方的）来到\Place{厄弗所}，遇见一些门徒}\BibleRef{宗 19:1}，\BibleQuote{就对他们说：你们信的时候，领受了圣神没有？他们回答说：我们连圣神的存在都没有听说过。\Person{保禄}说：那么你们受的是什么洗呢？他们说：是\Person{若翰}的洗礼。于是\Person{保禄}说：\Person{若翰}施的是悔改的洗礼，告诉百姓要信那在他以后来的，就是\Person{耶稣基督}。}\BibleRef{宗 19:2}因为他们连\Person{基督}都不信，这从他所说的话可以清楚看出：\BibleQuote{要信那在他以后来的。}他没有说\Person{若翰}的洗礼算不得什么，而是说它不完全。他并没有明说这一点，而是教导他们，于是许多人领受了圣神。\BibleQuote{他们听了这话，就因主\Person{耶稣}的名受洗。\Person{保禄}给他们覆手，圣神就降在他们身上；他们就说方言，也说预言。人数约有十二个}\BibleRef{宗 19:5}，因此他们很可能已经有了圣神，只是没有显出来。\BibleQuote{人数约有十二个。}\par

（重述。）\Quote{他们到了厄弗所，就把他们留在那里}\BibleRef{宗 18:19}：因为他不想带着他们同行，而把他们留在厄弗所。但他们后来住在格林多，保禄高度赞扬他们，并在写给罗马人的书信中问候他们。\BibleRef{罗 16:3}由此我似乎觉得，后来在尼禄时代，他们又回到了罗马，因为他们对那地方有感情，就是在克劳狄时代被驱逐出去的地方。\Quote{但他自己进了会堂。}在我看来，信友们仍然聚集在那里，因为他们没有立即离开他们。\Quote{当他们恳求他多住些日子，他不答应}\BibleRef{宗 18:20}，因为他急于赶往凯撒勒雅。\Quote{他到了凯撒勒雅}等，\Quote{走遍迦拉达和夫黎基雅地区，坚固众门徒。}（22-23节）在这些地区，他也只是再次\Emph{经过}，足以用他的临在来坚固他们。\Quote{有一个犹太人，名叫阿颇罗}等\BibleRef{宗 18:24}。因为他是一个觉醒的人，为这个目的周游各地。关于他，宗徒写道：\Quote{至于我们的弟兄阿颇罗}\BibleRef{格前 16:12}（\OriginalTerm{}{\GreekText{β}}）\Quote{阿桂拉和普黎史拉听了他的话}等\BibleRef{宗 18:26}。他留下他们在厄弗所不是无缘无故的，而是为了阿颇罗的缘故，圣神如此安排，使他能以更大的力量攻击（\OriginalTerm{}{\GreekText{ἑ}} \OriginalTerm{}{\GreekText{πιβἥναι}}）格林多。为什么他们对他没有采取行动，却攻击保禄？他们知道他是领袖，这个人的名声很大。\Quote{当他有意往阿哈雅去}\BibleRef{宗 18:27}，即凭信德，他一切都凭信德行事；\Quote{弟兄们就写信}等，全无嫉妒，全无恶眼。阿桂拉教导，或者不如说这个人让自己受教。他决定启程，他们就送信。\OriginalTerm{}{a}\Quote{因为他以很大的能力，公开驳倒犹太人}等\BibleRef{宗 18:28}。由此，他\Quote{公开}驳倒他们，显出了他的胆量；他论证的清晰，显出了他的能力；他从圣经中驳倒他们，显出了他的技巧（学识）。因为单有胆量而无能力，或单有能力而无胆量，都无济于事。\Quote{他以很大的能力驳倒}，经上说。（\OriginalTerm{}{\GreekText{β}}）\Quote{这事以后}等\BibleRef{宗 19:1}。但那些在厄弗所的人，怎么会有若翰的洗礼？可能他们在（若翰宣讲的）时候曾到耶路撒冷访问，甚至不认识耶稣。他没有对他们说：你们信耶稣吗？而是说：\Quote{你们领受了圣神没有？}\BibleRef{宗 19:2}他知道他们没有，但希望他们自己说出来，好让他们知道所缺少的，从而求问。\Quote{若翰固然用水施了洗}等\BibleRef{宗 19:4}。若翰从洗礼本身预言：他引导他们（看出）这就是若翰洗礼的意义。（a）\Quote{叫他们信那将要来的}：关于什么样的人？\Quote{我固然用水洗你们，但在我以后来的，要用圣神洗你们。}\BibleRef{玛 3:11}\Quote{保禄按手在他们头上，圣神便降在他们身上，他们就讲各种语言，也说预言。}\BibleRef{宗 19:6}（\OriginalTerm{}{\GreekText{β}}）这恩赐是双重的：语言和预言。由此显出一个重要的道理：若翰的洗礼是不完备的。他没有说\Quote{赦罪的洗礼}，而是说\Quote{悔改的洗礼}。那么，是什么情况呢？这些人没有圣神：他们不那样热心，甚至没有受教导。为什么阿颇罗没有受洗？在我看来，情况是这样的：这个人的胆量很大。\Quote{他殷勤教导有关耶稣的事}，但他需要更殷勤的教导。这样，虽然他并非完全认识，但因他的热忱，吸引了圣神，正如科尔乃略和他的同人一样。\par

或许许多人希望，啊，我们现在还有若翰的圣洗！但（即使有），许多人仍会疏忽于修德，并且可能会有人认为每个人是为了这个，而不是为了天国的缘故而追求德行。将有许多假先知：因为那时\BibleQuote{那些经得起考验的}不会十分\BibleQuote{显明。}\BibleRef{格前 11:19} 正如\BibleQuote{那些没有看见而相信的，是有福的}\BibleRef{若 20:29}，同样，那些没有看到神迹就相信的人也是如此。\BibleQuote{除非你们看见神迹，你们总不会相信。}基督说。\BibleRef{若 4:48} 因为我们没有因缺乏奇迹而损失什么，只要我们愿意留意自己。我们拥有美好事物的总和：通过圣洗，我们获得了罪过的赦免、圣化、圣神的参与、成为义子、永生。你们还想要什么？奇迹吗？但奇迹会终止（\OriginalTerm{但被废除}{\GreekText{ἀλλὰ} \GreekText{καταργεἵται}}）。你们有\Quote{信德、望德、爱德}这些常存之物；你们该追求这些，这些比奇迹更伟大。没有什么能与爱德相等。因为他说：\BibleQuote{最大的，是爱德。}\BibleRef{格前 13:5} 但现在，爱德处于危险中，因为只有名字留下来，而实际却无处可见，我们彼此分裂。那么，一个人该做什么来使我们重新合一呢？批评容易，但如何建立友谊，这才是需要研究的要点；我们如何聚集分散的肢体。即使我们有一个教会，或一个教义——但这还不是主要考虑；不，邪恶在于，我们在这些事上没有共融——正如宗徒所说，\BibleQuote{与众人平安相处}\BibleRef{罗 12:18}，相反，我们彼此不和。即使我们没有天天争斗，但你也不要只看这一点，而是要看这一点：我们也没有真诚而坚定不移的爱德。需要绷带和油。让我们牢记，爱德是基督门徒的标志；没有它，其他一切毫无用处；如果我们愿意，这是容易的事。是的，你们说，我们都知道这些，但如何着手才能实现呢？做什么，才能实现？怎样做，我们才能彼此相爱？首先，让我们除掉那些破坏爱德的事，然后我们才能建立它。不要有人心存怨恨，不要有人嫉妒，不要有人幸灾乐祸：这些是阻碍爱的事；那么，建立爱的事则是另一种。因为仅仅除掉阻碍是不够的；还必须提供建立的事。如今息辣告诉我们那些破坏（友谊）的事，却没有继续讲建立合一的事。他说：\BibleQuote{责骂、泄露秘密和阴险的伤害。}\BibleRef{德 22:27} 但说到那些时代的人，这些事也许可以提及，因为他们属于血肉之人；但在我们当中，但愿这些事连提都不提。我们不是从这些事中向你们提出劝诫，而是从另外的事。对我们来说，没有友谊就没有好事。即使有无数的好事，但如果没有友谊——无论是财富还是享乐——有什么益处呢？没有任何东西比得上友谊，即使在今世的事务中也是如此，正如没有什么比人们仇恨我们更糟糕一样。\BibleQuote{爱德遮盖许多罪过。}\BibleRef{伯前 4:8} 但仇恨，即使没有罪过，也会怀疑它们存在。不作敌人是不够的，还必须去爱。要想到，基督命令了，这就够了。甚至苦难也能缔造友谊，把人团结在一起。\Quote{那么，}你们说，\Quote{现在没有苦难时呢？请说，我们如何行动才能成为朋友？}我问你们，难道没有别的朋友吗？你们是怎样成为他们的朋友的？你们如何保持如此？首先，让人不要有敌人：这件事本身并非小事；让人不要嫉妒；不嫉妒的人是不可能被指责的。（\Emph{b}）那么，我们怎样才能感情热烈？什么使人爱慕？是人的容貌。那么，也让我们使自己的灵魂美丽，我们就彼此可爱了：因为当然，不仅需要去爱，也需要被爱。让我们先实现这一点，即我们被爱，另一点就容易了。怎么做才能被爱？让我们变得美丽，并这样做，使我们总有爱慕者。让人不以赚钱、得奴隶、置房产为务，而以被爱、有好名声为务。名声胜过大量财富。因为名声留存，财富消逝；名声可得，财富难求。因为得了坏名声的人很难摆脱它；但借着他的好名声，穷人可能很快致富。让一个人有一万塔冷通，另一个人有一百个朋友；后者比前者更富有资源。那么，我们不仅要做这事，还要以此为一种营生。\Quote{我们怎能做到呢？}你们说。\Quote{甘甜的口舌增多朋友，优雅的言语增添恩惠。}让我们有善于言辞的口和纯洁的品行。一个人不可能如此而不被人所知。\par

(\Emph{a}) 我们都居住在一个世界，以同样的果实为食。但这些是小事：藉着同样的圣事，我们分享同样的神粮。这确实是爱的理由！(\Emph{c}) 请留意，那些外人为了友谊想出了多少（诱因与诉求）：技艺或行业的共同点、邻里关系、亲属关系：但比这一切更有力的，是存在于我们之间的动力与纽带：这圣筵比（其他一切）更能促使我们羞愧地走向友善。但我们中许多人前来参加，却互不相识。也许有人说，这是因为人数太多。绝非如此；这只是我们自己的懒惰冷漠。（那时）有三千人 \BibleRef{宗 2:41}——有五千人 \BibleRef{宗 4:4}——他们却都一心一意：但现在各人不认识自己的兄弟，并且不羞于将责任推给人数，因为人数太多了！然而，拥有许多朋友的人对所有人都是不可战胜的：他比任何暴君更强。暴君有卫兵的保护，不如这个人有朋友的安全。此外，这个人比暴君更荣耀：因为暴君被自己的奴隶守护，而这个人被他的同伴守护；暴君是被不情愿且畏惧他的人守护，而这个人是被情愿且无所畏惧的人守护。在这里也可以看到一件奇妙的事——多在一中，一在多中。(\Emph{a}) 如同在一架竖琴上，声音不同，但和谐一致，它们共同奏出同一和谐与交响，(\Emph{c}) 我希望能带你们进入这样一座城，如果可能的话，在那里（所有人）应当同心合意：那时你们将看到超越一切竖琴与长笛和谐的更和谐的交响。(\Emph{b}) 但音乐家是爱的力量：正是它奏出甜美的旋律，(\Emph{d}) 唱着（同时）一首无一音符走调的曲调。这首曲调使天使和天使之主天主欢喜；这首曲调唤醒（听它）整个天上的听众；这首曲调甚至平息（邪恶）的激情——它甚至不容它们升起，而是深深的寂静。因为如在剧院中，当乐队演奏时，所有人都静静地听，没有噪音；同样在朋友中，当爱拨动琴弦时，所有激情都静止安睡，如同被驯服和瘫痪的野兽：正如哪里有仇恨，哪里就与此完全相反。但此刻让我们暂且不提敌意；让我们谈论友谊。即使你脱口而出一些匆忙的话，没有人会指责你，反而都原谅你；即使你做了（一些仓促的事），没有人会做最坏的解读，反而都给予谅解：每个人都乐于向跌倒者伸出援手，每个人都希望他站立。友谊确是一道坚不可摧的墙；一道连魔鬼自己都不能压倒的墙，更不用说人了。拥有许多朋友的人不可能陷入危险。（哪里有爱）就没有理由生气，只有愉悦的感觉；没有理由嫉妒；没有理由怨恨。看这个人，在属灵和属世的一切事物中，他都轻松地完成一切。那么，请问，谁能与这个人相比？这个人如同四面城墙环绕的城市，另一个人如同没有城墙的城市。——伟大的智慧，能够成为友谊的创造者！除去友谊，你就除去了所有，你就搅乱了所有。但如果友谊的相似有如此大的力量，那么现实本身又该是怎样的呢？那么，我恳求你们，让我们结交朋友，让每个人都以此为技艺。但是，瞧！你会说，我确实在努力，但别人没有。这对你来说是更大的赏报。诚然，你说，但这事更难。怎么难呢？我请问？看！我向你们作证并宣告，如果你们中有十人彼此联合，以此为业，如同宗徒以宣讲为业，先知以教导为业，那么我们将以结交朋友为业，赏报将是大的。让我们为自己制作君王的肖像。因为如果这是门徒的共同标记，我们就做了比能复活死人更伟大的工作。王冠和紫袍标记皇帝，若没有这些，即使他的衣服全是金的，皇帝也尚未显现。现在你们正在显明你们的族系。使别人成为你们的朋友，并使（他们成为）别人的朋友。没有人会因为被爱而愿意恨你。让我们学习色彩，它们是用什么成分调和的，这幅肖像是用什么（色调）构成的。让我们和蔼可亲：不要等我们的邻人先行动。不要说，如果我看到任何人退缩（不先主动），我就变得比他更糟：相反，当你看到这一点时，要抢先一步，消除他的不良情绪。你看见一个人患病，还要加重他的病情吗？这一点，最要紧的，我们要确保——\Quote{在尊敬上彼此争先，看别人比自己优越} \BibleRef{罗 12:10}，不要以为这会贬低你自己。如果你在尊荣上（让别人）占先，你就更尊荣了自己，吸引更高的荣誉到你自己身上。在所有场合，让我们把优先权让给别人。让我们心中不记念别人对我们做的恶，但如果行了什么善（我们只记住这些）。没有什么比和蔼的舌头、说好话的口、不自高自大的灵魂、轻视虚荣、藐视荣誉更能使人成为朋友。如果我们确保这些，我们就能对魔鬼的诡计不可战胜，并且严格修德，达到爱祂之人所应许的美善，藉着我们的主耶稣基督的恩宠和仁慈，祂和圣父及圣神一起，荣尊、主权、荣耀，现在以至永远，世世无穷。阿们。\par

\SourceRecord{Translated by J. Walker, J. Sheppard and H. Browne, and revised by George B. Stevens.；Nicene and Post-Nicene Fathers, First Series；Vol. 11.；Edited by Philip Schaff.；Buffalo, NY: Christian Literature Publishing Co.,；1889.；<http://www.newadvent.org/fathers/210140.htm>.}

% structure: p0001 source=h1 target=section reason=page-title

\section{论宗徒大事录 第41讲}

\BibleRef{宗 19:8-9}\par

\BibleQuote{他进入会堂，放胆讲论了三个月，辩论并劝服关于天主国的事。但有些人顽固不信，在群众面前毁谤道路，他就离开他们，将门徒分别出来，每天在\Person{提郎诺}的学堂里辩论。}\par

(\Emph{a}) 看他在各处强行进入会堂，并照此方式离开。因为每处，他都希望由他们（犹太人）给他机会。(\Emph{c}) 他希望由此将门徒分别出来，并且停止与他们（犹太人）聚会的开端由他们自己提出。他这样做并非没有原因，(\Emph{b}) 如同我刚才所说。从此以后，他\Quote{激起了他们的嫉妒}。因为外邦人乐意接受他，而犹太人见外邦人接受他便后悔了。(\Emph{a}) 这就是为什么他持续在他们中引起骚动，\Quote{三个月辩论并劝服关于天主国的事}；因为你不能因听到他\Quote{放胆讲论}就以为有何粗硬之处；他讲论的是好事，关于一个王国：有谁不愿听他呢？\Quote{但有些人顽固，毁谤道路。}他们称之为\Quote{道路}是恰当的；这确实是\OriginalTerm{道路}{the way}，引向天国。\BibleQuote{他离开了他们，将门徒分别出来，每天在\Person{提郎诺}的学堂里辩论。这样做了两年之久，以致凡住在\Place{亚细亚}的，都听见了主的圣言，犹太人和希腊人。} \BibleRef{宗 19:10} (\Emph{a}) 你注意到他的坚持带来了多大的效果吗？\Quote{犹太人和希腊人都听见了，(\Emph{c}) 凡住在\Place{亚细亚}的。}也正是为此，主以前不许他进入\Place{亚细亚} \BibleRef{宗 16:6}（以前）；等待的，据我看来，正是这个时机。（讲道集40，第245页）(\Emph{b}) \BibleQuote{天主借保禄的手行了异乎寻常的奇迹；甚至有人从他身上拿手巾或围裙放在病人身上，疾病便离开了他们，邪魔也出去了。}（11-12节）不仅触摸者（得到治愈），而且接受这些物品后，放在病人身上（也治愈了他们）。(\Emph{g}) \BibleQuote{信我的，} 基督说，\BibleQuote{要做比我所做的更大的事。} \BibleRef{若 14:12} 这就是他（在此话中）所指的，以及影子的奇迹。(\Emph{d}) \BibleQuote{那时，有几个流浪的犹太人，驱魔者，擅自向附有邪魔的人呼号主耶稣的名字，说：\InnerQuote{我因保禄所宣讲的耶稣，命令你们出来。}} \BibleRef{宗 19:13} 他们行事完全是出于交易！注意：流浪的，或游走的，犹太驱魔者。至于相信，他们毫无心意；但借那名号他们渴望驱魔。\Quote{因保禄所宣讲的耶稣。}且看保禄有了何等的名望！\BibleQuote{有七个儿子，是\Person{斯盖瓦}，一个犹太人，大司祭，也照样做了。邪魔回答他们说：\InnerQuote{耶稣我认识，保禄我也知道；但你们是谁？}那人身上有邪魔的，便跳在他们身上，制伏了他们，胜过了他们，以致他们赤身受伤地从那房子逃出去了。} \BibleRef{宗 19:14} 他们秘密行事，结果他们的无能当众暴露。(\Emph{f}) 可见，除非凭信德呼号，那名号并不能做什么。(\Emph{h}) 看他们如何用武器攻击自己！(\Emph{j}) 他们远未将耶稣视为伟大；不，他们必须加上保禄，认为他是大人物。这里令人惊奇的是，为何邪魔不与驱魔者的骗局合作，反而揭露他们，戳穿他们的把戏。在我看来，他（邪魔）是出于极怒：好比一个人处于绝境，被一个可怜虫出卖，便想将满腔怒火发泄在他身上。\Quote{耶稣我认识，保禄我也知道。}为了不让耶稣的名号受任何轻慢，（邪魔）首先承认（他），然后才被赋予许可。因为，为显示不是名号的软弱，而是完全由于那些人的欺诈，为何在保禄身上不发生同样的事？\Quote{他们赤身受伤地从那房子逃出去了。}他狠狠地打伤了他们的头，也许撕破了他们的衣服。(\Emph{e}) \BibleQuote{这事传遍了所有住在\Place{厄弗所}的犹太人、希腊人；众人都生敬畏，主耶稣的名号被尊大。许多信了的人，都来承认并述说自己所行的事。}（17-18节）既然他们借邪魔得了行这样事的权能，这样的结果理所当然。\BibleQuote{许多行巫术的人，把书集来，当众焚毁}——他们看见既然邪魔自己行这些事，这些书再无用处——\BibleQuote{并计算了书价，发现共值五万银币。天主的圣言如此强盛而得胜。}（19-20节）(\Emph{i}) \BibleQuote{他}（这样）\BibleQuote{辩论}，在\Person{提郎诺}的学堂里两年之久；那里有信徒，而且是极其（在信德上）长进的信徒。此外，（保禄）写信（给他们）如同写给大人一样。\par

（重述。）（\Emph{b}）\Quote{进入了会堂}等等。\BibleRef{宗 19:8}但为何\Quote{\OriginalTerm{放胆}{\GreekText{ἐπαρρησιάζετο}}}？这是说，他准备好面对危险，更公开地辩论，不隐藏教义。（\Emph{a}）\Quote{但有些人顽固不信，并在大众面前毁谤这道，保禄便离开他们，将门徒分别出来。}\BibleRef{宗 19:9}这是说，他阻止了他们的诽谤：他不愿激起他们的嫉妒，也不愿使他们陷入更多的争执。（\Emph{c}）由此我们也要学习，不要将自己置于诽谤者的路上，而要离开他们：他自己被诽谤时，并不以恶言相报。\Quote{他天天辩论}，借此赢得了许多人；尽管受到恶待和毁谤，他并没有完全断绝与他们的往来而疏远他们。（\Emph{e}）诽谤者被击败了。他们污蔑了教义本身；（因此）为了既不激起门徒的愤怒，也不* *他们，他退开了，表明他们在各处都同样拒绝了救恩。此时他甚至连辩护都不做了，因为外邦人在各处都已相信。\Quote{在提郎诺的学房：}并非他寻找这个地方，而是不加犹豫，哪里有学房，他就在那里讲论。（\Emph{d}）看，外来的试炼一结束，来自魔鬼的试炼就开始了。请注意犹太人的执迷不悟。他们看见他的衣服行奇迹，却毫不在意。还有什么比这更甚呢？但相反，它导致了恰恰相反的效果。倘若外邦人中有人不信，看见尘土本身产生这些效果，就让他信吧。（\Emph{f}）奇妙，信者之能力何等伟大！西满为了牟利而寻求圣神的恩宠，这些人为同样的目的也这样做。何等的顽梗！为什么保禄不斥责他们？那样会显得嫉妒，所以如此安排。同样的事也发生在基督身上\BibleRef{谷 9:36}：但那时那人没有被阻止，因为那是新秩序的起始：就如犹达斯也没有被阻止，而阿纳尼雅和撒斐辣却被击毙；许多犹太人甚至因反对基督也没受什么惩罚，而厄吕玛却瞎了眼。\Quote{因为我来，}基督说，\Quote{不是为审判世界，而是为叫世界得救。}\BibleRef{若 3:17}\Quote{有七个儿子}等等。\BibleRef{宗 19:14}看这些人的邪恶！他们仍是犹太人，却想利用那名字牟利。他们所行的一切都是为了光荣和利益。（\Emph{g}）看，在每一情形中，人们与其说因好事而皈依，不如说因可怕之事而皈依。在撒斐辣的事上，恐惧降临在教会，人不敢依附他们；这里他们接受手巾和围裙，得了医治；此后，他们才来承认自己的罪。（由此）魔鬼的能力在对付不信者时显得很大。因为为何他不说：\Quote{耶稣是谁？}他害怕，唯恐自己也受惩罚；但为了能对嘲弄他的人复仇，他这样做了：\Quote{耶稣，}他说，\Quote{我认识，}等等。他惧怕保禄。为何那些可怜的人不对他说：我们信？如果他们这样说，即如果他们认祂为主，他们会显得多么光辉啊？但相反，他们甚至说出了那些无意义的话：\Quote{指着保禄所宣讲的耶稣。}你注意到作者的宽容吗？他如何写史而不点名？这使得宗徒们可敬。\Quote{邪魔，}等等。\BibleRef{宗 19:15}，因为在斐理伯发生的事\BibleRef{宗 16:16}也已给了这些人一个教训。他提到名字和数目，以此给当时的人他写作的可信证据。为何他们巡回各处？为了牟利：当然不是为了传报圣言，这怎能是他们的目的？他们如何立刻奔跑，藉所受的苦宣讲？\Quote{甚至，}它说，\Quote{凡住在亚细亚的人都听见了主的话。}这岂不该使众人皈依？不必惊奇，因为没有什么能说服恶意。但来吧，让我们看看那些驱魔者的事，他们怀着何等邪恶的心意。为何在基督身上没有发生同样的事，这是另一个时候的探讨，不是现在，只除了这也安排得妥善有益。依我看来，他们这样做也是出于嘲弄，并且由于这（惩罚），无人敢随意称呼那名字。为何这促使他们告解？因为这是天主全知的最有力的论证，（因此）在魔鬼揭露他们之前，他们控告自己，害怕遭受同样的事。因为当他们的帮手魔鬼成为控告者时，除了以行为告解，还有什么希望呢？\par

但是请看我恳求你们，在行了如此多的神迹之后，短短时间内便发生了怎样的邪恶。人性就是这样：很快就忘记了。难道你们不记得我们中间所发生的事吗？去年天主不是震动了我们整座城吗？不是所有的人都跑去领洗吗？那些嫖客、柔弱的人和腐败的人不是离开了他们的住所和消磨时间的地方，改变而成为虔诚的人吗？但三天过去，他们又回到自己先前的邪恶之中。这是从哪里来的？来自极度的懒惰。如果当事情过去后（是这样），况且这些形象永久存留，结果尚且如此，这有什么稀奇呢？索多玛的命运——说，难道它（的后果）不是仍然存留吗？那么，住在它旁边的居民有没有变得更好呢？诺厄的儿子你们怎么说？他不就是那样（如所描述的）吗？他难道没有亲眼目睹如此巨大的毁灭，却仍然作恶吗？那么，当我们看见当这些事发生后，这些（厄弗所的）犹太人并不相信时，让我们不要惊奇，既然我们看到信德本身常常为他们转变为其反面，成为恶毒；例如，当他们说祂有魔鬼时，祂，天主子！你们难道现在不看见这些事，而且许多人像蛇一样，既无信又忘恩负义，这些人，像毒蛇一样，当享受了恩惠并被某些人温暖后，他们就刺伤他们的恩人？我们说这些，是怕有人惊奇：既然行了如此大的神迹，怎么不是所有人都归化了。因为看，在我们自己的时代，发生了那些有关殉道者巴比拉的事，那些有关耶路撒冷的事，那些有关圣殿被毁的事，但并非所有人都归化了。我何必说古老的事呢？我已经告诉你们去年发生的事，但没有人留意，反而渐渐坠落回原状。天一直高声宣告它有主人，我们所见的一切——我是指世界——都是工匠的作品，然而有些人却说并非如此。去年发生在德奥多禄身上的事，谁没有震惊呢？然而没有结果，只是他们一时变得虔诚，却又回到他们开始试图虔诚时的起点。犹太人也是如此。这就是先知论及他们所说的：\Quote{上主击杀他们时，他们就寻求他，回心转意，恳切寻觅天主。}\BibleRef{咏 77:34} 何必说那些对所有人都平常的事呢？多少人陷入疾病，多少人曾许诺如果被治愈就要做出如此大的改变，然而他们又变得和从前一样！这一点，如果没有别的，就表明我们拥有本性的自由意志——忽然间改变。如果邪恶是本性的，就不会如此：本性和必然的事，我们无法改变。\Quote{然而，} 你会说，\Quote{我们却会改变它们。因为我们岂不是看见有些人，他们有自然的视能，却因恐惧而盲目？}（是的——）因为这也是本性的：* * 如果没有另一种（必然的）本性介入的话；（因此）我们因恐惧而看不见是本性的；当更大的恐惧出现时，另一个便让位也是本性的。\Quote{那么，} 你会说，\Quote{如果正直之心确实是出于本性，但恐惧占了上风，就把它赶走了？} 那么，如果我将证明有些人甚至在那些恐惧中也不被带到正直之心，而是在这些恐惧中仍鲁莽行事呢？这是本性的吗？我要说古代的事吗？那么，最近的呢？多少人在那些恐惧中继续大笑、嘲笑，丝毫没有经历这类事？法老不是立即改变，又（迅速）回到他先前的邪恶吗？但在这里，好像（魔鬼）不认识祂，（驱魔者）反而加上说，\Quote{保罗所宣讲的，} 而他们本该说，\Quote{世界的救主，} \Quote{那复活的一位。} 由此他们表明他们确实知道，但他们不愿承认祂的光荣。因此魔鬼揭露了他们，跳到他们身上，说：\Quote{耶稣我认识，保罗我也认识。但你们是谁？} 所以你们不是信徒，而是你们说这话时滥用那名字。因此圣殿是荒凉的，工具容易克服。所以你们不是宣讲者；你们是我的，他说。魔鬼极其愤怒。宗徒们有能力对他们这样做，但他们还没有做。因为那有能力对魔鬼做这些事的人，对这些人本身更有能力。注意他们的容忍是如何显示出来的：他们被斥退的人做这些事，而他们所追求的魔鬼却做相反的事。\Quote{耶稣，} 他说，\Quote{我认识。} 羞愧吧，你们这些不认识（祂）的人。\Quote{保罗我也认识。} 说得好，\Quote{不要以为我因鄙视他们而做这些事。} 魔鬼的恐惧极大。为什么他不用这些话语而撕裂他们的衣服？因为那样他既可以发泄他的愤怒，又可以坚固欺骗。如我所说，他害怕那不可接近的力量，如果他没说这话，就不会有这种力量。但是请注意我们如何发现魔鬼处处（比犹太人）更正直，不敢反驳或控告宗徒或基督。在那里他们说：\Quote{我们认识你是谁}\BibleRef{玛 8:29}；\Quote{你为什么在时候未到前来折磨我们}\BibleRef{谷 1:24}；以及又，\Quote{我认识你是谁，你是天主子。} 而在这里，\Quote{这些人是至高天主的仆人}\BibleRef{宗 16:17}；以及又，\Quote{耶稣我认识，保罗我也认识。} 因为他们极其害怕并战栗于那些圣者。也许你们中有人听到这些事，希望自己拥有这种能力，使魔鬼不敢正视他，并认为这些圣人是幸福的，因为他们有这种能力。但要让他听基督说：\Quote{不要因为魔鬼屈服于你们而欢喜}\BibleRef{路 10:20}，因为祂知道所有人最欢喜于此，出于虚荣。因为如果你寻求那令天主喜悦的，和那对公众有益的，还有另一条更大的路。从魔鬼手中释放一个人，远不如从罪中拯救一个人伟大。魔鬼不会阻碍人进入天国，反而会合作，虽然不情愿，但依然合作，使拥有魔鬼的人更加清醒；但罪却将人赶出去。\par

但很可能有些人会说：\Quote{但愿天主永不让我以这种方式清醒！}我也不希望你这样，而是希望你以完全不同的方式，即因对基督的爱而行事；然而，如果，愿天主不许，你竟如此遭遇，那么，即使为此，我也要安慰你。所以，如果魔鬼不能使人（从天国）逐出，而是罪使人逐出，那么使人脱离罪便是更大的恩惠。\par

至此我们应学习解救我们的邻人，而且先于邻人，解救我们自己。我们要注意，免得我们也有魔鬼；我们要严格省察自己。罪比魔鬼更严重，因为魔鬼使人谦卑。你们没有看见那些附魔的人，在发作过后怎样垂头丧气、面带愁容、满含羞愧，甚至不敢抬头看人吗？多么奇怪的矛盾！那些人因自己所受的苦而羞愧，我们却因自己所行的事而不羞愧；他们因受害而惶恐，我们却因作恶而不惶恐。其实他们的状况并不值得羞耻，反而应得怜悯、温柔和宽容；甚至应得极大的钦佩和许多赞美，因为他们与那样的灵争战，却以感恩的心忍受一切。而我们的状况才真正是笑柄、羞耻、指责、纠正、惩罚、最坏的恶事、地狱之火，毫无怜悯可言。你们看见了吗，罪比魔鬼更坏？那些人从所受的苦中获得了双重的益处：第一，他们清醒过来，更加自制；第二，他们在此受了自己罪恶的惩罚，然后洁净地离开到他们的主人那里。因为关于这一点，我们常向你们讲过：那些在此受罚的人，如果感恩地忍受，很可能借此消除许多罪恶。而罪所带来的祸害是双重的：第一，我们冒犯了（天主）；第二，我们变得更坏。注意我所说的。罪带给我们的伤害不仅仅在于我们犯了罪；还有另一个更坏的：我们的灵魂养成了习惯。正如身体的情况——用例子说明会更清楚——就像一个人发了烧，不仅生病了，而且在病后变得更虚弱，即使久病之后恢复健康也是如此。同样，在罪中，即使我们可能恢复健康，但我们远远没有所需的坚强。比如一个粗暴辱骂的人：他是否因辱骂行为而受应得的报应？是的，但还有另一个更值得遗憾的坏事：他的灵魂变得更加不知羞耻。因为每犯一个罪，即使在罪已成事并停止之后，仍有一种毒液注入我们的灵魂。你们没有听见人们病愈后说：\Quote{我现在不敢喝水了}吗？那人确实恢复了健康；但疾病也给他造成了这种伤害。而那些（附魔）的人，虽然受苦，却感恩；我们呢，处境顺利，却亵渎天主，认为自己受了很大的亏待。你们会发现，健康富足的人比贫穷患病的人更常这样。因为魔鬼站在（附魔者）身旁，像一个凶残的刽子手，发出许多（威胁），像一个挥舞鞭子的教师，不让他们有任何松懈。即使有些人完全没有清醒，他们也不受惩罚；这可不是小事。正如愚人、疯子和儿童不被问责，这些人也不被问责；因为对于无意识状态下所做的事，没有人会残忍到去问责行事者。那么，我们这些犯罪的人，处境远不如那些有邪灵附身的人。我们确实不口吐白沫，不扭曲眼睛，不抽搐地挥舞双手；但这一点，但愿它发生在身体里而不是灵魂上！你们要我展示一个灵魂在冒泡沫、污秽，以及内心眼睛的扭曲吗？想想那些发怒、醉于愤怒的人；有比他们吐出的话语更污秽的形象吗？实在像喷出恶臭的唾液。正如附魔者不认识在场的人，这些人也不认识。他们的理智昏暗，眼睛歪斜，分不清谁是朋友、谁是敌人、谁值得尊敬、谁可鄙，而是毫无区别地看待所有人。那么，你们没有看见他们颤抖，就像那些人一样吗？但他们不倒地，你们说？不错，但他们的灵魂倒在地上，在那里抽搐；因为如果灵魂站立得稳，就不会陷入这种状况。或者，你们不认为，一个人醉于愤怒时所做和所说的一切，正表明灵魂卑贱地匍匐在地、完全失去自制吗？还有另一种比这更糟的疯狂。是什么呢？就是有些人甚至不允许自己发泄愤怒，反而在自己心里滋养着，好像一个拿着鞭子的刽子手，对他自己造成伤害，那就是记恨。因为他们怀有的恶意首先伤害他们自己。且不说将来的事，你们想，那个人在灵魂的鞭笞中每天想着如何向敌人报仇，他会忍受怎样的折磨？他先惩罚自己，受苦，膨胀着（压抑的激情），与自己争战，点燃自己。因为那火必须总是在你里面燃烧；你把热度升到如此之高，不让它减退，却以为自己是在伤害对方，其实你是在消耗自己，永远带着一个总是高度燃烧的火焰，不让你的灵魂有休息，而是永远处于狂怒之中，思想在纷乱和风暴中。还有什么比这种疯狂更痛苦：总是刺痛，总是肿胀发炎？因为记恨者的灵魂就是这样：当他们看到他们想报复的人时，立刻像被打了一拳；如果听见那人的声音，他们就畏缩颤抖；如果他们躺在床上，就会想象无数报复，绞死、折磨那个敌人；如果除此之外，他们还看到那人也出名，啊，他们所受的痛苦！宽恕他的冒犯，把自己从折磨中解放出来。为什么一直处于惩罚之中，只为了一次惩罚他、报复他？为什么要为自己建立一种消耗性的疾病？为什么当你的愤怒想要离开你时，你却把它留住？\BibleQuote{不要让它留到日落}，保禄说。\BibleRef{弗 4:26} 因为就像某种腐蚀或蛀虫一样，它甚至咬穿我们理解的根基。为什么把一头野兽关在你的肚腹里？让一条蛇或毒蛇躺在你的心中，也比愤怒和怨恨好：因为那些（蛇）很快会结束我们，而愤怒和怨恨永远停留，把毒牙钉入我们，注入毒液，释放一群痛苦的思想来攻击我们。你说：\Quote{他会嘲笑我，}你说，\Quote{他会轻视我！}可怜可悲的人，你宁愿被同仆嘲笑，而被主人憎恨吗？你宁愿被同仆轻视，而轻视你的主人吗？\par

被祂鄙视，你真受不了吗？但你难道不认为天主会因此愤怒，因为你嘲笑祂，因为你不遵行祂的命令，你就轻视了祂？至于你的敌人，他甚至不会嘲笑你，从这件事可以明显看出：你若追求报复，那是极大的嘲笑，极大的鄙视，因为这是心胸狭隘的记号；相反，你若宽恕他，则是极大的钦佩，因为这是灵魂伟大的记号。但你会说，他不知道这一点。让天主知道吧，好使你获得更大的赏报。因为祂说：\Quote{借给人，不指望偿还。}\BibleRef{路6:34}我们也应当这样对那些人行善，他们甚至不晓得有人正在为他们行善，免得他们因向我们回报赞美或任何其他的事，而减少我们的赏报。因为当我们从人那里一无所获时，我们就会从天主那里获得更大的赏报。还有什么比一个常发怒、总想报复的灵魂更可笑、更卑劣的呢？这种性情是懦弱的，是幼稚的。正如幼儿会对无生命的物体发怒，除非母亲拍打地面，他们才肯息怒；同样，这些人也总想报复那些得罪他们的人。由此可见，他们才是可笑的人；因为被情欲征服是幼稚理智的标志，而战胜情欲则是刚毅的象征。那么，当我们保持冷静时，并非我们，而是他们成了被嘲笑的对象。使人可鄙的不是被情欲压倒；真正可鄙的是，因害怕外界的嘲笑，以至于选择屈服于自己缠扰的情欲、得罪天主、并对自己进行报复。这些事确实可笑。让我们远离它们。假若有人说，那人害了我们无数恶事，却未受任何报应；让他说，这样他就可以再次疯狂攻击我们，而无所畏惧。其实，没有更好的方式能宣扬我们的美德；若他想要称赞我们，他也不会找到比那些看似辱骂的话更好的言词。但愿众人都这样议论我：\Quote{他是个温顺的可怜虫；人人都侮辱他，但他都忍受了；人人都践踏他，他却不为己复仇。}但愿他们再加上：\Quote{就算他想报复，他也做不到。}这样我就能得自天主的赞美，而非来自人。让他说，我们不为自己报仇是因为缺乏骨气。当天主洞悉一切时，这不会伤害我们；这只会使我们的宝藏更加安全。如果我们还要顾及他们的说法，我们就将失去一切。不要看他们说什么，而要看什么适合我们。可是，有人却说：\Quote{别让人嘲笑我。}有些人还以此自豪。哦，何等的愚昧！\Quote{没有人在得罪我之后嘲笑过我。}也就是说：\Quote{我已经报复了。}然而，正因你报复，你才该被嘲笑。这些话语是怎么在我们中间流传的？它们本是我们的耻辱和祸患，是对我们自身生活和纪律的颠覆！你这样说，简直是公然对抗天主。正是使你与天主相似的事——不报复自己——你竟认为这是可笑的！我们这样与天主作对，难道不该被自己人和外邦人嘲笑吗？我想讲一个古时发生的故事，不是关于忿怒，而是关于钱财。一个人有一块地，地里藏有宝藏，而主人不知情。他将这块地卖掉了。买主在挖掘栽种时发现了埋藏的宝藏，便来找卖主，想强迫他收回宝藏，声称自己买的是地，不是宝藏。卖主拒绝接受赠予，说：\Quote{地已归你，我已卖掉，与这宝藏毫无关系。}于是他们为此争吵起来，一人要赠送，一人坚持拒绝。他们遇到第三人，请他裁决，对他说：\Quote{这宝藏该归谁？}那人无法断定，但说可以解决他们的争执——他愿意（如果他们同意）自己成为宝藏的主人。于是他收下了宝藏，他们俩都心甘情愿地放弃了；结果后来他陷入了无尽的麻烦，并亲身体会到他们俩当初不沾手是明智的。对待忿怒也应如此：我们自己应当不追求报复，而那些得罪我们的人也应当追求补偿。但也许这些事也被视为可笑；因为当这种疯狂在人间广泛传播时，保持冷静的人反被嘲笑；在许多疯人中，不疯的人看上去反而像疯子。因此，我恳求你们，让我们从这疾病中痊愈，恢复理智，好摆脱这致命的激情，得以进入天国，赖祂独生子的恩宠和仁慈；愿光荣、权能、尊崇归于圣父，及圣子及圣神，从今世，及世世，以至于万世。阿们。\par

\SourceRecord{Translated by J. Walker, J. Sheppard and H. Browne, and revised by George B. Stevens.；Nicene and Post-Nicene Fathers, First Series；Vol. 11.；Edited by Philip Schaff.；Buffalo, NY: Christian Literature Publishing Co.,；1889.；<http://www.newadvent.org/fathers/210141.htm>.}

% structure: p0001 source=h1 target=section reason=page-title

\section{论宗徒大事录·第四十二篇讲道}

\BibleRef{宗 19:21,23}\par

\Quote{这些事结束后，保禄在圣神内决意，经过\Place{马其顿}和\Place{阿哈雅}后，前往耶路撒冷，说：\InnerQuote{我在那里以后，也必须看看罗马。}于是，他打发两位服事他的人，\Person{弟茂德}和\Person{厄辣斯托}，往\Place{马其顿}去；自己却暂留在\Place{亚细亚}。那时，关于这道起了不小的骚动。}\par

他派遣弟茂德和厄辣斯托往马其顿去，自己却留在厄弗所。在那城住了足够久之后，他希望再往别处去。但怎会一开始就选择往叙利亚去，现在又转回马其顿呢？经上说：\Quote{他有意}，\Quote{在圣神内}，表明一切（他所做的）并非靠自己的能力而行。现在他预言说，\Quote{我也必须看到罗马}：也许是为了安慰他们，因他并非远离他们，而是再次接近他们，并通过预言激发门徒的心灵。我认为此时正是他从厄弗所给格林多人写信的时候，说：\BibleQuote{我不愿意你们不知道，我们在亚细亚所遭遇的苦难。}\BibleRef{格后 1:8}。因为既然他已应许前往格林多，他为自己耽搁而辩解，并提及有关德默特琉事件的试炼。\Quote{关于这道，起了不小的骚动。}你看见声名（所获）了吗？经上说，他们反对；（然后）出现双重奇迹；（然后）再次危险：整个（历史）的脉络就是这样交替编织。\Quote{有一个银匠，名叫德默特琉，原是制造阿尔特米银龛的，使工匠们获利不少。}\BibleRef{宗 19:24}。\Quote{原是制造}，经上说，\Quote{阿尔特米的银龛。}银龛怎么可能做出来呢？也许是小盒子\OriginalTerm{小盒状物（可能指神龛的小模型）}{\GreekText{κιβώρια}}。在厄弗所，这（阿尔特米）受尊崇；因为当（\WorkTitle{厄弗所书讲道集序言}）他们的圣殿被烧时，他们非常悲伤，以至于禁止再提纵火者的名字。看，只要有偶像崇拜的地方，每次我们发现背后都是金钱。在前一例中是为了钱，在这个人的例子中，也是为钱。\BibleRef{宗 19:13}这不是因为他们的宗教，因为他们认为那处于危险中；不，是为了他们有利可图的手艺，怕它没有东西可做。要留意此人的恶意。他本人富有，对他而言损失不大；但对他们损失很大，因为他们贫穷，靠每日收入维持。然而，这些人什么也没说，只有他一个人说。注意：\Quote{他集合了那些同业工人，}与他们有共同事业，\Quote{说，诸位，你们知道，我们是从这手艺发财的}\BibleRef{宗 19:25}；然后他把危险带给他们，说我们正面临因这手艺而陷入饥饿的危险。\Quote{你们也看见也听见，不但在厄弗所，而且几乎在全亚细亚，这个保禄说服并引走了许多人，说，人手所做的不是神。这样不但我们这手艺有被藐视的危险，而且连伟大女神阿尔特米的圣殿也要被轻视，并且她的威仪也要被毁灭，全亚细亚和普世所敬拜的。众人听了这话，就满心恼怒，喊着说，大哉！厄弗所的阿尔特米！}\BibleRef{宗 19:26}。然而，他所说的话本足以带领他们皈依真宗教；但因为他们愚昧，就这样行事。因为如果这（保禄）人足够有能力引走所有人，而神明的崇拜处于危险，那么应当思想，这个人的天主该有多大，祂必将更多赐予你们所害怕的那些东西。起初（开头）他通过说\Quote{这个保禄引走了许多人，说，人手所做的不是神}已经抓住了他们的心。看，外教人为什么如此愤慨；因为他曾说\Quote{人手所做的不是神}。通篇，他针对他们的手艺说话。然后，他把最令他们伤心的事随后引出来。但是，与其他神明，他说，我们无关，但\Quote{伟大女神阿尔特米的圣殿也有被毁灭的危险}。然后，免得他似乎是出于利益才这样说，看他添加了什么：\Quote{全普世所敬拜的。}注意他如何显示保禄的能力更大，证明所有（他们的神）都是可怜卑微的，因为一个区区凡人，到处被驱逐，一个普通制帐篷的，竟有如此大的能力。注意他们的敌人对宗徒们的见证，即他们推翻了他们的崇拜。在吕斯特辣，他们带来\BibleQuote{花环和牛}\BibleRef{宗 14:13}。在这里他说，\Quote{我们这手艺有被藐视的危险。——你们充满了（一切）各处你们的道理。}\BibleRef{宗 5:28}犹太人论及基督时也这样说：\BibleQuote{你们看，全世界都跟他去了}\BibleRef{若 12:19}；以及\BibleQuote{罗马人必要来夺去我们的土地和人民}\BibleRef{若 11:48}。又在另一场合，\BibleQuote{这些扰乱天下的人也到这里来了。}\BibleRef{宗 17:6}。——\Quote{众人听了这话，就满心恼怒。}这种恼怒是因何而起？因听到关于阿尔特米和他们利益来源的话。\BibleQuote{喊着说，大哉！厄弗所的阿尔特米！满城一片混乱，他们齐心冲进剧场。}\BibleRef{宗 19:29}。\Quote{他们捉住了马其顿人加约和阿黎斯塔苛，保禄的旅伴，就拖去：}（此处）再次鲁莽，正如犹太人在雅松身上的所作所为；到处他们攻击他们。\BibleQuote{保禄正想进去见民众，门徒们不让他去。}\BibleRef{宗 19:30}，他们如此远离一切炫耀和贪图荣耀。\BibleQuote{有些亚细亚的首长，是他的朋友，派人到他那里，劝他不要冒险到剧场里去。}\BibleRef{宗 19:31}，到一个无序的民众和骚动之中。保禄听从，因为他不虚荣，也不野心勃勃。\BibleQuote{于是有人喊这个，有人喊那个：因为集会混乱。}\BibleRef{宗 19:32}。群众就是这样：像火落在燃料上一样鲁莽跟随；大多数人不知道为何聚集。\BibleRef{宗 19:32}\BibleQuote{他们从人群中拉出亚历山大，犹太人推他上前。}\BibleRef{宗 19:33}\BibleQuote{亚历山大摆手，要向民众分辩。}\BibleRef{宗 19:34}是犹太人推他上前；但按照上智安排，这人没有发言。\BibleQuote{但众人认出他是犹太人，就齐声喊了大约两个小时，大哉！厄弗所的阿尔特米。}\BibleRef{宗 19:34}真是幼稚的理解！好像他们害怕他们的崇拜会熄灭，就不停地喊叫。保禄在那里住了两年，看还有多少外教人！\BibleQuote{当书记官安抚了民众，说，厄弗所人啊，有谁不知道厄弗所城是伟大女神阿尔特米和自丢斯落下的神像的看守者呢？}\BibleRef{宗 19:35}。好像这事不明显似的。他首先用这话扑灭了他们的怒气。\Quote{以及自丢斯落下的神像。}另有被称为该名的另一\OriginalTerm{圣所}{\GreekText{ἱερόν}}。要么他指的是那块被烧的泥土，要么是她的神像。这是谎言。\Quote{既然这些事是无法辩驳的，你们应当安静，不要鲁莽行事。因为你们把这些人带来，他们既不是抢劫庙宇的，也不是亵渎你们女神的。}（36,37节）然而，他这一切是对民众说的；但为了使那些（工人）也变得合理，他说：\Quote{所以，如果德默特琉和他同业的工匠控告某人，法庭是开着的，并有代理人；让他们互相控告。但如果你们询问其他事情，应在合法集会中决定。因为我们有被控告今天骚动的危险，没有缘由，我们也就不能为这次聚集提出理由。}\BibleRef{宗 19:38}\Quote{合法集会}，他说，因为每月按律有三次集会；但这次是违法的。然后他又用\Quote{我们有被控告的}危险来恐吓他们。但让我们再看那些所说的话。\par

（综述。）\Quote{这些事结束后}，它说，\Quote{保禄在圣神内决意，经过\Place{马其顿}和\Place{阿哈雅}以后，要去\Place{耶路撒冷}}，说：\Quote{我在那里以后，也必须看看\Place{罗马}}。\BibleRef{宗 19:21} 他在这里不再以人的方式说话，或者，他决意经过那些地区，不再久留。为什么他打发\Person{弟茂德}和\Person{厄辣斯托}离开？关于这件事，我想他说：\Quote{因此，当我们不能再忍耐时，我们觉得独自留在\Place{雅典}是好的。他打发了，}它说，\Quote{两个服事他的人}\BibleRef{得前 3:1}，既是报告他的来临，也是使他们更热切。\Quote{他自己却在\Place{亚细亚}暂住。}\BibleRef{宗 19:22} 他大部分时间都待在\Place{亚细亚}；这有原因：那里，也就是哲学家的暴政。（后来）他也来再次向他们讲论。\Quote{大约那时候}等等。\BibleRef{宗 19:23}，因为迷信实在过分。(\Emph{a})\Quote{你们又看见又听见，}正在发生的结果如此明显——\Quote{不仅在此\Place{厄弗所}，而且几乎在全\Place{亚细亚}，这个\Person{保禄}用说服劝化了很多人，}不是用暴力：这是说服一座城的方式。然后，与他们切身相关的事：\Quote{那些不是神，而是人手所造的。}\BibleRef{宗 19:26} 他说，他颠覆了我们的手艺：(\Emph{e})\Quote{从这工作我们得了我们的财富。}他说服了人。他——一个无足轻重的人——是如何说服人的？如何胜过如此巨大的习惯力量？做了什么——说了什么？这不是一个\Person{保禄}（所能成就的），不是一个人所能成就的。仅仅这一点就够了：他说了\Quote{他们不是神。}现在，如果（异教宗教的）不敬如此容易被识破，它早就该被谴责了；如果它强大，它就不该这么快被推翻。(\Emph{b})因为，免得他们心里想（多么奇怪），一个人竟有这样的能力，如果一个人有能力成就这样的事，那么人就应该被那个人说服，他补充说：(\Emph{f})\Quote{不仅我们这手艺有危险被藐视，而且，}好像确实提出一个更大的考虑，\Quote{伟大女神\Person{狄安娜}的庙宇，}等等。(\Emph{c})\Quote{全亚细亚和普世都敬拜她。}\BibleRef{宗 19:27}(\Emph{g})\Quote{他们听见了，就充满怒气，喊着说：\InnerQuote{大哉！\Place{厄弗所}人的\Person{狄安娜}！}}\BibleRef{宗 19:28} 因为每座城都有它自己的神。(\Emph{d})他们想用自己的声音筑起一道屏障对抗天主圣神。这些希腊人真是孩子！(\Emph{h})他们的感觉好像用声音就能恢复对她的敬拜，并消除已发生的事！\Quote{全城，}等等。\BibleRef{宗 19:29} 看这无秩序的暴民！\Quote{当\Person{保禄}，}等等。\BibleRef{宗 19:30} \Person{保禄}当时想进去，以便向他们演说：因为他把所受的迫害当作教导的机会：\Quote{但门徒们不许他。}请注意，我们总是发现他们为他多么深谋远虑。起初他们把他带出来，免得他（本人）遭受致命打击；然而他们曾听他说过：\Quote{我也必须看看\Place{罗马}。}但这是天主的安排，他如此预先预言，免得他们因事件发生而困惑。但他们不愿他遭受任何坏事。\Quote{有些亚细亚的官长恳求他不要进入剧场。}知道他的热切，他们\Quote{恳求他：}众信徒都如此爱他。——\Quote{他们拉了\Person{亚历山大}，}等等。\BibleRef{宗 19:33} 这个\Person{亚历山大}，为什么想申诉？他被告了吗？不，而是要找一个机会，颠覆整个事情，并激起民众的愤怒。\Quote{但当他们知道他是犹太人时，众人同声大喊约两小时：\InnerQuote{大哉！\Place{厄弗所}人的\Person{狄安娜}！}}\BibleRef{宗 19:34} 你注意到这过分的愤怒了吗？好吧，文书带着责备说：\Quote{有哪一个人不知道\Place{厄弗所}城——}\BibleRef{宗 19:35}（切入要点）他们所害怕的。他说，是不是你们不敬拜她？他不说\Quote{不知道}\Person{狄安娜}，而是\Quote{我们的城，}说它一直敬拜她。\Quote{既然如此，这些事是不能反驳的。}\BibleRef{宗 19:36} 那么你们为什么还要质问它们，好像这些事不清楚？(\Emph{b})然后他温和地责备他们，表明他们毫无理由地聚集。他说，\Quote{也不要做任何鲁莽的事。}表明他们行事鲁莽。(\Emph{a})\Quote{因为你们把这些人带到这里，}等等。\BibleRef{宗 19:37} 他们想用宗教作为借口来掩盖他们自己的赚钱之事：(\Emph{c})不应该为私人的指控召开公众集会。因为他使他们陷入困境，不让他们有辩驳之词。他说，\Quote{没有原因，}\Quote{对于这次聚集，关于它}（事情）\Quote{我们不能提出报告。}\BibleRef{宗 19:40} 看这些不信者多么谨慎，多么聪明。他这样熄灭了他们的愤怒。因为正如它容易被点燃，也容易被熄灭。\Quote{他这样说了，}它说，\Quote{就解散了会众。}\BibleRef{宗 19:41}\par

你看见天主怎样容许试探，借此激动并唤醒门徒，使他们更加警醒吗？那么，我们不要在试探下沉沦；因为祂自己必\Quote{也开一条出路，使你们能够承担。}\BibleRef{格前 10:13} 没有什么比磨难更能缔结友谊，并使之如此牢固；没有什么更能凝聚和坚固信徒的灵魂；对于我们教师而言，没有什么比这更适时，好让听众听从我们所说的话。因为听众在安逸时是怠惰懒散的，甚至似乎对说话者感到厌烦；但当他在磨难和困苦中时，便极切渴望听讲。因为他内心苦恼时，四处寻求慰藉其困苦，而讲道所带来的安慰实在不小。\Quote{那么，}你会说，\Quote{犹太人呢？他们怎么会因心怯而不听呢？}唉，他们是犹太人，那些一向软弱可悲的人；而且，他们所受的磨难很大，但我们所说的是适度的磨难。因为请看：他们期望从周围的祸患中解脱出来，却陷入了无数更大的祸患；这对灵魂而言绝非普通的痛苦。磨难使我们割断对现世的眷恋，这表现在我们立刻希望死去，不再爱惜身体：这正是智慧的最大部分，即对现世生活无所贪恋，无所牵挂。遭受磨难的灵魂不愿操心许多事：它只渴望安息与宁静，满足于摆脱眼前的事物，即使此后别无他物。正如疲惫困苦的身体不愿沉溺于情欲或暴食，只愿安歇静卧；同样，被无数祸患折磨的灵魂，急切地渴望安息与平静。安逸的灵魂（容易）慌乱、惊恐、不安；而在磨难中则没有空虚，没有浪费；前者更显男子气概，后者更显稚气；前者更为稳重，后者更为轻浮。正如某种轻飘之物落入深水时，被颠簸飘荡；同样，灵魂陷入大喜乐时也是如此。再者，我们最大的过错源于过度的快乐，这是人人可见的。来吧，如果你愿意，让我们想象两座房子：一座是婚宴之地，另一座是哀悼之所；让我们想象进入其中，看看哪一所更胜一筹。啊，哀悼者的房屋将充满\OriginalTerm{哲学}{\GreekText{φιλοσοφίας}}，而婚宴者的房屋则充满猥亵。因为请看：羞辱的话语，无羁的笑声，更加放纵的动作，服饰和步态充满猥亵，言语尽是无稽和愚妄；总之，那里一切皆是戏弄，一切可笑。我并非说婚姻本身如此；愿天主禁止；而是婚姻的附属品。那时，本性在过度的狂欢中失常。在场的人不再是人类，而变成了畜类：有的像马嘶鸣，有的像驴踢踏：如此彻底的放纵，如此放荡不羁：没有严肃，没有高尚：（这是）魔鬼的排场，铙钹、笛子和充满奸淫、淫乱的歌曲。但在那哀悼的房屋里却非如此；那里一切都井然有序：如此安静，如此安息，如此镇定；毫无混乱，毫无过分；如果有人说话，他所说的每一句话都充满真正的哲学；而且奇妙的是，此时不仅男人，甚至仆人和妇女都像哲学家一样说话——因为忧伤的本性如此——当他们似乎安慰哀悼者时，实际上他们说出无数充满健全哲学的真话。那里首先有祈祷，愿苦难就此止住，不再加深；许多人安慰受苦者，并无数地述说许多同样有理由哀悼的人。\Quote{人算什么？}（他们问）进而严肃地审视我们的本性——\Quote{唉，那么，人算什么！}（接着）控诉（现世）生命及其无价值，提醒（彼此）将来的事，审判。（因此，从这两处场景）每个人回家：从婚宴回来，忧伤，因为自己未能享受同样的好运；从哀悼回来，轻松，因为自己未曾遭受同样的苦难，且内心的一切狂热都熄灭了。但你想要什么？我们要以监狱和剧场作为另一对比吗？因为前者是受苦之地，后者是享乐之所。让我们再考察。在监狱中，有严肃的心境；因为哪里有悲伤，就必然有严肃。从前富有且自负的人，现在甚至愿意让任何普通人与他交谈，恐惧和悲伤如同更强大的火焰落在他的灵魂上，软化了他的刚硬：于是他变得谦卑，面容忧愁，感受到人生的变迁，勇敢地承受一切。但在剧场中，一切恰恰相反——欢笑、粗鄙、魔鬼的排场、放荡、浪费时间、虚度光阴、谋划过度的情欲、研究奸淫、演习淫乱、接受放荡的教育、鼓励污秽、成为笑料、成为猥亵行为的榜样。监狱却非如此：在那里你会找到谦卑的心、劝勉、严肃的激励、对世俗之物的鄙视；这些都被践踏和唾弃，而恐惧如同教师管束孩童一样临到（那人身上），管束他尽一切本分。但如果你愿意，让我们用另一种方式考察。我希望你遇见一个刚从剧场回来的人，和另一个从监狱出来的人；你会看到前者的灵魂慌乱、不安，实际上被捆绑束缚；后者的灵魂你会看到开阔、自由，如展翅般轻盈。因为前者从剧场回来，被那里所见的女人束缚，背负着比铁还硬的锁链——他所见的场景、言语、姿态。而后者从监狱回来后，从一切（束缚）中释放，再不会因与他人的境况相比而认为自己遭受任何不幸。（想到）自己不在捆绑中，会使他此后永远感恩；他会轻视人间事务，因为看到许多富人在那里遭难，曾经有能力做许多大事的人，如今躺在那里被捆绑；如果他遭受任何不公正，他也会忍受；因为那里也有许多这样的例子：他会被引导思考将来的审判，并战栗，看到这里（地上的监狱）中是怎样，将来在那里也将怎样。因为正如一个人被关在这里的监狱中，同样，在那世界，在审判之前，在将来的日子之前也是如此。他对妻子、儿女和仆人将更加温和。\par

从戏院回来的人则不然：他看待妻子更加厌恶，对仆人生气，对子女刻薄，对众人粗暴。戏院给城市带来的祸害很大，的确很大，我们甚至不知道它们有多大。我们也要审视其他欢笑的场景吗？我指的是宴会，有奉承者、谄媚者和极尽奢华；相形之下，瘸子和盲人所在之处呢？如前所述，前者有醉酒、奢华和放荡，后者则相反。再看身体，当它血气旺盛、状态良好时，最容易染病；当它保持低落时则不然。那么让我向你们更清楚地说明我的意思：假设有一个身体，血液充盈，肌肉丰满，因生活优渥而肥胖；这个身体即使只是吃一点偶然的食物也容易引起发烧，只要它悠闲不动。但另一个身体，与饥饿和艰辛搏斗的：就不易被疾病制服，不易被疾病压倒。血液，虽然在我们体内可能是健康的，但往往因量过多而引发疾病；但如果量少，即使不健康，也容易排出。灵魂也是如此，过舒适奢华生活的灵魂，其冲动很快倾向罪恶：因为这样的灵魂邻近愚蠢、享乐、虚荣、嫉妒、阴谋和诽谤。看我们这座大城，何其庞大！罪恶从何而来？岂非来自富人？岂非来自享乐之人？是谁把人\Quote{拖}人\Quote{到法庭？}是谁挥霍财产？是那些悲惨和被遗弃的人，还是那些自以为是、享乐的人？一个受苦的灵魂不可能引发任何邪恶。\BibleRef{雅 2:6}保禄知道这益处：因此他说，\BibleQuote{磨难生忍耐，忍耐生老练，老练生望德，望德不叫人羞愧。}\BibleRef{罗 5:3}那么，让我们不要在磨难中消沉，反而凡事感谢，这样我们才能获得巨大收益，才能中悦那允许磨难的天主。磨难是大善：我们从自己的孩子学到这一点：因为没有磨难，孩子什么有用的也学不到。但我们比他们更需要磨难。因为如果在那里，当情欲尚安静时，责罚对他们有益，何况对我们呢？尤其我们被如此多的情欲所占据！不，我们比他们更需要老师：因为孩子的过失不会很大，但我们的过失却极大。我们的老师就是磨难。那么，我们不要故意招引它，但当它来到时，让我们勇敢地承受，因为它总是无数好事的缘由；好使我们既获得天主的恩宠，也获得为爱祂之人所预备的美善事物，在我们的主耶稣基督内，愿光荣、权能、尊威归于祂及圣父，偕同圣神，从今世直到永远，万世无疆。阿们。\par

\SourceRecord{Translated by J. Walker, J. Sheppard and H. Browne, and revised by George B. Stevens.；Nicene and Post-Nicene Fathers, First Series；Vol. 11.；Edited by Philip Schaff.；Buffalo, NY: Christian Literature Publishing Co.,；1889.；<http://www.newadvent.org/fathers/210142.htm>.}

% structure: p0001 source=h1 target=section reason=page-title

\section{\WorkTitle{论宗徒大事录 第四十三讲}}

\BibleRef{宗 20:1}\par

\BibleQuote{骚动平息以后，保禄叫门徒来，拥抱他们，就出发往马其顿去。}\par

骚动之后，需要很多的安慰。因此，他做了这事之后，就去了马其顿，然后去了希腊。因为经上说：\BibleQuote{他走遍那一带地方，用许多话劝勉他们以后，来到了希腊，在那里住了三个月。当犹太人设计害他，正要乘船往叙利亚去的时候，他就决意从马其顿回去。}（第2-3节）他又受犹太人迫害，往马其顿去。\BibleQuote{同他到亚细亚去的，有贝勒雅的\Person{索帕特尔}；得撒洛尼人\Person{阿黎斯塔恰}和\Person{塞滚德}；德尔贝的\Person{加约}和\Person{弟茂德}；亚细亚人\Person{提希苛}和\Person{特洛斐摩}。这些人先走，在特洛阿等我们。}（第4-5节）但他怎么称\Person{弟茂德}是\Quote{得撒洛尼人}呢？他的意思并非如此，而是\Quote{得撒洛尼人\Person{阿黎斯塔恰}、\Person{塞滚德}和\Person{加约}；德尔贝的\Person{弟茂德}}等等。他说这些人先到特洛阿去，为他预备道路。\BibleQuote{过了无酵节的日子，我们从斐理伯开船，五天到了特洛阿他们那里，在那里住了七天。}\BibleRef{宗 20:6}因为在我看来，他特意在大城市守节日。\BibleQuote{从斐理伯，}那里发生过监狱事件。这是他第三次来到马其顿，他对斐理伯人作了很高的见证，这也是他在那里住了一些日子的原因。\BibleQuote{在一周的第一天，门徒们聚在一起擘饼的时候，保禄因为他们次日就要起程，就向他们讲道，直讲到半夜。}\BibleRef{宗 20:7}那时正是五旬节期间。看一切如何都服从于讲道。经上说，那天也是主日。甚至在夜间他也没有沉默，反而因为安静而更加讲论。请注意他怎样做了长篇讲论，甚至超过了晚餐的时间。但魔鬼扰乱了宴席——虽然没有得逞——使听者沉睡，让他跌下去。\Quote{又有，}经上说，\BibleQuote{在他们聚集的楼上，有许多灯。有一个少年人，名叫\Person{厄乌提曷}，坐在窗台上，困倦沉睡；保禄讲论了许久，他因睡熟了，就从三楼坠下去，扶起来已经死了。保禄下去，伏在他身上，抱着他说：你们不要慌乱，因为他的灵魂还在他身上。他遂上去，擘开饼，吃了，又谈了许久，直到天亮，这才走了。他们把那少年人活活地带回来，就大为安慰。}\BibleRef{宗 20:8}但请你们看，剧场多么拥挤！奇迹是什么？\BibleQuote{他坐在窗台上，}在深夜。他们对听道的渴望如此之大！让我们羞愧吧！你们说：\Quote{啊，但那时是保禄在讲论。}是的，现在保禄也在讲论，或者说，那时和现在都不是保禄，而是基督，然而没有人愿意听。现在没有窗台，没有饥饿或睡眠的催促，我们却不愿听；这里没有拥挤的空间，也没有其他舒适。奇妙的是，他虽然是个少年，却不懒散冷漠；虽然被睡意压住，却没有离开，也不怕坠落的危险。他不是因为懒散而打盹，而是出于本性的需要。但请你们注意，我恳求你们，他们的热忱如此炽烈，以至他们在一个三楼聚会：因为他们还没有教堂。\BibleQuote{你们不要慌乱，}他说。他没有说：\Quote{他会复活，因为我要使他复活：}而是看他怎样谦逊地安慰他们：\BibleQuote{因为他的灵魂，他说，还在他身上。他上去，擘开饼，吃了。}这件事打断了讲论，却没有造成伤害。\Quote{他吃了以后，}经上说，\BibleQuote{又讲了许久，直到天亮，这才走了。}你们注意到晚餐的简朴吗？你们注意到他们如何整夜度过吗？他们的聚餐如此，以至于听者离开时头脑清醒，适合听道。而我们与狗有什么不同？你们注意到差别吗？\BibleQuote{他们把那少年人活活地带回来，}经上说，\BibleQuote{就大为安慰，}既因为看到他复活，也因为奇迹发生了。\BibleQuote{我们先上船，开往塔索斯，要在那里接保禄；因为他这样安排，自己打算步行。}\BibleRef{宗 20:13}我们常见保禄与门徒分离。看，他又自己步行：让他们走容易的路，自己选更痛苦的路。他步行，既为安排许多事，也为了训练他们忍受与他分离。\BibleQuote{他在塔索斯与我们相会，我们接他上船，来到米蒂利尼；从那里次日开船，到了希俄斯对面}——然后他们经过那岛——\BibleQuote{次日我们到了萨摩斯，在特罗基利翁停留，次日到了米利都。因为保禄已决定绕过厄弗所，免得在亚细亚浪费时间：他急忙，如果可能，要在五旬节到达耶路撒冷。}\BibleRef{宗 20:14}为什么这样匆忙？不是为了节日，而是为了群众。同时，他以此赢得了犹太人，因为他是一个尊重节日的人，希望甚至征服他的对手；同时他也传扬了圣言。因此，请看，由于众人都在，获得了多大的益处。但为了使厄弗所人的利益因此不受忽视，他以另一种方式安排。但我们再回顾一下所说过的话。\par

\Emph{（总结）} \BibleQuote{他拥抱了他们}，经上说，\BibleQuote{他动身往马其顿去。}\BibleRef{宗 20:1} 借此，他\OriginalTerm{再次安慰了他们}{\GreekText{ἀνεκτήσατο}}，给了他们许多安慰。\BibleQuote{他劝勉了}马其顿人，\BibleQuote{用许多话语，他来到了希腊。}\BibleRef{宗 20:2} 请注意，我们在各处都发现他藉着宣讲完成一切，而非藉着奇迹。\BibleQuote{我们开船了}等等。作者不断地向我们显示他急忙赶往叙利亚；原因是教会和耶路撒冷，但他仍然克制自己的愿望，以便也在那些地方整顿一切。特洛阿并不是一个大地方：为什么他们在那里停留了七天？也许从信徒的数量来看它是大的。他在那里过了七天后，第二天他整夜教导：他难以离开他们，他们也难以离开他。\BibleQuote{当我们聚在一起时}经上说，\BibleQuote{擘饼。}\BibleRef{宗 20:7} 正是在（擘饼的）时刻，讲论开始了*，延长了：表明他们饿了，这不是不合时宜的：因为主要目的（使他们聚在一起的）不是教导，他们聚在一起是为了擘饼；然而讲论开始了，他就延长了教导。请看，所有人也都分沾了保禄的餐桌。在我看来，他甚至在坐席时也讲论，教导我们要把所有其他事情都视为次要于此事。我恳请你们，想象那间房子，有灯光，有人群，保禄在中间讲论，甚至窗户都挤满了人：那是何等景象，去观看，去听那号角，去看那慈祥的面容！但他为什么在夜间讲论？因为经上说：\BibleQuote{他即将离开}，并且不再见到他们：虽然他并没有告诉他们这一点，他们太软弱（承受不了），但他确实告诉了其他人。同时，发生的奇迹会使他们永远记念那个晚上；所以跌倒反而成了老师的益处。听者的喜乐是大的，甚至被打断时更增加了。那个（年轻人）要责备所有对（话语）漫不经心的人，他的死只是因为这件事：他希望听保禄讲话。\BibleQuote{我们先上船去}等等。\BibleRef{宗 20:13} 为何作者说他们到了哪里，又去了哪里？首先表明他更悠闲地航行——这是出于人性的原因——并且驶过（一些地方）；也（出于同样的原因他告诉）他在哪里停留，驶过了哪些地方；（即）\BibleQuote{使他不至在亚细亚耗费时间。}\BibleRef{宗 20:16} 因为他若是到了那里，就不能驶过去；他不愿使那些会恳求他留下的人痛苦。\BibleQuote{他急忙，}经上说，\BibleQuote{如果可能，要在耶路撒冷过五旬节：}而这（如果留下）是不可能的。请注意，他也像其他人一样被感动。因此这一切发生，是为了使我们不认为他超乎人性：（因此）你看见他渴望（某物），急忙，在许多情况下不能达成（目标）：因为那些伟大的圣人与我们同有一样的本性；区别在于意志和目的，也正因此他们吸引了那伟大的恩宠。看，例如，他们用自己的经济安排多少事。\BibleQuote{我们不给任何冒犯} \BibleRef{格后 6:3} 给那些想要（被冒犯）的人，以及，\BibleQuote{使我们的职分不被指责。}看，无可指责的生活和另一方面谦让。这确实被称为经济，达到（它的）顶峰和高度。因为那超越基督诫命的人，另一方面比众人更谦卑。\BibleQuote{我成了一切人中的一切，}他说，\BibleQuote{为要赢得一切人。}\BibleRef{格前 9:22} 他也投身于危险中，正如他在另一处所说：\BibleQuote{在极多的忍耐中，在患难中，在困乏中，在苦难中，在鞭打中，在监禁中。}\BibleRef{格后 6:4} 他对基督的爱是大的。因为如果没有这个，其余一切，无论是（迁就的经济），还是无可指责的生活，还是置身于危险，都是多余的。\BibleQuote{谁软弱，}他说，\BibleQuote{而我不软弱？谁跌倒，而我不焦急？}\BibleRef{格后 11:29} 让我们效法这些话，为弟兄的缘故置身于危险。无论是火，还是刀剑，亲爱的，请你投身，为拯救你的肢体：投身，不要害怕。你是基督的门徒，祂为弟兄舍了性命：与保禄同作门徒，他选择了为敌人，为那些与他争战的人忍受无数苦难；你要满怀着热忱，效法梅瑟。他看见一人受委屈，就为他报仇；他轻视王室的奢华，为了那些受压迫的人成了逃亡者，流亡者，孤单和被遗弃；他在异乡度过时日；然而他没有责怪自己，也没有说：\Quote{这是什么事？我轻视了王权，及那一切尊荣和光荣：我选择了为受屈的人报仇，而天主却忽略了我：不但没有使我恢复从前的尊荣，反而四十年来我流落异乡。真的，我得到了美好的报偿，不是吗！}但他没有说也没有想任何这类的事。你也当如此：即便你因行善而受苦，即便（须等待）很长时间，不要生气，不要慌乱：天主必定会给你赏报。报偿越延迟，利息就越多。让我们有一个易于同情的心灵，让我们有一颗能与人同忧的心：没有无情的性格（\OriginalTerm{无情}{\GreekText{ὠμόν}}），没有人性缺乏。\par

纵然你不能提供任何缓解，仍要哭泣、呻吟、哀伤所发生的事：这也不是毫无益处。如果我们应当同情那些被天主公义惩罚的人，更应当同情那些被不义地虐待的人。（他们中的）\Quote{厄南，} 经上说，\Quote{没有出来为那靠近她的房屋哀悼} \BibleRef{米 1:11}：他们将受痛苦，\Quote{因为他们建了作嘲笑之物。} 再者，厄则克耳指责他们，因为他们没有为（受苦难者）哀悼。\BibleRef{则 16:2} 先知啊，你说什么？天主惩罚，我应当为那些受祂惩罚的人哀悼吗？确实如此，因为天主自己虽惩罚，也愿意这样，因为祂自己也不愿惩罚，甚至祂在惩罚时也悲伤。所以你不可因此欢喜。你将会说：\Quote{如果他们受的是公义的惩罚，我们就不该哀悼。} 为何？我们所应哀悼的正是这一点——他们竟配受惩罚。试想，当你看见你儿子接受烧灼或刀割，你岂不悲伤？你岂不对自己说：\Quote{这是什么？这是为健康而切割，以加速他的痊愈；这是为挽救他而烧灼？} 但尽管如此，当你听见他因无法忍受痛苦而呼喊时，你仍悲伤，而且痊愈的希望并不足以消除本性所受的震撼。这些人的情形也是这样：虽然受惩罚是为他们的益处，但我们仍应表现兄弟之情、为父之心。天主的惩罚就是切割和烧灼，但我们应当为此哭泣，因为他们有病，他们需要这种医治方式。如果有人为冠冕而受苦，你不要哀伤；例如保禄、伯多禄所受的苦：但若有人因受惩罚而承受公义，则要哭泣、呻吟。先知们就是如此行的；其中一位说：\Quote{哎！主啊，你要消灭以色列的残余吗？} \BibleRef{则 9:8} 我们看到杀人者、恶人受惩罚，我们感到痛苦，为他们哀伤。我们不可过度哲学化：要显示怜悯，为使我们得到怜悯；没有什么比这种美德更美好：没有什么比显示怜悯、善待同胞更能彰显人性。事实上，法律将惩罚的全部事务交给行刑者：他们迫使法官只宣判刑罚，然后召唤行刑者执行。确实如此，即使惩罚是公正的，但施行惩罚并非一个慷慨的（\OriginalTerm{哲士}{\GreekText{φιλοσόφου}}）灵魂的分内之事，而是需要另一种人：因为甚至天主也不亲手惩罚，而是借着天使。那么，天使是行刑者吗？绝对不是！我并非说他们如此，但他们是复仇之能者。当索多玛被毁灭时，完全是藉他们为工具；当埃及遭受审判时，也是藉着他们。因为经上说：\Quote{祂派遣，} \Quote{凶恶的天使在他们中间。} \BibleRef{咏 77:50} 但若有必要说明，天主亲自如此行事：祂派遣了圣子——（\Emph{b}）然而，\Quote{接纳你们的，就是接纳我；接纳我的，就是接纳那派遣我的。} \BibleRef{玛 10:40} （\Emph{a}）祂又说：\Quote{那时我必对天使说：把作恶的人收集起来，扔进火炉里。} \BibleRef{玛 13:30} 但有关义人，则不然。（\Emph{c}）又说：\Quote{捆起他的手和脚，把他丢在外面的黑暗中。} \BibleRef{玛 22:13} 注意，在那种情况下，祂的仆役服侍；但在行善时，祂亲自行善，亲自召唤：\Quote{来吧，我父所祝福的，继承为你们预备的国度。} \BibleRef{玛 25:34} 当事情是与亚巴郎交谈时，祂亲自到他那里；当要离开索多玛时，祂派遣祂的仆役，如同法官唤起那些行刑的人。\Quote{你在小事上忠信，我要委派你管理许多事。} \BibleRef{玛 25:21}；\Emph{I} (will make you)：但另一些事，不是祂亲自，而是祂的仆役捆绑。知道这些事，我们不要因受苦的人欢喜，反要哀伤：应为他们哀悼、哭泣，为由此也可获得赏报。但现在，许多人甚至因那些无辜受苦的人而欢喜。但我们却不可如此：要显示全部同情，使我们也蒙天主眷顾，藉祂独生子耶稣基督的恩宠和慈爱，愿光荣、权能、尊崇归于祂及圣父、圣神，从今世直到永远。阿们。\par

\SourceRecord{Translated by J. Walker, J. Sheppard and H. Browne, and revised by George B. Stevens.；Nicene and Post-Nicene Fathers, First Series；Vol. 11.；Edited by Philip Schaff.；Buffalo, NY: Christian Literature Publishing Co.,；1889.；<http://www.newadvent.org/fathers/210143.htm>.}

% structure: p0001 source=h1 target=section reason=page-title

\section{\WorkTitle{宗徒大事录讲道 第四十四篇}}

\BibleRef{宗 20:17-21}\par

\BibleQuote{他从\Place{米肋托}派人到\Place{厄弗所}，请教会的长老来。他们来到他那里，他便对他们说：你们知道，自从我进入\Place{亚细亚}的第一天起，我怎样与你们相处，始终以全然的谦卑和许多的眼泪，并因犹太人的阴谋所加于我的试探，服事主；我也没有隐瞒任何有益于你们的事，反而公开地并挨家挨户地教导你们，向犹太人和希腊人作证：对天主的悔改，以及对我们的主耶稣基督的信德。}\par

看哪，他急于航行经过，却并未忽略他们，反而为众人安排一切。他派人召来长老，借着他们向（厄弗所人）讲论；但令人钦佩的是，当他发现自己不得不说出一些关于自己的重大事情时，他力图尽量少提（\OriginalTerm{力求谦逊}{\GreekText{πειρἅτα} \GreekText{μετριάζειν}}）。\Quote{你们知道。}正如撒慕尔在将统治权交给撒乌耳时，在他们面前说：\BibleQuote{我岂从你们手中拿过什么吗？你们作证，天主也作证}\BibleRef{撒上 12:3}；（保禄也是如此）。达味在不被相信时也说：\BibleQuote{我正与羊群在一起，看守我父亲的羊；当熊来时，我用手把它吓跑}\BibleRef{撒上 17:34}；而且保禄自己也对格林多人说：\BibleQuote{我成了愚妄的人，是你们逼我的}\BibleRef{格后 12:11}。不仅如此，天主自己也如此行，并非在任何场合都谈论自己，而只在不受相信时才提及祂的恩惠。因此，请看保禄在此所作：首先他引述他们自己的见证；为使你不以为他的话只是自夸，他请听众亲自为他的话作证，因为他不可能在他们面前说谎。这就是老师的卓越之处：让自己的门徒作他功绩的见证。而且令人惊异的是，他说，他不是一天两天这样行。他愿意鼓舞他们面向未来，使他们能勇敢承受一切，既承受与他分离，也承受将要发生的考验——正如梅瑟和若苏厄的情形。请看他是如何开始的：\Quote{我如何始终与你们在一起，以全心的谦逊服事主。}请注意，这对统治者最相称的是：\BibleQuote{憎恶骄傲}\BibleRef{出 18:21}，梅瑟说；这个资格尤其切合统治者，因为他们几乎不可避免会变得傲慢。这（谦逊）是一切善行的根基，事实上基督说：\BibleQuote{神贫的人是有福的}\BibleRef{玛 5:3}。而且这里不是简单地说\Quote{以谦逊的心}，而是\Quote{以全心的谦逊}。因为谦逊有多种：言语上的和行为上的，对待统治者和对待被统治者。你们愿意我向你们提及几种谦逊吗？有些人向卑微的人谦卑，向高贵的人高傲：这不是谦逊的品质。有些人就是如此。然后，为免显得傲慢，他事先打下基础，消除那怀疑：他说：\Quote{如果我曾\InnerQuote{以全心的谦逊}行事，那么我现在所说的话并非出于傲慢。}然后是他的温和，总是以降尊纡贵使他们成为同伴。他说：\Quote{与你们}，他说，\Quote{我曾一同服事主}；他将善工与众人共享，无一归己。\Quote{什么？}（你可能会说）\Quote{难道他还能对天主傲慢吗？}然而有许多人对天主傲慢，但这个人甚至不对自己的门徒傲慢。这就是老师的功德：以自己的德行成就来塑造门徒的品格。然后是他的刚毅，在这方面他也非常简洁。他说：\Quote{经历许多眼泪，}和\Quote{因犹太人的谋算而临到我的试炼}。你看到他因他们的行为而悲伤吗？但这里他似乎也显示了他的同情：因为他为那些走向丧亡的人受苦，为他们本人；自己所受的，他甚至为此喜乐；因为他属于那伙人，他们\BibleQuote{因被算配为这名字受辱而喜乐}\BibleRef{宗 5:41}；他又说：\BibleQuote{如今我在为你们受苦中喜乐}\BibleRef{哥 1:24}；又说：\BibleQuote{我们这短暂轻微的苦难，正为我们成就极重无比永远的荣耀}\BibleRef{格后 4:17}。然而，他说这些话，是出于谦逊（\OriginalTerm{谦逊地}{\GreekText{μετριάζων}}）。但他在那里显示的刚毅，与其说在于勇敢，不如说在于忍耐。他说：\Quote{我曾受虐待，但与你们一同；而且事情中最痛苦的部分是：出于犹太人之手。}请注意，他在这里同时表现爱德与刚毅。我请你们留意教学的特色：他说：\Quote{我未曾保留任何事}，毫不吝啬的完全，毫不退缩的敏捷——\Quote{凡对你们有益的事}；因为有些事情他们不需要学习。因为隐瞒某些事就像吝啬，而说出一切事则是愚昧。所以他又加上：\Quote{凡对你们有益的事，我都显示给你们，并教导了你们}；不仅说了，也教导了；而且不是出于形式。要知道这就是他的意思，请看他说的话：\Quote{公开地，又挨家挨户}；以此表示极大的辛劳、极大的恳切和忍耐。\Quote{对犹太人和希腊人}。不仅是对你们。\Quote{见证：}这里指放胆宣讲；即使我们行不出善事，也必须宣讲；因为\Quote{见证}的意思是对不留意的人说话；所以\OriginalTerm{郑重见证}{\GreekText{διαμαρτύρασθαι}}这个词大多用于此。\BibleQuote{我呼天唤地作证}\BibleRef{申 4:26}，\OriginalTerm{我郑重见证}{\GreekText{διαμαρτύρομαι}}，梅瑟说；如今保禄自己\OriginalTerm{郑重见证}{\GreekText{Διαμαρτυρόμενος}}：\Quote{对犹太人和希腊人讲论悔改归向天主}。你见证什么？他们应当注意生活的方式：应当悔改，亲近天主。\Quote{对犹太人和希腊人}——因为犹太人不认识祂——他说，一方面因他们的行为，\Quote{悔改归向天主}，另一方面因他们不认识圣子，他又加上：\Quote{并信靠主耶稣}。那么，你说这些事是为达到什么目的？你提醒他们这些事是为达到什么目的？结果如何？你有什么事可指控他们？先激起他们的情感，然后他加上：\BibleQuote{现在，看，我受圣神催促，往耶路撒冷去，不知道在那里会遭遇什么事；但知道圣神在各城里向我作证说：有锁链和患难在等着我。但这些事都不能动摇我，我甚至不把我的性命为宝贵，为要欢欢喜喜地完成我的路程，并完成我从主耶稣所领受的职务，为天主恩宠的福音作证}\BibleRef{宗 20:22}。他为何说这些话？为预备他们时常准备面对危险，无论是可见的还是不可见的，并在一切事上顺服圣神。他表明他是为伟大的目标而离开他们。\Quote{但知道圣神，}他说，\Quote{在各城里向我作证说}——为表明他甘愿离去；免得你想那是束缚或必要，当他说：\Quote{受圣神催促——在各城里锁链和患难等着我}。然后他又加上：\Quote{我不把我的性命为宝贵，直到我完成了我的路程和从主耶稣所领受的职务。}\Quote{直到我欢欢喜喜地完成我的路程}，他说。你是否注意到这些话显然不是出自哀叹之人，而是出自不愿夸大（苦难）的人（\OriginalTerm{谦逊的}{\GreekText{μετριάζοντος}}），出自愿意教导那些听众，并与他们同情所发生之事的人。他并没有说：\Quote{我确实痛苦，但必须忍受}；而是说：\Quote{但这些事我都不能动摇，我也不认为我的性命宝贵}。这又是为什么？不是为夸耀自己，而是为教导他们，如同前者教导谦逊，后者教导刚毅和胆量：\Quote{我不以它为宝贵}，即\Quote{我不爱它胜于这事：我认为完成我的路程、作见证更为宝贵。}他并没有说\Quote{宣讲}、\Quote{教导}——而是说什么？\Quote{作见证}（\OriginalTerm{郑重见证}{\GreekText{διαμαρτύρασθαι}}）——天主的恩宠的福音。他即将说出一些\OriginalTerm{更沉重的话}{\GreekText{φορτικώτερον}}，即：\Quote{我从众人的血中是无罪的，因为在我这方面毫无欠缺}；他即将把全部重量和负担加在他们身上，所以先缓和他们的情绪，说：\Quote{现在，看，我知道你们不能再见到我的面了}。安慰是双重的：一方面\Quote{你们不能再见到我的面}，因为我心中与你们同在；另一方面并非只有他们（不能再见到他）：因为\Quote{你们所有的人，就是我曾四处宣讲天国的你们，不能再见到我的面了}。所以他完全可以（说）：\Quote{因此我请你们作证（读作\OriginalTerm{因此我作证}{\GreekText{διὸ} \GreekText{μαρτ}}，代替\OriginalTerm{郑重见证}{\GreekText{διαμαρτ}}）——既然我不再与你们同在——\BibleQuote{我与众人的血是无罪的}\BibleRef{宗 20:26}。}你看到他如何恐吓他们吗？当他们心灵烦恼痛苦时，他如何严厉地磨炼他们（\OriginalTerm{严加训诫}{\GreekText{ἐπιτρίβει}}）？但这是必要的。他说：\BibleQuote{因为我未曾回避，向你们宣示天主的全部计划}\BibleRef{宗 20:27}。那么，不讲的人就要负血债：也就是谋杀！没有比这更可怕的了。他表明他们若不做，也要负血债。所以，虽然他似乎在为自己辩护，其实是在恐吓他们。\BibleQuote{所以你们要谨慎，自己也要谨慎，也要照顾全群，圣神立你们为监督（或主教），牧养天主的教会（见注释3），就是祂用自己的血所取得的}\BibleRef{宗 20:28}。你看到了吗？他吩咐他们两件事。单单成功带领别人归正本身并无益处——因为他说：\BibleQuote{免得我传福音给别人，自己反被淘汰}\BibleRef{格前 9:27}；单单为自己努力也无益。因为这样的人是自私的，只求自己的利益，就像那个埋藏塔冷通的人。\Quote{谨慎自己}：他说这话，并非因为自己的得救比羊群的得救更宝贵，而是因为当我们谨慎自己时，羊群也受益。\Quote{圣神立你们为监督，牧养天主的教会}。看，你们的圣职来自圣神。这是一个约束；然后他说：\Quote{牧养主的教会}。看！另一个义务：教会是主的。第三个：\Quote{就是祂用自己的血所取得的}。这表明这事何等宝贵；危险非同小可，因为祂连自己的血都不吝惜。祂为使仇敌和好，甚至倾流了自己的血；而你，即使他们成了你的朋友，却不能留住他们。\BibleQuote{因为我知道，在我离去之后，将有凶暴的豺狼进入你们中间，不怜惜羊群}\BibleRef{宗 20:29}。他又从另一角度劝诫他们（\OriginalTerm{从另一角度劝诫}{\GreekText{ἐπιστρέφει}}），从将要发生的事入手：正如他说：\Quote{我们不是对抗血肉。}\Quote{在我离去之后，}他说，\BibleQuote{凶暴的豺狼将进入你们中间}\BibleRef{弗 6:12}；双重的恶，既因为他自己不在，又因为别人攻击他们。\Quote{既然你预先知道，为何还要离开？}他说，圣神催促我。他提到\Quote{豺狼}和\Quote{凶暴的，不怜惜羊群}；更糟的是，甚至\Quote{从你们自己中间}：这严重的事（就是），尤其是内战。这事极为严重，因为是\Quote{主的教会}：危险巨大，因为祂用血救赎了它；战争强大，且是双重的。\BibleQuote{就是从你们自己中间，也要有人起来，讲说悖谬的话，引诱门徒跟从他们}\BibleRef{宗 20:30}。\Quote{那么怎样？有什么安慰？}\BibleQuote{因此，你们要警醒，记念我三年之久，昼夜不停地流泪劝诫你们每一个人}\BibleRef{宗 20:31}。请看这里有多少强烈的措辞：\Quote{用眼泪}、\Quote{昼夜}和\Quote{每一个人}。因为他不是看见许多人时才参与工作，即便是为单个灵魂，他也能做一切（为那一个灵魂）。事实上，他就是这样使他们\OriginalTerm{紧密团结}{\GreekText{συνεκρότησεν}}。\Quote{我方面已做得足够：我留了三年}；他们有足够的建立，他说，足够的覆盖。\Quote{用眼泪}，他说。你看到眼泪是因这个缘故吗？恶人不悲伤：你悲伤；或许他也会悲伤。就像病人看到医生进食，也被激励同样行；同样在这里，当他看到你哭泣，他就被软化：他将成为一个善良而伟大的人。\par

（重述。）他说：\Quote{不知道，} \Quote{要临到我身上的事。}（第22-23节。）那么这就是你离开的原因吗？绝无此事；相反地，（我知道，）\Quote{锁链和患难在等着我。}我知道有考验，但不知道是什么种类：这更令人痛苦。\Quote{但这些事都不能动摇我} \BibleRef{见：宗 20:24}：因为不要以为我说这些是哀叹：因为\Quote{我不珍惜自己的性命。}他说这一切是为了提升他们的心志，劝勉他们不仅不要逃跑，还要勇敢地忍受。因此，他称此为\Quote{赛程}和\Quote{职务}，一方面，因其是竞赛而显示其荣耀，另一方面，显明其作为职务所当尽的本分。我是一名仆人；仅此而已。他先安慰他们，免得他们因他遭受如此恶待而悲伤，并告诉他们他\Quote{以喜乐}忍受了那些事，又显示出这些事的果实，然后（而非在此之前）他才引入那会令他们痛苦的事，以免压倒他们的心神。\Quote{现在看哪，}等等。\Quote{因此我向你们作证，我今天在众人的血上是纯洁的，因为我没有退缩，不将天主的全部计划宣告给你们} \BibleRef{见：宗 20:25}：* * * 那关乎当下之事的计划。\Quote{因为我知道这事，}等等。\BibleRef{见：宗 20:29}\Quote{那么，}有人会说，\Quote{你以为自己如此伟大吗？如果你离开，我们都要死吗？}他回答说：我不是说我的离开造成了这事；而是什么？将有另一种人起来反对你们：他不是说\Quote{因我离开}，而是说\Quote{在我离开之后：}即在他启程之后。——然而这事已经发生；将来（更会如此）更加严重。于是我们有了原因，\Quote{拉拢门徒跟随他们。} \BibleRef{见：宗 20:30}有异端，这就是原因，除此之外别无他因。接着也有安慰。但如果祂\Quote{赎买了}它\Quote{用自己的血，}祂必定会为它辩护。\Quote{昼夜不停地，}他说：\Quote{我流泪劝诫你们。} \BibleRef{见：宗 20:31}这话也很适用于我们的情况：虽然这话似乎特别指向教师，但同样也适用于门徒。因为，即使我昼夜说话、劝勉、流泪，但门徒若不听从，那又如何呢？因此他说：\Quote{我向你们作证：}因为他也说：\Quote{我在众人的血上是纯洁的，因为我没有回避向你们宣告。}（第26-27节。）那么，这就是做教师的本分：宣告、宣讲、教导、毫不退缩、昼夜劝勉；但如果一个人做了这一切，却毫无结果，你知道后果是什么。于是你有另一个辩护：\Quote{我在众人的血上是纯洁的。}不要以为这些话只对我们说：因为这讲话也是对你说的，叫你留心所讲的事，不要逃避听讲。我能做什么呢？看！我每天声嘶力竭地呼喊：\Quote{离开剧院：}许多人嘲笑我们：\Quote{停止发誓、停止贪婪：}我们劝诫无数，却无人听。但我不是在夜间讲道吗？我多么希望在夜间也这样做，在你们的餐桌上，如果一个人可能被分成一万块，以便与你们同在并讲道。但如果一周一次我们呼唤你们，你们却退缩，有些甚至不来这里，而那些来的人，离开时毫无收获——我们该怎么办？我知道许多人甚至嘲笑我们，说我们总是讲同样的事：我们因同样的内容多次重复而令你们厌烦。但这不怪我们，而应归咎于听众自己。因为那真正进步的人，总是乐于听同样的事；仿佛听到的是对他的赞美；但那不愿上进的人，甚至感到烦恼，即使只听了两次，也好像听了许多次。\par

他说：\BibleQuote{我无罪于众人的血。}\BibleRef{宗 20:26}这话\Person{保禄}说来合适，但我们却不敢说，因为我们自知有无数过犯。因此，对他那样时刻警惕、随时待命、为门徒的救恩忍受一切的人，说这话是合适恰当的；但我们却必须说\Person{梅瑟}的那句话：\BibleQuote{上主因你们的缘故向我发怒}\BibleRef{申 3:26}，因为你们也使我们陷入许多罪恶。因为当我们看到你们没有进步而灰心时，我们的大部分力量岂非被击倒？因为，试问，发生了什么事？看！靠着天主的恩宠，如今我们也过了三年的时间，虽非日夜劝诫你们，却常常每隔三天或七天这样做。又有什么结果呢？我们责备、斥责、哭泣、痛苦，虽不公开，却在心里。而那内心的泪水，远比外在的泪水更苦：因为外在的泪水确实给忧伤者带来某种缓解，而内在的泪水却加重忧伤，并将它紧紧束缚。因为当有悲伤的原因，却无法发泄，以免被人指责为好虚荣时，想想他受的是什么苦！若不是人们会指责我过分爱表现，你们会看见我每天泪流成泉：但我的卧室和独处的时刻可以作证。因为请相信我，我有时对自己的救恩感到绝望，但因了为你们哀悼，我甚至无暇哀叹自己的不幸：所以你们完全是我的一切。而且，无论我察觉你们在进步，那时因着极大的喜悦，我感不到自己的不幸；或是我见你们没有进步，因着那样的悲伤，我又抛开自己的忧虑：因着你们的好事而喜悦，尽管我自己有无数恶事；因着你们的痛苦而悲伤，尽管我自己有无数成功。因为当他的羊群被毁灭时，教师还有什么希望呢？他有什么样的生活，什么样的期待呢？他将以什么样的信心站在天主面前？他将说什么？因为假设他没有任何可被指控的，不必受罚，而是\BibleQuote{无罪于众人的血}，即便如此，他也会遭受无法治愈的悲伤：因为父亲们，即使不必为儿女的罪负责，也仍有悲伤和烦恼。而这于他们毫无益处，也不能保护他们（\OriginalTerm{保护}{\GreekText{προΐσταται}}）。\BibleQuote{因为他们为了你们的灵魂而警醒，如同要交账的人。}\BibleRef{希 13:17}这似乎是一件可怕的事：但对我来说，在你们的毁灭之后，这已不令我担忧。因为无论我是否交账，对我都没有益处。但愿你们得救，而我因你们而交账：你们得救，而我因没有尽我的本分而被责备！因为我的焦虑不是要你们通过我作为工具而得救，而只是要你们得救，无论通过什么人作为工具。你们不知道属灵生育的剧痛，何等强烈；怀有这种生育的人，宁愿自己被切成千万块，也不愿看到自己所生的人中有一个灭亡、败坏。我们用什么说服你们呢？确实没有别的证据，但凭着所做的事，在一切关于你们的事上，我们将为自己辩护。我们也能说，在一切事上我们\BibleQuote{没有退缩不向你们宣讲}全部真理：然而我们仍然悲伤；我们确实悲伤，这从我们制定的无数计划和策略可以明显看出。然而我们可能对自己说：这与我何干？我已经尽了我的本分，\BibleQuote{我无罪于他们的血}：但这不足以安慰。如果我们能剖开我们的心，将它展示给你们，你们会看到它里面何等宽广地容纳了你们，包括妇女、儿童和男人；因为爱德的力量如此之大，它使灵魂比天空更广阔。\BibleQuote{你们要接纳我们，} \Person{保禄}说：\BibleQuote{我们没有亏负过谁，你们在我们心中并不狭窄。}\BibleRef{格后 7:2}他将整个\Place{格林多}放在心里，并说：\BibleQuote{你们并不狭窄；你们也要心胸开阔。}\BibleRef{格后 6:13}但我自己不能说这话，因为我深知你们爱我，也接纳我。但我的爱或你们的爱有什么益处呢？要是属于天主的事在我们中间不兴旺？这更是悲伤的理由，是更糟损害（\OriginalTerm{损害}{\GreekText{λύμης}}，异文作\OriginalTerm{悲伤}{\GreekText{λύπης}}）的缘由。我没有什么可指责你们的：\BibleQuote{因为我为你们作证，如果可能的话，你们早已挖出自己的眼睛，送给我了。}\BibleRef{迦 4:15}\BibleQuote{我们不但愿意将天主的福音给你们，连自己的性命也愿意给你们。}\BibleRef{得前 2:8}我们被爱，我们也爱你们：但问题不在这里。让我们爱基督，\BibleQuote{因为第一条诫命是：你当爱主你的天主；第二条也相似：你当爱近人如你自己。}\BibleRef{玛 22:37}我们有了第二条，还需要第一条：需要第一条，极其需要，无论是我还是你们。我们有它，但不如所应有的。让我们爱祂：你们知道爱基督的人将获得多么大的赏报：让我们以灵魂的热忱爱祂，这样，藉着享有祂的恩惠，我们就能逃脱今生的惊涛骇浪，并堪当获得那为爱祂的人所预许的美好事物，藉着祂独生子、我们的主耶稣基督的恩宠和慈爱，愿光荣、权能、尊崇归于祂和圣父及圣神，从现在直到永远，万世无疆。阿们。\par

\SourceRecord{Translated by J. Walker, J. Sheppard and H. Browne, and revised by George B. Stevens.；Nicene and Post-Nicene Fathers, First Series；Vol. 11.；Edited by Philip Schaff.；Buffalo, NY: Christian Literature Publishing Co.,；1889.；<http://www.newadvent.org/fathers/210144.htm>.}

% structure: p0001 source=h1 target=section reason=page-title

\section{\WorkTitle{论宗徒大事录·讲道集第四十五篇}}

\BibleRef{宗 20:32}\par

\BibleQuote{现在，弟兄们，我把你们交托给天主和祂恩宠的圣言，祂能建立你们，并使你们在一切被圣化的人中得享嗣业。}\par

他在书信中所做的，在会议中说话时也同样做：从劝诫开始，以祈祷结束。因为当他引用\BibleQuote{凶暴的豺狼要进到你们中间}\BibleRef{宗 20:29}来使他们大为惊慌时，因此，为了不压倒他们，使他们失去镇定，请看他所给的安慰。他说，如同往常，\BibleQuote{现在，我把你们，弟兄们，托付给天主和他恩宠的道。}也就是，托付给祂的恩宠：是恩宠拯救。他不断提醒他们恩宠，使他们作为负债者更加热诚，并劝他们要有信心。\BibleQuote{这能建立你们。}他没有说建立，而是\Quote{建立起来}，表明他们已经被建立了。然后他提醒他们将来的希望，他说，\BibleQuote{赐给你们产业，在一切被圣化的人中。}再次劝诫：\BibleQuote{我没有贪图过任何人的金银或衣服。}\BibleRef{宗 20:33}他除去了万恶之根，贪爱金钱。他说\Quote{金银}，他没有说我没有拿，而是连\Quote{贪图}也没有。这并不是什么大事，但后面的事才是大事。\BibleQuote{你们自己也知道，这些手供应了我的需要，也供应了同我在一起的人。我凡事给你们示范，这样劳苦，你们应当扶助软弱的人。}\BibleRef{宗 20:34-35}看他劳碌工作，不仅是工作，而且是辛苦劳作。\BibleQuote{这些手供应了我的需要，也供应了同我在一起的人。}以此让他们羞愧。看他是多么配得他们。因为他没有说，你们应当表现出超越金钱，而是说什么？\BibleQuote{扶助软弱的人}——不是不加分辨地——\BibleQuote{并要听主所说的话：施予比领受更为有福。}因为恐怕有人以为这话是对他们说的，他把自己作为榜样，正如他在别处所说，\BibleQuote{给你们作榜样}\BibleRef{斐 3:17}，他补充了基督的话，主说，\BibleQuote{施予比领受更为有福。}他一边劝诫一边为他们祈祷：他既用行动表明——\BibleQuote{他说完这话，就跪下同众人一起祈祷，}\BibleRef{宗 20:36}——他不仅是祈祷，而且满怀着\OriginalTerm{深切的感动}{\GreekText{κατανύξεως}}：这是极大的安慰——又用说话表明，\BibleQuote{我把你们托付给主。}他们就大哭起来，抱着\Person{保禄}的颈项亲吻，最令他们伤心的，是他所说的话：他们将不再见到他的面容。他曾说，\BibleQuote{凶暴的豺狼要进来；}他曾说，\BibleQuote{我于众人的血是纯洁的；}然而最令他们伤心的却是这句，\Quote{他们将不再见到他，}因为正是这一点使战争变得惨烈。\BibleQuote{他们陪送我们，}经上说，\BibleQuote{直到船上。且发生了这样的事：我们与他们难分难舍，}——他们多么爱他，对他多么有感情——\BibleQuote{就开了船，直航到了\Place{科斯}，第二天到了\Place{罗得}，从那里到了\Place{帕塔辣}；并找着了一只往\Place{腓尼基}去的船，我们便上船出发。我们看见\Place{塞浦路斯}，就从左边驶过，航行到\Place{叙利亚}，在\Place{提洛}靠了岸}\BibleRef{宗 21:1}：他来到\Place{吕基雅}，离开\Place{塞浦路斯}后，航行到\Place{提洛}——\BibleQuote{因为船要在那里卸货。我们找着了门徒，就在那里住了七天；他们藉着圣神告诉\Person{保禄}，不要上\Place{耶路撒冷}去。}\BibleRef{宗 21:4}他们也预言了苦难。这是如此安排，使他们也说出这些话，以免有人认为\Person{保禄}无故说那些话，只是出于夸耀。在那里他们又彼此祈祷告别。\BibleQuote{过了这几天，我们就出发前行；他们众人同妻子儿女送我们直到城外。我们跪在岸上祈祷，彼此辞别后，我们上了船，他们就回家去了。我们从\Place{提洛}航行到了\Place{仆托肋买}，问候了弟兄们，同他们住了一天。第二天，我们\Person{保禄}的同伴出发，来到了\Place{凯撒勒雅}，进了传福音者\Person{斐理伯}的家，他是七人之一，同他住了。}\BibleRef{宗 21:5}到了\Place{凯撒勒雅}，经上说，我们同\Person{斐理伯}住下了，他是七人之一。\BibleQuote{这人共有四个女儿，都是童女，说预言。}\BibleRef{宗 21:9}但向\Person{保禄}预告的不是她们，虽然她们是女先知，而是\Person{阿加波}。\BibleQuote{我们住了多日，有一个先知从\Place{犹太}下来，名叫\Person{阿加波}。他来到我们这里，拿起\Person{保禄}的腰带，绑了自己的手和脚，说：圣神这样说：\Place{犹太人}要在\Place{耶路撒冷}这样捆绑这条腰带的主人，并把他交在外邦人手中。}\BibleRef{宗 21:10-11}那位从前预言饥荒的，他同一个人说：\BibleQuote{这条腰带的主人，将要这样被捆绑。}\BibleRef{宗 11:28}先知们常用这样的方式，将事件显现在眼前，当他们谈论被掳时——就如\Person{厄则克耳}所做的——这\Person{阿加波}也照样做了。这件事最令人痛心的部分是，\BibleQuote{把他交在外邦人手中。我们听了这些话，同那地方的人一起求他不要上\Place{耶路撒冷}去。}\BibleRef{宗 21:12}许多人甚至求他不要离开，但他仍不肯依从。\BibleQuote{\Person{保禄}回答说：你们为什么哭泣，使我心碎呢？}\BibleRef{宗 21:13}你明白吗？恐怕你们听见了那话，\BibleQuote{我因圣神被捆绑}\BibleRef{宗 20:22}，就以为这是出于不得已，或者他不知不觉陷进去，因此这些事被预告了。但他们哭泣，他安慰他们，为他们的眼泪而悲伤。因为他说：\BibleQuote{你们为什么哭泣，使我心碎？}再没有比这更充满爱意的了：因为他看见他们哭泣，他就悲伤，他自己对所受的试炼却毫不在意。\BibleQuote{我为主耶稣的名，不但被捆绑，就是死在\Place{耶路撒冷}也是愿意的。他既不肯听从，我们便不再劝，说：愿主的旨意成就。}\BibleRef{宗 21:13-14}你们这样做是亏负了我：难道我悲伤吗？他们便不再劝了，当他说\Quote{使我心碎}时。他说，我为你们哭泣，不是为自己的苦难；至于那些人，我甚至愿意为他们死。但让我们再回顾一下以上所说的话。\newline{}\par

（摘要。）\Quote{银子、金子或衣服，}等等。\EditorialNote{参见宗 20:33,34；格前 9；格后 11}可见，这事不仅发生在格林多——那些败坏门徒的人，也在亚细亚发生。但他在写信给厄弗所人时，从未以此责备他们。为什么？因为他没有遇到必须谈及此事的话题。但面对格林多人，他说：\BibleQuote{我的夸耀在阿哈雅地区没有被阻止。}\BibleRef{格后 11:10}他并没有说：你们没有给我；而是说：\Quote{我没有贪图银子、金子或衣服}，以免显得是他们的过失，因为他们没有给。他也没有说：我没有向任何人贪图生活必需品，以免再次显得是责备他们；但他巧妙地暗示了这一点，因为他不仅为自己，也为他人提供生计。看他如何热切地工作，\Quote{昼夜}讲论（对他人），\BibleQuote{含泪劝诫你们每一个人。}\BibleRef{宗 20:31}（这里）他又使他们畏惧：\Quote{我给你们显示了这一切，}他说：你们不能以无知为借口：\Quote{用行动向你们显示}，\Quote{应当这样辛劳，扶助软弱的人。}他并没有说接受是不好的，而是说不接受更好。因为他说：\Quote{你们要记住主耶稣说的话：施予比领受更有福。}\BibleRef{宗 20:35}这话在何处说的？也许是宗徒们以口传传统传下来的；或者从（记载的言论）中可以推断出来。事实上，在这里他既展示了面对危险的勇气，对属下的同情，勇敢的教导，谦卑，自愿的贫穷；但这里所展现的甚至超过了贫穷。因为如果他在（福音中）说：\BibleQuote{你若愿意成为成全的，去变卖你所有的，施舍给穷人}\BibleRef{玛 19:21}，那么当他不仅自己一无所取，还为他人提供生计，谁能与之相比？第一等是舍弃自己的财产；第二等是足以自给；第三等是还能供给他人；第四等是身为宣讲者并有权接受供给的人。所以他远胜于那些仅仅舍弃财产的人。\Quote{所以应当这样辛劳，扶助软弱的人：}这确实是同情软弱者；因为轻易地用他人的劳苦施与是容易的。\Quote{他们痛哭，扑在保禄的颈项上，}记载说，\Quote{和他亲吻。}\BibleRef{宗 20:37}他以\Quote{扑在他的颈项上}来显示他们的情感，仿佛作最后的拥抱，因他的讲论而产生的爱如此之大，爱情的魔力如此束缚他们。因为如果我们仅仅因彼此分离就叹息，虽然知道还会重逢，那么他们分离时的痛苦该是何等！我想保禄也流泪了。\Quote{我们勉强离开他们，}他说：用\Quote{勉强离开}显示这种分离的强烈。而且有理，否则他们决不可能出海。\Quote{我们直航来到科斯}是什么意思？换句话说，\Quote{我们没有绕行，也没有在其他地方停留。}然后\Quote{到了罗德。}\BibleRef{宗 21:1}看他如何赶路。\Quote{找到一只船要往腓尼基去，}\BibleRef{宗 21:2}可能他们乘坐的那只船在那里停泊；所以他们换乘另一只，因为没有找到去凯撒勒雅的船，而找到了去腓尼基的，就上了船继续航行，也离开了塞浦路斯和叙利亚；但\Quote{从左边绕过}不是说绕行，而是表示他们加快速度，也不到叙利亚去。\Quote{我们在提洛靠了岸。}\BibleRef{宗 21:3}然后与弟兄们同住了七天。现在他们接近耶路撒冷，不再奔跑。（\Emph{b}）\Quote{他们借着圣神告诉保禄，不要上耶路撒冷。}\BibleRef{宗 21:4}请注意，当圣神不禁止时，他就顺从。他们说：\Quote{不要冒险进入剧场，他就不去}\BibleRef{宗 19:31}；他们多次救他脱离危险，他就顺从；又一次他经窗户逃走；现在，尽管无数人恳求他，包括在提洛和凯撒勒雅的人，流泪并预言无数危险，他拒绝顺从。然而他们不仅预言了危险，而且\Quote{借着圣神说了。}如果圣神吩咐，他为何反对？\Quote{借着圣神，}即他们\Quote{借着圣神}知道了后果，就对他说：当然这不意味着他们的劝告是出于圣神。因为他们不仅（借着圣神）预言危险，而且（自己加上）他不该上去——出于爱惜他。但\Quote{我们过了这些日子，}即满了所定的日子，\Quote{就离别前行，他们众人同妻子儿女送我们到城外。}\BibleRef{宗 21:5}——请看恳求是多么强烈。他们又以祈祷告别。在托勒买也只住了一天，但在凯撒勒雅住了多日。\BibleRef{宗 21:6}（\Emph{a}）现在他们接近耶路撒冷，不再匆忙。因为请看，我请你们注意所有的天数。\Quote{过了无酵节，}他们\Quote{五天到了特洛阿}\BibleRef{宗 20:6}；在那里住了\Quote{七天；}总共十二天；然后到\Quote{塔索斯}、\Quote{米提肋乃}、\Quote{特洛基利翁}和\Quote{对岸的希约}、\Quote{撒摩}和\Quote{米肋托}\BibleRef{宗 20:13}；总共十八天。然后到\Quote{科斯}、\Quote{罗德}、\Quote{帕塔辣}，二十一天；再到\Quote{提洛}约五天；二十六天；在那里住了\Quote{七天}；三十三天；\Quote{托勒买}，三十四天；然后到\Quote{凯撒勒雅，住了多日}\BibleRef{宗 21:1}；然后此后，先知从那里促他们起行。（\Emph{c}）当保禄听说要受无数危险时，他就急忙，不是投身于危险，而是认为这是圣神的命令。（\Emph{e}）阿加波并没有说：\Quote{他们要捆绑}保禄，以免显得与保禄有默契，而是说：\Quote{这腰带的主人}\BibleRef{宗 21:11}——如此看来，他也有腰带。但当他们无法说服他——这就是为何他们哭泣——然后\Quote{他们住了口。}你注意到顺服吗？你注意到情感吗？\Quote{他们住了口，}记载说，\Quote{说：愿主的旨意成就。}\BibleRef{宗 21:12}（\Emph{g}）他们说：主自己会成就他眼中看为善的事。因为他们觉察到这是天主的旨意。否则保禄不会如此坚决——他向来在一切（其他）场合脱离危险。（\Emph{d}）\Quote{过了这些日子，}记载说，\Quote{我们收拾了行李}——即接受了旅行所需的（供应）——\Quote{便上耶路撒冷去。}\BibleRef{宗 21:15}\Quote{有凯撒勒雅的几位门徒也和我们同去，带我们到一个久为门徒的塞浦路斯人纳松家里，要我们与他同住。}\BibleRef{宗 21:16}\Quote{到了耶路撒冷，弟兄们欢欢喜喜地接待了我们。}\BibleRef{宗 21:17}（\Emph{f}）\Quote{带我们到，}记载说，\Quote{要与他同住}——不是到教会：因为此前一次\BibleRef{宗 15:4}，他们为法令上耶路撒冷时，住在教会里，但如今住在一个\Quote{久为门徒的}家里。这个说法表明宣讲已经进行了很长时间：由此我认为这《宗徒大事录》的作者概括了多年的历史，只叙述了最重要的事。（\Emph{h}）他们是如此不愿加重教会的负担，既然有别人接待他们；他们也如此不坚持自己的尊严。\Quote{弟兄们，}记载说，\Quote{欢欢喜喜地接待了我们。}犹太人中的事务现在十分和平：没有多少战争。\Quote{带我们到，}记载说，\Quote{要与他同住。}保禄是他所款待的客人。你们中也许有人说：是的，如果让我款待保禄作客，我会乐意并十分热切地这样做。看！你能够款待保禄的主作客，却不愿意：因为\BibleQuote{凡因我的名字接待一个这样的小孩子，}他说，\Quote{就是接待我。}\BibleRef{玛 18:5}兄弟越是最小的，基督越是通过他来到你这里。因为接待伟人，常常也出于虚荣；但接待卑微的人，纯粹是为了基督。你甚至能够款待基督的父作客，却不愿意：因为\BibleQuote{我作客，你们收留了我}\BibleRef{玛 25:35}；又说：\BibleQuote{你们为我最小的弟兄中的一个所做的，就是为我做的。}\BibleRef{玛 25:40}即使不是保禄，只要是一个信徒和兄弟，即使是最小的，基督也通过他来到你这里。打开你的家，接待他。\BibleQuote{谁接待一位先知，}他说，\Quote{必得先知的赏报。}\BibleRef{玛 10:41}因此，谁接待基督，必得那以基督为客的人的赏报。你不要不信他的话，而要相信。他自己说：通过他们，我到你这里来；为使你不致不信，他既为不接待者设立了惩罚，也为接待者设立了荣誉；因为若非他本人就是被尊荣和被侮辱者，他不会这样做。\Quote{你接待了我到我的住处，我要接待你到我父的国里；你解了我的饥饿，我解除你的罪过；你看见我被捆绑，我看见你被释放；你看见我作客，我使你成为天堂的公民；你给了我饼，我给你整个国度，使你承受并拥有它。}他没有说\Quote{接受}，而是说\Quote{承受}，这个词用于那些因所有权而拥有的人；如同我们说：\Quote{我承受了这。}你在暗中为我做了，我要公开宣扬：论到你的行为，我说是出于慷慨，但我的行为是出于债务。\Quote{因为你，}他说，\Quote{开了头，我就跟随而来：我不以承认我所受的恩惠为耻，也不隐瞒你救了我脱离什么——饥饿、赤身露体和漂泊。你看见我被捆绑，你必不见地狱之火；你看见我患病，你必不见折磨和惩罚。}哦，真正有福的手，服事于这样的善工，堪得服事基督！为基督的缘故进入监狱的脚，不畏火刑：没有捆绑的考验，那看见他被捆绑的手！你给他穿上衣服，你就穿上了救恩的衣服；你与他一同在监狱，你就与他一同在国度里，不以为耻，知道你去探望了他。圣祖不知道他在款待天使，却款待了他们。\BibleRef{创 18:3}我恳求你们，让我们羞愧：他正午坐在那里，在异地，没有产业，\Quote{连立足之地也没有}\BibleRef{宗 7:5}；他是外方人，外方人款待了外方人：因为他是天堂的公民。因此，即使在地上，他也并非外方人。我们比那位外方人更是外方人，如果我们不接待外方人。他没有家，他的帐幕就是接待之处。注意他的慷慨——他宰了一头牛犊，和了细面；注意他的热忱——他和他的妻子亲自操劳；注意他的谦逊——他俯伏请安。因为所有这些品质都应当在那个款待外方人的人身上——热忱、喜乐、慷慨。因为外方人的灵魂是羞愧的，感到不好意思；除非主人表现出极大的喜乐，他就会觉得被轻视而离开，这样接待比不接待更糟。所以他俯伏请安，用言语欢迎，用座位款待。因为谁会犹豫，知道这工作是为他做的？\Quote{但我们不是在异地。}如果我们愿意，就能效法他。有多少弟兄是外方人？有一个公共的房间，教会，我们称之为\Quote{接待室}。要好问\OriginalTerm{好问}{\GreekText{περιεργάζεσθε}}，坐在门口，亲自接待来者；即使你不愿意接他们到自己家里，至少以其他方式（接待他们），供应他们所需。\Quote{难道教会没有资源吗？}你会说。她有：但这与你何干？他们从教会公款中得到供养，这对你有好处吗？如果另一个人祈祷，你就不必祈祷吗？你为什么不说：\Quote{司铎们不祈祷吗？为什么我要祈祷？}\Quote{但我，}你会说，\Quote{给那不能在那里被接待的人。}给，即使只给那一个人：因为我们所关心的是，你无论如何要施与。听保禄说什么：\BibleQuote{好救济那些真正寡居的人，免得教会受累。}\BibleRef{弟前 5:16}无论你怎样做，只要做。但\Emph{我}说，不是\Quote{免得教会受累}，而是\Quote{免得你受累}；因为按这种逻辑，你会什么都不做，把一切都推给教会。这就是为什么教会设立了一个公共房间，免得你说这些话。\Quote{教会，}你说，\Quote{有土地，有钱，有收入。}难道她没有开支吗？我问；她没有日常支出吗？\Quote{无疑，}你会说。那么你为什么不帮助她有限的财力？我实在羞于说这些话；然而，若有人以为我所说的是为了利益，我并不强迫谁。在你的家里为自己造一个客房：在那里放一张床，放一张桌子和一盏灯台。\BibleRef{列下 4:10}因为士兵来了，你有为他们预备的房间，关心他们，供应一切，因为他们为你抵挡这世界有形的战争；而外方人却没有可住的地方，这岂不是荒谬吗？胜过教会。你想羞辱我们吗？这样做：在慷慨上超过我们：有一个房间，基督可以来；说：\Quote{这是基督的房间；这是为他设立的居所。}即使是地下室，简陋不堪，他也不嫌弃。\Quote{赤身露体、作客，}基督四处游走，他只需要一个栖身之所：即使只这一点，也给他。不要没有同情心，不要不人道；不要如此热衷于世事，而在灵性上如此冷漠。让你最忠心的仆人也负责此事，让他领进瘸腿的、乞丐和无家可归的人。我说这些是为了让你们羞愧。因为你们本当在家里最好的部分接待他们；但如果你们不愿这样做，即使在下层，即使在你养骡和仆人的地方，也在那里接待基督。也许你听了这话发抖。那么当你连这个也不做时又如何？看，我劝勉，看，我命令你们；让这成为一个要认真对待的事。但你们也许不愿这样？用其他方式做。有许多穷男人和穷女人：指定一个（这样的人）常住在那里：让穷人（作为你的同住者）如同你家的守卫：让他成为你的墙和篱笆，盾牌和枪。哪里有施舍，魔鬼就不敢靠近，也没有任何邪恶之物。我们不要忽视如此大的收益。但现在，有地方给马车，有地方给轿子\OriginalTerm{轿子}{\GreekText{βαστερνίοις}}；但给漂流无定的基督，连一个也没有！亚巴郎在他自己居住的地方接待了外方人；他的妻子站在仆人的地位，客人站在主人的地位。他不知道他在接待基督；不知道他在接待天使；所以若他知道，他会倾其所有。但我们，知道我们在接待基督，却不及那以为接待人的所表现的热忱。\Quote{但他们中有许多是骗子，不知感恩的人。}你会说。正因为这样，你的赏报更大，因为你为基督的名而接待。因为若你真知道他们是骗子，就不要接他们进你的家；但若你不知道，为什么轻易地指责他们？\Quote{因此我告诉他们去接待院。}但我们有什么借口，当我们连认识的人也不接待，反而对所有人关上大门？让我们的家成为基督的公共接待处：我们向他们要求的报酬，不是金钱，而是他们使我们的家成为基督的接待处；让我们四处奔跑，拖他们进来，夺取我们的战利品：我们所领受的恩惠大于我们所施与的。他没有叫你宰牛犊：给饥饿的人饼，给赤身露体的人衣服，给外方人住处。但为了你不以此为借口，有一个公共房间，就是教会的；把你钱投进去，这样你就接待了他们：因为（亚巴郎）也有通过仆人之手所做之事的赏报。\BibleQuote{他把牛犊交给仆人，仆人赶快去预备。}\BibleRef{创 18:7}他的仆人也是如此训练有素！他们跑，不像我们的仆人那样抱怨：因为他使他们虔诚。他带领他们出战，他们没有抱怨：如此有纪律。\BibleRef{创 14:14}因为他同等关心一切人，如同关心自己；他几乎像约伯那样说：\BibleQuote{我们原是在同一个胎中形成的。}\BibleRef{约 33:6}因此，我们也要为他们的得救着想，并把照顾仆人当作我们的职责，使他们成为善人；让我们的仆人也受到有关天主之事的教导。那么我们行善就不难了，如果我们有序地训练他们。正如在战争中，士兵纪律严明，将军就容易作战；但若不然，正相反；当水手也同心时，舵手就容易掌握舵绳；同样在这里。比如，如果你的仆人受过这样的训练，你就不易被激怒，不必指责，不会发怒，不需谩骂他们。甚至你可能敬畏你的仆人，如果他们值得赞扬，他们会成为你的助手，给你良好的建议。但从这一切中，一切使天主喜悦的事都将发生，这样全家将充满祝福，我们做了中悦天主的事，将从上主那里获得丰富的援助，愿我们众人皆能达到，靠我们的主耶稣基督的恩宠和仁慈，愿光荣、权能、尊崇归于他及父和圣神，从今世直到永远。阿们。\par

\SourceRecord{Translated by J. Walker, J. Sheppard and H. Browne, and revised by George B. Stevens.；Nicene and Post-Nicene Fathers, First Series；Vol. 11.；Edited by Philip Schaff.；Buffalo, NY: Christian Literature Publishing Co.,；1889.；<http://www.newadvent.org/fathers/210145.htm>.}

% structure: p0001 source=h1 target=section reason=page-title

\section{论宗徒大事录讲道第四十六篇}

\BibleRef{宗 21:18,19}\par

\BibleQuote{次日，\Person{保禄}同我们去见\Person{雅各伯}，所有的长老也都到了。他向他们问候后，便逐一细述\Person{天主}藉他的职务在外邦人中所行的事。}\par

这位是耶路撒冷的主教；保禄从前被派到他那里。这位（雅各伯）是主的兄弟，一位伟大而可钦佩的人。（默思他，经上说：）\Quote{保禄与我们一同进去。}注意主教的谦逊行为：\Quote{还有长老们}（在场）。保禄又向他们述说有关外邦人的事，不是出于虚荣——天主禁止——而是愿彰显天主的仁慈，并使他们充满大喜乐。\EditorialNote{宗 15} 请随之看：\Quote{他们听见了，}它说，\Quote{就光荣天主，}——不是称赞或钦佩保禄：因为他叙述的方式如此，将一切归于祂——\Quote{并对他说：\InnerQuote{兄弟，你看，犹太人中相信的有多少千。}}注意他们说话时的谦逊礼貌：\Quote{他们对他说：}不是（雅各伯）以主教身分权威地发言，而是他们把保禄视为同僚，与他意见一致；\Quote{兄弟，你看：}仿佛立刻且一开始就为自己道歉，并说：\Quote{我们并不愿如此。你看到事情的必要了吗？\InnerQuote{有多少千，}他们说，\InnerQuote{犹太人中有的}已经聚集了。}他们并不说：\Quote{我们使多少千成了慕道者，}而是说：\Quote{有的。而这些，}他们说，\Quote{都对于法律很热忱。} \BibleRef{宗 21:20} 有两个理由——他们的人数，以及他们的看法。因为即使他们数量不多，轻视他们也不对；即使他们很多，若并非都固守法律，也无需过于在意。然后还给出了第三个原因：\Quote{他们都，}它说，\Quote{已经得知你的事}——他们不说\Quote{听见了，}而是\OriginalTerm{他们受了教}{\GreekText{κατηχήθησαν}}，即他们已相信并受了教导，\Quote{就是你教导所有住在外的犹太人背弃梅瑟，叫他们不要给孩子行割损，也不要遵守习俗。} \BibleRef{宗 21:21} \Quote{那么怎么办呢？群众必须聚集，因为他们会听说你来了。所以你就照我们对你说的去做。}（22-23节）：他们说这些是建议，而非命令。\Quote{我们有四个人，身上都有愿；你带他们，并与他们一同行洁礼，替他们付出费用。}用行动而非言语来辩护——\Quote{使他们剃头，}它说，\Quote{所有人都知道，他们对你所传言的那些事都是子虚乌有，而你本人也循规蹈矩，遵守法律。}（23-24节）：他们不说\Quote{教导，}而是格外地说：\Quote{你本人也遵守法律。}因为当然，关键不在于他是否教导别人，而在于他自己是否遵守法律。\Quote{那么，}（他或许会说，）\Quote{如果外邦人知道了，我就会伤害他们。}怎么会呢？他们说，因为就连我们这些犹太人的教师，也给他们写信了。\Quote{至于信主的外邦人，我们已经写信议定，叫他们遵守这些事：只叫他们戒避祭偶像之物、血、勒死的牲畜和奸淫。} \BibleRef{宗 21:25} 此处带著一种劝诫的口吻（\OriginalTerm{劝诫地}{\GreekText{ἐντρεπτικὥς}}），正如他们说：\Quote{我们，}虽然我们是给犹太人的宣讲者，却命令了他们；那么你，虽然给外邦人的宣讲者，也该与我们合作。观察保禄：他并不说：\Quote{好吧，但可以提出提摩太，我曾给他行割损：好吧，但可以用我要说的话使他们满意。}而是顺从，做了这一切：因为实际上这样做是合宜的。因为采取有效措施为自己辩护是一回事，而在任何人都不知道的情况下做了这些事又是另一回事。甚至由他承担费用，这一步无可怀疑。\Quote{于是保禄带着那些人，第二天与他们一同行了洁礼，进了圣殿，宣布洁净期已满，直到为每人献上供物。} \BibleRef{宗 21:26} \Quote{宣布，}\OriginalTerm{公开宣告}{\GreekText{διαγγέλλων}}，即\OriginalTerm{宣告}{\GreekText{καταγγέλλων}}，公开通知：所以是他使自己显眼。\Quote{当七天将满时，从亚细亚来的犹太人}——因为（他的到达）大多与他们的时间吻合——\Quote{看见他在圣殿里，就煽动所有群众，下手拿住他，喊着说：以色列人，帮忙！这就是那在各处教训众人反对百姓、反对法律、反对这地方的人；而且还将希腊人带进圣殿，污秽了这圣地。}（27-28节）注意他们一贯的行径，我们随处可见他们多么骚动，无论有无理由都在当中吵闹。\Quote{因为他们先前看见厄弗所人特洛斐摩在城里与保禄在一起，就猜测保禄带进了圣殿。于是全城轰动，百姓一起跑来，抓住保禄，把他拉出圣殿，门立刻关了。}（29-30节）\Quote{以色列人，}它说，\Quote{帮忙！这就是那（教训）反对百姓、反对法律、反对这地方的人。}——最使他们困扰的事：圣殿和法律。保禄并不责怪宗徒们是他遭遇这些事的原因。\Quote{他们把他拉出圣殿，}它说，\Quote{关上了门。}因为他们想杀他，所以把他拖出来，以便更安全地下手。\Quote{当他们想要杀他时，有消息传到营里的千夫长那里，说全耶路撒冷都乱了。千夫长立即带着士兵和百夫长跑下去，冲到他们那里；他们看见千夫长和士兵，就停止了殴打保禄。于是千夫长上前拿住他，下令用两条铁链捆住他，并问他是什么人，做了什么事。群众中有的喊这个，有的喊那个。} \BibleRef{宗 21:31} 但千夫长下来救了他，并\Quote{下令用两条铁链捆住他：}（借此）平息群众的愤怒。\Quote{因不能确知实情，千夫长便下令把他带进营楼。当他走上台阶时，因群众的暴力，只得由士兵抬着；因为群众跟在后面喊着：除掉他！} \BibleRef{宗 21:34} \Quote{除掉他？}是什么意思？就是按罗马习俗所说的：把他拉到阵营去！\Quote{当保禄快要被带进营楼时，对千夫长说：我可以跟你说句话吗？} \BibleRef{宗 21:37} 在被抬上台阶时，他请求跟千夫长说话：注意他多么安静地做这事。\Quote{我可以跟你说句话吗？}他说。\Quote{千夫长说：你会说希腊话吗？那么你不是那先前作乱、带领四千凶徒往旷野去的埃及人吗？} \BibleRef{宗 21:38} 因为（这埃及人）是个叛乱的煽动者。关于这些，保禄为自己辩解，以及……\par

（综述。）\BibleQuote{因此，请做我们对你说的事，}等等（第23-24节）。他表明，\OriginalTerm{原则上}{\GreekText{προηγουμένως}}没有必要这样做——因此他们也得到了他的顺从——而是经济（上智安排）和迁就。\BibleQuote{至于外邦人，}等等。\BibleRef{宗 21:25}那么，这对宣讲并无妨碍，因为他们自己也为此立法。那么，在他责备伯多禄时，他并非\OriginalTerm{绝对地}{\GreekText{ἁπλῶς}}指控他做错了：因为他此时所做的，正是伯多禄彼时所做的，只是保持沉默，并确立了他的教义。\BibleRef{迦 2:11}他并没有说：为什么？教导外邦人是不对的。\BibleQuote{单单在那里还没有宣讲是不够的，还需要再多做一些，使那些人相信你遵守法律。这事属于迁就，不要惊慌。}他们没有更早地建议他这样做，直到他们先谈了经济（上智安排）和益处。\BibleQuote{而且，在耶路撒冷做这件事是能承受的。\InnerQuote{因此，你在这里做这件事}，以便你到外地时能做另一件事。}（\Emph{b}）\BibleQuote{第二天，}经上说，\BibleQuote{他带着他们，}\BibleRef{宗 21:26}：他没有拖延；因为当有经济（上智安排）时，这就是它的方式。（\Emph{a}）\BibleQuote{从亚细亚来的犹太人看见了他，}因为他们自然在那里住了几天，\BibleQuote{在圣殿里。}\BibleRef{宗 21:27}（\Emph{c}）注意此时显现的经济（上智安排）（\EditorialNote{第279页，注释1}）。在（相信的）犹太人被说服（关于他）之后，那些（亚细亚的）犹太人才攻击他，以便那些（相信的）犹太人不也攻击他。他们说：\BibleQuote{以色列人啊！}（救命！）仿佛是什么难以捕捉、难以征服的怪物落入了他们手中。\BibleQuote{所有的人，}他们说，\BibleQuote{他到处不住地教导；}不仅在这里。然后控诉更加因当前情况而加剧。\BibleQuote{更甚的是，}他们说，\BibleQuote{他玷污了圣殿，把希腊人带了进来。}\BibleRef{宗 21:28}然而在基督时代，有\BibleQuote{（希腊人）上来朝拜}\BibleRef{若 12:20}：不错，但这里指的是无意朝拜的希腊人。\BibleQuote{他们抓住了保禄，}等等。\BibleRef{宗 21:30}他们不再需要法律或法庭：他们也打了他。但他当时忍住不做辩护；他后来才辩护：有理由的；因为他们当时甚至不会听他。请问，他们为什么喊：\BibleQuote{除掉他？}\BibleRef{宗 21:36}他们怕他逃脱。看保禄如何顺从地对千夫长说话。\BibleQuote{我可以对你说话吗？那么你不是那个埃及人吗？}（第37-38节）。这个埃及人是个骗子和冒牌货，魔鬼期望通过他给（福音）投下阴影，并把基督和他的宗徒牵连到那些（冒牌货）的指控中：但他一无所成，反而真理更加辉煌，完全未被魔鬼的阴谋挫败，反而更加闪耀。因为如果没有冒牌货，而这些人（基督和他的宗徒）得胜，也许有人会抓住这一点；但当那些冒牌货确实出现时，这就是奇事。\BibleQuote{为了使，}（宗徒）说，\BibleQuote{那些经得起考验的人得以显明。}\BibleRef{格前 11:19}加玛里耳说：\BibleQuote{这些日子以前，特乌达起来。}那么，我们不要因异端存在而悲伤，因为假基督甚至在此以前和以后都试图攻击基督，意图遮蔽他，但每次我们都发现真理透明闪耀。先知们也是如此：有假先知，与他们对比，他们更加闪耀：正如疾病增强健康，黑暗增强光明，风暴增强平静。希腊人没有余地可说（我们的导师）是冒牌货和江湖骗子：因为那些（如此的人）被揭露了。梅瑟的情况也是如此：天主容忍了术士，特意使梅瑟不被怀疑是术士：他让他们教导所有人，法术在制造幻象方面能走多远：超越这一点，他们就不能欺骗，反而自己承认失败。冒牌货对\Emph{我们}无害，反而有益，如果我们用心思考。那么，你会说，如果我们和他们一起被普遍认为是一伙的呢？这评判不在我们中，而是在那些没有判断力的人中。我们不要过多关心多数人的评价，也不要过分在意。我们为天主而活，不是为人；我们在天上过我们的属灵生活，不是在地上；\Emph{那里}有我们劳苦的奖赏和赏报，我们从那里期待赞美，从那里期待冠冕。我们为人的忧虑到此为止——即我们不给他们把柄反对我们。但如果我们不提供把柄，那些人选择轻率无区别地控告我们，让我们笑，不要哭。\BibleQuote{你要在主和人面前，}你\BibleQuote{预备诚实的事}\BibleRef{格后 8:21}：如果你预备了诚实的事，那人讥笑，就不再为此担心。你在圣经中有榜样。因为，他说，\BibleQuote{我现在是讨人的喜欢，还是讨天主的喜欢？}\BibleRef{迦 1:10}又说：\BibleQuote{我们劝人，但在天主面前是显明的。}\BibleRef{格后 5:11}基督论及那些跌倒的人这样说：\BibleQuote{由他们吧！他们是瞎子领瞎子。}\BibleRef{玛 15:14}又说：\BibleQuote{几时众人都夸赞你们，你们是有祸的。}\BibleRef{路 6:26}又说：\BibleQuote{你们的光当在人前照耀，好使他们看见你们的善行，光荣你们在天之父。}\BibleRef{玛 5:16}又说：\BibleQuote{谁若使这些小子中的一个跌倒，倒不如拿一块驴拴的磨石，挂在他的颈上，沉在海的深处。}\BibleRef{玛 18:6}\par

这些言论并非矛盾，反而极其一致。因为当过犯在于我们时，我们就有祸了；但不在于我们时，就并非如此。再者，那藉着祂使\Quote{天主的名被亵渎。}\BibleRef{罗 2:24}的人有祸了。那么，如果我在某件事上做得对，但另一个人却亵渎，这与我何干？只与他有关，因为藉着他天主被亵渎了。\Quote{一个人怎么可能在某件事上做得对，却给其余的人留下把柄？}你们要我举例——从现今的时代，还是从古时？为了不轻易受惊（\OriginalTerm{不易受惊}{\GreekText{ψοφοδεεἵς}}），我们是否该直接谈论当前的问题？保禄在耶路撒冷犹太化了，但在安提约基雅却不然：他犹太化了，那些人就起了反感（第282页，注3），但那些人无权反感。据说他曾向尼禄的司酒官及其情妇致意：你们想，世人会因此怎样议论他？但他们无权如此。因为，如果他为了放荡或任何邪恶行为而吸引他们，那么反感是有道理的；但如果是为正直的生活，又有什么可反感的呢？让我提一件发生在我一位相识者身上的事。天主的愤怒曾降临到一座城，他当时非常年轻，属于执事品。主教那时不在，长老中也没有人关心此事，但他们不加分辨地在同一夜使大批群众同时受洗，那些人也不加分辨地接受了圣洗，全都一无所知：他将这些人分别召集，每百人或二百人一组，对他们讲道，不谈其他，只谈圣事，以致连未受洗的人也不得在场。许多人认为他这样做是贪图统治权。但他并不在意；然而他也没有继续做这件事很久，而是立刻停止了。那么，他是这绊倒人的原因吗？我认为不是。因为如果他无故这样做，人们或有理由归咎于他；同样，如果他继续这样做，也或有理由。因为当任何悦乐天主的事因他人反感而受阻时，我们便不应在意；但当我们不因他而不得不冒犯天主时，才是该留意的时候。譬如说，如果我们正在讲道，使醉汉羞愧（\OriginalTerm{嘲笑}{\GreekText{σκωπτόντων}}），有人因此反感——我就该停止讲道吗？请听基督说：\BibleQuote{你们也要走吗？}\BibleRef{若 6:67}因此，正确的做法是：既不要完全无视众人的软弱，也不要过分在意。我们岂不见医生也是如此：在可行的情况下，他们迁就病人的怪癖；但当放纵有害时，他们就不肯姑息。懂得中庸之道总是好的。许多人辱骂，因为一位美丽的贞女留了下来，他们便咒骂教导她的人。那么，他们因此就该停止吗？绝不。因为我们不仅要看是否有人反感，还要看他们是否理当反感，以及迁就他们是否对我们无害。保禄说：\BibleQuote{如果肉得罪我的弟兄，我就永远不吃肉。}\BibleRef{格前 8:13}这是有道理的：因为不吃肉对他无害。然而，如果我愿意\OriginalTerm{弃绝}{\GreekText{ἀποτάξασθαι}}世俗，这若使他反感，就不该在意他。你会问，这会使谁反感呢？据我所知，许多人。因此，当阻碍是无足轻重的事时，就让它去做。否则，如果我们只注意这一点，我们就有许多事要停止；反之，如果我们轻视一切反对，就会毁掉许多弟兄。正如保禄也预先顾虑到反感：\Quote{免得在这由我们经办的慷慨施舍中，}他说，因为消除恶意的猜疑对他并无损失。但当我们陷入这样的必要，即因他人反感而导致重大恶果时，我们就不要理会那人。他只能怪自己，我们不再负责，因为不可能在不伤害自己的情况下宽容他。有些人因某些信徒在异教庙宇中坐席而反感。坐席是不对的，因为不做此事并无害处。他们因伯多禄与外邦人同食而起反感。但他确实宽容了他们，但保禄却不如此。在所有情况下，我们理当遵行天主的法律，竭力不给人留下反感的把柄；这样我们自己不必承担责任，并能从天主那里获得怜悯，藉着祂独生子的恩宠和慈爱，愿荣耀、权能、尊崇归于圣父及圣神，自今至永远，永世无尽。阿们。\par

\SourceRecord{Translated by J. Walker, J. Sheppard and H. Browne, and revised by George B. Stevens.；Nicene and Post-Nicene Fathers, First Series；Vol. 11.；Edited by Philip Schaff.；Buffalo, NY: Christian Literature Publishing Co.,；1889.；<http://www.newadvent.org/fathers/210146.htm>.}

% structure: p0001 source=h1 target=section reason=page-title

\section{\WorkTitle{宗徒大事录}讲道 第47篇}

\BibleRef{宗 21:39-40}\par

\Quote{\Person{保禄}却说：\InnerQuote{我是犹太人，是\Place{基里基雅}的\Place{塔尔索}人，并不是一个无名小城的公民。我求你，准许我向百姓讲话。}得到许可后，\Person{保禄}就站在台阶上，向百姓挥手，大家就非常安静。他便用希伯来话对他们说：}\par

请注意，当他向外人讲论时，并不拒绝利用法律所提供的帮助。此处他以自己城邑的名字令千夫长心生敬畏。而在别处他又说：\Quote{我们本是罗马人，并没有被定罪，他们却公开地把我们下在监里。}\BibleRef{宗 19:37}因为千夫长说：\Quote{你就是那个埃及人吗？}他便立刻从他那种猜测中脱离出来；随后，为免人以为他否认自己的民族，他立即说：\Quote{我是犹太人\OriginalTerm{犹太人}{\GreekText{Ἰουδαῖος}}}——他是指自己的宗教。（\Emph{b}）那么，怎么呢？他没有否认（自己是基督徒）：决不是！因为他既是犹太人，也是基督徒，遵守着所当行的事；实在说，他在众人中最遵守法律，（\Emph{a}）正如他在别处自称：\Quote{在基督内，我是在法律之下}\BibleRef{格前 9:21}。这是什么意思，请问？（\Emph{c}）是指那信基督的人。而当与伯多禄辩论时，他说：\Quote{我们生来是犹太人……但我求你，准我向百姓说话}\BibleRef{迦 2:15}。这便证明他并未说谎，因为他以众人为证人。再看他是何等温和地说话。这也极有力地证明他无可指控，因他如此愿意为自己辩护，且渴望来与犹太百姓讲话。请看\OriginalTerm{一个准备好了的人}{\GreekText{τεταγμένον} \GreekText{ἄνδρα}}！请留意事情的眷顾安排：若非千夫长来到，若非他绑了他，他本不会但愿为自己辩护，也得不到那般的安静。\Quote{站在台阶上}。此外，地点也提供了便利，使他有一个高处向众人演讲——而且带着锁链！有什么场面能与此相比呢？看见保禄带着两条锁链，向众人演讲！他丝毫没有慌乱，丝毫不知所措；看见如此众多充满敌意的群众，长官站在旁边，他首先使他们停止发怒；然后，是何等明智地（这样做）。他在\WorkTitle{希伯来书}中所做的，在此也同样行之：首先以他们共同的母语声音吸引他们；然后以温和本身。\Quote{他用希伯来话对他们说：\InnerQuote{诸位弟兄各位父老，请听我现在对你们分诉。}}\BibleRef{宗 22:1}注意他的讲话，既无丝毫谄媚，又满含柔和。因为他不说\Quote{诸位大人}或\Quote{诸位主子}，而是说\Quote{弟兄们}，这是他们最喜欢听的词：\Quote{我不是外人}，他说，\Quote{也不是反对你们}。\Quote{诸位}，他说，\Quote{弟兄们和父老们}：这是尊敬的称呼，那是亲族的称呼。\Quote{请听}，他说，\Quote{我的}——他不说\Quote{教训}或\Quote{演讲}，而是\Quote{我现在的分诉}。他把自己置于恳求者的姿态。\Quote{他们听见他用希伯来话向他们说话，就更加安静了}\BibleRef{宗 22:2}。你是否注意到，使用同一种语言竟使他们降服？事实上，他们对那种语言有一种敬畏。也请注意他是如何为他的演讲铺路，这样开始：\Quote{我原是犹太人，生于基里基雅的塔尔索，却在这城里长大，在加玛里耳足前，按着祖传法律的严格方式受教，且对天主是热心的，如同你们今天一样。}\BibleRef{宗 22:3}\Quote{我原是……}，他说，\Quote{犹太人}：这是他们最爱听的；\Quote{生于基里基雅的塔尔索}。为免他们再以为他是外邦人，他补充了他的宗教：\Quote{却在这城里长大}（第282页注释4）。他表明他对敬拜是何等热心，因他离开自己那又大又远的家乡，为法律之故选择在这城里长大。请看他从一开始就依附法律。但他说这话，不仅是为向众人辩护，也是为表明他不是因人的意图而受引导去宣讲福音，而是因天主的能力：否则，受过如此教育，他不会突然改变。若他原是个普通人，或许会怀疑这点；但若他是最受法律约束的人之一，便不可能轻易改变，除非有强大的必要性。但也许有人会说：\Quote{在这城里长大并不能证明什么：因为若是来此经商，或另有原因呢？}因此他说：\Quote{在加玛里耳足前}；不仅是\Quote{由加玛里耳教导}，而是\Quote{在他足前}，表明他的恒心、勤勉、对听讲的热情，以及对那人的极大敬重。\Quote{按着祖传法律的严格方式受教}。不只说\Quote{法律}，而是\Quote{祖传法律}；表明他从起初便是如此，不仅仅是一个认识法律的人。这一切听起来似乎是为他们说话，但实际上却反对他们，因为他认识法律却离弃了它。\Quote{不错：但若你确实精确知道法律，却不捍卫它，甚至不爱它呢？}\Quote{是热心者}，他补充说：不仅（认识它）。既然他给了自己很高的赞誉，便又把这赞誉归给他们，加上：\Quote{如同你们今天一样}。因为他表明他们并不是出于人的目的，而是出于对天主的热情；这使他们高兴，先入为主地抓住了他们，且无害地赢得了他们。然后他也提出证据，说：\Quote{我曾按此道迫害他们直到死地，捆绑男女，下在监里。大司祭和整个长老团可以为我作证}（第4-5节）：\Quote{这如何显现？}他提出大司祭和长老们作为证人。他说：\Quote{是热心者，如同你们}（讲道集第十九篇第123页）：但他用行动表明，他超过了他们。因为我没有等待机会去捉拿他们：我既鼓动司铎，又外出行程；我不像你们只攻击男人，也扩展到女人：\Quote{捆绑男女，下在监里}。\par

这见证是无可反驳的；犹太人的（不信）已毫无借口。请看，他提出了多少证人——长老们、大司祭，以及城中的人。留意他的辩护，并非出于怯懦的恐惧（为自己申辩），不，而是为了教导和训诲。因为听者若不是石头，就必会感受到他所说的力量。因为直到此刻，他们自己就是证人；但其余的事却无证人：\BibleQuote{就是我从司祭长那里领了给弟兄们的书信，往大马士革去，要把那里的人捆起来，带到耶路撒冷受罚。当我前行，快到大马士革时，约在中午，忽然从天上有一道大光四面照着我。我就跌在地上，听见有声音向我说：扫禄，扫禄，你为什么迫害我？我回答说：主，你是谁？祂对我说：我是你所迫害的纳匝肋人耶稣。}（6、7、8节）既然如此，这些事本当凭先前的事就值得相信；否则他不会经历如此大的转变。你说，如果他只是编造一个漂亮的故事呢？回答我，他为何突然抛弃了所有的热忱？是为了寻求荣誉吗？然而他得到的恰恰相反。或许是安逸的生活？不，也不是。那么是为了别的什么？人的心思不能想出任何其他动机。于是，他让他们自己去推断，而只叙述事实。\BibleQuote{我走近大马士革时，}他说，\BibleQuote{约在中午。}请看那光何等强烈！你说，如果他只是编造故事呢？与他同行的人就是证人，他们曾搀扶他的手，见过那光。\BibleQuote{与我同行的人确实看见了那光，都害怕；却没有听见那对我说话者的声音。}\BibleRef{宗 22:9}但在另一处他说：\BibleQuote{听见声音，却看不见任何人。}\BibleRef{宗 9:7}这并不矛盾：原来有两个声音，保禄的声音和主的声音；在那处，作者指的是保禄的声音（讲道集十九，124页，注2）；事实上，（保禄）在此补充说：\BibleQuote{那对我说话者的声音。看不见任何人：}他没有说他们没有看见光，而是\Quote{没有人，}即\Quote{没有人在说话。}理应如此，因为惟有他堪当蒙受那声音。若他们也听见了，（奇迹）就不会如此伟大。因为心智粗陋的人更易被眼见说服，那些人也看见了光，就害怕了。事实上，那光在他们身上并未产生如在他身上的效果：因为甚至使他的眼睛瞎了；天主藉他所遭遇的事，也给了他们恢复视力的机会，只要他们愿意。在我看来，至少他们的不信是出于上智的安排，好使他们成为无可挑剔的证人。\BibleQuote{祂向我说，}经上说，\BibleQuote{我是你所迫害的纳匝肋人耶稣。}\BibleRef{宗 9:5}城市的名字（纳匝肋）也被加上，好让他们认出（此人）；此外，宗徒们也如此说。\EditorialNote{宗 2:22；4:10；10:38}祂亲自作证，他们正在迫害祂。\BibleQuote{与我同行的人看见了那光，都害怕；却没有听见那对我说话者的声音。我说：主，我当做什么？主对我说：起来，往大马士革去，在那里，将要告诉你你所被派定的一切事。我因那光的荣耀看不见，由同行的人搀着手，进了大马士革。有一个阿纳尼雅，按法律是虔诚的人，为住在那里的所有犹太人所称誉，来到我面前，站着说：扫禄兄弟，你看见吧！就在那时我抬头看见了他。进入城里去，}经上说，\BibleQuote{在那里，将要告诉你你所被派定的一切事。}\BibleRef{宗 22:10}看！又有一个证人。请看他也使那人成为无可挑剔的。\BibleQuote{有一个阿纳尼雅，}他说，\BibleQuote{按法律是虔诚的人，}——这远非外邦人的事！——\BibleQuote{为所有住在那里的犹太人所称誉}（那里）。\BibleQuote{就在同一时刻我恢复了视力。}接着是事实本身的见证。请看它是如何交织着人物与事实；这些人物既有犹太人，也有外邦人：司祭、长老，及他的同行者；事实则包括他所行的以及他所遭遇的：事实为事实作证，不仅止于人物。然后是阿纳尼雅，一个外邦人；然后是事实本身，视力的恢复；然后是一个伟大的预言。\BibleQuote{他说：我们祖先的天主拣选了你，叫你明白祂的旨意，并看见那义者。}\BibleRef{宗 22:14}\Quote{祖先的}说得好，为表明他们不是犹太人，而是法律的外邦人，且他们并非出于（法律的）热忱行事。\BibleQuote{叫你明白祂的旨意。}那么祂的旨意就是如此。请看，他如何在叙述的形式中进行教导。\BibleQuote{并看见那义者，听见他口中的声音。因为你要为祂向众人作证，见证你所看见所听见的。并且看见，}他说，\BibleQuote{那义者。}\BibleRef{宗 22:15}眼下他所说的不过如此：若祂是义的，他们就有罪。\BibleQuote{并听见他口中的声音。}请看他把这件事提到多高的层次！因为你要为祂作证——正是为此，因为你不至于背弃所见所闻（即\InnerQuote{背弃}）——\BibleQuote{你所看见的，和你所听见的：}藉着两种感官，他要求他的信德饱满——\BibleQuote{给众人。现在你为什么迟疑？起来，领受圣洗，洗去你的罪，呼求祂的名。}\BibleRef{宗 22:16}这里他讲了一件大事。因为他不说：\Quote{以祂的名受洗，}而是说：\BibleQuote{呼求基督的名。}这表明祂是天主：因为除了天主以外，呼求任何其他名字都是不合法的。然后他也表明自己并非被迫：因为\BibleQuote{我说，}他说，\BibleQuote{我该做什么？}没有一件事是没有证人的：不；他提出了全城的见证，因为他们曾看见他被搀着手带领。但请看预言的应验。\BibleQuote{给众人，}经上说。因为他确实成了祂的见证，且是应有的见证；藉他所受的苦、他所行的事、以及他所说的话。我们也应当成为这样的见证，不辜负我们所受托付的事：我所说的不仅是教义，也包括生活方式。\par

请注意：因为他看见了，因为他听见了，他向所有人作证，没有什么能阻止他。我们也在作证（现代文本作\Quote{听见了}），说有复活和无数美善之事：我们有责任向所有人作证。\Quote{是的，我们确实作证，}你会说，\Quote{而且相信。}何以？当你行为相反时？且说：若有人自称基督徒，而后背弃信仰，与犹太人同流，这证言能算数吗？绝不：因为人们渴望的是行为所表现的证言。同样，若我们说复活和无数美善之事存在，却轻视那些事物，偏爱眼前之物，谁会相信我们？所有人看的并非我们说甚麽，而是我们做甚麽。它说：\Quote{你要作见证，直到地极。}不仅是向友善者，也向不信者：因为作证的目的正在于此；不是说服知道的人，而是说服不知道的人。让我们成为可信赖的证人。但如何成为可信赖的呢？凭我们所过的生活。犹太人攻击了他：我们的情欲攻击我们，命令我们放弃证言。但我们不要服从它们：我们是来自天主的证人。（基督）被判定不是天主：祂派遣我们为祂作证。让我们作证并说服那些必须决断的人：若我们不作证，我们也要为他们的错误负责。但在世俗事务问案的法庭上，无人会接受一个充满无数恶行的证人，何况在此处，所考虑的事情如此（伟大）。\Emph{我们}说，我们听见了基督，并相信祂所许诺的事：他们说，用你们的行为证明吧：因为你们的生活却证明相反——你们不信。且说，我们该看看那些贪财之人、掠夺者、贪婪者吗？那些哀恸哭泣、建造忙碌于各样事务、仿佛永不会死的人？\Quote{你们不信自己会死，这事实如此明白明显：你们作证时，我们怎能相信你们？}因为有许多许多的人，其心态仿佛不会死。因为在漫长的老年中，他们着手建造和种植，何时会将死亡纳入考量？我们被召作证，却不能为自己所见的事作证，这将是我们的不轻惩罚。我们亲眼看见了天使，甚至比那些（有形地）看见他们的人更清楚。我们将成为（现代文本作\Quote{那么让我们成为}）基督的证人：因为不仅那些被称为\Quote{殉道者}（或证人，我们如此称呼）的是证人，我们自己也是。这就是为什么他们被称为殉道者，因为当被命令弃绝（信仰）时，他们忍受一切，为能说出真理；而我们，当被情欲命令弃绝时，不要被胜过。金子说：说基督不是基督。那么不要听它如同听天主，而要看轻它的吩咐。邪恶的情欲\Quote{自称认识天主，但在行为上却否认祂。}\BibleRef{铎 1:16}因为这不是作证，而是相反。并且，他人否认（祂）不足为奇；但我们这些被召作证的人否认祂，却是严重可憎的事：这对我们的事业伤害最大。祂说：\Quote{这要成为你们作证的凭据。}\BibleRef{路 21:13}，但（这是）当我们自己坚定站立的时候。如果我们都愿为基督作证，我们就很快能说服大部分外邦人。我亲爱的，所过的生活是一件大事。假设一个人野蛮如兽，公开因你的教义谴责你，他仍会私下认可、称赞、钦佩。且说，卓越的生活从何处而来？除了从在我们内运作的神圣能力，别无来源。\Quote{若有外邦人也如此品性，又如何？}若他们中有人在某处如此，部分出于本性，部分出于虚荣。你愿知道善良生活有何等的光芒、何等的说服力吗？许多异端者以此得胜，尽管他们的教义败坏，但大部分人因敬重其（有德）的生活而不去审查他们的教义；许多人甚至因教义谴责他们，却因生活敬重他们：这固然不对，但事实如此，他们就是这样感受（对待他们）的。这给我们信经的可敬条目带来了诽谤，这颠覆了一切，无人注重良善的生活：这是对信德的祸害。我们说基督是天主；我们提出无数别的论据，其中有一个，就是祂说服了所有人正当地生活：但这是少数情况。生活的恶劣是复活教义、灵魂不朽教义、审判教义的祸害；它也引来许多其他（虚假教义），命运、必然、否认天主眷顾。因为灵魂沉浸在无数恶习中，为了自我安慰，试图构想这些（虚假教义），以免因想到有审判以及善恶取决于我们自己的能力而感到痛苦。（这样的）生活产生无数罪恶，它使人成为野兽，甚至比野兽更无理性：因为野兽各自本性中的东西，它常聚集在一个人身上，颠覆一切。这就是为什么魔鬼引入了命运的教义：这就是为什么他说世界没有天主眷顾（Hom. ii. p. 15）：这就是为什么他提出善的本性和恶的本性的假说，以及（无始的）恶和（本质为）物质的假说：简言之，其余一切，为摧毁我们的生活。因为过这种生活的人，不可能从腐化的教义中自拔，也不可能保持健全的信德：而是必然不得不接受这一切。因为我个人认为，在生活不正的人中，很难找到谁不持有无数撒旦的计谋——例如，有\OriginalTerm{出生命运}{\GreekText{γένεσις}}，事情偶然发生，一切都是随意凑巧。因此，我恳求你们，让我们关心善的生活，以免接受邪恶的教义。加音受到的惩罚是（永远）呻吟颤抖。\BibleRef{创 4:14}恶人就是这样，他们内心意识到无数恶事，常从睡梦中惊醒，思绪充满混乱，眼睛充满惊慌；一切对他们都充满疑惧，一切令他们惊恐，灵魂充满痛苦的期待和怯懦的忧惧，被无力的恐惧和战栗所压抑。没有比这样的灵魂更懦弱、更空洞的了。如同疯子，没有自制之力。因为若能享有平静安宁，它本可认识自身的真正高贵。但当一切事物都吓唬它、使它陷入纷乱——梦境、言语、姿态、预感，混杂不清——它既如此烦恼惊惶，何时能审视自己呢？因此，让我们消除它的恐惧，打破它的束缚。因为即使没有别的惩罚，还有什么惩罚能超过这个——永远活在恐惧中，从不自信，从不安宁？因此，确知这些事，让我们保持平静，谨慎行善，持守健全的教义和正直的生活，使我们能毫无过失地度过今生，并能获得天主为爱祂的人所应许的美善事物，借着祂独生子的恩宠和仁慈，愿光荣、权能、尊崇归于祂，及于圣父和圣神，从今世直到永远。阿们。\par

\SourceRecord{Translated by J. Walker, J. Sheppard and H. Browne, and revised by George B. Stevens.；Nicene and Post-Nicene Fathers, First Series；Vol. 11.；Edited by Philip Schaff.；Buffalo, NY: Christian Literature Publishing Co.,；1889.；<http://www.newadvent.org/fathers/210147.htm>.}

% structure: p0001 source=h1 target=section reason=page-title

\section{论宗徒大事录讲道集 第四十八篇}

\BibleRef{宗 22:17-20}\par

\BibleQuote{及至我回到耶路撒冷，在圣殿里祈祷时，我陷入神魂超拔，看见主向我说：\InnerQuote{你快快离开耶路撒冷，因为他们不会接受你为我作的见证。}我说：\InnerQuote{主，他们知道，我曾在各会堂里监禁和鞭打那些信你的人；当你的殉道者斯德望的血流下时，我也站在一旁，赞同他的被杀，并看守那些杀他之人的衣服。}}\par

看他是怎样将自己投入（危险）的，他说：在那一神视之后，我\Quote{来到耶路撒冷。我神魂超拔}等等。再者，这件事没有见证；但请留意，见证来自结果。他说：\Quote{他们不会接受你的见证。}他们果然没有接受。然而，从理性的推断来看，本应猜想他们会毫无疑问地接纳他。因为我是那个与基督徒争战的人；所以他们理应接纳他。这里他确立了两件事：一是他们无可推诿，因为他们违背一切可能或理性的推断而迫害他；二是基督是天主，因为祂预言了与预期相反的事，并且不回顾过去，而是预知未来。那么，祂怎么说呢？\Quote{他要在外邦人、君王和以色列子民面前传扬我的名。}\BibleRef{宗 9:15}当然不是说服。此外，在其他场合我们发现犹太人被说服了，但这里却没有。在他们最应该被说服的地方，因为他们知道他从前（为他们）的热心，他们却没有被说服。\Quote{当你的殉道者斯德望的血……}等等。请看他的讲论又结束在哪里，即\OriginalTerm{强而有力的要点}{\GreekText{εἰς} \GreekText{τὸ} \GreekText{ἱσχυρὸν} \GreekText{κεφάλαιον}}：就是他迫害过，不仅迫害而且还杀害，甚至假若他有万手（\OriginalTerm{用无数的手}{\GreekText{μυρίαις} \GreekText{χερσὶν} \GreekText{ἀναιρών}}），他也会用它们来杀害斯德望。他提醒他们那种凶残的杀人精神（他和他们所）放纵的。于是当然他们最不会容忍他，因为这定了他们的罪；而预言也确实在实现：热情高昂，控诉激烈，犹太人自己成了基督真理的见证人！\Quote{他对我说：\InnerQuote{你去，因为我要派遣你到远方外邦人那里去。}他们直听到这句话，就扬声说：\InnerQuote{从地上除掉这样的人！因为他是不该活着的。}}（21、22节）犹太人不愿听他全部的宣讲，反而被怒气过度点燃，他们喊道，经上记着：\Quote{除掉他！因为他不该活着。他们喊叫，脱去衣服，向空中抛撒尘土，千夫长就吩咐将保罗带进营楼，且命人鞭打审问他，好知道他们为什么这样向他喊叫。}（23、24节）千夫长本应查明这些事是否属实——是的，犹太人自己也当如此——或者，若不属实，就应命人鞭打他，他却\Quote{吩咐用鞭子审问他，好知道他们为什么这样向他喧嚷}。然而，他本应从那些喊叫者那里了解，并询问他们是否抓住了所说的话中任何一点；但他却二话不说，任意妄为，以讨好他们为行事目的：因为他所看的不是如何行义，而是如何止息他们那不义的怒气。\Quote{他们用皮条绑上他时，保罗对站在旁边的百夫长说：\InnerQuote{你们鞭打一个尚未定罪的罗马人，这是合法的吗？}}\BibleRef{宗 22:25}保罗没有说谎，千万不要！因为他确实是罗马人；如果别无其他，他也会害怕（假装），以免被识破而遭受更严重的惩罚。（参阅苏埃托尼乌斯《\WorkTitle{克劳狄传}》第25节。）且请注意，他并不是直截了当地（\OriginalTerm{简单地}{\GreekText{ἁπλῶς}}）说，而是说：\Quote{你们鞭打……合法吗？}所提出的控诉有两项：既未经审问，他又是个罗马人。那时，他们视此为极大的特权：因为据说（只是）从哈德良时代起，所有人都被称为罗马人，但古时并非如此。他若被鞭打，就会受人轻视；但实际上，他却使他们（比他自己）更加恐惧。倘若他们鞭打了他，他们也会了结整个事件，甚或杀了他；但事实上，结果并非如此。请看天主是如何允许多（良好的结果）完全以人性的方式到来，无论在宗徒身上还是其他（人类）身上。请注意他们如何怀疑这事是个借口，以为保罗自称罗马人是撒谎：也许是从他的贫穷来推测。\Quote{百夫长听见这话，就去报告千夫长说：\InnerQuote{你要小心，这人是个罗马人。}千夫长就进前来问他说：\InnerQuote{告诉我，你是罗马人吗？}他说：\InnerQuote{是。}千夫长回答说：\InnerQuote{我用一大笔银子才取得这公民权。}保罗却说：\InnerQuote{我生来就是。}于是那些要审问他的人立刻离开了他；千夫长也知道他是罗马人，又因曾绑过他，也害怕起来。}\BibleRef{宗 22:26}——\Quote{但我}，他说，\Quote{是生来的。}那么他的父亲也是罗马人。这又有什么结果呢？他绑了他，又把他带到犹太人面前。\Quote{第二天，千夫长想要知道犹太人控告他的确实理由，就解开他的绑，并命令司祭长和全体公议会集合，带保罗下来，让他站在他们面前。}\BibleRef{宗 22:30}他现在不再向群众或百姓讲论。\Quote{保罗注视着公议会说：\InnerQuote{诸位弟兄，我在天主面前，行事为人都是凭着良心的纯全，直到今日。}}\BibleRef{宗 23:1}他的意思是：我自觉完全没有亏负你们，也没有做过任何配得这些捆绑的事。那么大司祭说了什么呢？非常公正、有领袖风范且温和：\Quote{大司祭阿纳尼雅命站在旁边的人打他的嘴。保罗便对他说：\InnerQuote{你这粉饰的墙，天主将要打你！你坐着按法律审判我，竟违反法律命人打我！}站在旁边的人说：\InnerQuote{你竟敢辱骂天主的大司祭？}保罗说：\InnerQuote{弟兄们，我不知道他是大司祭；因为经上记着：不可毁谤你百姓的首领。}}\BibleRef{宗 23:3}因为\Quote{我不知道他是大司祭}。有人说：那么他为什么为自己辩护，仿佛这是控告，还加上\Quote{不可毁谤你百姓的首领}？如果他不是首领，难道只因为不是就辱骂（他或任何）其他人就是对的吗？他自己说：\Quote{被人辱骂，我们就祝福；被人迫害，我们就忍受。}\BibleRef{格前 4:12}但在这里他却做了相反的事，不仅辱骂，而且还诅咒。这是勇敢的话，而非愤怒的话；他不愿在千夫长面前显得可鄙。因为假设千夫长本人已免除鞭打他，但他即将被交给犹太人，被他们的仆人殴打反而会使他更加大胆：这就是为什么保罗不攻击仆人，而攻击下令的人。但那句\Quote{你这粉饰的墙，你坐着按法律审判我？}（是）替代\Quote{你（自己）是个罪人}：仿佛在说：\Quote{（你）自己也该受无数鞭打。}请看他们的勇气使他们多么震惊；因为本来是要推翻整个事件，他们反倒称赞他了。（下文，第9节。）\Quote{因为经上记着}等等。他想表明他这样说，不是出于恐惧，也不是因为（阿纳尼雅）不配被称为这个，而是出于对法律在这点上的服从。而且我深信他确实不知道那是大司祭，因为他当时是经过长期间隔才回来，并不习惯与犹太人经常往来；况且看见他在许多人中间：因为大司祭已不再是一眼就可以认出的，他们人数众多且各不相同。所以，在我看来，他在此也是为着他针对他们的抗辩而说的：为表明他遵守法律；因此他（这样）为自己开脱。\par

(重申。) (\Emph{乙}) 但让我们回顾已经说过的话。 (\Emph{甲}) \Quote{当我再次来到耶路撒冷时，}等等。 \BibleRef{宗 22:17} 怎么回事？他身为犹太人，在那里长大受教，却没有留在那里？他也没有住在那里，除非他有意给自己制造无数机会：到处像流放者一样，从一个地方逃到另一个地方。 (\Emph{丙}) 他说：\Quote{当我在圣殿里祈祷时，}\Quote{我出神了。}（为要表明）这并非仅仅是想像的幻影，因此\Quote{当他祈祷时}（主）站在他旁边。他表明他逃跑并非出于对他们危险的恐惧，而是因为他们\Quote{不接受}他的\Quote{见证。} \BibleRef{宗 22:18} 但他为什么说\Quote{他们知道我囚禁了人？} \BibleRef{宗 22:19} 并非为了反驳基督，而是因为他想知道这件如此违背常理的事。然而基督并没有教导他（这事），只吩咐他离开，他就服从了：他是如此顺从。 经上说：\Quote{他们大声喊叫，}\Quote{说：除掉他！这个人不该活着。} \BibleRef{宗 22:22} 不，你们才是不该活着的人；不是他，他凡事服从天主。哦，恶徒和杀人犯！ 经上说：\Quote{他们抖了衣服，}\Quote{把尘土撒向空中} \BibleRef{宗 22:23}，为了激起更激烈的骚乱，因为他们想吓唬总督。请注意：他们不说指控是什么，因为事实上他们没有什么可指控的，只是想用喊叫来制造恐怖。 \Quote{千夫长下令，}等等，可他本应从控告者那里得知\Quote{他们为何如此喊叫反对他。当他们捆绑他时，等等。千夫长得知他是罗马人后，害怕了。} 那么，这不是假话。 \Quote{次日，因他要知道犹太人控告他的确切理由，等等，就把他带到公议会前。} \BibleRef{宗 22:24} 这本是他一开始就该做的。他带他进去，松开他。这件事尤其让犹太人不知所措。 经上说：\Quote{保禄，}\Quote{定睛看着他们。} 这表明他的勇气，以及这如何震慑了他们（\OriginalTerm{震慑}{\GreekText{τὸ} \GreekText{ἐντρεπτικόν}}）。 \Quote{然后大司祭阿纳尼雅。}等等。\EditorialNote{参见宗 23:1-2} 怎么，他说了什么冒犯的话？他为什么被打？哎呀，何等鲁莽，何等无耻！因此（保禄）斥责他：\Quote{天主必打击你，你这粉白的墙。} \BibleRef{宗 23:3} 因此（阿纳尼雅）自己也哑口无言，不敢说话：只有他身边的人受不了保禄的勇气。他们看见一个准备赴死的人……因为如果情况如此，（保禄）只需保持沉默，千夫长就会带走他，走开；他会把他交给他们。他既表明自己甘心受苦，又借此在他们面前为自己辩解——并非他想要向他们辩解——因为他甚至强烈谴责那些人——而是为了百姓的缘故。 \Quote{你们违背律法，竟吩咐打我吗？} 他这样说很有道理：因为杀死一个没有伤害（他们）的人，而且是无辜的人，就是违背律法。因为他所说的并非辱骂，除非有人把基督的话称为辱骂，当祂说：\Quote{祸哉，你们经师和法利塞人，因为你们像粉白的墙。} \BibleRef{玛 23:27} 不错，你会说：但如果他是在被打之前说的，那就表明不是愤怒，而是勇气。但我已经提到了原因。而且（照此看来）我们常常发现基督自己在被犹太人辱骂时\Quote{说话辱骂}他们；正如祂说：\Quote{不要想我要控告你们。} \BibleRef{若 5:45} 但这并非辱骂，绝对不可。看，他如何温柔地对他们说话：他说：\Quote{我不晓得，}\Quote{他是天主的大司祭}（4-5节）：并且（为了表明）他并非\OriginalTerm{装假}{\GreekText{εἰρωνεύεται}}，他补充说：\Quote{不可毁谤你百姓的官长。} 他甚至承认他仍是官长。 我们也学习温良，好使我们在这两方面都成全。因为必须仔细审视它们，以了解前者是什么，后者是什么：仔细，因为这些德性近旁都有相应的恶习：仅仅是鲁莽冒充勇气，仅仅是懦弱冒充温良：需要分辨它们，免得有恶习的人似乎有德性：这就像一个人以为他与女主人同居，却不知道那是女仆一样。那么，什么是温良，什么是仅仅的懦弱？当别人受冤屈，我们不挺身而出，反而保持沉默，这就是懦弱：当我们自己受虐待，我们忍受，这就是温良。什么是勇气？同样，当别人是我们为之争辩的对象时。什么是鲁莽？当我们在自己的事上愿意争斗时。因此，宽宏大量和勇气相伴，正如（仅仅的）鲁莽和（仅仅的）懦弱相伴。因为那为自己不怀恨的，几乎不能不为了别人而怀恨：那不为自己的事挺身而出的，几乎不能不为别人挺身而出。因为当我们的习惯性情没有激情时，它也能接纳德性。正如身体没有发烧就能接纳力量，同样，灵魂除非被激情败坏，就能接纳力量。这温良表明极大的力量；它需要一个慷慨而英勇的灵魂，一个极其崇高的灵魂。或者，你认为，忍受恶行而不被激怒是小事吗？确实，如果有人把为邻人挺身而出的性情称为男性勇气，他并不会错。因为那有力量克服如此强烈的激情（如自私），也会有力量敢于攻击别人。例如，有两种激情：懦弱和愤怒：如果你克服了愤怒，很明显你也克服了懦弱：但你通过温良获得了对愤怒的掌控：因此也对懦弱如此，你就会成为有丈夫气概的人。反之，如果你没有战胜愤怒，你就变得鲁莽好斗；但没有战胜这个，你也不能战胜恐惧；结果你也会成为懦夫：就像身体一样，如果虚弱，就很快被冷热征服：因为这就是病弱的体质，而良好的体质能经受一切变化。另外，心灵伟大是一种德性，旁边就是挥霍：节俭是一种德性，善于管理；旁边就是吝啬和小气。来，让我们再次比较德性和恶习。那么，挥霍的人不能被称为心灵伟大。他怎么能呢？被无数激情战胜的人，怎么能心灵伟大？因为这不是轻视金钱；这只是被其他激情支配：就像一个人，如果听从强盗的吩咐服从他们的命令，就不能自由。他的大量花费并非出于对金钱的藐视，而只是因为他不知道如何妥善处置：否则，如果既能保存又能花在享乐上，他会喜欢那样。但那把钱花在适当对象上的人，才是心灵崇高的人：因为真正的心灵崇高，是不受激情奴役，视金钱为无物。再者，节俭是好事：因为这样最善于管理的人，会以适当的方式花费，而不是随意没有管理。但吝啬与此不同。因为前者即使有紧急需要，也不动用本金；但后者将是前者的兄弟。那么，我们将把心灵伟大的人和明智的节俭者放在一起，同样也把挥霍者和吝啬者放在一起：因为这两者都是由于心灵渺小而如此，正如那些相反的人（由于相反）。因此，让我们不要称那只是花费的人为心灵崇高，而是称那正确花费的人；也不要称节俭的管理者为吝啬和小气，而是称那不合时宜地节省金钱的人。\par

那个\Quote{穿紫红袍和细麻衣的}富翁花费了何等巨额的财富？\BibleRef{路 16:19} 但他并不是一个心胸高尚的人：因为他的灵魂被无情的情感和无数的情欲所占据：那么它怎么能是伟大的呢？亚巴郎有一个伟大的灵魂，他为了款待客人，杀牛犊，并且在需要时，不仅不吝惜他的财产，甚至连性命也不顾惜。因此，如果我们看到一个人摆着丰盛的宴席，拥有他的娼妓和寄生虫，我们不要称他为心胸宽广的人，而是一个极其渺小的人。因为看，他受多少情欲的奴役和支配——饕餮、无度的快乐、谄媚：但那个被这么多情欲所占据，甚至无法逃脱其中任何一个的人，怎么能有人称他为宽宏大量呢？不，恰恰相反，我们最应该称他为渺小之人，正是当他花费最多的时候：因为他花费得越多，就越显示那些情欲的暴政：因为如果它们没有过分地控制他，他就不会过度花费。再者，如果我们看到一个人，什么都不给这样的人，却喂养穷人，帮助有需要的人，自己过着简朴的生活——我们就称他为极其高尚的人：因为轻视自己的安逸而关心别人的安逸，确实是伟大灵魂的标志。请告诉我，如果你看到一个人藐视所有的暴君，不把他们的命令当回事，却从他们的暴政中解救那些受压迫和受虐待的人；你难道不认为这是一个伟大的人吗？所以，让我们也这样评价这个人。情欲就是暴君：因此，如果我们轻视它们，我们就是伟大的；但如果我们还能帮助别人摆脱它们，我们就更伟大，因为我们不仅足够应付自己，还能帮助别人。但如果有人，按照暴君的命令，殴打他的另一个臣民，这是灵魂的伟大吗？当然不是：而是与他伟大程度成正比的极端奴役。现在也摆在我们面前\OriginalTerm{设于我们面前}{\GreekText{πρόκειται}}一个高尚而自由的灵魂：但那个浪荡子却命令他的情欲鞭打它：那么，那个鞭打自己的人，我们称他为高尚吗？绝不。那么，* *，但让我们看看什么是灵魂的伟大，什么是挥霍；什么是节俭，什么是吝啬；什么是温柔，和（什么）迟钝和怯懦；什么是勇敢，什么是冒失：好让我们分辨这些事情，得以中悦地度过此生，并得到所应许的美物，藉着我们的主耶稣基督的恩宠和慈爱，愿光荣归于祂，至于无穷之世。阿们。\par

\SourceRecord{Translated by J. Walker, J. Sheppard and H. Browne, and revised by George B. Stevens.；Nicene and Post-Nicene Fathers, First Series；Vol. 11.；Edited by Philip Schaff.；Buffalo, NY: Christian Literature Publishing Co.,；1889.；<http://www.newadvent.org/fathers/210148.htm>.}

% structure: p0001 source=h1 target=section reason=page-title

\section{宗徒大事录讲道集 第四十九篇}

\BibleRef{宗 23:6-8}\par

\BibleQuote{但当保禄察觉一部分是\OriginalTerm{撒杜塞人}{Sadducees}，另一部分是\OriginalTerm{法利塞人}{Pharisees}时，就在公议会中喊说：\InnerQuote{诸位仁人弟兄！我是法利塞人，是法利塞人的子孙；我是为了死人复活的望德与复活而受审。}他这话一出口，法利塞人与撒杜塞人之间就起了纷争，会众分裂了。因为撒杜塞人说没有复活，也没有天使，也没有神灵；法利塞人却两者都承认。}\par

他又如同普通人一样单纯地讲话，并不在所有场合都同样享有超自然的助益。\Quote{我是法利塞人，是法利塞人的儿子：}他在这里以及此后的话中，都希望分裂那对他怀有邪恶一致的群众。他在这里也没有说谎：因为他按祖先是法利塞人。\Quote{我是为了希望和死者的复活而受审。}因为他们不愿说出控告他的理由，所以他被迫自己说明。\Quote{但法利塞人，}经上说，\Quote{两者都承认。}然而有三件事：他怎么说是两者呢？\Quote{神和天使}被视为一个。当他站在他们一边时，他们便为他辩护。\Quote{于是起了大喊：法利塞党派的经师起来争辩说：我们在这人身上找不到恶；但}（什么）\Quote{若是一个神对他讲过话，或一个天使？}\BibleRef{宗 23:9}他们之前为何不为他辩护？你看见吗，当情感让步时，真理便被显明？他们说，若是一个天使或一个神对他讲话，罪在哪里？保禄没有给他们攻击他的把柄。\Quote{当起了大纷争时，千夫长害怕保禄会被他们扯碎，就命令兵士下去，从他们中间把他抢出来，带到营楼里。}\BibleRef{宗 23:10}千夫长害怕他被扯碎，既然他说自己是罗马人；这事并非没有危险。你看见吗，保禄有权宣称自己是罗马人？否则，千夫长现在也不会害怕。所以，兵士必须强行把他带走。但当那些恶徒看见一切都无济于事时，他们便亲自处理整个事情，正如他们先前想要做却被阻止的一样；他们的邪恶虽然受到如此多的阻碍，却仍未停止。然而有多少事情是天主上智安排的，为使他们能从愤怒中平静下来，并学习那些可能使他们恢复自己的事！但他们仍然攻击他。足以证明他清白的是，这个人即将被扯碎时却被救，并在如此大的危险中逃脱了一切。\Quote{当夜，主站在他旁边说：放心罢，保禄！因为你怎样在耶路撒冷为我作证，也必怎样在罗马为我作证。到了天亮，有些犹太人共谋，且发诅咒起誓说：若不杀了保禄，就不吃不喝。有四十多人这样同谋。}\BibleRef{宗 23:11}\Quote{他们发诅咒起誓，}经上说。看他们多么激烈和报复心切！\Quote{发诅咒起誓}是什么意思？那么那些人永远被控告，因为他们没有杀死保禄。而且四十多人一起。因为那民族的性格就是这样：当需要为一件好事合作时，连两个人也不能彼此一致；但当是为了坏事时，全体人民都去做。他们还让首领们作共犯。\Quote{他们来到司祭长和长老那里说：我们发了重誓，若不先杀了保禄，就不吃东西。所以现在你们同公议会通知千夫长，叫他明天带保禄下到你们这里来，好像你们要更详细查问他的事；我们在他走近之前，就预备好杀他。保禄姊妹的儿子听见他们设下埋伏，就来到营楼里告诉了保禄。保禄便叫来一个百夫长说：带这个青年到千夫长那里去，因为他有件事要告诉他。于是百夫长带他到千夫长那里说：囚犯保禄请我到他那𥚃，求我带这个青年到你这里来，他有话要告诉你。千夫长便拉着手，带他私下到一边，问他说：你有什么话要告诉我？他说：犹太人约定，要请求你明天带保禄下到公议会，好像要更详细查问他的事。但你不要听从他们，因为他们中有四十多人埋伏等着他，并且发了誓，若不先杀了他，就不吃不喝。现在他们预备好了，等候你的应允。于是千夫长放那青年走，吩咐他说：不要告诉人你把这些事报告了我。}\BibleRef{宗 23:14}他又因人的远见而得救。请看：保禄不让任何人知道这事，连百夫长也不让，免得事情泄露。百夫长来了，报告给千夫长。千夫长做得很好，他吩咐他保密，免得泄露；此外，他只在下令行事时吩咐百夫长们：这样保禄被解送到凯撒勒雅，为在那里也在一个更大的剧场和更显赫的听众前讲话：这样犹太人就不能说：\Quote{我们若见了保禄，就会相信——若听了他教训。}因此这个借口也被他们剪除了。\Quote{主，}经上说，\Quote{站在他旁边说：放心罢！因为你怎样在耶路撒冷为我作证，也必怎样在罗马为我作证。}（然而）即使他显现给他后，又让他藉人的手段得救。人们也大可惊奇保禄；他没有惊慌失措，也没有说：\Quote{怎么？这是什么事？难道我受基督欺骗了？}但他相信；然而，因为他相信，他并没有因此睡觉；不；他凭人的智慧在自己能力范围内的事，他没有放弃。\Quote{他们发诅咒起誓：}那些人藉着诅咒给自己加了一种约束。\Quote{若不杀了保禄，就不吃不喝。}看，禁食是谋杀之母！正如黑落德因他的誓言给自己加上了那种约束，这些也一样。因为魔鬼的（方式）就是这样：表面上藉着虔诚，暗地里设下陷阱。\Quote{他们来到司祭长那里，}等等。然而他们本该来到（千夫长那里），提出控诉，召集法庭：因为这些不是司铎的行为，而是匪首的行为，这些不是统治者的行为，而是暴徒的行为。他们还想败坏统治者；但这是上智的安排，为使他（千夫长）也知道他们的阴谋。因为不仅因为他们无话可说，也因为他们秘密的企图，他们证明自己是恶的。很可能在（保禄走后）司祭长们来见（千夫长）提出请求，并且蒙羞。当然他既不愿拒绝也不愿答应他们的请求。他如何相信（那青年的故事）？他这样做是由于已经发生的事；因为他们也可能这样做。看他们的邪恶：他们几乎是强迫司祭长们：因为他们自己承担了如此大的事，并置身于全部危险中，那么那些人更应该做这么多。你看见吗，保禄被外邦人视为清白，正如基督也被比拉多视为清白。看他们的毒计归于无有：他们交出他，要杀死并定他的罪；但结果正好相反：他既得救，又被视为清白。若非如此，他就会被扯碎；若非如此，他就会丧命，被定罪。千夫长不仅把他从（对他的）冲击中救出，也从许多其他（暴力）中：你看他如何成为他的助手，以致他毫无危险地被如此大的军队安全带走。\Quote{他叫来两个百夫长说：预备两百士兵，七十骑兵，两百长枪手，在夜里第三时辰出发，往凯撒勒雅去；并预备牲口，叫保禄骑上，安全解送到总督斐理斯那里。千夫长写了这样的信：革老丢里息雅向总督斐理斯大人请安。这人被犹太人拿住，快要被他们杀死；我带领兵士前来，救了他出来，因为知道他是罗马人。我想要知道他们控告他的原因，就带他下到他们的公议会去。我发现他被控告，是关乎他们法律的问题，并没有什么该死或该监禁的罪名。后来有人把犹太人要害这人的计谋告诉我，我便立时解他到你那里去，又吩咐控告他的人，到你面前告他。愿你平安。}\BibleRef{宗 23:23}看这信是怎样为他辩护——因为它说：\Quote{我找不到该死的事，}但与其说是控告他，不如说是控告他们。\Quote{快要被他们杀死：}他们多么想置他于死地。首先，\Quote{我带领兵士前来，救了他出来：}然后，\Quote{我带他下到他们那里：}即便如此，他们也没有找到什么可以控告他的；当他们本应对前一事感到惧怕和羞愧时，他们又企图杀他，以致他的事更加明朗。\Quote{他的控告者，}他说，\Quote{我已打发到你那里去：}为在那更严格审查这些事的法庭上，他能被证明无罪。\par

（概要。）让我们回顾上面所说的话。\Quote{我}，他说，\Quote{是法利塞人}；然后，为免显得谄媚，他补充说：\Quote{我被审讯，正是为了死者的希望和复活。}\BibleRef{宗 23:6} 借此控诉与诽谤，他为自己辩护。\Quote{因为撒杜塞人确实}等等。撒杜塞人对无形体之物一无所知，甚至可能对天主也一无所知；他们是如此\OriginalTerm{粗鄙}{\GreekText{παχεἵς}}：因此他们也不愿相信有复活。\Quote{并且经师们}等等。看，千夫长也听见法利塞人已宣告他无罪，并且（手稿和版本作\OriginalTerm{他已判决}{\GreekText{ἐψηφίσατο}}，即\Quote{他已判决}）做出有利于他的判决，于是更有信心地将他从危险中带走。而且（保禄）所说的一切都充满\OriginalTerm{哲学}{\GreekText{φιλοσοφίας}}。\BibleQuote{当夜，主站在他旁边}等等。请看何其大的安慰！首先祂赞美他：\Quote{你既在耶路撒冷为我作证}；然后祂不让他担心前往罗马的旅途结果不确定：因为在那里，祂说，你不会单独离开（\OriginalTerm{单独}{\GreekText{μόνος}}，校勘本和版本作\OriginalTerm{单独}{\GreekText{μόνον}}），但你也会拥有这全部的宣讲勇气。由此显明，他不仅得救，而且是为在大城市中得大荣冠而如此。但为何主不在他陷入危险之前显现给他？因为在患难中天主总是安慰我们；因为祂显得更被人渴慕，即使在危险中祂也锻炼和训练我们。此外，他当时无拘无束，安然无恙；但现在巨大的危险正等候着他。\Quote{我们已发下重誓}，他们说，\Quote{不吃饭不喝酒。}\BibleRef{宗 23:14} 这一切热忱是什么？\Quote{好将他带下来}，经上说，\Quote{到你们这里，好似要更详细查问他的案情。}\BibleRef{宗 23:15} 他不是已经两次向你们讲话了吗？他不是说过他是法利塞人吗？除此外你们还要什么？他们是如此鲁莽，无所畏惧，不惧法庭，不惧法律：他们的胆大妄为到了无以复加的地步。他们既宣告了他们的意图，也宣布了实现的方式。\Quote{保禄姐妹的儿子听到了这事。}\BibleRef{宗 23:16} 这是天主的眷顾，因为他们未觉察这事会被听到。保禄如何应对？他不惊慌，但看出这是天主的作为：并将一切交托给祂，于是自己摆脱了（进一步的担忧：）\Quote{叫了一个百夫长来}等等。\BibleRef{宗 23:17} 他报告了这阴谋，他被相信了；他得救了。如果他已被宣告无罪，为何（千夫长）还要送控诉者？为的是查问更严格：为使人更完全被洗清。\par

天主的安排方式便是如此：正是那伤害我们的事物，也藉由同一事物使我们受益。例如\Person{若瑟}：他的主母企图毁灭他；她看似在策划他的毁灭，但藉着她的策划，却使他处于安全之境；因为那间关着这头（如野兽般的）女人的屋子，与监牢相比，反倒是一个兽穴，而监牢却显得温和。\BibleRef{创 39:1} 因为他在那里时，虽受敬重和奉承，却常恐惧，唯恐主母加害于他；这恐惧比任何监牢更甚：但在控告之后，他反倒安全和平安，脱离了那头野兽、她的淫荡和毁灭他的阴谋；因为与受苦的人为伍，好过与疯狂的主母相伴。他在狱中得安慰，因他为贞洁而陷入其中；而在主母家中，他害怕灵魂遭受致命打击：世上没有什么比一个恋慕的女子对不肯（回应她）的青年更令人烦恼的了；没有比（这样的女子）更可憎、更凶残的了；世间一切束缚与这相比都算轻省。所以事实并非他进了监牢，而是他脱离了监牢。她使他的主人成为他的仇敌，却使天主成为他的朋友；使他与那真正的主人更亲密；她将他逐出家中的管家的职位，却使他成为那主人的密友。再者，他的兄弟们出卖了他\BibleRef{创 37:18} ；却使他从与仇敌同住一屋、从嫉妒和极深的敌意、从每日谋害他的阴谋中得释放；他们使他远离那些恨他的人。因为还有什么比被迫与嫉妒自己的兄弟同住一屋、被猜疑、成为恶计的目标更糟的呢？因此，当他们和她各自图谋自己的目的时，天主眷顾却为这义人成就了极其不同的伟大后果。当他受尊荣时，他处于危险；当他受屈辱时，他反倒安全。太监不记念他，这真是好事，好使他得救的时机更为荣耀；好让一切归功于天主眷顾，而非人的恩惠\BibleRef{创 40:23} ；好让\Person{法老}在迫不得已的时刻，将他带出来；好让国王释放他出监，不是赐予恩惠，而是接受恩惠。\BibleRef{创 41:40} 这不应是奴仆的赠礼，而是国王被迫如此行；应该彰显他内有怎样的智慧。因此太监忘记了他，好让\Place{埃及}不忘记他，好让国王不忽视他。若当时他被释放，很可能他就想返回本国；因此他被无数的约束留住，先是为主人，其次为监牢，第三为治理王国，好让这一切由天主眷顾成就。如同一匹骏马急于奔向同伴，天主将他留在那里，为荣耀的缘故。因为他渴望见到父亲，并解除他的痛苦，这从他召唤父亲到\Place{埃及}可以明显看出。\BibleRef{创 45:9}\par

让我们看看恶意谋算的其他例子，它们如何对我们有利——不仅通过它们最终的报应，而且通过它们在当时恰恰好为我们带来益处？\Person{若瑟}的叔父（\Person{厄撒乌}）曾怀恶念对付他的父亲（\Person{雅各伯}），并把他赶出故乡。结果如何？\BibleRef{创 27:41} \Person{厄撒乌}反而（借此）使\Person{雅各伯}远离了危险，因为\Person{雅各伯}也（借此）得以安全。\Person{厄撒乌}使\Person{雅各伯}成为一个更智慧、更优秀的人（\OriginalTerm{更有智慧的}{\GreekText{φιλοσοφώτερον}}）；并且成为了他得到那个梦境的因由。\BibleRef{创 28:12} 但你会说，他却在异乡为奴？是的，但他到了自己亲属那里，娶了新娘，并在他岳父面前显得配得上。\BibleRef{创 29:23} 但岳父也欺骗了他？是的，但这同样对他有利，使他成为众多孩子的父亲。但\Person{厄撒乌}内心怀有恶意反击他？没错，但即便如此也对他有利，使他因此得以返回本国；因为如果他境遇良好，就不会如此渴望回家。但\Person{拉班}克扣了他的工钱？确实，但他却因此得到了更多。\BibleRef{创 31:7} 由此可见，在这些人的历史中，有越多的人谋害他们，他们的事业就越兴旺。如果\Person{雅各伯}没有先娶了长女，他就不会这么快成为这么多孩子的父亲；他将在无子嗣中度过漫长的岁月，像他的妻子一样哀伤。因为他的妻子确实有理由哀伤，因为她没有成为母亲\BibleRef{创 30:1-2}；但\Person{雅各伯}却有安慰，因此他也拒绝了她的请求。再者，如果\Person{拉班}没有克扣他的工钱，他就不会渴望看见自己的祖国；那人品性的更高境界（\OriginalTerm{智慧}{\GreekText{φιλοσοφία}}）就不会显现出来，他的妻子们也不会对他更加依恋。因为请看她们所说的：\BibleQuote{他吞噬了我们，吞没了我们的钱财。}\BibleRef{创 31:15} 所以这成为了巩固她们对他爱慕的手段。此后，他拥有的不只是妻子，还有（忠心的）奴隶；他被她们所爱：这是任何财富都无法比拟的；因为没有任何东西，绝对没有任何东西，比被妻子如此爱慕并爱慕她更宝贵。经上说：\Quote{与丈夫和睦的妻子}\BibleRef{德 25:1}，又说：\Quote{和睦相爱的夫妻。}智者把这事列为使人有福的事之一；因为哪里有这事，那里就有一切财富，一切繁荣；反之，哪里没有它，那里的一切便毫无益处，一切都出错，只有烦恼和混乱。因此，让我们超过一切追求这事。追求金钱的人，不追求这事。让我们追求那些能持久不变的事物。让我们不要向富人求妻，免得她一方财富过多产生骄傲，免得那骄傲毁掉一切。难道你们没有看见天主做了什么？祂如何使女人处于隶属地位？\BibleRef{创 3:16} 你们为什么忘恩负义，为什么毫无知觉？天主凭本性赐给你们的恩惠，你们不要破坏它原有的帮助。因此，我们求妻不是为了她的财富，而是为了我们得到一位生活的伴侣，为了生育子女的既定秩序。天主赐予女人，不是要她带来金钱，而是要她做一名助手。但带来金钱的女人，不但不是妻子，反而成了自我意志的设立者（\OriginalTerm{阴谋家}{\GreekText{ἐπίβουλος}}），一位女主人——她可能成了野兽而不是妻子——她以为自己有权凭借财富趾高气扬。没有比以此方式致力于发财的男人更可耻的。如果财富本身充满诱惑，对于这样得来的财富，我们还有什么可说？因为你们绝不可着眼于那些罕见、特例、违背事物理性而成功的人的个例；正如我们在其他事上，不应盯着人们可能享有的好处或偶然的成功，这些超乎常规：而应从事物的本身理性来看，看看这事是否充满无尽的烦恼。你不仅使自己陷入不名誉的境地，而且如果碰巧你先妻子离世，你还留下子女贫困，使他们蒙羞；并且你给她提供了多得多的机会去联系第二个新郎。难道你没有看见许多女人以这为借口再婚——为了不被轻视；她们希望有个男人来管理她们的财产？那么，让我们不要为了金钱而带来如此大的邪恶；让我们摒弃所有（这样的目标），寻求美丽的灵魂，使我们也能获得爱。这就是极大的财富，这就是巨大的宝藏，这就是无尽的美好事物：愿我们众人靠着我们的主\Person{耶稣基督}的恩宠和慈爱达到这境界，祂与圣父及圣神同享荣耀、主权、尊威，自今至永远，及世之世。阿们。\par

\SourceRecord{Translated by J. Walker, J. Sheppard and H. Browne, and revised by George B. Stevens.；Nicene and Post-Nicene Fathers, First Series；Vol. 11.；Edited by Philip Schaff.；Buffalo, NY: Christian Literature Publishing Co.,；1889.；<http://www.newadvent.org/fathers/210149.htm>.}

% structure: p0001 source=h1 target=section reason=page-title

\section{第50讲：论\WorkTitle{宗徒大事录}}

\BibleRef{宗 23:31-33}\par

\BibleQuote{于是兵士们依照所命令的，夜间带着保禄到了安提帕特。次日，他们留下骑兵同他前往，自己返回了营寨。他们到了凯撒勒雅，把信呈给总督，也把保禄引到他面前。}\par

如同一位君王由他的卫兵护送一样，这些人就这样运送\Person{保禄}；人数也如此众多，且是在夜间，因为害怕百姓的怒火。那么，你会说，既然他们把他带出城，他们就停止暴行了吗？不，的确没有。但是（千夫长）不会如此谨慎地为了他的安全而遣送他，除非他自己没有发现他有什么过错，并且知道敌人的杀人心肠。\Quote{当总督读了那封信后，他问他是哪一省的人。当他知道他是\Place{基里基雅}人时，他说：\InnerQuote{等你的控告者也来了，我才听你。}} \Person{里息雅}已经为他辩护了；但（犹太人）企图预先赢得听者。\Quote{并命令把他看守在\Person{黑落德}的总督府里}（第34、35节）：\Person{保禄}再次被囚。\Quote{过了五天，大司祭\Person{阿纳尼雅}同一些长老下来了。}看，尽管这样，他们仍不罢休；虽然被无数障碍所阻，他们还是来了，只是为了在这里也蒙羞。\Quote{还有一个辩士，一个叫\Person{特土罗}的。}而哪里需要\Quote{一个辩士？这些人也向总督告了\Person{保禄}。}\BibleRef{宗 24:1}看这个人从最初（\Emph{b}）就企图以赞美预先赢得判官。\Quote{等\Person{保禄}被传来时，\Person{特土罗}开始控告他说：\InnerQuote{因着你，我们得享大平安，且因你的明察，为这民族成就了极美的事，我们总是、处处、最尊贵的\Person{斐理斯}，都满怀感谢地接受。}}（第2、3节）然后，因为他有许多话要说，他就略过其余的事：\Quote{但是，为了不再烦扰你，我求你仁慈地听我们几句话。因为我们发现这个人是个祸害，在天下众犹太人中间煽动叛乱。}（\Emph{a}）他愿意把他当作一个革命者和煽动叛乱者交出去。然而，也许可以回答说，是你们做了这事。（\Emph{c}）再看，他如何激起判官想要惩罚，因为他这里有机会制服这个翻天覆地的人！好像他们做成了一件功德，他们大肆宣扬：\Quote{发现了这个家伙，}等等，\Quote{煽动叛乱者，}他们说，\Quote{在天下众犹太人中间。}（如果他曾是这种人），他们会宣布他为民族的恩人和救主！\Quote{还是纳匝肋人教派的一个首领。}（第4、5节）他们认为这很可能成为一项耻辱——\Quote{纳匝肋人的：}并以此企图损害他——因为\Place{纳匝肋}是个卑微的地方。而且，\Quote{我们找到了他，}他们说：看他们多么恶毒地诽谤他：（找到他），好像他一直在避开他们，他们好不容易才捉住他：虽然他在圣殿里已经七天了！\Quote{他还曾试图亵渎圣殿；我们捉住他，[并想按我们的法律审问。]}\BibleRef{宗 24:6}看他们甚至侮辱法律；诚然，是像法律那样，就是殴打、杀害、埋伏！然后是对\Person{里息雅}的控告：他们认为他无权干预，他过度自信地把他从我们手中夺走：[ \Quote{但千夫长\Person{里息雅}来到我们中间，用极大的暴力把他从我们手中夺走，命令他的控告者到你这里来]：由你来审问他们，便可知道我们控告他的这一切事。犹太人也附和，说这些事是真的。}\BibleRef{宗 24:7}那么\Person{保禄}说什么呢？\Quote{然后\Person{保禄}，在总督示意他说话之后，回答说：\InnerQuote{我晓得你多年来作这民族的公正判官，所以我更加欣然为自己分诉。}}\BibleRef{宗 24:10}这不是谄媚的话，他证明判官的公正：不，奉承反而是那辩士的话里所有的：\Quote{因着你，我们得享大平安。}如果是这样，那么你们为什么叛乱呢？\Person{保禄}所求的是公正。\Quote{知道你是一位公正的判官，我欣然，}他说，\Quote{为自己分诉。}然后他也用时间长短来加强这一点：他（作判官）\Quote{多年。因为使你知道，从我上\Place{耶路撒冷}去朝拜到现在，不过才十二天。}\BibleRef{宗 24:11}这是什么意思？（意思是）\Quote{我不能立刻引起骚动。}因为控告者在\Place{耶路撒冷}拿不出什么（证据），看他所说的话：\Quote{在天下众犹太人中间。}因此，\Person{保禄}在这里有力地吸引他——\Quote{去朝拜，}他说，\Quote{我上来了，}我远非煽动叛乱——并强调这一点作为强点。\Quote{他们并没有看见我在圣殿里与人争辩，也没有鼓动群众，无论是在会堂里，还是在城里。}\BibleRef{宗 24:12}这其实是事实。控告者确实用了\Quote{首领}这个词，好像这是一场斗殴和叛乱；但请看\Person{保禄}如何温和地回答。\Quote{但我向你承认，我是按照他们称为异端的道，这样事奉我祖先的天主，相信法律和先知书上所记的一切事；并对天主怀着希望，这也是他们自己所承认的，就是将来死人复活，无论义人还是不义的人都要复活。}（第14、15节）控告者把他分离（当作外人），但他与法律认同，如同他们中的一员。\Quote{为此，}他说，\Quote{我操练自己，总要对天主、对人良心无愧。过了多年，我带着施舍和祭品到我民族那里去。正在行这事的时候，他们在圣殿中找到我，我已经洁净了，没有聚众，也没有骚乱。}（第16、17、18节）那么你为什么上来呢？什么把你带到这里？他说：去朝拜；去施舍。这不是一个结党者的行为。然后他也排除他们的人：\Quote{但是，}他说，（发现我的人是）\Quote{一些从\Place{亚细亚}来的犹太人，他们本该在你面前，如果他们有反对我的事，就该提出异议。否则，就让这些人自己说说，他们在我站在公议会面前时，发现我有什么恶行，除非是为了这一句话，就是我站在他们中间喊说：我今天在你们面前受审，是为死人复活。}（第19、20、21节）因为这就是极其充足的辩护，不是逃避控告者，而是预备向众人交账。\Quote{为死人复活，}他说，\Quote{我今天被审问。}而他绝口不提他将要说的事，就是他们如何密谋陷害他、如何暴力拘押他、如何埋伏他——因为这些事已经由千夫长粗略述说了——但\Person{保禄}，虽然有危险，却不这样：不，他保持沉默，只为自己辩护，虽然他有极多话可以说。（\Emph{b}）\Quote{在其中}（施舍），他说，\Quote{他们找到我正在圣殿里行洁净礼。}那么他如何亵渎了呢？因为同一个人既洁净自己、又崇拜、并为此而来，然后又亵渎，这是不合理的。这使人猜测他的事业是正义的，因为他没有陷入冗长的演说。我想，他也因此取悦了判官（也就是），使他的辩护简洁：（\Emph{d}）因为在他之前的\Person{特土罗}做了一篇长篇大论。（\Emph{f}）这也是温和的证明，即当一个人有很多话可说时，为了不烦人，他只说很少的话。（\Emph{c}）但是，让我们再看看所说的事。\par

（复述。）\Quote{于是兵士们}等。\BibleRef{宗 23:31}（a）保禄带着这么多军队来到，也使他在凯撒勒雅出了名。\Quote{但是，}特图罗说，\Quote{为免我过于烦琐，}（e）表明（斐理斯）确实觉得他烦琐（\OriginalTerm{被妨碍}{\GreekText{ἐ} \GreekText{γκόπτεται}}）：\Quote{我恳求你，}他没有说\Quote{请听此事}，而是\Quote{请宽容听我们说。}\BibleRef{宗 24:4}大概是为了讨好，他才这样安排自己的演说。（g）\Quote{我们看这个人是个瘟疫，在所有犹太人中煽动叛乱，}\BibleRef{宗 24:5}那么，若他在别处（而不在此）做了这事，又怎能呢？不，他说；在我们当中他也亵渎了圣殿；\Quote{试图，}他说，\Quote{亵渎它：}但他没有说明如何亵渎。\Quote{我们也把他拿住了。}等。\Quote{但是千夫长}等。他如此夸大千夫长所做的事，而看他如何淡化控告者自己的部分。\Quote{我们拿住了他，}他说，\Quote{并愿意按我们的法律审判他。}\BibleRef{宗 24:6}他表明他们不得不来到外邦法庭是一种磨难，若非千夫长强迫他们，他们本不会麻烦他；且千夫长与这事无关，却强行拿住了这人：因为事实上，这些错事是针对我们的，法庭本该在我们这里。因为这就是此意，请看下文：\Quote{用极大的暴力，}\BibleRef{宗 24:7}他说。因为这种行为就是暴力。\Quote{从他你可以知道。}他既不敢控告他（千夫长）——因为那人确实是宽容的——也没有完全放过。接着，为免他显得在撒谎，他引出保禄自己作为他的控告者。\Quote{从他，你审问他，就可以知道这一切事了。}\BibleRef{宗 24:8}接着，作为所说之事的证人，这些控告者，同样的人既是证人又是控告者：\Quote{犹太人也附和}等。\BibleRef{宗 24:9}但保禄说：\Quote{我知道你多年作公正的法官。}\BibleRef{宗 24:10}那么，他既不是陌生人、外侨或叛乱分子，因他认识这位法官已多年。他恰当加上\Quote{公正}这一称呼，以使他（斐理斯）不看大司祭、百姓或控告者的面子。看，他如何没有让自己陷入辱骂，尽管受到了强烈的挑衅。\Quote{相信，}他说，\Quote{会有复活：}现在一个相信复活的人，绝不会做这样的事——\Quote{这（复活）他们自己也承认。}\BibleRef{宗 24:15}他不是说他们相信\Quote{先知书所写的一切：}是他相信一切，而不是他们：但如何是\Quote{一切}，需要长篇论证。他绝口不提基督。在这里他说\Quote{相信}，就（实际上）引入了与基督有关的事；目前他专讲复活，这教义也是他们共同的，并且消除了任何叛乱的嫌疑。关于他上耶路撒冷的原因，\Quote{我来了，}他说，\Quote{是为给我的民族带来赈济和奉献。}\BibleRef{宗 24:17}那么，我怎会扰乱那些我长途跋涉来为之奉献的人呢？\Quote{没有聚众，也没有喧嚷。}\BibleRef{宗 24:18}处处他都消除叛乱的指控。他也做得很好，向那些来自亚细亚的控告者挑战，\Quote{他们应该到你面前来控告}等，但他也不拒绝这个：\Quote{否则，}他说，\Quote{就让这些本人说。关于死人复活，}等。（第19、20、21节）实际上，他们从起初就因他宣讲复活而大为困扰。这一点证明后，与基督有关的事也容易引入了，即祂复活了。\Quote{什么恶事，}他说，\Quote{他们在公会里从我身上找到了。}\BibleRef{宗 4:2}他说的是：审讯不是在私下进行的。我所说的这些事是真实的，那些控告我的人可以作证。\Quote{我常存，}他说，\Quote{无愧的良心，对天主、对人。}\BibleRef{宗 24:16}这便是德性的顶峰，即使在人们面前我们也不给他们攻击的把柄，并留心在天主面前保持无愧。\Quote{我在公会里喊叫，}他说。他也显示了他们的暴力。他们不能说：你以赈济为借口做了这些事；因为（它）\Quote{没有聚众，也没有喧嚷：}尤其是关于此事作了调查之后，再没有发现什么。你注意到他的节制吗，尽管有危险？你注意到他如何约束舌头不恶言伤人，他只寻求一件事，即为自己解脱控告，而不是要控诉他们，除非在自卫中不得已而为之？正如基督也说过：\Quote{我没有鬼，我尊敬我的父；你们却侮辱我。}\BibleRef{若 8:49}\par

我们该效法他，因为他也是基督的效法者。如果他在面对甚至行凶杀人的敌人时，没有对他们说任何冒犯的话，那么我们在辱骂和凌辱中暴怒，称我们的敌人为恶棍、可憎之徒，又该得到什么宽恕呢？我们竟然有敌人，又该得到什么宽恕呢？你们没有听见吗？尊荣他人就是尊荣自己。的确如此：但我们却羞辱自己。你控告某人辱骂了你：那你为何使自己处于同样的控告之下？为何要打击自己？保持无欲，保持无伤：不要因想打击他人而使自己受伤。我们灵魂的骚动还不够多吗？那骚动即使无人煽动也会兴起——例如，无耻的欲望、忧伤和痛苦，以及诸如此类——我们还要再堆上另一堆吗？你会说，当人受到侮辱和谩骂时，怎能忍受呢？我问你，又怎能不忍呢？言语会造成伤口吗？言语会在我们的身体上留下伤痕吗？那么，对我们有何伤害呢？所以，只要我们愿意，我们就能忍受。我们要为自己立下一条不忧伤的法律，我们就能忍受：我们要对自己说：\Quote{这不是出于敌意，而是出于软弱}——因为那确实是出于软弱；因为，为了证明这不是出于敌意或性情的恶意，而是出于软弱，对方也本想克制（他的愤怒），尽管他遭受了无数次的冤屈。只要我们心里有这个念头，即这是出于软弱，我们就能忍受，同时宽恕冒犯者，并努力不让自己陷入其中。我请问你们所有在场的人：你们是否愿意能够拥有这样的理性平静，以忍受那些侮辱你们的人？我想是的。那么，他是不情愿地侮辱了人；他宁愿不这样做，但他做了，被他的情欲所迫：管束你自己吧。你们没有看见附魔者（发作时）的情形吗？正如他们那样，他也一样：与其说是出于敌意，不如说是出于软弱（他才那样做）：忍受吧。至于我们——我们动怒，与其说是由于侮辱本身，不如说是由于我们自己；否则，当狂人向我们施以同样的侮辱时，我们为何能忍受呢？再者，如果侮辱我们的是我们的朋友，我们也能忍受；或者是我们的上司，我们也能忍受：那么，在这三种情况——朋友、狂人和上司——中我们都能忍受，而当对方与我们平级或低于我们时，我们却不能忍受，这岂不是荒谬吗？我常说过：这只是一时的冲动，某种突然冲昏我们头脑的东西：让我们稍加忍耐，我们就能忍受整个事情。侮辱越大，冒犯者越软弱。你知道什么时候我们应该悲伤吗？当我们侮辱了别人，而对方保持沉默时：因为那时他是坚强的，我们是软弱的；但如果情况相反，你甚至应当喜乐：你已被加冕，被宣布为胜利者，甚至没有进入竞技，没有忍受日晒、炎热和尘土，没有与对手扭打并让他靠近你；仅仅是你一个意愿，或坐或站，你就得到了一个伟大的冠冕：一个远大于那些（角斗者）赢得的冠冕：因为打倒一个迎面而来的敌人，远不如克服愤怒的投枪那样伟大。你胜利了，甚至没有让他靠近你，你打倒了自己内心的情欲，杀死了被激起的野兽，平熄了愤怒的狂澜，如同一个优秀的牧人。这场战斗本是一场内乱，一场内战。因为，正如那些从外围攻的人设法使（被围者）陷入内乱，然后战胜他们；同样，侮辱者除非激起我们内心的情欲，否则不能战胜我们：除非我们自己点燃火焰，他就没有力量。让愤怒的火星存在于我们之内，以便在适当的时刻准备点燃，但不是针对我们自己，也不是让我们陷入无数祸患。你们没有看见房屋里的火是如何被隔离起来，不随意到处乱扔，既不在稻草中，也不在亚麻布中，也不在可能的地方，以免风吹时危险，引起火焰吗？但是无论女仆拿着灯，还是厨子生火，都有许多嘱咐，不要在有风的地方、靠近木墙或在夜间做这事：但当夜晚来临时，我们熄灭火，恐怕当我们睡着且无人帮忙时，它烧起来，将我们全部烧毁。对愤怒也当如此：不要让它四处散落在我们的思想中，而应让它存在于心灵的某个深处，使那从反对我们之人的话语中升起的风不容易吹到它，而是让它从我们自己——我们知道如何适度安全地激怒它——那里承受（激起它的）风。如果它承受从外来的风，它就不懂得节制；它会把一切烧毁：常常当我们睡着时，这风会吹到它，将一切烧尽。因此，让它以这样的方式在我们那里（安全保存），只点燃一盏灯：因为当愤怒被恰当管理时，确实能点燃一盏灯：让我们拿着火炬对付那些冤枉别人的人，对付魔鬼。不要让火星随意在任何地方，也不要被扔来扔去；我们把它安全地保存在灰烬之下：在谦卑的思想中让它沉睡。我们并非总需要它，但在需要制服、使之柔顺、软化顽硬、确证灵魂时才需要。愤怒和暴躁的情欲造成了多少罪恶！更严重的是，当我们分离后，我们不再有能力重新和好，而是等待别人（来做这事）：每个人都羞愧，不好意思自己回去与对方和好。看，他并不羞愧于分离和分开；不，他带头作恶：但出来修补撕裂之处，他却羞于去做：这就像一个人不介意切断肢体，却羞于再把它接上一样。你说什么，人啊？你受了很大的伤害，自己造成了争吵？那么，你就应当首先去和好，因为你提供了原因。但对方做了错事，他是敌意的原因？那么，正因如此，你也必须去做，使人们更加钦佩你，使你在这件事上也获得头奖：正如你不是敌意的起因，你也不是敌意延长的原因。也许对方，由于自己内心意识到无数恶事，感到羞耻并脸红。但他傲慢吗？尤其为此，你应当毫不犹豫地跑去见他：因为如果他身上的病是双重的——既傲慢又愤怒——那么你刚才提到恰恰是你应当先去见他的理由，因为你是健康的人，是能够看见的人：而他是在黑暗中：愤怒和虚伪的骄傲就是这样。但你这不受这些影响、身体健康的人，去见他吧——你这位医生，去见病人吧。有哪位医生说，因为某人生病，我就不去见他吗？不，这恰恰是为什么要去的原因，当他们看到病人不能来见他们时。因为对于那些能够（来）的人，他们不太在意，仿佛他们不是病得很重；但对那些卧病在家的人则不然。骄傲和愤怒，难道不比任何疾病更糟糕吗？难道其中一个不像高烧，另一个不像发炎肿胀的身体吗？想一想，发烧发炎是什么样子：去见他，扑灭那火，因为靠着天主的恩宠你能做到，去，像用水一样降温。\Quote{但是，}你会说，\Quote{如果我这样做，他只会更加傲慢呢？}这与你无关：你已尽了你的本分，让他自己交账吧：不要让我们的良心谴责我们，说这事的发生是由于我们没有做应做的事。圣经说：\BibleQuote{你这样做，} \BibleQuote{就是把炭火堆在他头上。}\BibleRef{罗 12:20} 尽管如此，圣经还是命我们去和好、行善——不是要堆炭火，而是要（我们的敌人）知道将来的后果，好被当前的善意所安抚，使他战栗，使他惧怕我们的善行甚于我们的敌意，我们的友谊甚于我们的恶意。因为人通过表示愤怒和敌意，并不如通过行善和表示善意那样伤害他的仇敌。因为通过愤怒，他既伤害了自己，也可能在某种程度上伤害了对方；但通过行善，他却把炭火堆在了对方头上。\Quote{那么，}你会说，\Quote{为了不堆炭火，就不该这样做（\Emph{b}），而应使敌意持续更久。}绝非如此：不是你造成这结果，而是他野蛮的性情造成的。因为如果你向他行善、尊敬他、愿意和好，而他仍坚持敌意，那就是他自己点燃了火，把他自己的头烧了；你是无罪的。不要想比天主更仁慈（\Emph{b}），或者，如果你愿意，你也不能，连一点点也不能。怎能呢？圣经说：\BibleQuote{天离地有多高，} \BibleQuote{我的谋略就离你们的谋略有多高}\BibleRef{依 45:8}；又说：\BibleQuote{你们纵然邪恶，尚且知道把好东西给你们的儿女，} \BibleQuote{何况你们的天父}\BibleRef{玛 7:11}？但实际上，这些话只是借口和搪塞。我们不要在天主的诫命上支吾其事。\Quote{我们怎么支吾呢？}你会说。他说：\BibleQuote{你这样做，就是把炭火堆在他头上}；而你说，我不喜欢这样做。（\Emph{a}）但你是否愿意以另一种方式堆炭火，即堆在你自己的头上呢？因为实际上，愤怒正是这样做的：（\Emph{c}）因为你将遭受无数祸患。（\Emph{e}）你说：\Quote{我害怕为我的敌人，因为他对我造成了很大的伤害：}实际上，你是这么说吗？但你为什么会有敌人？但你为什么仇恨你的敌人？你为那伤害你的人害怕，却不为自己害怕吗？但愿你能关心你自己！不要以这样的目的去做（善行）；或者，即使以此为目的，也要去做。但你根本不做。我不对你说：\BibleQuote{你将把炭火堆在他头上}；不，我说另一件更重要的事：你尽管去做吧。因为保禄这样说，只是召叫你们（只要）在对报复的盼望中结束敌意。因为我们本性如同野兽般野蛮，除非期待某种报复，否则我们不肯爱仇敌，所以他（保禄）提供这个如同给野兽的一块饼。但对于宗徒们，（主）没有这样说，而是说什么？\BibleQuote{好使你们成为你们在天之父的子女。}\BibleRef{玛 5:45} 此外，施恩者和受恩者之间不可能保持敌意。这就是为什么保禄这样说。为什么，你在言语上装作高尚慷慨的原则，但在行为上却连普通的节制都不遵守呢？（听起来）很好：你不喂养他，因为怕因此把炭火堆在他头上；那么，你饶了他？那么，你爱他，你这样做是为了这个目的？天主知道，你这样说时是否有这个目的，还是把这番话当作借口和搪塞来敷衍我们。你关心你的敌人，你怕他被惩罚，那么你当先熄灭了你的愤怒吧？因为爱到如此程度，以致为了对方的好处而忽略自己的利益，这样的人就没有敌人。（那时）你才可以这样说。我们还要在不可儿戏、不可推诿的事情上儿戏多久呢？所以我恳求你们，让我们除去这些借口；不要轻视天主的法律：好使我们能蒙主喜悦，度过今生，并获得所应许的美善，因我们的主耶稣基督的恩宠和慈爱，愿光荣、权能、尊崇归于他和圣父及圣神，自今至永远，世世无尽。阿们。\par

\SourceRecord{Translated by J. Walker, J. Sheppard and H. Browne, and revised by George B. Stevens.；Nicene and Post-Nicene Fathers, First Series；Vol. 11.；Edited by Philip Schaff.；Buffalo, NY: Christian Literature Publishing Co.,；1889.；<http://www.newadvent.org/fathers/210150.htm>.}

% structure: p0001 source=h1 target=section reason=page-title

\section{\BibleBook{}{宗徒大事录}讲道 第五十一讲}

\BibleRef{宗 24:22-23}\par

\Quote{斐理斯听见了这些事，对这道有更确切的知识，就故意拖延他们，并说：\InnerQuote{等千夫长里息雅下来，我再彻底审断你们的事。}他又命令一个百夫长看守保禄，但要让他自由，并且不许禁止他的任何亲友来服事他或到他那里去。}\par

请看，许多人在长时间内进行了多么仔细的调查，以免有人说审判是草率进行的。因为，当那演说家提到吕息雅时，说他\Quote{用暴力带走他，斐理斯}，他说，\Quote{拖延了他们。既知道那道}：也就是说，他故意拖延他们；不是因为想知道，而是想摆脱犹太人。为了他们的缘故，他不愿释放他；惩罚他又不可能；那会太露骨了。\Quote{并且让他得自由，也不禁止任何熟人服事他。} 他也完全免了他的控告。然而，为讨他们喜欢，他拘留了他，此外，因期望得到钱财，他召见了保禄。\BibleQuote{过了几天，斐理斯和他的妻子杜息拉——她是犹太人——一同来，就打发人请保禄来，听他讲论关于在基督内的信德。保禄讲论文义德、节操（即节制或贞洁）和将来的审判时，斐理斯感到恐惧，回答说：\InnerQuote{你现在去吧，等我有方便的时候，再叫你来。}他也希望保禄给他钱，好释放他，因此屡次叫他来，和他谈论。过了两年，颇尔基约·斐斯托接了斐理斯的任；斐理斯愿意向犹太人讨好，就把保禄留在监里。}\BibleRef{宗 24:24} 请看这些所写的事与真理多么接近。但他屡次叫他来，不是因为他钦佩他，也不是因为他称赞所说的话，也不是因为他愿意相信，而是为什么？经上说：\Quote{期望}，\Quote{给他钱。} 请看他在这里毫不掩饰法官的心态。\Quote{因此他打发人去叫他}等。而且，如果他定了他的罪，就不会这样做，也不会愿意听一个被定罪和品行恶劣的人。并且请注意保禄，尽管与一位长官辩论，他并没有说任何可能取悦和娱乐的话，而是（经上说，他\Quote{辩论}）\Quote{关于义德、将来的审判}以及复活。他的话语如此有力，甚至使总督感到恐惧。这人被另一个人接替了职位，他留下保禄作囚犯；然而他不应该这样做；他本应了结此事；但他为了讨好他们而留下他。然而他们如此焦急，以至于再次恳求法官。但是，他们从未如此顽固地反对过任何一位宗徒；在那里，当他们攻击时，立刻又停止了。所以，他得以有上智的安排离开耶路撒冷，得以与这些野兽周旋。尽管如此，他们仍要求把他提到那里受审。\BibleQuote{现在斐斯托到了省里，过了三天，就从凯撒勒雅上耶路撒冷去。大司祭和犹太人的首领向他控告保禄，并恳求他，求他开恩反对保禄，叫人把他提到耶路撒冷来，在路上埋伏要杀他。}\BibleRef{宗 25:1} 这时，天主的眷顾介入，不让总督这样做；因为他刚上任，自然愿意讨好他们；但天主不容许他。\BibleQuote{斐斯托却回答说：\InnerQuote{保禄应押在凯撒勒雅，我自己快要往那里去。}又说：\InnerQuote{你们中间有能力的，可以和我一同下去，若在这人身上有什么恶事，就控告他。}他在他们那里住了不过十天，就下到凯撒勒雅；第二天坐在审判席上，吩咐把保禄带上来。}\BibleRef{宗 25:4} 但他们在下来后，立即无耻地、更加激烈地提出控诉；并且由于他们无法在法律方面定他的罪，他们就照常挑起了关于凯撒的问题，正如他们对基督所做的那样。他们诉诸于此的事实是明显的，因为保禄为自己辩护是针对冒犯凯撒的罪名。\BibleQuote{当保禄被带来时，从耶路撒冷下来的犹太人围着他，提出许多严重的控诉，都是他们不能证实的。保禄为自己辩护说：\InnerQuote{我既没有干犯犹太人的法律，也没有干犯圣殿，更没有干犯凯撒。}但斐斯托愿意向犹太人讨好，回答保禄说：\InnerQuote{你愿意上耶路撒冷去，在那里在我面前受审判吗？}}\BibleRef{宗 25:7} 因此他也讨好了犹太人、全民族和那城。既然如此，保禄也使他恐惧，使用人的武器来自卫。\BibleQuote{保禄却说：\InnerQuote{我站在凯撒的审判台前，这是我应当受审的地方；我并没有得罪犹太人，你自己也知道得很清楚。我若行了不义，犯了什么该死的罪，就是死我也不辞；但若这些人所控告我的都是虚妄，就没有人能把我交给他们。我上诉凯撒。}}\BibleRef{宗 25:10-11} 有人可能会说：既然主曾告诉他：\BibleQuote{你也要在罗马为我作证}，\BibleRef{宗 23:11} 他为何好像不信似的这样做？断乎不可！不，他这样做，是因为他深信不疑。因为若是仗着那个宣布而胆大妄为，将自己投入无数危险，并说：\Quote{让我们看看天主是否还能这样拯救我。} 那就是试探天主。但保禄不是这样；不，他尽自己的本分，竭尽全力，将一切交托给天主。他也平静地责备总督：\Quote{如果，他说，我是罪犯，你做得对；但如果不是，你为什么要把我交出去？}\Quote{没有人}他说，\Quote{能把我当作牺牲品。} 他使总督害怕，以致即使他愿意，也不能把他牺牲给他们；同时，作为对他们的借口，他可以说保禄上诉了。\BibleQuote{于是斐斯托与议会议论之后，回答说：\InnerQuote{你上诉凯撒吗？你就到凯撒那里去。}过了几天，阿格黎帕王和贝勒尼切来到凯撒勒雅，向斐斯托问安。}\BibleRef{宗 25:12-13} 请注意，他把这事告诉阿格黎帕，以便有更多的听众，包括王、军队和贝勒尼切。然后，他为自己的无罪辩护。\BibleQuote{他们在那里住了多日，斐斯托将保禄的案件向王陈述说：\InnerQuote{这里有一个人，是斐理斯留下的囚犯。我在耶路撒冷的时候，犹太人的司祭长和长老告发他，要求定他的罪。我回答他们说：\InnerQuote{罗马人没有把人送来判刑之前，先叫被告与原告对质，并准许他为自己辩护的惯例。}因此，当他们来到这里，我没有耽误，第二天就坐在审判席上，吩咐把那人带上来。控告他的人起来，并没有提出我所预料的那类罪状；但有一些关于他们自己的宗教和一位名叫耶稣的死人，保禄却说他还活着的争论。我对这些争论感到疑惑，就问他是否愿意上耶路撒冷去在那里受审。当保禄上诉请求保留给奥古斯都审理时，我下令把他看管起来，等我把他解送凯撒。}阿格黎帕对斐斯托说：\InnerQuote{我也愿意亲自听听这人。}他说：\InnerQuote{明天你就可以听他。}}\BibleRef{宗 25:14} 请注意，对犹太人的指责不仅来自保禄，也来自总督。他说：\Quote{渴望对他作出判决。} 我对他们说，令他们羞愧的是：\Quote{罗马人的规矩不是}在给予人自辩的机会之前\Quote{把人牺牲掉。}但我确实给了他（这样的机会），却找不出他有什么过失。他说：\Quote{因为我对这些问题感到疑惑。}他也对自己的错误有所掩饰。于是那人想要见他。\Emph{乙}但让我们再看看所说的话。\par

（复述。）\Quote{且斐理斯}等等。\BibleRef{宗 24:22}试看，总督们处处设法摆脱犹太人给他们的麻烦，常常被迫违背正义行事，寻求拖延的借口；显然他拖延案件并非出于无知，而是明知故为。他的妻子也与总督一同听讲。\BibleRef{宗 24:24}这在我看来显示极大的尊重。因为若不是他对保禄极为推崇，他不会带妻子一同出席听审。在我看来她也渴望此事。试看保禄如何立即谈论，不仅关乎信德，不仅关乎罪过的赦免，也关乎实践的责任。\Quote{你暂且去吧，}他说，\Quote{等我有机会时再叫你来。}\BibleRef{宗 24:25}试看他的心硬：听了这些话，\Quote{竟指望从他手中得钱！}\BibleRef{宗 24:26}不仅如此，甚至在与他交谈之后——因为他的任期将尽——他把他留下捆绑，\Quote{愿意讨好犹太人}\BibleRef{宗 24:27}：可见他不仅贪财，也爱虚荣。可悲啊，你怎能期望从一个宣讲相反道理的人那里得钱？他没有得到钱，这从留下捆绑可以明显看出；若他收到了钱，他早就释放他了。\Quote{关于节操，}经上说，他辩论，而对方却渴望从宣讲这些事的人那里得钱！他不敢开口索要，因为邪恶如此；但他存着希望。\Quote{过了两年，}等等，因此他自然要讨好他们，因为他已在那地做总督这么久。\Quote{斐斯托到了省里，}等等。\EditorialNote{ch. 25:1, 2}一开始，司祭们就来到他面前，他们本会毫不犹豫地前往凯撒勒雅，除非他立即上到那里，因为他一到他们就来找他。他花了十天时间，我想是为了让那些想用贿赂腐蚀他的人有机会。保禄却在监里。\Quote{他们请求他，}经上说，\Quote{派人将保禄解来：}如果他是该死的人，他们为何当作恩惠来求？然而他们的阴谋甚至在他面前也暴露了，以致在（向阿格黎帕）讲述这事时，他说：\Quote{想要定他的罪。}他们想促使他立即判决，因为害怕保禄的口才。你怕什么？你为何如此匆忙？事实上，\Quote{要把他留住}\BibleRef{宗 25:4}这句话表明了这一点。他想逃走吗？\Quote{所以，}他说，\Quote{你们中间有能的人，就同他一起下去控告他。}\BibleRef{宗 25:5}又是控告者，又是凯撒勒雅，又是保禄被带出来。他一来到，立即\Quote{坐在审判台上}\BibleRef{宗 25:6}，极其匆忙；他们如此催逼，如此促使他。当他还不认识犹太人，也未经验到他们对他的尊敬时，他回答得正确；但如今他在耶路撒冷住了十天，他也想讨好他们（牺牲保禄）；然后，也要收留保禄，\Quote{你愿意，}他说，\Quote{在那里由我审判这些事吗？}\BibleRef{宗 25:9}我不是把你交给他们——但事实正是如此——他把选择权留给保禄自己，想用这种尊敬使他屈服：因为判决权在他手中，而他在此地没有被定罪，却要把他带回那里，这太露骨了。\Quote{保禄却说：我站在凯撒的审判台前，}等等。\BibleRef{宗 25:10}他没有说\Quote{我不愿意}，以免使审判官更加激烈，而是（此处）又显出他的大无畏：他们自己一旦把我赶出去，就想以此定我的罪，要表明我冒犯了凯撒：我愿在他的审判台前受审，在那受伤害者本人的台前。\Quote{我没有对犹太人做任何不义的事，你也知道得很清楚。}此时他责备了他，因为他也想牺牲保禄讨好犹太人；然后，另一方面，他缓和了话语的严厉：\Quote{我若有不法的事，或做了任何该死的事，我不拒绝死。}我给自己下判决。因为与无畏的言辞一起，必须有正义的理由，才能叫对方惭愧。\Quote{但若这些控告我的事毫无根据，没有人}——无论他多么愿意——\Quote{没有人可以把我讨好他们而牺牲掉。}他没有说\Quote{我不该受死}，也没有说\Quote{我该被释放}，而是说\Quote{我准备在凯撒面前受审}。同时，他也因想起那梦而更有信心上诉。\BibleRef{宗 23:11}他没有说\Quote{你不能}，而是说\Quote{也没有任何其他人可以牺牲我}，以免冒犯他。\Quote{斐斯托与议会商量之后}——你看见他如何寻求讨好他们吗？因为这就是偏袒——\Quote{商量之后，}经上说，\Quote{就说：你已上诉凯撒了吗？你必到凯撒那里去。}\BibleRef{宗 25:12}请看他的审判又如何延长，而对他的阴谋成了宣讲的机会：使他能轻易且安全地被解往罗马，无人图谋害他；因为他单纯地去那里，与因这样的案件去那里并不相同。事实上，这正是使犹太人在那里聚集的原因。\BibleRef{宗 28:17}然后，又过了一些时间，他在耶路撒冷停留，好让你知道，虽然时间过去，对他的毒计毫无作用，因为天主不允许。但这阿格黎帕王，也是一位黑落德，是雅各伯时代之后的那一位，所以这是第四位（黑落德）。请看他的仇敌如何违反己意与他合作。为使听众众多，阿格黎帕产生了听讲的愿望；而且他不只是听听，而是大张旗鼓。请看这多么大的辩护\OriginalTerm{辩护}{\GreekText{απολογιαν}}！斐斯托如此写道，犹太人的凶残被公开揭露：因为当总督说这些话时，他是最无可怀疑的证人；因此犹太人也被他定案。因为当所有人都判了他们败诉时，那时，且非更早，天主降罚于他们。但请看：里息雅定他们败诉，斐理斯定他们败诉，斐斯托定他们败诉——虽然他想讨好他们——阿格黎帕也定他们败诉。还有什么？法利塞人——连他们也定了自己败诉。\Quote{没有恶事，}斐斯托说，\Quote{像我所想的；他们没有任何控告能成立。}\BibleRef{宗 25:18}然而他们确实控告了：不错，但他们没有证明；因为他们的毒计和胆大包天的阴谋使人猜测如此，但审讯并未得出此类结果。\Quote{关于一位已死的耶稣，}他说。\BibleRef{宗 25:19}他自然地说\Quote{一位}（耶稣），因为他是一个官员，不关心这些事。\Quote{至于我自己，不知如何就这些事进行调查}\BibleRef{宗 25:20}——当然，审查这些事超过了法官的职责。你若不知，为何把他解往耶路撒冷？但另一位不愿屈尊：不，\Quote{到凯撒那里}（他说）；事实上，他们正是以凯撒的事控告他。你听见上诉了吗？听见犹太人的阴谋了吗？听见他们的党派精神了吗？这一切都激起他听讲的愿望：他满足了他们的愿望，保禄更出名了。因为正如我所言，这就是（仇敌的）毒计。若不是这样，这些统治者没有一位会屈尊听他，也没有一位会如此安静沉默地听。他似乎是在教导，似乎在辩护；但毋宁说，他是在非常有序地公开发表演说。那么，我们不要以为别人陷害我们是什么痛苦的事。只要我们不自害，就没有人能害我们：或者，人可以这样做，但他们不能伤害我们；不，甚至对我们有极大的益处：因为自受不受恶是在我们自己。看！我作证，并大声宣告，比号角的声音更刺耳——若可能登高呼喊，我也不会退缩——基督徒，地上所有人类都不能伤害他。我为何说人类？连那恶神自己、暴君、魔鬼也不能，除非人伤害自己；无论任何人做什么，都是徒劳。因为正如天使在地上，没有人能伤害他，同样，一个人也不能伤害另一个人。但反过来，只要他是善的，他自己也不能伤害别人。还有什么能与此相比？既不能受伤害，也不能伤害别人？因为这不比前者小：不愿伤害别人。那人是一种天使，是的，像天主。因为天主正是如此；只不过祂本质如此，而这人是藉道德选择：既不能受伤害（对两者而言），也不能伤害别人。但这\Quote{不能}，不要以为是由于能力不足——相反，缺乏能力——不，我说的是（道德上）不可容许之事\OriginalTerm{（道德上）不可容许的事}{\GreekText{τὸ} \GreekText{ἀνενδεκτόν}}。因为（神圣）本性既不能受伤害，也不能伤害别人：因为这伤害本身就是伤害。因为我们伤害自己，无非是藉伤害别人，而我们最大的罪正是因伤害自己而成为的。所以，为此理由基督徒也不能被伤害，即因为他也不能伤害人。但如何藉伤害别人而伤害自己，来，让我们仔细审察这句论断。假设一个人冤枉别人、侮辱、欺骗；那么他伤害了谁？岂不首先是自己？这对每个人都是清楚的。因为对那人，损失在金钱；对自己，则在灵魂，以致毁灭和惩罚。再假设另一个人嫉妒：他岂不伤害了自己？因为不义的本性正是如此：它首先给它的作者带来不可估量的伤害。\Quote{是的，但对别人也有？}不错，但不足挂齿；或者，甚至毫厘——不，反而有益于他。因为假设——此事大多在这类例子中——假设有一个穷人，拥有很少财产，只够糊口，另有一个富人，有钱有势，然后富人夺去穷人的财产，剥光他，使他挨饿，他自己则用不当取来的东西享乐：他不但没有伤害那人——甚至有益于他，而他自己不但没有得益，反而受了伤害。怎能不是这样？首先，被恶良心折磨，日日自责，且被众人谴责；然后，在将来的审判中。那人又如何得益？因为受恶并勇敢忍受，就是大利：因为受恶是除罪，是修练哲学，是德行的操练。让我们看看两人中哪一个境况不好，这人或那人。因为一个若是心灵有秩序的人，会勇敢忍受；另一个则日日处于持续的颤栗和疑虑中；那么哪一个受伤害，这人或那人？\Quote{你胡说，}你说：\Quote{因为当一个人没东西吃，被迫哀叹、感到极其悲惨，或者来乞讨而一无所获，这岂不毁了灵魂和身体？}不，是你胡说：因为我要用事实证明。请问，难道没有富人感到自己悲惨吗？那么，贫穷是他悲惨的原因吗？\Quote{但他没饿死。}那又怎样？惩罚更大，因为拥有财富却如此行。因为财富不使人坚强，贫穷不使人软弱：否则，没有一个富人活得悲惨，也没有一个穷人（不）咒骂命运。但我将藉此证明你确实是胡说。保禄在贫穷中还是在财富中？他挨饿还是不？你可以听他自己说：\Quote{在飢渴之中。}\BibleRef{格后 11:27}先知们挨饿吗？他们也曾艰苦度日。\Quote{你又提保禄，又提先知，十个二十个人。}但我从哪里举例子？\Quote{从众人中给我看一些勇敢忍受恶的人。}但罕见者总是如此；然而，你若愿意，让我们按其本质审查此事。让我们看谁的忧虑更大更尖锐，谁的更易承受。一个忧虑必需的粮食，另一个从那种忧虑中解脱，却忧虑无数的事。富人不惧怕饥饿，但惧怕其他事；常常甚至惧怕生命本身。穷人并非不为食物焦虑，但他没有其他焦虑，他有安全、有平静、有保障。\par

如果伤害他人不是邪恶，而是良善，那为何我们感到羞耻？为何我们遮掩面孔？为何被责备时，我们烦恼不安？如果被伤害不是一件好事，那为何我们以此自傲、夸耀，并为自己辩护？\newline{}你想知道为何这比那更好吗？观察那些处于一种状况和另一种状况中的人。为何有法律？为何有法庭？为何有刑罚？岂不正是因为那些人，如同患病和不健康的人吗？但你或许会说，快乐是巨大的。\newline{}我们先不谈将来：让我们看现在。还有谁比一个遭受如此猜疑的人更糟？还有谁更不稳定？还有谁更不健康？他不是总处于船难之中吗？即使他做了任何公正的事，也不被信任，因为所有人因他的（伤害）能力而定他的罪：因为凡与他同住的人都是他的控告者；他不能享受友谊：因为没有人会轻易选择成为一个有如此品性之人的朋友，恐怕因对他的看法而牵连自己。如同对待野兽，所有人都远离他；如同对待瘟疫、仇敌、杀人犯和自然的敌人，他们躲避不义之人。\newline{}如果那个伤害他人的人碰巧被带到法庭，他甚至不需要原告，他的品性代替任何原告定他的罪。被伤害的人则不然；他有所有人站在他一边，同情他，伸出援助之手：他立于安全之地。\newline{}如果伤害他人是良善和安全的事，让任何人承认他是不义的；但如果他不敢这样做，那他为何还追求它作为一件好事呢？\newline{}但让我们在我们自身中看看，如果同样的事发生在那里，会带来什么邪恶：（我的意思是，）如果我们内在的任何部分或功能越过了其适当界限，攫取其他部分的职能。因为如果脾脏，如果它愿意，离开了它的适当位置，并夺取了属于其他器官的部分连同其自己的部分，这难道不是疾病吗？我们体内的湿气，如果充满每个地方，难道不是水肿和痛风吗？这难道不是既毁灭自己，也毁灭其他吗？同样，如果胆汁寻求更宽的空间，如果血液扩散到每个部分。\newline{}但在灵魂中，愤怒、贪欲和其余一切，如果食物超过其适当量度，又会怎样？同样在身体中，如果眼睛想要摄取更多，或看到比分配给它的更多，或接受比适当更强的光。但如果，当光是好的时，眼睛却因选择看见比正确更多而毁灭：那么在一件邪恶的事上会是怎样？如果耳朵接收（太）大的声音，感官就被震聋：如果思想推理超越自身的事物，它就被压倒：一切过度都会破坏一切。因为这正是\OriginalTerm{贪婪}{\GreekText{πλεονεξία}}，想要拥有比规定和分配更多的。\newline{}关于金钱也是如此：当我们非要给自己加上（比适当）更多的负担时，虽然我们没有觉察到，却在自己里面滋养了一头野兽，给我们造成极大的伤害；拥有很多，却仍需要很多，我们使自己纠缠于无数的忧虑，给魔鬼提供许多对付我们的把柄。然而，对于富人，魔鬼甚至不需要费力，因为他们自己的生意事务必定会毁灭他们。\newline{}因此，我恳求你们戒绝这些事物的贪欲，使我们能够逃脱恶者的陷阱，并抓住德行，获得永恒的美善，藉着我们的主耶稣基督的恩宠和慈爱，祂和圣父及圣神，一同永生永王。阿们。\par

\SourceRecord{Translated by J. Walker, J. Sheppard and H. Browne, and revised by George B. Stevens.；Nicene and Post-Nicene Fathers, First Series；Vol. 11.；Edited by Philip Schaff.；Buffalo, NY: Christian Literature Publishing Co.,；1889.；<http://www.newadvent.org/fathers/210151.htm>.}

% structure: p0001 source=h1 target=section reason=page-title

\section{\WorkTitle{宗徒大事录讲道集 52}}

\BibleRef{宗 25:23}\par

\BibleQuote{第二天，阿格黎帕和王后贝勒尼切大张威仪而来，进了厅堂，同千夫长和城里的显要人物在一起；非斯都一声令下，保禄便被带上来。}\par

看哪，为保禄聚集了怎样的听众！总督召集了他所有的卫队，同国王和千夫长，以及城中的首要人物前来。经上说：\Quote{与城中的首要人物。}然后保禄被带上来，看他如何被宣布为胜利者。斐斯托自己宣告他无罪，因为斐斯托说什么？\Quote{斐斯托说：阿格里帕王和一切与我们同在这里的人，你们看这个人，为了他，犹太人的全体群众曾在耶路撒冷和这里向我恳求，喊叫说他不该再活着。但我查明他没有犯什么该死的罪，而且他自己已向奥古斯都上诉，我便决定把他解去。我没有什么确切的事可向我主奏报，因此我带他到你们面前，尤其在你阿格里帕王面前，经过审查后，我或许有所可写。因为在我看来，解送一个囚犯而不指明他所犯的罪案，是不合理的。} \BibleRef{宗 25:24} 注意他怎样控告他们，同时宣告他无罪。哦，多么丰富的辩护！经过这么多次的审查，总督仍找不到定他罪的办法。他们说他是该死的。为此他也说：\Quote{我查明，}他说，\Quote{他没有犯什么该死的罪。……我没有什么确切的事可向我主奏报。}这也证明保禄的无瑕，因为法官关于他竟无话可说。\Quote{因此我带他出来，}他说，\Quote{到你们面前，因为在我看来，解送一个囚犯而不指明他所犯的罪案，是不合理的。}犹太人使自己和他们首领陷入多大的困境！那么怎样？\Quote{阿格里帕对保禄说：准你为自己辩护。} \BibleRef{宗 26:1} 国王出于极想听他的渴望，准他说话。但保禄立刻放胆直言，不是谄媚，而是因此说他是有福的，即因为（阿格里帕）知道一切。\Quote{保禄便伸起手来，为自己辩护说：阿格里帕王，我今天在你面前，为犹太人所控告我的一切事辩护，我认为自己是有福的，尤其因为你晓得犹太人的一切问题；因此我求你耐心听我。}（26:2-3）然而，如果他自觉有罪，他应该害怕在知道一切实情的人面前受审；但这是良心清白的一个标记，不退缩于对情况有准确了解的法官，反而欢喜，并自称有福。\Quote{我求你，}他说，\Quote{耐心听我。}因为他将要延长讲论，并谈一些关于自己的事，为此他先提出恳求，然后说：\Quote{我自幼年起始，在耶路撒冷本族中最初的生活，凡犹太人都知道；他们从起初就知道我，如果他们愿意作证，我是按我们教中最严格的教派过了法利塞人的生活。}（26:4-5）那么，我年轻时既然被众人这样作证，怎么会成为煽动者？其次从他的教派：\Quote{按我们教中最严格的教派，}他说，\Quote{我生活过。}\Quote{那么，如果教派确实值得钦佩，而你是恶的，又如何？}关于这点我也请众人作证——关于我的生活和行为。\Quote{如今我站着受审，是因为对天主向我们祖辈所应许的盼望；我们十二支派，昼夜切切地事奉天主，盼望达到这应许。阿格里帕王，为了这盼望，我被犹太人控告。为什么你们中间竟认为天主使死人复活是不可思议的呢？} \BibleRef{宗 26:6-8} 他提出了关于复活的两种论证：一种是从先知来的论证，但他没有特别举出任何先知，而是举出犹太人自己持有的教义本身；另一种更强有力的论证是从事实来的——（特别是基督亲自与他交谈过）。他用（其他）论证为此作准备，详细叙述他先前的狂乱。然后，他以高度称赞犹太人的口吻说：\Quote{昼夜，}他说，\Quote{事奉（天主）企盼达到。}所以，即使我并非无可指责的生活，我也不该为了这个（教义）受审：——\Quote{阿格里帕王，为了这盼望，我被犹太人控告。}然后又一个论证：\Quote{为什么你们中间竟认为天主使死人复活是不可思议的呢？}因为，如果这样的意见不存在，如果他们不是在这些教条中长大，而是现在初次被引入，也许有人不会接受这话。然后他叙述自己怎样迫害：这也帮助了论证；他请司祭长和\Quote{外邦城市}作证，并且他听见主对他说：\Quote{用脚踢刺棍是难的，}并显示天主的仁慈，虽然受迫害，祂却显现（给人），不仅对我行了这恩惠，也派遣我作老师给其他人；而且显示当时他所听到的预言现在已经应验：\Quote{我救你脱离百姓和外邦人的手，我现在派遣你到他们那里去。}他显明这一切，说：\Quote{我本人曾以为应当多方反对纳匝肋人耶稣的名。这事我就在耶路撒冷作了；从司祭长得了权柄，把许多圣人关在监里；他们被杀时，我也投票反对他们。在各会堂里，我多次用刑强逼他们说亵渎的话；我因对他们极其疯狂，甚至追逼到外邦的城市。为此我带着权柄和司祭长的使命往大马士革去时，王啊，约在中午，我在路上看见从天上来的光，比太阳还亮，照射着我并与我同行的人。我们都仆倒在地，我便听见有声音用希伯来话向我说：扫禄，扫禄，你为什么迫害我？用脚踢刺棍是难的！我说：主啊，你是谁？主说：我就是你所迫害的耶稣。但是请你起来，站定；因为我显现给你，为要立你作仆人和见证人，见证你所看见的事，以及我将要显现给你的事。我要救你脱离百姓和外邦人的手，我现在派遣你到他们那里去，开他们的眼，使他们从黑暗转向光明，从撒殚的权下转向天主，好使他们因信德获得罪过的赦免。} \BibleRef{宗 26:9-18} ——注意他多么温和地讲论——天主，他说，对我说了这话：\Quote{好使他们因信德获得罪过的赦免，并在由信德而圣化的人中得基业。}（26:18 部分）他说，我被这些事说服了，这异象吸引我归向祂，以致我不再延迟。\Quote{阿格里帕王，我没有违反天上的异象，而是首先向大马士革的人，然后向耶路撒冷和犹太全地，再向外邦人宣讲，叫他们悔改归向天主，并作与悔改相称的事。}（26:19-20）那么，我既教导别人关于最卓越的生活方式，我自己怎么会成了煽动和争乱的祸首？\Quote{为此犹太人在圣殿里拿住我，想要杀我。然而我得了天主的帮助，直到今日还站立，向卑贱和尊贵的人作证，所说的不外是先知们和梅瑟所说未来必要发生的事。}（26:21-22）请看他的话何等不谄媚，又如何将一切归于天主。然后他的勇敢——但我现在也不停止；以及确凿的根据——因为我从先知引出问题：\Quote{基督是否应受苦难；}然后复活和应许：\Quote{祂首先从死人中复活，要将光明传给百姓和外邦人。} \BibleRef{宗 26:23} 斐斯托看见这勇敢，说了什么？因为保禄始终向国王说话——他有点恼怒，向他说：\Quote{保禄，你疯了！}因为，\Quote{当他这样辩护时，斐斯托大声说：保禄，你疯了！许多学问使你疯狂了。} \BibleRef{宗 26:24} 然后保禄说什么？他温和地说：\Quote{斐斯托大人，我不是疯狂，而是讲论真理和清醒的话。} \BibleRef{宗 26:25} 然后他也暗示为何他离开他而向国王说话：\Quote{因为王知道这些事，我也向他坦白说话；我深信这些事没有一件能瞒过他，因为这不是在角落里作的。} \BibleRef{宗 26:26} 他显示（国王）完全知道这一切；同时，几乎等于对犹太人说：你们原应当知道这些事——因为他补充说：\Quote{因为这不是在角落里作的。} \BibleRef{宗 26:26} \Quote{阿格里帕对保禄说：\OriginalTerm{少时}{\GreekText{᾿} \GreekText{Εν} \GreekText{ὀλίγῳ}} 你劝我作基督徒。} \BibleRef{宗 26:27} \OriginalTerm{在少时}{\GreekText{ἐ} \GreekText{ν} \GreekText{ὀλίγῳ}} 是什么？就是 \GreekText{\GreekText{παρὰ} \GreekText{μικρόν}}（差一点）。\Quote{保禄说：我祈求天主，\OriginalTerm{无论在少时或在多时}{\GreekText{καὶ} \GreekText{ἐν} \GreekText{ὀλίγῳ} \GreekText{καὶ} \GreekText{ἐν} \GreekText{πολλᾥ}}（无论在少时或在多时），} 对我说来，不是\Quote{在少时}（而是\Quote{在多时}）；他不是简单地祈求，他祈求（不简短，而是）宽宏地——\Quote{不但你，而且今天听我的所有人，都像我一样。}然后他加上：\Quote{除了这些锁链以外；}然而这是光荣的事；不错，但看他们对这事的看法，因此他说：\Quote{除了这些锁链以外。} \BibleRef{宗 26:27}\par

（摘要。）\Quote{次日}等。\BibleRef{宗 25:23}自保禄行使上诉权以来，犹太人便放弃纠缠。然后，为他预备的剧场也颇为壮观：\Quote{大有排场}地出场。\Quote{斐斯托斯说}等。全体犹太人——不是某些人，而其他人不如此——\Quote{在耶路撒冷，也在这里}都说：\Quote{他不该再活着。}\BibleRef{宗 25:24}\Quote{我虽查明}等。这显示他上诉于凯撒是正当的，因为即使他们没有重要事由控告他，但那边的人却疯狂敌视他，他满有理由去见凯撒。\Quote{经你们审断以后，}他说，\Quote{好叫我得些材料上书。}请注意此事如何屡经考验。因此，犹太人应为为保禄的辩护而感谢自己，这辩护也传到罗马的人耳中。看他们如何成为自己邪恶与保禄德性的不情愿的传讯者，甚至传到皇帝本人：因此保禄被带走时，比不受捆绑而去更享荣名；因为他不是作为骗子或欺骗者，在这么多法官宣告他无罪之后，才被带往那里。因此，在与他一同出生的人中，他摆脱了一切指控，不仅于此，还以脱离一切嫌疑的姿态出现在罗马。\Quote{保禄说}等。\BibleRef{宗 26:1}他没有说：\Quote{何故如此？我已一次上诉凯撒；我已受审多次；何时才了结？}但他做了什么？他又准备答辩，并且是在最通晓此事的人面前；充满胆量，既然他们不是定他罪行的法官；但尽管如此，既然那道谕旨依然有效：\Quote{你必须到凯撒那里去}，他便答辩并给予完整回答，\Quote{关于一切事}，而不只是零碎地提及一两件。他们控告我煽动叛乱，控告我传异端，控告我亵渎圣殿：\Quote{关于这一切我为自己答辩}：这些事与我的行径不符，控告者自己就是见证：\Quote{我自幼的生活}等。\BibleRef{宗 26:4}这与先前他说的\Quote{一个热忱者}\BibleRef{宗 22:3}相同。并且当全体百姓在场时，他便挑战他们的见证：不是在法庭前，而是在里息雅前，又在这里，当更多人时；而在那次审理中，他不需要太多自辩，因为里息雅的书信已为他开脱。\Quote{所有的犹太人都知道，}他说，\Quote{那些从起初就认识我的。}他没有说自己的生活怎样，而是交托于他们自己的良心，并将全部重点放在他的派别上，因为他若是一个邪恶品行不良（\OriginalTerm{邪恶和坏}{\GreekText{πονηρὸς} \GreekText{καὶ} \GreekText{μοχθηρός}}）的人，就不会选择那个派别。\Quote{但为了这希望}（手抄本与版本作\OriginalTerm{派别}{\GreekText{αἱρέσεως}}）他说，\Quote{我站立并受审。}（6、7节）这希望在他们之中也受尊重，因为他们为此祈祷，为此朝拜，为要达到这希望：我所显示的正是同一件事。那么，为达此目的而行一切事，却迫害相信同一事的人，岂不是疯狂？\Quote{我本人曾以为，}即我决定，\Quote{应做许多事反对纳匝肋人耶稣的名。}\BibleRef{宗 26:9}我不是基督的门徒：我是在反对他的人的阵营中。因此，他是可信的证人，因为他这个曾做无数事、攻打信徒、劝他们亵渎、煽动一切人（城市、统治者）反对他们，并自动这样做的人，竟突然被改变。然后又是见证人，即与他同行的；接着他显示他有什么正当理由被说服：既因那光，又因先知，又因结果，又因已发生的事。请看，他如何既从先知、又从这些细节，向他们确证这证据。为免他显得在宣扬新奇之事，尽管他有伟大之事可说，他仍再次托庇于先知，并将此作为议题提出。这因实际应验而更可信；但既然只有他看见了（基督），他便再次从先知那里取得证据。请注意他在法庭和集会中如何说话不同：在那里他说：\Quote{你们杀了他}；但在这里并不如此，免得更激起他们的怒气：但他说\Quote{默西亚必须受难}来表明同一事。他如此为他们开脱：因为先知们说这事。因此你们也该接受其余的部分。既然他提到了异象，便无所畏惧地继续说及由此而来的善果。\Quote{使他们从黑暗中转向光明，从撒殚的权下归向天主。我向你显现，正是为此}\BibleRef{宗 26:16}，并非为惩罚，而是为立你为宗徒。他显示不信者所得之害：\Quote{撒殚，黑暗}；信者所得之福：光明、天主、\Quote{众圣徒的基业。因此，阿格黎帕王啊，}等。（19、20节）他不仅劝告他们悔改，还劝他们活出值得钦佩的生命。请看，外邦人处处被接纳而与百姓（以色列）相连：因为当时在场的人中有外邦人。\Quote{我见证，}他说，\Quote{无论对大小}\BibleRef{宗 26:22}，即对尊贵与卑微。这也是对士兵说的。请注意：他离开被告的位置，而担任教师的角色——因此斐斯托斯对他说：\Quote{你疯了}——但为了不显得自己是教师，他引入先知和梅瑟：\Quote{默西亚是否必须受难，他是否首先从死者中复活，向百姓和外邦人宣报光明。}\BibleRef{宗 26:23}\Quote{斐斯托斯便大声说}——怀着愤怒与不快——\Quote{保禄，你疯了。}保禄怎么说？\Quote{我并没有疯}等。\Quote{这事，}他说，\Quote{不是在角落里做的。}（25、26节）这里他谈到十字架、复活：那教义已传遍世界每一部分。\Quote{阿格黎帕王，}他说，\Quote{你信吗}——他不是说复活，而是——\Quote{先知吗？}\BibleRef{宗 26:27}然后他抢先说：\Quote{我知道你信。}\OriginalTerm{以少许}{\GreekText{᾿} \GreekText{Εν} \GreekText{ὀλίγῳ}}（即几乎）\Quote{你几乎说服我作一个基督徒。}\BibleRef{宗 26:28}保禄不懂短语\GreekText{\GreekText{ἐ} \GreekText{ν} \GreekText{ὀλίγῳ}}的意思；他以为它意为\OriginalTerm{由于少许}{\GreekText{ἐ} \GreekText{ξ} \GreekText{ὀλιγου}}（即花很少代价或麻烦），因此他也如此回答；他是这样不学无文。他没有说：\Quote{我不希望}，而是说：\Quote{我祈祷，不仅是你，还有所有听见的人。}请注意他的话何等无谄媚。——\Quote{我祈望今天他们都成为像我一样，除了这些锁链。}\BibleRef{宗 26:29}他，以锁链为荣，视之为金链的人，却为这些人祈求免除锁链：因为他们心志尚软弱，他这样说乃是迁就。因为有什么比那些锁链更好？这些锁链他在书信中总视为贵重，说：\Quote{保禄，为基督耶稣的囚犯}\BibleRef{弗 3:1}；又说：\Quote{我为这缘故被这锁链捆锁}\BibleRef{宗 28:20}，\Quote{但天主的道却不被捆锁}；以及\Quote{甚至到锁链，如同作恶的。}\BibleRef{弟后 2:9}惩罚是双重的。因为若他真的如此被绑，且是为他的好处，这事或许带着某种安慰；但现在他被绑，既\Quote{如同作恶的}，又面对极坏后果；然而他对这一切全不在意。\par

这便是为天上之爱所振翅的灵魂。因为若那些珍惜污秽的（世人所称之）情爱者，视无物为荣耀或贵重，唯独那些能遂其私欲的事物，他们视之为荣耀尊贵，且其情妇于他们即是一切；那么，那些为这天上之爱所俘获的人，岂不更将\OriginalTerm{代价}{\GreekText{τὰ} \GreekText{ἐπιτίμια}}视为无物。但若你们不明白我所说的话，不足为奇，因我们对此神圣智慧尚不熟练。因为若有人为基督之爱的火焰所俘虏，他将变成如同独居地上之人，对荣耀或耻辱全然不在意；正如若他独居，便一无所顾，此处亦然。至于试炼，他如此蔑视它们，无论鞭打或监禁，如同那承受这些苦楚的身体是别人的而非自己的，或如同他拥有钻石般的身体；而对于今生的甜美事物，他如此嗤笑，如此无感，如同我们对尸体无感，因我们自己已死。他远离被任何情欲所俘虏，如同经火炼净的纯金毫无杂质。因为正如苍蝇不会冲入火焰，反而飞离，同样，情欲不敢靠近这人。但愿我能从我们中间举出这一切的实例：但既然我们找不到，便不得不求助于这位保禄。那么请观察他如何对待全世界。他说：\BibleQuote{世界于我已被钉在十字架上，} \BibleQuote{我于世界亦然} \BibleRef{迦 6:14}：我于世界已死，世界于我亦死。又说：\BibleQuote{不再是我生活，而是基督在我内生活} \BibleRef{迦 2:20}。为使你们看见他如同独处，如此看待现世，请听他亲口说：\BibleQuote{我们并不注目那看得见的，而是那看不见的} \BibleRef{格后 4:18}。你说什么？回答我。然而你所说的恰恰相反；你看见不可见之物，却看不见可见之物。你所获得的这样的眼睛，正是基督所赐的眼睛：因为正如这些肉眼确实看见可见之物，却看不见不可见之物；同样，那些（天上的眼睛）则相反：没有人既看见不可见之物，又看见可见之物；没有人看见可见之物，又看见不可见之物。难道我们不也是如此吗？因为当我们收敛心神，思想任何不可见的事物时，我们的眼界便超越地上之物。让我们轻视荣耀：情愿被嘲笑，也不愿被赞扬。因为被嘲笑的人毫无损伤，而被赞扬的人却大受损害。让我们不要看重那些使人恐惧的事物，而要如同我们对待孩童那样做：即，若我们看见有人恐吓孩童，我们并不钦佩那人；因为事实上，凡恐吓者，只恐吓孩童；若是对成人，他不能恐吓他。正如那些（游戏中）恐吓孩童的人，要么翻起眼皮，要么扭曲面容，但以自然和蔼的眼神，他们便不能做到；同样，这些人通过扭曲他们的\OriginalTerm{理智的洞察力}{\GreekText{τὸ} \GreekText{διορατικὸν} \GreekText{τἥς} \GreekText{διανοίας}}来做到。因此，对温和且心灵美丽的人，无人会害怕；相反，我们所有人都尊重他，尊敬他，崇敬他。你们岂不见，那引起恐惧的人也为我们众人所憎恨厌恶吗？对于仅能引起恐惧的事物，我们还有什么不回避的呢？岂非如此对待野兽、声响、景象、地方、空气如黑暗？因此，若人畏惧我们，我们不要以为这是大事。因为首先，无人真正因我们而恐惧；其次，（即使有）也不是什么大事。德性是一大善：且看何其大。无论我们认为构成德性的手段多么可悲，我们仍钦佩德性本身，并以为拥有它的人是有福的。因为谁不认为那忍耐受苦的人是有福的呢，即使贫穷及类似事物看似可悲？因此，当德性透过那些看似可悲的事物发出光芒时，请看这是何等超卓！人啊，你因为拥有权力就自大吗？这是怎样的权力？请说，它是因任命而授予的吗？（若是如此），你从人领受了权力：你应从内心任命自己。因为统治者并非那被称为统治者的人，而是那真正是统治者的人。因为正如国王不能使一个人成为医生或演说家，同样也不能使一个人成为统治者：因为不是（皇帝的）诏书或名分使人成为统治者。因为，若你愿意，让某人开设一间药铺，让他也有学徒，让他也有器械和药品，并让他探访病人：这些足够使他成为医生吗？绝非如此：还需要技艺，没有它，这些事物不仅无益，反而有害：因为若不是医生，最好连药品也不拥有。不拥有药品的人，既不救人也害不了人；但拥有药品的人，若不知如何使用，便害人：因为治愈之力不仅在于药品的性质，还在于施用者的技艺；若没有技艺，一切皆毁。统治者也是如此：他拥有的工具是声音、愤怒、刽子手、流放、荣誉、礼物和赞扬；他也拥有法律作为药品；他也拥有人民作为病人；他实践技艺的场所是法庭；他拥有士兵作为学徒：因此，若他不懂治愈之道，这一切对他毫无益处。法官是灵魂的医生，不是身体的医生：但若肉身的医术尚需如此谨慎，那么灵魂的医术更需谨慎，因为灵魂比身体更重要。因此，并非仅仅拥有统治者的名分就是统治者：因为其他人也被称为伟大的名字：如保禄、伯多禄、雅各伯、若望：但名字并不使他们成为所称呼的那样，正如我的名字不能使我成为（若望那样）；我确实与那位有福者同名（\OriginalTerm{同名者}{\GreekText{ὁ} \GreekText{μώνυμος}}），但并非同义（\OriginalTerm{并非同义词}{\GreekText{οὐ} \GreekText{μὴν} \GreekText{συνώνυμος}}），我不是若望，只是被如此称呼。同样，他们不是统治者，只是被称为统治者。但那些没有这些附随物的人却是统治者，正如一位医生，即使他实际上不行医，但若他心中有医术，他就是医生。那些能管理自己的人，才是统治者。因为有四样事物：灵魂、家庭、城市、世界；它们形成一种\OriginalTerm{逐步进展}{\GreekText{ὁ} \GreekText{δᾥ} \GreekText{προβαίνει}}。因此，那要监督一个家庭并妥善管理它的人，必须先使自己的灵魂有序；因为他的家庭就是灵魂：但若他不能管理自己的家庭，那里只有一个灵魂，他自己是主人，他常与自己同在，又如何管理他人？那能够管理自己灵魂，使一部分统治、另一部分服从的人，他也能管理家庭；而能管理家庭的人，也能管理城市；若能管理城市，也能管理世界。但若他不能管理自己的灵魂，又如何能管理世界？我说这些话，是为使我们不为统治职位所激动；为使我们知道什么是统治：因为（所谓的）统治并非统治，而是嘲笑的对象、纯粹的奴役，以及许多其他名称。请告诉我，统治者的特质是什么？岂非帮助臣民并使他们受益？那么若不然，又当如何？那尚未帮助自己的人，如何帮助他人？那在自己灵魂中拥有无数情欲暴政的人，如何根除他人的情欲？再者，关于所谓的\Quote{奢华}（\Quote{奢华}）或享乐生活：真正的奢华或享乐并非（所谓的）那样，而是另一回事。因为正如我们已证明，统治者不是被称为统治者的人，而是另有其人（拥有超越名分的实质），同样，真正享乐的人也是另一种人。因为\Quote{奢华}或享乐生活似乎确实是享受快乐和满足肚腹：然而并非如此，而是相反：它在于拥有一个值得称赞的灵魂，并处于快乐状态。因为让一个人吃喝纵欲，然后让他遭受忧虑和沮丧：这个人能说是处于快乐状态吗？因此，不是吃喝，而是处于快乐之中，才构成真正的奢华或享乐生活。让一个人只吃干面包，却内心充满喜乐：这不是快乐吗？那么，这便是真正的奢华。让我们看看，这幸福归于谁——是富人，还是非富人？不全归于这一边，也不全归于那一边，而是归于那些妥善管理自己灵魂，以致没有许多悲伤缘由的人。这样的生活在哪里能找到呢？因为我看见你们全都急切渴望听到这无悲伤的生活是什么。那么，首先要你们承认：这就是快乐，这就是真正的奢华——没有烦恼引起不安；不要向我索求肉食、酒、酱料、丝袍和丰盛的筵席。但若我证明，除这一切之外，有这样的生活近在咫尺，那么你就接受这快乐和这生活吧：因为大部分痛苦之事源于我们未能恰当计算。那么，谁会有最多的悲伤——是那毫不关心这些事物的人，还是那关心它们的人？是那惧怕变化的人，还是那不怕的人？是那恐惧嫉妒、猜忌、诬告、阴谋、毁灭的人，还是那远离这些恐惧的人？是那需要许多事物的人，还是那无所需求的人？是那作无数主人奴隶的人，还是那不作任何人的奴隶的人？是那需要许多事物的人，还是那自由的人？是那惧怕一个主人的人，还是那惧怕无数暴君的人？那么，这里的快乐更大。因此，让我们追求这快乐，不要为现世之事激动：让我们嘲笑生活中的一切浮华，并处处实践节制，使我们能以无痛苦的方式度过此生，并得到所许诺的福分，赖我们的主耶稣基督的恩宠和慈爱，偕同他，圣父及圣神，一同接受光荣、权能、尊崇，从今直到永远，世世无尽。阿们。\par

\SourceRecord{Translated by J. Walker, J. Sheppard and H. Browne, and revised by George B. Stevens.；Nicene and Post-Nicene Fathers, First Series；Vol. 11.；Edited by Philip Schaff.；Buffalo, NY: Christian Literature Publishing Co.,；1889.；<http://www.newadvent.org/fathers/210152.htm>.}

% structure: p0001 source=h1 target=section reason=page-title

\section{\WorkTitle{宗徒大事录} 讲道 第五十三篇}

\BibleRef{宗 26:30-32}\par

\BibleQuote{当他说完这话，王和总督，以及\Person{贝勒尼切}，并与他们同坐的人都起来；退到一边，彼此谈论说，这个人并没有犯什么该死或该绑的罪。\Person{阿格黎帕}对\Person{斐斯托斯}说，这人若没有向\Person{凯撒}上诉，早已可以释放了。}\par

请看，他们又是如何再次宣告他无罪的；在说了\Quote{你疯了}\BibleRef{宗 26:24}之后，他们开释了他，认为他不仅不该死，甚至不该受捆绑。事实上，要不是他向凯撒上诉，他们本会完全释放他。但这事出于上智的安排，好让他带着锁链离去。他说：\Quote{至于锁链，如同恶人一般。}\BibleRef{弟前 2:9}因为如果他的主\Quote{被列于罪犯之中}\BibleRef{谷 15:28}，何况他呢？但正如主不与他们同流合污，保禄也是如此。因为奇妙之处正在于此：与这样的人混杂在一起，却不从他们受到任何伤害。\BibleQuote{既议定要我们坐船往意大利去，便将保禄和别的囚犯，交给御营里一个名叫犹里约的百夫长。上了一只阿得辣米雅船，要开船沿着亚细亚一带地方前行；有马其顿的得撒洛尼人阿黎斯塔苛同我们在一起。第二天，我们到了漆冬。}\BibleRef{宗 27:1}请看阿黎斯塔苛陪伴保禄多么远。阿黎斯塔苛在场有良好而有益的目的，因为他会将一切报告带回马其顿。\Quote{犹里约宽待保禄，准他往朋友那里去，获得照应。}请看作者也没有隐瞒这一点，即保禄希望\Quote{获得照应}。\BibleQuote{从那里又开了船，因为没有顺风，就沿着塞浦路斯背风而行。}\BibleRef{宗 27:4}又是考验，又是逆风。请看圣人们的生活就是这样交织着：逃出法庭，又遭遇船难和风暴。\BibleQuote{我们沿着基里基和旁非利亚海行，到了里基雅的米辣。百夫长在那里找到一只亚历山大里亚的船，要开往意大利去，便叫我们上了那船。}\BibleRef{宗 27:5-6}经上说：\Quote{一只亚历山大里亚的船。}很可能，前一只船上的人会把保禄的遭遇传到亚细亚，而这些（亚历山大里亚船的人）也会在里基雅传扬同样的事。请看天主并没有改变或扰乱自然的秩序，而是让他们迎着不顺的风航行。但即便如此，奇迹仍发生了。为了让他们安全航行，他没有让他们进入（开阔的）大海，而是总在靠近陆地航行。\BibleQuote{一连多日，船行得很慢，好不容易才到了克尼笃对面，因为风不顺，我们就沿着克里特背风而行，从撒耳摩讷对面经过；好不容易才到了一个地方，名叫佳港，附近有拉撒雅城。过了很多日子，行船已经危险，因为禁食节已经过了，保禄便劝告他们。}\BibleRef{宗 27:7-9}这里所说的\Quote{禁食节}，我想是指犹太人的禁食节。因为他们在五旬节之后很久才离开那里，所以到克里特海岸时大约已是仲冬。这也是不小的奇迹，即他们因保禄的缘故也得救。\BibleQuote{保禄劝告他们说：诸位同人，我看这次航行，不但货物和船要受灾，连我们的性命也要大受损失。但百夫长更信从船长和船主，不听保禄所说的话。又因为这港口不便过冬，多数人主张离开那里，也许能到腓尼司去过冬；腓尼司是克里特的一个港口，朝向西南和西北。那时，南风缓缓吹来，他们以为目的达到了，便起锚沿着克里特航行。但不久，一股名为\Quote{欧拉基罗}的飓风从岛上扑来。船被风卷住，不能与风抗争，我们只好任船漂流。}\BibleRef{宗 27:10-15}因此保禄劝告他们留下，并预言了后果。但他们匆忙行事，又因该地不便，希望到腓尼司去过冬。请看事件的天主眷顾安排：起初，\Quote{南风缓缓吹来，他们以为目的达到了}，就解开船缆出发了；然后当飓风袭来时，他们任由船被吹走，好不容易才得救。\BibleQuote{我们到了一个名叫克劳达的小岛旁边，费了大力才收住了小艇。把小艇拉上来后，就用绳索捆绑船底；又怕在叙尔提斯搁浅，便落下船帆任船漂流。我们被风浪颠簸得很厉害，第二天他们开始抛货；第三天他们亲手把船上的器具也扔了。一连多日，不见太阳，也没有星辰，风暴依然凶猛，我们得救的希望终于破灭了。众人多日没有进食，保禄便站在他们中间说：诸位同人，你们本该听我的话，不该从克里特开船，以致招来这场损害和损失。}\BibleRef{宗 27:16-21}在这样大的风暴之后，他并没有以侮辱的口吻对他们说话，而是希望无论如何他们将来会相信他。因此他也引述已经发生的事，作为他将要说的话真实性的证据。\BibleQuote{现在我劝你们放心吧，你们中没有一个会丧命的，只损失这只船。因为今夜，我所归属并事奉的天主的使者站在我身边说：保禄，不要害怕，你必须在凯撒面前受审；并且天主把你同船的人都赐给了你。所以诸位同人，请放心吧！我信天主必会这样成就，正如对我所说的。不过我们必要搁浅在一个岛上。}\BibleRef{宗 27:22-26}他预言了两件事：他们必要搁浅在一个岛上，并且虽然船会沉没，但船上的人都会得救——他说这话不是出于猜测，而是出于预言——以及他\Quote{必须在凯撒面前受审}。他所说的\Quote{天主把你们都赐给了你}，并非出于自夸，而是为了赢得同船的人：他说这话，不是为了让他们觉得自己欠了他，而是为了让他们相信他所说的话。\Quote{天主把你们都赐给了你}，好比说：他们确实该死，因为他们不听你的话；但这事是出于对你的恩宠。\BibleQuote{到了第十四天夜间，我们在亚得里亚海漂流时，约在午夜，水手们以为靠近了陆地；便探测水深，得到二十寻；再往前走一点，又探测，得到十五寻。他们怕撞在石头上，便从船尾抛下四个锚，盼望着天亮。水手们想弃船逃走，便把小艇放到海里，假装要从船头抛锚的样子。保禄对百夫长和士兵们说：这些人若不留在船上，你们就不能得救。于是士兵们砍断小艇的绳索，任它漂去。}\BibleRef{宗 27:27-32}水手们原想逃走，不信所说的话；但百夫长相信保禄。因为保禄说：\Quote{这些人若不留下，你们就不能得救。}这样说，并非为此，而是为了制止他们，使预言不致落空。请看他们如何仿佛在教堂里，被保禄冷静的举止所教导，他如何从危险的中心拯救他们。出于天主的眷顾，保禄一度不被相信，以便经过事实的证明后，他能被相信：结果正是如此。他又劝他们吃点东西，他们照做了；他先吃，为要说服他们，不仅用言语，也用行动表明风暴对他们无害，反而有益于他们的灵魂。\BibleQuote{天快亮时，保禄劝众人都吃点东西，说：你们一直挨饿，不吃什么，已经十四天了。}\BibleRef{宗 27:33}你们说，他们怎么不吃东西，什么也没吃？他们怎么受得了？恐惧占据了他们，不让他们产生对食物的欲望，因为他们处于极度危险之中；但他们无暇顾及食物。\BibleQuote{所以我劝你们吃点东西，这是为你们的益处，因为连你们头上的一根头发也不会失落。说完这话，就拿起饼来，在众人面前感谢了天主，然后掰开开始吃。于是众人都放了心，也吃了。}\BibleRef{宗 27:34}他们看见自己的生命得救已无问题，便都放了心。\BibleQuote{船上一共有二百七十六人。他们吃饱了，便将麦子抛在海里，为减轻船的重量。天亮时，他们认不出那地方，但看见一个有沙滩的海湾，便决定尽可能把船靠岸。于是砍断锚，把它弃在海里，同时松开舵绳，扬起前帆，顺着风向岸边驶去。}\BibleRef{宗 27:37}他们\Quote{向着岸边驶去}，把舵柄交给了风：因为他们通常不是这样做的。他们松开了帆索，即帆，随波逐流。\BibleQuote{却落在两海交汇的地方，船搁了浅，船头胶住不动，船尾被波浪的猛力冲坏了。}\BibleRef{宗 27:41}因为当有强风时，这就是后果：船尾承受了（风暴的）主要冲击。\BibleQuote{士兵们提议要把囚犯杀死，免得有人游泳逃走。}\BibleRef{宗 27:42}魔鬼又想阻碍预言，他们本打算杀死一些人，但百夫长为要救保禄，不准他们这样做：百夫长对保禄如此依恋。\BibleQuote{但百夫长愿意救保禄，便阻止他们这样做，命令会游泳的人先跳下海，上岸去；其余的人有的用木板，有的用船上的碎木。这样，众人就都安全上了岸。}\BibleRef{宗 27:43-44}\BibleQuote{他们上岸后，才知道那岛名叫默里达。}\BibleRef{宗 28:1}你看出风暴的好处了吗？那么，风暴临到他们，并非他们被舍弃的标志。所发生的事是由于季节的缘故；但更奇妙的是，在这样的季节，他们竟从危险中得救：他得救，并因他的缘故，其余的人也得了救，而且是在亚得里亚海。总共有二百七十六人：这也不是小事，如果他们真的信了。航行是在不合时宜的季节。可以推测，他们会问航行原因，并知道一切。航行如此漫长也不是徒然的；它给保禄提供了施教的机会。\par

（重述。）保禄说：\BibleQuote{我看出这次航行将有损伤和亏损。}\BibleRef{宗 27:10} 请注意这话是何等谦虚。为避免他看似在预言，而只是猜测之辞，他说：\BibleQuote{我看出}。因为若一开始就说这话，他们不会接受。事实上，他在此前的场合确实预言了，后来也一样，并说（那里）：\BibleQuote{我所事奉的天主}，引导他们。那么，为何不是\Quote{损伤}他们的\Quote{性命}呢？本该如此，但天主使他们安然度过。就事情本身而言，他们本已丧亡，但天主阻止了。为表明他并非出于猜测而这样说，船主——一个在此事上有经验的人——却说了相反的话 \BibleRef{宗 27:11}，所以保禄的建议绝不是出于猜测。此外，那地方（船主所说的）也暗示了这点，\Quote{不便停泊}；显然大多数人出于猜测而建议 \BibleRef{宗 27:12} 如他们所行，而非保禄。接着，风暴猛烈，黑暗深沉；为使他们不忘却，船只也破碎，谷物被抛出，一切尽失，使他们此后无法再厚颜无耻。因此船只破碎，他们的灵魂被紧紧绷住。而且，风暴与黑暗对他获得听众的信任贡献不小。由此可见，百夫长如何遵从他的吩咐，甚至放手让小船漂走并毁掉它。若水手们尚未听从他的命令，但后来还是听从了：事实上，这是一群鲁莽的人。\BibleRef{宗 27:13} \BibleQuote{诸位，你们本该听从我，}等等。\BibleRef{宗 27:21} 在灾难中责备人，通常不会有好效果；但当他告诉他们灾难中还有什么（临到），然后预言好事，那时他才被接纳。因此他先攻击他们，当时\Quote{得救的一切希望都已断绝}：免得有人说，什么事也没发生。他们的恐惧也作证。而且，这地方是个考验，因为是在亚得里亚海，再加上他们长期绝食。他们正处于死亡之中。那时已是第十四天，他们未进食，什么也没有吃。所以他说：\BibleQuote{我劝你们用饭，因为这与你们有益}\BibleRef{宗 27:34}，你们该吃，免得饿死。请注意，他在这一切之后谢恩，坚固了他们。因为这表明他确信必得救。\Emph{(b)} \BibleQuote{于是众人都放心，也用了饭。}\BibleRef{宗 27:36} 不仅如此，此后他们如此将一切挂虑都交给保禄，甚至把谷物也抛了出去 \BibleRef{宗 27:37}，人数如此之多。\Emph{(a)} 二百七十六人 \BibleRef{宗 27:38}：他们从哪里得到食物？\Emph{(a)} 看他们如何尽人事，以及保禄如何没有禁止他们。\Emph{(c)} \BibleQuote{到了天亮，}等等，\BibleQuote{他们松开了舵绳。}\BibleRef{宗 27:39-40} 船在白天破碎，使他们不至于因恐惧而完全崩溃；好叫你们看见预言成为事实。\BibleQuote{兵士的计谋，}等等。\BibleRef{宗 27:42} 你们注意到，在这方面他们也归了保禄吗？因为为了他，百夫长没有让他们被杀。我看那些人真是罪恶昭彰，甚至宁愿杀死他们：但有的游上岸，有的抱着木板，全都得救，预言得以应验；（这预言）虽非因时间长久而庄重——因为他并非早在多年以前就说出，而是紧贴事情本身的性质——（仍然是预言），因为一切都超乎希望。他们是通过自己被救而认识保禄是谁。但有人会说：他为什么没有保全船？为使他们明白自己逃脱了多大的危险；并且整件事不靠人的帮助，而完全依靠天主的手，无需船只而拯救他们。这样看来，义人虽处风暴、海上、深渊中，也不受任何可怕之事，反而连他人一同拯救。若有一艘船遇险沉没，囚犯因保禄的缘故得救，那么家中有一位圣人是何等大事！因为也有许多风暴袭击我们，比这些（自然）风暴更可怕，但他也能使我们得救，只要我们像船上那些人一样服从圣人，遵行他们的吩咐。因为他们不仅得救，自己也促使别人相信（\OriginalTerm{带来信心}{\GreekText{πίστιν} \GreekText{εἰσήνεγκαν}}）。圣人虽受捆绑，却行出比自由人更大的事。看这里的情况正是如此。自由的百夫长需要被捆绑的囚犯；熟练的舵手需要那位非舵手——倒不如说，需要那位真正的舵手。因为他掌舵的不是这种（地上）的船，而是全世界的教会，他从也是海的主的那一位学得；（他掌舵）不靠人的技艺，而靠圣神的智慧。在这船上有许多船难、许多波浪、邪恶的势力，\BibleQuote{内有争斗，外有恐惧}\BibleRef{格后 7:5}，所以他才是真正的舵手。看我们整个生命，正像这次航行。因为我们时而遇到和善，时而遇到风暴；有时因缺乏谋略，有时因懒惰，陷于无数灾祸；由于我们不听从保禄，一心要去他不许我们去的地方。因为保禄如今也与我们同航，只是不像当时那样被捆绑；他现在也劝诫我们，对那些在这海上航行的人说：\BibleQuote{你们要谨慎，因为我离开之后，将有凶暴的豺狼进入你们中间}\BibleRef{宗 20:29}；又说：\BibleQuote{末世将有危险时期，人将专爱自己、贪爱钱财、自夸。}\BibleRef{弟后 3:2} 这比一切风暴更可怕。因此，让我们留在他吩咐的地方——在信仰中，在安全的港湾；让我们听从他的话，而不是听从我们内的舵手，即我们自己的理性。不要立刻按这理性所建议的去做；也不要按船主的建议，而是按保禄的建议：他经历过许多这样的风暴。让我们不要通过经验学习（付出代价），而是在经验之前\Quote{避免伤害和亏损}。听他说：\BibleQuote{那些想发财的人，陷于诱惑。}\BibleRef{弟前 6:9} 因此，让我们服从他；否则，看看他们因不听从他的劝告而受了什么苦。他在另一处又讲明什么造成船难：\BibleQuote{有人}他说，\BibleQuote{在信仰上翻了船。但你要持守你所学的和所确信的。}\BibleRef{弟前 1:19} 让我们服从保禄：即使我们在风暴中，也必脱离危险；即使我们十四天没有食物，即使得救的希望已离开我们，即使我们在黑暗和迷雾中，遵行他的命令，就必脱离危险。让我们把整个世界看作一艘船，在其中，恶人和有无数恶习的人，有些是统治者，有些是看守，有些是义人如保禄，有些是囚犯，即被罪恶捆绑的人；如果我们按保禄所吩咐的去做，我们就不会在捆绑中灭亡，反而得释放；天主也将我们赐给他。你们不认为罪恶和情欲是沉重的捆绑吗？因为被捆绑的不只是手，而是整个人。告诉我，当一个人拥有许多钱财，却不用不花，紧锁不放，他岂不被吝啬捆绑得比任何囚犯更紧？那不能折断的锁链！再者，当一个人屈服于命运（的信仰），他岂不是也被别的锁链捆绑？屈服于（时辰的）观察呢？屈服于征兆呢？这些岂不比一切捆绑更痛苦？再者，当一个人屈服于不理智的情欲和爱时，谁能为你折断这些锁链？需要天主的帮助才能解开。但既有捆绑又有风暴，试想危险有多大！其中哪一个不足以毁灭？饥饿、风暴、船上人的邪恶、季节的不适宜？然而，保禄的荣耀在这一切面前屹立不倒。现在也是如此：让我们将圣人留在我们身边，就不会有风暴；或者，即使有风暴，也会有极大的平静安宁，脱离危险。因为那位寡妇有了圣人为友，她孩子的死亡就被解除了，她得回了活的儿子。\BibleRef{列上 17:17} 圣人足迹所踏之处，必无痛苦之事；若发生这样的事，也是为了考验我们，并彰显天主更大的光荣。让你的房屋地面习惯于被这样的脚踩踏，邪灵就不会在那里行走。因为正如哪里有香气，哪里就没有恶臭；同样，哪里有圣油，哪里邪灵就被窒息，它使附近的人喜乐、欢欣、灵魂舒畅。哪里有荆棘，哪里就有野兽；哪里有款待，哪里就没有荆棘：因为施舍进入，比任何镰刀更锐利地铲除荆棘，比任何火更猛烈。不要害怕：（恶者）害怕圣人的足迹，如同狐狸害怕狮子。因为经上说：\BibleQuote{义人像狮子一样勇敢。}\BibleRef{箴 28:1} 让我们把这些狮子带进我们的房屋，一切野兽就逃跑；狮子不需要吼叫，只要发声。因为狮子吼叫驱散野兽，远不如义人的祈祷驱散邪灵那样有效：只要他一开口，它们就退缩。你也许会说，如今哪里找得到这样的人？处处都有，只要我们相信、寻找、尽力。告诉我，你曾在何处寻找？你何时着手此事？你何时以此为务？若不寻找，不要奇怪你找不到。因为\BibleQuote{寻找的，就找到}\BibleRef{玛 7:7}，不寻找的则不然。要听那些住在旷野的人的话：抛开你的金银财宝：（这样圣人）在世上各处都能找到。即使你未能在家里接待这样一位，却可到他那里去，与他同住，住在他的居所，这样你就能获得并享有他的祝福。接受圣人的祝福是件大事：让我们小心获得它，藉着他们的祈祷，得到天主的怜悯，透过祂独生子的恩宠和慈爱，愿光荣、权能、尊崇归于祂和圣父及圣神，从今世直到永远，世世无尽。阿们。\par

\SourceRecord{Translated by J. Walker, J. Sheppard and H. Browne, and revised by George B. Stevens.；Nicene and Post-Nicene Fathers, First Series；Vol. 11.；Edited by Philip Schaff.；Buffalo, NY: Christian Literature Publishing Co.,；1889.；<http://www.newadvent.org/fathers/210153.htm>.}

% structure: p0001 source=h1 target=section reason=page-title

\section{宗徒大事录第五十四篇讲道}

\BibleRef{宗 28:1}\par

\BibleQuote{那些土人待我们非常友善，因为他们生了火，又接待了我们每一个人，因为当时有雨，又因为天气寒冷。当保禄拾了一捆树枝，放在火上时，有一条毒蛇因热逸出，咬住了他的手。}\par

\BibleQuote{显示了}，他说，\BibleQuote{对我们无比的慈爱——那些土人}（他们虽是土人）——\BibleQuote{生起了火}：否则即使得救也无用，如果寒冷的天气必定摧毁他们。然后保禄拾了树枝，放在火上。看他多么活跃；我们注意到，他从不无故行奇迹，只在必要时才做。在风暴中，当有原因时，他预言，不是为了显示预言，而在这里，他首先拾取树枝——不为炫耀，只为使他们得以保全，并得些温暖。然后一条毒蛇\BibleQuote{咬住了他的手。那些土人看见毒蛇悬在他手上，就彼此说：这人必定是凶手，虽然从海里逃生了，但正义还不许他活着。}\BibleRef{宗 28:4}这事也得以发生，好让他们既看见这事，又说出这想法，以便结果出现时，无人怀疑这奇迹。看他们（对受难者）的善意，说这话时（不是大声，而是）彼此间——也看即使是土人也明确表达的自然判断，他们如何不无故定罪。这些人也观看，为要更加惊奇。\BibleQuote{他便把毒蛇抖进火里，并未受伤。他们等着看他肿胀，或突然倒毙；但等了许久，见他并无损害，就改变了想法，说他是神。}（五、六节）据说他们预期他会倒毙；但看到他没有发生类似之事，就说他是神。又见\BibleRef{宗 14:11}，这些人的另一个过分之举。\BibleQuote{在那地带，岛长名唤部百流，他的田地；他接待我们，友善地款待了我们三天。恰好部百流的父亲卧病，患热病和痢疾；保禄进去，祈祷后，按手在他身上，医好了他。}（七、八节）看，又一位好客之人部百流，他富有且产业广大：他未见过什么，纯粹出于怜悯他们的不幸，就接待并照料了他们。因此他配得善待：所以保禄作为回报，\BibleQuote{医好了他。这事之后，岛上其他患病的人也来了，都得了医治；他们也用许多尊荣尊敬我们；当我们启程时，他们给我们装了所需之物}（九、十节），无论我们众人还是其余的人。看他们脱离风暴后，并未更加怠惰，而是因保禄的缘故得到何等慷慨的款待：他们在那里住了三个月，都得供养。看这一切都为保禄而行，为使囚犯、士兵和百夫长相信。即使他们冰冷如石，从他们听见他给出的劝告、听见他作出的预言、知道他行的奇迹、以及因他而获得的供养，也必对他有了极高的看法。看，当判断正确，不为私欲所蔽时，它如何立刻得正见，作出公断。\BibleQuote{过了三个月，我们乘一艘在岛上过冬的亚历山大船启航，船以汎斯达和珀鲁克斯为记。在息辣谷撒登陆，我们停了三天。从那里绕行，到了勒基雍；次日南风刮起，我们次日到了颇提约里；在那里找到弟兄们，被邀请与他们同住七天；于是我们向罗马进发。从那里，弟兄们听见我们的事，便来到亚比乌市场和三馆迎接我们；保禄见了他们，感谢天主，勇气倍增。}\BibleRef{宗 28:11}宣讲已传到息辣地方：看它如何遍传（甚至到那些地区）：在颇提约里也找到一些人；还有人来迎接他们。弟兄们的热诚如此，毫不因保禄被捆锁而烦恼。但也看保禄本人如何有人性感受。因为经上说：\BibleQuote{他见了弟兄们，勇气倍增。}虽然行了那么多奇迹，但连亲眼所见也使他增添了（信心）。由此我们得知，他既有人性的安慰，也有相反的感受。\BibleQuote{我们进了罗马，保禄获准独自居住，由一个兵士看守。}\BibleRef{宗 28:16}他获准独居。这也是他备受尊敬的有力证明：显然他们并未将他与其余人同列。\BibleQuote{过了三天，他召集了犹太人的首领们。}过了三天，他召集了犹太人首领，以免他们耳中先有偏见。他与他们有何相干？否则他们很可能控告他。然而，他并不在意这个；他关心的是教导，不让自己所说的话冒犯他们。\par

(摘要。) \BibleQuote{野蛮人}等。\BibleRef{宗 28:2} 那时犹太人看见他们所行的许多奇迹，就迫害和骚扰（保禄）；但那些没有看见任何奇迹的野蛮人，仅仅因他的不幸而善待他。—— \BibleQuote{毫无疑问，}他们说，\BibleQuote{这人是个凶手。}\BibleRef{宗 28:4} 他们并非简单地宣告判决，而是说：\BibleQuote{毫无疑问，}（即）任何人都看得出，\BibleQuote{复仇，}他们说，\BibleQuote{不容他活命。}那么，他们也持守关于天主眷顾的道理，这些野蛮人比那些哲学家更有智慧，因为那些哲学家不承认天主眷顾的恩惠扩展到\InnerQuote{月亮之下}的事物；而这些野蛮人却认为天主无所不在，并且一个（有罪的）人虽能逃脱许多（危险），但最终必定难逃惩罚。他们并不立即攻击他，却因他的不幸而一时间尊敬他：也不公开宣布他们的猜测，而是\InnerQuote{在彼此之间}说：\InnerQuote{一个凶手；}因为锁链使他们起了这样的疑心。\BibleQuote{他们表示了不小的友善，}但（他们中的一些人）也是囚犯。让那些说不要向监狱中的人行善的人羞愧吧：让这些野蛮人使我们羞愧；因为他们不知道这些人是谁，仅仅因为他们在不幸中（就行善）：他们只意识到这些人是人，因此认为他们有人性的权利要求。经上说：\BibleQuote{他们等了很久，以为他会死。}\BibleRef{宗 28:6} 但当他抖动手，把毒蛇甩掉时，他们看见了并感到惊讶。这个奇迹并不是突然发生的，而是经过了一段时间，\InnerQuote{他们看了很久之后，}明显没有欺骗，没有\OriginalTerm{急促}{\GreekText{συναρπαγή}}。经上说：\BibleQuote{颇理约友善地款待了他们}\BibleRef{宗 28:7}：一共二百七十六人。想想他的好客带来了多么大的收获：他不是出于不得已，也不是不情愿，而是看作一种收益，款待了他们三天；之后他得到了报偿，自然更加尊敬保禄，因为其他人也得了治愈。经上说：\BibleQuote{他也以许多荣誉尊敬了我们}\BibleRef{宗 28:10}：并不是他得到了工钱，绝对不可；但正如经上所写：\BibleQuote{工人配得他的食物。当我们离开时，他们给我们装上了所需之物。}\BibleRef{玛 10:10} 显然，他们这样接待了他们，也接受了宣讲的圣言：因为不能设想，在整整三个月里，他们会受到如此厚待，除非这些人深信不疑，并在此表现出（他们皈依的）果实：因此从这里我们可以看到一个有力的证据，证明相信的人很多。这足以使（保禄）在那些（他的同船旅客）眼中建立信誉。注意，在这次航行中，他们无处靠岸一个城市，而是（被抛）到一个岛上，度过了整个冬天（在那里或航行中）——那些人，我说的是他的同船旅客，这时正在接受信德的训练。 (a) \BibleQuote{过了三个月，我们乘一艘在岛上过冬的亚历山大船起航，船的旗号是卡斯托尔和波鲁克斯。}\BibleRef{宗 28:11} 大概这是画在船上的：他们是那么沉迷于他们的偶像。(d) \BibleQuote{南风一吹，我们第二天就到了\Place{波楚里}：在那里我们找到了弟兄们，被邀请与他们同住七天；然后我们就往罗马去。}\BibleRef{宗 28:13-14} (b) 注意他们暂停片刻，又继续赶路。(e) \BibleQuote{从那里，弟兄们听到我们的消息，就出来迎接我们，直到\Place{阿丕约市场}和\Place{三馆}}\BibleRef{宗 28:15}：不畏惧危险。(c) 因此保禄这时受到如此尊重，以至于他甚至可以独居一处：因为即使在这之前他们善待他，现在更会如此。(g) 经上说：\BibleQuote{他被允许独自居住，由一名兵士看守。}\BibleRef{宗 28:16} 这样就不可能在那边对他设什么阴谋——因为现在不可能有叛乱兴起。所以事实上他们不是在拘禁保禄，而是在保护他，免得不愉快的事情发生：因为现在在这样大的城市里，皇帝也在那里，保禄又上诉，不可能做出违反秩序的事。确实如此，总是通过那些似乎反对我们的事，一切都为我们转为有利。\InnerQuote{与那兵士}——因为他是保禄的看守。\BibleQuote{他召集了犹太人的首领}\BibleRef{宗 28:17}，对他们讲论，他们既离开反对，又被他责骂，却不敢说什么：因为他们不被允许按自己的意愿处理他的事。因为这是一件奇妙的事：不是通过那些似乎对我们安全的事，而是通过它们的反面，一切都为我们成为有益。为使你学到这一点——法郎命令将婴孩丢进河里。\BibleRef{出 1:22} 若非婴孩被丢弃，梅瑟就不会得救，不会在王宫长大。当他安全时，他并不受尊荣；当他被暴露时，他却受了尊荣。但天主这样做，是为显示祂丰富的资源和设计。那个犹太人威胁他说：\BibleQuote{你要杀我吗？}\BibleRef{出 1:14} 这对他也有益处。这是天主眷顾的安排，为使他在旷野中看到那个异象，为使合适的时间得以完成，使他在旷野中学习智慧，并在那里安全生活。在犹太人对他的一切阴谋中，同样的事发生：他就变得更加显赫。同样在亚郎的事上；他们起来反对他，因此使他更加显赫\EditorialNote{户 16-17}：为使他被祝圣成为无可置疑，为使他将来也因\OriginalTerm{铜版}{\GreekText{τὥν} \GreekText{πετάλων} \GreekText{τοὕ} \GreekText{χαλκοὕ}}而被钦佩。你当然知道这段历史，所以我略过叙述。如果你愿意，让我们从头再来看这些同样的例子。\par

\Person{加音}杀了他的弟弟，但这反倒有益于弟弟：因为请听圣经说：\BibleQuote{你弟弟的血由地上向我喊冤。}\BibleRef{创 4:10} 又在另一处说：\BibleQuote{那血说话比亚伯尔的血更好。}\BibleRef{希 12:24} 他使弟弟免于未来的不确定，增加了他的赏报：我们全都由此得知天主多么爱他。他究竟受了什么损害？丝毫没有，因为他更早得着了结局。请说，那些死得较慢的人得了什么好处？什么也没有：因为过好日子不在于年岁的多少，而在于善用生命。三位青年被投入火窑，借此他们变得更光荣；\Person{达尼尔}被投入狮坑，从此他更受称扬。（达 3 章和 6 章）你们看见，各种试炼即使在今生也带来极大益处，更何况来世呢？至于恶意，就如一个人拿着芦苇去与火搏斗：似乎是在打火，却使火更旺，而自己白白消耗。因为恶人的恶意成了德性的养料和显耀的机会：由于天主将不义转为善用，我们的品格便更加发光。同样，当魔鬼做成这类事时，它使忍耐者更显荣光。\newline{}   你或许会说：\Quote{但\Person{亚当}的情况并非如此，相反，他更蒙羞辱？} 不，天主恰恰在这件事上将（那恶者的）恶意转为善用：但若（\Person{亚当}）因此而更坏，那是他害了自己：因为别人加于我们的伤害能成为我们的大益，但自己加于自己的伤害却不能。事实上，正因为受别人伤害时我们忧伤，而受自己伤害时却不忧伤，所以天主表明：无辜受苦者得荣耀，而自害者受亏损，好叫我们勇敢忍受前者，却不至于后者。何况，那整件事是\Person{亚当}自己的作为。你为何听从了女人？\BibleRef{创 3:6} 她劝你违背（天主）时，你为何不拒绝她？你实在是你自己的原因。否则，若魔鬼是原因，那么凡受诱惑的都应灭亡；但若并非都灭亡，则（灭亡的）原因在于我们自己。\Quote{但是，}你或许会说：\Quote{按照这个道理，凡受诱惑的都应当成功。} 不：因为原因在于我们自己。\Quote{照此说来，有些人灭亡时魔鬼竟毫无关系。} 是的：事实上许多人灭亡了，魔鬼并未参与：因为魔鬼并不促成所有（我们的恶行）；不，很多也来自我们自身的懒惰：即便魔鬼在某处作为原因，那也是我们给了机会。请说，魔鬼为何在\Person{犹达斯}身上得逞？当\BibleQuote{撒殚进入了他}\BibleRef{若 13:27} 时，你或许会说。是的，但请听原因：是因为\BibleQuote{他是个贼，掌管钱囊，常取其中所存。}\BibleRef{若 12:6} 是他自己给魔鬼敞开了大门，所以并非魔鬼将开端放在我们里面，而是我们接受并邀请他。\Quote{但是，}你或许会说：\Quote{如果没有魔鬼，邪恶就不会这么大。} 不错，但那样我们的惩罚就无可宽恕了；然而如今，亲爱的，我们的惩罚更轻，假若我们自行作恶，刑罚将不可忍受。请说，如果\Person{亚当}没有受任何劝诱而犯了他所犯的罪，谁能把他从危险中救出？\Quote{但他就不会犯罪，}你或许会说？你凭什么这样说？因为那如此不坚定、如此迟钝、如此易于接受这种劝告的人，即使没有劝诱，更会变成（他所变成的那样）。什么魔鬼挑起了\Person{若瑟}的兄弟们嫉妒？如果我们警醒，魔鬼反而成为我们得荣耀的原因。看，\Person{约伯}陷入如此无助的痛苦，他损失了什么？\Quote{不要提这个例子，}你或许会说：\Quote{（\Person{约伯}没有损失，）但软弱的人会损失。} 是的，软弱的人会损失，即使没有魔鬼。\Quote{但是程度更大，}你或许会说：\Quote{当有魔鬼的力量与他一同作用时。} 不错，但当他因魔鬼的作用而犯罪时，他受罚更轻；因为不是所有罪的刑罚都相同。不要自欺：魔鬼并非我们受害的原因，只要我们警醒；相反，他所做的，是唤醒我们出离睡眠；他所做的，是使我们保持警觉。让我们暂且思考这些：假设没有野兽，没有反常的气候；没有疾病，没有疼痛，没有忧伤，也没有诸如此类的事，人将会变成什么？不是人，而是猪，沉溺于贪食醉酒，不受任何困扰。但如今，挂虑和焦虑是哲学的训练和纪律，是最好的锻炼方法。请说，假如一个人养在宫中，没有痛苦、没有挂虑、没有焦虑，也没有恼怒或失败的原因，凡他心思所想的，都让他去做，让他成功，所有人都服从他：（看看）这样的人岂不比任何野兽更无理性？但如今，我们的逆境和苦难有如磨刀石，磨砺我们。因此，穷人大多比富人更聪明，因为他们被许多波涛冲击。同样，身体若是懒惰不动，便多病丑陋；但若锻炼，劳苦，则更健美；我们发现灵魂也是如此。铁若不用，便生锈，若被锤炼，则明亮发光；灵魂也是如此，被不断运动。这些逆境正是保持灵魂运动的因素。技艺若不活动便衰亡；但灵魂只有不事事顺遂时才活动；它因逆境而活动。如果没有逆境，就没有什么能激动它：假如一切都已经现成地美好，技艺便没有锻炼的机会。同样，如果一切事物都容易理解，灵魂便没有发挥的机会；如果它到处被抱着走，就会变得丑陋。你们岂不见，我们劝告乳母不要习惯总是抱着孩子，免得养成习惯而变得无助？正是因此，在父母眼底下长大的孩子身体软弱，因为娇惯，过分姑息反而损害健康。适度的痛苦是好事；挂虑是好事；缺乏是好事；因为它们使我们坚强：它们的反面也是好事；但每种东西若过度就会毁灭我们；一种使人松弛，另一种（过度紧张）使人破裂。你们岂不见，\Person{基督}也正是这样训练祂的门徒？如果他们需要这些，我们更需要。既然我们需要，就让我们不要忧伤，反而在磨难中喜乐。因为这些都是药剂，针对我们的创伤，有的苦涩，有的温和；但任何一种单独使用都无益。因此，让我们为这一切感谢天主：因为祂不是随意让这些发生，而是为了我们灵魂的益处。所以，让我们显示感恩，向祂谢恩，赞美祂，勇敢忍受，思念这只是暂时的，心神向往将来的事，好使我们既轻看现今，又配得将来的美物，因祂独生子的恩宠和仁慈，圣子与圣父及圣神同享荣耀、权能、尊崇，从今世直到永远。阿们。\par

\SourceRecord{Translated by J. Walker, J. Sheppard and H. Browne, and revised by George B. Stevens.；Nicene and Post-Nicene Fathers, First Series；Vol. 11.；Edited by Philip Schaff.；Buffalo, NY: Christian Literature Publishing Co.,；1889.；<http://www.newadvent.org/fathers/210154.htm>.}

% structure: p0001 source=h1 target=section reason=page-title

\section{宗徒大事录讲道 第五十五篇}

\BibleRef{宗 28:17}\par

\BibleQuote{到了第三天，保禄召集犹太人的首领。他们来了，保禄对他们说：诸位仁人弟兄，我虽没有做什么反对人民或祖先规矩的事，却从耶路撒冷被囚禁交到罗马人手中。他们审问了我，愿意释放我，因为我身上没有致死的罪。但是犹太人反对，我被迫上告凯撒，并非有什么要控告我的民族。因此，我请你们来，要见你们，和你们谈话，因为我是为了以色列的盼望才被这条锁链捆锁的。}\par

他愿意为自己和别人开脱；为自己，免得他们控告他，并因此伤害自己；为别人，免得显得一切都是他们做的。因为当时很可能有一种传言，说他是被犹太人交出来的；这足以使他们惊慌。于是他针对这一点，为自己辩护。有人可能说：\Quote{那么，他们无缘无故把你交出来，这合理吗？}罗马总督们为我作证，他们愿意放我走。\Quote{那么他们为什么没有放你走？}\Quote{当犹太人反对时，}他说。请注意他是如何减轻他们对他的指控的。如果他愿意加重事态，他本可以利用这些来更严厉地对待他们。因此他说：\Quote{我被迫上告凯撒；}所以他的整个讲话都是宽恕性的。那么，你这样做是为了控告他们吗？不，他说：\Quote{并非有什么要控告我的民族，}而是为了躲避危险。因为\Quote{我是为了你们才被这锁链捆锁的。}他说，我对你们毫无敌意。于是他们也被他的话感动，也为自己的同族道歉：\Quote{他们对他说：我们没有从犹太地接到关于你的信，也没有弟兄中有人来报告或说过你什么坏话。}\BibleRef{宗 28:21}无论是通过信件还是通过人，他们都没有说你什么坏话。不过，我们愿意听你自己说说：\Quote{但我们愿意听你自己谈谈你的看法。}\BibleRef{宗 28:22}然后他先表明自己的态度。\Quote{因为关于这教派，我们知道它到处受人反对。}他们约定了日子，有许多人到他的住处来。他给他们讲解并证明天主的国，从梅瑟法律和先知书起，从早到晚劝他们信从耶稣。\Quote{有些相信了所说的话，有些却不信。}（第23、24节）他们没有说我们反对，而是说\Quote{它被人反对。}然后他没有立即回答，而是给了他们一天的时间。他们来了，他讲话，如经上所记：\Quote{从梅瑟法律和先知书起。有些人信了，有些人不信。}当他们彼此不合时，他们就离开了，保禄说了一句话：\Quote{圣神藉依撒意亚先知向我们的祖先说得真好，说：你到这百姓那里去，说：你们听了又听，却不明白；看了又看，却不察觉。因为这百姓的心迟钝了，耳朵发沉，眼睛闭着，免得眼睛看见，耳朵听见，心里明白，回转过来，我就医治他们。}\BibleRef{宗 28:25}但当他们离开时，彼此对立，于是他才责备他们，不是要责备那些不信的，而是要坚固那些信的人。\Quote{依撒意亚说得好，}他对他们说。因此，这奥秘赐给了外邦人。如果他们抗拒，那也不足为奇，这从一开始就预言了。然后他又激起他们对那些外邦人的嫉妒。\Quote{所以你们应当知道，天主的救恩已传给外邦人，他们将要听从。}他说了这些话后，犹太人便离开了，彼此激烈争辩。保禄在自己租来的房子里住了整整两年，凡来见他的人，他都接待，宣讲天主的国，并教导关于主耶稣基督的事，毫无阻碍，没有人禁止他。阿们。\BibleRef{宗 28:28}这表明他当时享有自由：在罗马毫无阻碍，而他曾在犹太地受阻；他在那里住了两年教导。此后呢？\par

(重述提纲。) (d) \Quote{他们审断了我，}他说，\Quote{在我身上找不出什么罪状} \BibleRef{宗 28:18}。那些本应解救他的人，反而\Quote{把他交在罗马人手里}。如此之过分，* * 因为他们无权定罪，却把他当作囚犯交了出来。\Quote{我并非有什么要控告我本族的人} \BibleRef{宗 28:19}，我才来的。请看言辞何等友善：\Quote{我本族}——他并不视他们为外人。他不说\Quote{我不控告}，而是说\Quote{我没有什么可控告的}，尽管他曾遭受他们如此多的恶待。但这些他一句也不提，也不使他的言辞带有攻击性；他也不显得是出于恩惠而宽恕他们。因为关键在于，要表明是那些人把他当作囚犯交给罗马人，而那些本应定他罪的人却没有这样做。\Quote{因此，}他说，\Quote{我请求见你们} \BibleRef{宗 28:20}：为使人无权控告我，说一些\OriginalTerm{所想到的}{\GreekText{τὰ} \GreekText{παριστάμενα}}话，以为我逃脱了他们的手而来此，不是为了给他人带来祸患，而是为了逃避祸患。\Quote{我不得已上诉凯撒。}请观察他们也对他说得更温和了。他们说：\Quote{我们请求你。}并想为那些（在耶路撒冷的）人开脱。而他们本应控告那些人，却为他们辩护；事实上，他们替那些人开脱的举动正成了对他们的控告。因为这本身就是证明，他们知道自己极其理亏。倘若他们有把握，他们无论如何都会这么做，好使他不能按自己的方式陈述事实，此外他们也不愿前来。而他们多次尝试，表明了* * \Quote{至于这教派，我们知道，}他们说，\Quote{是到处被反对的。}（第21、22节）不错，但（人们）也到处被说服（事实上，这里也是如此），\Quote{有些人被说服了，有些人不相信。他们定了日子，}等等\BibleRef{宗 28:23}。再看，他又一次不是藉着奇迹，而是藉着法律和先知使他们无言，我们总是发现他这样做。其实他本可行奇迹，但那样就不再是信德的事了。事实上，从法律和先知讲论本身就是一个伟大的记号。然后，为叫你不以为怪（他们不信），他引用了那预言说：\Quote{你们听了，却不明白，}较那时更甚；\Quote{你们看了，却不察觉} \BibleRef{宗 28:26}，较那时更甚。这话不是为前者说的，而是为不信的人说的。那么，预言说那些信了的，又怎样呢？\Quote{你去对这百姓说}（即对那不信的百姓）。他说这话并非为了侮辱他们，而是为了消除障碍。\Quote{所以你们要知道，}他说，\Quote{天主的救恩已传给外邦人。他们，}他说，\Quote{也要听。} \BibleRef{宗 28:28} 那么你为什么向我们讲论？你不知道这事吗？是的，但为使你们被说服，也为我可为我自己辩护，不给任何人把柄。（c）不信的人退去了。但请看，他们现在不设谋害他。因为在犹太地，他们有一种暴政。那么，为什么天主的眷顾定意要他往那里去，而主却曾说过：\Quote{你赶快离开耶路撒冷？} \BibleRef{宗 22:18} 为要既显明他们的恶毒，又使基督的预言应验——他们不愿听他说话；也为叫众人知道，他准备承受一切，而且这事为在犹太地的人是一种安慰：因为在那里（弟兄们）也正遭受许多严重的苦难。但若他在宣讲犹太教义时已如此受苦，那么若他宣讲基督荣耀的教义，他们怎能容忍他呢？当他\Quote{行洁净礼} \BibleRef{宗 21:26}时，他已令人无法忍受，那么他宣讲时又怎能被容忍呢？你们指控他什么？你们听见了什么？他没有说过这类话。他只是被看见，便激起了所有人反对他。那么他自然被分别出来归向外邦人，自然被远远差遣；在那里也注定要向外邦人讲论。他先招呼犹太人，然后向他们展示事实后，才来到外邦人。\BibleRef{宗 24:18} \Quote{圣神说得真好，}等等。但这句话\Quote{圣神说}并不稀奇：因为天使也被说成是说主所说的话；但祂（圣神）不是这样。当人说到天使所说的话时，人不会说\Quote{天使说得真好}，而是说\Quote{主说得真好}。\Quote{圣神说得真好：}等于说，你们不信的不是我。但天主从一开始就预知这事。\Quote{他放胆无阻地讲论} \BibleRef{宗 28:31}：因为放胆讲论可能，却受阻碍。他的胆量毫无阻拦；事实上他也未受阻地讲论。\Quote{讲论，}它说，\Quote{有关天主国的事：}注意，没有一件是可感觉的、属于现在的事。但在两年后他的事，我们怎么说？（作者）让听众渴求更多：外邦作者也是这样做的（在他们的作品中），因为知道一切会使读者乏味厌倦。或者他这样做，是因为他无权凭己知来展示。注意天主眷顾的安排：\Quote{我多次被阻延，不能到你们那里去……多年来我切切想望到你们那里去。} \BibleRef{罗 15:22} 但他用希望喂养他们。\Quote{我急于往西班牙去，}他说，\Quote{并希望在旅途中见到你们，且得你们送行，首先在某种程度上与你们相交。} \BibleRef{罗 15:24} 关于这件事，他说：\Quote{我带着基督福音的圆满祝福而来，同你们一同安息} \BibleRef{罗 15:29}；又\Quote{我现在往耶路撒冷去，供给圣徒} \BibleRef{罗 15:25}：这正是他在这里所说的：\Quote{我来是为我本族行施舍。} \BibleRef{宗 24:17} 你注意到他并非预见一切——那神圣而尊贵的头脑，那比诸天更高的人，那灵魂能把握一切的人，那居首位的人——保禄？那人的名字，为认识他的人，足以激发灵魂，使人警醒，驱散一切沉睡！罗马迎接了他，他被绑着，从海上来，从船难中得救——并从谬误的船难中得救。如同一位打过海战并得胜的皇帝，他进入了那最帝国的城邑。他离冠冕更近了。罗马迎接了他被绑，却看见他加冕并被宣布为得胜者。在那里他曾说：\Quote{我要同你们一同安息；}但这却是一次新旅程的开始，他添上胜利于胜利，一个不可战胜的人。格林多留了他两年，亚细亚三年，这城这次留了他两年；第二次他又进入这城，那时他也被完成了。这样他当时逃脱了，并充满了整个世界，如此结束了生命。你为什么想了解这两年后的事？那些也是同样的：锁链、折磨、争斗、监禁、埋伏、诬告、死亡，日复一日。你只看见一小部分吗？无论你看见多少，其余都是如此。如同天空，你若看见它的一部分，无论你往哪里去，你都会看见它如此；如同太阳，你虽只看见它的一部分光芒，也可推测其余：保禄也是如此。他的宗徒大事录你只见了一部分；全部都是如此，充满危险。他是一个天空，内有正义的太阳，不是像我们所见的太阳；所以那人比天空本身更好。你以为这是一件小事吗——当你说\Quote{宗徒，}立刻人人想到他（是这样的）；当你说\Quote{洗者，}立刻人人想到若翰？拿什么与他的言语相比？与海，或与海洋？但没有什么能与它们相比。\par

比这（海）更丰沛的是（他的）江河；更纯洁、更深邃；以致人若称保禄的心既是海又是天，为纯洁而言是海，为深邃而言是天，也不为过。他是海，其航行者不是那些从一城航行到另一城的人，而是那些从地到天的人：若有人航行在这海上，必会旅途顺利。在这海上，不是风，而是代替风的圣神和神圣的圣神吹送航行其上的灵魂：这里没有波浪，没有礁石，没有怪物：一切都是平静的。这是一个比港湾更平静安稳的海，没有苦涩的咸水，而是一道清澈的泉源，比 * * 更甘甜，比太阳更明亮更清澈：这海没有宝石，没有像我们这里的紫色染料，却有远胜这些的宝藏。凡愿意下到这海中的人，不需要各样东西，不需要油，而需要大量的\OriginalTerm{慈爱}{\GreekText{φιλανθρωπίας}}：他将在其中找到天国里的一切美物。他甚至能成为君王，拥有整个世界，得享极大的尊荣；航行在这海上的人绝不会遭遇海难，反而会通晓一切。但正如不熟悉这（我们可见的）海的人（尝试潜水时）会被淹死，在那另一个海上也是如此：这正是异端者的情况，当他们尝试超出自己能力之事时。因此应当知道其深度，否则就不要冒险。若我们要航行在这海上，让我们整装待发。\Quote{我不能，}他说，\Quote{对你们说话如同对属神的人，却如同对属血肉的人。}\BibleRef{格前 3:1}让没有忍耐的人不要航行在这海上。让我们为自己预备船只，即热忱、恳切、祈祷，好能平安渡过这海。因为这确实是活水。就如一个人若得了火舌，那熟悉保禄的人也得了这样的火舌；就如一个人若有锋利的剑，这样的人也变得不可战胜。要理解保禄的话，也需要纯洁的生活。因此他也说：\Quote{你们成了需要奶的人，因为你们听觉迟钝。}\BibleRef{希 5:11}因为确实有听觉的软弱。正如软弱的胃不能消化健康的食物，同样，一个变得肿胀、发热、无力、松弛的灵魂，不能接受圣神的话语。听门徒们说：\Quote{这话生硬，谁能听得下去呢？}\BibleRef{若 6:60}但如果灵魂强壮健康，一切就最容易，一切就轻松：它变得更崇高、更轻快，更能翱翔并提升到高处。因此，知道这些后，让我们使自己的灵魂进入健康状态：让我们效法保禄，仿效那高贵、那金刚石般的灵魂：好使我们在他生命的脚步中前进，能够航行过今世的海，抵达那没有波浪的港湾，并获得那为爱祂的人所预备的美物，藉着我们的主耶稣基督的恩宠和慈爱，愿光荣、权能、尊崇归于祂，并与圣父及圣神，从今世直到永远，世世无尽。阿们。\par

\SourceRecord{Translated by J. Walker, J. Sheppard and H. Browne, and revised by George B. Stevens.；Nicene and Post-Nicene Fathers, First Series；Vol. 11.；Edited by Philip Schaff.；Buffalo, NY: Christian Literature Publishing Co.,；1889.；<http://www.newadvent.org/fathers/210155.htm>.}
